% Options for packages loaded elsewhere
\PassOptionsToPackage{unicode}{hyperref}
\PassOptionsToPackage{hyphens}{url}
\PassOptionsToPackage{dvipsnames,svgnames,x11names}{xcolor}
\documentclass[
  letterpaper,
  11pt]{article}
\usepackage{xcolor}
\usepackage[margin=1.05in]{geometry}
\usepackage{amsmath,amssymb}
\setcounter{secnumdepth}{5}
\usepackage{iftex}
\ifPDFTeX
  \usepackage[T1]{fontenc}
  \usepackage[utf8]{inputenc}
  \usepackage{textcomp} % provide euro and other symbols
\else % if luatex or xetex
  \usepackage{unicode-math} % this also loads fontspec
  \defaultfontfeatures{Scale=MatchLowercase}
  \defaultfontfeatures[\rmfamily]{Ligatures=TeX,Scale=1}
\fi
\usepackage{lmodern}
\ifPDFTeX\else
  % xetex/luatex font selection
  \ifLuaTeX
    \usepackage{luaotfload}
    \directlua{luaotfload.add_fallback("mainfontfallback",{
      "DejaVuSans:mode=node"
    })}
  \fi
  \setmainfont[,RawFeature={fallback=mainfontfallback}]{Latin Modern Roman}
  \setmonofont[Scale=0.85]{DejaVu Sans Mono}
  \setmathfont[]{Latin Modern Math}
\fi
% Use upquote if available, for straight quotes in verbatim environments
\IfFileExists{upquote.sty}{\usepackage{upquote}}{}
\IfFileExists{microtype.sty}{% use microtype if available
  \usepackage[]{microtype}
  \UseMicrotypeSet[protrusion]{basicmath} % disable protrusion for tt fonts
}{}
\makeatletter
\@ifundefined{KOMAClassName}{% if non-KOMA class
  \IfFileExists{parskip.sty}{%
    \usepackage{parskip}
  }{% else
    \setlength{\parindent}{0pt}
    \setlength{\parskip}{6pt plus 2pt minus 1pt}}
}{% if KOMA class
  \KOMAoptions{parskip=half}}
\makeatother
\usepackage{longtable,booktabs,array}
\usepackage{caption}
\captionsetup[table]{skip=6pt}
\newcounter{none} % for unnumbered tables
\usepackage{calc} % for calculating minipage widths
% Correct order of tables after \paragraph or \subparagraph
\usepackage{etoolbox}
\makeatletter
\patchcmd\longtable{\par}{\if@noskipsec\mbox{}\fi\par}{}{}
\makeatother
\setlength{\emergencystretch}{3em} % prevent overfull lines
\providecommand{\tightlist}{%
  \setlength{\itemsep}{0pt}\setlength{\parskip}{0pt}}
\usepackage{mathtools}
\usepackage{amssymb}
\providecommand{\lt}{<}
\providecommand{\gt}{>}
\providecommand{\Bbb}{\mathbb}
\providecommand{\bold}{\mathbf}
\providecommand{\llbracket}{\lBrack}
\providecommand{\rrbracket}{\rBrack}
\usepackage{fancyhdr}
\pagestyle{fancy}\fancyhf{}\fancyhead[L]{\small Four-dimensional quantum Yang--Mills theory: the proof, outlined}\fancyhead[R]{\small\thepage}
\setlength{\headheight}{14pt}
\setlength{\emergencystretch}{3em}
\sloppy
\AtBeginDocument{\renewcommand{\setminus}{\mathbin{\backslash}}}
\usepackage{bookmark}
\IfFileExists{xurl.sty}{\usepackage{xurl}}{} % add URL line breaks if available
\urlstyle{same}
% fallback for those not using the hyperref driver hyperxmp:
\makeatletter
\@ifundefined{xmpquote}{\newcommand{\xmpquote}[1]{#1}}{}
\makeatother
\hypersetup{
  pdftitle={Four-dimensional quantum Yang--Mills theory: the proof{,} outlined},
  colorlinks=true,
  linkcolor={[RGB]{0,0,114}},
  filecolor={Maroon},
  citecolor={Blue},
  urlcolor={[RGB]{0,0,114}},
  pdfcreator={LaTeX via pandoc}}

\title{Four-dimensional quantum Yang--Mills theory: the proof, outlined}
\usepackage{etoolbox}
\makeatletter
\providecommand{\subtitle}[1]{% add subtitle to \maketitle
  \apptocmd{\@title}{\par {\large #1 \par}}{}{}
}
\makeatother
\subtitle{A natural-language account of the Lean proof
(OpenAI-YM-20261028, commit 3456e5e)}
\author{}
\date{11 November 2026}

\begin{document}
\maketitle

{
\setcounter{tocdepth}{2}
\tableofcontents
}
\section{The proof in one reading}\label{the-proof-in-one-reading}

This is a substantial overview, about 10,000 words. The contents links
locate individual topics; the
\href{https://github.com/scottnarmstrong/yang-mills-informalization/blob/main/statement.md}{exact
statement} fixes the scope of every claim.

\subsubsection{The theorem}\label{the-theorem}

Let \(G\subset U(m)\) be a closed subgroup whose Lie algebra
\(\mathfrak g\) is simple and nonabelian, for instance \(SU(N)\) with
\(N\ge2\); \(G\) need not be connected. The main theorem is
\texttt{YangMills.jointExistence}; its precise statement is in
\href{https://github.com/scottnarmstrong/yang-mills-informalization/blob/main/statement.md}{statement.md}.
It asserts that there is one family of Wilson lattice gauge theories for
\(G\) on four-dimensional tori, consisting of:

\begin{itemize}
\tightlist
\item
  spacings \(a_N=1/D_N\downarrow0\) and physical side lengths
  \(L_p\uparrow\infty\);
\item
  bare couplings \(\beta_{N,p}\) with
  \(\beta_{N,p}/\log(1/a_N)\to b_\ast>0\) uniformly in \(p\);
\item
  insertions \(I_{N,p}(v,y)\) indexed by tokens \(v\), each a continuous
  gauge-invariant function of the links in a stencil of fixed size in
  lattice units around the site \(y\).
\end{itemize}

The family has the following property. For every strictly increasing
sequence \(p(k)\) of volume indices there are a further subsequence
\(p(q(k))\) and a continuum hierarchy satisfying (1)--(5) below.

\textbf{(1) Lattice origin.} Let
\(\mathrm{Conf}_n\subset(\mathbb R^4)^n\) be the open set of
configurations of distinct points. For every \(n\), every list of tokens
\(v\) and every \(f\in C_c^\infty(\mathrm{Conf}_n)\),

\[\lim_{k\to\infty}\,\lim_{N\to\infty}\,a_N^{4n}\sum_{y}f(\hat y)\,\Big\langle\prod_{i=1}^nI_{N,p(q(k))}(v_i,y_i)\Big\rangle_{N,p(q(k))}=\int f(x)\,S^{\mathrm{mix}}(\pi v;x)\,dx .\]

Here the sum runs over \(n\)-tuples of sites, \(\hat y\) are their
physical positions, and \(\langle\cdot\rangle_{N,p}\) is the Wilson
measure. - The inner limit runs over the whole sequence of spacings, at
every volume whose fundamental box contains the support of \(f\); this
holds for all large \(k\). - The test functions are joint functions of
the \(n\) points and vanish near all coincidences. - One subsequence
serves all \(n\), \(v\) and \(f\).

\textbf{(2) Phases.} On \(\mathrm{Conf}_n\),
\(S^{\mathrm{mix}}=\int_XS^\xi\,d\nu(\xi)\). Here \(X\) is a compact
metric space, \(\nu\) is a probability measure of full support, and
\(\xi\mapsto S^\xi(a;x)\) is continuous. All correlations \(S^s\),
\(s\in X\sqcup\lbrace\mathrm{mix}\rbrace\), are smooth and permutation
symmetric on \(\mathrm{Conf}_n\).

\textbf{(3) Reconstruction.} A selected phase \(\xi_\ast\) carries a
Wightman theory: - a strongly continuous unitary representation of the
proper orthochronous Poincaré group; - spectrum in the closed forward
cone; - a vacuum that is unique up to scalars among
translation-invariant vectors; - fields that are operator-valued
tempered distributions on the invariant domain spanned by field
monomials applied to the vacuum; - local commutativity and faithful
field labels.

Its Wightman functions are boundary values, on product test functions,
of the continuation of \(S^{\xi_\ast}\) that is holomorphic in the
ordered complex times.

\textbf{(4) Operator product expansion.} One system of coefficients,
with finitely many local fields of each degree, satisfies finite-order
OPE bounds in the mixed state and in every phase.

\textbf{(5) Nonfreeness.} Let \(e_4\) be the time direction. A field
\(\Theta\) of degree 4 with vanishing vacuum expectation satisfies

\[\lim_{r\downarrow0}\,r^8\log^2(1/r)\,G_\Theta(r)=\frac{3\dim_{\mathbb R}\mathfrak g}{\pi^4},\qquad G_\Theta(r)=S^{\xi_\ast}\big(\Theta(-\tfrac r2e_4)\,\Theta(\tfrac r2e_4)\big).\]

Moreover, \(G_\Theta\) agrees on no interval \((0,\delta)\) with a
finite Wick profile, that is, a finite sum of products of free
mass-shell kernels
\(\int E^q h(\mathbf p)(2E)^{-1}e^{-rE}\,d^3\mathbf p\).

The lattice data are fixed before the volume sequence. The refinement
and all continuum objects may depend on that sequence.

\textbf{The stronger statement.} The theorem is a projection of a longer
statement, \texttt{YangMills.constructed}. That statement holds for
every value \(b>0\) of the penalty parameter of the block-spin kernel;
the public theorem takes \(b=1\). The construction supplies the
following facts; the slope value is established in the proof of
regularity, rather than stated as a separate conjunct of
\texttt{YangMills.constructed}: - The slope is
\(b_\ast=\kappa_G=11C_{\mathfrak g}/(24\pi^2)\). - A token is a pair
\((d,c)\), with \(c\) in the space \(V_d\) of gauge-invariant
polynomials of weight at most \(d\) in the curvature and its covariant
derivatives (curvature has weight 2, and each derivative adds 1). - The
field spaces are \(F_d=V_d/N_d\), where \(N_d\) consists of the labels
whose mixed correlations with all separated spectators vanish. The label
space of the relativistic theory is the directed union of the \(F_d\),
so there is no finite field multiplet. - \(\Theta\) is the centred trace
of a constructed conserved stress tensor. - Every phase carries an
irreducible Wightman theory, with low-energy bounds that are uniform in
the phase. - The correlations extend to distributions across coincident
points; uniqueness is not asserted. - After finite-dimensional changes
of frame, the rescaled correlations converge at short distance to those
of \(\dim\mathfrak g\) free Maxwell fields. - Nested and integrated
forms of the OPE hold, as do further nonfreeness tests.

This overview describes the proof of the stronger statement.

\textbf{What is not claimed.} - No mass gap is proved; the available
spectral bounds in every phase are polynomial. - The infinite-volume
limit is not shown to be unique or independent of the subsequence. - The
phases are not shown to agree, and \(X\) is not shown to be a point. -
The mixed state is not shown to cluster. - The limit is not shown to be
independent of the lattice data. - The two limits are taken in a fixed
order. No diagonal limit, exchange of limits or rate of convergence is
asserted. - No phase is shown to be a limit of lattice states; the
lattice limit is the mixture. - The public theorem asserts nothing at
coincident points and no Euclidean axioms. The construction proves those
axioms internally. - There is no continuum measure on connections, no
Wilson-loop observable or area law, and no particle or scattering
theory. - Only finite Wick profiles are excluded. Generalized free
fields with a continuous mass distribution, infinite sums and general
composites of Gaussian fields are not excluded. - The OPE is a
finite-order bound and is not compared with perturbation theory.

\textbf{Normalizations.} On \(\mathfrak g\), put
\(\langle X,Y\rangle_{\mathfrak g}=-\tfrac1m\mathrm{Re}\,\mathrm{tr}(XY)\),
and define \(C_{\mathfrak g}>0\) by
\(\mathrm{tr}(\mathrm{ad}X\,\mathrm{ad}Y)=-C_{\mathfrak g}\langle X,Y\rangle_{\mathfrak g}\).
Two facts fix the coupling convention: -
\(1-\tfrac1m\mathrm{Re}\,\mathrm{tr}\,e^X=\tfrac12|X|_{\mathfrak g}^2+O(|X|^4)\);
- each unoriented plaquette is counted once in the Wilson action.

Hence \(\beta=1/g_0^2\) for the continuum action
\(\frac1{2g_0^2}\int\sum_{\mu<\nu}|F_{\mu\nu}|^2_{\mathfrak g}\). The
constructed couplings satisfy, uniformly in \(p\),

\[\frac{\beta_{N,p}}{\log(1/a_N)}\longrightarrow\kappa_G=\frac{11C_{\mathfrak g}}{24\pi^2}=2b_0,\qquad b_0=\frac{11C_{\mathfrak g}}{48\pi^2}.\]

For \(SU(N)\) in the defining representation, \(C_{\mathfrak g}=2N^2\)
and \(\kappa_G=11N^2/(12\pi^2)\). In that case \(\beta\) is the usual
Wilson coupling \(2N/g^2\), and the law above is the one-loop running
\(1/g^2\simeq(11N/24\pi^2)\log(1/a)\) of Gross--Wilczek and Politzer.

\subsubsection{The strategy}\label{the-strategy}

\textbf{Order of limits.} One lattice family is fixed first.
Correlations of renormalized insertions are taken to the limit
\(a_N\to0\) at fixed volume, along the whole sequence of spacings. The
volume then tends to infinity along a subsequence of the prescribed one.
The renormalization group (RG) is used only for the first limit. Its
levels run from the Wilson lattice (level \(N\)) to the torus of side
\(L_p\) at unit spacing (level \(0\)). The bare coupling is tuned so
that the effective coupling at level \(0\) is one fixed number
\(\bar\beta\), the same for all \(N\) and \(p\). No RG step is taken
above the unit scale; infinite volume is reached by compactness, from
local bounds that are uniform in the volume.

\textbf{Ultraviolet construction (Parts I--III).} At each level the
density is written as two pieces: - an explicit gauge-invariant action,
whose normalized curvature-squared coefficient is the running coupling
\(\beta_{N,p,m}\); - polymer activities: regular ones, controlled where
the curvature is small, and forced ones, which carry the large-field
regions.

One block-spin step integrates the fine links against an exactly
normalized, gauge-covariant probability kernel: - In the small-field
region this is a Gaussian integral around a stationary background field,
organized by a cluster expansion. - Large-field regions are kept and
paid for by a lower bound on the quadratic energy of the step. - The
marginal term is absorbed into the coupling. The coupling changes by the
one-loop amount \(\kappa_G\log L_m\) up to an error, and block ratios
that grow slowly along the levels make the accumulated error
\(o(\log(1/a_m))\).

The same steps transport inserted observables. This second layer
produces what later parts use: - chessboard bounds on renormalized
insertions and on their difference quotients, uniform in the volume; -
convergence at fixed volume, obtained from fixed-level convergence and
uniform observable tails. The construction also supplies analytic
estimates in a tuning parameter.

\textbf{Euclidean limit (Part IV).} - Chessboard bounds make the
correlations locally bounded, uniformly in the volume. - A
transfer-matrix representation and complex interpolation make them
equicontinuous, and Arzelà--Ascoli gives the subsequence. - Difference
quotients of insertions are close to insertions of higher degree. This
yields derivative maps on the field spaces and kernels that are smooth
on \(\mathrm{Conf}_n\). - The limit inherits translations, the signed
permutations of the axes and reflection positivity. It does not inherit
continuous rotations or clustering.

\textbf{Symmetry (Part VI).} - Lattice integration by parts gives shell
identities for a conserved tensor. Its normalization is fixed by
short-distance comparison with free Maxwell theory. - Reflection
positivity controls the flux at infinity, and analyticity extends the
Ward identities to all separated configurations. - The infinitesimal
rotations integrate to an \(O(4)\) action, again by comparison with free
Maxwell theory.

\textbf{Phases and reconstruction (Part V).} - The zero-energy subspace
of the Osterwalder--Schrader (OS) space of the mixed state carries a
commutative \(C^\ast\)-algebra; its proof of commutativity uses the
\(O(4)\) action. The characters of this algebra are the phases. - Each
phase is reconstructed as a Wightman theory, from bounds on operators
inserted in thin time slabs. Lorentz covariance comes from the Ward
identities. - A low-energy spectral bound, uniform in the phase and
derived from locality, upgrades vacuum uniqueness from almost every
phase to every phase.

\textbf{Identification (Part VII).} - Along sequences of scales tending
to zero, the theory converges, after changes of frame, to free Maxwell
theory. - The one-loop coefficient produces an inverse-logarithmic law
for the centered stress-tensor trace two-point function, and that law
excludes the specified finite Wick profiles. - The RG also yields the
OPE.

The logical order is not the order of the parts. Part VI uses the OS
spaces of the mixed state from Part V and the free comparison from Part
VII. Part V uses the \(O(4)\) action from Part VI. Part VIII assembles
everything for one family and one limit table.

\subsubsection{Part I. The lattice
theory}\label{part-i.-the-lattice-theory}

\emph{Unit:
\href{https://github.com/scottnarmstrong/yang-mills-informalization/blob/main/outline/units/I-1.md}{I.1
--- Lattice, symmetries, gauge groups and the Wilson measure}.}

On \(\Lambda_\ell=(\mathbb Z/\ell\mathbb Z)^4\), a configuration assigns
\(U(y,\mu)\in G\) to each positively oriented link. With \(h\) the
normalized Haar measure, the Wilson measure is

\[d\mu_{\ell,\beta}(U)=Z^{-1}\exp\Big[-\beta\sum_y\sum_{\mu<\nu}\big(1-\tfrac1m\mathrm{Re}\,\mathrm{tr}\,U_{y;\mu\nu}\big)\Big]\prod_{(y,\mu)}dh\big(U(y,\mu)\big).\]

The energy lies in \([0,12\ell^4]\), so this is a probability measure
for every real \(\beta\). It is invariant under gauge transformations,
translations and the signed permutations of the axes.

\textbf{Lattice schedule.} - Block ratios are odd integers \(L_m\ge3\),
with \(D_{m+1}=D_mL_m\) and \(a_m=1/D_m\). - The level-\(m\) torus at
volume \(p\) has \(L_pD_m\) sites per side, with \(4\mid L_p\). - Odd
ratios make the coarse lattice compatible with the signed permutations.
- Supports are measured by cover numbers by balls of radius \(B\), which
obey a coarsening inequality uniform in the volume.

\textbf{Group facts.} On \(G\) the unit proves: - a real-analytic
logarithm near \(1\); - an equivariant nearest-point retraction onto
\(G\), whose derivative is the orthogonal projection onto
\(\mathfrak g\); - a curvature map \(F(u)\in\mathfrak g\) with
\(|F(u)|\le\lVert u-1\rVert/\sqrt m\) and
\(F(vuv^{-1})=\mathrm{Ad}(v)F(u)\).

These are qualitative statements about single group elements;
quantitative control comes in Part II.

\textbf{Role of simplicity.} Simplicity enters through one fact: every
\(\mathrm{Ad}\)-invariant symmetric bilinear form on \(\mathfrak g\) is
a multiple of \(\langle\cdot,\cdot\rangle_{\mathfrak g}\). Together with
lattice symmetry, this leaves a single marginal coupling. The same fact
is used in the one-loop computation and in the uniqueness of the stress
tensor.

\textbf{Reflection positivity.} The Wilson measure is reflection
positive for every \(\beta\ge0\), for reflections through planes of
sites and through planes between sites (the argument of Osterwalder and
Seiler). This is proved in Part IV for the original lattice measure, not
for an effective action, and it is used only through the chessboard
inequality.

\subsubsection{Part II. Fields, insertions and polymer
activities}\label{part-ii.-fields-insertions-and-polymer-activities}

\emph{Units:
\href{https://github.com/scottnarmstrong/yang-mills-informalization/blob/main/outline/units/II-1.md}{II.1
--- Field labels and their lattice representatives};
\href{https://github.com/scottnarmstrong/yang-mills-informalization/blob/main/outline/units/II-2.md}{II.2
--- Polymer activities, norms, the Wilson family, the construction
record, reflection-positive cells}.}

\textbf{Labels.} In radial gauge, a connection near a point is described
by its Taylor coefficients, and the coefficient of order \(j\) has
weight \(j+2\). \(V_d\) is the space of \(\mathrm{Ad}(G)\)-invariant
complex polynomials of weight at most \(d\) in these coefficients.
Equivalently, it consists of the gauge-invariant polynomials of weight
at most \(d\) in the curvature and its covariant derivatives. It is
finite-dimensional, and \(1\in V_d\subset V_{d+1}\).

\textbf{Lattice representatives.} The basic lattice variables are
``slots'': plaquette curvatures transported to a base point along walks.
- For each \(d\), finitely many slots have linearizations with no common
kernel on the jets of weight at most \(d\). - The slot map then has an
analytic \(\mathrm{Ad}\)-equivariant left inverse near \(0\); it is
chosen non-constructively. - Compose a basis of \(V_d\) with the
order-\(d\) Taylor polynomial of this inverse, and average over the
signed permutations. For each label \(c\) this gives a gauge-invariant
polynomial \(\Phi_c\) in the link entries of a fixed box stencil.

Two properties follow: - On links obtained by parallel transport of
polynomial radial-gauge potentials at spacing \(h\), \(\Phi_c\)
reproduces the classical label to order \(h^d\). - Suppose the
plaquettes near the stencil have chord at most \(t\). Then a component
with \(f\ge1\) curvature factors, and its derivatives, are \(O(t^f)\).
Conversely, a function with such bounds determines a label with
uniformly bounded coefficients.

\textbf{Renormalized insertions.} The insertion of \(v\in V_d\) at a
site \(o\) is

\[O_{N,p}(v,o)=a_N^{-d}\,\Phi^{(d)}_{S_{1/\theta_N}\tau^{-1}v}(o).\]

\begin{itemize}
\tightlist
\item
  \(S_t\) multiplies the coefficient of a component with \(f\) curvature
  factors by \(t^f\).
\item
  \(\theta_N\) is the small-field threshold of the finest level.
\item
  \(\tau\) is a linear automorphism of \(V_d\). It commutes with
  conjugation and with the signed permutations, and it fixes the unit.
\end{itemize}

The RG defines \(\tau\). The running label coordinates of an insertion
obey a transfer relation from level \(N\) down to level \(m\), and
\(\tau\) is their value at a reference level \(r(d)\) fixed in advance.
So each insertion is normalized to have running coordinates exactly
\(v\) at level \(r(d)\). All labels of weight at most \(d\), including
the unit, may mix; no triangular structure is imposed. The large-field
part of the inserted density goes into reserved forced polymers, so the
construction never divides by the vacuum density.

\textbf{Equation relations.} Haar invariance gives the lattice
Schwinger--Dyson identity
\(\langle\langle A,\nabla_eS_W\rangle\,S\rangle_\beta=0\) for \(S\)
independent of the link \(e\). Products of such forces with slot
polynomials span the linearized free-Maxwell equation ideal, up to terms
of higher weight. They give lattice observables whose correlations with
distant spectators vanish exactly. Two facts are established later,
through the comparison with free theory: these relations become null in
the limit, and they exhaust \(N_d\).

\textbf{Polymers, budgets and marks.} - A polymer activity has a finite
support of sites. - Regular activities are anchored at a site and
measured by suprema of derivatives over small-field configurations. -
Forced activities carry the large-field regions, together with an
inventory of flagged sites. - Activities are controlled by budgets
summed through each site,
\(\sum_{X\ni x}e^{a\,\mathrm{cov}_B(X)}\lVert K_X\rVert\le\text{budget}\).
- Insertions carry marks: commuting variables \(\lambda\) with
\(\lambda^2=0\), one for each label or site that may occur. Exponentials
of marked quantities are finite sums, and logarithms are finite
recursions in the number of marks. - The weighted norm
\(\lVert a\rVert_w=\sum_x\prod_{\lambda\in x}w_\lambda\,|a(x)|\) is
submultiplicative, so bounds do not grow with the number of marks. - A
tree-graph bound without factorials controls connected clusters of
overlapping supports.

\subsubsection{Part III. The ultraviolet
construction}\label{part-iii.-the-ultraviolet-construction}

\paragraph{The block-spin step and the background
field}\label{the-block-spin-step-and-the-background-field}

\emph{Units:
\href{https://github.com/scottnarmstrong/yang-mills-informalization/blob/main/outline/units/III-1.md}{III.1
--- The block-spin step and the symmetric gauge fixing};
\href{https://github.com/scottnarmstrong/yang-mills-informalization/blob/main/outline/units/III-6.md}{III.6
--- The stationary branch: background field of the effective action};
\href{https://github.com/scottnarmstrong/yang-mills-informalization/blob/main/outline/units/III-3.md}{III.3
--- Free lattice kernels and covariances}.}

At level \(m\), relative to product Haar measure, the density is

\[\rho_m(U)=\exp\big(v|\Lambda_m|-\beta\,\mathcal K(U)-\mathcal L(U)-\mathcal E(U)\big)\,\mathcal G(U).\]

\begin{itemize}
\tightlist
\item
  \(\mathcal K\) is built from lattice paths: a plaquette term, smeared
  curvature squares, and cubic and quartic polynomials in transported
  curvatures. It is averaged over the signed permutations and normalized
  so that its quadratic symbol at zero momentum is the identity; this
  fixes the meaning of \(\beta\).
\item
  \(\mathcal L\) is a loop term independent of \(\beta\).
\item
  \(\mathcal E\) is the sum of the regular activities times small-field
  indicators.
\item
  \(\mathcal G\) is the gas of forced activities.
\item
  At the finest level, \(\rho_N=e^{-\beta_{N,p,N}S_W}\).
\end{itemize}

\textbf{The kernel.} One step is
\(\rho_m(V)=\int k_\beta(V,U)\,\rho_{m+1}(U)\,dU\). Given the fine
configuration \(U\), a finite tag \(t\) is drawn with probability
\(\lambda(U,t)\). Each coarse link is then drawn independently, with
Haar density

\[Z_\beta^{-1}\exp\Big(-\beta\kappa\big(1-\tfrac1m\mathrm{Re}\,\mathrm{tr}\,V_bQ_b(U,t)^{-1}\big)\Big),\]

where \(\kappa>0\) is the penalty. The predictor \(Q_b\) is built in two
steps: 1. On a patch of radius \((J+3)L\) around each coarse site, a
gauge frame is chosen. - If all route holonomies on the patch are close
to \(1\), the frame is the unique minimizer of an averaged holonomy
discrepancy. - Otherwise it is a uniformly drawn eligible axial gauge;
such gauges exist for every configuration. 2. Fine transports of length
\(L\), read in these frames, are combined by a nonlinear (minimizing)
average with a \(J\)-fold convolution weight, with a random fallback.

Properties of the kernel: - Haar invariance gives
\(\int k_\beta(V,U)\,dV=1\) exactly. - The kernel is exactly covariant
under fine gauge transformations and the lattice symmetries. - Hence
\(\int\rho_m\) equals the Wilson partition function at every level, and
the step is defined without any small-field assumption.

\textbf{Background field.} For the integration, the links of a spanning
forest are set to \(1\), and each remaining link is \(\exp(a)\) times a
fixed lift of the coarse field \(x\). The background is the critical
branch \(a=\rho(x)\) of the classical Taylor model of the action
(quadratic, cubic and quartic terms) plus the blocking penalty, at fixed
\(x\). The steps are: - The Hessian at the flat field is a positive
``bed'' form plus the penalty. - Each non-tree link is bounded by the
curvature in its block plus the response of the predictor. With this,
the Hessian is uniformly positive in the vertical directions, with a
constant independent of the volume and the forest. - The analytic
implicit function theorem then gives \(\rho\) as a locally unique
analytic germ with explicit linear part. Statements with a uniform
radius and localized jets give control independent of the volume. - The
branch is not shown to be a global minimizer, and this is not used. -
Changing the forest multiplies the Hessian determinant by the square of
a Jacobian and the Haar measure by the Jacobian itself. So the critical
value and the one-loop term do not depend on the forest.

\textbf{Free kernels.} The free covariances are written in Fourier space
through the actual averaging filters, which are entire and
multiplicative under composition of blockings. Aliased momenta carry
extra powers of \(1/L\). - All constants (decay rates, strip widths,
lower bounds) are fixed before the block ratio \(L\). - The quadratic
forms decay exponentially and are bounded below uniformly in the volume.
- An exact identity on each pair of tori splits a quadratic form into a
coarse part and a fluctuation part (a Schur complement).

\paragraph{Gaussian integration and the cluster
expansion}\label{gaussian-integration-and-the-cluster-expansion}

\emph{Units:
\href{https://github.com/scottnarmstrong/yang-mills-informalization/blob/main/outline/units/III-4.md}{III.4
--- Gaussian integration: change of Gaussian sector, determinants,
window law};
\href{https://github.com/scottnarmstrong/yang-mills-informalization/blob/main/outline/units/III-5.md}{III.5
--- Cluster expansion of the polymer gas};
\href{https://github.com/scottnarmstrong/yang-mills-informalization/blob/main/outline/units/III-8.md}{III.8
--- Connected activities and backward dependence cones}.}

\textbf{Change of variables.} Active links are written in exponential
coordinates around the background. Inactive links stay in Haar measure.
- The change of variables to a Gaussian integral with covariance \(C\)
and mean \(m(Y)\) is an exact identity. It holds for any
positive-definite \(C\), keeps all Jacobians, and uses no smallness
assumption. - The covariance actually used is the inverse of the
gauge-fixed Gram form on the active links, not the inverse of the Wilson
quadratic form. Its precision is bounded below uniformly in the volume.
- Window law: the global inverse is approximated by inverses on windows
of radius \(r\), with error \(Ce^{-vr}\) uniform in the volume. -
\(\log\det\) separates into local terms and far-field terms with
summable errors.

\textbf{Cluster expansion.} The gas is
\(\Xi(w)=\sum_s\prod_{i\in s}w_i\), summed over families of polymers
with pairwise disjoint supports and real weights. With
\(B=\lceil(\log\hat\beta)^2\rceil\), the smallness condition is

\[\sum_{i:\,x\in X_i}|w_i|\,e^{(a+\delta)\,\mathrm{cov}_B(X_i)}\le\gamma\ \ \text{for every site }x,\qquad 4(2B+1)^4\gamma/\delta\le q<1 .\]

Under this condition: - \(\log\Xi\) is an absolutely convergent sum of
connected clusters, bounded through spanning trees (Penrose's partition
scheme); - removing every polymer that meets a set \(Y\) changes
\(\log\Xi\) by at most \(|Y|\gamma/(1-q)\), uniformly in the volume.

So the vacuum logarithm is extensive, but a local change costs only its
size. A second, graph-free argument based on a deletion recursion gives
bounds of the same type.

\textbf{Connected activities and dependence cones.} - Fixed-order
connected coefficients come from the heat equation
\(\partial_tW=\tfrac12(\Gamma:\nabla^2W+\Gamma:\nabla W\,\nabla W)\) for
\(W_t=\log\mathbb E_{t\Gamma}e^V\). A rooted \(\ell^1\) norm with weight
\(e^{\sigma\,\mathrm{diam}}\) makes contraction and joining lossless, so
these coefficients are polynomially bounded uniformly in the volume.
There is no resummation. - Pulling a coarse observable back through the
block kernel gives a function of the fine links in a window of fixed
radius around its support. The backward cone therefore grows by a
bounded physical width per step. - For disjoint supports, the pull-back
factorizes exactly.

\paragraph{Large fields and one nonlinear
step}\label{large-fields-and-one-nonlinear-step}

\emph{Unit:
\href{https://github.com/scottnarmstrong/yang-mills-informalization/blob/main/outline/units/III-10.md}{III.10
--- One full nonlinear renormalization step: bounds on irrelevant
terms}.}

At reference coupling \(\hat\beta\) the small-field threshold is

\[t=\hat\beta^{-1/2}P(\hat\beta),\qquad P(\hat\beta)=\exp\big((\log\log\hat\beta)^2\big),\]

well above the Gaussian fluctuation scale \(\hat\beta^{-1/2}\).

\textbf{Seeds.} A seed is one of: - a fine plaquette with chord at least
\(t/4\); - a coarse link on which the coarse field departs from the
predictor by more than a fixed multiple of \(t\); - a flagged coarse
site.

\textbf{Energy lower bound.} For every store \(A\) of seeds,

\[\kappa\,t^2\,\sharp A\ \le\ \text{quadratic energy of the step integrand, after subtracting the stationary predictor}.\]

The energy is a sum over five channels, one of which is the blocking
penalty. The proof assigns charges to fine cells with bounded
multiplicity, and controls the linearization defects on cells without
seeds.

\textbf{Why this suffices.} Since \(\beta t^2\simeq P^2\), each seed
gains a factor \(e^{-cP^2}\). \(P^2\) grows faster than any power of
\(\log\hat\beta\), so this gain beats: - the polylogarithmic number of
connected seed sets through a site; - any per-seed loss of a fixed power
of \(\hat\beta\), or of order \(P^z\) with \(z<2\).

The decomposition into large-field sectors is exact: large-field sectors
are integrated, not discarded.

\textbf{One step.} The step takes inputs at level \(m+1\): regular norms
at most \(\hat\beta^{-11/10}\), forced norms at most \(e^{-P^{1/5}}\),
and the coupling in the band \([\hat\beta/2,2\hat\beta]\). For \(L\)
large, and then \(\hat\beta\) large, it produces: - site-summable
budgets for the new regular activities; - the bound
\(\exp\big((C\,\mathrm{bill}^5-\beta\mu t^2)\sharp A\big)\) for each
large-field sector, where \(\mathrm{bill}=1+B+(\log\hat\beta)^2\); -
bounds on the probabilities of the remaining catastrophic events.

The marginal and vacuum terms are extracted, and the level-\(m\) bounds
restored, by the reset below.

\paragraph{The marginal coupling and the
schedule}\label{the-marginal-coupling-and-the-schedule}

\emph{Units:
\href{https://github.com/scottnarmstrong/yang-mills-informalization/blob/main/outline/units/III-7.md}{III.7
--- Matching the Wilson energy to fourth order};
\href{https://github.com/scottnarmstrong/yang-mills-informalization/blob/main/outline/units/III-9.md}{III.9
--- The marginal flow: the running coupling and its reset};
\href{https://github.com/scottnarmstrong/yang-mills-informalization/blob/main/outline/units/III-2.md}{III.2
--- Choosing the constants: the order of choice and the RG schedule}.}

\textbf{Matching and the one-loop coefficient.} - At the finest level, a
path representation is constructed whose jets of order at most 4 agree
with those of the Wilson energy in exponential coordinates. - At later
levels, the cubic and quartic parts are represented in curvature
variables by quasi-local lifts (a cost-controlled Poincaré lemma). - In
each step, the coefficient \(\alpha_{N,m}\) of curvature squared
produced by the one-loop term satisfies
\(|\alpha_{N,m}+\kappa_G\log L_m|\le D\), with \(D\) fixed before \(L\).

The value of \(\kappa_G\) comes from an explicit momentum integral. The
comparison symbol is

\[h(k)=4C_{\mathfrak g}\Big(\frac{k_1^2}{|k|^6}-\frac{4k_1^2k_0^2}{|k|^8}-\frac1{|k|^4}\Big).\]

Its angular average is
\(4C_{\mathfrak g}(\tfrac14-\tfrac16-1)|k|^{-4}=-\tfrac{11}3C_{\mathfrak g}|k|^{-4}\),
and \(\int_{1/L<|k|\le1}|k|^{-4}\,d^4k=2\pi^2\log L\). Hence
\((2\pi)^{-4}\int_{1/L<|k|\le1}h=-\kappa_G\log L\). Comparing the step's
Hessian with \(h\) up to \(O(|k|^{-3})\) gives the bound on
\(\alpha_{N,m}\).

\textbf{The reset.} After the step, the blocked density is rewritten
exactly, at every configuration including large fields: - a new coupling
\(\beta_{N,p,m}=\beta_{N,p,m+1}+\alpha_{N,m}+\delta\); - a new vacuum
constant; - reorganized activities: flat values and isotropic quadratic
parts are subtracted, compensating path terms are added, and mismatches
of the small-field masks go into the forced gas.

The correction \(\delta\) is the isotropic Hessian coefficient of the
regular activities: - By simplicity and hyperoctahedral symmetry, their
quadratic part at a point is \(\delta\) times the identity on plaquette
orientations. - \(|\delta|\le1\), and \(\delta\) is fixed by a scalar
equation with a unique root.

Altogether, with \(\omega_j\) a drift allowance attached to regularity
index \(j\),

\[\big|\beta_{N,p,m+1}-\beta_{N,p,m}-\kappa_G\log L_m\big|\le\omega_{J_{m+1}} .\]

\textbf{Endpoint.} For real \(|b|<r_\ast\) there is a trajectory with
level-\(0\) coupling \(\bar\beta+b\). Its bare coupling is unique within
the construction, and the trajectory continues analytically to complex
\(|b|<r_\ast\), with \(r_\ast\) independent of \(N\) and \(p\). The
family uses \(b=0\), so \(\beta_{N,p,0}=\bar\beta\) for all \(N\) and
\(p\). An intermediate-value shooting argument also produces such
trajectories; the analytic version is the one used for convergence at
fixed volume below.

\textbf{Running.} Put \(s_m=\log(1/a_m)\) and
\(\hat\beta_m=\bar\beta+\kappa_Gs_m\). Induction from level \(0\) gives
\(|\beta_{N,p,m}-\hat\beta_m|\le\sum_{k\le m}\omega_{J_k}\), and the
schedule makes this \(o(s_m)\):

\[\beta_{N,p,m}=\bar\beta+\kappa_G\log(1/a_m)+o\big(\log(1/a_m)\big)\qquad\text{uniformly in }N\ge m\text{ and }p.\]

At \(m=N\) this is the slope law with \(b_\ast=\kappa_G\). Since
\(\bar\beta\) is chosen large, every level stays in the band
\([\hat\beta_m/2,2\hat\beta_m]\). The statement is leading-logarithmic
only: there is no two-loop term.

\textbf{Why the block ratios grow.} The one-loop error per step is
bounded but does not decrease with \(L\). For the accumulated error to
be \(o(s_m)\), the allowance per unit of logarithmic scale must tend to
zero. The schedule therefore imposes two conditions: -
\(\omega_{J_{m+1}}\le J_{m+1}^{-2}\log L_m\), which forces
\(\log L_m\to\infty\); - \(\log L_m/s_m\to0\), which keeps each step
small compared with the total scale.

The ratios are constant while the regularity index is fixed,
\(L_m=\ell_{J_{m+1}}\). The regularity index \(J\) is nondecreasing,
rises by at most one per level, tends to infinity toward the fine end,
and stays constant on long windows of levels. One sequence \((L_m)\)
serves every cutoff \(N\).

\textbf{Order of choices.} The constants are chosen in a fixed order: 1.
the kernel with its penalty, and for every index \(j\) its regularity
profile (derivative order, spatial power, support exponent, widths, path
counts, drift allowance), together with the one-loop error constant
\(D\); 2. all block ratios \(\ell_j\), so that the allowances and the
logarithms of the combinatorial costs are at most \(j^{-2}\log\ell_j\);
3. the smallness exponents; 4. the entry couplings and the volume
conditions; 5. \(\bar\beta\); 6. the assignment of regularity profiles
to levels.

Each degree \(d\) has a reference level \(r(d)\) with regularity index
at least \(d+3\). Beyond \(r(d)\) the derivative order exceeds \(d\) and
the spatial power is at least \(d\). All choices precede any observable
or extraction order, so every degree is served by the same family.

\paragraph{Renormalized observables, uniformly in the
volume}\label{renormalized-observables-uniformly-in-the-volume}

\emph{Units:
\href{https://github.com/scottnarmstrong/yang-mills-informalization/blob/main/outline/units/III-11.md}{III.11
--- Comparison with the free Gaussian measure: interacting rank and
determinant budgets};
\href{https://github.com/scottnarmstrong/yang-mills-informalization/blob/main/outline/units/III-12.md}{III.12
--- Bounds uniform in the volume and convergence as the spacing tends to
zero};
\href{https://github.com/scottnarmstrong/yang-mills-informalization/blob/main/outline/units/III-15.md}{III.15
--- Cost accounting for observables and the remainder comparison};
\href{https://github.com/scottnarmstrong/yang-mills-informalization/blob/main/outline/units/III-16.md}{III.16
--- Analytic toolkit}.}

\textbf{The chessboard root.} The quantity propagated for observables is
the chessboard root. For a function \(F\) of the links in one cell of a
tiling of the torus by cells of side \(a_m\),

\[\lVert F\rVert_T=\Big|\,\mathbb E_{\beta_{N,p}}\Big[\prod_{\text{cells}}\text{reflected copies of }F\Big]\Big|^{1/\sharp \text{cells}} .\]

The main bound reads as follows. For every degree \(d\), radius \(R\)
and \(\varepsilon>0\) there is \(A\) such that, for every \(m\ge r(d)\),
there is \(N_0\) independent of \(p\) with

\[\lVert O_{N,p}(v,o)\rVert_{T_m(o)}\le A\,a_m^{-d-\varepsilon}\qquad(N\ge N_0,\ \text{all }p,\ \lVert v\rVert\le R).\]

It is proved by running the insertion through the same RG steps;
\(\varepsilon\) absorbs losses smaller than any power. Reflection
positivity enters only afterwards, through the chessboard inequality.
That inequality makes the root a seminorm and bounds correlations of
separated observables by products of roots.

\textbf{No loss of \(\log\beta\) per site.} A root over \(c\) copies
tolerates a fixed cost per copy, but not a factor \(\log\beta\) per
site. The determinant budget avoids this by a gauge count on a terminal
torus with \(c=L^4\) sites: - A tree gauge leaves \(4c-(c-1)=3c+1\) free
links. - In the numerator, \(3c-3\) of them are integrated one at a
time, each being the only new link of a chosen plaquette. Each integral
gives at most \(b\,E^{\dim G}\), with \(E\propto g=\beta^{-1/2}\). The
four links that wrap around the torus, one per axis, are left over. -
The normalization is bounded below by confining all \(3c+1\) free links
to Haar balls of radius \(g/4\). - The ratio is therefore at most
\(\exp(c\,\pi-4\dim G\,\log g)\), where the cost \(\pi\) per site is
independent of \(g\). - Since the number of copies is \(c\), the root
carries only \(g^{-4\dim G/c}\).

\textbf{Field content.} Free Gaussian detectors give uniform rank
detection:
\(\lVert\pi x\rVert\le Z\,\lVert\mathrm{probe}(\pi x)\rVert\), where
\(\pi\) is the projection onto quotient labels. The limit field spaces
are isomorphic, degree by degree, to those of free Maxwell theory,
compatibly with conjugation and reflections.

\textbf{Difference quotients.} - A normalized lattice difference of an
insertion of degree \(d\) equals an insertion \(O(x)\) of degree
\(d+|l|\), up to a term whose root tends to zero. - Write
\(x=\pi x+(x-\pi x)\). The part \(x-\pi x\) lies in the Maxwell equation
ideal and has small root by the lattice equations of motion. The part
\(\pi x\) is bounded by the rank detector. - So difference quotients
satisfy root bounds of the same form. By the chessboard inequality,
lattice correlations with any number of differences are bounded at
\(r\)-separated points, uniformly in \(p\).

These bounds are propagated through the RG; they do not come from
differentiating a sup norm.

\textbf{A rank-free trace bound.} Let \(X\ge0\) be a matrix and let the
insertions \(b_a\) be separated by gaps of more than \(2r\) factors of
\(X\). Moving \(m\le r\) factors from one gap to another changes
\(\mathrm{tr}\prod_ab_aX^{Y_a}\) by at most \((\mathrm{cost}/r)\,m\).
The cost is a product of trace norms with no dimension factor. The proof
is a Hölder inequality for complex exponents, obtained from Hadamard's
three-lines theorem, followed by a Cauchy estimate. By quadrature it
extends to continuous positive kernels. Part IV uses it for
equicontinuity.

\textbf{Bookkeeping.} The numerical conditions of all steps can be met
in one fixed order: constants before \(L\), then \(L\), then the later
constants, then \(\beta\) large. This holds uniformly across regularity
profiles.

\paragraph{Convergence at fixed volume and the construction
run}\label{convergence-at-fixed-volume-and-the-construction-run}

\emph{Units:
\href{https://github.com/scottnarmstrong/yang-mills-informalization/blob/main/outline/units/III-12.md}{III.12};
\href{https://github.com/scottnarmstrong/yang-mills-informalization/blob/main/outline/units/III-13.md}{III.13
--- Extension to all labels: one common Wilson family};
\href{https://github.com/scottnarmstrong/yang-mills-informalization/blob/main/outline/units/III-14.md}{III.14
--- The construction run: flat representations and the ultraviolet
rate}.}

\textbf{The obstacle.} At fixed volume, two cutoffs \(N\) and \(N'\)
have different bare couplings, each fixed by the condition
\(\beta_{N,p,0}=\bar\beta\). Continuity in a fixed bare coupling
therefore compares nothing.

\textbf{Fixed-level convergence and tails.} The construction has a
one-parameter family of trajectories: - The level-\(0\) coupling is
\(\bar\beta+ws\), with a real tuning parameter \(s\in[-1,1]\). For each
\(s\) the bare coupling is unique within the recursion. - The bare
coupling and effective data have analytic dependence on the tuning
parameter. The estimates on that dependence are separate from the Cauchy
argument below. - For fixed \(j\), convergence of the trajectory and
stability of the observable transfers give convergence of the effective
insertions and their polymer parts. Uniform convergence passes through
the Haar integrals, and a positive lower bound on the normalized
partition function gives convergent correlation approximations. - The
omitted part is bounded, uniformly in the cutoff, by
\(2F_0e^{-B_j/64}\), where
\(B_j=\lceil(\log\hat\beta_j)^2\rceil\to\infty\). - Choose \(j\) to make
this tail small, then take the cutoffs large enough that their
level-\(j\) approximations are close. This proves the Cauchy property
first at lattice-fitted configurations; equicontinuity extends it to all
configurations distinct modulo the period. See
\href{https://github.com/scottnarmstrong/yang-mills-informalization/blob/main/outline/units/III-12.md\#45-convergence-in-the-mesh}{III.12,
§§4.5--4.6}.

Combined with the difference bounds, this gives at each fixed volume: -
limit kernels smooth away from periodic coincidences; - convergence of
difference quotients to derivatives; - convergence of smeared moments
along the whole sequence of spacings.

\textbf{The family used by the public theorem.}
\href{https://github.com/scottnarmstrong/yang-mills-informalization/blob/main/outline/units/VIII-1.md}{VIII.1}
assembles the construction for one family, using the bounds and
convergence results described in III.9--III.12.
\href{https://github.com/scottnarmstrong/yang-mills-informalization/blob/main/outline/units/III-13.md}{III.13}
describes an alternative coordinated family with root bounds, an
analytic shooting family and Maxwell rank comparisons. Its existence
statements do not identify that family with the one used by
\texttt{YangMills.constructed}; properties from the two constructions
should not be combined without such an identification.

\textbf{The construction run (III.14).} III.14 packages the construction
run. It also proves quantitative convergence of the order-zero quadratic
part of the effective action (its position-space Hessian kernel) at each
fixed level \(i\), at rate \((a_N/a_i)^c\). - The difference between two
cutoffs is written in curvature variables as \(d^{\mathsf T}Ad\), with
\(A\) defined modulo kernels annihilated by this map. - So the curvature
precision, which has zero modes, is never inverted. - This rate concerns
the quadratic germ only, not the full effective action or the final
limits.

\subsubsection{Part IV. Infinite volume}\label{part-iv.-infinite-volume}

\emph{Unit:
\href{https://github.com/scottnarmstrong/yang-mills-informalization/blob/main/outline/units/IV-1.md}{IV.1
--- Infinite volume: extraction of a subsequence and the Euclidean
axioms of the limit}.}

\textbf{Local bounds.} - Reflection positivity gives the chessboard
inequality \(|\langle\prod_cF_c\rangle|\le\prod_c\lVert F_c\rVert_T\)
for every \(\beta\ge0\). Log-convexity of roots under translation
handles cells not centred at the insertion. - Put each insertion in its
own cell of side \(a_m<r\). Then \(|C_{N,p}(v;x)|\le A^n\) whenever each
pair of points has some coordinate at distance at least \(r\), uniformly
in \(p\) and for large \(N\). - Near collisions, the same argument gives
\(|S(v;x)|\le Cr^{-4n(1+\max d_i)}\) when all points are at mutual
distance at least \(r\).

\textbf{Equicontinuity.} Let \(\mathrm{Gen}_n\) be the set of
configurations whose points have pairwise distinct coordinates on every
axis. 1. Integrating out the links along one axis gives a positive
transfer kernel \(K\). A correlation becomes
\(\mathrm{Tr}\prod_kB_kK^{Y_k}/\mathrm{Tr}\,K^M\), with slab kernels
\(B_k\) and gaps \(Y_k\). 2. Moving one point shifts units between two
gaps. The trace is entire in the size of the shift. 3. The rank-free
Hölder bound controls the trace on a disc, and a Cauchy estimate gives a
Lipschitz bound. 4. The trace norms of the slab kernels are themselves
chessboard roots of \(2q\) alternating copies, so the bound is uniform
in the volume.

The result is \(|C(x)-C(y)|\le16nA^n|x-y|/R\) on compact subsets of
\(\mathrm{Gen}_n\) whose coordinate gaps exceed \(R\). No derivative
bound is used.

\textbf{Extraction.} Arzelà--Ascoli is applied twice, over countably
many degrees, basis labels and compact sets: first to extract a spacing
subsequence that works for all volumes, then a volume subsequence. This
gives limit kernels on \(\mathrm{Gen}_n\). The complement
\(\mathrm{Conf}_n\setminus\mathrm{Gen}_n\) is a null set, and the
kernels are locally bounded on \(\mathrm{Conf}_n\), so dominated
convergence gives the smeared limits.

\textbf{Derivatives and smoothness.} - Exact summation by parts in the
lattice positions moves a difference quotient onto the test function. -
The ultraviolet estimates show that the lattice difference quotient of
an insertion is chessboard-close to an insertion of one degree higher. -
A dimension count, from the rank comparison, identifies the kernel of
the limiting projections with the null labels. - Finite-dimensional
compactness then gives fixed derivative maps
\(\partial_a:F_d\to F_{d+1}\) with

\[\partial_{x_i^a}S(\dots,v_i,\dots;x)=S(\dots,\partial_av_i,\dots;x)\qquad\text{in }\mathcal D'(\mathrm{Conf}_n).\]

Iteration and a Fourier argument give
\(S\in C^\infty(\mathrm{Conf}_n)\), with derivatives polynomially
bounded near collisions.

\textbf{What the limit inherits.} The limit inherits: - translations; -
the group of signed permutations of the axes, acting on points and
labels; - permutation symmetry, reality and the unit label; - reflection
positivity at coordinate planes, for point and smeared tests.

It does not inherit continuous rotations, clustering or uniqueness.

\textbf{Removing the spacing refinement.} The extraction refines both
the spacing and the volume sequences. At each admissible fixed volume,
the smeared limit along the whole spacing sequence exists (Part III) and
agrees with the subsequential one. This removes the spacing refinement
from the final statement.

\subsubsection{Part V. Phases and relativistic
reconstruction}\label{part-v.-phases-and-relativistic-reconstruction}

\emph{Units:
\href{https://github.com/scottnarmstrong/yang-mills-informalization/blob/main/outline/units/V-1.md}{V.1
--- Decomposition of the limit into phases};
\href{https://github.com/scottnarmstrong/yang-mills-informalization/blob/main/outline/units/V-2.md}{V.2
--- Osterwalder--Schrader reconstruction of a Wightman theory};
\href{https://github.com/scottnarmstrong/yang-mills-informalization/blob/main/outline/units/V-3.md}{V.3
--- Reconstruction for every phase and the reconstructed theory}.}

\textbf{The OS space of the mixture.} Positive-time field words, modulo
the null space of \(\langle\theta w,w'\rangle_S\) and completed, form a
Hilbert space. It carries: - a contraction semigroup \(e^{-tH}\) with
\(H\ge0\); - a unitary representation of the spatial translations.

The joint spectrum lies in a half-space. No positive spectral gap is
asserted. The zero-energy space may be large. Its projection \(\Pi\) is
the strong limit of \(e^{-tH}\) as \(t\to\infty\).

\textbf{Thin slabs.} A word \(w\) placed in a thin time slab defines an
operator \(B\) on the OS space. To bound it, consider the train
\(v_{n+1}=B^\ast Bv_n\), which contains \(n\) reflected copies of \(w\)
and an arbitrary exterior word. - The lattice bounds give
\(\lVert v_n\rVert^2\le B_vA^{4n}\). Here \(A\) is polynomial in the
word data and in the inverse buffer width, and \(B_v\) depends on the
exterior word but not on \(n\). - A power lemma removes \(B_v\): if
\(\lVert(BA)^nv\rVert\le K_vD^{2n}\) for all \(n\), then
\(\lVert Av\rVert\le D\lVert v\rVert\).

So the exterior word is paid once and disappears under the \(n\)-th
root.

\textbf{Phases.} - The compressions \(M(w)=\Pi\,B_w|_{\ker H}\) are
independent of the width and placement of the slab. -
\(\langle\Omega,M(w_1)\cdots M(w_k)\Omega\rangle\) is the limit of \(S\)
as the clusters are moved apart in time order.

The compressions commute: 1. Place four clusters on the rays
\(n(ue_4+ve_1)\), with \((u,v)=(0,0),(1,2),(2,1),(3,3)\). 2. Exchanging
the time axis with a spatial axis reverses the order of the middle pair.
3. \(O(4)\) covariance and permutation symmetry of \(S\) then give
\(\langle\Omega,M_xM_yM_wM_z\Omega\rangle=\langle\Omega,M_xM_wM_yM_z\Omega\rangle\).
4. By density, \([M_y,M_w]=0\).

The resulting phase decomposition: - The \(M(w)\) generate a commutative
unital \(C^\ast\)-algebra. Its Gelfand spectrum is \(X\), which is
compact and metrizable. - \(\nu\) represents the vacuum state, and it
has full support because the vacuum is separating. - The phase
correlations are \(S^\xi(w)=\xi(M(w))\). They are continuous in \(\xi\),
and \(\int S^\xi\,d\nu=S^{\mathrm{mix}}\). - Reflection positivity,
\(O(4)\) invariance, the derivative and stress identities and uniform
growth bounds hold for every \(\xi\). They are proved as operator
identities before any character is applied. - Every phase has the same
null spaces as the mixture. - A variance argument over a countable total
family of words gives a one-dimensional zero-energy space for
\(\nu\)-almost every phase. - The selected phase \(\xi_\ast\) is one
such phase, chosen by this existence argument; no further condition is
imposed. The trace law of Part VII is proved for the phase so chosen.

\textbf{Reconstruction.} For each phase, slab operators bounded by
\(C\,p^{-A}\) (with \(p\) the slab half-width) define

\[F_n(z;\mathbf x)=\big\langle\Omega,\,B(v_1)\,e^{i(z_1-z_2)H}U(\mathbf x_1-\mathbf x_2)\,B(v_2)\cdots\Omega\big\rangle .\]

This function is holomorphic for
\(\mathrm{Im}\,z_1<\dots<\mathrm{Im}\,z_n\), and it equals the Euclidean
correlation at imaginary times. The Wightman theory is obtained in five
steps: 1. \textbf{Boundary values.} Bounds of order \(k^{-m}\), when all
height gaps exceed \(k\), give tempered boundary values \(W_n\). They
converge against every joint Schwartz test as the vector of heights
tends to zero through all ordered vectors. 2. \textbf{Positivity and
fields.} This mode of convergence allows two independently regularized
vectors to be compared. So \(W\) is a limit of Gram matrices of OS
vectors: it is positive, and the Wightman space embeds isometrically
into the OS space. Fields act on finite words modulo the null space by
prepending a letter. This gives the invariant domain, the formal
adjoints and the cyclic vacuum. 3. \textbf{Covariance and spectrum.} The
rotation Ward identities of Part VI, in planes containing the time axis,
continue to infinitesimal boost identities. These integrate to Lorentz
covariance without any spectral assumption. The spectral measure of a
Wightman vector equals that of its OS image, so it is supported in
\(p^0\ge0\); covariance then moves the support into the forward cone. 4.
\textbf{Locality.} Only adjacent transpositions are needed. For a
spacelike adjacent pair, both orderings have holomorphic germs,
constructed with complex boosts. The germs agree to all orders on the
Euclidean slice, because the Euclidean kernel is symmetric and smooth
away from coincidences. 5. \textbf{Faithfulness.} A label whose Wightman
functions vanish has a vanishing tube function, hence a vanishing
Euclidean kernel.

\textbf{Every phase.} The low-energy bound uses only locality and the
spectrum condition. Take \(A\) a finite sum of field words with
compactly supported tests. - The commutator of time-translated copies of
\(A\) vanishes for times at most \(1\) and spatial separations beyond a
radius \(R\) independent of \(\xi\). - Average over unit time, pair with
a spatial-momentum bump of radius \(2u\) (whose Fourier transform is
\(O(u^3)\) in three space dimensions), and use \(|\mathbf p|\le E\) on
the forward cone. This bounds the spectral weight of each shell
\(u/2<E\le u\) by \(O(u^2)\), and dyadic summation gives the first bound
below.

\[\mu_{A\Omega_\xi}(0<E\le u)\le M_Au^2\ \ (0<u\le1),\qquad 0\le\big\langle A\Omega_\xi,(e^{-sH_\xi}-P_{0,\xi})A\Omega_\xi\big\rangle\le M_As^{-2}\ \ (s\ge1).\]

The constant \(M_A\) is independent of \(\xi\), by compactness of \(X\).

Vacuum uniqueness then extends from almost every phase to every phase: -
The function
\(\xi\mapsto\langle A\Omega_\xi,e^{-sH_\xi}A\Omega_\xi\rangle-|\langle\Omega_\xi,A\Omega_\xi\rangle|^2\)
is continuous. - For fixed \(A\) and \(s\), the set of phases where this
function lies in \([0,M_As^{-2}]\) is closed by continuity. It contains
the full-measure set of phases with unique vacuum; full support of
\(\nu\) makes this closed set all of \(X\). - Letting \(s\to\infty\)
gives a unique vacuum in every phase.

The comparison map from the Wightman space to the OS space is
surjective, proved independently by pushing time-ordered vectors apart.
This transfers vacuum uniqueness to the OS space. Each phase theory is
irreducible. Its Wightman distributions are continuous in \(\xi\), and
their \(\nu\)-average is the boundary value of the continuation of
\(S^{\mathrm{mix}}\). The bounds allow spectrum down to zero and give
only polynomial decay \(s^{-2}\) in Euclidean time, not exponential
clustering.

\subsubsection{Part VI. Ward identities, rotations and the stress
tensor}\label{part-vi.-ward-identities-rotations-and-the-stress-tensor}

\emph{Units:
\href{https://github.com/scottnarmstrong/yang-mills-informalization/blob/main/outline/units/VI-1.md}{VI.1
--- Ward identities, the rotation action and the conserved stress
tensor};
\href{https://github.com/scottnarmstrong/yang-mills-informalization/blob/main/outline/units/VI-2.md}{VI.2
--- The free Maxwell stress tensor and the local Ward charge}.}

The lattice has only the hypercubic symmetries. The proof obtains the
generators of rotations from stress fluxes, and the group action by
comparison with free Maxwell theory. It uses the OS spaces of the mixed
state in the coordinate directions and the free comparison of Part VII.
It uses no decay of correlations.

\textbf{Free Maxwell theory.} Consider \(\dim\mathfrak g\) free Maxwell
fields, with field spaces the \(G\)-invariant polynomials in on-shell
curvature jets. - The centred, real, symmetric tensors of degree 4 that
are covariant under signed permutations form a three-dimensional space.
- For simple \(\mathfrak g\), the divergence maps this space onto a
two-dimensional space. So the conserved tensors form a line, spanned by
the Maxwell stress tensor. - For every Killing field \(v\), the flux of
\(\sum_bv_bT_{ab}\) through a shell around an insertion equals the Lie
derivative of the insertion. For elementary fields this follows from
\(-\Delta G_0=\delta\); for polynomials, from the Leibniz rule and Wick
contraction. - An explicit three-point function fixes the sign and
normalization of the Ward constant.

\textbf{Conserved tensor.} There are scaling sequences
\(r_k\downarrow0\) with frames \(U_k\), isomorphisms for large \(k\).
Along them the rescaled interacting divergence converges to the free
one. - Surjectivity is an open condition, so the interacting divergence
is also onto, and its kernel is one-dimensional. - A nonzero conserved
\(J\) in this kernel converges, after rescaling, to a multiple of the
free stress tensor.

\textbf{Shell identities.} Lattice integration by parts in the link
variables, with ultraviolet control of the error terms, gives shell
identities. Take \(\chi\) equal to \(1\) near \(x_0\) and vanishing near
the spectators:

\[-\sum_a\int\partial_a\chi(y)\,S\big(J_{ab}(y)\,B(x_0)\cdots\big)\,dy=S\big((p_bB)(x_0)\cdots\big).\]

With weight \(K(y-x_0)\), for a rotation generator \(K\), the same
construction gives maps \(Y_K\). Here \(p_b\) and \(Y_K\) are linear
maps on labels, not yet identified with derivatives and rotations.
Moving the centre of the shell gives \(\partial p=p\,\partial\) and
\(\partial_bY_K+\sum_cK_{cb}\,p_c=Y_K\partial_b\).

\textbf{Flux at infinity.} - For rotations, cut off the rotation field
so that it is compactly supported and equal to the exact field on a box
far larger than the configuration. - Conservation and symmetry cancel
the cross terms. - The remaining exterior contribution is bounded
through spectral measures of the OS semigroups along and transverse to a
coordinate line, and it tends to zero. - This gives the angular identity
for configurations on one coordinate line, together with a line-sum
identity, and hence transversality of the label maps.

\textbf{Normalization.} In free theory, the finitely many relations
satisfied by \(p\) and \(\partial\) at a given degree have a
one-dimensional solution space. Injectivity is an open condition, so the
same holds along the scaling sequence. Hence \(p_b=c\,\partial_b\), with
\(c\) independent of the degree and \(c\ne0\). The tensor \(T=J/c\)
satisfies

\[-\sum_a\int\partial_a\chi(y)\,S^s\big(T_{ab}(y)\,B(x_0)\cdots\big)\,dy=\frac{\partial}{\partial x_0^b}\,S^s\big(B(x_0)\cdots\big)\]

in the mixed state. By the operator identities of Part V, it holds in
every phase. Since \(\partial\chi\) vanishes near every insertion, no
contact terms arise.

\textbf{All separated configurations.} Positivity of Gram functions,
together with the Ward values already obtained, gives analytic-vector
bounds \(\lVert f^{(m)}\rVert\le C(\rho/2)^{-m}m!\). The angular Ward
functional is therefore real-analytic along walks that move one point at
a time. Starting from configurations on a line, it vanishes on all
configurations whose points have pairwise distinct first coordinates,
and by continuity on all of \(\mathrm{Conf}_n\).

\textbf{Integration to \(O(4)\).} - The bracket defect
\([Y'_u,Y'_v]-Y'_{[u,v]}\) (with \(Y'=Y/c\)) satisfies a summed
identity. A rigidity statement, proved in free theory and transferred by
the comparison, forces the defect to vanish. - The infinitesimal action
converges to the Maxwell one:
\(U_k^{-1}Y'_KU_k\to\mathrm{Lie}^{\mathrm{Max}}_K\). - A perturbation of
spectral projections of the Casimir operators gives exact intertwiners
\(V_d\). The action on \(F_d\) is
\(R_d(k)=V_dR^{\mathrm{Max}}_d(k)V_d^{-1}\). - Since
\(t\mapsto S(R(e^{tK})v;e^{tK}x)\) has zero derivative, and reflections
act through the real structure, \(S(R(k)v;kx)=S(v;x)\) for all
\(k\in O(4)\). - Averaging \(J/c\) over \(O(4)\) gives a conserved,
real, symmetric, covariant tensor that satisfies both the translation
and the rotation Ward identities.

\subsubsection{Part VII. Identification as Yang--Mills
theory}\label{part-vii.-identification-as-yangmills-theory}

\paragraph{Free Maxwell fields at short
distance}\label{free-maxwell-fields-at-short-distance}

\emph{Units:
\href{https://github.com/scottnarmstrong/yang-mills-informalization/blob/main/outline/units/VII-1.md}{VII.1
--- Free Maxwell comparison theory};
\href{https://github.com/scottnarmstrong/yang-mills-informalization/blob/main/outline/units/VI-2.md}{VI.2
--- Interacting stress comparison};
\href{https://github.com/scottnarmstrong/yang-mills-informalization/blob/main/outline/units/VII-6.md}{VII.6
--- Scaling variants and fixed-field estimates}. The comparison is
assembled in
\href{https://github.com/scottnarmstrong/yang-mills-informalization/blob/main/outline/units/VIII-1.md}{VIII.1}.}

\textbf{The reference theory.} It consists of \(\dim\mathfrak g\) free
Maxwell fields built from \(\Gamma(x)=1/(4\pi^2|x|^2)\). - On the
lattice, corrected covariances of finite-difference jets converge to
\(\partial_a\partial_b\Gamma\) away from coincidences. - Wick moments
satisfy the same Gaussian recursion on the lattice and in the continuum,
so correlations of quotient labels converge. - A triangular basis built
from the Kelvin transform gives strict reflected Gram positivity for
every label outside the ideal \(I_d\) generated by the free Maxwell
equations and the Bianchi identities. - So the free null space is
exactly \(I_d\), and the free field space \(M_d=P_d/I_d\), where \(P_d\)
is the corresponding space of polynomials in free curvature jets, is
detected by finitely many correlations.

\textbf{The comparison.} There are scales \(r_k\downarrow0\) and linear
frames \(U_k\) on the field spaces with the following property. The
correlations of the mixed state, rescaled by \(r_k\) and transformed by
\(U_k\), converge to the free Maxwell correlations: - with all
derivatives; - uniformly on compact subsets of \(\mathrm{Conf}_n\); -
for all degrees along one sequence.

The stress tensor converges to the Maxwell stress tensor. An upper bound
on \(\dim F_d\) makes the frames isomorphisms for large \(k\), so
\(F_d\cong M_d\).

The frames are necessary: the claim is not that a fixed normalization of
each composite field has a limit along every sequence of radii. The
schedule contains a real parameter that moves the scale at which
observables are read, and the estimates are uniform in it. Tuning this
parameter, every sequence of radii tending to zero has a subsequence
along which the comparison holds.

\paragraph{Running coupling}\label{running-coupling}

\emph{Unit:
\href{https://github.com/scottnarmstrong/yang-mills-informalization/blob/main/outline/units/VII-5.md}{VII.5
--- Asymptotic freedom of the running coupling}.}

With \(g_{N,p}(m)=\beta_{N,p,m}^{-1/2}\), Part III gives

\[g_{N,p}(m)\,\sqrt{\log(1/a_m)}\longrightarrow\frac1{\sqrt{2b_0}}\qquad(m\to\infty),\]

uniformly in \(N\ge m\) and \(p\). This is the leading-logarithmic
running of the coupling of the construction's effective actions. No
two-loop term, no \(\Lambda\) parameter and no comparison with another
renormalization scheme is claimed. For \(SU(N)\), the conventionally
normalized coupling \(g_{\mathrm{std}}=\sqrt{2N}\,g\) satisfies
\(g_{\mathrm{std}}\sqrt{\log(1/a_m)}\to(11N/24\pi^2)^{-1/2}\).

\paragraph{The trace law and
nonfreeness}\label{the-trace-law-and-nonfreeness}

\emph{Units:
\href{https://github.com/scottnarmstrong/yang-mills-informalization/blob/main/outline/units/VII-2.md}{VII.2
--- The trace of the stress tensor and the trace anomaly};
\href{https://github.com/scottnarmstrong/yang-mills-informalization/blob/main/outline/units/VII-3.md}{VII.3
--- The theory is not free: Gram positivity and exclusion of free-field
profiles}.}

\textbf{Leading form of the trace.} The classical Maxwell and
Yang--Mills stress tensors are traceless. The quantum trace comes from
the scale dependence of the coefficient of curvature squared. - The
proof pulls the trace of the lattice stress observable back through the
RG steps. - Fourth-order matching identifies the leading part of the
trace as the derivative of the effective action with respect to scale. -
The one-loop coefficient enters through a momentum-shell estimate of the
kind used in Part III.

Along the scaling sequences, with \(\beta_k\) the running coupling at
the RG level that reads scale \(r_k\), and \(O_{F^2}\) the normalized
curvature-squared label,

\[\beta_k\,r_k^4\,U_k^{-1}\Theta\longrightarrow4d_2\,O_{F^2},\qquad d_2=-\frac{\kappa_G}4,\qquad\frac{\log(1/r_k)}{\beta_k}\longrightarrow\frac1{\kappa_G}.\]

In perturbative terms, this is the familiar
\(\Theta\simeq\frac{\beta(g)}{2g^3}F^2\), with
\(\beta(g)/(2g^3)=-b_0/2=d_2\) at one loop. The trace is a
curvature-squared field suppressed by
\(g^2(r)\simeq1/(\kappa_G\log(1/r))\).

\textbf{The two-point law.} Here
\(\Theta=\sum_\mu T_{\mu\mu}-S^{\xi_\ast}\big(\sum_\mu T_{\mu\mu}\big)\mathbf1\),
in the selected phase, which has a unique vacuum. 1. The free two-point
function of \(O_{F^2}\) is \(3\dim\mathfrak g/(\pi^4|x|^8)\). 2. So at
lattice-fitted configurations along the scaling sequence,
\(r_k^8\log^2(1/r_k)\langle\Theta\Theta\rangle\to(4d_2/\kappa_G)^2\cdot3\dim\mathfrak g/\pi^4=3\dim\mathfrak g/\pi^4\).
3. Rotation invariance moves the pair onto the time axis. 4. The
spectral representation in the phase gives
\(G_\Theta(r)=\int e^{-rE}\,d\mu(E)\) with a positive measure. So
\(G_\Theta\) is nonincreasing, which sandwiches \(G_\Theta(r)\) between
its values at fitted points. 5. If the limit failed along some sequence
of radii, the flexible scaling above would give a subsequence along
which it holds, a contradiction.

Hence the limit exists over all radii:

\[\lim_{r\downarrow0}\,r^8\log^2(1/r)\,G_\Theta(r)=\frac{3\dim_{\mathbb R}\mathfrak g}{\pi^4}.\]

The amplitude does not depend on the value of the one-loop coefficient.
\(\kappa_G\) enters once through the trace and once through the running,
and \((4d_2/\kappa_G)^2=1\) identically. So the test detects the
inverse-logarithmic suppression and the free normalization of \(F^2\),
not the value of \(b_0\); that value is established by the running law
of Part III.

\textbf{Exclusion of free profiles.} Consider a finite Wick profile
\(W(r)=\sum_\alpha c_\alpha\prod_\beta K_{\alpha\beta}(r)\). Each factor
\(K(r)=\int E^q\,h(\mathbf p)(2E)^{-1}e^{-rE}\,d^3\mathbf p\) has
\(E=\sqrt{\mathsf m^2+|\mathbf p|^2}\) and \(h\) homogeneous. - \(W\)
has an asymptotic expansion at \(r=0\) in the functions
\(r^a(\log r)^k\), with \(a\in\mathbb Z\) and \(k\ge0\). - Along
\(r=2^{-n}\), \(r^8\log^2(1/r)\,W(r)\) is then a finite sum of terms
\(c\,\rho^n n^{k+2}\), which can converge only to \(0\). - Since the
limit for \(G_\Theta\) is positive, \(G_\Theta\) differs from \(W\) at
points arbitrarily close to \(0\).

This excludes exactly the finite class of profiles in the statement, and
it concerns only \(G_\Theta\) on the time axis.

The stronger statement adds further tests: - a Wightman-side Wick test;
- Maxwell--Fock tests. Consider any isometric embedding of the selected
phase's Hilbert space into a Fock space of transverse photons that
preserves the vacuum and translations and has transverse symmetric
image. No such embedding makes \(\Theta\) local with respect to the Weyl
operators.

\paragraph{Operator product expansion and contact
extensions}\label{operator-product-expansion-and-contact-extensions}

\emph{Unit:
\href{https://github.com/scottnarmstrong/yang-mills-informalization/blob/main/outline/units/VII-4.md}{VII.4
--- Operator product expansion and contact extensions}.}

\textbf{The lattice statement.} The product of insertions in one cell
equals a single insertion of the same total degree, plus a remainder
with small chessboard root. The proof is an induction over the RG levels
between the cluster scale and the observation scale: the part of degree
at most \(D\) is kept, and the rest is paid into the remainder.

\textbf{The continuum coefficients.} They are read off rather than given
by a formula. - A prescription fixes, degree by degree, biorthogonal
systems of fields and detectors. The detectors are finite combinations
of mixed-state correlations of the cluster with fixed, separated
exterior words. - The prescription is averaged over \(O(4)\) and
symmetrized for reality. - The coefficient \(c_{j,\kappa}(a;x,o)\) is a
coordinate of the cluster's new component of degree \(j\) in a real
basis.

Consider an \((r,M,\eta)\)-cluster: points within \(Mr\) of \(o\) and at
mutual distance at least \(\eta r\), carrying fields of total degree
\(w\). Take spectators at distance at least \(c\) from \(o\) and from
each other. Then, for all sufficiently small \(r\),

\[|c_{j,\kappa}(a;x,o)|\le r^{\,j-w-\varepsilon},\qquad
\Big|S^s(a,x;b,y)-\sum_{j\le D}\sum_\kappa c_{j,\kappa}(a;x,o)\,S^s(O_{j,\kappa},o;b,y)\Big|\le r^{\,D+1-w-\varepsilon}.\]

\textbf{Uniformity and limits of the statement.} - The coefficients do
not depend on the spectators \(b\) and \(y\), on the state \(s\), or on
\(D\). - The threshold for \(r\) depends on the fields, on
\(M,\eta,c,\varepsilon,D\), and in the first bound on the channel, but
it is uniform in \(s\), \(x\), \(o\) and \(y\). - Uniformity in the
spectators comes from the train argument of Part V, in which the
exterior word is paid once. - Uniformity in the state comes from
bounding the error in the zero-energy algebra, whose norm dominates
every character and the mixture. - The loss \(\varepsilon\) absorbs
logarithms. There is no logarithmic refinement, so the bound is weaker
than the perturbative OPE. - Uniqueness of the coefficients is not
asserted. Neither convergence as \(D\to\infty\) nor agreement with
perturbation theory is claimed. - Nested versions over finite trees of
clusters, and integrated versions against measures of bounded variation,
also hold.

\textbf{Contact extensions.} These use Taylor subtractions over forests
of collapsing clusters on dyadic shells. The OPE remainder makes the
shell sums converge. For every state and every list of labels this gives
a distribution of finite order on \((\mathbb R^4)^n\) that agrees with
the correlation on \(\mathrm{Conf}_n\). Not established: uniqueness,
reflection positivity on overlapping supports, contact Ward identities,
real-time products, and convergence of the lattice correlations on the
diagonals.

\subsubsection{Part VIII. Assembly}\label{part-viii.-assembly}

\emph{Unit:
\href{https://github.com/scottnarmstrong/yang-mills-informalization/blob/main/outline/units/VIII-1.md}{VIII.1
--- Assembly of the main theorem}.}

The construction run depends only on \(G\) and the penalty \(b\). It is
chosen before any extraction order or observable. For each extraction
order (a spacing sequence and a volume sequence), the assembly composes
further refinements: - one refinement for the limits of the running
coupling; - one that produces the limit table, with the Euclidean
hierarchy, stress tensor, phases and reconstruction of Parts IV--VI.

To these it adds the reconstructions of all phases, the OPE
prescription, the nested and contact results and the Maxwell comparison.
Everything is transported to a common refinement, so all objects refer
to one table.

The public theorem is the case \(b=1\), with the full spacing sequence
and the given volume sequence; its \(q\) is the volume part of the
refinement.

No new estimate enters, but identifying all objects with one table is a
real constraint. Separate existence results for a lattice limit, a
Wightman theory, an OPE and a trace law would not show that they concern
the same correlation functions. The formal construction supplies the
hypotheses under which the limit theory is first proved. The outlines
have not explained every supplier in full; the
\href{https://github.com/scottnarmstrong/yang-mills-informalization/blob/main/outline/units/VIII-1.md\#7-what-remains-unexplained}{assembly
unit} records the limits of its source reading.

\subsubsection{Relation to earlier work, and where the difficulty
lies}\label{relation-to-earlier-work-and-where-the-difficulty-lies}

\textbf{Classical sources.} - Balaban's program is the model for the
ultraviolet analysis: block-spin transformations, gauge fixing around a
background field, fluctuation covariances built from averaging
operators, the small-field/large-field decomposition, and
renormalization of the marginal coupling. - The classical cluster
expansion (Glimm--Jaffe, Brydges, Kotecký--Preiss) organizes the polymer
activities. - Reflection positivity (Osterwalder--Seiler) and chessboard
estimates (Glimm--Jaffe--Spencer, Fröhlich--Israel--Lieb--Simon) connect
lattice estimates with bounds on observables. - Osterwalder--Schrader
and Glimm--Jaffe reconstruction turns Euclidean correlations into
Wightman fields. - The slope is the one-loop coefficient of
Gross--Wilczek and Politzer.

\textbf{Departures from Balaban's scheme.} - The block kernel is an
exactly normalized, gauge-covariant Markov kernel with random tags, not
a linear averaging operator in a fixed background gauge. - Integration
uses a spanning-forest gauge. - The block ratios grow along the levels.
- The threshold carries the factor \(e^{(\log\log\hat\beta)^2}\). - The
construction includes an analytic family in a tuning parameter. The unit
outlines do not establish a comprehensive comparison with Balaban's
papers.

\textbf{Beyond the ultraviolet analysis}, the proof does the following:
- One schedule and one family serve composite fields of every degree.
Each field has a fixed normalization, root bounds uniform in the volume,
and derivatives controlled by rank detection, with no coupling-dependent
loss per site. - Convergence at fixed volume is proved despite tuning
that depends on the cutoff. - Equicontinuity in infinite volume comes
from a rank-free trace inequality, without derivative bounds. -
Commutativity of the zero-energy algebra is proved, not assumed. -
Vacuum uniqueness in every phase follows from a uniform low-energy bound
that replaces a mass gap or a cluster argument. - Lorentz covariance
comes from Ward identities. - Rotation invariance is restored without
decay of correlations, and the group action is obtained by comparison
with free Maxwell theory. - The trace law holds as a limit over all
radii. - The OPE is derived from the RG, with one system of coefficients
for the mixture and every phase.

\textbf{Where the difficulty lies.} 1. \textbf{Gauge degeneracy and
large fields} must be controlled with a bounded cost per site and per
seed. This is done by the gauge count of the determinant budget, the
energy lower bound for seeds, and the exactness of the kernel. 2.
\textbf{Effective-action estimates must become estimates for
observables}: for every composite field, its derivatives and its
fixed-volume limit, inside one family. 3. \textbf{The passage from a
lattice mixture to relativistic phases} has no continuous symmetry and
no clustering to start from. It is handled by Ward identities with free
comparison, the commutative zero-energy algebra, and spectral bounds
based on locality.

The order of choice of constants and the transport of estimates through
the construction account for a substantial part of the formal
development. Its length alone does not measure mathematical novelty.

\subsubsection{Steps whose informal account is
incomplete}\label{steps-whose-informal-account-is-incomplete}

The theorem is proved in Lean, and the conditional statements of the
limit theory are discharged within its proof. The unit outlines,
however, describe the following steps only at the level of statements,
or with readings marked tentative. A written proof will have to supply
them.

\begin{itemize}
\tightlist
\item
  \textbf{One-loop identification (III.7, VII.2).} The slope bound is
  first proved for an auxiliary functional and then transferred to the
  completed Hessian of the actual step. The revised
  \href{https://github.com/scottnarmstrong/yang-mills-informalization/blob/main/outline/units/III-7.md\#7-what-remains-unexplained}{III.7}
  explains the determinant factorization and the transfer of the
  coefficient; several subsidiary real-quotient, complex-strip, tadpole,
  bubble and alias estimates remain inputs at statement level. The trace
  amplitude does not test this value.
\item
  \textbf{Covariance and Hessian (III.4).} The Gaussian covariance is
  the inverse of the gauge-fixed Gram form on the active links. Whether
  it equals the inverse Hessian of the effective action at the
  background is left open.
\item
  \textbf{Nonlinear step.}

  \begin{itemize}
  \tightlist
  \item
    The correspondence between the outputs of the one-step bounds
    (III.10) and the norms restored by the reset (III.9) is not written
    out.
  \item
    Several hypotheses of the step were read only as statements.
  \item
    The cluster expansion is proved for real weights (III.5), while the
    fixed-volume argument uses complex tuning.
  \end{itemize}
\item
  \textbf{Fixed-volume convergence (III.12, IV.1).}
  \href{https://github.com/scottnarmstrong/yang-mills-informalization/blob/main/outline/units/III-12.md\#7-what-remains-unexplained}{III.12}
  now explains fixed-level convergence, the uniform tail bound and the
  Cauchy argument; the underlying trajectory-convergence and multiscale
  estimates remain imported inputs.
  \href{https://github.com/scottnarmstrong/yang-mills-informalization/blob/main/outline/units/IV-1.md}{IV.1}
  locates the upgrade to the full spacing sequence, while leaving the
  upstream construction of its fixed-volume kernels outside that unit's
  review.
\item
  \textbf{Shell identities and the flux at infinity (VI.1, VI.2).} The
  lattice derivation of the shell identities is given only at statement
  level. The exterior-flux estimate was not read. The descent of the
  rotation action to exact intertwiners is marked tentative.
\item
  \textbf{Reconstruction (V.2).} The boundary values, the spectrum
  condition and the locality argument are described from statements;
  their mechanisms were confirmed in the Lean but not written out in the
  unit outlines.
\item
  \textbf{Contact extensions and nested OPE (VII.4).} The shell
  estimates and the nested rates are described from docstrings.
\item
  \textbf{Field relations (II.1).} It is not shown directly that the
  lattice equation relations exhaust the null space. The identification
  \(F_d\cong M_d\) comes from the comparison with free theory.
\end{itemize}

\subsubsection{Reading guide}\label{reading-guide}

The
\href{https://github.com/scottnarmstrong/yang-mills-informalization/blob/main/statement.md}{exact
statement} separates the public theorem from the stronger construction.
The
\href{https://github.com/scottnarmstrong/yang-mills-informalization/blob/main/outline/README.md}{architecture}
describes what each part hands to the next. For the stated inputs and
remaining gaps in this account, follow the relevant unit's ``What it
uses'' and ``What remains unexplained'' sections. The
\href{https://github.com/scottnarmstrong/yang-mills-informalization/blob/main/outline/units/VIII-1.md}{assembly
unit} describes how the construction is put together.

{\def\LTcaptype{none} % do not increment counter
\begin{longtable}[]{@{}
  >{\raggedright\arraybackslash}p{(\linewidth - 2\tabcolsep) * \real{0.5000}}
  >{\raggedright\arraybackslash}p{(\linewidth - 2\tabcolsep) * \real{0.5000}}@{}}
\toprule\noalign{}
\begin{minipage}[b]{\linewidth}\raggedright
Question
\end{minipage} & \begin{minipage}[b]{\linewidth}\raggedright
Units to read
\end{minipage} \\
\midrule\noalign{}
\endhead
\bottomrule\noalign{}
\endlastfoot
What are the fields, and how are their normalizations fixed? &
\href{https://github.com/scottnarmstrong/yang-mills-informalization/blob/main/outline/units/II-1.md}{II.1},
\href{https://github.com/scottnarmstrong/yang-mills-informalization/blob/main/outline/units/III-12.md}{III.12},
\href{https://github.com/scottnarmstrong/yang-mills-informalization/blob/main/outline/units/IV-1.md}{IV.1},
\href{https://github.com/scottnarmstrong/yang-mills-informalization/blob/main/outline/units/VII-1.md}{VII.1} \\
How does the ultraviolet analysis relate to Balaban's? &
\href{https://github.com/scottnarmstrong/yang-mills-informalization/blob/main/outline/units/III-1.md}{III.1},
\href{https://github.com/scottnarmstrong/yang-mills-informalization/blob/main/outline/units/III-3.md}{III.3},
\href{https://github.com/scottnarmstrong/yang-mills-informalization/blob/main/outline/units/III-4.md}{III.4},
\href{https://github.com/scottnarmstrong/yang-mills-informalization/blob/main/outline/units/III-6.md}{III.6},
\href{https://github.com/scottnarmstrong/yang-mills-informalization/blob/main/outline/units/III-5.md}{III.5},
\href{https://github.com/scottnarmstrong/yang-mills-informalization/blob/main/outline/units/III-10.md}{III.10} \\
Where does the one-loop coefficient enter? &
\href{https://github.com/scottnarmstrong/yang-mills-informalization/blob/main/outline/units/III-7.md}{III.7},
\href{https://github.com/scottnarmstrong/yang-mills-informalization/blob/main/outline/units/III-9.md}{III.9},
\href{https://github.com/scottnarmstrong/yang-mills-informalization/blob/main/outline/units/VII-5.md}{VII.5},
\href{https://github.com/scottnarmstrong/yang-mills-informalization/blob/main/outline/units/VII-2.md}{VII.2} \\
Why does one family serve every observable and every limit? &
\href{https://github.com/scottnarmstrong/yang-mills-informalization/blob/main/outline/units/III-2.md}{III.2},
\href{https://github.com/scottnarmstrong/yang-mills-informalization/blob/main/outline/units/II-1.md}{II.1},
\href{https://github.com/scottnarmstrong/yang-mills-informalization/blob/main/outline/units/III-12.md}{III.12},
\href{https://github.com/scottnarmstrong/yang-mills-informalization/blob/main/outline/units/VIII-1.md}{VIII.1};
\href{https://github.com/scottnarmstrong/yang-mills-informalization/blob/main/outline/units/III-13.md}{III.13}
treats an alternative family \\
How are bounds made uniform in the volume? &
\href{https://github.com/scottnarmstrong/yang-mills-informalization/blob/main/outline/units/III-11.md}{III.11},
\href{https://github.com/scottnarmstrong/yang-mills-informalization/blob/main/outline/units/III-12.md}{III.12},
\href{https://github.com/scottnarmstrong/yang-mills-informalization/blob/main/outline/units/III-16.md}{III.16} \\
How is infinite volume reached without an infrared RG? &
\href{https://github.com/scottnarmstrong/yang-mills-informalization/blob/main/outline/units/III-16.md}{III.16},
\href{https://github.com/scottnarmstrong/yang-mills-informalization/blob/main/outline/units/IV-1.md}{IV.1} \\
Why are there phases, and why does every phase reconstruct? &
\href{https://github.com/scottnarmstrong/yang-mills-informalization/blob/main/outline/units/V-1.md}{V.1},
\href{https://github.com/scottnarmstrong/yang-mills-informalization/blob/main/outline/units/V-2.md}{V.2},
\href{https://github.com/scottnarmstrong/yang-mills-informalization/blob/main/outline/units/V-3.md}{V.3};
\href{https://github.com/scottnarmstrong/yang-mills-informalization/blob/main/outline/units/VII-4.md}{VII.4}
for the train bounds \\
How are rotations and the stress tensor recovered? &
\href{https://github.com/scottnarmstrong/yang-mills-informalization/blob/main/outline/units/VI-2.md}{VI.2},
\href{https://github.com/scottnarmstrong/yang-mills-informalization/blob/main/outline/units/VI-1.md}{VI.1},
with
\href{https://github.com/scottnarmstrong/yang-mills-informalization/blob/main/outline/units/VII-1.md}{VII.1}
and
\href{https://github.com/scottnarmstrong/yang-mills-informalization/blob/main/outline/units/VII-6.md}{VII.6}
for the free comparison \\
What makes the theory nonfree? &
\href{https://github.com/scottnarmstrong/yang-mills-informalization/blob/main/outline/units/VII-2.md}{VII.2},
\href{https://github.com/scottnarmstrong/yang-mills-informalization/blob/main/outline/units/VII-3.md}{VII.3},
with
\href{https://github.com/scottnarmstrong/yang-mills-informalization/blob/main/outline/units/VII-6.md}{VII.6}
for arbitrary short-distance scales \\
What does the OPE say, and what is known at coincident points? &
\href{https://github.com/scottnarmstrong/yang-mills-informalization/blob/main/outline/units/VII-4.md}{VII.4},
then
\href{https://github.com/scottnarmstrong/yang-mills-informalization/blob/main/outline/units/VIII-1.md}{VIII.1}
for the assembly \\
\end{longtable}
}

\section{The architecture}\label{the-architecture}

This page describes how OpenAI's Lean proof of the construction of
four-dimensional Yang--Mills theory is organised: its eight parts, what
each part hands to the next, and the units that make up each part
(\href{https://github.com/scottnarmstrong/yang-mills-informalization/blob/main/outline/units.md}{detailed
map}).

It is read from the Lean statements and documentation and from the
dependency graph. The mathematical framing follows the Yang--Mills
group's
\href{https://github.com/yang-mills-group/YM/blob/main/docs/assessment-1.pdf}{assessment}.
Lean names link to the fixed source commit
\href{https://github.com/yang-mills-group/OpenAI-YM-20261028/tree/3456e5e75cf8d4cd558955a5c7638fba2d68f464}{\texttt{3456e5e}}.

\subsubsection{The order of limits}\label{the-order-of-limits}

The construction uses Wilson lattice theories on tori of side \(L_p\),
with lattice spacings \(a_N\) and couplings \(\beta_{N,p}\) satisfying
\(\beta_{N,p}/\log(1/a_N)\to 11C_G/(24\pi^2)\). The public theorem
leaves the positive limiting constant unspecified; see the
\href{https://github.com/scottnarmstrong/yang-mills-informalization/blob/main/statement.md}{exact
statement}.

Correlations of renormalized local observables are first taken to the
limit \textbf{\(N\to\infty\) at fixed volume}, along the whole sequence
of spacings. They are then taken to the limit \textbf{\(p\to\infty\)},
along a subsequence that may be extracted from any prescribed sequence
of volumes. The public form of this statement is
\href{https://github.com/yang-mills-group/OpenAI-YM-20261028/blob/3456e5e75cf8d4cd558955a5c7638fba2d68f464/PublicWilson.lean\#L85-L90}{\texttt{WilsonFamily.Converges}}.

The renormalization group is used only for the first limit. The second
is obtained by compactness.

\subsubsection{The eight parts}\label{the-eight-parts}

{\def\LTcaptype{none} % do not increment counter
\begin{longtable}[]{@{}
  >{\raggedright\arraybackslash}p{(\linewidth - 4\tabcolsep) * \real{0.3333}}
  >{\raggedright\arraybackslash}p{(\linewidth - 4\tabcolsep) * \real{0.3333}}
  >{\raggedright\arraybackslash}p{(\linewidth - 4\tabcolsep) * \real{0.3333}}@{}}
\toprule\noalign{}
\begin{minipage}[b]{\linewidth}\raggedright
part
\end{minipage} & \begin{minipage}[b]{\linewidth}\raggedright
name
\end{minipage} & \begin{minipage}[b]{\linewidth}\raggedright
units
\end{minipage} \\
\midrule\noalign{}
\endhead
\bottomrule\noalign{}
\endlastfoot
I & Setting &
\href{https://github.com/scottnarmstrong/yang-mills-informalization/blob/main/outline/units/I-1.md}{I.1} \\
II & Data of the construction &
\href{https://github.com/scottnarmstrong/yang-mills-informalization/blob/main/outline/units/II-1.md}{II.1},
\href{https://github.com/scottnarmstrong/yang-mills-informalization/blob/main/outline/units/II-2.md}{II.2} \\
III & Ultraviolet construction & III.1--III.16 \\
IV & Infinite volume &
\href{https://github.com/scottnarmstrong/yang-mills-informalization/blob/main/outline/units/IV-1.md}{IV.1} \\
V & Phases and reconstruction &
\href{https://github.com/scottnarmstrong/yang-mills-informalization/blob/main/outline/units/V-1.md}{V.1},
\href{https://github.com/scottnarmstrong/yang-mills-informalization/blob/main/outline/units/V-2.md}{V.2},
\href{https://github.com/scottnarmstrong/yang-mills-informalization/blob/main/outline/units/V-3.md}{V.3} \\
VI & Symmetry &
\href{https://github.com/scottnarmstrong/yang-mills-informalization/blob/main/outline/units/VI-1.md}{VI.1},
\href{https://github.com/scottnarmstrong/yang-mills-informalization/blob/main/outline/units/VI-2.md}{VI.2} \\
VII & Identification as Yang--Mills &
\href{https://github.com/scottnarmstrong/yang-mills-informalization/blob/main/outline/units/VII-1.md}{VII.1},
\href{https://github.com/scottnarmstrong/yang-mills-informalization/blob/main/outline/units/VII-2.md}{VII.2},
\href{https://github.com/scottnarmstrong/yang-mills-informalization/blob/main/outline/units/VII-3.md}{VII.3},
\href{https://github.com/scottnarmstrong/yang-mills-informalization/blob/main/outline/units/VII-4.md}{VII.4},
\href{https://github.com/scottnarmstrong/yang-mills-informalization/blob/main/outline/units/VII-5.md}{VII.5},
\href{https://github.com/scottnarmstrong/yang-mills-informalization/blob/main/outline/units/VII-6.md}{VII.6} \\
VIII & Assembly &
\href{https://github.com/scottnarmstrong/yang-mills-informalization/blob/main/outline/units/VIII-1.md}{VIII.1} \\
\end{longtable}
}

\paragraph{Ultraviolet units}\label{ultraviolet-units}

\begin{itemize}
\tightlist
\item
  Blocking, constants and Gaussian estimates:
  \href{https://github.com/scottnarmstrong/yang-mills-informalization/blob/main/outline/units/III-1.md}{III.1},
  \href{https://github.com/scottnarmstrong/yang-mills-informalization/blob/main/outline/units/III-2.md}{III.2},
  \href{https://github.com/scottnarmstrong/yang-mills-informalization/blob/main/outline/units/III-3.md}{III.3},
  \href{https://github.com/scottnarmstrong/yang-mills-informalization/blob/main/outline/units/III-4.md}{III.4}.
\item
  Cluster expansion, stationary field, matching and dependence:
  \href{https://github.com/scottnarmstrong/yang-mills-informalization/blob/main/outline/units/III-5.md}{III.5},
  \href{https://github.com/scottnarmstrong/yang-mills-informalization/blob/main/outline/units/III-6.md}{III.6},
  \href{https://github.com/scottnarmstrong/yang-mills-informalization/blob/main/outline/units/III-7.md}{III.7},
  \href{https://github.com/scottnarmstrong/yang-mills-informalization/blob/main/outline/units/III-8.md}{III.8}.
\item
  Coupling reset, nonlinear step and observable limits:
  \href{https://github.com/scottnarmstrong/yang-mills-informalization/blob/main/outline/units/III-9.md}{III.9},
  \href{https://github.com/scottnarmstrong/yang-mills-informalization/blob/main/outline/units/III-10.md}{III.10},
  \href{https://github.com/scottnarmstrong/yang-mills-informalization/blob/main/outline/units/III-11.md}{III.11},
  \href{https://github.com/scottnarmstrong/yang-mills-informalization/blob/main/outline/units/III-12.md}{III.12}.
\item
  Alternative coordinated family, construction data and analytic tools:
  \href{https://github.com/scottnarmstrong/yang-mills-informalization/blob/main/outline/units/III-13.md}{III.13},
  \href{https://github.com/scottnarmstrong/yang-mills-informalization/blob/main/outline/units/III-14.md}{III.14},
  \href{https://github.com/scottnarmstrong/yang-mills-informalization/blob/main/outline/units/III-15.md}{III.15},
  \href{https://github.com/scottnarmstrong/yang-mills-informalization/blob/main/outline/units/III-16.md}{III.16}.
\end{itemize}

The limit theory (Parts IV--VI) is written in the code around numbered
\href{https://github.com/scottnarmstrong/yang-mills-informalization/blob/main/outline/work-items.md}{work
items}; the units of those parts follow them.

\subsubsection{What each part hands to the
next}\label{what-each-part-hands-to-the-next}

\paragraph{Ultraviolet → infinite
volume}\label{ultraviolet-infinite-volume}

Let \(G\) be a compact matrix group with simple nonabelian Lie algebra,
and \(b>0\) a parameter of the block-spin kernel. The ultraviolet
construction provides a Wilson family with the following properties.

\begin{enumerate}
\def\labelenumi{\arabic{enumi}.}
\tightlist
\item
  \textbf{Specification.} The family satisfies its normalization,
  covariance and budget conditions, and its couplings obey the
  asymptotic-freedom law above.
\item
  \textbf{Chessboard bound, uniform in the volume.} For each degree
  \(d\) and \(\varepsilon>0\), every renormalized local observable at
  block scale \(m\) satisfies a chessboard bound of order
  \(a_m^{-d-\varepsilon}\), uniformly in the volume, once the lattice is
  fine enough. A chessboard bound is a bound on the root of a product of
  reflected copies.
\item
  \textbf{Differences.} The same holds for finite-difference quotients,
  so correlation kernels at \(r\)-separated points are bounded uniformly
  in the volume.
\item
  \textbf{Fixed volume.} At each fixed volume the correlation kernels
  converge along the whole sequence of spacings, to smooth limits.
\item
  \textbf{Corrected fields.} The same estimates hold for the corrected
  fields.
\end{enumerate}

For the subsequential extraction, the infinite-volume part uses 1--3 and
5, together with the reflection-positivity inequality on cells, which
holds for every \(\beta\ge0\). Its conclusion: for every increasing
sequence of volumes there is a subsequence along which all smeared
correlations converge to kernels that are smooth away from coincident
points. The fixed-volume convergence in item 4 then removes the spacing
refinement; see
\href{https://github.com/scottnarmstrong/yang-mills-informalization/blob/main/outline/units/IV-1.md}{IV.1}.
These two limits correspond to Theorems 4.1 and 5.1 of the group's
assessment.

A caution about the code. About three quarters of the code filed under
``full volume'' (folder
\href{https://github.com/yang-mills-group/OpenAI-YM-20261028/tree/3456e5e75cf8d4cd558955a5c7638fba2d68f464/YM/FullVolume}{\texttt{YM/FullVolume}})
belongs to the ultraviolet construction: it proves the volume-uniform
bounds. The genuinely infinite-volume argument is short, about 7,000
lines.

\paragraph{Infinite volume → phases, reconstruction and
symmetry}\label{infinite-volume-phases-reconstruction-and-symmetry}

The limit satisfies a Euclidean \emph{hierarchy}:

\begin{itemize}
\tightlist
\item
  translation homogeneity;
\item
  reflection positivity, including for smeared joint test functions;
\item
  smoothness on separated configurations, with controlled derivatives;
\item
  a compatible calculus of derivatives.
\end{itemize}

Reconstruction also needs a uniform bound on products of fields in thin
time slabs.

\paragraph{Reconstruction →
identification}\label{reconstruction-identification}

For every phase, a completed Wightman theory: unique vacuum, a
light-cone bound, heat-kernel decay, and a trivial weak commutant. The
heat bound is polynomial; no mass gap follows.

\paragraph{Symmetry → identification}\label{symmetry-identification}

A conserved stress tensor and its Ward identities. Both are compared,
along sequences of scales tending to zero, with the corresponding
objects for free Maxwell fields.

\paragraph{Into the assembly}\label{into-the-assembly}

A single Wilson family serves every part of the public theorem, from
extraction to the trace test;
\href{https://github.com/scottnarmstrong/yang-mills-informalization/blob/main/outline/units/VIII-1.md}{VIII.1}
describes the assembly. The construction run exists for every allowed
\(G\) and every \(b>0\).
\href{https://github.com/scottnarmstrong/yang-mills-informalization/blob/main/outline/units/III-13.md}{III.13}
describes an alternative coordinated family; its properties require a
separate identification before they can be used for this family.

\subsubsection{The work items of the limit
theory}\label{the-work-items-of-the-limit-theory}

Parts IV--VI are written in the code as about seventy numbered work
items. Each states one obligation, a ``proof request'' in the code's
vocabulary, and proves it in its own file.

{\def\LTcaptype{none} % do not increment counter
\begin{longtable}[]{@{}
  >{\raggedright\arraybackslash}p{(\linewidth - 6\tabcolsep) * \real{0.2500}}
  >{\raggedright\arraybackslash}p{(\linewidth - 6\tabcolsep) * \real{0.2500}}
  >{\raggedright\arraybackslash}p{(\linewidth - 6\tabcolsep) * \real{0.2500}}
  >{\raggedright\arraybackslash}p{(\linewidth - 6\tabcolsep) * \real{0.2500}}@{}}
\toprule\noalign{}
\begin{minipage}[b]{\linewidth}\raggedright
family
\end{minipage} & \begin{minipage}[b]{\linewidth}\raggedright
items
\end{minipage} & \begin{minipage}[b]{\linewidth}\raggedright
part
\end{minipage} & \begin{minipage}[b]{\linewidth}\raggedright
content
\end{minipage} \\
\midrule\noalign{}
\endhead
\bottomrule\noalign{}
\endlastfoot
D & 6 & IV & the limit passes to derivatives of the correlation kernels;
the Euclidean hierarchy \\
B & 18 & V & time-slab operators, decomposition into phases, reflection
positivity and smoothness for each phase \\
C & 29 & V & Osterwalder--Schrader → Wightman: continuation in time,
boost covariance, gluing of tubes, locality \\
A & 17 & VI & so(4) brackets, integration to the rotation action, the
averaged conserved stress tensor, boosts, comparison with free Maxwell
fields \\
\end{longtable}
}

These families interleave. D supplies the limit hierarchy; B begins with
its OS spaces and slab operators. A uses those spaces and the free
Maxwell comparison to obtain rotations and Ward identities. Those
symmetry results enter both the phase decomposition in B and the boost
and covariance arguments in C. The part numbers organize the subjects;
they do not give a linear dependency order. See
\href{https://github.com/scottnarmstrong/yang-mills-informalization/blob/main/summary.md\#the-strategy}{the
strategy} and the
\href{https://github.com/scottnarmstrong/yang-mills-informalization/blob/main/outline/work-items.md}{work-item
introductions}.

The same mechanism appears in the ultraviolet construction: families
E--I cover the energy chain of one renormalization step, its stationary
point and Hessian, label polynomials, and support bookkeeping.

See
\href{https://github.com/scottnarmstrong/yang-mills-informalization/blob/main/outline/work-items.md}{work-items.md}
for the complete list.

\subsubsection{What does not follow the order of the
parts}\label{what-does-not-follow-the-order-of-the-parts}

The Lean is not written in this order everywhere. Some definitions are
needed before the part that explains their properties:

\begin{itemize}
\tightlist
\item
  the reflection-positive cells and chessboard roots;
\item
  the record of the whole construction run;
\item
  the basic types of the renormalization group and of the labels.
\end{itemize}

Some sections are filed far from their mathematical home. Two examples:

\begin{itemize}
\tightlist
\item
  the rotation-invariance estimate, in folder \texttt{Next/ROTATIONS};
\item
  some renormalization-group pieces filed with the trace-anomaly code.
\end{itemize}

None of this involves new mathematics. A written account should simply
state these definitions early.

\subsubsection{How reliable is this}\label{how-reliable-is-this}

\begin{itemize}
\tightlist
\item
  Reachable declaration sets for seventeen named results were compared
  with OpenAI's compiled record. This comparison does not verify every
  edge or the mathematical interpretation of the interfaces
  (\href{https://github.com/scottnarmstrong/yang-mills-informalization/blob/main/method.md\#1-the-dependency-graph}{method}).
\item
  The plain-language interfaces summarize statements and documentation.
  The
  \href{https://github.com/scottnarmstrong/yang-mills-informalization/blob/main/outline/units.md}{unit
  index} gives each outline's review status; its local limitations
  describe which arguments remain at statement level. The
  \href{https://github.com/scottnarmstrong/yang-mills-informalization/blob/main/statement.md}{exact
  statement} distinguishes the public theorem from additional
  construction conclusions.
\item
  \href{https://github.com/scottnarmstrong/yang-mills-informalization/issues}{Report
  a correction} with the page, heading and source reference.
\end{itemize}

\section{The units}\label{the-units}

\subsection{I.1. Lattice, symmetries, gauge groups and the Wilson
measure}\label{i.1.-lattice-symmetries-gauge-groups-and-the-wilson-measure}

\textbf{Part:} I, Setting. \textbf{Lean:}
\href{https://github.com/yang-mills-group/OpenAI-YM-20261028/tree/3456e5e75cf8d4cd558955a5c7638fba2d68f464/YM/Lattice}{\texttt{YM/Lattice}}
(1,215 lines),
\href{https://github.com/yang-mills-group/OpenAI-YM-20261028/tree/3456e5e75cf8d4cd558955a5c7638fba2d68f464/YM/Gauge}{\texttt{YM/Gauge}}
(1,131),
\href{https://github.com/yang-mills-group/OpenAI-YM-20261028/tree/3456e5e75cf8d4cd558955a5c7638fba2d68f464/YM/Space}{\texttt{YM/Space}}
(430),
\href{https://github.com/yang-mills-group/OpenAI-YM-20261028/tree/3456e5e75cf8d4cd558955a5c7638fba2d68f464/YM/Corona}{\texttt{YM/Corona}}
(405),
\href{https://github.com/yang-mills-group/OpenAI-YM-20261028/tree/3456e5e75cf8d4cd558955a5c7638fba2d68f464/YM/Mink}{\texttt{YM/Mink}}
(363),
\href{https://github.com/yang-mills-group/OpenAI-YM-20261028/tree/3456e5e75cf8d4cd558955a5c7638fba2d68f464/YM/Wilson}{\texttt{YM/Wilson}}
(157). \textbf{Size:} 3,701 lines in 595 declarations used by the main
theorem.

\subsubsection{1. What this unit proves}\label{what-this-unit-proves}

Throughout, \(G\subset U(N)\) is a closed subgroup, \(N\ge1\),
\(\ell\ge1\) is integral, and
\(\Lambda_\ell=(\mathbb Z/\ell\mathbb Z)^4\). Configurations assign
\(U(x,i)\in G\) to links \(x\to x+e_i\); transport multiplies these
variables, using inverses on backward steps.

\textbf{Proposition I.1.1 (torus geometry).} Every site has a unique
lift \(r_\ell(x)\in(-\ell/2,\ell/2]^4\cap\mathbb Z^4\), and
\(|\Lambda_\ell|=\ell^4\). The metric
\(|x-y|_\ell=\max_i|r_\ell(x-y)_i|\) is translation and
signed-permutation invariant; \(|[y]_\ell|_\ell\le|y|_\infty\) for every
\(y\in\mathbb Z^4\). \emph{Status:} theorem. \emph{Lean:}
\href{https://github.com/yang-mills-group/OpenAI-YM-20261028/blob/3456e5e75cf8d4cd558955a5c7638fba2d68f464/YM/Lattice/Site.lean\#L61-L68}{\texttt{lift\_eq\_iff}},
\href{https://github.com/yang-mills-group/OpenAI-YM-20261028/blob/3456e5e75cf8d4cd558955a5c7638fba2d68f464/YM/Lattice/Site.lean\#L103-L105}{\texttt{separation\_triangle}}.

\textbf{Proposition I.1.2 (transport and symmetries).} On
\(\Lambda_\ell\) or \(\mathbb Z^4\), let \(A=(\rho,b)\),
\(Ax=\rho x+b\), with \(\rho=(\pi,\epsilon)\in B_4\) a signed
permutation, \((\rho x)_{\pi(i)}=\epsilon_ix_i\), \(b\in\Lambda\).
Pullback assigns image-link transport and satisfies \(1^*=\mathrm{id}\),
\(B^*A^*=(A\circ B)^*\) and \((A^*U)(x;w)=U(Ax;\rho w)\). For every
\(g:\Lambda\to G\), \(U^g(x,i)=g(x)U(x,i)g(x+e_i)^{-1}\) gives
\[U^g(x;w)=g(x)U(x;w)g(x_{\rm end})^{-1}.\] Transport depends only on
links with endpoints visited by \(w\). For \(p=(x;i<j)\) put
\(U_p=U(x;e_i,e_j,-e_i,-e_j)\). Affine maps injectively send plaquettes
to image squares \(A_\sharp p\), bijectively on the torus; every
conjugation- and inversion-invariant \(\phi:G\to M\) satisfies
\(\phi((A^*U)_p)=\phi(U_{A_\sharp p})\). Gauge invariance means
\(F(U^g)=F(U)\) for all \(g\); locality on \(X\) means dependence only
on links with both endpoints in \(X\). Periodic extension
\(E_bU(y,i)=U([y]_\ell+b,i)\) respects transport, gauge transformations
and locality, and \(E_b((\rho,a)^*U)=(\rho,0)^*E_{\rho b+a}U\).
\emph{Status:} theorem. \emph{Lean:}
\href{https://github.com/yang-mills-group/OpenAI-YM-20261028/blob/3456e5e75cf8d4cd558955a5c7638fba2d68f464/YM/Lattice/Path.lean\#L160-L167}{\texttt{travel\_pull}},
\href{https://github.com/yang-mills-group/OpenAI-YM-20261028/blob/3456e5e75cf8d4cd558955a5c7638fba2d68f464/YM/Lattice/Path.lean\#L141-L143}{\texttt{travel\_gauge}},
\href{https://github.com/yang-mills-group/OpenAI-YM-20261028/blob/3456e5e75cf8d4cd558955a5c7638fba2d68f464/YM/Lattice/Faces.lean\#L128-L143}{\texttt{Face.send\_price}},
\href{https://github.com/yang-mills-group/OpenAI-YM-20261028/blob/3456e5e75cf8d4cd558955a5c7638fba2d68f464/YM/Lattice/Path.lean\#L227-L235}{\texttt{expandAt\_pull}}.

\textbf{Proposition I.1.3 (spacings and blocks).} For
\(m,p\in\mathbb N\), odd integers \(L_m\ge3\) and strictly increasing
positive integers \(P_p\) divisible by \(4\), set \(D_0=1\),
\(D_{m+1}=D_mL_m\), \(a_m=D_m^{-1}\), \(\ell(m,p)=P_pD_m\),
\(s_m=\log(1/a_m)\). Then \(a_m\downarrow0\) strictly, \(s_m\to\infty\),
and \(a_m\ell(m,p)=P_p\); placements \(a_mr_{\ell(m,p)}\) lie in
\((-P_p/2,P_p/2]^4\) and are preserved by subdivisions. If
\(\log L_m/s_m\to0\) for \(m\ge1\), some \(c>0\) satisfies
\(a_{m+1}\ge ca_m^2\) for every \(m\).

For every integer \(L\ge1\), \(\ell'=\ell L\), the equivariant injective
homomorphism \(\iota x=[Lr_\ell(x)]_{\ell'}\) satisfies
\(r_{\ell'}\iota=Lr_\ell\). The bijection
\((x,y)\mapsto\iota x+[r_L(y)]_{\ell'}\) defines parent and residue
maps. For odd \(L\), parent is \(B_4\)-equivariant with cubic fibres of
radius \((L-1)/2\) centred at \(\iota x\). For odd \(L\),
\(\mathrm{par}(z+e_i)=\mathrm{par}(z)+e_i\) if
\(r_L(\mathrm{res}\,z)_i=(L-1)/2\), and equals \(\mathrm{par}(z)\)
otherwise. For all fine sites \(z,z'\),
\[|z-\iota\mathrm{par}\,z|_{\ell'}\le L/2,\qquad L|\mathrm{par}\,z-\mathrm{par}\,z'|_\ell\le|z-z'|_{\ell'}+L.\]
\emph{Status:} schedule definition, chosen in III.2; consequences
theorems. \emph{Lean:}
\href{https://github.com/yang-mills-group/OpenAI-YM-20261028/blob/3456e5e75cf8d4cd558955a5c7638fba2d68f464/YM/Lattice/Scale.lean\#L28-L34}{\texttt{Schedule}},
\href{https://github.com/yang-mills-group/OpenAI-YM-20261028/blob/3456e5e75cf8d4cd558955a5c7638fba2d68f464/YM/Lattice/Scale.lean\#L133-L162}{\texttt{mesh\_inequality}},
\href{https://github.com/yang-mills-group/OpenAI-YM-20261028/blob/3456e5e75cf8d4cd558955a5c7638fba2d68f464/YM/Lattice/Subdivision.lean\#L103-L111}{\texttt{Ratio.assemble\_bijective}},
\href{https://github.com/yang-mills-group/OpenAI-YM-20261028/blob/3456e5e75cf8d4cd558955a5c7638fba2d68f464/YM/Corona/Spend.lean\#L32-L36}{\texttt{Corona.foreshorten}}.

\textbf{Proposition I.1.4 (cover numbers).} Let
\(\mathrm{halo}_M(D)=\{y:\exists x\in D,\ |x-y|_\ell\le M\}\), and let
\(\mathrm{cov}_M(X)\) minimize \(|D|\) subject to
\(X\subset\mathrm{halo}_M(D)\) and halo connectivity by steps of
sup-distance at most \(1\); empty sets are allowed. For all sets and
\(M\in\mathbb N\), cover numbers are monotone and subadditive for
intersecting sets,
\[|X|\le(2M+1)^4\mathrm{cov}_M(X),\qquad |x-y|_\ell\le(4M+2)\mathrm{cov}_M(X)+2M\quad(x,y\in X).\]
A halo is connected exactly when its centres are chained by steps at
most \(2M+1\), and \(|\mathrm{halo}_M(D)|\le(2M+1)^4|D|\). For
\(b,k,m\in\mathbb N\), \(b\ge1\),
\(\mathrm{cov}_b(\mathrm{halo}_{kb}\{x\})\le(2k+1)^4\); a scale-\(mb\)
cover \(D\) gives \(\mathrm{cov}_b(X)\le(2m+3)^4|D|\). For subdivision
ratio \(L\ge1\), \(X\subset\Lambda_{\ell L}\), \(M,B,s\in\mathbb N\),
\(B\ge1\), \(M\le2B\),
\[L\mathrm{cov}_B(\mathrm{halo}_{sB}(\mathrm{par}\,X))\le(L+10\mathrm{cov}_M(X))(2s+13)^4.\]
\emph{Status:} theorem. \emph{Lean:}
\href{https://github.com/yang-mills-group/OpenAI-YM-20261028/blob/3456e5e75cf8d4cd558955a5c7638fba2d68f464/YM/Lattice/Support.lean\#L69-L71}{\texttt{coverNumber}},
\href{https://github.com/yang-mills-group/OpenAI-YM-20261028/blob/3456e5e75cf8d4cd558955a5c7638fba2d68f464/YM/Corona/Spend.lean\#L73-L105}{\texttt{Corona.collapse}}.

\textbf{Proposition I.1.5 (cost and Haar measure).} \(G\) is compact and
second countable. Put \(\sigma(u)=1-N^{-1}\mathrm{Re}\,\mathrm{tr}\,u\),
\(\chi(u)=\|u-1\|_F/\sqrt N\). Both are continuous, conjugation- and
inversion-invariant, vanish at \(1\), and satisfy \(\chi^2=2\sigma\),
\(0\le\sigma,\chi\le2\), \(\chi(uv)\le\chi(u)+\chi(v)\). Haar
probability \(\mu_G\) is invariant under inversion and both
translations. \emph{Status:} theorem. \emph{Lean:}
\href{https://github.com/yang-mills-group/OpenAI-YM-20261028/blob/3456e5e75cf8d4cd558955a5c7638fba2d68f464/YM/Gauge/MatrixGroup.lean\#L117-L126}{\texttt{chord\_price}},
\href{https://github.com/yang-mills-group/OpenAI-YM-20261028/blob/3456e5e75cf8d4cd558955a5c7638fba2d68f464/YM/Gauge/MatrixGroup.lean\#L193-L200}{\texttt{probHaar\_rightInv}}.

\textbf{Proposition I.1.6 (Lie algebra and Casimir).}
\(\mathfrak g=\{X\in M_N(\mathbb C):e^{tX}\in G\ \forall t\in\mathbb R\}\)
is a closed real Lie subalgebra of skew-Hermitian matrices, invariant
under conjugation. Every curve \(f:\mathbb R\to G\) differentiable at
\(0\) with \(f(0)=1\) has \(f'(0)\in\mathfrak g\). The invariant inner
product \(-N^{-1}\mathrm{tr}(XY)\) identifies \(\mathfrak g\) with
colours \(\mathbb R^d\), \(d=\dim\mathfrak g\), \(c\mapsto X_c\), with
\(\|X_c\|_F=\sqrt N|c|\). Projection satisfies
\(X_{\Pi A}=P_{\mathfrak g}A\), \(|\Pi A|\le\|A\|_F/\sqrt N\);
conjugation induces orthogonal \(\mathrm{Ad}\), and \(\mathrm{ad}(x)y\)
represents \([X_x,X_y]\).

If \(\mathfrak g\) is nonabelian with no nonzero proper ideals, every
symmetric real bilinear form satisfying
\(B(\mathrm{ad}(x)y,z)=-B(y,\mathrm{ad}(x)z)\) for all \(x,y,z\) is a
scalar multiple of the inner product. Writing
\(C_{\mathfrak g}=-d^{-1}\sum_i\mathrm{tr}(\mathrm{ad}(b_i)^2)\) in the
orthonormal colour basis, for all colours \(x,y\),
\[-\mathrm{tr}(\mathrm{ad}(x)\mathrm{ad}(y))=C_{\mathfrak g}\langle x,y\rangle,\qquad C_{\mathfrak g}>0,\qquad \kappa=11C_{\mathfrak g}/(24\pi^2)>0.\]
\emph{Status:} theorems; constants definitions. \emph{Lean:}
\href{https://github.com/yang-mills-group/OpenAI-YM-20261028/blob/3456e5e75cf8d4cd558955a5c7638fba2d68f464/YM/Gauge/Lie.lean\#L126-L135}{\texttt{lie}},
\href{https://github.com/yang-mills-group/OpenAI-YM-20261028/blob/3456e5e75cf8d4cd558955a5c7638fba2d68f464/YM/Gauge/Casimir.lean\#L99-L125}{\texttt{invariant\_scalar}},
\href{https://github.com/yang-mills-group/OpenAI-YM-20261028/blob/3456e5e75cf8d4cd558955a5c7638fba2d68f464/YM/Gauge/Casimir.lean\#L140-L158}{\texttt{trace\_adjoint}},
\href{https://github.com/yang-mills-group/OpenAI-YM-20261028/blob/3456e5e75cf8d4cd558955a5c7638fba2d68f464/YM/Gauge/Casimir.lean\#L160-L172}{\texttt{casimir\_pos}}.

\textbf{Proposition I.1.7 (small balls).} For every \(0<r\le1\),
\(\mu_G\{\chi\le r\}\ge(r/4)^{2N^2}\). \emph{Status:} theorem.
\emph{Lean:}
\href{https://github.com/yang-mills-group/OpenAI-YM-20261028/blob/3456e5e75cf8d4cd558955a5c7638fba2d68f464/YM/Gauge/Size.lean\#L80-L130}{\texttt{shell\_ge}}.

\textbf{Proposition I.1.8 (local group geometry).} There exist open
neighbourhoods \(V\ni1\), \(W\ni0\) and
\(\log_G:M_N(\mathbb C)\to\mathbb R^d\), real-analytic on \(V\), with
\(\log_G(1)=0\), \(\log_G(e^{X_c})=c\) for \(c\in W\) and
\(e^{X_{\log_Gu}}=u\) for \(u\in G\cap V\). Every matrix \(X\) has a
Frobenius-nearest point \(u\in G\), satisfying
\(u^*X-1\in\mathfrak g^\perp\). Nearest points respect left and right
translations by \(G\) and are unique in some tube about \(G\); near
\(1\) this retraction is analytic with strict derivative
\(P_{\mathfrak g}\). For every \(\varepsilon>0\), some \(\delta>0\)
gives \[\|u-v-P_{\mathfrak g}(X-Y)\|_F\le\varepsilon\|X-Y\|_F\] whenever
\(\|X-1\|_F,\|Y-1\|_F<\delta\) and \(u,v\) are nearest to \(X,Y\).
Curvature \(F(u)=\Pi((u-u^{-1})/2)\) is continuous, odd under inversion,
satisfies \(F(vuv^{-1})=\mathrm{Ad}(v)F(u)\), \(F(1)=0\),
\(|F(u)|\le\chi(u)\), and has derivative \(c\) along \(e^{tX_c}\) at
\(0\). \emph{Status:} theorem. \emph{Lean:}
\href{https://github.com/yang-mills-group/OpenAI-YM-20261028/blob/3456e5e75cf8d4cd558955a5c7638fba2d68f464/YM/Gauge/Chart.lean\#L164-L175}{\texttt{plog\_right}},
\href{https://github.com/yang-mills-group/OpenAI-YM-20261028/blob/3456e5e75cf8d4cd558955a5c7638fba2d68f464/YM/Gauge/Nearest.lean\#L172-L181}{\texttt{nearest\_approximate}},
\href{https://github.com/yang-mills-group/OpenAI-YM-20261028/blob/3456e5e75cf8d4cd558955a5c7638fba2d68f464/YM/Gauge/Curvature.lean\#L78-L96}{\texttt{leading\_curvature}}.

\textbf{Proposition I.1.9 (Wilson law).} For product Haar \(\mu_\ell\)
and plaquette holonomies \(U_p\), set
\[E_\ell=\sum_{x,i<j}\sigma(U_{(x;i<j)}),\qquad Z_\ell(\beta)=\int e^{-\beta E_\ell}d\mu_\ell,\qquad d\mathbb P_{\ell,\beta}=Z_\ell(\beta)^{-1}e^{-\beta E_\ell}d\mu_\ell.\]
For every real \(\beta\), \(E_\ell\) is continuous,
\(0\le E_\ell\le12\ell^4\), \(Z_\ell(\beta)>0\), and
\(\mathbb P_{\ell,\beta}\) is a probability measure integrating
continuous or bounded measurable complex functions. \emph{Status:}
definition with theorems. \emph{Lean:}
\href{https://github.com/yang-mills-group/OpenAI-YM-20261028/blob/3456e5e75cf8d4cd558955a5c7638fba2d68f464/YM/Wilson/Law.lean\#L87}{\texttt{law}},
\href{https://github.com/yang-mills-group/OpenAI-YM-20261028/blob/3456e5e75cf8d4cd558955a5c7638fba2d68f464/YM/Wilson/Law.lean\#L76-L80}{\texttt{partition\_pos}}.

\textbf{Proposition I.1.10 (invariance).} For every real \(\beta\),
affine lattice map \(A\) and \(g:\Lambda_\ell\to G\), pullback and gauge
transformation preserve \(\mu_\ell\), \(E_\ell\) and
\(\mathbb P_{\ell,\beta}\). \emph{Status:} theorem. \emph{Lean:}
\href{https://github.com/yang-mills-group/OpenAI-YM-20261028/blob/3456e5e75cf8d4cd558955a5c7638fba2d68f464/YM/Wilson/Symmetry.lean\#L93-L95}{\texttt{law\_pull}},
\href{https://github.com/yang-mills-group/OpenAI-YM-20261028/blob/3456e5e75cf8d4cd558955a5c7638fba2d68f464/YM/Wilson/Symmetry.lean\#L97-L99}{\texttt{law\_gauge}}.

\textbf{Proposition I.1.11 (continuum symmetries).} \(B_4\to O(4)\)
respects coordinate actions; skew-matrix exponentials lie in \(O(4)\).
With \(x=(\vec x,x_t)\), set \(\mathsf m=x_t^2-|\vec x|^2\),
\(\overline V_+=\{|\vec x|\le x_t\}\), \(V_+=\{|\vec x|<x_t\}\). The
topological group
\(\mathcal L=\{g\in GL(4,\mathbb R):\det g=1,\ \mathsf m\circ g=\mathsf m,\ g\overline V_+=\overline V_+\}\)
preserves Lebesgue measure and \(V_+\) and contains time-fixing elements
of \(SO(4)\). For \(R(\vec x,x_t)=(-\vec x,x_t)\),
\(g^\vee=RgR=(g^{-1})^T\) satisfies
\(\langle gx,g^\vee p\rangle=\langle x,p\rangle\). For every unit
spatial vector \(e\), the boost \(B_e(r)\) is a continuous one-parameter
subgroup of \(\mathcal L\), with generator
\((\vec x,x_t)\mapsto(x_te,\vec x\cdot e)\). The closed cone is
self-dual and equals \(\{x:(g^\vee x)_t\ge0\ \forall g\in\mathcal L\}\).
For \(x\in V_+\) and \(p\in\overline V_+\setminus\{0\}\),
\(\langle x,p\rangle>0\). The Poincaré group
\(\mathcal L\ltimes\mathbb R^4\), acting by \(x\mapsto\Lambda x+b\),
preserves Lebesgue measure diagonally on every \((\mathbb R^4)^n\).
Distinct-point configurations form an open set; pairwise distinct
coordinates on every axis, or positive pairwise sup-separation, imply
distinct points. \emph{Status:} definitions and theorems. \emph{Lean:}
\href{https://github.com/yang-mills-group/OpenAI-YM-20261028/blob/3456e5e75cf8d4cd558955a5c7638fba2d68f464/YM/Mink/Lorentz.lean\#L32-L44}{\texttt{Mink.Proper}},
\href{https://github.com/yang-mills-group/OpenAI-YM-20261028/blob/3456e5e75cf8d4cd558955a5c7638fba2d68f464/YM/Mink/Boost.lean\#L135-L158}{\texttt{Lorentz.from\_half}},
\href{https://github.com/yang-mills-group/OpenAI-YM-20261028/blob/3456e5e75cf8d4cd558955a5c7638fba2d68f464/YM/Mink/Poincare.lean\#L32-L34}{\texttt{Poincare}}.

\textbf{Proposition I.1.12 (cover boxes and routes).} For every
\(r\in\mathbb N\), \(B_r=\{-r,\ldots,r\}^4\) is \(B_4\)-invariant with
\((2r+1)^4\) points. For cover sites \(x,y\), axis routes from \(x\) to
\(y\) end at \(y\), give the reverse route from \(y\) to \(x\) when axis
order is reversed, and transform covariantly under affine maps. For
\(a\in\mathbb Z^4\) and an axis set \(J\), slabs
\(\{y\in B_r:y_i=a_i\ (i\notin J)\}\) decompose into pinned layers when
an axis \(i\notin J\) is freed; their sums decompose accordingly for
every function into a commutative additive monoid. Pinned layers
\(z,w\in\{-r,\ldots,r\}\) correspond by translation \((w-z)e_i\). For
every cover configuration and integer \(d\),
\(U(x;\mathrm{straight}_j(d+1))=U(x;\mathrm{straight}_j(d))U(x+de_j,j)\),
where \(\mathrm{straight}_j(d)\) has displacement \(de_j\).
\emph{Status:} theorem. \emph{Lean:}
\href{https://github.com/yang-mills-group/OpenAI-YM-20261028/blob/3456e5e75cf8d4cd558955a5c7638fba2d68f464/YM/Lattice/Box.lean\#L74-L80}{\texttt{route\_end}},
\href{https://github.com/yang-mills-group/OpenAI-YM-20261028/blob/3456e5e75cf8d4cd558955a5c7638fba2d68f464/YM/Lattice/Foliation.lean\#L62-L71}{\texttt{slab\_sum}}.

\subsubsection{2. What it uses}\label{what-it-uses}

\begin{itemize}
\tightlist
\item
  No inputs from other units.
\item
  Standard results: Haar uniqueness, product measures, matrix
  exponential, analytic inverse function theorem, spectral theorem,
  finite-dimensional volume scaling, Minkowski geometry.
\item
  Local radii depend on \(G\); the mesh constant depends on the
  schedule.
\end{itemize}

\subsubsection{3. Strategy}\label{strategy}

The spacing schedule separates physical periods from mesh sizes; centred
residues define the metric. Connected covers turn set size into path
length, allowing coarsening by sparse sampling. Image-link transport
reduces Wilson invariance to Haar invariance and conjugation/inversion
invariance of the cost.

The exponential definition of the Lie algebra supplies tangent
directions. Normal complements and the inverse function theorem give
local coordinates and nearest points; simplicity forces invariant
symmetric forms to be scalar.

\subsubsection{4. Steps}\label{steps}

\begin{enumerate}
\def\labelenumi{\arabic{enumi}.}
\tightlist
\item
  Centred division gives the metric and subdivision; odd residue windows
  give parent equivariance.
  \href{https://github.com/yang-mills-group/OpenAI-YM-20261028/blob/3456e5e75cf8d4cd558955a5c7638fba2d68f464/YM/Lattice/Subdivision.lean\#L154-L157}{\texttt{Ratio.parent\_signed}}.
\item
  Traverse a fine cover by a path, sample parents every \(LB\) steps,
  then tile the resulting coarse halo; this proves the coarsening bound.
  \href{https://github.com/yang-mills-group/OpenAI-YM-20261028/blob/3456e5e75cf8d4cd558955a5c7638fba2d68f464/YM/Corona/Spend.lean\#L73-L105}{\texttt{Corona.collapse}}.
\item
  Walk induction gives transport covariance; gauge factors telescope.
  Reversal and change of starting corner give plaquette inversion and
  conjugation.
  \href{https://github.com/yang-mills-group/OpenAI-YM-20261028/blob/3456e5e75cf8d4cd558955a5c7638fba2d68f464/YM/Lattice/Path.lean\#L141-L143}{\texttt{travel\_gauge}},
  \href{https://github.com/yang-mills-group/OpenAI-YM-20261028/blob/3456e5e75cf8d4cd558955a5c7638fba2d68f464/YM/Lattice/Faces.lean\#L128-L143}{\texttt{Face.send\_price}}.
\item
  For the curve of I.1.6, approximate \(e^X\) by \(f(1/n)^n\) to prove
  the tangent criterion. Isolation of normal exponentials gives the
  local logarithm.
  \href{https://github.com/yang-mills-group/OpenAI-YM-20261028/blob/3456e5e75cf8d4cd558955a5c7638fba2d68f464/YM/Gauge/Lie.lean\#L98-L111}{\texttt{tangent\_generates}},
  \href{https://github.com/yang-mills-group/OpenAI-YM-20261028/blob/3456e5e75cf8d4cd558955a5c7638fba2d68f464/YM/Gauge/Chart.lean\#L71-L91}{\texttt{isolates}}.
\item
  Minimize distance and use normality with local coordinates to identify
  nearest points. The symmetric operator associated with an invariant
  form has ideal eigenspaces; simplicity gives the Casimir identity.
  \href{https://github.com/yang-mills-group/OpenAI-YM-20261028/blob/3456e5e75cf8d4cd558955a5c7638fba2d68f464/YM/Gauge/Nearest.lean\#L128-L139}{\texttt{nearest\_is\_focus}},
  \href{https://github.com/yang-mills-group/OpenAI-YM-20261028/blob/3456e5e75cf8d4cd558955a5c7638fba2d68f464/YM/Gauge/Casimir.lean\#L99-L125}{\texttt{invariant\_scalar}}.
\item
  Pack ambient matrix balls for I.1.7; normalize the positive Wilson
  density and use link permutations, inversions and two-sided
  translations for invariance.
  \href{https://github.com/yang-mills-group/OpenAI-YM-20261028/blob/3456e5e75cf8d4cd558955a5c7638fba2d68f464/YM/Gauge/Size.lean\#L80-L130}{\texttt{shell\_ge}},
  \href{https://github.com/yang-mills-group/OpenAI-YM-20261028/blob/3456e5e75cf8d4cd558955a5c7638fba2d68f464/YM/Wilson/Symmetry.lean\#L75-L91}{\texttt{law\_preserved}}.
\item
  Boosts and Cauchy--Schwarz establish the cone characterizations.
  \href{https://github.com/yang-mills-group/OpenAI-YM-20261028/blob/3456e5e75cf8d4cd558955a5c7638fba2d68f464/YM/Mink/Geometry.lean\#L113-L129}{\texttt{cone\_test}},
  \href{https://github.com/yang-mills-group/OpenAI-YM-20261028/blob/3456e5e75cf8d4cd558955a5c7638fba2d68f464/YM/Mink/Boost.lean\#L135-L158}{\texttt{Lorentz.from\_half}}.
\end{enumerate}

\subsubsection{5. Where the difficulty
lies}\label{where-the-difficulty-lies}

Local group geometry must handle disconnected groups. Coarsening is
uniform in torus size. Real simplicity suffices without irreducibility
of the defining representation. Reflections shift link base points;
even-side lift ties prevent unrestricted placement covariance.

\subsubsection{6. Notation}\label{notation}

\([y]_\ell\): reduction; \(r_\ell\): centred lift; \(\mathrm{par}\):
block parent; \(\|\cdot\|_F\): Frobenius norm; \(P_{\mathfrak g}\): real
Frobenius orthogonal projection; \(\mathrm{ad}\): Lie bracket in colour
coordinates. Time is coordinate \(3\).

\subsubsection{7. What remains
unexplained}\label{what-remains-unexplained}

Helpers outside the six principal folders remain uncovered. The
condition \(M\le2B\) for the consecutive cover scales in III.5 and
III.10 remains to be traced. Exact bibliographic theorem locations
remain unchecked.

Detailed notes: the full account of this unit, with proof sketches for
every step, is on the
\href{https://github.com/scottnarmstrong/yang-mills-informalization/blob/main/outline/notes/I-1.md}{detailed
notes page}.

\subsection{II.1. Field labels and their lattice
representatives}\label{ii.1.-field-labels-and-their-lattice-representatives}

\textbf{Part:} II, Data of the construction. \textbf{Lean:} work items
G01--G07:
\href{https://github.com/yang-mills-group/OpenAI-YM-20261028/tree/3456e5e75cf8d4cd558955a5c7638fba2d68f464/YM/Labels}{\texttt{YM/Labels/}},
\href{https://github.com/yang-mills-group/OpenAI-YM-20261028/tree/3456e5e75cf8d4cd558955a5c7638fba2d68f464/YM/Source}{\texttt{YM/Source/}},
\href{https://github.com/yang-mills-group/OpenAI-YM-20261028/tree/3456e5e75cf8d4cd558955a5c7638fba2d68f464/YM/Work}{\texttt{YM/Work/}}
(G01--G07),
\href{https://github.com/yang-mills-group/OpenAI-YM-20261028/tree/3456e5e75cf8d4cd558955a5c7638fba2d68f464/YM/Next/FOUNDATIONS}{\texttt{YM/Next/FOUNDATIONS/}},
\href{https://github.com/yang-mills-group/OpenAI-YM-20261028/tree/3456e5e75cf8d4cd558955a5c7638fba2d68f464/YM/Push/RELAY_MARK}{\texttt{YM/Push/RELAY\_MARK/}},
\href{https://github.com/yang-mills-group/OpenAI-YM-20261028/tree/3456e5e75cf8d4cd558955a5c7638fba2d68f464/YM/Extra/SOURCE_HISTORIES}{\texttt{YM/Extra/SOURCE\_HISTORIES/}},
\href{https://github.com/yang-mills-group/OpenAI-YM-20261028/tree/3456e5e75cf8d4cd558955a5c7638fba2d68f464/YM/Push/LOCAL_SPLITTING}{\texttt{YM/Push/LOCAL\_SPLITTING/}},
\href{https://github.com/yang-mills-group/OpenAI-YM-20261028/tree/3456e5e75cf8d4cd558955a5c7638fba2d68f464/YM/Push/RELAY_DIFFERENCE}{\texttt{YM/Push/RELAY\_DIFFERENCE/}},
\href{https://github.com/yang-mills-group/OpenAI-YM-20261028/tree/3456e5e75cf8d4cd558955a5c7638fba2d68f464/YM/Planned/ATOM_GUARDED_LAYOUT}{\texttt{YM/Planned/ATOM\_GUARDED\_LAYOUT/}},
\href{https://github.com/yang-mills-group/OpenAI-YM-20261028/tree/3456e5e75cf8d4cd558955a5c7638fba2d68f464/YM/Further/LOCAL_LOG_JETS}{\texttt{YM/Further/LOCAL\_LOG\_JETS/}},
\href{https://github.com/yang-mills-group/OpenAI-YM-20261028/tree/3456e5e75cf8d4cd558955a5c7638fba2d68f464/YM/LocalH}{\texttt{YM/LocalH/}}.
\textbf{Size:} 10,805 lines in the main-theorem closure.

\subsubsection{1. What this unit proves}\label{what-this-unit-proves-1}

Throughout, \(G\subset U(n)\) is compact, \(n\ge1\).

\textbf{Proposition II.1.1 (labels and frames).} Let \(\mathrm{Tay}_j\)
consist of \(4\)-tuples of symmetric \((j+1)\)-linear maps
\((\mathbb R^4)^{j+1}\to\mathfrak g\) satisfying
\(\sum_\mu z_\mu A_\mu(z,\ldots,z)=0\). A jet is \(J=(J_j)_{j\ge0}\);
coordinates of \(J_j\) have weight \(j+2\). For every \(d\ge0\), the
space \(V_d\) of \(\mathrm{Ad}(G)\)-invariant complex jet polynomials of
weight \(\le d\) is finite-dimensional and determined by evaluation on
real jets. It splits by weight and factor number \((w,f)\); conjugation
and \(O(4)\) preserve grades. Its zero-factor part is \(\mathbb C1\).
There exist nested real pure-grade bases \((b_i)\) of grades
\((w_i,f_i)\) with \(b_0=1\). \emph{Status:} definitions and theorems.
\emph{Lean:}
\href{https://github.com/yang-mills-group/OpenAI-YM-20261028/blob/3456e5e75cf8d4cd558955a5c7638fba2d68f464/YM/Labels/Basis.lean\#L107-L111}{\texttt{frames}}.

\textbf{Proposition II.1.2 (holonomies and sections).} For every
\(d,J,h,x,\mu\), the equation \(u'=u\,hA^{(d)}_\mu(h(x+te_\mu))\),
\(u(0)=1\), has a unique solution in \(G\), where
\(A^{(d)}_\mu(z)=\sum_{j+2\le d}J_{j,\mu}(z,\ldots,z)\). Its time-one
endpoint \(U^{(d)}_{J,h}\) satisfies \(U_{J,h}=U_{\mathrm{sc}_hJ,1}\)
and \(U_{J,0}=1\). A section \((S,\Phi)\) maps coordinates linearly to
link-entry polynomials on finite \(S\). Soundness requires
signed-permutation-invariant \(S\ni0\), gauge invariance,
\(\Phi_{e_0}=1\), \(\Phi_{\bar c}=\overline{\Phi_c}\),
\(\Phi_c(s^*U)=\Phi_{s\cdot c}(U)\), and equality of the \(C^d\) jets of
\(\Phi_c(U_{J,h})\) and \(p_c(\mathrm{sc}_hJ)\) at \(0\), for all
\(c,J\). \emph{Status:} definitions; ODE theorems. \emph{Lean:}
\href{https://github.com/yang-mills-group/OpenAI-YM-20261028/blob/3456e5e75cf8d4cd558955a5c7638fba2d68f464/YM/Source/ODE.lean\#L162-L168}{\texttt{Hol.exists\_ode}};
\href{https://github.com/yang-mills-group/OpenAI-YM-20261028/blob/3456e5e75cf8d4cd558955a5c7638fba2d68f464/YM/Source/Local.lean\#L140-L147}{\texttt{SectionData.Sound}}.

\textbf{Proposition II.1.3 (primary sections).} For every holonomy datum
and coherent frames there are sound sections \(\Sigma_d\) whose column
\(i\) is a polynomial in finitely many transported plaquette curvatures,
every monomial having at least \(f_i\) factors. For every
\(k,P\in\mathbb N\) there is \(A\ge0\) such that, for every budget \(b\)
of order \(k\), power \(P\), side and placement, column \(i\) is
\(b\)-regular with norm \(\le At_b^{f_i}\). For the fixed primary
sections, nondecreasing \(\pi(D)\ge1\) bounds these costs for
\(d,k,P\le D\), and \(c_d\ge1\) bounds all column values globally.
\emph{Status:} theorem. \emph{Lean:}
\href{https://github.com/yang-mills-group/OpenAI-YM-20261028/blob/3456e5e75cf8d4cd558955a5c7638fba2d68f464/YM/Integrated.lean\#L216-L219}{\texttt{Integrated.source\_sections}}.

\textbf{Proposition II.1.4 (readout).} Fix holonomies, \(d\) and a
frame. Every \(b\)-regular \(f\) on \(X\subset\Lambda_\ell\), with
\(k_b\ge d\), has a unique read label \(p\in V_d\) at every \(o\): test
jets agree through order \(d\) if \(10\mathrm{dist}(x,o)<\ell\) for all
\(x\in X\); otherwise \(p=f(1)1\). If also \(P_b\ge d\) and
\(\lVert f\rVert_b\le c\), then
\(\lVert S_{t_b}\mathrm{coord}(p)\rVert_1\le Cc\), with \(C\ge1\)
depending only on holonomies and frame. A sound placed section reads
\(p_c\) if \(20(\lvert x_\mu\rvert+1)\le\ell\) throughout its stencil.
\emph{Status:} theorem. \emph{Lean:}
\href{https://github.com/yang-mills-group/OpenAI-YM-20261028/blob/3456e5e75cf8d4cd558955a5c7638fba2d68f464/YM/Next/FOUNDATIONS/Readout.lean\#L238-L247}{\texttt{exists\_read}};
\href{https://github.com/yang-mills-group/OpenAI-YM-20261028/blob/3456e5e75cf8d4cd558955a5c7638fba2d68f464/YM/Next/FOUNDATIONS/TestBounds.lean\#L271-L290}{\texttt{price\_of\_readout}}.

\textbf{Proposition II.1.5 (equation sections).} Let \(\mathcal E_d\) be
the labels vanishing on all jets of radial potentials whose curvatures
satisfy the free Bianchi and divergence equations, degree by degree. For
every holonomy datum, \(d\) and frame, there exist a sound section with
the slot structure and cost of II.1.3 and an automorphism \(R\)
commuting with conjugation and signed permutations, such that
\(\Phi_{\mathrm{coord}(Rp)}\) is micro for every \(p\in\mathcal E_d\).
Micro means continuous and, for some \(r,M\), orthogonal under every
Wilson expectation with \(\ell\ge M\), every placement, and
\(\beta\ne0\), to every bounded measurable spectator using only links
with both endpoints farther than \(r\) from the placement. The choices
exist simultaneously along coherent frames. \emph{Status:} unconditional
theorem. \emph{Lean:}
\href{https://github.com/yang-mills-group/OpenAI-YM-20261028/blob/3456e5e75cf8d4cd558955a5c7638fba2d68f464/YM/Push/RELAY_MARK/Portfolio.lean\#L101-L105}{\texttt{Portfolio.build}}.

\textbf{Proposition II.1.6 (observables).} A record carries a history,
holonomies, frames, sections, reference levels and transfer families.
Define
\[O_{N,p}(v,o)=a_N^{-d}\Phi^{(d)}_{S_{1/\theta_N}\tau^{-1}v}(o).\] For
all records this is linear, continuous, local and literally translation
covariant; soundness gives gauge invariance. Under \(\mathrm{Makes}\)
and \(r(d)\le N\), it respects reality and affine hypercubic pull-backs,
satisfies \(O_{N,p}(a_{r(d)}^de_0,o)=1\), and has reference-level
readout \(v\). Affine covariance is
\(O(v,A_{s,b}o;U)=O(s^{-1}\cdot v,o;A_{s,b}^*U)\). Here
\(\mathrm{Makes}\) requires selection, soundness, \(d+3\le J(r(d))\),
\(k_{b_m}\ge d+1\), \(P_{b_m}\ge d\) for \(m\ge r(d)\),
\(20(\lvert x_\mu\rvert+1)\le\ell_{N,p}\) on \(S_d\) when \(r(d)\le N\),
column bounds \(\kappa_N\theta_N^{f_i}\) and based transfers.
\emph{Status:} theorems; \(\mathrm{Makes}\) holds for VIII.1's family.
\emph{Lean:}
\href{https://github.com/yang-mills-group/OpenAI-YM-20261028/blob/3456e5e75cf8d4cd558955a5c7638fba2d68f464/YM/Source/Construction.lean\#L44-L62}{\texttt{Family.Makes}};
\href{https://github.com/yang-mills-group/OpenAI-YM-20261028/blob/3456e5e75cf8d4cd558955a5c7638fba2d68f464/YM/Source/Construction.lean\#L42}{\texttt{Family.O}}.

\textbf{Proposition II.1.7 (transfers).} With terminal linear
automorphism \(\tau\), a based family for \((\Sigma,r)\) requires
\(r\le N\) and, for every \(v\) and \(r\le m\le N\), linear coordinates
and remainders on fixed supports, regular remainders on \(X_i\ni0\) with
\(g_i(1)=0\) and vanishing test derivatives through \(d\) when
nonwrapping, and measurable gauge-invariant forced functions local in
their supports and flags, vanishing unless eligible and carrying \(0\),
and remainder bound \((\sqrt{w_m},\kappa_m\lVert x_m^\circ\rVert_1)\). A
bound \((\kappa,c)\) means summable regular and forced weighted masses
\(\le a\) and \(\le\kappa b\), respectively, for \(a,b\ge0\),
\(a+b\le c\), dominating regular norms and forced absolute values
pointwise. It requires \(x_N=v\), zero terminal regular remainder,
\(\mathcal D_N=\rho_N\Phi_{S_{1/\theta_N}v}\), and
\(\mathcal D_m=L_m^dT^{(m)}_{\beta_{N,p,m+1}}\mathcal D_{m+1}\). Also
\(\lVert x_{m+1}\rVert_1\le\kappa_{m+1}\lVert x_m\rVert_1\),
\(\lVert x_m\rVert_1\le\kappa_{m+1}L_m^d\lVert x_{m+1}\rVert_1\),
\(x_r=\tau v\), \(\tau\) respects reality/symmetry, and
\(x_m(a_N^de_0)=a_m^de_0\). For selected histories,
\(\int\mathcal D_m(E_m(\tau^{-1}v))=a_m^dZ_{N,p}\langle O_{N,p}(v,0)\rangle\).
\emph{Status:} specification, supplied by III.10; integral identity is a
theorem. \emph{Lean:}
\href{https://github.com/yang-mills-group/OpenAI-YM-20261028/blob/3456e5e75cf8d4cd558955a5c7638fba2d68f464/YM/Source/Insertion.lean\#L89-L114}{\texttt{TransferData.Based}}.

\textbf{Proposition II.1.8 (birth).} Let \(\rho\) be lawful on
\(\Lambda_\ell\) with \(w>0\), threshold \(\theta\),
\(o\in\Lambda_\ell\), \(Y=Y_o\), and \(\Sigma\) sound with
\(\lvert\Phi_{e_i}(o;\cdot)\rvert\le B\), \(B\ge1\). Suppose distinct
reserved zero-activity forced addresses exist for every finite
\(\varnothing\ne I\subseteq X\) and auxiliary index, with support \(X\),
inventory \(I\) and everywhere-satisfied auxiliary test. If
\[B\theta^{-f_{\max}}\mathrm{mass}(Y,1)(e^{\lvert Y\rvert w}-1)\le\sqrt w,\]
then, for every reservation index, there is linear \(E(v)\) with fixed
supports, \(x=v\), \(g=0\), Format and bound
\((\sqrt w,\lVert v^\circ\rVert_1)\), satisfying
\(\mathcal D(E(v))=\rho\Phi_{S_{1/\theta}v}\) pointwise. The forced part
is precisely the large-field supplement. The guarded variant
(\texttt{Guard.summon}, used at level N with w=w\_N) additionally
requires a selected history. It requires the auxiliary tests only when
the support has exactly the prescribed flagged inventory, together with
a separate zero-activity assumption. The section bound and fee condition
are unchanged. \emph{Status:} theorem. \emph{Lean:}
\href{https://github.com/yang-mills-group/OpenAI-YM-20261028/blob/3456e5e75cf8d4cd558955a5c7638fba2d68f464/YM/Extra/SOURCE_HISTORIES/Birth.lean\#L116-L161}{\texttt{launched}}.

\textbf{Proposition II.1.9 (numerical conditions).} Let \(C_j\ge1\), let
\((\mathsf K_j)\) be nondecreasing with \(\mathsf K_j\ge1\), and a stage
profile satisfy \(j\le k_j,P_j\), \(\pi(\max(k_j,P_j))\le\bar\kappa_j\),
\((C_j+1)\Omega_d\le\mathsf K_j\) for \(d\le j\), and
\(64\mathsf K_{j+1}^2\le\bar\kappa_j\). There exist \(n_j,\beta_j^*\)
such that every admissible level system of this profile with stage block
factors \(\ge n_j\) and \(\hat\beta_m\ge\beta^*_{J(m+1)}\) satisfies
\(J(m)\le k_{b_m},P_{b_m}\), \(\pi(\max(k_{b_m},P_{b_m}))\le\kappa_m\),
latch for \(d\le J(m)\), and
\[64\mathsf K_{J(m)+1}^2\le\kappa_m,\qquad10^5\mathsf K_{j+1}^6(1+\kappa_{m+1})^2\le L_m^{1/8},\]
\[\eta(1+\kappa_{m+1})\Omega_d\le L_m^{-3/8},\qquad(C_j+\eta\kappa_{m+1})\Omega_d\le\mathsf K_j,\]
where \(j=J(m+1)\), \(\eta=L_m^{-2/5}\), \(d\le J(m)\). Latch requires
sufficient side length and II.1.8's fee bound at every placement, with
\(B=c_d\), \(\theta=\theta_m\), \(w=w_m\). Such plans exist by enlarging
only stage constants whenever \(j\le k_j,P_j\). For fixed \(d\) and
support exponent, latch holds at sufficiently large coupling. These
conditions and \(d+3\le J(r(d))\) imply the section clauses of
\(\mathrm{Makes}\). \emph{Status:} theorems. \emph{Lean:}
\href{https://github.com/yang-mills-group/OpenAI-YM-20261028/blob/3456e5e75cf8d4cd558955a5c7638fba2d68f464/YM/Extra/SOURCE_HISTORIES/Reserve.lean\#L76-L102}{\texttt{Doors.afford}};
\href{https://github.com/yang-mills-group/OpenAI-YM-20261028/blob/3456e5e75cf8d4cd558955a5c7638fba2d68f464/YM/Extra/SOURCE_HISTORIES/Issue.lean\#L38-L45}{\texttt{Issue.refer}}.

\textbf{Proposition II.1.10 (comparison).} Diagonal sign changes,
translations and conjugation preserve centred root tests of interior
functions. If \(F_0\) is interior to a centred level-\(k\) cell with
root \(\le A\), and \(T\) has side divisible by that cell side and
satisfies the seam inequality, every translate at distance at least
three cell sides from its boundary has root \(\le A\). Under
\(\mathrm{Makes}\), fix \(\gamma>0\) and an extraction order. If
finitely many eventually literal-translate families have centred roots
\(\le e^{-\gamma s_N}\) eventually at arbitrarily large fixed levels
\(k\le N\), with interior central support, then combinations with
placement-independent coefficients eventually bounded by
\(e^{\varepsilon s_N}\) for every \(\varepsilon>0\) are close with bound
\(y=e^{-\gamma s_N/2}\to0\): for every \(j,m\in\mathbb N\) there is
\(k\ge m\) such that eventually \(k\le N\) and, for every level-\(k\)
link tiling \(T\) and
\(o\in T.\mathrm{margin}(0,\lfloor T.\mathrm{side}/(8(j+1))\rfloor)\),
the combination is interior to cell \(0\) with root \(\le y\).
\emph{Status:} theorems. \emph{Lean:}
\href{https://github.com/yang-mills-group/OpenAI-YM-20261028/blob/3456e5e75cf8d4cd558955a5c7638fba2d68f464/YM/Push/RELAY_DIFFERENCE/Placement.lean\#L232-L270}{\texttt{offset}};
\href{https://github.com/yang-mills-group/OpenAI-YM-20261028/blob/3456e5e75cf8d4cd558955a5c7638fba2d68f464/YM/Push/RELAY_DIFFERENCE/Combination.lean\#L65-L93}{\texttt{gather\_close}}.

\subsubsection{2. What it uses}\label{what-it-uses-1}

\begin{itemize}
\tightlist
\item
  I.1: Wilson/Haar measure and lattice symmetries.
\item
  II.2: regularity, norms, lawful polymers, tilings.
\item
  III.1: mass-preserving blocking.
\item
  III.2/III.9: profiles and selected histories.
\item
  III.10--III.11: transfer estimates.
\item
  III.16: Taylor calculus.
\item
  IV.1: fused extraction order.
\end{itemize}

Standard inputs: matrix ODEs, finite-group averaging, analytic inverse
function theorem, polarization, norm equivalence, Haar integration by
parts, reflection positivity.

\subsubsection{3. Strategy}\label{strategy-1}

Invert a finite sampling of transported plaquette curvatures to realize
jet labels. Readout reverses this construction. Link integration by
parts realizes equation labels; marking reserved polymers restores the
large-field complement of the insertion.

\subsubsection{4. Steps}\label{steps-1}

\begin{enumerate}
\def\labelenumi{\arabic{enumi}.}
\tightlist
\item
  Slots detect truncated jets: shrinking plaquette circulations recover
  curvature, hence radial potentials. \emph{Lean:}
  \href{https://github.com/yang-mills-group/OpenAI-YM-20261028/blob/3456e5e75cf8d4cd558955a5c7638fba2d68f464/YM/Work/G03.lean\#L212-L224}{\texttt{Work.G03.answer}}.
\item
  Taylor expansion of an equivariant inverse gives sections; each slot
  factor costs one threshold factor. \emph{Lean:}
  \href{https://github.com/yang-mills-group/OpenAI-YM-20261028/blob/3456e5e75cf8d4cd558955a5c7638fba2d68f464/YM/Work/G05.lean\#L113-L139}{\texttt{Work.G05.answer}};
  \href{https://github.com/yang-mills-group/OpenAI-YM-20261028/blob/3456e5e75cf8d4cd558955a5c7638fba2d68f464/YM/Work/G07.lean\#L79-L102}{\texttt{Work.G07.answer}}.
\item
  Rescaled test links, polarization and finite-dimensional norm
  equivalence give uniform readout. \emph{Lean:}
  \href{https://github.com/yang-mills-group/OpenAI-YM-20261028/blob/3456e5e75cf8d4cd558955a5c7638fba2d68f464/YM/Next/FOUNDATIONS/TestBounds.lean\#L234-L265}{\texttt{coeff\_bound}}.
\item
  Force lifts avoid the varied link; a unitriangular change absorbs
  higher-grade errors. \emph{Lean:}
  \href{https://github.com/yang-mills-group/OpenAI-YM-20261028/blob/3456e5e75cf8d4cd558955a5c7638fba2d68f464/YM/Push/RELAY_MARK/Portfolio.lean\#L60-L97}{\texttt{weaveBuild}}.
\item
  Splicing the large-field supplement into reserved addresses gives
  exact birth. \emph{Lean:}
  \href{https://github.com/yang-mills-group/OpenAI-YM-20261028/blob/3456e5e75cf8d4cd558955a5c7638fba2d68f464/YM/Extra/SOURCE_HISTORIES/Purchase.lean\#L42-L78}{\texttt{Attach.paid}}.
\item
  Seam log-convexity extends centred bounds to offsets. \emph{Lean:}
  \href{https://github.com/yang-mills-group/OpenAI-YM-20261028/blob/3456e5e75cf8d4cd558955a5c7638fba2d68f464/YM/Push/RELAY_DIFFERENCE/Placement.lean\#L232-L270}{\texttt{offset}}.
\end{enumerate}

\subsubsection{5. Where the difficulty
lies}\label{where-the-difficulty-lies-1}

Readout constants must avoid volume dependence. Equation lifts must
retain independence of the varied link. Lawfulness alone supplies no
reserved addresses. Higher grade orders factor number first, then
weight. Stencils are finite but nonexplicit.

\subsubsection{6. Notation}\label{notation-1}

\(\mathrm{sc}_hJ=(h^{j+2}J_j)_j\); \((S_tc)_i=t^{f_i}c_i\);
\(p_c=\sum_ic_ib_i\); \(c^\circ=c-c_0e_0\); \(Y=o+S\).
\(a_m,\theta_m,w_m,\kappa_m,J(m)\) are history parameters;
\(\Omega_d=1+2\sum_{s,i}\lVert s\cdot e_i\rVert_1\) sums over signed
permutations and frame indices. \(\mathcal D\) is the first-order
inserted density. Budget regularity means locality, measurability, gauge
invariance and \(C^{k_b}\) exponential-coordinate regularity on the
small-field set. The threshold satisfies \(t_b\in(0,1]\). Its norm
bounds all ordered left derivatives, including values, in directions
\(\lvert A_e\rvert\le t_b(1+\mathrm{dist}(e_-,o))^{P_b}\). Bounds
include the small-field boundary. A slot is
\(\mathrm{Ad}(U_\omega)F_P\), where \(F_P\) projects
\((U_P-U_P^{-1})/2\) onto \(\mathfrak g\). Eventually means volume
first, then mesh.

\subsubsection{7. What remains
unexplained}\label{what-remains-unexplained-1}

Identification of the auxiliary equation sections with limit null labels
remains tentative. Compatibility with all later transfer estimates
belongs to III.2 and III.10--III.15.

Detailed notes: the full account of this unit, with proof sketches for
every step, is on the
\href{https://github.com/scottnarmstrong/yang-mills-informalization/blob/main/outline/notes/II-1.md}{detailed
notes page}.

\subsection{II.2. Marker algebra, polymer combinatorics and Gaussian
estimates}\label{ii.2.-marker-algebra-polymer-combinatorics-and-gaussian-estimates}

\textbf{Part:} II, Data of the construction. \textbf{Lean:}
\href{https://github.com/yang-mills-group/OpenAI-YM-20261028/tree/3456e5e75cf8d4cd558955a5c7638fba2d68f464/YM/Marker}{\texttt{YM/Marker/}},
\href{https://github.com/yang-mills-group/OpenAI-YM-20261028/tree/3456e5e75cf8d4cd558955a5c7638fba2d68f464/YM/Ac}{\texttt{YM/Ac/}},
\href{https://github.com/yang-mills-group/OpenAI-YM-20261028/blob/3456e5e75cf8d4cd558955a5c7638fba2d68f464/YM/Ovens/Coffee.lean}{\texttt{YM/Ovens/Coffee.lean}},
\href{https://github.com/yang-mills-group/OpenAI-YM-20261028/blob/3456e5e75cf8d4cd558955a5c7638fba2d68f464/YM/Further/SHELL_BIRTH_49/Fusing/Coal.lean}{\texttt{YM/Further/SHELL\_BIRTH\_49/Fusing/Coal.lean}},
\href{https://github.com/yang-mills-group/OpenAI-YM-20261028/tree/3456e5e75cf8d4cd558955a5c7638fba2d68f464/YM/Extra/MULTISOURCE_TRAJECTORIES}{\texttt{YM/Extra/MULTISOURCE\_TRAJECTORIES/}},
\href{https://github.com/yang-mills-group/OpenAI-YM-20261028/blob/3456e5e75cf8d4cd558955a5c7638fba2d68f464/YM/Push/GRID_APPROXIMATION/Nil.lean}{\texttt{YM/Push/GRID\_APPROXIMATION/Nil.lean}},
\href{https://github.com/yang-mills-group/OpenAI-YM-20261028/tree/3456e5e75cf8d4cd558955a5c7638fba2d68f464/YM/MergeProof}{\texttt{YM/MergeProof/}},
\href{https://github.com/yang-mills-group/OpenAI-YM-20261028/tree/3456e5e75cf8d4cd558955a5c7638fba2d68f464/YM/Next/RANK}{\texttt{YM/Next/RANK/}},
\href{https://github.com/yang-mills-group/OpenAI-YM-20261028/blob/3456e5e75cf8d4cd558955a5c7638fba2d68f464/YM/Vac/Treasury.lean}{\texttt{YM/Vac/Treasury.lean}}.
\textbf{Size:} 3,691 lines in 793 declarations used by the main theorem.

\subsubsection{1. What this unit proves}\label{what-this-unit-proves-2}

\textbf{Proposition II.2.1 (marker algebra).} For finite \(V\) and a
downward-closed family \(R\subseteq2^V\) containing \(\emptyset\) and
singletons, \(\mathcal A_R=\mathbb C^R\) is a commutative unital
algebra: \(\delta_x\delta_y=\delta_{x\cup y}\) for disjoint \(x,y\) with
\(x\cup y\in R\), and zero otherwise; \(1=\delta_\emptyset\). Constant
coefficients multiply. If \(a(\emptyset)=0\), then \(a^k(x)=0\) for
\(|x|<k\), hence \(a^k=0\) for \(k>|V|\). \emph{Status:} theorem.
\emph{Lean:}
\href{https://github.com/yang-mills-group/OpenAI-YM-20261028/blob/3456e5e75cf8d4cd558955a5c7638fba2d68f464/YM/Marker/Console.lean}{ring
structure},
\href{https://github.com/yang-mills-group/OpenAI-YM-20261028/blob/3456e5e75cf8d4cd558955a5c7638fba2d68f464/YM/Marker/Truncated.lean\#L26-L27}{\texttt{Poly.stops}}.

\textbf{Proposition II.2.2 (norm).} For positive \(w_v\), set
\(w^x=\prod_{v\in x}w_v\), \(\|a\|_w=\sum_{x\in R}w^x|a(x)|\) and
\(\mathrm{Exp}(a)=\sum_{k=0}^{|V|}a^k/k!\). Then \(\|1\|_w=1\),
\(\|c\delta_x\|_w=w^x|c|\), \(\|ab\|_w\le\|a\|_w\|b\|_w\), and
\(\|\mathrm{Exp}(a)\|_w\le e^{\|a\|_w}\). A countable family with
\(\|f_i\|_w\le b_i\) and \(\sum_i b_i<\infty\) has a coefficientwise sum
of norm at most \(\sum_i b_i\). \emph{Status:} theorem. \emph{Lean:}
\href{https://github.com/yang-mills-group/OpenAI-YM-20261028/blob/3456e5e75cf8d4cd558955a5c7638fba2d68f464/YM/Marker/Weight.lean\#L139-L146}{\texttt{multiplication}},
\href{https://github.com/yang-mills-group/OpenAI-YM-20261028/blob/3456e5e75cf8d4cd558955a5c7638fba2d68f464/YM/Marker/Truncated.lean\#L51-L58}{\texttt{budgeted}}.

\textbf{Proposition II.2.3 (connected polynomial).} For zero-constant
\(a,b\), \(\mathrm{Exp}(a+b)=\mathrm{Exp}(a)\mathrm{Exp}(b)\) and
\(\mathrm{Exp}(-a)\mathrm{Exp}(a)=1\). For nonempty \(x\in R\),
\(\mathrm{Exp}(c\delta_x)=1+c\delta_x\). Every \(F(\emptyset)=1\) has a
unique zero-constant \(C\) with \(\mathrm{Exp}(C)=F\); \(C\) depends
continuously on \(F\). \emph{Status:} theorem. \emph{Lean:}
\href{https://github.com/yang-mills-group/OpenAI-YM-20261028/blob/3456e5e75cf8d4cd558955a5c7638fba2d68f464/YM/Next/RANK/Exponential.lean\#L56-L88}{\texttt{Exponential.pass\_add}},
\href{https://github.com/yang-mills-group/OpenAI-YM-20261028/blob/3456e5e75cf8d4cd558955a5c7638fba2d68f464/YM/Next/RANK/Exponential.lean\#L179-L180}{\texttt{connected\_spec}}.

\textbf{Proposition II.2.4 (tagging; repacking).} For finite \(s\) and
\(\tau:\iota\to R\), a nonzero tagged product requires pairwise disjoint
tags with union in \(R\); it equals
\((\prod_{i\in s}c_i)\delta_{\bigcup_i\tau(i)}\), and weights factor. If
\(\sigma:\iota\to\kappa\), admissibility is any property of finite slot
sets, and \(\sum_i\|F_i\|_w<\infty\), then
\[\sum_{\substack{s\text{ finite}:\ \sigma|_s\text{ injective}\\\sigma(s)\text{ admissible}}}\prod_{i\in s}F_i=\sum_{\substack{S\text{ finite}\\S\text{ admissible}}}\prod_{k\in S}\sum_{\sigma(i)=k}F_i,\]
with summable terms. \emph{Status:} theorem. \emph{Lean:}
\href{https://github.com/yang-mills-group/OpenAI-YM-20261028/blob/3456e5e75cf8d4cd558955a5c7638fba2d68f464/YM/Extra/MULTISOURCE_TRAJECTORIES/Algebra.lean\#L57-L69}{\texttt{tagged\_prod}},
\href{https://github.com/yang-mills-group/OpenAI-YM-20261028/blob/3456e5e75cf8d4cd558955a5c7638fba2d68f464/YM/Extra/MULTISOURCE_TRAJECTORIES/Packets.lean\#L114-L124}{\texttt{Packets.repack}}.

\textbf{Proposition II.2.5 (components).} For arbitrary finite supports
\(X_i\) and finite label set \(s\), overlap-chain classes partition
\(s\) into nonempty connected sets, called sparks, with disjoint
support-unions. Conversely, disjointly supported sparks recover exactly
these classes. Removing \(i\in s\) from a spark \(s\) leaves components
each meeting \(X_i\). \emph{Status:} theorem. \emph{Lean:}
\href{https://github.com/yang-mills-group/OpenAI-YM-20261028/blob/3456e5e75cf8d4cd558955a5c7638fba2d68f464/YM/Ac/Nodes.lean\#L79-L101}{\texttt{separates}},
\href{https://github.com/yang-mills-group/OpenAI-YM-20261028/blob/3456e5e75cf8d4cd558955a5c7638fba2d68f464/YM/Ac/Nodes.lean\#L103-L113}{\texttt{excise}}.

\textbf{Proposition II.2.6 (tree bound).} Let \(M\in\mathbb N\),
\(a,\gamma,f_i\ge0\) and \((2M+1)^4\gamma\le1\). If, at every site
\(x\), \(\sum_{i:x\in X_i}f_i e^{(a+1)\mathrm{cov}_M(X_i)}\) is summable
and at most \(\gamma\), then
\[\sum_{\substack{s\text{ spark}\\x\in X(s)}}e^{a\mathrm{cov}_M(X(s))}\prod_{i\in s}f_i\le e^\gamma-1,\]
and this family is summable; labels are not repeated. \emph{Status:}
theorem. \emph{Lean:}
\href{https://github.com/yang-mills-group/OpenAI-YM-20261028/blob/3456e5e75cf8d4cd558955a5c7638fba2d68f464/YM/Ac/Tree.lean\#L89-L127}{\texttt{Ac.forest}}.

\textbf{Proposition II.2.7 (two overlaps).} Give each label finite
mentions and walls, and mark some labels dirty. Within a selection, a
label is hot when mention-connected to dirt; its fuel comprises mentions
and, when hot, walls in disjoint copies of their ambient sets. Fuel
components retain their internal hotness. They bijectively decompose
finite selections into nonempty fuel-connected sets with disjoint
mention-unions and disjoint wall-unions among dirty-containing
components. If walls are lattice supports and mention overlap implies
wall overlap inside such a connected set, it is wall-connected.
\emph{Status:} theorem. \emph{Lean:}
\href{https://github.com/yang-mills-group/OpenAI-YM-20261028/blob/3456e5e75cf8d4cd558955a5c7638fba2d68f464/YM/Further/SHELL_BIRTH_49/Fusing/Coal.lean\#L210-L240}{\texttt{Kiln.undo}},
\href{https://github.com/yang-mills-group/OpenAI-YM-20261028/blob/3456e5e75cf8d4cd558955a5c7638fba2d68f464/YM/Further/SHELL_BIRTH_49/Fusing/Coal.lean\#L242-L253}{\texttt{Kiln.attends}}.

\textbf{Proposition II.2.8 (fusion).} Fix a gauge group \(G\),
four-dimensional periodic lattice \(\Lambda\), activity parameters and
regular/forced addresses (forced inventories nonempty and contained in
their supports, with extra tests) and coefficients \(E_i,K_i\), plus
logarithms \(\ell_i\) on nonempty supports \(Y_i\), indexed by
\(\mathbb N\). Patches are forced supports or \(Y_i\). Pellets are
finite patch-connected selections containing a forced label, with
disjoint forced supports; their reaches are \(K_i\) and
\(e^{\ell_i}-1\). Pellets admit a bijective enumeration. Fusion
preserves regular activities and assigns each pellet its union support,
union forced inventory, conjunction of forced eligibility tests, and
eligibility indicator times its reach-product. \emph{Status:}
definitions; enumeration theorem. The density, bound, covariance and
parametric continuity results belong to III.10. \emph{Lean:}
\href{https://github.com/yang-mills-group/OpenAI-YM-20261028/blob/3456e5e75cf8d4cd558955a5c7638fba2d68f464/YM/Ovens/Coffee.lean\#L78}{\texttt{fused}},
\href{https://github.com/yang-mills-group/OpenAI-YM-20261028/blob/3456e5e75cf8d4cd558955a5c7638fba2d68f464/YM/Ovens/Coffee.lean\#L121-L133}{\texttt{encoding}}.

The digest condition requires unfused validity with bounds \(\delta,b\);
logarithms supported on patches meeting the large-field set, uniformly
summably bounded, measurable, gauge invariant, local in internal links
and flags, and affine covariant. Fusion bounds require
\(\kappa,\gamma\ge0\), \((2B+1)^4\gamma\le1\), and a
configuration-independent \(f\ge0\) with
\(|\mathrm{reach}|\le e^{-\kappa}f\) on forced labels and
\(|\mathrm{reach}|\le f\) otherwise, satisfying II.2.6's hit condition
with cover size \(B\) and exponent \(e_P+1\). \emph{Lean:}
\href{https://github.com/yang-mills-group/OpenAI-YM-20261028/blob/3456e5e75cf8d4cd558955a5c7638fba2d68f464/YM/Ovens/Coffee.lean\#L81-L90}{\texttt{Digests}},
\href{https://github.com/yang-mills-group/OpenAI-YM-20261028/blob/3456e5e75cf8d4cd558955a5c7638fba2d68f464/YM/Ovens/Coffee.lean\#L94-L98}{\texttt{Buys}}.

\textbf{Proposition II.2.9 (switching).} For finite \(\iota\),
\(\mathrm{val}_x(g)(r)=\sum_{i\subseteq r}x^{|i|}g(i)\) is invertible
for \(x>0\). A model satisfies \(\mathrm{val}_x(g_x)=f_x\) on \([0,1]\)
and \(g_x\to g_0\); its jet is \(g_0\). A family is paid at cost \(d\)
when models share a bound \(|g_x^u(i)|\le b(i)\) with \(b\ge0\),
\(\sum_i b(i)\le d\). Costs add under finite sums and multiply under
finite products; summable costs permit infinite sums and commuting jets
with sums. If modeled \(f_x(r)=f_x(r\setminus\{v\})\) for every \(r\)
and \(x\in[0,1]\), then \(\mathrm{jet}(f)(i)=0\) whenever \(v\in i\);
pointwise measurability of paid families passes to jets. If
\(\|\rho_v\|_w\le1\), then
\(\sum_i\mathrm{jet}(f^u)(i)\prod_{v\in i}\rho_v\) has a common
coefficient majorant of weighted norm at most \(d\). \emph{Status:}
theorem. \emph{Lean:}
\href{https://github.com/yang-mills-group/OpenAI-YM-20261028/blob/3456e5e75cf8d4cd558955a5c7638fba2d68f464/YM/MergeProof/Switch.lean\#L33-L42}{\texttt{entire}},
\href{https://github.com/yang-mills-group/OpenAI-YM-20261028/blob/3456e5e75cf8d4cd558955a5c7638fba2d68f464/YM/MergeProof/Jetty.lean\#L13-L19}{\texttt{Accent.crystal}}.

\textbf{Proposition II.2.10 (Gaussian calculus).} For complex symmetric
\(P,Q\), define \(\mathcal G_Qp=p(X+\partial_Q)1\) and
\(\mathbb E_{Q,m}p=(\mathcal G_Qp)(m)\). These satisfy Wick recursion,
\(\mathcal G_P\mathcal G_Q=\mathcal G_{P+Q}\), and covariance transport
\(Q\mapsto AQA^T\) under substitutions with finitely supported rows. A
real random vector \(\xi\) under a probability law with all polynomial
moments and
\(\mathbb E[\xi_vq(\xi)]=\mathbb E[(\partial_{Q_{v\cdot}}q)(\xi)]\) for
every \(v,q\) has expectation \(\mathbb E[q(m+\xi)]=\mathbb E_{Q,m}q\)
for every real \(m\). If \(B\ge1\), \(|m_v|,|Q_{vw}|\le B\),
\(\deg p\le d\), then \(|\mathbb E_{Q,m}p|\le\|p\|_1\pi(B,d)\). Two such
pairs differing entrywise by at most \(E\ge0\) have expectation
difference at most \(\|p\|_1E\pi(B,d)\). \emph{Status:} theorem.
\emph{Lean:}
\href{https://github.com/yang-mills-group/OpenAI-YM-20261028/blob/3456e5e75cf8d4cd558955a5c7638fba2d68f464/YM/Next/RANK/Gaussian.lean\#L96-L102}{\texttt{warms}},
\href{https://github.com/yang-mills-group/OpenAI-YM-20261028/blob/3456e5e75cf8d4cd558955a5c7638fba2d68f464/YM/Next/RANK/ParticleMoment.lean\#L61-L78}{\texttt{payment}},
\href{https://github.com/yang-mills-group/OpenAI-YM-20261028/blob/3456e5e75cf8d4cd558955a5c7638fba2d68f464/YM/Next/RANK/Common.lean\#L124-L126}{\texttt{Common.differs}}.

\textbf{Proposition II.2.11 (polynomial steps).} Let \(Q\) be symmetric
and \(H_v\) linear. If \(C\ge1\), \(\|H_v\|_1,|Q_{vw}|\le C\),
\(\deg p\le d\), then \(T_{H,Q}p=(\mathcal G_Qp)\circ H\) has degree at
most \(d\) and norm at most \(\|p\|_1(2(d+2)C)^d\). Two such pairs
differing by at most \(\varepsilon\ge0\) in these norms have output
difference bounded by \(\|p\|_1\varepsilon(2(d+2)C)^d\). Zero
cross-covariance between variables of \(P,R\) gives \(T(PR)=T(P)T(R)\).
\emph{Status:} theorem. \emph{Lean:}
\href{https://github.com/yang-mills-group/OpenAI-YM-20261028/blob/3456e5e75cf8d4cd558955a5c7638fba2d68f464/YM/Next/RANK/Herald.lean\#L119-L124}{\texttt{crux}},
\href{https://github.com/yang-mills-group/OpenAI-YM-20261028/blob/3456e5e75cf8d4cd558955a5c7638fba2d68f464/YM/Next/RANK/Herald.lean\#L125-L126}{\texttt{crossing}},
\href{https://github.com/yang-mills-group/OpenAI-YM-20261028/blob/3456e5e75cf8d4cd558955a5c7638fba2d68f464/YM/Next/RANK/Crust.lean\#L59}{\texttt{bridge}}.

Weighted coefficient quantities \(\sum_\beta r^{w\cdot\beta}|c_\beta|\),
for \(w:V\to\mathbb N\), \(r\ge0\), are submultiplicative, and norms for
\(r>0\). If \(0\le\varepsilon\le1\), \(|Q_{vw}|\le\varepsilon\),
\(\deg p\le d\), then
\(\|\mathcal G_Qp-p\|_1\le\|p\|_1\varepsilon(2(d+2))^d\). For linear
\(f_v\) and polynomials \(A_v\) without terms below degree two,
\(\|f_v\|_1,\|A_v\|_1\le B\), \(B\ge1\), \(0<t\le1\), symmetric \(K\)
with \(0\le\varepsilon\le1\), \(|K_{vw}|\le\varepsilon\), and
\(\deg P\le n\),
\[\|\Phi_t(P)-(\mathcal G_KP)\circ f\|_1\le\|P\|_1\big(ntB(2B)^n+\varepsilon(2(n+2))^nB^n\big),\]
where \(\Phi_t:X_v\mapsto f_v+t^{-1}A_v(tX)\). A measurable
\(g\in[0,1]\) with \(\int(1-g)d\mu\le\eta\) is an \(\eta\)-guard. For
probability \(\mu\) and measurable \(f\) with
\(\int|f|^2d\mu\le B<\infty\), \(|\int gf\,d\mu-\int f\,d\mu|\) is at
most \(\varepsilon'B+\eta/\varepsilon'\) for every \(\varepsilon'>0\).
Under II.2.10's probability-law hypotheses and entry bounds,
\(\int|p(m+\xi)|^2\le\|p\|_1^2\pi(B,2d)\). For finite \(V\),
\(d\in\mathbb N\), \(\delta>0\), there exist finite nonempty
\(T\subseteq[-\delta,\delta]^V\) and \(C\ge0\) with
\(\|p\|_1\le C\max_{x\in T}|p(x)|\) for every \(\deg p\le d\).
\emph{Lean:}
\href{https://github.com/yang-mills-group/OpenAI-YM-20261028/blob/3456e5e75cf8d4cd558955a5c7638fba2d68f464/YM/Next/RANK/Pin.lean\#L66-L84}{\texttt{LiftMark.stock}},
\href{https://github.com/yang-mills-group/OpenAI-YM-20261028/blob/3456e5e75cf8d4cd558955a5c7638fba2d68f464/YM/Next/RANK/Guard.lean\#L108-L125}{\texttt{escaped}},
\href{https://github.com/yang-mills-group/OpenAI-YM-20261028/blob/3456e5e75cf8d4cd558955a5c7638fba2d68f464/YM/Next/RANK/SpectrumReal.lean\#L56-L81}{\texttt{Spindle.molds}}.

\textbf{Proposition II.2.12 (scales).} For \(s>e\), put
\(y=\log\log s\), \(\Pi=e^{y^2}\), \(\theta=s^{-1/2}\Pi\),
\(h=e^{y^{3/2}}\), \(j=h^{1/5}\), \(M=\lceil(\log s)^2\rceil\),
\(W=M^2\). For every fixed \(a,p>0\),
\[y^2\ll j\ll h\ll\Pi^a\ll s^p,\qquad W\sim(\log s)^4.\] For fixed
\(a>0\), \(k\in\mathbb N\), \(c\in\mathbb R\), eventually
\(W^k,h\le\Pi^a\), \(c\log W\le ay^2\), \(y(s)^2\le aW(s/2),aj(s/2)\)
and \(j(s)\le h(s/2)\). Eventually \(s\theta^2=\Pi^2\),
\(s\theta^3,s\theta^4\le1\), and \(2\Pi^v\theta^2\le1\) for every
\(v\le1\). \emph{Status:} theorem. \emph{Lean:}
\href{https://github.com/yang-mills-group/OpenAI-YM-20261028/blob/3456e5e75cf8d4cd558955a5c7638fba2d68f464/YM/Vac/Treasury.lean\#L83-L108}{\texttt{uniform\_loop}},
\href{https://github.com/yang-mills-group/OpenAI-YM-20261028/blob/3456e5e75cf8d4cd558955a5c7638fba2d68f464/YM/Next/RANK/Reign.lean\#L110-L122}{\texttt{rate\_fire}}.

\textbf{Proposition II.2.13 (numerical feasibility).} Fix
\(m,Q\in\mathbb N\), \(e,U,V\ge0\), \(a>0\), \(\alpha\ge1\), \(L>1\),
and \(G(u)=e^{-ah(u)^2}\). Eventually, uniformly for \(u,v\in[s/2,s]\),
\((2W(u)+1)^4\Pi(s)^{-1}\le1\), \(G(u)\le\Pi(s)^{-3/4}e^{-7j(v)/10}\)
and
\((\alpha+1)^{\lceil10\log M(s)/\log L\rceil}\Pi(s)^{-3/4}\le\Pi(s)^{-1/2}\).
For fixed \(k,n\in\mathbb N\), \(A\ge1\), \(\phi=1/(4(nk+k+1))\),
eventually
\[n\Pi^{-2}(A+\Pi^{-2})^n\pi(\Pi^\phi,nk)+A^n\Pi^{-2}\pi(\Pi^\phi,nk)\le\Pi^{-3/2}.\]
\emph{Status:} theorems and numerical definitions. \emph{Lean:}
\href{https://github.com/yang-mills-group/OpenAI-YM-20261028/blob/3456e5e75cf8d4cd558955a5c7638fba2d68f464/YM/Next/RANK/Stepped.lean\#L81-L89}{\texttt{on\_time}},
\href{https://github.com/yang-mills-group/OpenAI-YM-20261028/blob/3456e5e75cf8d4cd558955a5c7638fba2d68f464/YM/Next/RANK/Planck.lean\#L39-L62}{\texttt{collected}}.

For numerical records with \(t,g,H_0\in\mathbb N\),
\(b,c,f,F\in\mathbb R\), \(K\ge1\), \(c\ge1\), \(F\ge0\), \(f+c\le-8\),
\(g\ge6\), \(U=V\), \(f=1/16-U\), \(2V+2\le g-1\), put
\(C_d=(2(d+2)K)^d\). If \(b\ge0\), \(L>1\) and \(\log L\ge2(b+2)\log q\)
for \(q\in\{\alpha+2,C_t,C_{2t}\}\), then at
\(H=\lceil(b+2)y^2/\log L\rceil\) eventually
\((\alpha+2)^H,C_t^H,C_{2t}^H,H\le\Pi\) and
\(\Pi^{2c+1}+\Pi^{2c+2}\le\Pi^{2c+3}\). Consequently
\((\alpha+2)^r\Pi^{-g}\le\Pi^{-4}\) for \(0\le r\le H\). Such block-size
bounds hold for all sufficiently large \(L\); parameters precede \(L\),
which precedes \(s\). \emph{Lean:}
\href{https://github.com/yang-mills-group/OpenAI-YM-20261028/blob/3456e5e75cf8d4cd558955a5c7638fba2d68f464/YM/Next/RANK/Odometer.lean\#L56-L61}{\texttt{racetrack}}.

Put \(A_Q=(2Q+1)^4\) and \(R_Q(z)=e^{zA_Q}\). For
\(B,\varepsilon,w\ge0\), \(D\in\mathbb N\), define
\[q_0=w\pi(B,D),\quad d_0=wDB^D\varepsilon,\]
\[q_1=d_0+\varepsilon w^2\pi(B,2D)+\varepsilon+w\varepsilon\pi(B,D),\]
\[q_2=2q_0d_0+\varepsilon w^4\pi(B,4D)+2\varepsilon+w^2\varepsilon\pi(B,2D),\quad\mathcal Q=q_1+q_2+(q_0+q_1)q_1+q_1q_0.\]
For these records, if \(F\ge g+10Z+10\), \(Z=2c+8t+13\), eventually
\(4R_Q(e+2)\mathcal Q(\Pi,\Pi^{-F},\Pi^{2c+3},2t)\le\Pi^{-g}\).
\emph{Lean:}
\href{https://github.com/yang-mills-group/OpenAI-YM-20261028/blob/3456e5e75cf8d4cd558955a5c7638fba2d68f464/YM/Next/RANK/Purchased.lean\#L49-L56}{\texttt{funding}}.

\subsubsection{2. What it uses}\label{what-it-uses-2}

\begin{itemize}
\tightlist
\item
  I.1: lattice supports, connected cover numbers, cardinality and
  intersecting-union bounds.
\item
  III.1: activity format, eligibility, weighted bounds and gas.
\item
  Kotecký--Preiss induction; nilpotent exponentials; Boolean Möbius
  inversion; Gaussian integration by parts and Wick recursion;
  finite-dimensional norm equivalence.
\end{itemize}

\subsubsection{3. Strategy}\label{strategy-2}

Square-zero markers encode allowed products; their exponential separates
connected contributions. Connected-component decomposition and a
Kotecký--Preiss induction control polymer sums. Switching transfers
scalar coefficient bounds to markers. Algebraic Gaussian convolution
controls polynomial comparisons, and separated scales absorb their
fixed-degree losses.

\subsubsection{4. Steps}\label{steps-2}

\begin{enumerate}
\def\labelenumi{\arabic{enumi}.}
\item
  Construct the connected polynomial by
  \(C_{k+1}=C_k+F-\mathrm{Exp}(C_k)\): each correction removes another
  marker order.
  \href{https://github.com/yang-mills-group/OpenAI-YM-20261028/blob/3456e5e75cf8d4cd558955a5c7638fba2d68f464/YM/Next/RANK/Exponential.lean\#L133-L149}{\texttt{rolls}}.
\item
  Remove one root per component. The extra unit in the hit exponent pays
  each attachment; the cardinality bound closes induction.
  \href{https://github.com/yang-mills-group/OpenAI-YM-20261028/blob/3456e5e75cf8d4cd558955a5c7638fba2d68f464/YM/Ac/Table.lean\#L81-L106}{\texttt{lookup}},
  \href{https://github.com/yang-mills-group/OpenAI-YM-20261028/blob/3456e5e75cf8d4cd558955a5c7638fba2d68f464/YM/Ac/Tree.lean\#L89-L127}{\texttt{Ac.forest}}.
\item
  Restrict hotness to fuel components to recover the two-overlap
  decomposition.
  \href{https://github.com/yang-mills-group/OpenAI-YM-20261028/blob/3456e5e75cf8d4cd558955a5c7638fba2d68f464/YM/Further/SHELL_BIRTH_49/Fusing/Coal.lean\#L149-L162}{\texttt{hot\_comp}},
  \href{https://github.com/yang-mills-group/OpenAI-YM-20261028/blob/3456e5e75cf8d4cd558955a5c7638fba2d68f464/YM/Further/SHELL_BIRTH_49/Fusing/Coal.lean\#L210-L240}{\texttt{Kiln.undo}}.
\item
  Invert switching triangularly and pass convergent models to zero.
  \href{https://github.com/yang-mills-group/OpenAI-YM-20261028/blob/3456e5e75cf8d4cd558955a5c7638fba2d68f464/YM/MergeProof/Switch.lean\#L58-L63}{\texttt{Models.exact}}.
\item
  Apply Wick recursion inductively to moments and differences.
  \href{https://github.com/yang-mills-group/OpenAI-YM-20261028/blob/3456e5e75cf8d4cd558955a5c7638fba2d68f464/YM/Next/RANK/Common.lean\#L121-L123}{\texttt{width}},
  \href{https://github.com/yang-mills-group/OpenAI-YM-20261028/blob/3456e5e75cf8d4cd558955a5c7638fba2d68f464/YM/Next/RANK/Common.lean\#L124-L126}{\texttt{differs}}.
\item
  Compare logarithms of the scale functions and intersect the eventual
  inequalities, uniformly on \([s/2,s]\).
  \href{https://github.com/yang-mills-group/OpenAI-YM-20261028/blob/3456e5e75cf8d4cd558955a5c7638fba2d68f464/YM/Next/RANK/Bird.lean\#L145-L158}{\texttt{mint}}.
\end{enumerate}

\subsubsection{5. Where the difficulty
lies}\label{where-the-difficulty-lies-2}

Cover subadditivity requires intersecting unions. Hotness depends on the
selection. Switching needs convergence, beyond boundedness. Numerical
feasibility supplies no analytic trajectory estimates.

\subsubsection{6. Notation}\label{notation-2}

\(P=(r,q,e_P,B,t_P)\) has \(r,q,B\in\mathbb N\), \(r\ge6\), \(q\ge2\),
\(B\ge1\), \(e_P>0\), \(0<t_P\le1\); \(e_P\) is its support exponent.
Polynomials are complex; degrees are total and lie in \(\mathbb N\).
\(\partial_bX_v=b_v\). \(X(s)=\bigcup_{i\in s}X_i\); \(\mathrm{cov}_M\)
is the least number of radius-\(M\) balls whose union is connected and
contains the support. \(\|p\|_1\) sums absolute coefficients;
\(\pi(B,d)=(2B(d+1))^d\). \(f\ll g\) means \(f/g\to0\).

\subsubsection{7. What remains
unexplained}\label{what-remains-unexplained-2}

III.11 must construct restarted representations and verify analytic
hypotheses along trajectories. The auxiliary regrouping and
switching-cost declarations' full role remains to be traced.

Detailed notes: the full account of this unit, with proof sketches for
every step, is on the
\href{https://github.com/scottnarmstrong/yang-mills-informalization/blob/main/outline/notes/II-2.md}{detailed
notes page}.

\subsection{III.1. The block-spin step and the symmetric gauge
fixing}\label{iii.1.-the-block-spin-step-and-the-symmetric-gauge-fixing}

\textbf{Part:} III, Ultraviolet construction (renormalization group).
\textbf{Lean:}
\href{https://github.com/yang-mills-group/OpenAI-YM-20261028/tree/3456e5e75cf8d4cd558955a5c7638fba2d68f464/YM/RG}{\texttt{YM/RG/}}
(Block, Patch, Measurable, Push, Masks, Activities, PathActions,
Bounded, History),
\href{https://github.com/yang-mills-group/OpenAI-YM-20261028/tree/3456e5e75cf8d4cd558955a5c7638fba2d68f464/YM/Branch}{\texttt{YM/Branch/}},
\href{https://github.com/yang-mills-group/OpenAI-YM-20261028/tree/3456e5e75cf8d4cd558955a5c7638fba2d68f464/YM/Pressure}{\texttt{YM/Pressure/}},
\href{https://github.com/yang-mills-group/OpenAI-YM-20261028/tree/3456e5e75cf8d4cd558955a5c7638fba2d68f464/YM/PREC}{\texttt{YM/PREC/}},
\href{https://github.com/yang-mills-group/OpenAI-YM-20261028/tree/3456e5e75cf8d4cd558955a5c7638fba2d68f464/YM/PCalc}{\texttt{YM/PCalc/}},
\href{https://github.com/yang-mills-group/OpenAI-YM-20261028/tree/3456e5e75cf8d4cd558955a5c7638fba2d68f464/YM/Tree}{\texttt{YM/Tree/}},
\href{https://github.com/yang-mills-group/OpenAI-YM-20261028/tree/3456e5e75cf8d4cd558955a5c7638fba2d68f464/YM/Exterior}{\texttt{YM/Exterior/}},
and the small packages \texttt{YM/RG17/},
\texttt{YM/Push/CORRECTED\_COLUMNS/},
\texttt{YM/Planned/CALIBRATED\_MIXED\_HISTORY/},
\texttt{YM/UVNext/ANNULUS\_COMPARISON/}, \texttt{YM/FullVolume/Train/},
\texttt{YM/SectorTools/ACTUAL\_SCALAR\_LEDGER/}. \textbf{Size:} 4,219
lines (693 declarations) in the main-theorem closure.

\subsubsection{1. What this unit proves}\label{what-this-unit-proves-3}

Throughout, \(N\ge1\), \(G\subset U(N)\) is closed, \(c\ge1\), \(f=cL\),
\(L\ge3\) is odd, \(J\ge4\), \(R=(J+3)L\), and
\(K=(\kappa,\delta,\varepsilon,a)>0\). Branch conditions mean uniqueness
of the centre-normalized minimizing frame when every route holonomy in
every centre axial gauge has chord \(<\delta\) (B1), uniqueness of the
minimizing predictor when its minimum price is \(<\varepsilon\) (B2),
and nonempty eligible axial tags for every configuration (B3).

\textbf{Proposition III.1.1 (Markov kernel).} Under (B1)--(B3), for
every real \(\beta\), \(\rho_\beta\) is nonnegative, Borel, bounded and
satisfies \(\int\rho_\beta(U,V)d\mu_c(V)=1\) for every \(U\). For every
Banach-valued \(F\in L^1(\mu_f)\),
\(\mathcal P_\beta F(V)=\int\rho_\beta(U,V)F(U)d\mu_f(U)\) exists for
every \(V\), is integrable, depends linearly on \(F\), and preserves its
integral. Positivity and integrability in \(V\) need no branch
conditions. \emph{Status:} theorem. \emph{Lean:}
\href{https://github.com/yang-mills-group/OpenAI-YM-20261028/blob/3456e5e75cf8d4cd558955a5c7638fba2d68f464/YM/RG/Block.lean\#L203-L208}{\texttt{isKernel}}.

\textbf{Proposition III.1.2 (covariance).} Under (B1)--(B3), for every
\(\beta,U,V\), fine gauge \(k\), signed axis permutation \(s\) and
coarse site \(b\),

\[\rho_\beta(k\cdot U,(k\circ\mathrm{embed})\cdot V)=\rho_\beta(U,V)=\rho_\beta(s_{Lb}U,s_bV),\qquad\mathrm{embed}(x)=Lx.\]

A calendar satisfying branch conditions on every patch has block control
and Markov kernels at every step. \emph{Status:} theorem. \emph{Lean:}
\href{https://github.com/yang-mills-group/OpenAI-YM-20261028/blob/3456e5e75cf8d4cd558955a5c7638fba2d68f464/YM/Branch/Covariant.lean\#L56-L67}{\texttt{covary}};
\href{https://github.com/yang-mills-group/OpenAI-YM-20261028/blob/3456e5e75cf8d4cd558955a5c7638fba2d68f464/YM/Branch/Covariant.lean\#L50-L54}{\texttt{channel}}.

\textbf{Proposition III.1.3 (uniform branch constants).} Let \(d\) be as
in III.1.5 and \(r_G>0\) a unique-nearest-point tubular radius of \(G\).
For every \(\kappa>0\), \(0<\delta\le d/\sqrt N\),
\(0<\varepsilon\le r_G^2/(2N)\) and \(a\ge16\), (B1)--(B3) hold for
every patch. Every configuration and axis ordering admit an eligible
root. \emph{Status:} theorem. \emph{Lean:}
\href{https://github.com/yang-mills-group/OpenAI-YM-20261028/blob/3456e5e75cf8d4cd558955a5c7638fba2d68f464/YM/Branch/Pick.lean\#L137-L142}{\texttt{Axial.viable}};
\href{https://github.com/yang-mills-group/OpenAI-YM-20261028/blob/3456e5e75cf8d4cd558955a5c7638fba2d68f464/YM/Branch/Uniform.lean\#L18-L56}{\texttt{RG.frame\_unique}};
\href{https://github.com/yang-mills-group/OpenAI-YM-20261028/blob/3456e5e75cf8d4cd558955a5c7638fba2d68f464/YM/Branch/Uniform.lean\#L58-L63}{\texttt{RG.predictor\_unique}}.

\textbf{Proposition III.1.4 (axial gauge).} For every list \(l\) of
distinct axes, \(r\in\mathbb N\), \(o\in B_r\) and configuration \(U\)
on \(\mathbb Z^4\), put \(S=\{y\in B_r:y_j=o_j\ (j\notin l)\}\). Some
\(x\in S\) gives an axial gauge \(g_y=T_U(x\to y;l)\) with
\(E_S(g\cdot U)\le4|l|(2r+1)F_S(U)\). \emph{Status:} theorem.
\emph{Lean:}
\href{https://github.com/yang-mills-group/OpenAI-YM-20261028/blob/3456e5e75cf8d4cd558955a5c7638fba2d68f464/YM/Branch/Pick.lean\#L80-L105}{\texttt{Axial.existence}}.

\textbf{Proposition III.1.5 (synchronization uniqueness).} There exists
\(d(G)\in(0,1/4]\) such that, for all finite nonempty \(n,\Lambda\),
data \(W:n\times n\times\Lambda\to G\) with \(W_{xxk}=1\),
\(W_{yx\sigma(k)}=W_{xyk}^{-1}\) for a permutation \(\sigma\), and
\(\|W_{xyk}-1\|\le d\) have this property: any two global minimizers
over \(G^n\) of \(\mathrm{avg}_{x,y,k}\mathrm{pr}(h_xW_{xyk}h_y^{-1})\)
differ by one global left multiple. \emph{Status:} theorem. \emph{Lean:}
\href{https://github.com/yang-mills-group/OpenAI-YM-20261028/blob/3456e5e75cf8d4cd558955a5c7638fba2d68f464/YM/Branch/Synchronize.lean\#L232-L272}{\texttt{Branch.align}}.

\textbf{Proposition III.1.6 (frame comparison).} There exists \(A(G)>0\)
such that, for every patch, \(K\), configuration \(W\),
positive-frame-mass tag, realized frame \(h\), and gauge
\(H:\mathbb Z^4\to G\), some \(q\in G\) satisfies

\[\sum_{y\in B_R}\mathrm{ch}(q^{-1}h_yH_y^{-1})\le A(2R+1)[E_{B_R}(H\cdot W)+a(2R+1)F_{B_R}(W)].\]

No branch conditions are needed. \emph{Status:} theorem. \emph{Lean:}
\href{https://github.com/yang-mills-group/OpenAI-YM-20261028/blob/3456e5e75cf8d4cd558955a5c7638fba2d68f464/YM/Pressure/Gauge.lean\#L135-L162}{\texttt{Pressure.synchronized}}.

\textbf{Proposition III.1.7 (common gauge).} Fix \(G,K\) and \(j\ge4\);
set \(\ell=2j+15\), \(D=16\ell\), \(C=A\ell(D+a\ell)\) and
\(R'=(j+7)L\). For every step with \(J=j\), fine \(U\), coarse \(x\),
and tag usable at \(x\), there exist \(H:\mathbb Z^4\to G\) and
\(p_y\in G\) for \(y\in B_2\) such that, for every such \(y\),

\[E_{B_{R'}}(H\cdot U^{(x)})\le DL\,\mathrm{port}(x,U),\qquad\sum_{i\in B_R}\mathrm{ch}(p_y^{-1}h_i^{(x+y)}H_{i+Ly}^{-1})\le CL^2\mathrm{port}(x,U).\]

Every positive-tag-mass tag is usable everywhere. \emph{Status:}
theorem. \emph{Lean:}
\href{https://github.com/yang-mills-group/OpenAI-YM-20261028/blob/3456e5e75cf8d4cd558955a5c7638fba2d68f464/YM/Pressure/Cost.lean\#L72-L122}{\texttt{Pressure.planting}}.

\textbf{Proposition III.1.8 (coarse curvature).} With III.1.7's
constants and \(C_8=8\max(1,\varepsilon^{-1})(2C+D)\), for every step
with \(J=j\), configurations \(U,V\), coarse site \(x\), tag \(t\)
usable at \(x\), and plane \(p\),

\[\mathrm{ch}(Q(U,t)_{(x,p)})\le C_8L^{-2}\mathrm{port}(x,U),\qquad\mathrm{ch}(V_{(x,p)})\le\mathrm{fuel}(Q,V)_{(x,p)}+C_8L^{-2}\mathrm{port}(x,U).\]

\emph{Status:} theorem. \emph{Lean:}
\href{https://github.com/yang-mills-group/OpenAI-YM-20261028/blob/3456e5e75cf8d4cd558955a5c7638fba2d68f464/YM/Pressure/Cost.lean\#L210-L252}{\texttt{Pressure.qdc}}.

\textbf{Proposition III.1.9 (densities and histories).} Representations
evaluate to an exponential of vacuum minus action and masked regular
activities, times the gas. Flags mark sites touched by plaquettes of
chord above threshold. Reading links and flags in \(X\) requires only
links in \(X\) and corners of plaquettes meeting \(X\). The gas equals
\(1\) without flags. Every lawful representation (path bounds and valid
activities) has bounded Borel evaluation. Selected histories start at
the bare Wilson density and obey exact pushes; every selected history
satisfies \(\int\mathrm{eval}_{N,p,m}d\mu=Z_{N,p}(\beta_{N,p,N})\) for
all \(m\le N\). \emph{Status:} definitions, with four consequences
proved; quantitative construction belongs to III.3--III.10. \emph{Lean:}
\href{https://github.com/yang-mills-group/OpenAI-YM-20261028/blob/3456e5e75cf8d4cd558955a5c7638fba2d68f464/YM/RG/Masks.lean\#L108-L110}{\texttt{UsesFinal.allReads}};
\href{https://github.com/yang-mills-group/OpenAI-YM-20261028/blob/3456e5e75cf8d4cd558955a5c7638fba2d68f464/YM/RG/Activities.lean\#L135-L147}{\texttt{gas\_one}};
\href{https://github.com/yang-mills-group/OpenAI-YM-20261028/blob/3456e5e75cf8d4cd558955a5c7638fba2d68f464/YM/RG/History.lean\#L195-L198}{\texttt{Rep.envelope}};
\href{https://github.com/yang-mills-group/OpenAI-YM-20261028/blob/3456e5e75cf8d4cd558955a5c7638fba2d68f464/YM/RG/History.lean\#L242-L249}{\texttt{History.partition\_eq}}.

\textbf{Proposition III.1.10 (tube identity).} For every tube, path
cutoff, real thresholds and \(U,Q,V\), the six works satisfy

\[W_F+W_I+W_C+W_N+W_L+W_D=C_f+P-C_c,\]

where \(C_f,C_c\) are truncated crown sums and
\(P=\sum_b\kappa\,\mathrm{ch}(V_bQ_b^{-1})^2/2\). \emph{Status:}
theorem. \emph{Lean:}
\href{https://github.com/yang-mills-group/OpenAI-YM-20261028/blob/3456e5e75cf8d4cd558955a5c7638fba2d68f464/YM/Pressure/Rows.lean\#L145-L163}{\texttt{Tube.locates}}.

\textbf{Proposition III.1.11 (energy per seed).} Estimator data require:
for every sufficiently large \(L\) and then sufficiently large \(W\),
some \(\kappa_0>0\) gives \(\kappa_0t^2|A|\le\mathrm{Work}(A)\) at all
sufficiently large reference couplings, uniformly over compatible tubes,
forests and plots with trivial trunks and \(c\ge100(1+\varrho)\). For a
directed history satisfying ringing, each \(J_0\le i\le j\le i+1\)
admits \(\mu>0\) such that \(\mu t_{m+1}^2|A|\le\mathrm{Work}(A)\) for
all \(N,m<N,p\) with \((J_{m+1},J_m)=(j,i)\) and every plot at
\((\hat\beta_m,F(j,i))\) with trivial trunks. \emph{Status:}
requirement, proved in III.10; transfer theorem here. \emph{Lean:}
\href{https://github.com/yang-mills-group/OpenAI-YM-20261028/blob/3456e5e75cf8d4cd558955a5c7638fba2d68f464/YM/Pressure/Seeds.lean\#L51-L58}{\texttt{Stores.Plot}};
\href{https://github.com/yang-mills-group/OpenAI-YM-20261028/blob/3456e5e75cf8d4cd558955a5c7638fba2d68f464/YM/Pressure/Seeds.lean\#L129-L131}{\texttt{Teller.Grants}};
\href{https://github.com/yang-mills-group/OpenAI-YM-20261028/blob/3456e5e75cf8d4cd558955a5c7638fba2d68f464/YM/Pressure/Seeds.lean\#L166-L181}{\texttt{Testing.notation\_ticket}}.

\textbf{Proposition III.1.12 (kinetic specification).} Fix levels, a
quadratic horizon, Haar-chart density and positive stage constants.
Kinetic data assign shifts \(\alpha\), path data and free-link maps
\(\psi\). Stages and factors agree with the fixed choices; rows attach
to the horizon and satisfy stage bounds. At level \(N\), the unmasked,
unclipped action is analytic, matches Wilson derivatives through order
four, and has zero loop coefficients. For every \(m<N\), volume, forest
and tag, let \(\mathcal I_t\) be the fine unmasked, unclipped
unit-coupling action without loops plus predictor penalty, on fields
with trunks \(1\) and free links \(\exp(b_e)(\iota\exp u)_e\), where
\(\iota\) copies coarse links across block boundaries. Requirements, for
every free field \(v\), are

\[D_b^2\mathcal I_t(0,0)[v,v]\ge\eta\sum_{e\ {\rm free}}|v_e|^2,\qquad D_b\mathcal I_t(u,\psi(u))=0\]

near zero: \(\psi(0)=0\), both maps are analytic, \(D\psi(0)\) is the
prescribed background adjusted to tree gauge minus the linearized block
copy, and the analytic coarse action matches \(\mathcal I_t(u,\psi(u))\)
through order four. Writing \(H_t(u)=D_b^2\mathcal I_t(u,\psi(u))\) in
physical coordinates,

\[T_t=\mathcal L_f^\circ(U_F(u,\psi(u)))+\tfrac12\log\det H_t-\log\prod_{e\ {\rm free}}J(\psi(u)_e),\qquad D=\alpha\mathcal Q_c^\circ(\exp u)+\mathcal L_c^\circ(\exp u).\]

\(\mathcal L^\circ,\mathcal Q^\circ\) denote loop and quadratic actions;
\(J\) is the Haar-chart density. Both \(T_t,D\) are analytic,
\(\det H_t>0\) near zero, and

\[|D^2T_t(0)[a,b]-D^2D(0)[a,b]|\le|\Lambda_c|\Theta\Lambda e^{-\gamma c}\]

for real plane waves with direction-colour amplitudes bounded by \(1\);
\(\Lambda\) is the fine loop bound. Also
\(|\alpha+\kappa_{\rm slope}\log L_m|+1\le\mathrm{drift}_{m+1}\).
Directed histories require paid levels, a sufficient horizon, matching
recipe and local Haar-chart formula. Moves require exact push,
continuation, retention, lawful coarse data, forced-activity bound
\(\exp[-g_j\mathrm{monitor}_{m+1}^2]\), and
\(|\beta_m-\beta_{m+1}-\alpha|\le1\). \emph{Status:} definitions;
conditional constructors for simple nonabelian \(\mathfrak g\) belong to
III.2. \emph{Lean:}
\href{https://github.com/yang-mills-group/OpenAI-YM-20261028/blob/3456e5e75cf8d4cd558955a5c7638fba2d68f464/YM/PCalc/Loop.lean\#L93-L115}{\texttt{Kin.Planned}};
\href{https://github.com/yang-mills-group/OpenAI-YM-20261028/blob/3456e5e75cf8d4cd558955a5c7638fba2d68f464/YM/PCalc/Band.lean\#L49-L55}{\texttt{Directed.Move}}.

\textbf{Proposition III.1.13 (tree gauge).} For odd \(L\), block forests
exist. Inverse tree transport gives a gauge equal to \(1\) at roots,
depending only on trunks, and trivializing them. Every
real-Banach-valued integrable \(F\) invariant under gauges trivial at
roots satisfies
\(\int Fd\mu_f=\int F(\mathrm{drape}\,s)d\mu_{\rm free}(s)\). For a
kernel satisfying III.1.2(a) and gauge-invariant integrable \(F\), the
same reduction applies to \(\mathcal P_\beta F(V)\). \emph{Status:}
theorem. \emph{Lean:}
\href{https://github.com/yang-mills-group/OpenAI-YM-20261028/blob/3456e5e75cf8d4cd558955a5c7638fba2d68f464/YM/Tree/Cut.lean\#L81-L111}{\texttt{Forest.pack}};
\href{https://github.com/yang-mills-group/OpenAI-YM-20261028/blob/3456e5e75cf8d4cd558955a5c7638fba2d68f464/YM/Tree/Cut.lean\#L114-L120}{\texttt{Forest.pushing}}.

\textbf{Proposition III.1.14 (linearization).} Under (B1)--(B3), for
every tag \(t\) and coarse link \(b\), \(v\mapsto Q_b((\exp v_e)_e,t)\)
on fine \(\mathfrak g\)-valued \(1\)-forms equals \(1\) at \(v=0\) and
is real analytic nearby, using minimizing branches. Its tag-independent
derivative is \(\mathcal Qv=\mathcal Mv+d_0\mathcal Fv\), and
\(\mathcal Qd_0\varphi=d_0(\varphi\circ\mathrm{embed})\) for every fine
\(0\)-form. Frames have derivative
\(h_y^{(x)\prime}=\tau_x(y)-\tau_x(0)\). The kernel is analytic near
\((1,1)\) for fixed \(\beta\); neighbourhoods may depend on the step.
\emph{Status:} theorem. \emph{Lean:}
\href{https://github.com/yang-mills-group/OpenAI-YM-20261028/blob/3456e5e75cf8d4cd558955a5c7638fba2d68f464/YM/PREC/Derivative.lean\#L38-L68}{\texttt{Push.derivative}};
\href{https://github.com/yang-mills-group/OpenAI-YM-20261028/blob/3456e5e75cf8d4cd558955a5c7638fba2d68f464/YM/PREC/Derivative.lean\#L147-L160}{\texttt{Push.prices}}.

\subsubsection{2. What it uses}\label{what-it-uses-3}

\begin{itemize}
\tightlist
\item
  I.1: lattice geometry, Haar invariance, Wilson energy and
  nearest-point projection.
\item
  III.16: offset lengths; III.2: schedules and directed histories;
  III.10: energy estimates.
\item
  Tubular neighbourhood theorem, analytic implicit function theorem,
  Fubini, Jensen, continuity of unique compact minimizers.
\end{itemize}

\subsubsection{3. Strategy}\label{strategy-3}

Minimize the average route price to fix centre-normalized frames,
transport across blocks, and average the connectors. The predictor \(Q\)
is the nearest group point to their weighted matrix mean on its
minimizing branch. The kernel averages
\(\prod_b Z_\beta^{-1}e^{-\beta\kappa\mathrm{pr}(V_bQ_b^{-1})}\) over
tags, with \(Z_\beta=\int_Ge^{-\beta\kappa\mathrm{pr}(g)}dg\). Unique
minimizers handle small fields; symmetric sampling of
curvature-controlled axial frames and connectors handles the remaining
configurations. A common macro-box gauge compares neighbouring frames.

\subsubsection{4. Steps}\label{steps-3}

\begin{enumerate}
\def\labelenumi{\arabic{enumi}.}
\tightlist
\item
  \textbf{Uniform cutoffs.} Normalize minimizers by nearest centroid
  points; tangential cancellation makes feedback contract. Slab
  induction gives eligible axial roots.
  \href{https://github.com/yang-mills-group/OpenAI-YM-20261028/blob/3456e5e75cf8d4cd558955a5c7638fba2d68f464/YM/Branch/Synchronize.lean\#L232-L272}{\texttt{align}};
  \href{https://github.com/yang-mills-group/OpenAI-YM-20261028/blob/3456e5e75cf8d4cd558955a5c7638fba2d68f464/YM/Branch/Pick.lean\#L80-L105}{\texttt{existence}}.
\item
  \textbf{Kernel.} Tag probabilities and Haar noise each have mass one;
  uniqueness and tag permutations give covariance.
  \href{https://github.com/yang-mills-group/OpenAI-YM-20261028/blob/3456e5e75cf8d4cd558955a5c7638fba2d68f464/YM/RG/Block.lean\#L197-L201}{\texttt{density\_total}};
  \href{https://github.com/yang-mills-group/OpenAI-YM-20261028/blob/3456e5e75cf8d4cd558955a5c7638fba2d68f464/YM/Branch/Covariant.lean\#L50-L54}{\texttt{channel}}.
\item
  \textbf{Curvature.} Synchronization costs \(L^2\), while
  \(w\le L^{-4}\); four padded links give the plaquette bound.
  \href{https://github.com/yang-mills-group/OpenAI-YM-20261028/blob/3456e5e75cf8d4cd558955a5c7638fba2d68f464/YM/Pressure/Cost.lean\#L129-L179}{\texttt{indicated}};
  \href{https://github.com/yang-mills-group/OpenAI-YM-20261028/blob/3456e5e75cf8d4cd558955a5c7638fba2d68f464/YM/Pressure/Cost.lean\#L210-L252}{\texttt{qdc}}.
\item
  \textbf{Local analysis.} The implicit function theorem gives
  derivatives; Haar invariance removes trunks. Adjointness cancels cross
  terms and squares in the tube identity.
  \href{https://github.com/yang-mills-group/OpenAI-YM-20261028/blob/3456e5e75cf8d4cd558955a5c7638fba2d68f464/YM/PREC/Pred.lean\#L129-L187}{\texttt{scheme}};
  \href{https://github.com/yang-mills-group/OpenAI-YM-20261028/blob/3456e5e75cf8d4cd558955a5c7638fba2d68f464/YM/Tree/Cut.lean\#L81-L111}{\texttt{pack}};
  \href{https://github.com/yang-mills-group/OpenAI-YM-20261028/blob/3456e5e75cf8d4cd558955a5c7638fba2d68f464/YM/Pressure/Rows.lean\#L145-L163}{\texttt{Tube.locates}}.
\end{enumerate}

\subsubsection{5. Where the difficulty
lies}\label{where-the-difficulty-lies-3}

Uniqueness must remain uniform in patch size. Both branches must
preserve area scaling. Global regularity is Borel; analyticity is local.

\subsubsection{6. Notation}\label{notation-3}

\(B_r=\{-r,\ldots,r\}^4\); \(\mathrm{ch}(u)=\|u-1\|/\sqrt N\),
\(\mathrm{pr}(u)=1-\mathrm{Re\,tr}(u)/N\); \(\|\cdot\|\) is Frobenius,
\(|v|=\|v\|/\sqrt N\). \(\mu\) denotes product probability Haar measure;
\(E_X,F_X\) sum link and plaquette chords. \(U^{(x)}\) is the
root-centred periodic lift; \(\mathrm{port}=F_{B_{R'}}(U^{(x)})\).
Usable means positive frame masses on \(x+B_2\); fuel sums four link
discrepancies with multiplicity. \(w\) is the \(J\)-fold block-uniform
convolution. \(\tau\) averages negative route integrals, \(\mathcal M\)
averages straight length-\(L\) integrals,
\(\mathcal Fv(x)=\tau_x(0)-\sum_yw_y\tau_x(y)\), and \(d_0\) is forward
difference. Here
\(\tau_x(y)=-\mathrm{avg}_{z\in B_R,o\in S_4}\int_{Lx+y\to Lx+z;o}v\).

\(\mathrm{monitor}_m=e^{(\log\log\hat\beta_m)^2}\);
\(t_m=\mathrm{temp}(\hat\beta_m)\),
\(\mathrm{temp}(z)=z^{-1/2}e^{(\log\log z)^2}\);
\(t=\mathrm{temp}(\hat\beta+\kappa_{\rm slope}\log L)\),
\(\varrho=\lceil(\log\hat\beta)^2\rceil\),
\(\kappa_{\rm slope}=11C_A/(24\pi^2)\).

A store \(A\) comprises fine plaquettes, coarse links and coarse sites;
\(|A|\) sums their cardinalities. \(\mathrm{Work}(A)\) sums local works
whose fine bases' parents or coarse bases lie within distance
\(2\varrho\) of its locations, with cutoff
\(\lfloor\varrho/10^5\rfloor\). \(F(j,i)>0\) is the assigned wire
parameter.

\subsubsection{7. What remains
unexplained}\label{what-remains-unexplained-3}

Energy lower bounds and kinetic construction belong to later units. The
uses of \(B_2\) rather than \(B_1\), and \(J\ge4\) rather than
\(J\ge1\), remain unspecified here.

Detailed notes: the full account of this unit, with proof sketches for
every step, is on the
\href{https://github.com/scottnarmstrong/yang-mills-informalization/blob/main/outline/notes/III-1.md}{detailed
notes page}.

\subsection{III.2. Choosing the constants: the order of choice, the
level schedule and the shooting
argument}\label{iii.2.-choosing-the-constants-the-order-of-choice-the-level-schedule-and-the-shooting-argument}

\textbf{Part:} III, Ultraviolet construction (renormalization group).
\textbf{Lean:}
\href{https://github.com/yang-mills-group/OpenAI-YM-20261028/tree/3456e5e75cf8d4cd558955a5c7638fba2d68f464/YM/Price}{\texttt{YM/Price/}}
(the constants and the level schedule, 406 lines),
\href{https://github.com/yang-mills-group/OpenAI-YM-20261028/tree/3456e5e75cf8d4cd558955a5c7638fba2d68f464/YM/Shoot}{\texttt{YM/Shoot/}}
(shooting, 228 lines),
\href{https://github.com/yang-mills-group/OpenAI-YM-20261028/tree/3456e5e75cf8d4cd558955a5c7638fba2d68f464/YM/Integration/RG05Preconstruction}{\texttt{YM/Integration/RG05Preconstruction/}},
\href{https://github.com/yang-mills-group/OpenAI-YM-20261028/tree/3456e5e75cf8d4cd558955a5c7638fba2d68f464/YM/Integration/NativeSelection}{\texttt{YM/Integration/NativeSelection/}}
and other files of
\href{https://github.com/yang-mills-group/OpenAI-YM-20261028/tree/3456e5e75cf8d4cd558955a5c7638fba2d68f464/YM/Integration}{\texttt{YM/Integration/}},
\href{https://github.com/yang-mills-group/OpenAI-YM-20261028/tree/3456e5e75cf8d4cd558955a5c7638fba2d68f464/YM/RG08/Calendar}{\texttt{YM/RG08/Calendar/}},
and parts of
\href{https://github.com/yang-mills-group/OpenAI-YM-20261028/tree/3456e5e75cf8d4cd558955a5c7638fba2d68f464/YM/RG08}{\texttt{YM/RG08/}},
\href{https://github.com/yang-mills-group/OpenAI-YM-20261028/tree/3456e5e75cf8d4cd558955a5c7638fba2d68f464/YM/RG12}{\texttt{YM/RG12/}},
\href{https://github.com/yang-mills-group/OpenAI-YM-20261028/tree/3456e5e75cf8d4cd558955a5c7638fba2d68f464/YM/FullVolume}{\texttt{YM/FullVolume/}}
(in particular \texttt{Roll/}) and
\href{https://github.com/yang-mills-group/OpenAI-YM-20261028/tree/3456e5e75cf8d4cd558955a5c7638fba2d68f464/YM/TraceExtension}{\texttt{YM/TraceExtension/}}.
\textbf{Size:} 18,156 lines in the main-theorem closure.

\subsubsection{1. What this unit proves}\label{what-this-unit-proves-4}

\textbf{Proposition III.2.1 (order of choice).} For simple \(G\),
choose:

\begin{enumerate}
\def\labelenumi{\arabic{enumi}.}
\tightlist
\item
  A sheet: kernel \(K\) with branches on every patch; first stage
  \(j_0\ge4\); floors \(f_j\in\mathbb N\); stage records with a clip of
  relaxation \(r_j\), derivative orders \(k_j\ge6\), spatial powers
  \(q_j\ge2\), support exponents \(a_j>0\), widths \(w_j>0\), path
  counts \(\varpi_j>0\), drift allowances \(\omega_j\ge0\), column
  bounds \(\gamma_j\ge1\), \(r_jw_j\ge12\); nondecreasing \(k,q,a\).
\item
  Nondecreasing odd \(\ell_j\ge\max(3,f_j)\), with, for \(j\ge j_0\),
  \[\omega_j/\log\ell_j\le\min(\kappa/4,j^{-2}),\qquad\log\mathcal C_j/\log\ell_j\le j^{-2},\qquad\mathcal C_j=1+\varpi_j+w_j^{-1}+\gamma_j+r_j+a_j.\]
\item
  Positive nonincreasing \(\theta_j\le1/j\) for \(j\ge j_0\).
\item
  Unrestricted delays: real entry couplings \(\beta^{\rm ent}_j\), real
  barrier-dependent entry scales \(\sigma^{\rm ent}_j(\bar\beta)\),
  natural shells \(S_j\) and exponents \(g_j\).
\item
  A barrier \(\bar\beta>1\), with \(\bar\beta\ge\beta^{\rm ent}_{j_0}\),
  with \(x^{-1/2}e^{(\log\log x)^2}\le1\) for every \(x\ge\bar\beta\).
\item
  A wait: real readiness scales \(\sigma^{\rm rdy}_j\), windows
  \(h_j:\mathbb R\to\mathbb N\) with \(h_j(s)/s\to0\) as \(s\to\infty\)
  for each fixed \(j\), and natural \(n_{\min}\).
\end{enumerate}

Delays and waits are ordered componentwise. Factors exist for every
sheet; \(\theta\) can lie below any prescribed positive sequence;
barriers exist above any prescribed real. \emph{Status:} definitions and
three existence theorems. \emph{Lean:}
\href{https://github.com/yang-mills-group/OpenAI-YM-20261028/blob/3456e5e75cf8d4cd558955a5c7638fba2d68f464/YM/Price/Intent.lean\#L26-L35}{\texttt{Sheet}},
\href{https://github.com/yang-mills-group/OpenAI-YM-20261028/blob/3456e5e75cf8d4cd558955a5c7638fba2d68f464/YM/Price/Intent.lean\#L79-L84}{\texttt{Rolls.exists\_rolls}},
\href{https://github.com/yang-mills-group/OpenAI-YM-20261028/blob/3456e5e75cf8d4cd558955a5c7638fba2d68f464/YM/Price/Intent.lean\#L98-L111}{\texttt{Protection.exists\_margin}},
\href{https://github.com/yang-mills-group/OpenAI-YM-20261028/blob/3456e5e75cf8d4cd558955a5c7638fba2d68f464/YM/Price/Length.lean\#L19-L37}{\texttt{Bar.exists\_above}}.

\textbf{Proposition III.2.2 (schedule).} For simple \(G\) and any
III.2.1 data, there are \(J_m\), \(L_m=\ell_{J_{m+1}}\),
\(n_p=4(M+p+1)\) with natural \(M\), retaining the prescribed data, such
that:

\begin{itemize}
\tightlist
\item
  \(J_0=J_1=j_0\); \(J\) is nondecreasing, increases by at most one, and
  diverges; \(k_{J_m},q_{J_m},a_{J_m}\) are nondecreasing.
\item
  For every \(m\), \(\beta^{\rm ent}_{J_{m+1}}\le\hat\beta_m\); also
  \(\sigma^{\rm ent}_{J_{m+1}}(\bar\beta)\le s_m\) if \(J_{m+1}>j_0\).
\item
  For every \(m,p\), with \(j=J_{m+1}\): \(n_p\ge n_{\min}\),
  \(S_j(1+B_m+B_{m+1})^{g_j}\le n_pe^{s_m}\),
  \(100(B_m+1)\le n_pe^{s_m}\), \(100(R_m+1)\le n_pe^{s_{m+1}}\).
\item
  Each \(j\ge j_0\) has some \(k\) with \(\sigma^{\rm rdy}_j\le s_k\),
  \(h_j(s_k)\le k\), and \(J_i=j\) for \(|i-k|\le h_j(s_k)\).
\item
  \(\hat\beta_m>1\), \(t_m\le1\), \(0<\theta_{J_m}\le1/J_m\);
  \(\omega_{J_{m+1}}/\log L_m\) and
  \(\log\mathcal C_{J_{m+1}}/\log L_m\) are at most \(J_{m+1}^{-2}\);
  \(\log L_m/s_m\to0\). Kernels are well posed, lattice symmetric and
  gauge covariant.
\end{itemize}

Every positive integer divides \(n_p\) infinitely often. \emph{Status:}
theorem; divisibility variant in III.10. \emph{Lean:}
\href{https://github.com/yang-mills-group/OpenAI-YM-20261028/blob/3456e5e75cf8d4cd558955a5c7638fba2d68f464/YM/Price/Waits.lean\#L137-L239}{\texttt{Wait.ready}},
\href{https://github.com/yang-mills-group/OpenAI-YM-20261028/blob/3456e5e75cf8d4cd558955a5c7638fba2d68f464/YM/PhysicalNext/RG_NEXT_COMPLETE_REPEATABLE_NONLINEAR_STEP/Pane/Divisors.lean\#L26-L135}{\texttt{Wait.engine}}.

\textbf{Proposition III.2.3 (shooting).} Let \(G\) be simple; fix a
level system with \(\omega_{J_{m+1}}\le(\kappa/4)\log L_m\) and a recipe
prescribing paths \(P^{\mathcal R}_{N,m}\) and start, retention and
transition tests. Admissibility requires these paths, the level budget
point and clip, width \(w_{J_m}\), path count \(\varpi_{J_m}\), loop cap
\(2M_m^{\theta_{J_m}}\), allowances
\(\hat\beta_m^{-11/10},e^{-M_m^{1/5}}\), and retention. Continuity
requires continuous \(\beta,v\), fixed activity address sets, and
differences tending to zero in summed budget masses of regular
derivative bounds through budget order on the healthy domain and forced
uniform bounds on all configurations.

Assume, for every \(N,p\), a continuous admissible Wilson start family
on \(\mathfrak B_N\) with coupling equal to its parameter, \(v=0\),
density equal to the Wilson weight at that coupling, passing the start
test. Assume also that every continuous admissible level-\((m+1)\)
family on any \([l,u]\) with \(l\le u\), \(m<N\), with couplings in
\(\mathfrak B_{m+1}\), admits a continuous admissible level-\(m\) family
whose density is its block-spin push-forward at the input coupling,
passing the transition test, with
\(|\beta'-\beta-\kappa\log L_m|\le\omega_{J_{m+1}}\).

Then for all \(N,p\) a selected history exists: admissible
representations in \(\mathfrak B_m\), Wilson top with \(v=0\), these
transitions, \(\beta_{N,p,0}=\bar\beta\), paths independent of \(p\).
For every \(\varepsilon>0\), some \(m_0\) gives
\(|\beta_{N,p,m}-\hat\beta_m|\le\varepsilon s_m\) for all
\(N\ge m\ge m_0\), all \(p\). \emph{Status:} conditional theorem;
hypotheses discharged in III.10. \emph{Lean:}
\href{https://github.com/yang-mills-group/OpenAI-YM-20261028/blob/3456e5e75cf8d4cd558955a5c7638fba2d68f464/YM/Shoot/Race.lean\#L67-L82}{\texttt{Recipe.races}}.

\textbf{Proposition III.2.4 (joint purchase).} Assume simple \(G\),
kernel \(K\) branched on every patch, first stage \(\ge4\),
nondecreasing \(k_j\ge\max(6,j)\), \(q_j\ge\max(2,j)\), \(a_j\ge100\),
nondecreasing \(Q:\mathbb N\to\mathbb N\) with \(Q_j\ge1\), records
\(\Omega,\Delta\), degree-\(300\) tableau at \(\xi\) and contraction
condition
\href{https://github.com/yang-mills-group/OpenAI-YM-20261028/blob/3456e5e75cf8d4cd558955a5c7638fba2d68f464/YM/FullVolume/Sym/Purse.lean\#L29}{\texttt{sphere}}
at \((\xi,\omega_v,\omega_t)\), with \(\xi\ge2\), \(0<\omega_v\le1\),
\(\omega_t>0\), and prediction \(\pi\) prepaid at every sheet and price
map.

There exist a sheet \(s\) preserving the prescribed kernel, first stage,
derivative orders, powers and support exponents, price map \(T\), real
sequence \(Y\), nondecreasing \(K'\ge1\) and floor \(\tau\ge f^s\). For
every floor \(k\ge\tau\), factors exist; for every block-factor sequence
\(\ell\) for \(s\) with floors \(k\) there is \(U\); for every wage
\(w\ge U\), postage rate and exchange datum there are post-\(L\) costs;
for every positive \(\theta^0\) there is protection
\(\theta\le\theta^0\); for every menu, auxiliary record and delay
\(T_0\) there is \(V\ge T_0\); for every \(V'\ge V\), barrier and wait
\(W\) there is \(Z\ge W\); for every \(W'\ge Z\) and level system
meeting III.2.2 with divisible periods there are a directed datum and
factor \(X\) recording these choices, satisfying \(\pi(X)\).

Wages, exchanges, menus and auxiliary records equal the prescribed data;
postage rates do not exceed the prescribed rates. The same factor
satisfies the III.12 conditions
\href{https://github.com/yang-mills-group/OpenAI-YM-20261028/blob/3456e5e75cf8d4cd558955a5c7638fba2d68f464/YM/FullVolume/Sym/Prices.lean\#L84}{\texttt{Degree}},
\href{https://github.com/yang-mills-group/OpenAI-YM-20261028/blob/3456e5e75cf8d4cd558955a5c7638fba2d68f464/YM/FullVolume/Sym/Shorn.lean\#L32-L35}{\texttt{Savvy}},
\href{https://github.com/yang-mills-group/OpenAI-YM-20261028/blob/3456e5e75cf8d4cd558955a5c7638fba2d68f464/YM/FullVolume/C/Ribbon.lean\#L29}{\texttt{Roofing}},
\href{https://github.com/yang-mills-group/OpenAI-YM-20261028/blob/3456e5e75cf8d4cd558955a5c7638fba2d68f464/YM/FullVolume/Florid/Ready.lean\#L56}{\texttt{PaysW}},
\href{https://github.com/yang-mills-group/OpenAI-YM-20261028/blob/3456e5e75cf8d4cd558955a5c7638fba2d68f464/YM/FullVolume/Luna/Hours.lean\#L43}{\texttt{Spire}},
\href{https://github.com/yang-mills-group/OpenAI-YM-20261028/blob/3456e5e75cf8d4cd558955a5c7638fba2d68f464/YM/FullVolume/Auvi/Entries.lean\#L31}{\texttt{Wild}},
\href{https://github.com/yang-mills-group/OpenAI-YM-20261028/blob/3456e5e75cf8d4cd558955a5c7638fba2d68f464/YM/FullVolume/Spiral/Rent.lean\#L25-L27}{\texttt{Tower.Sav}}
with
\href{https://github.com/yang-mills-group/OpenAI-YM-20261028/blob/3456e5e75cf8d4cd558955a5c7638fba2d68f464/YM/FullVolume/Fares/Query.lean\#L45}{\texttt{Canyon}},
\href{https://github.com/yang-mills-group/OpenAI-YM-20261028/blob/3456e5e75cf8d4cd558955a5c7638fba2d68f464/YM/FullVolume/Artes/Events.lean\#L34-L39}{\texttt{Cadre}},
\href{https://github.com/yang-mills-group/OpenAI-YM-20261028/blob/3456e5e75cf8d4cd558955a5c7638fba2d68f464/YM/MergeProof/North/Tibetan.lean\#L60}{\texttt{AccordV}}
with \(Q\) (III.11), top kinetic data (III.9), amplification (III.1),
and a bandage with extras \((\Omega,\Delta,b)\) (III.15). Support
exponents are \(\ge100\);
\(e^{(a_j+2)(2\,\mathrm{span}_{ij}+13)^4}\le\ell_i\) for
\(j\le i\le j+1\). The bandage also supplies the conditional family
package
\href{https://github.com/yang-mills-group/OpenAI-YM-20261028/blob/3456e5e75cf8d4cd558955a5c7638fba2d68f464/YM/FullVolume/Unlatch.lean\#L25-L34}{\texttt{Unlatched}}
with a vacuous additional conclusion, under the family hypotheses of
III.2.5. \emph{Status:} theorem. \emph{Lean:}
\href{https://github.com/yang-mills-group/OpenAI-YM-20261028/blob/3456e5e75cf8d4cd558955a5c7638fba2d68f464/YM/Integration/RG05Preconstruction/Fares.lean\#L26-L85}{\texttt{richFaresWithPrediction}}.

\textbf{Proposition III.2.5 (plentiful construction).} Under III.2.4 and
\(Q_i\ge2\), its conclusion additionally ensures: every compatible
trajectory, reference levels with \(d+3\le J_{t(d)}\), section transfers
and consecutive gauge doors satisfying family consistency, door
compatibility and prescribed transfer bounds from level \(t(d)\)
whenever \(t(d)\le N\) give, for every extraction order, the complete
weak-continuum requirement (III.11) and trace gallery (VII.2). Floors,
delays and waits may increase; the prediction persists. \emph{Status:}
theorem. \emph{Lean:}
\href{https://github.com/yang-mills-group/OpenAI-YM-20261028/blob/3456e5e75cf8d4cd558955a5c7638fba2d68f464/YM/Integration/RG05Preconstruction/Bardo.lean\#L26-L81}{\texttt{plentifulWithPrediction}},
\href{https://github.com/yang-mills-group/OpenAI-YM-20261028/blob/3456e5e75cf8d4cd558955a5c7638fba2d68f464/YM/FullVolume/Unlatch.lean\#L25-L34}{\texttt{Unlatched}}.

\textbf{Proposition III.2.6 (prediction).} Let \(\pi\) be a prediction.
For simple \(G\), a sheet with the perceived-estimate package of
detailed Lemma 4.16 at fixed tariff and scope, and \(\tau>0\), suppose:
for every choice of targets with the canonical chart, expenses with
deficit \(-7/16\), programs, initial wages and baton there is a floor
\(u\); for every floor \(l\ge u\), factors, wages, duration bounds,
register, postage, scouts, protection, step records with
\(\varepsilon\le1/4\) on transitions with output stage at least the
first stage, post-\(L\) costs and delay \(V\), there is \(W\ge V\); for
every later delay, barrier, wait, pricing record, directed datum and
remaining factor data, shelter and step conditions and divisible periods
imply \(\pi\) on the assembled factor. Then the complete native purchase
of detailed Lemma 4.17 retains \(\pi\) on the emitted factor, on level
systems with divisible periods, maximizing its floor with \(u\) and
inserting \(W\) before its booth and final delay. \emph{Status:}
theorem; used with vacuous prediction. \emph{Lean:}
\href{https://github.com/yang-mills-group/OpenAI-YM-20261028/blob/3456e5e75cf8d4cd558955a5c7638fba2d68f464/YM/Integration/NativeSelection/FundedCompanySelection.lean\#L49-L106}{\texttt{company}},
\href{https://github.com/yang-mills-group/OpenAI-YM-20261028/blob/3456e5e75cf8d4cd558955a5c7638fba2d68f464/YM/Integration/NativeSelection/FundedCompanySelection.lean\#L30-L42}{\texttt{Prepaid}}.

\textbf{Proposition III.2.7 (complete run).} Given simple \(G\), penalty
\(b>0\), records \(\Omega,\Delta\), custody request \(\mathcal H\),
eventually paid \(\Lambda\), and a demand assigning every sheet and
price map nonempty choices \(c\), floors \(e_c:\mathbb N\to\mathbb N\)
and real arrays
\(F_c(n,i),F'_c[\ell](https://github.com/scottnarmstrong/yang-mills-informalization/blob/main/outline/units/n,i)\),
there is a run with a campaign, penalty \(b\), preselection \(s,q,c\),
and floor \(n\ge\max(f^s,e_c)\). It uses \(s\) with floors \(n\) and the
price map derived from \(q\); both arrays are paid as in III.2.8(b),
using actual \(\ell\), and
\(F'_c[\ell](https://github.com/scottnarmstrong/yang-mills-informalization/blob/main/outline/units/0,J_{m+1})\le\hat\beta_m\)
for all \(m\). Orders satisfy \(k_j,q_j\ge j\); \(\Lambda\), treasury
funding, common-period funding, pillars, the funded, news and read-out
records, custody and sky hold. Any universally prepaid prediction can
persist. \emph{Status:} theorem. \emph{Lean:}
\href{https://github.com/yang-mills-group/OpenAI-YM-20261028/blob/3456e5e75cf8d4cd558955a5c7638fba2d68f464/YM/Integration/NativePreselectedPostnumericPeriodRun.lean\#L99-L114}{\texttt{selectRunWithTreasuryPeriod\_preselected\_quote\_orders\_postnumeric}}.

\textbf{Proposition III.2.8 (persistent payments).}

\begin{enumerate}
\def\labelenumi{(\alph{enumi})}
\item
  Eventually paid means: for every sheet, factors, protection and delay
  \(V\), some \(X\ge V\) ensures the property for every delay
  \(W\ge X\), every barrier and wait, and every system meeting their
  entrance conditions. Conjunction preserves it.
\item
  For every real array \(F(n,i)\), eventually paid are
  \(1\le1+\log\log\hat\beta_m\) at all levels and: for every \(n\),
  eventually in \(m\), \(|F(n,i)|\le1+\log\log\hat\beta_m\) for all
  \(i\le J_{m+1}\). For every real sequence \(F\),
  \(F(J_{m+1})\le\hat\beta_m\) at all levels is eventually paid.
\item
  Every countably encoded family of requests (stage floors and costs
  depending on the transition, block geometry and stage record
  \(X_{J_m}\) at step \(m\)) has a common floor \(\mathrm{pre}\) fixed
  before factors. If \(\ell_j\ge\mathrm{pre}(j)\) at every stage, every
  delay can be enlarged so every later delay, any barrier and wait, and
  every admitted system satisfy, for \(\mathrm{code}(\alpha)\le J_n\),
  \(m\ge n\), \(\mathrm{pre}_\alpha(J_{m+1})\le L_m\) and step cost
  \(\le1+\log(2+\log\hat\beta_m)\).
\item
  For simple \(G\) and every kernel, a plan fixes charge names and
  dictionaries of finite words and spectra for every request before
  factors. Above its floor, delays can be enlarged persistently to fund
  requests and the inverse charge from level zero. \emph{Status:}
  definition in (a); theorems. \emph{Lean:}
  \href{https://github.com/yang-mills-group/OpenAI-YM-20261028/blob/3456e5e75cf8d4cd558955a5c7638fba2d68f464/YM/R2I/Campaign/Clock.lean\#L36-L38}{\texttt{EventuallyPaid}},
  \href{https://github.com/yang-mills-group/OpenAI-YM-20261028/blob/3456e5e75cf8d4cd558955a5c7638fba2d68f464/YM/Integration/NativeDiagonalClockRequest.lean\#L52-L79}{\texttt{eventuallyPaid}},
  \href{https://github.com/yang-mills-group/OpenAI-YM-20261028/blob/3456e5e75cf8d4cd558955a5c7638fba2d68f464/YM/RG08/Calendar/Base.lean\#L68-L77}{\texttt{minDelay}},
  \href{https://github.com/yang-mills-group/OpenAI-YM-20261028/blob/3456e5e75cf8d4cd558955a5c7638fba2d68f464/YM/Integration/CalendarTreasuryPurchase.lean\#L42-L64}{\texttt{buy}}.
\end{enumerate}

\textbf{Proposition III.2.9 (analytic estimates).} Norms are
componentwise exponentially weighted displacement \(\ell^1\) norms,
summing colour indices for multilinear kernels; colour spaces are fixed
and finite, and constants may depend on them. Crosses map coarse fields
to fine fields at block size \(L\). Vanishing orders concern symbols at
zero momentum.

\begin{itemize}
\tightlist
\item
  For \(\sigma,\tau>0\), real \(C\), natural \(d\), some
  \(0<\eta\le\sigma\), \(K\ge0\) give
  \(\|A^{\mathsf T}RB\|_\tau\le KML^{2-d}\) for positive odd \(L\) with
  \(\tau\le L\eta\), \(\|A\|_\tau\le C/L\),
  \(\|\nabla^sB\|_\tau\le C/L^{1+|s|}\) through order \(d\), \(M\ge0\),
  \(\|R\|_\sigma\le M\), and vanishing order \(d\). For one-form rows,
  \(d=3\) gives \(ML^{-1/2}\) for all sufficiently large odd \(L\).
  \href{https://github.com/yang-mills-group/OpenAI-YM-20261028/blob/3456e5e75cf8d4cd558955a5c7638fba2d68f464/YM/RG08/Flat/Pullback.lean\#L56-L110}{\texttt{pullback\_gain}},
  \href{https://github.com/yang-mills-group/OpenAI-YM-20261028/blob/3456e5e75cf8d4cd558955a5c7638fba2d68f464/YM/RG08/Flat/KStep.lean\#L130-L146}{\texttt{large\_contraction}}.
\item
  For natural \(r,d\), \(\sigma,w>0\), \(C\ge0\), some
  \(0<\eta\le\sigma\), \(K\ge0\), fixed before \(L\), give transfer
  bound \(KBL^4/L^{r+1+d}\) at width \(w\). Assume a real
  \((r+1)\)-linear kernel of width \(\sigma\), bound \(B\), supported on
  tuples with first site zero, vanishing order \(d\) in independent
  momenta; positive odd \(L\), \((r+2)w\le L\eta\); crosses with
  \(\|\nabla^iQ\|_{(r+2)w}\le C/L^{1+|i|}\) through order \(d\);
  transfer through \(r+1\) copies of \(Q\), followed by symmetrization
  and recentering.
  \href{https://github.com/yang-mills-group/OpenAI-YM-20261028/blob/3456e5e75cf8d4cd558955a5c7638fba2d68f464/YM/RG08/Ballast/Slide.lean\#L44-L90}{\texttt{flow\_soft}}.
\item
  For finite coordinates \(I\), natural \(n\), \(\sigma>0\), some
  \(0<\eta\le\sigma\), \(K\ge0\) ensure: every real kernel
  \(f:\mathbb Z^I\to\mathbb R\) of weighted bound \(A\ge0\), vanishing
  order \(n\), has real kernels \(g_b\) of bound \(KA\), width \(\eta\),
  with
  \(\widehat f(p)=\sum_b\prod_{a=1}^n(e^{ip_{b(a)}}-1)\widehat g_b(p)\)
  for all real \(p\), where \(b:\{1,\ldots,n\}\to I\).
  \href{https://github.com/yang-mills-group/OpenAI-YM-20261028/blob/3456e5e75cf8d4cd558955a5c7638fba2d68f464/YM/RG08/Flat/Differences.lean\#L46-L92}{\texttt{real\_difference\_division}}.
\item
  Fix step scenes with input coupling tending to infinity, leaves,
  cutoff parameter, radii satisfying the logarithm-chart conditions and
  fluctuations \(z_i\); put \(H_i=\frac12\log\beta_i^{\rm out}\),
  \(\eta_{*,i}=e^{-H_i}\). Assume fluctuation bounds, leaf weights and
  third-derivative bounds \(D_i\) have upper exponential rate zero
  relative to \(H_i\); eventually leaves are masked regular logarithmic
  terms whose threshold-squared expressions have zero first derivative
  at zero and derivatives through order three bounded by \(D_i\) on the
  unit ball; \(\pm z_i\) are healthy; both cutoffs equal one near the
  origin in the two deformation parameters. Here upper rate \(a\) means
  an eventual bound \(e^{rH_i}\) for every \(r>a\). Writing \(P_i\) for
  the paved expansion, for every \(\varepsilon>0\), eventually
  \(\|D_y^2(P_i(z_i;y)+P_i(-z_i;y))|_{y=0}\|\le\eta_{*,i}^{2-\varepsilon}\).
  \href{https://github.com/yang-mills-group/OpenAI-YM-20261028/blob/3456e5e75cf8d4cd558955a5c7638fba2d68f464/YM/RG12/Shadows/Shape.lean\#L54-L70}{\texttt{mirror\_price}}.
\item
  Fix a representation, degree-\(d\) frame and section readable at the
  origin for the representation's budget. Two prescribed coefficient
  presentations with equal section densities, remainder bounds \(a,b\)
  and the same positive forced parameter satisfy
  \(\|x_P-x_Q\|_{\rm coeff}\le K_{F,H}e^{-a_*m}(a+b)\) for every natural
  \(m\) with \(10(m(4B+2)+2B)<n\). Regular terms have supports
  containing zero, vanish at identity and have zero local degree-\(d\)
  projection on nonwrapping supports; nonzero forced supports contain
  zero and satisfy eligibility, support, gauge-invariance and
  measurability conditions.
  \href{https://github.com/yang-mills-group/OpenAI-YM-20261028/blob/3456e5e75cf8d4cd558955a5c7638fba2d68f464/YM/TraceExtension/Quince/Unum.lean\#L11-L42}{\texttt{unique\_part}}.
\end{itemize}

Further estimates concern the post-block Hessian
(\href{https://github.com/yang-mills-group/OpenAI-YM-20261028/blob/3456e5e75cf8d4cd558955a5c7638fba2d68f464/YM/RG12/Spins/Muse.lean\#L35-L59}{\texttt{design}}),
cubic and quartic coefficients
(\href{https://github.com/yang-mills-group/OpenAI-YM-20261028/blob/3456e5e75cf8d4cd558955a5c7638fba2d68f464/YM/FullVolume/Canvas/Reservoir.lean\#L84-L118}{\texttt{tanker}}),
and trace insertion
(\href{https://github.com/yang-mills-group/OpenAI-YM-20261028/blob/3456e5e75cf8d4cd558955a5c7638fba2d68f464/YM/TraceExtension/Zn/Arc.lean\#L16-L33}{\texttt{cloche}}).
\emph{Status:} theorems under the displayed hypotheses; only statements
of these supplementary estimates are included in the notes.

\subsubsection{2. What it uses}\label{what-it-uses-4}

\begin{itemize}
\tightlist
\item
  I.1: \(\kappa>0\); III.1: calendars, histories, branches; III.9:
  reference quantities and Wilson startup; III.10: one-step construction
  and divisible periods.
\item
  III.11--III.12, III.15, VII.2: purchased conditions; III.14: demand
  and external records; III.16: strip and Laurent norms and jets.
  Free-kernel and stationary-branch inputs: III.3 and III.6 (probable).
\item
  Intermediate value theorem, compactness, logarithmic asymptotics,
  countable diagonal enumeration.
\end{itemize}

The coupling shift with error \(\omega\) is an input to shooting,
supplied by III.10
(\href{https://github.com/yang-mills-group/OpenAI-YM-20261028/blob/3456e5e75cf8d4cd558955a5c7638fba2d68f464/YM/PhysicalNext/RG_NEXT_COMPLETE_REPEATABLE_NONLINEAR_STEP/Graces/Eagles.lean\#L33-L50}{\texttt{Graces.winnow}}).
Positivity of \(\kappa\) comes from
\href{https://github.com/yang-mills-group/OpenAI-YM-20261028/blob/3456e5e75cf8d4cd558955a5c7638fba2d68f464/YM/Gauge/Casimir.lean\#L174-L176}{\texttt{slopeConstant\_pos}}.

\subsubsection{3. Strategy}\label{strategy-4}

Choose stage data before factors, then costs, protection, delays,
barrier and wait. Maxima combine floors; persistent increases combine
delays and waits. Long plateaux provide comparison windows; exponential
side growth pays volume demands. Shooting narrows continuous Wilson
families into successive coupling bands until the level-zero coupling
equals \(\bar\beta\). The shooting windows are
\([\bar\beta+(\kappa/2)s_m,\bar\beta+2\kappa s_m]\subset\mathfrak B_m\).

\subsubsection{4. Steps}\label{steps-4}

\begin{enumerate}
\def\labelenumi{\arabic{enumi}.}
\tightlist
\item
  \textbf{Exhaustion.} Long plateaux meet entry and window thresholds;
  one period offset pays the finite first plateau.
  \href{https://github.com/yang-mills-group/OpenAI-YM-20261028/blob/3456e5e75cf8d4cd558955a5c7638fba2d68f464/YM/Price/Waits.lean\#L41-L84}{\texttt{Program.exhaust}},
  \href{https://github.com/yang-mills-group/OpenAI-YM-20261028/blob/3456e5e75cf8d4cd558955a5c7638fba2d68f464/YM/Price/Waits.lean\#L137-L239}{\texttt{Wait.ready}}.
\item
  \textbf{Narrowing.} Endpoint drift inequalities straddle the next
  window; first and last hitting points give a subinterval mapping
  entirely into it. Iterate downward.
  \href{https://github.com/yang-mills-group/OpenAI-YM-20261028/blob/3456e5e75cf8d4cd558955a5c7638fba2d68f464/YM/Shoot/Reach.lean\#L81-L90}{\texttt{gates}},
  \href{https://github.com/yang-mills-group/OpenAI-YM-20261028/blob/3456e5e75cf8d4cd558955a5c7638fba2d68f464/YM/Shoot/Reach.lean\#L18-L49}{\texttt{narrow}},
  \href{https://github.com/yang-mills-group/OpenAI-YM-20261028/blob/3456e5e75cf8d4cd558955a5c7638fba2d68f464/YM/Shoot/Reach.lean\#L112-L154}{\texttt{steps}}.
\item
  \textbf{Running law.} The coupling error is bounded by
  \(\sum_{k=1}^m\omega_{J_k}=o(s_m)\) using \(J_k\to\infty\).
  \href{https://github.com/yang-mills-group/OpenAI-YM-20261028/blob/3456e5e75cf8d4cd558955a5c7638fba2d68f464/YM/Shoot/Race.lean\#L23-L53}{\texttt{trajectory}}.
\item
  \textbf{Diagonal payment.} Raise entry couplings above double
  exponentials of finite partial sums of array entries. Since
  \(J_m\to\infty\), every fixed row is eventually included.
  \href{https://github.com/yang-mills-group/OpenAI-YM-20261028/blob/3456e5e75cf8d4cd558955a5c7638fba2d68f464/YM/Integration/NativeDiagonalClockRequest.lean\#L52-L79}{\texttt{eventuallyPaid}}.
\item
  \textbf{Composition.} Combine sheet and native suppliers, joint and
  trace purchases and persistent payments by maxima of floors and
  transitivity of delay and wait bounds. Choose the schedule only after
  these increases, then assemble the run.
  \href{https://github.com/yang-mills-group/OpenAI-YM-20261028/blob/3456e5e75cf8d4cd558955a5c7638fba2d68f464/YM/FullVolume/Roll/Fiscal.lean\#L97-L169}{\texttt{sailsW}},
  \href{https://github.com/yang-mills-group/OpenAI-YM-20261028/blob/3456e5e75cf8d4cd558955a5c7638fba2d68f464/YM/Integration/RG05Preconstruction/Bardo.lean\#L26-L81}{\texttt{plentifulWithPrediction}},
  \href{https://github.com/yang-mills-group/OpenAI-YM-20261028/blob/3456e5e75cf8d4cd558955a5c7638fba2d68f464/YM/Integration/RG05Preconstruction/PeriodConstruction.lean\#L30-L47}{\texttt{engageWithCustodyFactor}}.
\item
  \textbf{Analytic gains.} Move differences from vanishing kernels onto
  crosses: their two bounds give \(L^{-2}\), fine-site summation gives
  \(L^4\), and \(d\) differences give \(L^{-d}\). Reflection cancels odd
  terms; local projection leaves exponentially small wrapping
  contributions.
  \href{https://github.com/yang-mills-group/OpenAI-YM-20261028/blob/3456e5e75cf8d4cd558955a5c7638fba2d68f464/YM/RG08/Flat/Pullback.lean\#L56-L110}{\texttt{pullback\_gain}},
  \href{https://github.com/yang-mills-group/OpenAI-YM-20261028/blob/3456e5e75cf8d4cd558955a5c7638fba2d68f464/YM/RG12/Shadows/Shape.lean\#L54-L70}{\texttt{mirror\_price}},
  \href{https://github.com/yang-mills-group/OpenAI-YM-20261028/blob/3456e5e75cf8d4cd558955a5c7638fba2d68f464/YM/TraceExtension/Quince/Unum.lean\#L11-L42}{\texttt{unique\_part}}.
\end{enumerate}

\subsubsection{5. Where the difficulty
lies}\label{where-the-difficulty-lies-4}

Factor-dependent costs require later persistent thresholds. Shooting
needs full-family continuity and band membership throughout each
interval. The \(\kappa/4\) drift bound separates endpoints; \(j^{-2}\)
controls accumulated error. Each row is paid eventually; the zeroth
post-\(L\) row needs separate all-level payment.

\subsubsection{6. Notation}\label{notation-4}

\(m=0\) is coarsest, \(N\) finest; \(p\) indexes volume, \(j\) stages.
\(s_m=\sum_{i<m}\log L_i\), \(\kappa=11C_G/(24\pi^2)\),
\(\hat\beta_m=\bar\beta+\kappa s_m\),
\(\mathfrak B_m=[\hat\beta_m/2,2\hat\beta_m]\),
\(M_m=e^{(\log\log\hat\beta_m)^2}\), \(t_m=\hat\beta_m^{-1/2}M_m\),
\(B_m=\lceil(\log\hat\beta_m)^2\rceil\), \(R_m=(J_{m+1}+3)L_m\). In
read-out, \(n,B,a_*\) are side, cover size, support exponent;
\(K_{F,H}\) depends on the frame and holonomy tests.
\(\mathrm{span}_{ij}\) is transition support; a prediction is a
predicate on the factor recording all choices. Stage records, protection
and admitted entrance conditions are
\href{https://github.com/yang-mills-group/OpenAI-YM-20261028/blob/3456e5e75cf8d4cd558955a5c7638fba2d68f464/YM/RG/History.lean\#L45-L61}{\texttt{Stage}},
\href{https://github.com/yang-mills-group/OpenAI-YM-20261028/blob/3456e5e75cf8d4cd558955a5c7638fba2d68f464/YM/Price/Intent.lean\#L88-L92}{\texttt{Protection}}
and
\href{https://github.com/yang-mills-group/OpenAI-YM-20261028/blob/3456e5e75cf8d4cd558955a5c7638fba2d68f464/YM/Price/Intent.lean\#L145-L162}{\texttt{Paid}}.

\subsubsection{7. What remains
unexplained}\label{what-remains-unexplained-4}

Purchased predicates need comparison with their owning units. Underlying
analytic suppliers and dictionaries lack independent verification here;
drift normalization against III.9 remains unresolved. Odd-cancellation
scenes, coefficient and readability interfaces and auxiliary parameters
need fuller mathematical accounts. Supplementary Hessian,
coefficient-extraction and trace-insertion proofs remain unexplained.

Detailed notes: the full account of this unit, with proof sketches for
every step, is on the
\href{https://github.com/scottnarmstrong/yang-mills-informalization/blob/main/outline/notes/III-2.md}{detailed
notes page}.

\subsection{III.3. Free lattice kernels and
covariances}\label{iii.3.-free-lattice-kernels-and-covariances}

\textbf{Part:} III, Ultraviolet construction (renormalization group).
\textbf{Lean:}
\href{https://github.com/yang-mills-group/OpenAI-YM-20261028/tree/3456e5e75cf8d4cd558955a5c7638fba2d68f464/YM/Next/FREE_KERNELS}{\texttt{YM/Next/FREE\_KERNELS/}}
(5,450 lines),
\href{https://github.com/yang-mills-group/OpenAI-YM-20261028/tree/3456e5e75cf8d4cd558955a5c7638fba2d68f464/YM/MultiNext/SIGNED_SOURCE_CONTINUATION}{\texttt{YM/MultiNext/SIGNED\_SOURCE\_CONTINUATION/}}
(1,674),
\href{https://github.com/yang-mills-group/OpenAI-YM-20261028/tree/3456e5e75cf8d4cd558955a5c7638fba2d68f464/YM/Further/FREE_COVARIANCE_KP}{\texttt{YM/Further/FREE\_COVARIANCE\_KP/}}
(924),
\href{https://github.com/yang-mills-group/OpenAI-YM-20261028/tree/3456e5e75cf8d4cd558955a5c7638fba2d68f464/YM/UVNext/PRICED_KERNELS}{\texttt{YM/UVNext/PRICED\_KERNELS/}}
(533),
\href{https://github.com/yang-mills-group/OpenAI-YM-20261028/tree/3456e5e75cf8d4cd558955a5c7638fba2d68f464/YM/Further/CORRECTED_LIFT_COEFFICIENTS}{\texttt{YM/Further/CORRECTED\_LIFT\_COEFFICIENTS/}}
(248),
\href{https://github.com/yang-mills-group/OpenAI-YM-20261028/tree/3456e5e75cf8d4cd558955a5c7638fba2d68f464/YM/Push/CUTOFF_PROJECTORS}{\texttt{YM/Push/CUTOFF\_PROJECTORS/}}
(139),
\href{https://github.com/yang-mills-group/OpenAI-YM-20261028/tree/3456e5e75cf8d4cd558955a5c7638fba2d68f464/YM/ProofPorts/GRID_RENEWED_RESERVES}{\texttt{YM/ProofPorts/GRID\_RENEWED\_RESERVES/}}
(30). \textbf{Size:} 8,998 lines in the main-theorem closure.

\subsubsection{1. What this unit proves}\label{what-this-unit-proves-5}

For fine/coarse levels, primes denote coarse data and
\(V=(\sqrt c\,1,R)\), where \(c\) is the direct clip coefficient. The
diagonal identities (D) and roll identities (R) are \[\begin{aligned}
 d^*HdW&=Q^*d^*H'd,&QW&=1-\kappa^{-1}d^*H'd,\\
 T^*d&=dQ,&TD'd&=DdW;\\
 M_f^*M_f+M_n^*M_n+O^*O&=1,&\sqrt\kappa M_n^*&=V'd,\\
 M_fV'd&=VdW,&M_nV'd&=\kappa^{-1/2}d^*H'd,\\
 OV'd&=0,&M_f^*Vd&=\sqrt\kappa M_n^*Q.
\end{aligned}\] The defect identity (Δ) is
\[H'(H'^{-1}-T^*H^{-1}T-\kappa^{-1}dd^*)H'=d_2^*Pd_2.\] Under (D),(R),
for all fine/coarse potentials \(A,B\), setting \(Z=A-WB\) gives
\[\langle dA,HdA\rangle+\kappa|QA-B|^2=\langle dB,H'dB\rangle+\langle dZ,HdZ\rangle+\kappa|QZ|^2.\]

\textbf{Proposition III.3.1 (catalogue).} For every kernel choice \(K\),
there exist sequences \(\alpha_j,\eta_j,\gamma_j>0\),
\(\Omega_j\in\mathbb R\), \(\lambda_j\ge0\), \(b_j\in\mathbb N\),
\(j\in\mathbb N\), depending only on \(K\). For every history of index
\(j\ge4\) and clip coefficient \(0<c\le\gamma_j\), its explicit metric
data \((D,D^{-1},H,H^{-1},R)\) lie in \(\mathcal S(\alpha_j,\Omega_j)\),
with \[H^{-1}=D^{-1}+\lambda_j(d_2D^{-1})^*(d_2D^{-1}).\] On every
torus, \(DD^{-1}=HH^{-1}=1\), \(c1+R^*R=H\), \(\hat H(0)=1\), \(Hd=Dd\);
\(H,D\) are symmetric with bounds \(\eta_j1\le H,D\le\Omega_j1\). For
every gauge group, \(d^*Hd\) is hypercubic invariant; empty histories
give \(D=1\).

For every flight \((j,t)\), \(j,t\ge4\), \(t\in\{j-1,j\}\), and
difference word \(|u|\le j-2\),
\[\nabla_uW\in\mathcal S(\alpha_j,\Omega_jL^{-1-|u|}),\qquad\nabla_uT\in\mathcal S(\alpha_j,\Omega_jL^{-2-|u|}).\]
Its rolled kernel
\(\mathcal P_{\lambda_j,\lambda_t}\in\mathcal S(\alpha_j,\Omega_j)\); if
\(L\ge b_j\), its symbol is Hermitian with
\(\eta_j1\le\hat{\mathcal P}\le\Omega_j1\) at every real momentum.
\emph{Status:} theorem. \emph{Lean:}
\href{https://github.com/yang-mills-group/OpenAI-YM-20261028/blob/3456e5e75cf8d4cd558955a5c7638fba2d68f464/YM/Next/FREE_KERNELS/Priced.lean\#L67-L93}{\texttt{pledges}}.

\textbf{Proposition III.3.2 (families).} For every gauge group \(G\) and
kernel choice \(K\), a catalogue as above exists such that every sheet
\(\mathfrak S\) with kernel \(K\) admits a second table
\((\alpha'_j,\Omega'_j,\eta_j)\). For every trajectory with that kernel
and the sheet's stage profile satisfying, at every level and step
respectively, \[c_i\le\gamma_{J(i)},\qquad L_i\ge b_{J(i+1)},\] there
are pre-repair and output families simultaneously for all \(N\ge0\),
\(0\le i\le N\). They share the explicit history metrics with amendments
\(\lambda_{J(i)}\) and flight kernels \(W,T\) of index \((J(i+1),J(i))\)
for \(i<N\); bare histories are empty and \(D^{N,N}=1\). Before repair,
III.3.1's bounds and symmetry hold, as do (D) and (Δ) on the schedule
torus pairs of each step, with \(P\in\mathcal S(\alpha_j,\Omega_j)\)
Hermitian and bounded below by \(\eta_j\) on \(\mathrm{ran}\,\sigma_2\)
at real momenta, \(j=J(i+1)\). After repair, (D),(R) hold there; metric
and difference bounds use the second table, retaining \(\eta_j\), and
\[M_n,O\in\mathcal S(\alpha'_j,\Omega'_j),\qquad M_f\in\mathcal S(\alpha'_j,\Omega'_jL_i^{-2}).\]
\emph{Status:} theorem under the credit condition (C1), (C2) as
hypothesis. The variant \texttt{before\_paid} (assigned to III.4)
discharges it from the III.9 sheet and the III.2 payment record.
\emph{Lean:}
\href{https://github.com/yang-mills-group/OpenAI-YM-20261028/blob/3456e5e75cf8d4cd558955a5c7638fba2d68f464/YM/Next/FREE_KERNELS/Rack.lean\#L134-L145}{\texttt{retained\_support}}.

\textbf{Proposition III.3.3 (all tori).} Every family satisfying
III.3.2's output conditions, for any table, satisfies (D),(R) for every
\(N\), \(m<N\) and positive sides \(c,f\) with \(f=cL_m\), using the
same kernels and level clips. \emph{Status:} theorem. \emph{Lean:}
\href{https://github.com/yang-mills-group/OpenAI-YM-20261028/blob/3456e5e75cf8d4cd558955a5c7638fba2d68f464/YM/UVNext/PRICED_KERNELS/Tori.lean\#L109-L157}{\texttt{complete\_roll}}.

\textbf{Proposition III.3.4 (rolled positivity).} For fixed kernel
choice \(K\) and integers \(j,t\ge4\), there exist \(C_0,B>0\) depending
only on these data. For every \(\lambda\ge C_0\), \(\mu\ge0\), there
exist \(\eta,\sigma,\Omega>0\) such that every flight \((j,t)\) has
\(\mathcal P_{\mu,\lambda}\in\mathcal S(\sigma,\Omega)\); if
\(\mu B\le L^2\), its symbol is Hermitian with
\(\eta1\le\hat{\mathcal P}\le\Omega1\) at all real momenta.
\emph{Status:} theorem. \emph{Lean:}
\href{https://github.com/yang-mills-group/OpenAI-YM-20261028/blob/3456e5e75cf8d4cd558955a5c7638fba2d68f464/YM/Next/FREE_KERNELS/Eject.lean\#L99-L140}{\texttt{Flight.frontier}}.

\textbf{Proposition III.3.5 (filters).} Define
\(\mathrm{rail}_M(z)=M^{-1}\sum_{s=0}^{M-1}e^{isz}\) for positive
integers \(M\),
\(\mathrm{plank}_M(z)=M^{-1}\sum_{|s|\le(M-1)/2}e^{isz}\) for odd \(M\),
and \(a_L(p)=\prod_a\mathrm{plank}_L(p_a)\). For all sizes where defined
and all \(z\in\mathbb C\), both line filters satisfy
\(F_L(Mz)F_M(z)=F_{LM}(z)\) and \(|F_M(z)|\le e^{M|\mathrm{Im}\,z|}\);
plank is even, \(2\pi\)-periodic and real on \(\mathbb R\). Absolute
\(c_*,c_1>0\) give \(|\mathrm{rail}_M(z)|\ge c_*\) when
\(M|z|\le3\pi/2\), and \(|\mathrm{plank}_M(z)|\ge e^{-c_1M/L}\) for real
\(z\), \(M\le L\), \(L|z|\le3\pi/2\). Absolute \(B\) bounds both filters
at \((q_\ell)_a\), \(q_\ell=(p+2\pi\ell)/L\), by \(B/(1+|\ell_a|)\) for
odd \(L\) and \(|\mathrm{Re}\,p_a|<5\pi/4\),
\(|\mathrm{Im}\,p_a|<\pi/4\). \emph{Status:} theorem. \emph{Lean:}
\href{https://github.com/yang-mills-group/OpenAI-YM-20261028/blob/3456e5e75cf8d4cd558955a5c7638fba2d68f464/YM/Next/FREE_KERNELS/Lines.lean\#L68-L76}{\texttt{plank\_bind}},
\href{https://github.com/yang-mills-group/OpenAI-YM-20261028/blob/3456e5e75cf8d4cd558955a5c7638fba2d68f464/YM/Next/FREE_KERNELS/Lines.lean\#L162-L196}{\texttt{principal\_plank}},
\href{https://github.com/yang-mills-group/OpenAI-YM-20261028/blob/3456e5e75cf8d4cd558955a5c7638fba2d68f464/YM/Next/FREE_KERNELS/Lines.lean\#L200-L245}{\texttt{alias\_rail}}.

\textbf{Proposition III.3.6 (curvature covariance).} Fix
\(\alpha,\eta>0\), \(\Omega\in\mathbb R\). A real 2-form kernel \(f\) is
firm if \(f\in\mathcal S(\alpha,\Omega)\), \(\hat f(0)=1\), and
\(\eta|\omega|^2\le\langle\omega,f_n\omega\rangle\le\Omega|\omega|^2\)
on every torus; symmetry is not assumed. On
\(U=\{p\ne0:|p_a|<3\pi/2\}\subset\mathbb R^4\), set
\[\Phi_f=\sigma_1^*\hat f\sigma_1+\sigma_0\sigma_0^*,\qquad G_f=\sigma_1\Phi_f^{-1}\sigma_1^*,\qquad G_\infty=s|p|^{-2}s^*,\]
where \(s\) wedges with \(ip\). There exist \(b>0\) and \(C_n\ge0\) for
every \(n\in\mathbb N\), depending only on \((\alpha,\Omega,\eta)\) and
additionally on \(n\) for \(C_n\), such that for every firm \(f\),
\(p\in U\) and complex 1-form \(x\),
\[\mathrm{Re}\langle x,\Phi_fx\rangle\ge b|p|^2|x|^2,\quad\|D^nG_f\|\le C_n|p|^{-n},\quad\|D^n(G_f-G_\infty)\|\le C_n|p|^{1-n}.\]
For scalar \(f\in\mathcal S(\alpha,\Omega)\), \(\hat f\in C^\infty\)
with \(\|D^n\hat f\|\le A_n\Omega\) everywhere, \(A_0=1\),
\(A_n=(n/\alpha)^n\) for \(n\ge1\). \emph{Status:} theorems.
\emph{Lean:}
\href{https://github.com/yang-mills-group/OpenAI-YM-20261028/blob/3456e5e75cf8d4cd558955a5c7638fba2d68f464/YM/Further/FREE_COVARIANCE_KP/KP/Core.lean\#L115-L125}{\texttt{tank\_full}},
\href{https://github.com/yang-mills-group/OpenAI-YM-20261028/blob/3456e5e75cf8d4cd558955a5c7638fba2d68f464/YM/Further/FREE_COVARIANCE_KP/KP/Alloy.lean\#L110-L120}{\texttt{gem\_quote}},
\href{https://github.com/yang-mills-group/OpenAI-YM-20261028/blob/3456e5e75cf8d4cd558955a5c7638fba2d68f464/YM/Further/FREE_COVARIANCE_KP/KP/Plain.lean\#L101-L133}{\texttt{radial\_prices}}.

\textbf{Proposition III.3.7 (nonlinear calculus).} \emph{Status:}
theorems.

\begin{enumerate}
\def\labelenumi{(\alph{enumi})}
\item
  For \(b,k\in\mathbb N\), \(k\le b\), there exists \(C=C(b,k)\ge1\). If
  \(L\ge1\), \(M\ge0\), and \(p:\mathbb Z^4\to E\), with \(E\) a real
  normed space, satisfies
  \(\|\nabla_up(x)\|\le ML^{-1-|u|}(1+L^{-1}|x|_1)^b\) for every \(x\)
  and \(|u|\le k\), then its degree-at-most-\(i\) Newton pieces
  \(P_i(x)=\sum_{|\nu|=i}\binom{x}{\nu}\nabla^\nu p(0)\) obey
  \(\|P_i(x)\|\le CML^{-1-i}(1+|x|_1)^b\) for \(i<k\); the remainder
  \(p-\sum_{i<k}P_i\) obeys the same bound with \(i=k\). \emph{Lean:}
  \href{https://github.com/yang-mills-group/OpenAI-YM-20261028/blob/3456e5e75cf8d4cd558955a5c7638fba2d68f464/YM/MultiNext/SIGNED_SOURCE_CONTINUATION/TownHall.lean\#L31-L90}{\texttt{capital}}.
\item
  Let \(d\in\mathbb N\) and \(f:E\to F\) be \(C^d\) at \(0\) between
  real normed spaces, with finitely many continuous linear
  \(p_v:E\to E\), \(\sum_vp_v=1\), integer weights \(Y_v\ge1\). If all
  derivatives through \(d\) of \(f(\sum_vt^{Y_v}p_vz)\) vanish at \(0\)
  for every \(z\), then \(D^mf(0)(x_1,\ldots,x_m)=0\) for all
  \(m\in\mathbb N\) whenever integers \(k_i\ge1\) satisfy
  \(\sum_i k_i\le d\) and \(p_vx_i=0\) for \(Y_v>k_i\). \emph{Lean:}
  \href{https://github.com/yang-mills-group/OpenAI-YM-20261028/blob/3456e5e75cf8d4cd558955a5c7638fba2d68f464/YM/MultiNext/SIGNED_SOURCE_CONTINUATION/Weighted.lean\#L120-L147}{\texttt{cancel}}.
\item
  Let \(b,d\in\mathbb N\), \(b\le d\), \(n\ge2\), \(o\in\Lambda_n\), and
  \(X\subset\Lambda_n\) finite with \(10\,\mathrm{dist}_\infty(x,o)<n\)
  for every \(x\in X\) (torus distance). Suppose
  \(u:\mathrm{Config}(G,\Lambda_n)\to\mathbb C\) is gauge invariant,
  depends only on links with both endpoints in \(X\), and \(u\circ\exp\)
  is \(C^b\) at \(0\). If \(u(U_{J,o}(t))\) has zero derivatives through
  \(b\) for every radial polynomial connection jet through weight \(d\)
  (degree \(r-1\) has weight \(r\ge2\)), then
  \(D^m(u\circ\exp)(0)(w_1,\ldots,w_m)=0\) for all \(m\in\mathbb N\),
  integers \(p_i\ge1\), \(\sum_i p_i\le b\), whenever the link field
  \(w_i\) is polynomial of degree \(\le p_i-1\) in centred initial-site
  coordinates on those links; other links are unrestricted. Here
  \(U_{J,o}(t)\) denotes based holonomies with weight-\(r\) components
  multiplied by \(t^r\). \emph{Lean:}
  \href{https://github.com/yang-mills-group/OpenAI-YM-20261028/blob/3456e5e75cf8d4cd558955a5c7638fba2d68f464/YM/MultiNext/SIGNED_SOURCE_CONTINUATION/PMOM.lean\#L43-L61}{\texttt{wiped}}.
\item
  For every \(j\in\mathbb N\), every family of homogeneous
  degree-\((j+1)\) polynomials \(Q_a:\mathbb R^4\to\mathfrak g\)
  (\(a=1,\ldots,4\)) decomposes as \(Q_a=K_a+\partial_ar\), where
  \(r=(j+2)^{-1}\sum_ax_aQ_a\) has degree \(j+2\) and \(K\) has degree
  \(j+1\) with \(\sum_ax_aK_a=0\). \emph{Lean:}
  \href{https://github.com/yang-mills-group/OpenAI-YM-20261028/blob/3456e5e75cf8d4cd558955a5c7638fba2d68f464/YM/MultiNext/SIGNED_SOURCE_CONTINUATION/Averaging.lean\#L84-L114}{\texttt{locale}}.
\item
  There exists \(C(G)>0\): if \(0\le b\le1\), \(\tau\ge C(G)b^2\), \(v\)
  is a Lie-algebra link field with \(|v_x|\le b\) on links within
  distance \(5\) of \(Y\), and plaquettes touching \(Y\) have chord
  distance \(\le\tau\) from \(1\) under \(\exp(v)\), those bounds
  persist under \(\exp(sv)\) for all \(0\le s\le1\). \emph{Lean:}
  \href{https://github.com/yang-mills-group/OpenAI-YM-20261028/blob/3456e5e75cf8d4cd558955a5c7638fba2d68f464/YM/MultiNext/SIGNED_SOURCE_CONTINUATION/Safeguard.lean\#L148-L177}{\texttt{strings}}.
\end{enumerate}

\subsubsection{2. What it uses}\label{what-it-uses-5}

\begin{itemize}
\tightlist
\item
  I.1: lattice, gauge group and signed symmetries; III.1: block
  averages, clips and trajectories; III.2: sheets and payment.
\item
  III.6: positivity, square roots and linear algebra F06--F09; III.10:
  repair and rigidity; III.16: Fourier kernels and weighted norms.
\item
  Fourier/Bloch decomposition, Schur complements, lattice Koszul
  exactness, Newton interpolation and radial gauge. The penalty
  \(\kappa>0\) comes from \(K\).
\end{itemize}

\subsubsection{3. Strategy}\label{strategy-5}

Solve the Gaussian block-spin problem by alias sums. Cancel the singular
principal-alias poles before inversion; smoothing and geometric dilation
along histories give uniform strip bounds. Amend inverse metrics away
from exact forms to make the defect positive, then construct the
isometry blocks without changing \(W,T\) or the quadratic forms on
potentials.

\subsubsection{4. Steps}\label{steps-5}

\begin{enumerate}
\def\labelenumi{\arabic{enumi}.}
\tightlist
\item
  History filters have uniform alias decay and a principal-alias lower
  bound; geometric dilation controls products of filters.
  \href{https://github.com/yang-mills-group/OpenAI-YM-20261028/blob/3456e5e75cf8d4cd558955a5c7638fba2d68f464/YM/Next/FREE_KERNELS/Train.lean\#L169-L184}{\texttt{Train.smooth\_upper}}.
\item
  Lagrange's identity gives the diagonal curvature metric from the
  auxiliary covariance; cancelling poles gives analytic invertible
  symbols.
  \href{https://github.com/yang-mills-group/OpenAI-YM-20261028/blob/3456e5e75cf8d4cd558955a5c7638fba2d68f464/YM/Next/FREE_KERNELS/Bend.lean\#L188-L194}{\texttt{Train.bend\_firm}}.
\item
  Alias algebra gives (D), including zero momentum.
  \href{https://github.com/yang-mills-group/OpenAI-YM-20261028/blob/3456e5e75cf8d4cd558955a5c7638fba2d68f464/YM/Next/FREE_KERNELS/Primal.lean\#L63-L118}{\texttt{Flight.tangent}}.
\item
  Koszul exactness factors the defect through \(d_2\); amendments
  replace its 3-form factor by \(C+\lambda-\mu S^*S\), with
  \(S^*S=O(L^{-2})\).
  \href{https://github.com/yang-mills-group/OpenAI-YM-20261028/blob/3456e5e75cf8d4cd558955a5c7638fba2d68f464/YM/Next/FREE_KERNELS/Eject.lean\#L144-L181}{\texttt{Flight.rebased}}.
\item
  Fix the catalogue, impose the credit bounds and apply F06--F09 to
  construct (R); periodisation limits extend identities to all tori.
  \href{https://github.com/yang-mills-group/OpenAI-YM-20261028/blob/3456e5e75cf8d4cd558955a5c7638fba2d68f464/YM/Bid/Craft.lean\#L107-L166}{\texttt{Bid.pinning}},
  \href{https://github.com/yang-mills-group/OpenAI-YM-20261028/blob/3456e5e75cf8d4cd558955a5c7638fba2d68f464/YM/UVNext/PRICED_KERNELS/Echo.lean\#L105-L116}{\texttt{Echo.allAcross}}.
\item
  Coercivity and the resolvent identity give the covariance comparison;
  weighted Taylor expansion and radial decomposition give the nonlinear
  vanishing results.
  \href{https://github.com/yang-mills-group/OpenAI-YM-20261028/blob/3456e5e75cf8d4cd558955a5c7638fba2d68f464/YM/Further/FREE_COVARIANCE_KP/KP/Alloy.lean\#L110-L120}{\texttt{gem\_quote}},
  \href{https://github.com/yang-mills-group/OpenAI-YM-20261028/blob/3456e5e75cf8d4cd558955a5c7638fba2d68f464/YM/MultiNext/SIGNED_SOURCE_CONTINUATION/Routes.lean\#L63-L127}{\texttt{covers}}.
\end{enumerate}

\subsubsection{5. Where the difficulty
lies}\label{where-the-difficulty-lies-5}

Bounds must hold uniformly over arbitrary history lengths and block
sizes. Positivity is established at real momenta before thinning the
analytic strip. Harmonic forms require defect identities on all tori.
Repair constants depend on direct clip coefficients; cross-kernel decay
is measured on the coarse lattice.

\subsubsection{6. Notation}\label{notation-5}

\(\Lambda_n=(\mathbb Z/n\mathbb Z)^4\); \(d_r\) acts on real
\(r\)-cochains, \(d=d_1\), \(\sigma_r=\hat d_r\), and \(*\) is adjoint.
\(R\in\mathcal S(\alpha,\Omega)\) means, for every entry,
\(\sum_ye^{\alpha|y|_1}|R_{ab}(y)|\le\Omega\). \(\hat R,R_n\) denote
Fourier symbols and periodisations; aliases use centred representatives.
\(J(i)\) is the stage index and \(c_i\) the direct clip coefficient.
\(D^n\) denotes full Fréchet derivatives with multilinear operator norm;
\(\nabla_u\) denotes fine forward differences. The cochain inner product
counts cells. Histories comprise odd blocks \(L\ge3\), stage indices
\(J\ge4\) dropping by at most one; a flight \((j,t)\) adds one block to
an index-\(j\) history, giving index \(t\in\{j-1,j\}\). \(Q\) averages
potentials, \(W\) minimises the quadratic action, \(T\) transports
curvature, and \(M_f,M_n,O\) are the isometry blocks; \(C,S\) denote the
unamended 3-form defect and transport.
\href{https://github.com/yang-mills-group/OpenAI-YM-20261028/blob/3456e5e75cf8d4cd558955a5c7638fba2d68f464/YM/PREC/Stacks.lean\#L70-L75}{\texttt{Towers}}.

\subsubsection{7. What remains
unexplained}\label{what-remains-unexplained-5}

The iterated Feynman-gauge interpretation remains probable. Gauge and
harmonic zero modes, Gaussian determinants and integration-variable
changes require further treatment in III.4 and III.10. F06--F07 and
analytic coefficient extraction are inputs here. III.3.6 asserts no
position-space continuum limit or contact-term statement.

Detailed notes: the full account of this unit, with proof sketches for
every step, is on the
\href{https://github.com/scottnarmstrong/yang-mills-informalization/blob/main/outline/notes/III-3.md}{detailed
notes page}.

\subsection{III.4. Gaussian integration: change of Gaussian sector,
determinants, window
law}\label{iii.4.-gaussian-integration-change-of-gaussian-sector-determinants-window-law}

\textbf{Part:} III, Ultraviolet construction (renormalization group).
\textbf{Lean:} about 40 work packages. The main ones are
\href{https://github.com/yang-mills-group/OpenAI-YM-20261028/tree/3456e5e75cf8d4cd558955a5c7638fba2d68f464/YM/SectorNext/GAUSSIAN_SECTOR_CHANGE}{\texttt{YM/SectorNext/GAUSSIAN\_SECTOR\_CHANGE/}}
(exact change of variables, gap, windows),
\href{https://github.com/yang-mills-group/OpenAI-YM-20261028/tree/3456e5e75cf8d4cd558955a5c7638fba2d68f464/YM/AnalyticBridge/RELATIVE_LOGDET_TRACE}{\texttt{YM/AnalyticBridge/RELATIVE\_LOGDET\_TRACE/}}
and
\href{https://github.com/yang-mills-group/OpenAI-YM-20261028/tree/3456e5e75cf8d4cd558955a5c7638fba2d68f464/YM/PhysicalNext/AFTERNOON_INTERACTING_RANK_COMPARISON}{\texttt{YM/PhysicalNext/AFTERNOON\_INTERACTING\_RANK\_COMPARISON/}}
(determinants and the polymer form of a sector),
\href{https://github.com/yang-mills-group/OpenAI-YM-20261028/tree/3456e5e75cf8d4cd558955a5c7638fba2d68f464/YM/AnalyticBridge/ANNULUS_RATIO_INSERTIONS}{\texttt{YM/AnalyticBridge/ANNULUS\_RATIO\_INSERTIONS/}},
\href{https://github.com/yang-mills-group/OpenAI-YM-20261028/tree/3456e5e75cf8d4cd558955a5c7638fba2d68f464/YM/SectorTools/PHYSICAL_GAUSSIAN_INSERTION}{\texttt{YM/SectorTools/PHYSICAL\_GAUSSIAN\_INSERTION/}}
and
\href{https://github.com/yang-mills-group/OpenAI-YM-20261028/tree/3456e5e75cf8d4cd558955a5c7638fba2d68f464/YM/Primitives/CONNECTED_GAUSSIAN_CUMULANTS}{\texttt{YM/Primitives/CONNECTED\_GAUSSIAN\_CUMULANTS/}}
(Gaussian moments),
\href{https://github.com/yang-mills-group/OpenAI-YM-20261028/tree/3456e5e75cf8d4cd558955a5c7638fba2d68f464/YM/AnalyticBridge/SECTOR_AFFINE_JETS}{\texttt{SECTOR\_AFFINE\_JETS/}},
\href{https://github.com/yang-mills-group/OpenAI-YM-20261028/tree/3456e5e75cf8d4cd558955a5c7638fba2d68f464/YM/AnalyticBridge/SECTOR_COVARIANCE_JETS}{\texttt{SECTOR\_COVARIANCE\_JETS/}}
and
\href{https://github.com/yang-mills-group/OpenAI-YM-20261028/tree/3456e5e75cf8d4cd558955a5c7638fba2d68f464/YM/AnalyticBridge/SELECTED_ORDERED_KERNELS}{\texttt{SELECTED\_ORDERED\_KERNELS/}}
(jets),
\href{https://github.com/yang-mills-group/OpenAI-YM-20261028/tree/3456e5e75cf8d4cd558955a5c7638fba2d68f464/YM/AnalyticBridge/LOCAL_GAUSSIAN_LAWS}{\texttt{LOCAL\_GAUSSIAN\_LAWS/}},
\href{https://github.com/yang-mills-group/OpenAI-YM-20261028/tree/3456e5e75cf8d4cd558955a5c7638fba2d68f464/YM/AnalyticBridge/GAUSSIAN_FIELD_ALIGNMENT}{\texttt{GAUSSIAN\_FIELD\_ALIGNMENT/}},
\href{https://github.com/yang-mills-group/OpenAI-YM-20261028/tree/3456e5e75cf8d4cd558955a5c7638fba2d68f464/YM/Push/FREE_RESOLVENT}{\texttt{YM/Push/FREE\_RESOLVENT/}}
and
\href{https://github.com/yang-mills-group/OpenAI-YM-20261028/tree/3456e5e75cf8d4cd558955a5c7638fba2d68f464/YM/RG08/Gaussian}{\texttt{YM/RG08/Gaussian/}}
(locality, infinite-volume kernels),
\href{https://github.com/yang-mills-group/OpenAI-YM-20261028/tree/3456e5e75cf8d4cd558955a5c7638fba2d68f464/YM/SectorNext/LARGE_FIELD_REMAINDERS}{\texttt{YM/SectorNext/LARGE\_FIELD\_REMAINDERS/}},
and the constant-producing statements in \texttt{Next/STATIONARY},
\texttt{MergeProof} and \texttt{PhysicalNext/V9\_JOINT\_ENDPOINT\_LOAD}.
\textbf{Size:} 37,956 lines in 4,737 declarations of the main-theorem
closure.

\textbf{Conventions.} \(x=\hat\beta_m\) in applications;
\(R(x)=\lceil(\log x)^2\rceil\), \(t(x)=x^{-1/2}e^{(\log\log x)^2}\),
\(\beta_F(x,L)=x+s_G\log L\), \(t_L(x)=t(\beta_F(x,L))\),
\(g=\beta^{-1/2}\) and \(r_x=\lfloor R(x)/60\rfloor\). Write \(\|A\|_s\)
for the maximum weighted row and column sums with weight
\(e^{s\,\mathrm{dist}(i,j)}\), using coarse positions. Matrix
inequalities are quadratic-form inequalities. ``For large \(x\)'' means
beyond a threshold depending only on previously fixed constants.

\subsubsection{1. What this unit proves}\label{what-this-unit-proves-6}

\textbf{Setting.} One block-spin step \(S\) maps a fine torus
\(\Lambda_f\) (side \(f=Lc\), block size \(L\) odd) to a coarse torus
\(\Lambda_c\). The push-forward of the density at fixed coarse field
\(V\) is decomposed (III.1, III.9) into sectors. A sector is given by:

\begin{itemize}
\tightlist
\item
  a store \(A\): the large-field data, made of cost cells, aisles and
  flagged sites, with a partition of unity \(\mathrm{bill}_A\) given by
  an allocation \(D\);
\item
  a choice \(k\) of gauge-fixing tags;
\item
  a spanning forest \(F\) (tree gauge).
\end{itemize}

The integration variables are the free links of \(F\), with product Haar
measure. In each sector the unit splits the free links as follows
(\href{https://github.com/yang-mills-group/OpenAI-YM-20261028/blob/3456e5e75cf8d4cd558955a5c7638fba2d68f464/YM/SectorNext/GAUSSIAN_SECTOR_CHANGE/Shell.lean\#L38-L40}{\texttt{chill}}):

\begin{itemize}
\tightlist
\item
  \textbf{frozen links} \(\mathcal F_A\): a parent block of an endpoint
  lies within distance \(10R(x)\) of \(A\). These are kept as Haar
  variables \(Y\).
\item
  \textbf{active links} \(\mathcal A_A\): all others. These are written
  \(u_e=\exp(a_e)\,\nu_e\), with \(a_e\in\mathfrak g\) and
  \(\nu=\iota V\) the copy of \(V\) on the fine lattice
  (\href{https://github.com/yang-mills-group/OpenAI-YM-20261028/blob/3456e5e75cf8d4cd558955a5c7638fba2d68f464/YM/PCalc/Chart.lean\#L35-L36}{\texttt{replica}}).
\end{itemize}

We write \(n_a=|\mathcal A_A|\dim\mathfrak g\).

On selected stencils, the step energy contains a sum of affine squares
\(Q_Y(a)=\tfrac12\lvert Ma+s(Y)\rvert^2\)
(\href{https://github.com/yang-mills-group/OpenAI-YM-20261028/blob/3456e5e75cf8d4cd558955a5c7638fba2d68f464/YM/SectorNext/GAUSSIAN_SECTOR_CHANGE/Remainder.lean\#L37}{\texttt{guide}}):

\begin{itemize}
\tightlist
\item
  \(M\) is a linear map from active coordinates to rows of linearised
  plaquette curvatures and block constraints;
\item
  \(s(Y)\) depends on the frozen links.
\end{itemize}

From these the unit defines three objects:

\begin{itemize}
\tightlist
\item
  the \textbf{precision} \(W=M^{\top}M\)
  (\href{https://github.com/yang-mills-group/OpenAI-YM-20261028/blob/3456e5e75cf8d4cd558955a5c7638fba2d68f464/YM/SectorNext/GAUSSIAN_SECTOR_CHANGE/Remainder.lean\#L27}{\texttt{Rows.gram}});
\item
  the \textbf{covariance} \(C=W_{\mathcal W}\), the symmetrised windowed
  inverse of \(W\), with windows of coarse radius \(R(x)/20\)
  (\href{https://github.com/yang-mills-group/OpenAI-YM-20261028/blob/3456e5e75cf8d4cd558955a5c7638fba2d68f464/YM/SectorNext/GAUSSIAN_SECTOR_CHANGE/Chosen.lean\#L69-L70}{\texttt{fog}});
\item
  the \textbf{mean} \(m(Y)=-(W^{\mathrm{per}})^{\top}M^{\top}s(Y)\)
  (\href{https://github.com/yang-mills-group/OpenAI-YM-20261028/blob/3456e5e75cf8d4cd558955a5c7638fba2d68f464/YM/SectorNext/GAUSSIAN_SECTOR_CHANGE/Chosen.lean\#L72-L74}{\texttt{meanLog}}).
  This is a window approximation to the minimiser
  \(-W^{-1}M^{\top}s(Y)\); its gradient need not vanish.
\end{itemize}

\textbf{Proposition III.4.1 (exact Gaussian change of variables in one
sector).} Fix the following data:

\begin{itemize}
\tightlist
\item
  a step \(S\) and a tube \(T\) satisfying the funding conditions of a
  teller: the prescribed kernel and stage, the multiscale metric bounds,
  the prescribed fine and coarse cutoffs, bounds on both path families,
  and compatible free metrics and towers;
\item
  a representation \(R\) of the fine density whose paths satisfy their
  bounds and whose activities satisfy the validity bounds, with
  \(\beta=\beta_R\);
\item
  \(x\), \(V\), \(k\), \(A\), \(F\), an allocation \(D\) and a cutoff
  profile \(\chi\);
\item
  a chart ratio \(J\), meaning:

  \begin{itemize}
  \tightlist
  \item
    \(J_{\rm sc}>0\) and \(\rho_J>0\), and \(J\) is measurable;
  \item
    \(J\) is real-analytic and positive on a ball
    \(B_{\rho_J}\subset\mathfrak g\);
  \item
    \(J(0)=1\);
  \item
    \(\exp\) is injective on \(B_{\rho_J}\);
  \item
    \(J_{\rm sc}\,J(a)\,da\) is mapped by \(\exp\) onto normalised Haar
    measure restricted to \(\exp B_{\rho_J}\);
  \item
    \(\exp B_{\rho_J}\) is open;
  \end{itemize}
\item
  a number \(d>0\).
\end{itemize}

Let \(C\) be a matrix on the active coordinates and \(m\) a map from
frozen arrays to active coordinates. Assume they are \emph{permitted}
(\href{https://github.com/yang-mills-group/OpenAI-YM-20261028/blob/3456e5e75cf8d4cd558955a5c7638fba2d68f464/YM/SectorNext/GAUSSIAN_SECTOR_CHANGE/Change.lean\#L95-L103}{\texttt{Permitted}}):

\begin{itemize}
\tightlist
\item
  \(C\succ0\), \(\beta>0\), \(m\) measurable, and \(4d<\rho_J\);
\item
  the \emph{standing condition}
  (\href{https://github.com/yang-mills-group/OpenAI-YM-20261028/blob/3456e5e75cf8d4cd558955a5c7638fba2d68f464/YM/SectorNext/GAUSSIAN_SECTOR_CHANGE/Shell.lean\#L51-L55}{\texttt{Standing}}):
  whenever \(\mathrm{bill}_A\ne0\) and the gauge tags are eligible,
  every active link equals \(e^{a_e}\nu_e\) with
  \(\lvert a_e\rvert\le d\).
\end{itemize}

Let
\(\Phi_{A,k}=\mathrm{bill}_A\cdot\mathrm{exposed}_R\cdot\mathrm{enters}_{T,\beta,k,V}\)
be the sector integrand. Here \(\mathrm{exposed}_R\) is the interacting
density of \(R\), and \(\mathrm{enters}\) is the tagged block-spin
kernel times \(e^{-\beta\cdot\text{gap energy}}\). Then the following
hold
(\href{https://github.com/yang-mills-group/OpenAI-YM-20261028/blob/3456e5e75cf8d4cd558955a5c7638fba2d68f464/YM/SectorNext/GAUSSIAN_SECTOR_CHANGE/Change.lean\#L127-L135}{\texttt{Licenses}}):

\begin{itemize}
\tightlist
\item
  \begin{enumerate}
  \def\labelenumi{(\alph{enumi})}
  \tightlist
  \item
    \(\Phi_{A,k}\) is integrable on the slice;
  \end{enumerate}
\item
  (b) for every \(Y\), \(\Psi(Y,\cdot)\in L^1(\gamma_C)\), and
  \(\Psi\in L^1(\mathrm{Haar}^{\mathcal F_A}\otimes\gamma_C)\);
\item
  (c) the identity
  \[\int\Phi_{A,k}\,d\mathrm{Haar}^{F}=\mathrm{pref}_S(\beta)\,J_{\rm sc}^{\,\lvert\mathcal A_A\rvert}\,(2\pi g^2)^{n_a/2}(\det C)^{1/2}\int\mathrm{Haar}^{\mathcal F_A}(dY)\int\gamma_C(dz)\,\Psi(Y,z).\]
\end{itemize}

Here \(\gamma_C=N(0,C)\). Set \(a=m(Y)+gz\), and let \(U\) be the
configuration with frozen links \(Y\), active links \(e^{a_e}\nu_e\) and
tree links \(1\). Then
(\href{https://github.com/yang-mills-group/OpenAI-YM-20261028/blob/3456e5e75cf8d4cd558955a5c7638fba2d68f464/YM/SectorNext/GAUSSIAN_SECTOR_CHANGE/Change.lean\#L63-L71}{\texttt{dyed}})
\[\Psi(Y,z)=\tilde J_d(a)\,\chi_d(a)\,\mathrm{bill}_A\,\mathrm{lik}_k(U)\,\mathrm{exposed}_R(U)\,\exp\!\Big(-\beta\,\overline{\mathrm{res}}(Y,U,a)-\beta Q_Y(m(Y))+\eta\cdot z+\tfrac12 z\cdot(C^{-1}-W)z\Big).\]
In this formula:

\begin{itemize}
\tightlist
\item
  \(\eta=-g^{-1}(Wm(Y)+M^{\top}s(Y))\);
\item
  \(\overline{\mathrm{res}}\) is the orientation-averaged step energy
  minus the affine squares;
\item
  \(\tilde J_d(a)\) is the product of the extended link Jacobians, each
  equal to \(J(a_e)\) for \(\lvert a_e\rvert\le 2d\); \(\chi_d\) is the
  product cutoff, equal to \(1\) when every \(\lvert a_e\rvert\le d\).
\end{itemize}

The identity holds for every permitted \((C,m)\) without quantitative
locality assumptions.

\emph{Status:} theorem.

\emph{Lean:}
\href{https://github.com/yang-mills-group/OpenAI-YM-20261028/blob/3456e5e75cf8d4cd558955a5c7638fba2d68f464/YM/SectorNext/GAUSSIAN_SECTOR_CHANGE/Change.lean\#L137-L192}{\texttt{authorized\_change}}.

\textbf{Proposition III.4.2 (licence of every committed RG step; the
transfer identity).} Let \(\mathscr R\) be a directed RG run (III.9)
whose delays satisfy the rate conditions and the per-stage chart
conditions
(\href{https://github.com/yang-mills-group/OpenAI-YM-20261028/blob/3456e5e75cf8d4cd558955a5c7638fba2d68f464/YM/SectorNext/GAUSSIAN_SECTOR_CHANGE/Handoff.lean\#L57-L71}{\texttt{Starcash}}).
Let \(m<N\), let \(p\) be a period, and let \(R\) be a representation at
level \(m+1\) on its selected history with
\(\beta_R\in[\hat\beta_{m+1}/2,\,2\hat\beta_{m+1}]\). Write:

\begin{itemize}
\tightlist
\item
  \(S\) for the step from level \(m+1\) to level \(m\);
\item
  \(x=\hat\beta_m\) and \(\beta=\beta_R\);
\item
  \(j=J(m+1)\) and \(i=J(m)\) for the stage indices;
\item
  \(d=a_{ji}\,t(\hat\beta_{m+1})\), and \(w\) for the allocation width
  at this stage pair.
\end{itemize}

Then:

\begin{itemize}
\tightlist
\item
  \(w>0\), \(d>0\), \(\beta>0\), \(E_{ji}>0\), \(2d<\rho_{\rm first}\)
  and \(4d<\rho_J\);
\item
  \(\lvert\log J(t)\rvert\le E_{ji}\lvert t\rvert^2\) for
  \(\lvert t\rvert\le4d\);
\item
  there is \(\kappa>0\) such that for all allocations \(D\) of width
  \(w\), forests \(H\) and profiles \(\chi\):

  \begin{itemize}
  \tightlist
  \item
    \begin{enumerate}
    \def\labelenumi{(\roman{enumi})}
    \tightlist
    \item
      writing \(I_A\) for the signed sector column, its contribution to
      the push-forward, and \(L_A\) for the nonnegative column with
      absolute fine density, after removing the selected packet energy
      and kernel prefactor, both have bounded Borel envelopes and
      \(|I_A(V)|\le\mathrm{pref}(\beta)e^{-\beta\kappa t_L(x)^2|A|}L_A(V)\)
      (\href{https://github.com/yang-mills-group/OpenAI-YM-20261028/blob/3456e5e75cf8d4cd558955a5c7638fba2d68f464/YM/Push/SECTOR_MEASURE/Presented.lean\#L151-L155}{\texttt{PortionCost}});
    \end{enumerate}
  \item
    (ii) for all \(A,k,V\), the chart is admitted
    (\href{https://github.com/yang-mills-group/OpenAI-YM-20261028/blob/3456e5e75cf8d4cd558955a5c7638fba2d68f464/YM/SectorNext/GAUSSIAN_SECTOR_CHANGE/Stamped.lean\#L36-L42}{\texttt{Admitted}});
    when \(\mathrm{flags}(A)=\mathrm{flags}_{t(x)}(V)\), also
    \(C=W_{\mathcal W}\succ0\), \(W\succ0\) and the selected predictor
    rows have analytic derivatives with a common bound; and the
    conclusions of III.4.1 hold with \(C=W_{\mathcal W}\) and the
    windowed mean for every store (for wrong flags trivially, since
    \(\mathrm{bill}_A\equiv0\));
  \item
    (iii) for every \(V\):
    \[\mathcal S_\beta(\rho_R)(V)=\sum_A\mathrm{sieve}(V)\sum_k\mathrm{pref}\,J_{\rm sc}^{\lvert\mathcal A_A\rvert}(2\pi g^2)^{n_a/2}(\det C_{A,k,V})^{1/2}\iint\Psi\,d\gamma_{C}\,d\mathrm{Haar}.\]
  \end{itemize}
\end{itemize}

Here \(\mathcal S_\beta\) is the block-spin push-forward, \(\rho_R\) is
the fine density, and
\(\mathrm{sieve}(V)=\exp(-\beta\cdot\text{orientation-averaged coarse energy of }V)\).

\emph{Status:} theorem, conditional on the run, history and band
hypotheses. III.4.12 supplies the rate and chart conditions; III.9
supplies the selected history.

\emph{Lean:}
\href{https://github.com/yang-mills-group/OpenAI-YM-20261028/blob/3456e5e75cf8d4cd558955a5c7638fba2d68f464/YM/SectorNext/GAUSSIAN_SECTOR_CHANGE/Handoff.lean\#L132-L182}{\texttt{committed}},
\href{https://github.com/yang-mills-group/OpenAI-YM-20261028/blob/3456e5e75cf8d4cd558955a5c7638fba2d68f464/YM/SectorNext/GAUSSIAN_SECTOR_CHANGE/Stamped.lean\#L124-L137}{\texttt{displaced\_wrong}}.

\textbf{Proposition III.4.3 (gap and locality of the actual precision
and covariance).} Fix \(G\), a patch \(P=(L,J)\) and a teller whose
kernel satisfies the branch condition. There are \(v,\eta,b,E>0\) and
\(K\ge0\) (post-\(L\) constants) with the following property. For large
\(x\), take any step with patch \(P\) and coarse side
\(c\ge100(1+R(x))\), any funded tube, and any \(A,F,k,V\) with
\(\mathrm{flags}(A)=\mathrm{flags}_{t(x)}(V)\). Then:

\begin{itemize}
\tightlist
\item
  \begin{enumerate}
  \def\labelenumi{(\roman{enumi})}
  \tightlist
  \item
    \(\eta\lvert a\rvert^2\le a\cdot Wa\) for all active \(a\)
    (\href{https://github.com/yang-mills-group/OpenAI-YM-20261028/blob/3456e5e75cf8d4cd558955a5c7638fba2d68f464/YM/SectorNext/GAUSSIAN_SECTOR_CHANGE/Gap.lean\#L242-L264}{\texttt{Gap.not\_flat}});
  \end{enumerate}
\item
  (ii) \(W\) has exponential off-diagonal decay in the Schur norm:
  \(\sup_i\sum_j e^{v\,\mathrm{dist}(i,j)}\lvert W_{ij}\rvert\le K\),
  and the same for columns, where \(\mathrm{dist}\) is the coarse
  distance between the parent blocks of the coordinates;
\item
  (iii) \(\lVert C-W^{-1}\rVert_0\le E\,e^{-br_x}\), and \(C\succ0\);
\item
  (iv) \(C\le K'\) for some post-\(L\) constant \(K'\)
  (\href{https://github.com/yang-mills-group/OpenAI-YM-20261028/blob/3456e5e75cf8d4cd558955a5c7638fba2d68f464/YM/AnalyticBridge/SHARP_NORMALIZED_TAIL/Apertures.lean\#L43-L55}{\texttt{open\_price}}).
\end{itemize}

The constants do not depend on the volume, the store, the forest or the
field.

\emph{Status:} theorem.

\emph{Lean:}
\href{https://github.com/yang-mills-group/OpenAI-YM-20261028/blob/3456e5e75cf8d4cd558955a5c7638fba2d68f464/YM/SectorNext/GAUSSIAN_SECTOR_CHANGE/Chosen.lean\#L102-L143}{\texttt{opaqueCost}}.

\textbf{Proposition III.4.4 (window law).} Let \(\sigma,\eta>0\),
\(K\ge0\), and fix a nonnegative function \(B:\mathbb R\to\mathbb R\).
There are \(u,E,g>0\) such that for every sufficiently large integer
\(r\) the following holds, uniformly in the torus and the finite
coordinate set. Let the coordinate positions \(p_j\) satisfy
\(\sum_j e^{-v\,\mathrm{dist}(i,p_j)}\le B(v)\) for every coarse site
\(i\) and every \(v>0\)
(\href{https://github.com/yang-mills-group/OpenAI-YM-20261028/blob/3456e5e75cf8d4cd558955a5c7638fba2d68f464/YM/SectorNext/GAUSSIAN_SECTOR_CHANGE/Windows.lean\#L125}{\texttt{Revenue}}).
Let \(A\) be a finite symmetric matrix with:

\begin{itemize}
\tightlist
\item
  Schur decay
  \(\lVert e^{\sigma\,\mathrm{dist}}A\rVert_{\rm Schur}\le K\);
\item
  spectral gap \(\eta\);
\item
  windows \(\mathcal W_i\) that contain the \(r\)-ball around \(i\).
\end{itemize}

Let \(L\) be the windowed inverse,
\(L_{\cdot i}=(A|_{\mathcal W_i})^{-1}e_i\) extended by zero, and
\(A_{\mathcal W}=\tfrac12(L+L^{\top})\)
(\href{https://github.com/yang-mills-group/OpenAI-YM-20261028/blob/3456e5e75cf8d4cd558955a5c7638fba2d68f464/YM/SectorNext/GAUSSIAN_SECTOR_CHANGE/Windows.lean\#L70}{\texttt{clouded}}).
Then:

\begin{itemize}
\tightlist
\item
  \(\lVert A_{\mathcal W}-A^{-1}\rVert_{\rm Schur}\le Ee^{-ur}\);
\item
  \(A_{\mathcal W}\ge g=\eta/(2(K+1)^2)\) and \(A_{\mathcal W}\succ0\).
\end{itemize}

A strengthened form
(\href{https://github.com/yang-mills-group/OpenAI-YM-20261028/blob/3456e5e75cf8d4cd558955a5c7638fba2d68f464/YM/AnalyticBridge/GAUSSIAN_FIELD_ALIGNMENT/Relative.lean\#L72-L118}{\texttt{same\_purchase}})
gives a rate \(s\in(0,\sigma]\) and a constant \(K_1\) such that \(A\),
\(A^{-1}\), \(L\) and \((A_{\mathcal W})^{-1}\) all have \(s\)-weighted
Schur norm \(\le K_1\), and
\(\lVert e^{s\,\mathrm{dist}}(L-A^{-1})\rVert_{\rm Schur}\le K_1e^{-sr}\).

\emph{Status:} theorem. The constants depend only on
\((\sigma,\eta,K,B)\).

\emph{Lean:}
\href{https://github.com/yang-mills-group/OpenAI-YM-20261028/blob/3456e5e75cf8d4cd558955a5c7638fba2d68f464/YM/SectorNext/GAUSSIAN_SECTOR_CHANGE/Windows.lean\#L184-L239}{\texttt{Schur.windows}}.

\textbf{Proposition III.4.5 (log-determinants as local path series).}
Let \(A\succ0\) with windows \(\mathcal W\) (\(i\in\mathcal W_i\)). Let
\(P_u\) be the column-wise windowed inverse of \(A+uI\), and
\(\Delta_u=(A+uI)P_u-I\). Assume the following bounds
(\href{https://github.com/yang-mills-group/OpenAI-YM-20261028/blob/3456e5e75cf8d4cd558955a5c7638fba2d68f464/YM/AnalyticBridge/RELATIVE_LOGDET_TRACE/Walls.lean\#L79-L89}{\texttt{Priced}}),
at a fixed nonnegative Schur rate, with \(C\ge1\) and \(0\le q<1\), for
every \(u\ge0\):

\begin{itemize}
\tightlist
\item
  \(\lVert A\rVert\le C\);
\item
  \(\lVert(A+u)^{-1}\rVert,\ \lVert P_u\rVert\le C/(1+u)\);
\item
  \(\lVert\Delta_u\rVert\le q/(1+u)\);
\item
  \(\lVert S\rVert\le q\), where \(S=AA_{\mathcal W}-1\).
\end{itemize}

Assume also \(A_{\mathcal W}\succ0\). Then
\[\log\det A=\sum_i\Big(\ell_i+\sum_{n\ge0}r_{i,n}\Big),\qquad \log\det(1+S)=\sum_i\sum_{n\ge0}\frac{(-1)^n(S^{n+1})_{ii}}{n+1},\qquad \log\det A_{\mathcal W}=-\log\det A+\log\det(1+S),\]
with:

\begin{itemize}
\tightlist
\item
  \(\ell_i=\int_0^\infty[(1+u)^{-1}-(P_u)_{ii}]\,du=(\log A|_{\mathcal W_i})_{ii}\),
  the single-window letter;
\item
  \(r_{i,n}=-\int_0^\infty(P_u(-\Delta_u)^{n+1})_{ii}\,du\) and
  \(\int\lvert\cdot\rvert\le Cq^{n+1}\).
\end{itemize}

All series converge absolutely. Moreover
(\href{https://github.com/yang-mills-group/OpenAI-YM-20261028/blob/3456e5e75cf8d4cd558955a5c7638fba2d68f464/YM/AnalyticBridge/RELATIVE_LOGDET_TRACE/Alphabet.lean\#L196-L208}{\texttt{labeled}}),
the matrix products expand into closed index paths. A term \(r_{i,n}\)
has one \(P_u\) factor and \(n+1\) defect factors, hence \(n+2\) jumps;
the corresponding trace-series term has \(n+1\) jumps. The absolute
values, summed over the intermediate indices, are bounded by
\(Cq^{n+1}\) and \(q^{n+1}/(n+1)\), respectively.

\emph{Status:} theorem.

\emph{Lean:}
\href{https://github.com/yang-mills-group/OpenAI-YM-20261028/blob/3456e5e75cf8d4cd558955a5c7638fba2d68f464/YM/AnalyticBridge/RELATIVE_LOGDET_TRACE/Series.lean\#L203-L227}{\texttt{Priced.chains}}.

For the relative expansions, the sector regularity conditions are: a
lawful representation; matching fine paths and cutoff; a local chart
with \(d>0\), \(4d<\rho_J\) and \(2d<\rho_{\rm first}\); positive
coupling and sector width; \(0<t_L(x)\le1\); output threshold \(t(x)\);
and the tube's geometric, local-readability, predictor-derivative,
covariance-positivity and spectral-gap conditions
(\href{https://github.com/yang-mills-group/OpenAI-YM-20261028/blob/3456e5e75cf8d4cd558955a5c7638fba2d68f464/YM/PhysicalNext/AFTERNOON_INTERACTING_RANK_COMPARISON/Opacity/Desk.lean\#L41-L54}{\texttt{Chips}},
\href{https://github.com/yang-mills-group/OpenAI-YM-20261028/blob/3456e5e75cf8d4cd558955a5c7638fba2d68f464/YM/PhysicalNext/AFTERNOON_INTERACTING_RANK_COMPARISON/Exposure.lean\#L40-L46}{\texttt{Boosted}}).

\textbf{Proposition III.4.6 (relative normalisation of a sector and its
localisation).} Let \(K_F\) be the flat kinetic precision: the Gram
matrix of the move at \(V=1\) with empty store, on all free links of
\(F\).

\begin{enumerate}
\def\labelenumi{(\roman{enumi})}
\item
  If the fine path family satisfies its path bounds and the kernel
  satisfies the patch's branch condition, \(\det K_F\) does not depend
  on the forest or on the gauge tags
  (\href{https://github.com/yang-mills-group/OpenAI-YM-20261028/blob/3456e5e75cf8d4cd558955a5c7638fba2d68f464/YM/SectorNext/GAUSSIAN_SECTOR_CHANGE/Bore.lean\#L303-L305}{\texttt{Native.shared}}).
\item
  Assume:
\end{enumerate}

\begin{itemize}
\tightlist
\item
  the tube is funded, and its fine path family and cutoff are those of
  \(R\);
\item
  the precision bounds of Lemma 4.21 of the detailed notes holds;
\item
  \(J\) is a local chart ratio, \(\beta>0\), and the threshold of the
  output budget point is \(t(x)\).
\end{itemize}

Then on every sector point \(p=(A,F,k,V,Y,z)\) with matching flags, the
positive density factor, with the sector masks and polymer product kept
outside it, factorises as
\[\text{freight}(p)=\text{loading}(V)\,e^{\Lambda(p)},\qquad \text{loading}(V)=\text{banker}(V)\,\mathrm{pref}(\beta)\,e^{v_R\lvert\Lambda_f\rvert}\,\mathrm{spare}(\beta)^{\lvert F^{\rm free}\rvert}(\det K_F)^{-1/2}>0.\]
Here
\(\text{freight}=J_{\rm sc}^{\lvert\mathcal A\rvert}(2\pi g^2)^{n_a/2}(\det C)^{1/2}e^{\frac12 z\cdot C^{-1}z}\,\tilde J\cdot(\text{noise})\cdot e^{\mathrm{exponent}_R}\),
and \(\mathrm{spare}(\beta)=(2\pi g^2)^{\dim\mathfrak g/2}J_{\rm sc}\)
is the Gaussian normalisation per link. The exponent is
\[\Lambda(p)=-\beta\,\mathcal E(p)+\eta\cdot z+\tfrac12z\cdot(C^{-1}-W)z+\kappa(p)+\sum_e\log\tilde J(a_e)+(\text{tails}),\]
\[\kappa(p)=-\lvert\mathcal F_A\rvert\log\mathrm{spare}(\beta)-\tfrac12\big(\log\det W-\log\det K_F\big)+\tfrac12\log\det(WC).\]
Here \(\mathcal E\) is the orientation-averaged row energy, including
the subtraction of the affine energy at \(a\) and its addition at \(m\):
\[\mathcal E=E_{\rm noise}(U,V)+E_{\rm fine}^{\rm cut}(U)-E_{\rm coarse}^{\rm cut}(V)-Q_Y(a)+Q_Y(m).\]
The tails are the remaining fine exponent after removing its volume
constant and restoring its cut energy, together with the coarse cutoff
correction.

\begin{enumerate}
\def\labelenumi{(\roman{enumi})}
\setcounter{enumi}{2}
\tightlist
\item
  Assume in addition the sector regularity conditions stated above, a
  nonnegative polymer exponent, and both path-counting bounds of Lemma
  4.21 of the detailed notes. Then
  \(\kappa(p)=-\lvert\mathcal F_A\rvert\log\mathrm{spare}+\sum_e\mathrm{sect}_e+I+\mathrm{Sh}\)
  (\href{https://github.com/yang-mills-group/OpenAI-YM-20261028/blob/3456e5e75cf8d4cd558955a5c7638fba2d68f464/YM/PhysicalNext/AFTERNOON_INTERACTING_RANK_COMPARISON/Opacity/Shards.lean\#L34-L39}{\texttt{book\_price}}),
  where:
\end{enumerate}

\begin{itemize}
\tightlist
\item
  \(\mathrm{sect}_e=-\tfrac12\sum_{\rm colours}[\ell_e(\text{actual})-\ell_e(\text{flat})]\),
  the difference of the single-window letters of III.4.5 for the
  identity-completed full precision at \((A,V)\) and at
  \((\emptyset,1)\);
\item
  \(I\) is the corresponding difference of path letters;
\item
  \(\mathrm{Sh}=\tfrac12\log\det(1+S)\).
\end{itemize}

\begin{enumerate}
\def\labelenumi{(\roman{enumi})}
\setcounter{enumi}{3}
\tightlist
\item
  Regrouping
  (\href{https://github.com/yang-mills-group/OpenAI-YM-20261028/blob/3456e5e75cf8d4cd558955a5c7638fba2d68f464/YM/PhysicalNext/AFTERNOON_INTERACTING_RANK_COMPARISON/Opacity/Enigma.lean\#L36-L48}{\texttt{rate\_pass}})
  gives \(\kappa=\sum_{o}\mathrm{rent}_o+\mathrm{far}+I+\mathrm{Sh}\):
\end{enumerate}

\begin{itemize}
\tightlist
\item
  \(o\) runs over the cores of \(A\) (the components formed by overlap
  of the \(500R(x)\)-halos of its seeds);
\item
  \(\mathrm{rent}_o\) is the sum of
  \(\mathrm{sect}_e-\log\mathrm{spare}\cdot1_{e\in\mathcal F_A}\) over
  links within \(200R(x)\) of \(o\);
\item
  \(\mathrm{far}\) is the same sum over links farther than \(200R(x)\)
  from \(A\).
\end{itemize}

\emph{Status:} theorem.

\emph{Lean:}
\href{https://github.com/yang-mills-group/OpenAI-YM-20261028/blob/3456e5e75cf8d4cd558955a5c7638fba2d68f464/YM/PhysicalNext/AFTERNOON_INTERACTING_RANK_COMPARISON/Opacity/Prism.lean\#L102-L168}{\texttt{Dyed.cardinal\_spin}},
\href{https://github.com/yang-mills-group/OpenAI-YM-20261028/blob/3456e5e75cf8d4cd558955a5c7638fba2d68f464/YM/PhysicalNext/AFTERNOON_INTERACTING_RANK_COMPARISON/Opacity/Prism.lean\#L50-L54}{\texttt{Dyed.keptEarly}},
\href{https://github.com/yang-mills-group/OpenAI-YM-20261028/blob/3456e5e75cf8d4cd558955a5c7638fba2d68f464/YM/PhysicalNext/AFTERNOON_INTERACTING_RANK_COMPARISON/Opacity/Prism.lean\#L56-L63}{\texttt{Dyed.spoolLift}}.

\textbf{Proposition III.4.7 (polymer form of a sector, with large-field
suppression).} Fix:

\begin{itemize}
\tightlist
\item
  a teller and a patch \(D\) satisfying the branch condition;
\item
  a determinant-jet fund and a local chart ratio \(J\);
\item
  a polymer exponent \(e\ge0\), a chart-width multiplier \(h\ge0\),
  constants \(a,\alpha,\tau>0\) with \(2\alpha\le a/2\), and bound
  constants \(C_0\ge0\), \(g_0,b_0,d_0>0\), \(\eta_0\in\mathbb R\).
\end{itemize}

Then for large \(x\) the following holds. Let \(y=\mathrm{retime}(x)\),
\(z=\beta_F(x,D.L)\) and \(\mathrm{hot}(s)=e^{(\log\log s)^{3/2}}\).
Eventually \(x\le y\le x^2\) and \(t(y)=x^{-1/2}\mathrm{hot}(x)\). Take
a step with patch \(D\), a sector datum with chart \(J\), and positive
cutoff parameter \(W_0\), subject to the following assumptions:

\begin{itemize}
\tightlist
\item
  the representation is lawful; the tube has the same fine paths and
  cutoff as the representation and satisfies the funded, geometric and
  chart-admission conditions; the chart radius \(d=h\,t_{D.L}(y)\)
  satisfies \(d>0\), \(4d<\rho_J\) and \(2d<\rho_{\rm first}\); the
  output threshold is \(t(y)\);
\item
  the step satisfies the free-metric size condition and
  \(c\ge100(1+R(y))\), and its fine and coarse path families have a
  common path bound;
\item
  \(\beta\in[z/2,2z]\); the fine threshold \(s=z^{-1/2}\mathrm{hot}(z)\)
  also lies in \([t_{D.L}(y)/2,2t_{D.L}(y)]\);
\item
  on each eligible packet the row energy \(\mathcal E\) is at least
  \(a s^2\) times its size, whenever the sector cutoffs are nonzero;
\item
  the support exponents satisfy the reserve inequality with parameters
  \((e,\tau)\); the budget radius \(B_0\) satisfies
  \(25R(y)+30\le B_0\le\mathrm{walls}(y)\), and the fine budget radius
  is at most \(2B_0\);
\item
  the lawful-representation bound can be chosen with activity constant
  \(\Gamma\le\beta_F(y,D.L)^{1/8}\) and smallness parameter
  \(\delta\le y^{-1/2}\), and the fine paths satisfy the required
  row-counting bound;
\item
  the determinant-jet bounds at halo \(20\), exponential order \(2\) and
  derivative order \(1\), and the precision and two path-counting bounds
  of Lemma 4.21 of the detailed notes hold at \(y\), with constants
  \((g_0,C_0,\eta_0,b_0,d_0)\);
\item
  uniformly over flag-matched sectors and frozen arrays, \(\|C\|_0\) is
  bounded by a common constant, every component of \(m\) and of
  \(Mm+s(Y)\) has absolute value at most \(1\), and
  \(|\eta_i|\le y^{-4}\); on the chart-cutoff support, \(|z_i|\le y^2\).
\end{itemize}

For such data there is a decomposition
(\href{https://github.com/yang-mills-group/OpenAI-YM-20261028/blob/3456e5e75cf8d4cd558955a5c7638fba2d68f464/YM/PhysicalNext/AFTERNOON_INTERACTING_RANK_COMPARISON/Opacity/Plotter.lean\#L19-L28}{\texttt{Plate}})
\[\Lambda(p)=\sum_{o\in\mathrm{cores}(A)}\mathrm{base}_o(p)+\sum_X\mathrm{aux}_X(p)\]
whenever the flags match and the product of the fine-field, block,
gauge-tag and chart cutoffs is nonzero. The bounds below hold for every
flag-matched \(p\), including points outside that support. In it:

\begin{itemize}
\tightlist
\item
  \(\mathrm{base}_o\) depends only on the data in the
  \(500R(y)\)-neighbourhood of \(o\);
\item
  \(\mathrm{aux}_X\) depends only on the data near the local support
  \(X\);
\item
  both are gauge-invariant and lattice-covariant;
\item
  \(\mathrm{base}_o\le-\alpha\,\mathrm{hot}(z)^2\lvert o\rvert\), where
  \(\lvert o\rvert\) counts cost cells, aisles and flags;
\item
  \(\lvert e^{\mathrm{aux}_X}-1\rvert\le f_X\), with \(f\ge0\)
  satisfying the polymer bound
  \(\sum_{X\ni y'}f_Xe^{e\,N_{B_0}(X)}\le y^{-1/48}\) at every site
  \(y'\), where \(N_{B_0}\) is the cover number at the output budget
  radius
  (\href{https://github.com/yang-mills-group/OpenAI-YM-20261028/blob/3456e5e75cf8d4cd558955a5c7638fba2d68f464/YM/Ac/Tree.lean\#L12-L13}{\texttt{Hit}}).
\end{itemize}

\emph{Status:} theorem.

\emph{Lean:}
\href{https://github.com/yang-mills-group/OpenAI-YM-20261028/blob/3456e5e75cf8d4cd558955a5c7638fba2d68f464/YM/PhysicalNext/AFTERNOON_INTERACTING_RANK_COMPARISON/Opacity/Gush.lean\#L96-L134}{\texttt{foreshadows}},
with hypotheses
\href{https://github.com/yang-mills-group/OpenAI-YM-20261028/blob/3456e5e75cf8d4cd558955a5c7638fba2d68f464/YM/PhysicalNext/AFTERNOON_INTERACTING_RANK_COMPARISON/Opacity/Gush.lean\#L82-L95}{\texttt{Helps}}.

\textbf{Proposition III.4.8 (moments under the sector law).} Fix \(G\).
For every covariance cap \(L>0\) and \(k\in\mathbb N\) there is
\(D\ge1\) with the following property. Take:

\begin{itemize}
\tightlist
\item
  any finite index set \(I\) and any probability space \((\Omega,\mu)\);
\item
  any \(C\succeq0\) on \(\mathbb R^{I\times\dim\mathfrak g}\) with
  \(C\le L\);
\item
  any finite family \((r_i)_{i\in Y}\) of measurable insertions with
  \(w_i,d_i\ge0\),
  \(\lvert r_i(\omega,z)\rvert\le w_ie^{d_i(1+\lvert z_{a_i}\rvert^2+\lvert z_{b_i}\rvert^2)}\),
  and \(\sum_{i:\,x\in\{a_i,b_i\}}d_i\le c=1/(24L)\) at every coordinate
  site \(x\);
\item
  any \(\rho,H\ge0\) and sets \(P,J\subset I\);
\item
  any measurable \(v\) with
  \(\lvert v\rvert\le H(1+\sum_{i\in P}\lvert z_i\rvert^2)^k\) that
  vanishes unless \(\lvert z_i\rvert\ge\rho\) for all \(i\in J\).
\end{itemize}

Then \(f=v\prod_{i\in Y}r_i\in L^1(\mu\otimes\gamma_C)\) and
\[\int\lvert f\rvert\,d(\mu\otimes\gamma_C)\le H\,D\,(1+\lvert P\rvert\dim\mathfrak g)^{2k+1}\,e^{-c\lvert J\rvert\rho^2}\,2^{\dim\mathfrak g\,\lvert J\rvert}\prod_{i\in Y}\big(e^{c}\,4^{\dim\mathfrak g}w_i\big).\]
In a sector, take \(\mu=\mathrm{Haar}^{\mathcal F_A}\) and
\(C=W_{\mathcal W}\); III.4.3 supplies a uniform covariance cap.
Moreover:

\begin{itemize}
\tightlist
\item
  (\href{https://github.com/yang-mills-group/OpenAI-YM-20261028/blob/3456e5e75cf8d4cd558955a5c7638fba2d68f464/YM/AnalyticBridge/ANNULUS_RATIO_INSERTIONS/Ember.lean\#L69-L111}{\texttt{r4}})
  For measurable \(q_i\) on the same product space, let
  \(\lvert q_i\rvert\le x_i^2(1+\lvert z_{a_i}\rvert^2+\lvert z_{b_i}\rvert^2)\)
  with \(0\le x_i\le1\) and \(\sum_{i\ni\text{site}}8x_i\le c\). Then
  \(\int(\prod_i(e^{q_i}-1))^4\le(\prod_iMx_i)^4\), with \(M=M(G,L)\).
\item
  (\href{https://github.com/yang-mills-group/OpenAI-YM-20261028/blob/3456e5e75cf8d4cd558955a5c7638fba2d68f464/YM/AnalyticBridge/ANNULUS_RATIO_INSERTIONS/Echo.lean\#L191-L223}{\texttt{oasis}})
  Fix a patch, a teller satisfying its branch condition and a covariance
  cap \(L>0\). For large \(x\), every funded step with that patch and
  \(c\ge100(1+R(x))\), matching flags, \(0<d\le1\),
  \(2d<\rho_{\rm first}\) and
  \(\beta\in[\beta_F(x,L_{\rm block})/2,2\beta_F(x,L_{\rm block})]\) has
  the fourth-moment insertion bounds above. The single-site exponents
  come from \(\eta\cdot z\) and the pair exponents from
  \(\tfrac12z\cdot(C^{-1}-W)z\). Their weights are \(De^{-uR(x)/480}\)
  for sites and exponentially decaying in distance for pairs; the
  constants are fixed before the threshold in \(x\), which may depend on
  the covariance cap.
\end{itemize}

\emph{Status:} theorem.

\emph{Lean:}
\href{https://github.com/yang-mills-group/OpenAI-YM-20261028/blob/3456e5e75cf8d4cd558955a5c7638fba2d68f464/YM/AnalyticBridge/ANNULUS_RATIO_INSERTIONS/Arc.lean\#L189-L253}{\texttt{same\_law}}.

\textbf{Proposition III.4.9 (Gaussian insertions supported on selected
coordinates).} Fix \(p\ge1\), a finite jet order, \(G\), a patch and a
teller satisfying the branch condition. There are post-\(L\) constants
\(K,M,s,b>0\) and \(B\ge0\) such that, for large \(x\), the following
holds. Take a funded step with that patch, \(c\ge100(1+R(x))\) and
matching flags. Fix any \(g,d,\theta>0\), frozen array and slot set
\(U\). Restrict the linear and quadratic correction to \(U\), obtaining
\(t\cdot z_U+\tfrac12z_U\cdot\mathcal Jz_U\), where \(t=\eta\) and
\(\mathcal J=C^{-1}-W\). Its tilted law is
\(N(\mathcal A t_U,\mathcal A)\), with
\(\mathcal A=(C^{-1}-\mathcal J_U)^{-1}\succ0\). The likelihood ratio of
the \(U\)-marginals satisfies
\[\mathbb E_{\gamma_{C,U}}\big[(dN_U/d\gamma_{C,U})^{p}\big]\le\exp\!\Big(2p\big(\lvert U\rvert\ell^2/(4\theta)+2K\lvert U\rvert(\theta+\varepsilon/2)\big)\Big)\]
whenever \(Kp(\theta+\varepsilon/2)\le1/4\). The parameters are:

\begin{itemize}
\tightlist
\item
  \(\varepsilon=Ke^{-bR(x)}\);
\item
  \(\ell=g^{-1}M^2e^{-sr_x}B\,(t(x)+d)\).
\end{itemize}

The restriction is on \(z\); coefficients may depend on all background
and frozen data.

\emph{Status:} theorem.

\emph{Lean:}
\href{https://github.com/yang-mills-group/OpenAI-YM-20261028/blob/3456e5e75cf8d4cd558955a5c7638fba2d68f464/YM/SectorTools/PHYSICAL_GAUSSIAN_INSERTION/Consumer.lean\#L167-L199}{\texttt{Physical.blocks}},
built on
\href{https://github.com/yang-mills-group/OpenAI-YM-20261028/blob/3456e5e75cf8d4cd558955a5c7638fba2d68f464/YM/SectorTools/PHYSICAL_GAUSSIAN_INSERTION/Insertion.lean\#L120-L172}{\texttt{insert\_letter}}.

\textbf{Proposition III.4.10 (jets of the block coefficients and of the
covariance).} Fix \(G\), a patch \(P\), a teller satisfying the branch
condition, and an order \(n\). There are \(r,d,v,\eta,\lambda_0,C_*>0\)
such that for large \(x\) the following holds. Take every funded step
with patch \(P\) and \(c\ge100(1+R(x))\), every \(A,F,V\) with matching
flags, every \(k\), and every logarithmic perturbation \(B\) of the
background with \(\lVert B\rVert\le r\). The store, forest and tags
remain fixed as \(B\) varies. Then:

\begin{itemize}
\item
  the move matrix \(M(B)\) is real-analytic;
\item
  \(M\), \(W=M^{\top}M\) and \(C\) satisfy weighted Schur bounds up to
  order \(n\). For every list \(\ell\) of at most \(n\) coordinate
  directions, with coarse positions \(p_b\), and output and input
  positions \(p_i,p_j\),
  \[\prod_{b\in\ell}e^{d\,\mathrm{dist}(p_b,p_i)}\sum_j e^{d\,\mathrm{dist}(p_i,p_j)}\lvert\partial_\ell X_{ij}(B)\rvert\le C_*,\]
  and the analogous column bound holds with the external weights
  anchored at \(p_j\);
\item
  \(W(B)\ge\eta I\), \(C(B)\ge\lambda_0 I\) and \(C(B)\succ0\);
\item
  for every \(u\ge0\), the full, principal and windowed resolvents have
  the same marked jet bounds with cap \(C_* /(1+u)\) at order zero and
  \(C_* /(1+u)^2\) at positive orders;
\item
  the window defect has marked jet cap \(C_*e^{-vR(x)}/(1+u)\), and the
  difference between the windowed and full resolvents has cap
  \(C_*e^{-vR(x)}/(1+u)^2\).
\end{itemize}

The selected predictor rows are real-analytic, with \(k\)-jets bounded
by constants depending only on \((G,P,K,k)\), on a ball of fixed radius
(\href{https://github.com/yang-mills-group/OpenAI-YM-20261028/blob/3456e5e75cf8d4cd558955a5c7638fba2d68f464/YM/AnalyticBridge/SECTOR_COVARIANCE_JETS/Sounds.lean\#L69-L101}{\texttt{sounds}},
\href{https://github.com/yang-mills-group/OpenAI-YM-20261028/blob/3456e5e75cf8d4cd558955a5c7638fba2d68f464/YM/AnalyticBridge/SECTOR_AFFINE_JETS/Branch.lean\#L140-L219}{\texttt{chosen\_branch}}).
The precision \(W\) is also real-analytic; covariance and resolvent
conclusions are finite-order jet bounds.

\emph{Status:} theorem.

\emph{Lean:}
\href{https://github.com/yang-mills-group/OpenAI-YM-20261028/blob/3456e5e75cf8d4cd558955a5c7638fba2d68f464/YM/AnalyticBridge/SECTOR_AFFINE_JETS/Covariance.lean\#L146-L210}{\texttt{assembly}},
\href{https://github.com/yang-mills-group/OpenAI-YM-20261028/blob/3456e5e75cf8d4cd558955a5c7638fba2d68f464/YM/AnalyticBridge/SECTOR_COVARIANCE_JETS/Sky.lean\#L62-L77}{\texttt{fromTeller}},
\href{https://github.com/yang-mills-group/OpenAI-YM-20261028/blob/3456e5e75cf8d4cd558955a5c7638fba2d68f464/YM/AnalyticBridge/SECTOR_COVARIANCE_JETS/Roof.lean\#L39-L62}{\texttt{coverall}},
\href{https://github.com/yang-mills-group/OpenAI-YM-20261028/blob/3456e5e75cf8d4cd558955a5c7638fba2d68f464/YM/PhysicalNext/PHYSICAL_JOINT_ENDPOINT_LOAD_F16/DLog/Fund.lean\#L120-L127}{\texttt{stores}}.

\textbf{Proposition III.4.11 (affine squares in the large-field
annulus).} Fix a teller, a block size \(L\), \(W_0>0\),
\(K\in\mathbb R\) and an integer halo \(D\ge0\). There are \(C\ge1\) and
an eventually nonnegative function \(e(x)\) with
\((1+R(x))^n\beta_F(x,L)e(x)\to0\) for every \(n\in\mathbb N\), such
that for large \(x\) the following holds. Take a funded step of block
size \(L\) with \(c\ge100(1+R(x))\), stores \(O,A\), a forest \(F\),
fields \(U,V\) and tags \(k\). Assume \(U\) is in tree gauge. In the
\((D+5)R(x)\)-halo of \(O\), assume the following local sector
conditions: fine plaquettes whose support meets that halo and which
belong to \(A\) have toll at least \(t_L(x)/4\), those outside \(A\)
have toll at most \(t_L(x)/2\); block bonds whose support meets that
halo and which belong to \(A\) have discrepancy at least \(W_0t_L(x)\),
those outside have discrepancy at most \(2W_0t_L(x)\); the coarse flags
match \(V\); and the chosen frames have positive mass. Then there is a
direction \(a\), vanishing on the tree links, such that for every pair
of windows \(H,Q\) with \(H_{\rm face}\subseteq Q_{\rm face}\) and
\(H_{\rm noise}\subseteq Q_{\rm noise}\), and with the face, noise and
coarse rows of \(Q\) in \(\{p\notin A^{(2R)},\ p\in O^{(DR)}\}\):

\begin{itemize}
\tightlist
\item
  \(Q\) is charted by \(a\) with constant \(C\);
\item
  the affine series on \(H\) converge, the required curvature and
  predictor derivatives exist with the stated linearisation, and
  \(\mathrm{work}_Q\ge\mathrm{squares}_H(a)-e(x)\lvert O\rvert\);
\item
  for every other \(a'\) charting \(H\) with constant \(K\), the same
  convergence, derivative statements and bound hold with \(a'\).
\end{itemize}

Here \(|O|\) is the number of cost cells, aisles and flags. The
comparison with \(a'\) uses uniqueness of the logarithm on the selected
stencils.

\emph{Status:} theorem.

\emph{Lean:}
\href{https://github.com/yang-mills-group/OpenAI-YM-20261028/blob/3456e5e75cf8d4cd558955a5c7638fba2d68f464/YM/SectorNext/LARGE_FIELD_REMAINDERS/Annulus.lean\#L36-L127}{\texttt{annulus\_pay}},
with
\href{https://github.com/yang-mills-group/OpenAI-YM-20261028/blob/3456e5e75cf8d4cd558955a5c7638fba2d68f464/YM/SectorNext/LARGE_FIELD_REMAINDERS/Clean.lean\#L110-L208}{\texttt{cleans\_at}}
and
\href{https://github.com/yang-mills-group/OpenAI-YM-20261028/blob/3456e5e75cf8d4cd558955a5c7638fba2d68f464/YM/SectorNext/LARGE_FIELD_REMAINDERS/CollarData.lean\#L131-L189}{\texttt{reacts}}.

\textbf{Proposition III.4.12 (production of the step constants).}
Assume:

\begin{itemize}
\tightlist
\item
  \(G\) is simple;
\item
  the kernel choices \(K\) satisfy the branch condition for every patch;
\item
  the initial stage is at least \(4\); the jet orders, position orders
  and support exponents are nondecreasing, the support exponents are
  positive, and at stage \(i\) the two orders are at least \(\max(6,i)\)
  and \(\max(2,i)\), respectively;
\item
  costs \(C_i\ge1\) are given.
\end{itemize}

Then there are pre-\(L\) data with kernel \(K\) and the prescribed
orders, two families of path prices, free-field and slope bounds, and a
cost bound, with the following property. For every tariff and scope
there is a sequence of block-size floors. After any increase of those
floors and any choice of block sizes above them, there is a stagewise
Hessian spectral gap \(\eta\). For every compatible post-\(L\) cost
choice (spectral gap \(\eta\), activity cap at least \(1\), and slope
bound below the chosen cap), protection and initial delays, the delays
can be enlarged so that the rate and chart conditions hold for all
further enlargements. On every paid run with these data there are
stationary path families, including the full stationary estimates in the
retained family. For each adjacent stage pair \((j,i)\) above the
initial stage, the per-step chart conditions hold with \(a_{ji},E_{ji}\)
chosen after \(L_j\)
(\href{https://github.com/yang-mills-group/OpenAI-YM-20261028/blob/3456e5e75cf8d4cd558955a5c7638fba2d68f464/YM/SectorNext/GAUSSIAN_SECTOR_CHANGE/Stamped.lean\#L139-L168}{\texttt{sponsored\_permit}}).
\href{https://github.com/yang-mills-group/OpenAI-YM-20261028/blob/3456e5e75cf8d4cd558955a5c7638fba2d68f464/YM/SectorNext/GAUSSIAN_SECTOR_CHANGE/Handoff.lean\#L99-L116}{\texttt{issue\_cosmos}}
replaces each entry value by the maximum of its old value and the
thresholds for \((j,j)\) and \((j,j-1)\), preserving the rate
conditions.

\emph{Status:} theorem under these kernel and order assumptions.

\emph{Lean:}
\href{https://github.com/yang-mills-group/OpenAI-YM-20261028/blob/3456e5e75cf8d4cd558955a5c7638fba2d68f464/YM/MergeProof/Seaside.lean\#L96-L154}{\texttt{sailsV}},
using
\href{https://github.com/yang-mills-group/OpenAI-YM-20261028/blob/3456e5e75cf8d4cd558955a5c7638fba2d68f464/YM/Next/STATIONARY/FullPlanned.lean\#L59-L66}{\texttt{geared}}.

\subsubsection{2. What it uses}\label{what-it-uses-6}

\begin{itemize}
\tightlist
\item
  \textbf{III.1:} block-spin maps, tree gauge, tags and predictors
  (\href{https://github.com/yang-mills-group/OpenAI-YM-20261028/blob/3456e5e75cf8d4cd558955a5c7638fba2d68f464/YM/RG/Block.lean\#L22-L25}{\texttt{StepData}});
  the branch condition requires unique selected minimisers and an
  eligible axial gauge.
\item
  \textbf{II.2, III.5, III.8:} representations, activity bounds,
  large-field allocations and polymer summability
  (\href{https://github.com/yang-mills-group/OpenAI-YM-20261028/blob/3456e5e75cf8d4cd558955a5c7638fba2d68f464/YM/RG/History.lean\#L181-L183}{\texttt{Rep.Lawful}}).
\item
  \textbf{III.10:} funded tubes, work rows, cutoffs, path bounds and
  compatible free metrics
  (\href{https://github.com/yang-mills-group/OpenAI-YM-20261028/blob/3456e5e75cf8d4cd558955a5c7638fba2d68f464/YM/Pressure/Seeds.lean\#L101-L109}{\texttt{Funds}}).
\item
  \textbf{III.3, III.6--III.7:} free-field estimates, stationary paths
  and fourth-order matching
  (\href{https://github.com/yang-mills-group/OpenAI-YM-20261028/blob/3456e5e75cf8d4cd558955a5c7638fba2d68f464/YM/Next/STATIONARY/Level/Recurse.lean\#L301-L305}{\texttt{FullPaths}}).
\item
  \textbf{III.2, III.9:} ordered constants, directed runs, selected
  histories and coupling bands.
\item
  \textbf{I.1, III.16:} Haar measure, exponential charts, analytic jets
  and Gaussian measures.
\item
  \textbf{Standard results:} Gaussian affine transport, Combes--Thomas,
  the resolvent integral for \(\log\), Neumann series, Gaussian
  integration by parts and analytic-strip Fourier decay.
\end{itemize}

Constants are chosen before \(L\), after \(L\), then before the
threshold in \(x\). Packet-energy, support-exponent, chart, noise and
representation assumptions in III.4.7 remain separate inputs.

\subsubsection{3. Strategy}\label{strategy-6}

Keep links near the store as Haar variables and put the remaining links
in exponential coordinates. Split off affine squares, shift by the
windowed mean and divide by a Gaussian density; the linear and quadratic
errors remain in the integrand. A local Poincaré inequality gives the
precision's spectral gap, and decay estimates make its windowed inverse
positive and exponentially close to the full inverse.

Expand relative log-determinants into window and path terms, then
collect terms near large-field cores. Packet energy absorbs their
logarithmic cost; the remainder has a small polymer norm. Gaussian
moments control exponential insertions, marked jets control background
derivatives, and annulus estimates recover affine squares away from the
store.

\subsubsection{4. Steps}\label{steps-6}

\begin{enumerate}
\def\labelenumi{\arabic{enumi}.}
\item
  \textbf{Chart transport and Gaussian scaling.} Under III.4.1's support
  and chart hypotheses, extended Jacobians and cutoffs preserve the
  sector integral and its integrability. Translate and dilate the
  Lebesgue variables, then expand the affine square to obtain the exact
  Gaussian formula.
  \href{https://github.com/yang-mills-group/OpenAI-YM-20261028/blob/3456e5e75cf8d4cd558955a5c7638fba2d68f464/YM/SectorNext/GAUSSIAN_SECTOR_CHANGE/Change.lean\#L137-L192}{Lean:
  \texttt{authorized\_change}}.
\item
  \textbf{Precision and windows.} Under III.4.3's hypotheses, a linkwise
  gauge choice compares predictor derivatives with the flat field, and
  the local Poincaré estimate absorbs the errors. Combes--Thomas bounds
  the window defect, giving III.4.4 and positivity.
  \href{https://github.com/yang-mills-group/OpenAI-YM-20261028/blob/3456e5e75cf8d4cd558955a5c7638fba2d68f464/YM/SectorNext/GAUSSIAN_SECTOR_CHANGE/Gap.lean\#L242-L264}{Lean:
  \texttt{Gap.not\_flat}},
  \href{https://github.com/yang-mills-group/OpenAI-YM-20261028/blob/3456e5e75cf8d4cd558955a5c7638fba2d68f464/YM/SectorNext/GAUSSIAN_SECTOR_CHANGE/Windows.lean\#L184-L239}{\texttt{Schur.windows}}.
\item
  \textbf{Admission along the run.} Under the funding and size
  hypotheses, axial gauge bounds active-link chords by a constant times
  \(t_L(x)\), so the logarithmic chart gives the standing condition.
  With III.4.2's run and band assumptions, enlarge stage thresholds,
  apply the sector identity to matching flags, and use vanishing for
  wrong flags to recover the push-forward.
  \href{https://github.com/yang-mills-group/OpenAI-YM-20261028/blob/3456e5e75cf8d4cd558955a5c7638fba2d68f464/YM/SectorNext/GAUSSIAN_SECTOR_CHANGE/Stamped.lean\#L45-L119}{Lean:
  \texttt{matinee}},
  \href{https://github.com/yang-mills-group/OpenAI-YM-20261028/blob/3456e5e75cf8d4cd558955a5c7638fba2d68f464/YM/SectorNext/GAUSSIAN_SECTOR_CHANGE/Handoff.lean\#L132-L182}{\texttt{committed}}.
\item
  \textbf{Relative determinants.} Identity completion puts actual and
  flat precisions on the same coordinates. Under III.4.5--III.4.6's
  assumptions, the resolvent integral and Neumann series give window
  letters and closed paths, which are regrouped near cores.
  \href{https://github.com/yang-mills-group/OpenAI-YM-20261028/blob/3456e5e75cf8d4cd558955a5c7638fba2d68f464/YM/AnalyticBridge/RELATIVE_LOGDET_TRACE/Series.lean\#L203-L227}{Lean:
  \texttt{Priced.chains}},
  \href{https://github.com/yang-mills-group/OpenAI-YM-20261028/blob/3456e5e75cf8d4cd558955a5c7638fba2d68f464/YM/PhysicalNext/AFTERNOON_INTERACTING_RANK_COMPARISON/Opacity/Enigma.lean\#L36-L48}{\texttt{rate\_pass}}.
\item
  \textbf{Core suppression.} Under III.4.7's hypotheses, packet energy
  dominates core normalisation costs; the remaining terms have powers
  \(y^{-1/8}\), \(y^{-1/24}\) or exponential decay. Clamping preserves
  core values on the cutoff support, and symmetrisation gives the
  \(y^{-1/48}\) bound.
  \href{https://github.com/yang-mills-group/OpenAI-YM-20261028/blob/3456e5e75cf8d4cd558955a5c7638fba2d68f464/YM/PhysicalNext/AFTERNOON_INTERACTING_RANK_COMPARISON/Opacity/Gush.lean\#L96-L134}{Lean:
  \texttt{foreshadows}}. \emph{(proof read; some supporting path and
  interaction bounds statement only)}
\item
  \textbf{Gaussian moments and local insertions.} Under
  III.4.8--III.4.9's assumptions, exponential quadratic moments and
  polynomial moments control insertions and tails. A correction
  supported on selected coordinates has a likelihood ratio determined by
  those marginals.
  \href{https://github.com/yang-mills-group/OpenAI-YM-20261028/blob/3456e5e75cf8d4cd558955a5c7638fba2d68f464/YM/AnalyticBridge/ANNULUS_RATIO_INSERTIONS/Arc.lean\#L189-L253}{Lean:
  \texttt{same\_law}},
  \href{https://github.com/yang-mills-group/OpenAI-YM-20261028/blob/3456e5e75cf8d4cd558955a5c7638fba2d68f464/YM/SectorTools/PHYSICAL_GAUSSIAN_INSERTION/Insertion.lean\#L120-L172}{\texttt{insert\_letter}}.
  \emph{(moment and insertion proofs read; marginal-support assertion
  statement only)}
\item
  \textbf{Jets.} Under III.4.10's hypotheses, transplant predictor
  stencils to a fixed torus and use a local gauge to obtain uniform
  analytic rows. Perturbation preserves the precision's spectral gap and
  supplies the weighted inverse and window jet bounds.
  \href{https://github.com/yang-mills-group/OpenAI-YM-20261028/blob/3456e5e75cf8d4cd558955a5c7638fba2d68f464/YM/AnalyticBridge/SECTOR_COVARIANCE_JETS/Sounds.lean\#L69-L101}{Lean:
  \texttt{sounds}},
  \href{https://github.com/yang-mills-group/OpenAI-YM-20261028/blob/3456e5e75cf8d4cd558955a5c7638fba2d68f464/YM/AnalyticBridge/SECTOR_AFFINE_JETS/Covariance.lean\#L146-L210}{\texttt{assembly}}.
  \emph{(predictor and precision proofs read; transport and cut-error
  bounds statement only)}
\item
  \textbf{Annulus squares and constants.} Under III.4.11's hypotheses,
  clean logarithmic coordinates give cubic row errors; summing the rows
  and discarding unselected positive squares proves the annulus bound.
  Under III.4.12's assumptions, choose stationary paths and block sizes,
  then enlarge delays and chart thresholds.
  \href{https://github.com/yang-mills-group/OpenAI-YM-20261028/blob/3456e5e75cf8d4cd558955a5c7638fba2d68f464/YM/SectorNext/LARGE_FIELD_REMAINDERS/Annulus.lean\#L36-L127}{Lean:
  \texttt{annulus\_pay}},
  \href{https://github.com/yang-mills-group/OpenAI-YM-20261028/blob/3456e5e75cf8d4cd558955a5c7638fba2d68f464/YM/MergeProof/Seaside.lean\#L96-L154}{\texttt{sailsV}}.
  \emph{(assembly and constant selection proofs read; clean-field
  construction and row recovery statement only)}
\end{enumerate}

\subsubsection{5. Where the difficulty
lies}\label{where-the-difficulty-lies-6}

The spectral gap requires a different gauge comparison for each link;
its smallness radius comes from continuity. Inverse bounds spend decay
rate while remaining uniform in volume and sector. Relative determinants
compare different active dimensions, and the large-field threshold must
absorb their core costs while preserving polymer summability.

\subsubsection{6. Notation}\label{notation-6}

{\def\LTcaptype{none} % do not increment counter
\begin{longtable}[]{@{}
  >{\raggedright\arraybackslash}p{(\linewidth - 4\tabcolsep) * \real{0.3333}}
  >{\raggedright\arraybackslash}p{(\linewidth - 4\tabcolsep) * \real{0.3333}}
  >{\raggedright\arraybackslash}p{(\linewidth - 4\tabcolsep) * \real{0.3333}}@{}}
\toprule\noalign{}
\begin{minipage}[b]{\linewidth}\raggedright
Symbol
\end{minipage} & \begin{minipage}[b]{\linewidth}\raggedright
Lean
\end{minipage} & \begin{minipage}[b]{\linewidth}\raggedright
Meaning
\end{minipage} \\
\midrule\noalign{}
\endhead
\bottomrule\noalign{}
\endlastfoot
\(\gamma_C\) &
\href{https://github.com/yang-mills-group/OpenAI-YM-20261028/blob/3456e5e75cf8d4cd558955a5c7638fba2d68f464/YM/SectorNext/GAUSSIAN_SECTOR_CHANGE/Gaussian.lean\#L31}{\texttt{Gaussian.cloud}}
& Centred Gaussian \(N(0,C)\). \\
\(\mathcal F_A,\mathcal A_A\) &
\href{https://github.com/yang-mills-group/OpenAI-YM-20261028/blob/3456e5e75cf8d4cd558955a5c7638fba2d68f464/YM/SectorNext/GAUSSIAN_SECTOR_CHANGE/Shell.lean\#L38-L40}{\texttt{chill}}
& Frozen and active free links. \\
\(W=M^{\top}M\) &
\href{https://github.com/yang-mills-group/OpenAI-YM-20261028/blob/3456e5e75cf8d4cd558955a5c7638fba2d68f464/YM/SectorNext/GAUSSIAN_SECTOR_CHANGE/Remainder.lean\#L27}{\texttt{Rows.gram}}
& Precision of the affine squares. \\
\(C=W_{\mathcal W}\) &
\href{https://github.com/yang-mills-group/OpenAI-YM-20261028/blob/3456e5e75cf8d4cd558955a5c7638fba2d68f464/YM/SectorNext/GAUSSIAN_SECTOR_CHANGE/Chosen.lean\#L69-L70}{\texttt{fog}}
& Symmetrised windowed inverse. \\
\(m(Y)\) &
\href{https://github.com/yang-mills-group/OpenAI-YM-20261028/blob/3456e5e75cf8d4cd558955a5c7638fba2d68f464/YM/SectorNext/GAUSSIAN_SECTOR_CHANGE/Chosen.lean\#L72-L74}{\texttt{meanLog}}
& Windowed mean. \\
\(\kappa(p),\Lambda(p)\) &
\href{https://github.com/yang-mills-group/OpenAI-YM-20261028/blob/3456e5e75cf8d4cd558955a5c7638fba2d68f464/YM/PhysicalNext/AFTERNOON_INTERACTING_RANK_COMPARISON/Opacity/Prism.lean\#L50-L54}{\texttt{keptEarly}}
& Relative log-normalisation and full log-integrand. \\
\end{longtable}
}

\(N_{B_0}(X)\) is the cover number; store size counts cost cells, aisles
and flags.

\subsubsection{7. What remains
unexplained}\label{what-remains-unexplained-6}

The simultaneous supply of III.4.7's packet-energy and support-exponent
assumptions on the selected run, and the transfer of stationary paths to
work rows, remain unexplained. Several supporting kernel, counting,
transport and clean-field proofs were read only as statements. The
free-kernel comparison is for flat sectors without frozen links under
compatibility assumptions; no direct comparison for arbitrary
interacting sectors was located. Downstream reset and repeatable-step
results and the additional assigned packages remain unchecked; precise
bibliographic references remain open.

Detailed notes: the full account of this unit, with proof sketches for
every step, is on the
\href{https://github.com/scottnarmstrong/yang-mills-informalization/blob/main/outline/notes/III-4.md}{detailed
notes page}.

\subsection{III.5. Cluster expansion of the polymer
gas}\label{iii.5.-cluster-expansion-of-the-polymer-gas}

\textbf{Part:} III, Ultraviolet construction (renormalization group).
\textbf{Lean:} work packages
\href{https://github.com/yang-mills-group/OpenAI-YM-20261028/tree/3456e5e75cf8d4cd558955a5c7638fba2d68f464/YM/Push/ABSOLUTE_SUMS}{\texttt{YM/Push/ABSOLUTE\_SUMS/}}
(3,624 lines),
\href{https://github.com/yang-mills-group/OpenAI-YM-20261028/tree/3456e5e75cf8d4cd558955a5c7638fba2d68f464/YM/AnalyticNext/CONNECTED_ACTIVITY_BOUNDS}{\texttt{YM/AnalyticNext/CONNECTED\_ACTIVITY\_BOUNDS/}}
(3,589),
\href{https://github.com/yang-mills-group/OpenAI-YM-20261028/tree/3456e5e75cf8d4cd558955a5c7638fba2d68f464/YM/Push/SECTOR_MEASURE}{\texttt{YM/Push/SECTOR\_MEASURE/}}
(1,153),
\href{https://github.com/yang-mills-group/OpenAI-YM-20261028/tree/3456e5e75cf8d4cd558955a5c7638fba2d68f464/YM/SectorNext/VACUUM_NORMALIZATION}{\texttt{YM/SectorNext/VACUUM\_NORMALIZATION/}}
(935),
\href{https://github.com/yang-mills-group/OpenAI-YM-20261028/tree/3456e5e75cf8d4cd558955a5c7638fba2d68f464/YM/Progress/STATIONARY_SCENES}{\texttt{YM/Progress/STATIONARY\_SCENES/}}
(445). \textbf{Size:} 9,746 lines in 1,097 declarations used by the main
theorem.

\subsubsection{1. What this unit proves}\label{what-this-unit-proves-7}

All propositions are \textbf{theorems} under the stated hypotheses.
Infinite nonnegative bounds include summability.

\textbf{Proposition III.5.1.} For finite labels, reflexive symmetric
incompatibility \(R\) and real \(w\), put
\(\ell_m=\sum_{z\in\iota^{m+1}}\phi_R^T(z)\prod_aw_{z_a}/(m+1)!\). If
\(|\ell_m|\le Cq^m\) for every \(m\ge0\), \(C\ge0\), \(0<q<1\), then
\(\Xi(w)=\exp(\sum_m\ell_m)\).
\href{https://github.com/yang-mills-group/OpenAI-YM-20261028/blob/3456e5e75cf8d4cd558955a5c7638fba2d68f464/YM/Push/ABSOLUTE_SUMS/FiniteLog.lean\#L66-L145}{\texttt{Gas.exact\_polynomial}}

\textbf{Proposition III.5.2.} On every finite graph with a distinguished
vertex, the absolute alternating sum over connected spanning subgraphs
is bounded by its number of spanning trees.
\href{https://github.com/yang-mills-group/OpenAI-YM-20261028/blob/3456e5e75cf8d4cd558955a5c7638fba2d68f464/YM/Push/ABSOLUTE_SUMS/Graph.lean\#L113-L148}{\texttt{Graph.pay}}

\textbf{Proposition III.5.3.} For arbitrary labels, supports
\(X_i\subseteq\Lambda_n\), \(B\in\mathbb N\), \(w\ge0\), \(a\ge0\),
\(R\subseteq\) overlap and \(\mathrm{Ready}(w;a,\gamma,\delta,q)\),
every site \(x\) and \(m\ge0\) satisfy
\[\sum_{k\ge m}\sum_{z\in\iota^{k+1}:x\in X(z)}\frac{|\phi_R^T(z)|\prod_aw_{z_a}}{(k+1)!}e^{a\mathrm{cov}_B(X(z))}\le\frac{\gamma q^m}{1-q}.\]
The same bound holds for connected finite sets through \(x\) of at least
\(m+1\) labels, weighted by \(\prod_iw_ie^{a\mathrm{cov}_B(X_i)}\). If
\(Q\subseteq\iota\) has total inflated weight
\(\sum_{i\in Q}w_ie^{(a+\delta)\mathrm{cov}_B(X_i)}\le b\), the
order-\(m\) cluster sum marked by \(Q\), with individual covering
weights, is \(\le bq^m\).
\href{https://github.com/yang-mills-group/OpenAI-YM-20261028/blob/3456e5e75cf8d4cd558955a5c7638fba2d68f464/YM/Push/ABSOLUTE_SUMS/Clusters.lean\#L224-L243}{\texttt{Ursell.log\_hit\_tail}}
\href{https://github.com/yang-mills-group/OpenAI-YM-20261028/blob/3456e5e75cf8d4cd558955a5c7638fba2d68f464/YM/Push/ABSOLUTE_SUMS/Clusters.lean\#L148-L183}{\texttt{Ursell.cluster\_at}}

\textbf{Proposition III.5.4.} For the same supports, real \(w\),
reflexive symmetric \(R\subseteq\) overlap, \(a\ge0\) and
\(\mathrm{Ready}(|w|;a,\gamma,\delta,q)\), the gas converges absolutely
with absolute sum \(\le e^{n^4\gamma}\); its cluster series converges
absolutely to \(\log\Xi(w)\), with absolute sum \(\le n^4\gamma/(1-q)\)
and order-\(m\) bound \(n^4\gamma q^m\). In particular \(\Xi(w)>0\).
\href{https://github.com/yang-mills-group/OpenAI-YM-20261028/blob/3456e5e75cf8d4cd558955a5c7638fba2d68f464/YM/Push/ABSOLUTE_SUMS/Log.lean\#L114-L149}{\texttt{Gas.price\_log}}

\textbf{Proposition III.5.5.} Under III.5.4, with nonempty supports and
overlap incompatibility, deleting all polymers meeting any
\(Y\subseteq\Lambda_n\) gives \(\Xi(w^Y)=\Xi(w)e^{-p}\). \(p\) sums
clusters meeting \(Y\); their absolute tail from order \(m\) is
\(\le|Y|\gamma q^m/(1-q)\) for every \(m\ge0\).
\href{https://github.com/yang-mills-group/OpenAI-YM-20261028/blob/3456e5e75cf8d4cd558955a5c7638fba2d68f464/YM/Push/ABSOLUTE_SUMS/Absence.lean\#L69-L95}{\texttt{Gas.restores\_log}}

\textbf{Proposition III.5.6.} On one lattice, let unforced \((X_i,u_i)\)
and forced \((Y_k,d_k)\) have nonempty supports, \(u,d\ge0\),
\(R\subseteq\) overlap. Assume \(A\ge0\), \(\delta'>0\), \(t\ge1\),
\(0\le\sigma<1\), \(\mu\ge0\), \(\eta>0\), \(q'<1\),
\(\mathrm{Ready}(u;A+\delta',\gamma,\delta,q)\),
\(\mathrm{Hit}_B(d;A+\delta',c)\), and
\[\frac{te^\mu c+t\gamma/(1-q)}{1-\sigma}\le\eta\le\sigma,\qquad4\mathrm{ball}_B\eta/\delta'\le q'.\]
Atoms are forced polymers and unforced cluster tuples with absolute
cluster weights. For every site \(x\) and \(m\in\mathbb N\), connected
atom multisets through \(x\), containing a forced atom and total
multiplicity \(\ge m\), have total weight
\(e^{A\mathrm{cov}_B(\cup)}\prod_\alpha w_\alpha^{n_\alpha}/n_\alpha!\)
at most \(t^{-m}e^{-\mu}\eta/(1-q')\).
\href{https://github.com/yang-mills-group/OpenAI-YM-20261028/blob/3456e5e75cf8d4cd558955a5c7638fba2d68f464/YM/Push/ABSOLUTE_SUMS/Omissions.lean\#L48-L100}{\texttt{OmitClusters.restore}}

\textbf{Proposition III.5.7.} For a step of block side \(L\),
\(l,B\in\mathbb N\), connected packets with
\(0\le p_A\le Ce^{(s-g)|A|}\), \(a,C\ge0\), and
\(\rho=K_{\rm step}^2e^{s-g+a(2l+3)^4}\le q<1\) satisfy, for every
coarse site \(x\) and \(m\in\mathbb N\),
\[\sum_{A:x\in D(A),\,|A|\ge m}p_Ae^{a\mathrm{cov}_B(D(A))}\le K_{\rm hit}Cq^{\max(1,m)}/(1-q).\]
This includes variants supported within domes whose total packet weight
obeys the same envelope.
\href{https://github.com/yang-mills-group/OpenAI-YM-20261028/blob/3456e5e75cf8d4cd558955a5c7638fba2d68f464/YM/Push/ABSOLUTE_SUMS/Charge.lean\#L66-L105}{\texttt{Packets.packet\_income}}
\href{https://github.com/yang-mills-group/OpenAI-YM-20261028/blob/3456e5e75cf8d4cd558955a5c7638fba2d68f464/YM/Push/ABSOLUTE_SUMS/Charge.lean\#L124-L175}{\texttt{Packets.packet\_variants\_tail}}

For \(\varepsilon,\delta,\gamma_{\rm root}>0\), \(a,C,e_s,\mu_s\ge0\),
assume the packet condition at
\((a+2\varepsilon+\delta,s+2\varepsilon e_s+\mu_s)\) with ratio
\(q_p<1\), optional polymers \(w_i,\lambda_i\ge0\) satisfying
\(\mathrm{Hit}_B(we^{2\varepsilon\lambda};a+2\varepsilon+\delta,\gamma_{\rm opt})\),
and \(K_{\rm hit}Cq_p/(1-q_p)+\gamma_{\rm opt}\le\gamma_{\rm root}\),
\(4\mathrm{ball}_B\gamma_{\rm root}/\delta\le q<1\). Connected mixed
sets through \(x\), of size \(\ge m+1\), seed count \(\ge r\) and charge
\(\mathrm{ch}\ge T\), satisfy
\[\sum_s\prod_{\alpha\in s}w_\alpha(1+\mathrm{ch}(s))^ke^{a\mathrm{cov}_B(D(s))}\le k!\varepsilon^{-k}e^\varepsilon e^{-\varepsilon T-\mu_sr}\gamma_{\rm root}q^m/(1-q)\]
for every \(m,k,r\in\mathbb N\), \(T\in\mathbb R\). Charge sums
individual covering numbers and lengths; packet length is \(e_s|A|\).
Disconnected packets have weight zero.
\href{https://github.com/yang-mills-group/OpenAI-YM-20261028/blob/3456e5e75cf8d4cd558955a5c7638fba2d68f464/YM/Push/ABSOLUTE_SUMS/Mixed.lean\#L153-L189}{\texttt{Mixed.mixed\_hit}}

\textbf{Proposition III.5.8.} Fix the group, \(L,l\), \(\mu,u>0\),
\(a,D\ge0\), \(d\in\mathbb N\), \(z<2\). Eventually in \(x\), whenever
\(s\le F\), \(g\ge M\), \(0\le C\le e^F\), where
\(F=D\mathrm{bill}^dP(\beta_F)^z\), \(M=\mu uP(\beta_F)^2\), the packet
ratio is \(\le q=e^{-M/2}\) and
\(K_{\rm hit}Cq^{\max(1,m)}/(1-q)\le e^{-M\max(1,m)/4}\) for every
\(m\). Every \(\beta\ge u\beta_F\) gives \(\beta\mu t^2\ge M\).
\href{https://github.com/yang-mills-group/OpenAI-YM-20261028/blob/3456e5e75cf8d4cd558955a5c7638fba2d68f464/YM/Push/ABSOLUTE_SUMS/Delay.lean\#L165-L214}{\texttt{Delay.fits}}

For \(\varepsilon,\delta>0\), \(\beta\ge u\beta_F\), \(C,e_s\ge0\),
\(C\le e^F\), \(s+2\varepsilon e_s\le F\),
\(0\le\gamma_{\rm opt}\le\delta/(16\mathrm{ball}_B)\), the mixed
numerical conditions eventually hold with \(g=\beta\mu t^2\),
\(\mu_s=M/2\), \(q=1/2\),
\(\gamma_{\rm root}=e^{-M/8}+\gamma_{\rm opt}\). Actual weight bounds
remain hypotheses. For every estimator, choose \(L\) sufficiently large,
then \(W\) sufficiently large: some \(\mu>0\), chosen before
\(l,d,u,a,D,z\), gives the point charge for all sufficiently large \(x\)
and funded steps of coarse side \(\ge100(1+B)\), uniformly in volume.
\href{https://github.com/yang-mills-group/OpenAI-YM-20261028/blob/3456e5e75cf8d4cd558955a5c7638fba2d68f464/YM/Push/ABSOLUTE_SUMS/Buy.lean\#L67-L103}{\texttt{Buy.quoted}}
\href{https://github.com/yang-mills-group/OpenAI-YM-20261028/blob/3456e5e75cf8d4cd558955a5c7638fba2d68f464/YM/Push/ABSOLUTE_SUMS/Grants.lean\#L52-L71}{\texttt{Grant.order}}

\textbf{Proposition III.5.9.} Let \(J\) be finite,
\(\emptyset\ne X_i\subseteq J\), \(v\ge0\), \(w\ge1\) submultiplicative
on overlapping unions, \(a\ge0\), \(0\le\delta\le1-e^{-a}\), and
\(\sum_{i:x\in X_i}w(X_i)v_ie^{a|X_i|}\le\delta\) for every \(x\). For
every real \(|u|\le v\), the disjoint-support gas satisfies
\(Z_u(B)>0\), \(Z_u(B)=e^{-\sum_{A\subseteq B}\lambda_u(A)}\),
\(|\lambda_u|\le\lambda_{-v}\) and
\(\sum_{A\ni x}w(A)\lambda_{-v}(A)\le a\). Deleting any \(Q\subseteq J\)
multiplies \(Z_u(J)\) by \(e^l\), \(|l|\le a|Q|\).
\href{https://github.com/yang-mills-group/OpenAI-YM-20261028/blob/3456e5e75cf8d4cd558955a5c7638fba2d68f464/YM/SectorNext/VACUUM_NORMALIZATION/Gas.lean\#L217-L228}{\texttt{ReLog.catalog}}

For countably many local, gauge-invariant Borel \(E_i\) bounded by
\(v_i\) on unflagged supports, mask flagged supports. If
\(M\in\mathbb N\), \(A,a\ge0\), \(v\ge0\), \(a(2M+1)^4\le1\),
\(\delta\le1-e^{-a}\) and
\(\sum_{i:x\in X_i}e^{(A+1)\mathrm{cov}_M(X_i)}v_i\le\delta\) at every
site, the masked gas has these identities and a field-independent
envelope \(|\lambda_P(V)|\le g(P)\),
\(\sum_{P\ni x}e^{A\mathrm{cov}_M(P)}g(P)\le a\). Raw coefficients are
local, gauge-invariant Borel; masked coefficients vanish on flagged
regions.
\href{https://github.com/yang-mills-group/OpenAI-YM-20261028/blob/3456e5e75cf8d4cd558955a5c7638fba2d68f464/YM/SectorNext/VACUUM_NORMALIZATION/Bulk.lean\#L157-L183}{\texttt{ReLog.intake}}

\textbf{Proposition III.5.10.} For finite real matrices \(A>0\),
\(C=C^T\), \(S=AC-I\), and \(\sum_j|S_{ij}|\le q<1\), \(q\ge0\), both
\(\det C,\det(I+S)>0\), and
\(\log\det(I+S)=\sum_{k\ge1}(-1)^{k+1}\mathrm{tr}(S^k)/k\) absolutely.
\href{https://github.com/yang-mills-group/OpenAI-YM-20261028/blob/3456e5e75cf8d4cd558955a5c7638fba2d68f464/YM/SectorNext/VACUUM_NORMALIZATION/Round.lean\#L149-L155}{\texttt{Round.exact\_trace}}

\textbf{Proposition III.5.11.} Fix a funded step, tube, estimator,
threshold \(t\), \(W,\kappa>0\), adapted lawful representation, coupling
\(\beta\ge0\), tree-gauge forest and Borel nonnegative allocation
summing to one. Require global cell chords \(\ge t/4\) on listed seeds,
\(\le t/2\) otherwise; wire chords \(\ge Wt\) on listed bonds,
\(\le2Wt\) otherwise; and equality of packet sites with coarse flags
wherever allocated weight is nonzero. Assume the point charge
\(\mathcal E_A\ge\kappa t^2|A|\) on every local sector on every
tree-gauge slice; local tests use \(D_{5B}(A)\) and positive frame mass.
Then, for every \(V\), \(\mathrm{push}_\beta\rho_R(V)=\sum_Ac_A(V)\) and
\(|c_A(V)|\le\pi_\beta e^{-\beta\kappa t^2|A|}\ell_A(V)\), with bounded
Borel sector outputs and remainders. The point charge comes from III.10
or the committed-scale hypothesis of III.1.
\href{https://github.com/yang-mills-group/OpenAI-YM-20261028/blob/3456e5e75cf8d4cd558955a5c7638fba2d68f464/YM/Push/SECTOR_MEASURE/Presented.lean\#L157-L210}{\texttt{Presented.portionCost}}

\textbf{Proposition III.5.12.} For finite symmetric \(M\ge rI>0\) with
lattice positions, weighted row and column sums \(\le C\) at rate
\(\eta\ge0\), \(0\le\theta\le1\), \(\theta C\le r/2\), every \(u\ge0\)
satisfies
\[e^{\theta\eta d(a_i,a_j)}|(M+uI)^{-1}_{ij}|\le[\min(1,r/2)(1+u)]^{-1}.\]
Hypotheses persist under principal restriction.
\href{https://github.com/yang-mills-group/OpenAI-YM-20261028/blob/3456e5e75cf8d4cd558955a5c7638fba2d68f464/YM/Progress/STATIONARY_SCENES/Green.lean\#L188-L251}{\texttt{Green.inverse}}

\textbf{Proposition III.5.13.} For finite real symmetric \(A\ge rI>0\),
fixed positive weighted row/column bounds at positive rate and counting
bound \(\mathrm{card}\{j:d(y,d_j)\le\rho\}\le U(1+\rho)^4\), windows
containing every index within distance \(R\) of their addressed index
have, uniformly for large \(R\), inverse residual and inverse difference
\(\le Ke^{-cR}<1\) for some \(K,c>0\) at a positive rate. Full/windowed
inverses are bounded by \(K\) throughout \(tA+(1-t)I\), \(0\le t\le1\).
At \(t=1\), the symmetrized windowed inverse has positive uniform lower
bounds, as does its inverse.
\href{https://github.com/yang-mills-group/OpenAI-YM-20261028/blob/3456e5e75cf8d4cd558955a5c7638fba2d68f464/YM/AnalyticNext/CONNECTED_ACTIVITY_BOUNDS/Birthday.lean\#L106-L178}{\texttt{Win.Promise.receive}}

\textbf{Proposition III.5.14.} Fix an estimator and patch satisfying
branch conditions. For all large \(x\), every funded step with that
patch, coarse side \(\ge100(1+B)\), forest, packet and field with
matching flags has completed precision
\(G=D^TD+I_{\rm inactive}\ge rI\), uniform counting constant \(U\), and
weighted row/column bound \(K\) at half the upper path-decay rate, for
some \(r,U,K>0\). Endpoint windows of radius \(\lfloor B/4\rfloor\)
satisfy III.5.13 with error
\(K_{\rm win}e^{-c_{\rm win}\lfloor B/4\rfloor}<1\).
\href{https://github.com/yang-mills-group/OpenAI-YM-20261028/blob/3456e5e75cf8d4cd558955a5c7638fba2d68f464/YM/AnalyticNext/CONNECTED_ACTIVITY_BOUNDS/Bedrock.lean\#L142-L158}{\texttt{Aff.bedrock}}
\href{https://github.com/yang-mills-group/OpenAI-YM-20261028/blob/3456e5e75cf8d4cd558955a5c7638fba2d68f464/YM/AnalyticNext/CONNECTED_ACTIVITY_BOUNDS/Warsaw.lean\#L25-L42}{\texttt{Aff.install\_promises}}

\subsubsection{2. What it uses}\label{what-it-uses-7}

\begin{itemize}
\tightlist
\item
  I.1, II.2: lattice covering bounds and connected supports.
\item
  III.1, III.2: step data, funding, branches, lawful representations and
  committed-scale energy assumptions.
\item
  III.10: energy grant; III.6: gauge-fixed Poincaré inequality.
\item
  Mayer expansion, Penrose partition, lattice-animal counting, Möbius
  inversion, Schur test, Combes--Thomas conjugation.
\end{itemize}

\subsubsection{3. Strategy}\label{strategy-7}

Component decomposition identifies the finite-gas logarithm. Tree
estimates give geometric cluster tails; absolute convergence extends the
identity to arbitrary labels. Seed energy dominates packet entropy.
Deletion gives another logarithm; exponential conjugation and the
spectral gap control windowed inverses.

\subsubsection{4. Steps}\label{steps-7}

\begin{enumerate}
\def\labelenumi{\arabic{enumi}.}
\tightlist
\item
  For gas polynomial \(P\) and cluster series \(L\), splitting the
  distinguished component gives \(P'=L'P\), hence \(Pe^{-L}=1\).
  \href{https://github.com/yang-mills-group/OpenAI-YM-20261028/blob/3456e5e75cf8d4cd558955a5c7638fba2d68f464/YM/Push/ABSOLUTE_SUMS/Coef.lean\#L129-L138}{\texttt{Coef.recurrence}}
\item
  Penrose intervals cancel free edges; parent factorials cost at most
  \(4^mm!\), absorbed by the reserve exponent.
  \href{https://github.com/yang-mills-group/OpenAI-YM-20261028/blob/3456e5e75cf8d4cd558955a5c7638fba2d68f464/YM/Push/ABSOLUTE_SUMS/Graph.lean\#L113-L148}{\texttt{Graph.pay}}
  \href{https://github.com/yang-mills-group/OpenAI-YM-20261028/blob/3456e5e75cf8d4cd558955a5c7638fba2d68f464/YM/Push/ABSOLUTE_SUMS/Weigh.lean\#L192-L228}{\texttt{Trees.trees\_afford}}
\item
  Count packets by spanning walks and compare their entropy with
  \(P(\beta_F)^2\).
  \href{https://github.com/yang-mills-group/OpenAI-YM-20261028/blob/3456e5e75cf8d4cd558955a5c7638fba2d68f464/YM/Push/ABSOLUTE_SUMS/Packet.lean\#L260-L269}{\texttt{Packets.stores\_count}}
  \href{https://github.com/yang-mills-group/OpenAI-YM-20261028/blob/3456e5e75cf8d4cd558955a5c7638fba2d68f464/YM/Push/ABSOLUTE_SUMS/Delay.lean\#L165-L214}{\texttt{Delay.fits}}
\item
  Deletion induction bounds Möbius logarithm pieces by negative-envelope
  pieces.
  \href{https://github.com/yang-mills-group/OpenAI-YM-20261028/blob/3456e5e75cf8d4cd558955a5c7638fba2d68f464/YM/SectorNext/VACUUM_NORMALIZATION/Deletion.lean\#L224-L233}{\texttt{ReLog.swept}}
\item
  Sum allocations, then extract the orientation-averaged point charge.
  \href{https://github.com/yang-mills-group/OpenAI-YM-20261028/blob/3456e5e75cf8d4cd558955a5c7638fba2d68f464/YM/Push/SECTOR_MEASURE/Oriented.lean\#L267-L297}{\texttt{Oriented.turn\_charge}}
\item
  Poincaré and background comparison give the spectral gap; conjugation
  gives decay, and window residuals occur beyond \(R\).
  \href{https://github.com/yang-mills-group/OpenAI-YM-20261028/blob/3456e5e75cf8d4cd558955a5c7638fba2d68f464/YM/AnalyticNext/CONNECTED_ACTIVITY_BOUNDS/Bedrock.lean\#L72-L113}{\texttt{Aff.surplus}}
  \href{https://github.com/yang-mills-group/OpenAI-YM-20261028/blob/3456e5e75cf8d4cd558955a5c7638fba2d68f464/YM/AnalyticNext/CONNECTED_ACTIVITY_BOUNDS/Windows.lean\#L146-L169}{\texttt{Windows.window\_price}}
\end{enumerate}

\subsubsection{5. Where the difficulty
lies}\label{where-the-difficulty-lies-7}

The bounds retain pins, covering weights and repeated labels. Global
seed tests survive orientations. The sector remainder has no
volume-uniform bound here.

\subsubsection{6. Notation}\label{notation-7}

\(\Lambda_n=(\mathbb Z/n\mathbb Z)^4\), \(n\ge1\); \(d\) is maximum
distance; \(\mathrm{cov}_B\) is connected ball-cover number;
\(\mathrm{ball}_B=(2B+1)^4\). \(\mathrm{Hit}_B(w;a,\gamma)\) means every
site sum of \(w_ie^{a\mathrm{cov}_B(X_i)}\) is \(\le\gamma\).
\(\mathrm{Ready}\) adds \(\gamma,\delta>0\),
\(4\mathrm{ball}_B\gamma/\delta\le q<1\), using exponent \(a+\delta\).

\(\Xi\) sums compatible finite sets; \(\phi^T\) is the connected
spanning alternating sum on the occurrence graph (edges require mutual
\(R\)); singleton coefficient \(1\); \(X(z)=\cup_aX_{z_a}\). Order \(m\)
means \(m+1\) occurrences. Seeds are fine cells, coarse bonds or sites;
packets are finite seed sets; \(D_r(A)\) is the radius-\(r\)
neighbourhood of seed locations (projected cell vertices, bond
endpoints, or site). Skins are individual radius-\(lB\) neighbourhoods;
\(D(A)\) is their union, connectedness their overlap connectivity.
\(\tau_L=L^4\mathrm{card}\{\text{plaquette orientations}\}+5\),
\(K_{\rm hit}=(2lB+3)^4\tau_L\), \(K_{\rm step}=(4lB+5)^4\tau_L\).

\(\beta_F=x+b\log L\), \(B=\lceil(\log x)^2\rceil\),
\(P(s)=e^{(\log\log s)^2}\), \(t=\beta_F^{-1/2}P(\beta_F)\),
\(\mathrm{bill}=1+B+(\log\beta_F)^2\). \(\lambda_u\) denotes Möbius
pieces of \(-\log Z_u\). \(\mathcal E_A\) assesses energy anchored in
\(D_{2B}(A)\) with truncated walks. \(\pi_\beta\) is the inverse product
of single-bond normalizations; \(\ell_A\) is the remaining absolute
sector integral. \(D\) is the transported flat row matrix, using
truncated curl/path rows and linearized block constraints, replica
transport and centred path gauges in every orientation. Keep row centres
outside \(D_{8B}(A)\) and nontrunk columns with projected endpoints
outside \(D_{10B}(A)\); normalize by the square root of the orientation
count. Its range is \(\lfloor B/100000\rfloor+R_{\rm pred}+2\). Weighted
sums at rate \(\eta\) use factors \(e^{\eta d(a_i,a_j)}\). A windowed
inverse uses column \(j\) of the principal inverse on \(P(j)\), extended
by zero.

\subsubsection{7. What remains
unexplained}\label{what-remains-unexplained-7}

Imported energy and Poincaré proofs, supporting counts and local
Gaussian comparisons remain inputs. Raw logarithm coefficients are
local; masked coefficients additionally inspect flags. The transported
flat matrix is not identified with the nonlinear derivative.

Detailed notes: the full account of this unit, with proof sketches for
every step, is on the
\href{https://github.com/scottnarmstrong/yang-mills-informalization/blob/main/outline/notes/III-5.md}{detailed
notes page}.

\subsection{III.6. The stationary branch: background field of the
effective
action}\label{iii.6.-the-stationary-branch-background-field-of-the-effective-action}

\textbf{Part:} III, Ultraviolet construction (renormalization group).
\textbf{Lean:}
\href{https://github.com/yang-mills-group/OpenAI-YM-20261028/blob/3456e5e75cf8d4cd558955a5c7638fba2d68f464/YM/Bid/Station.lean}{\texttt{YM/Bid/Station.lean}},
\href{https://github.com/yang-mills-group/OpenAI-YM-20261028/blob/3456e5e75cf8d4cd558955a5c7638fba2d68f464/YM/Bid/Craft.lean}{\texttt{YM/Bid/Craft.lean}},
\href{https://github.com/yang-mills-group/OpenAI-YM-20261028/tree/3456e5e75cf8d4cd558955a5c7638fba2d68f464/YM/Push/STATIONARY_UNIFORM}{\texttt{YM/Push/STATIONARY\_UNIFORM/}},
\href{https://github.com/yang-mills-group/OpenAI-YM-20261028/tree/3456e5e75cf8d4cd558955a5c7638fba2d68f464/YM/Planned/STATIONARY_JET_PORT}{\texttt{YM/Planned/STATIONARY\_JET\_PORT/}},
\href{https://github.com/yang-mills-group/OpenAI-YM-20261028/tree/3456e5e75cf8d4cd558955a5c7638fba2d68f464/YM/Extra/ITERATION_NO_LOSS}{\texttt{YM/Extra/ITERATION\_NO\_LOSS/}}.
\textbf{Size:} 21,329 lines in the main-theorem closure, including
attached packages.

\subsubsection{1. What this unit proves}\label{what-this-unit-proves-8}

\textbf{Setting.} \(G\) is the gauge group and \(\mathfrak g\) its Lie
algebra. In dimension four, \(f=Lc\), with odd \(L\ge3\) and patch order
\(J\ge4\). Sums use counting measure and physical colour coordinates. A
step \(S\) has kernel choices \(K\), penalty \(\kappa>0\), tagged
predictors \(Q_k\), and radius \(R_S=(J+7)L\). ``Larger'' means unique
selected minimizing frames and predictors, nonempty fallback axial
frames, \(100R_S\le f\), and \(c>1\). A forest \(H\) has one
block-centred rooted tree per block, height at most \(4L\);
\(H^{\rm fr}\) denotes free bonds.

Here \(x\) is a coarse \(\mathfrak g\)-valued bond field and \(a\) a
free-bond field. The axis field is \(U^H(x,a)=e^a\mathcal R(e^x)\) on
free bonds and \(1\) on tree bonds; \(\mathcal R\) copies coarse links
onto crossing fine bonds. With the symmetrized Wilson, row-square, cubic
and quartic main term \(\mathcal M_p\),
\[\mathcal T^H_k(x,a)=\mathcal M_p(W;U^H(x,a))+\kappa\sum_b\mathrm{price}\big(e^{x_b}Q_k(U^H(x,a))_b^{-1}\big),\qquad\mathrm{price}(U)=\|U-1\|^2/(2c_G^2).\]
A critical branch has \(\rho(0)=0\), is analytic at \(0\), and satisfies
\(D_a\mathcal T^H_k(x,\rho(x))=0\) near \(0\) for every tag, with every
tested action jointly analytic there. For a local exponential-chart Haar
density \(J_{\mathscr J}\), its thermal function is
\[\Theta^H_k(x)=\mathrm{loop}_W(U^H(x,\rho(x)))+\tfrac12\log\det\mathrm{Hess}_a\mathcal T^H_k(x,\rho(x))-\sum_{e\in H^{\rm fr}}\log J_{\mathscr J}(\rho(x)_e).\]

\textbf{Proposition III.6.1 (tested second jet; F03).} Let \(S\) be
larger, \(H\) a forest, \(B\) a bed with floor \((a,A)\) for clip \(p\)
on \(\Lambda_f\), and \(W\) have bounds \((p,\sigma,C,\Lambda)\), be
attached to \(B\), with \(B\) sworn. Assume main-term analyticity along
analytic families through the identity. For every tag \(k\),
\(\mathcal T^H_k\) is analytic at \(0\), its value and gradient vanish,
and for all \((x,a),(x',a')\),
\[D^2\mathcal T^H_k(0)[(x,a),(x',a')]=\langle ds,M_Bds'\rangle+\kappa\langle x-Ps,x'-Ps'\rangle,\qquad s=\iota x+a,\quad s'=\iota x'+a'.\]
Free-bond fields are extended by zero on tree bonds. \emph{Status:}
theorem. \emph{Lean:}
\href{https://github.com/yang-mills-group/OpenAI-YM-20261028/blob/3456e5e75cf8d4cd558955a5c7638fba2d68f464/YM/Bid/Station.lean\#L78-L81}{\texttt{Bid.F03}}.

\textbf{Proposition III.6.2 (main-term analyticity; F02).} If \(W\) has
bounds \((p,\sigma,C,\Lambda)\), \(X\) is a finite-dimensional real
normed space, \(x_0\in X\), and \(A(y)\) is a configuration family with
\(A(x_0)=1\) and link matrices analytic at \(x_0\), then
\(\mathcal M_p(W;A(y))\) is analytic at \(x_0\). \emph{Status:} theorem;
supplies III.6.1's analyticity assumption. \emph{Lean:}
\href{https://github.com/yang-mills-group/OpenAI-YM-20261028/blob/3456e5e75cf8d4cd558955a5c7638fba2d68f464/YM/Bid/Station.lean\#L73}{\texttt{Bid.F02}}.

\textbf{Proposition III.6.3 (vertical gap; F01).} Assume shell
comparison E02. For every \((G,L,J,K,a)\) with \(a>0\) there is
\(\eta>0\) such that, for all sides \(c,f\), larger \(S\) with these
parameters, forests \(H\), clips \(p\), beds \(B\) with floor \((a,A)\)
for clip \(p\) on \(\Lambda_f\) for some \(A\), and path data \(W\)
whose tested actions have the analyticity and jets of III.6.1,
\[D_a^2\mathcal T^H_k(0,0)[v,v]\ge\eta\sum_e|v_e|^2\] for every tag
\(k\) and free-bond field \(v\). One may take
\(\eta=\min(a,\kappa)/(1+\beta+\gamma)\), where
\(\beta=2C_{\rm E02}^2L^4\cdot4\cdot6|X|^2\), \(\gamma=2L^4\), and
\(X=\{-R_S,\ldots,R_S\}^4\). Thus \(\eta\) is independent of volume,
forest, \(B\) beyond its floor, and \(W\). \emph{Status:} theorem
conditional on E02, proved in III.10. \emph{Lean:}
\href{https://github.com/yang-mills-group/OpenAI-YM-20261028/blob/3456e5e75cf8d4cd558955a5c7638fba2d68f464/YM/Bid/Station.lean\#L68-L71}{\texttt{Bid.F01}}.

\textbf{Proposition III.6.4 (branch, tangent and value; F04).} Let \(S\)
be larger, \(H\) a forest, \(B,C\) beds, \(p\) a clip, \(W\) a path
datum, and \(T\) towers for \((S,B,C)\) satisfying diagonal identities.
Assume every tested action has the analyticity and jets of III.6.1, a
vertical gap \(a>0\), and a floor \((v,u)\) for \(C\) on \(\Lambda_c\)
for some clip \(q\) and real \(v,u\). There is a critical branch with
\[D\rho(0)=\mathrm{rel}_H,\qquad D^2[\mathcal T^H_k(\cdot,\rho(\cdot))](https://github.com/scottnarmstrong/yang-mills-informalization/blob/main/outline/units/0)[x,x']=\langle dx,M_Cdx'\rangle.\]
For every tag the stationary value is analytic with zero value and
gradient. \emph{Status:} theorem giving a germ. \emph{Lean:}
\href{https://github.com/yang-mills-group/OpenAI-YM-20261028/blob/3456e5e75cf8d4cd558955a5c7638fba2d68f464/YM/Bid/Station.lean\#L86-L90}{\texttt{Bid.F04}}.

\textbf{Proposition III.6.5 (forest independence; F05).} Let \(S\) be
larger, \(H,H'\) forests, \(W\) have bounds \((p,\sigma,C,\Lambda)\),
and \(\rho,\rho'\) be critical branches with positive gaps \(a,b\),
respectively. Given a local chart ratio \(\mathscr J\), for every tag
and \(x\) near \(0\) the stationary values and thermal functions agree
between forests. \emph{Status:} theorem; thermal equality has additive
constant zero. \emph{Lean:}
\href{https://github.com/yang-mills-group/OpenAI-YM-20261028/blob/3456e5e75cf8d4cd558955a5c7638fba2d68f464/YM/Bid/Station.lean\#L94-L99}{\texttt{Bid.F05}},
\href{https://github.com/yang-mills-group/OpenAI-YM-20261028/blob/3456e5e75cf8d4cd558955a5c7638fba2d68f464/YM/Work/F05.lean\#L42-L57}{\texttt{Work.F05.same\_thermal}}.

\textbf{Proposition III.6.6 (assembled branch).} For every
\((G,L,J,K,a)\) with \(a>0\) there is \(\eta>0\): for every larger \(S\)
with these parameters, forest \(H\), beds \(B,C\), clips \(p,q\), path
datum \(W\) and towers \(T\), assume diagonal identities for
\((S,B,C)\), floors \((a,A)\) for \((B,p)\) on \(\Lambda_f\) and
\((v,u)\) for \((C,q)\) on \(\Lambda_c\) for some real \(A,v,u\), bounds
\((p,\sigma,C_0,\Lambda)\) for \(W\), attachment to \(B\), and \(B\)
sworn. Then III.6.3 holds with \(\eta\) and there is a branch satisfying
III.6.4. \emph{Status:} theorem with no work-item assumptions.
\emph{Lean:}
\href{https://github.com/yang-mills-group/OpenAI-YM-20261028/blob/3456e5e75cf8d4cd558955a5c7638fba2d68f464/YM/Bid/Station.lean\#L102-L106}{\texttt{Bid.Scene}},
\href{https://github.com/yang-mills-group/OpenAI-YM-20261028/blob/3456e5e75cf8d4cd558955a5c7638fba2d68f464/YM/Integrated.lean\#L229-L232}{\texttt{Integrated.stationary\_scene}}.

\textbf{Proposition III.6.7 (volume-independent ball).} (a) Assign
finite horizontal and vertical coordinate sets \(\iota,\tau\) to torus
sites by \(\ell\), with multiplicity at most \(K_0\), using sup norms
and periodic max-distance. Fix \(C_0,a_0,r,m>0\), integer \(n\ge3\).
Suppose \(f:\mathbb R^{\iota\sqcup\tau}\to\mathbb R\) is analytic on
\(\|z\|_\infty<r\) and, for every such \(z\), \(1\le k\le n\),
coordinate tuple and \(1\le j\le k\),
\[|\partial_{q_1}\cdots\partial_{q_k}f(z)|\le C_0e^{-a_0\sum_i\mathrm{dist}(\ell(q_j),\ell(q_i))}.\]
If \(\nabla_\tau f(0)=0\) and \(D^2_{\tau\tau}f(0)\ge mI\), there are
\(0<r_{\rm hor}\le r_{\rm zone}\le r_{\rm safe}<r\), depending only on
\((K_0,C_0,a_0,r,m)\), and a map \(g:\mathbb R^\iota\to\mathbb R^\tau\)
with \(g(0)=0\) such that for \(\|x\|_\infty<r_{\rm hor}\), \(g\) is
analytic at \(x\), has norm at most \(r_{\rm zone}/2\), and is the
unique vertical critical point in the closed \(r_{\rm zone}\)-ball.

\begin{enumerate}
\def\labelenumi{(\alph{enumi})}
\setcounter{enumi}{1}
\tightlist
\item
  Fix \(G\), real \(\mu,C\), order \(J\), choices \(K\), and integer
  \(n\ge3\). Auxiliary constants can be chosen before \(L\) so that (a)
  applies to every clip \(p\), real \(\Lambda\), larger \(S\) of order
  \(J\) with choices \(K\), forest \(H\), and \(W\) with bounds
  \((p,\mu,C,\Lambda)\). One branch serves all tags; radii depend on the
  auxiliary constants, \(p,L\), but not \(c,f,H,\Lambda\) or \(W\)
  beyond its bounds. Coarse coordinates are located at fine block roots.
  \emph{Status:} theorem. \emph{Lean:}
  \href{https://github.com/yang-mills-group/OpenAI-YM-20261028/blob/3456e5e75cf8d4cd558955a5c7638fba2d68f464/YM/Push/STATIONARY_UNIFORM/ImplicitRadius.lean\#L179-L215}{\texttt{Receipt.throughout}},
  \href{https://github.com/yang-mills-group/OpenAI-YM-20261028/blob/3456e5e75cf8d4cd558955a5c7638fba2d68f464/YM/Push/STATIONARY_UNIFORM/Paths.lean\#L253-L294}{\texttt{Before.graph}}.
\end{enumerate}

\textbf{Proposition III.6.8 (localized derivatives).} For every integer
\(N\ge0\), assume III.6.7(b), choose auxiliary constants at order
\(N+3\), and take any branch satisfying its analytic, size and
uniqueness conclusions. For every tag, on the closed joint
\(r_{\rm safe}\)-ball the vertical Hessian is analytic, \(\ge m/2\),
with localized coefficient derivatives through \(N+1\). Every principal
sub-block shifted by \(tI\), \(t\ge0\), is invertible: its inverse is
exponentially localized with constant \(O((1+t)^{-1})\), has localized
derivatives through \(N+1\) and \(\ell^2\) norm at most
\((m/2+t)^{-1}\); its log determinant is analytic with localized
derivatives of orders \(1,\ldots,N+1\), bounded proportionally to
\((1+t)^{-1}\). On the horizontal ball, \(g\) and its graph have
localized derivatives of orders \(1,\ldots,N+2\), the classical action
of orders \(1,\ldots,N+3\), and the Hessian along the branch has the
same block package. Constants depend only on the corresponding bounds
and radii. No volume-independent value bound for the action or log
determinant is asserted. \emph{Status:} theorem. \emph{Lean:}
\href{https://github.com/yang-mills-group/OpenAI-YM-20261028/blob/3456e5e75cf8d4cd558955a5c7638fba2d68f464/YM/Planned/STATIONARY_JET_PORT/InUse.lean\#L116-L145}{\texttt{Counter.service}},
\href{https://github.com/yang-mills-group/OpenAI-YM-20261028/blob/3456e5e75cf8d4cd558955a5c7638fba2d68f464/YM/Planned/STATIONARY_JET_PORT/Pack.lean\#L65-L84}{\texttt{BlockPack}}.

\textbf{Proposition III.6.9 (one iteration level).} Assume III.2's
pre-\(L\) constants, warm-up and paid schedule; a horizon of beds
\(M_{N,m}\) and towers with prescribed strip/derivative bounds, floors
on every torus, diagonal/stack identities and identity diagonal at
cutoff
(\href{https://github.com/yang-mills-group/OpenAI-YM-20261028/blob/3456e5e75cf8d4cd558955a5c7638fba2d68f464/YM/PREC/Stacks.lean\#L132-L139}{\texttt{Suffices}});
sworn beds at the periods used; harmonic towers (zero-momentum two-form
coefficient independent of fine block residue); a levelwise gap quote
supplied by
\href{https://github.com/yang-mills-group/OpenAI-YM-20261028/blob/3456e5e75cf8d4cd558955a5c7638fba2d68f464/YM/Extra/ITERATION_NO_LOSS/Post.lean\#L99-L106}{\texttt{Iteration.post}};
a local chart ratio; and the
\href{https://github.com/yang-mills-group/OpenAI-YM-20261028/blob/3456e5e75cf8d4cd558955a5c7638fba2d68f464/YM/Extra/ITERATION_NO_LOSS/Gates.lean\#L79-L96}{\texttt{fresh-data\ hypotheses}}.
These supply bounded fresh cubic, quartic and loop atoms for every
admissible datum and family of branches satisfying (i) below,
classical/thermal matchings on chosen forests/tags per period (possibly
different choices), and cutoff cubic/quartic atoms.

For every \(m<N\) and path datum \(T\) funded at \(m+1\), there are
\(R\) funded at \(m\), branches \(\rho_{p,F}\) for every period and
forest, and \(\alpha\in\mathbb R\), such that:

\begin{itemize}
\tightlist
\item
  \begin{enumerate}
  \def\labelenumi{(\roman{enumi})}
  \tightlist
  \item
    Each branch satisfies III.6.6 with the level-\((m+1)\) gap, tower
    map and bed \(M_{N,m}\).
  \end{enumerate}
\item
  (ii) \(|\alpha+\kappa_G\log L_m|+1\le\mathrm{drift}_{m+1}\).
\item
  (iii) For every period, forest and tag, the stationary tested action
  for \(T\) with the level-\((m+1)\) clip has the same jets through
  order four as \(\mathcal M_q(R;e^x)\), with level-\(m\) clip \(q\).
\item
  (iv) Every thermal function is analytic with positive Hessian
  determinant near zero. For all tags and real parts \(w_1,w_2\) of
  complex plane waves with every axis/colour coefficient of modulus at
  most one,
  \[|D^2\Theta^F_k(0)[w_1,w_2]-D^2(\alpha\mathcal M_q^{(2)}(R)+\mathrm{loop}_R)(0)[w_1,w_2]|\le|\Lambda_c|\Gamma\lambda_{m+1}e^{-\gamma c}.\]
\end{itemize}

Here \(\Gamma,\gamma>0\) are chosen after \(L\); \(R,\alpha\) serve all
periods and forests. \(\mathcal M^{(2)}\) is the full symmetrized
Wilson-plus-row-square term, including higher-order dependence. There is
also a funded datum at cutoff \(N\) whose main term has Wilson-energy
jets through order four. \emph{Status:} theorem under these hypotheses.
\emph{Lean:}
\href{https://github.com/yang-mills-group/OpenAI-YM-20261028/blob/3456e5e75cf8d4cd558955a5c7638fba2d68f464/YM/Extra/ITERATION_NO_LOSS/Mechanism.lean\#L115-L150}{\texttt{Iteration.extend}},
\href{https://github.com/yang-mills-group/OpenAI-YM-20261028/blob/3456e5e75cf8d4cd558955a5c7638fba2d68f464/YM/Extra/ITERATION_NO_LOSS/Mechanism.lean\#L35-L48}{\texttt{Iteration.peak}}.

\subsubsection{2. What it uses}\label{what-it-uses-8}

\begin{itemize}
\tightlist
\item
  III.1: block-spin data, forest gauges, root-trivial gauge invariance,
  tower identities and Haar-preserving forest changes
  (\href{https://github.com/yang-mills-group/OpenAI-YM-20261028/blob/3456e5e75cf8d4cd558955a5c7638fba2d68f464/YM/PCalc/Profiles.lean\#L57-L77}{\texttt{Profiles.spending}}).
\item
  III.10: E02 shell comparison
  (\href{https://github.com/yang-mills-group/OpenAI-YM-20261028/blob/3456e5e75cf8d4cd558955a5c7638fba2d68f464/YM/Work/E02.lean\#L132-L244}{\texttt{Work.E02.answer}});
  III.3: free kernels; III.2: constant schedule; III.9: one-loop slope
  \(\kappa_G\).
\item
  I.1, II.1--II.2: lattice forms, group geometry, path bounds; III.16:
  analytic derivatives and jets.
\item
  Analytic implicit/inverse function theorems, Weierstrass normal
  convergence, change of variables, Banach contraction, Combes--Thomas
  estimates; Bałaban's background-field method.
\end{itemize}

\subsubsection{3. Strategy}\label{strategy-8}

Forest gauge removes blockwise gauge directions; curvature and the
predictor penalty control the remaining links. The vertical gap gives a
stationary germ, with tangent and quadratic value identified by the
towers. Under forest changes, the Hessian's squared Jacobian cancels the
Haar-density ratio in the thermal function.

Exponential inverse-Hessian decay supplies a volume-independent sup-norm
bound for Newton contraction. Differentiated stationarity and resolvent
identities control localized derivatives. The iteration joins
transported atoms to supplied fresh classical and one-loop data.

\subsubsection{4. Steps}\label{steps-8}

\begin{enumerate}
\def\labelenumi{\arabic{enumi}.}
\tightlist
\item
  \textbf{Analyticity and jets.} Exponential summability gives normally
  convergent holomorphic row/atom series; quadratic limits and the floor
  identity give III.6.1 after chart composition.
  \href{https://github.com/yang-mills-group/OpenAI-YM-20261028/blob/3456e5e75cf8d4cd558955a5c7638fba2d68f464/YM/Work/F02.lean\#L141-L143}{\texttt{Work.F02.answer}},
  \href{https://github.com/yang-mills-group/OpenAI-YM-20261028/blob/3456e5e75cf8d4cd558955a5c7638fba2d68f464/YM/Work/F03.lean\#L159-L175}{\texttt{Work.F03.answer}}.
  \emph{(proof read; metric and cubic/quartic limits: statement only)}
\item
  \textbf{Gap.} Linearize E02, square and sum link bounds, then use the
  floor and penalty to control vertical directions.
  \href{https://github.com/yang-mills-group/OpenAI-YM-20261028/blob/3456e5e75cf8d4cd558955a5c7638fba2d68f464/YM/Work/F01.lean\#L266-L311}{\texttt{Work.F01.answer}}.
\item
  \textbf{Branch and tangent.} Tags agree near identity; apply the
  analytic implicit function theorem, then differentiate stationarity
  and use
  \(D^2\mathcal T(0)[(y,b),(x,\mathrm{rel}_Hx)]=\langle dy,M_Cdx\rangle\)
  to identify the tangent and stationary quadratic value.
  \href{https://github.com/yang-mills-group/OpenAI-YM-20261028/blob/3456e5e75cf8d4cd558955a5c7638fba2d68f464/YM/Work/F04.lean\#L149-L159}{\texttt{Work.F04.answer}}.
\item
  \textbf{Forest comparison.} Transport critical graphs by root-trivial
  gauge transformations; Haar change of variables cancels the
  determinant Jacobian.
  \href{https://github.com/yang-mills-group/OpenAI-YM-20261028/blob/3456e5e75cf8d4cd558955a5c7638fba2d68f464/YM/Work/F05.lean\#L60-L67}{\texttt{Work.F05.answer}}.
\item
  \textbf{Uniform radius.} Combes--Thomas decay gives summable inverse
  rows; Newton contraction yields the ball, using a direct shell
  estimate and Wilson positivity without beds or attachment.
  \href{https://github.com/yang-mills-group/OpenAI-YM-20261028/blob/3456e5e75cf8d4cd558955a5c7638fba2d68f464/YM/Push/STATIONARY_UNIFORM/Resolvent.lean\#L171-L208}{\texttt{inverse\_sealed}},
  \href{https://github.com/yang-mills-group/OpenAI-YM-20261028/blob/3456e5e75cf8d4cd558955a5c7638fba2d68f464/YM/Push/STATIONARY_UNIFORM/ImplicitRadius.lean\#L179-L215}{\texttt{Receipt.throughout}},
  \href{https://github.com/yang-mills-group/OpenAI-YM-20261028/blob/3456e5e75cf8d4cd558955a5c7638fba2d68f464/YM/Push/STATIONARY_UNIFORM/Paths.lean\#L253-L294}{\texttt{Before.graph}}.
  \emph{(proof read; direct shell/action-derivative estimates: statement
  only)}
\item
  \textbf{Localized jets.} Differentiate \(D(M^{-1})=-M^{-1}(DM)M^{-1}\)
  and \(D\log\det M=\mathrm{tr}(M^{-1}DM)\), control the index sums by
  exponential convolution, and compose with the stationary graph.
  \href{https://github.com/yang-mills-group/OpenAI-YM-20261028/blob/3456e5e75cf8d4cd558955a5c7638fba2d68f464/YM/Planned/STATIONARY_JET_PORT/Lines.lean\#L56-L108}{\texttt{inverse\_lines}},
  \href{https://github.com/yang-mills-group/OpenAI-YM-20261028/blob/3456e5e75cf8d4cd558955a5c7638fba2d68f464/YM/Planned/STATIONARY_JET_PORT/InUse.lean\#L116-L145}{\texttt{Counter.service}}.
  \emph{(proof read for inverse induction and assembly; convolution and
  composition: statement only)}
\item
  \textbf{Iteration.} Join fresh atoms to inherited atoms, whose cubic
  and quartic charges carry factors \(L^{-2}\) and \(L^{-4}\); use the
  schedule for the drift bound. Match jets on reference forests and
  transfer the classical and thermal matchings to every forest and tag.
  \href{https://github.com/yang-mills-group/OpenAI-YM-20261028/blob/3456e5e75cf8d4cd558955a5c7638fba2d68f464/YM/Extra/ITERATION_NO_LOSS/Mechanism.lean\#L115-L150}{\texttt{Iteration.extend}}.
  \emph{(proof read for assembly; inherited-atom bounds and matching:
  statement only)}
\end{enumerate}

\subsubsection{5. Where the difficulty
lies}\label{where-the-difficulty-lies-8}

Curvature alone leaves flat and gauge directions. The gap is uniform in
volume; no lower bound uniform in \(L\) is asserted. A spectral gap
alone gives no uniform sup-norm radius: exponential summability is also
needed. The conclusions concern local stationarity and ball uniqueness,
without global constrained minimization.

\subsubsection{6. Notation}\label{notation-8}

\begin{itemize}
\tightlist
\item
  \(P\): predictor derivative; \(\iota\): linearized replica; \(d\):
  lattice exterior derivative.
\item
  A bed \(B\) consists of translation-invariant plaquette kernels, with
  metric \(M_B\), roots \(R_B\) and diagonal kernel \(D_B\)
  (\href{https://github.com/yang-mills-group/OpenAI-YM-20261028/blob/3456e5e75cf8d4cd558955a5c7638fba2d68f464/YM/PREC/Stacks.lean\#L35-L40}{\texttt{Bed}}).
  A floor \((a,A)\) includes \(M_B=c_pI+R_B^\top R_B\), metric/diagonal
  inverse identities, \(\widehat M_B(0)=I\), \(M_Bd=D_Bd\), symmetry,
  and quadratic bounds \(aI\le M_B,D_B\le AI\)
  (\href{https://github.com/yang-mills-group/OpenAI-YM-20261028/blob/3456e5e75cf8d4cd558955a5c7638fba2d68f464/YM/PREC/Stacks.lean\#L61-L68}{\texttt{Floor}}).
  Attachment identifies row coefficients, summed by offset and plaquette
  type, with \(R_B\); swornness means signed-coordinate-permutation
  invariance of \(\langle dA,M_BdA\rangle\)
  (\href{https://github.com/yang-mills-group/OpenAI-YM-20261028/blob/3456e5e75cf8d4cd558955a5c7638fba2d68f464/YM/PREC/Stacks.lean\#L44-L46}{\texttt{Attached}},
  \href{https://github.com/yang-mills-group/OpenAI-YM-20261028/blob/3456e5e75cf8d4cd558955a5c7638fba2d68f464/YM/Bid/Station.lean\#L54-L55}{\texttt{Sworn}}).
\item
  Towers satisfy
  \href{https://github.com/yang-mills-group/OpenAI-YM-20261028/blob/3456e5e75cf8d4cd558955a5c7638fba2d68f464/YM/PREC/Stacks.lean\#L82-L86}{\texttt{Towers.Diagonal}}.
  With \(K_B=d^\top M_Bd\) and linearized block average \(\bar Q\),
  their map obeys \(K_B\mathcal W=\bar Q^\top K_C\),
  \(\bar Q\mathcal W=1-\kappa^{-1}K_C\);
  \(\mathrm{rel}_H=(\widehat{\mathcal W}_H-\iota)|_{H^{\rm fr}}\), where
  \(\widehat{\mathcal W}_H\) is its forest-gauge transform.
\item
  A clip \(p\) carries \(c_p>0\). Path data \(W\) contain row series of
  transported plaquette variables and invariant loop, cubic and quartic
  atoms. Bounds \((p,\sigma,C,\Lambda)\) require exponential-length
  summability, row/cubic/quartic bounds \(C\), loop bound \(\Lambda\),
  row normalization \(c_p\delta_{ij}+\sum_k r_{ki}r_{kj}=\delta_{ij}\),
  and summable vanishing loop zeroth and first offset moments in every
  offset component, for every pair of plaquette types and colours
  (\href{https://github.com/yang-mills-group/OpenAI-YM-20261028/blob/3456e5e75cf8d4cd558955a5c7638fba2d68f464/YM/RG/PathActions.lean\#L173-L185}{\texttt{Bounds}}).
  Here \(r_{ij}\) sums row-\(i\) coefficients of plaquette type \(j\).
  ``Funded'' further specifies root-kernel rows on fixed coordinate
  paths, level bounds and weighted cubic/quartic caps
  (\href{https://github.com/yang-mills-group/OpenAI-YM-20261028/blob/3456e5e75cf8d4cd558955a5c7638fba2d68f464/YM/Extra/ITERATION_NO_LOSS/Gates.lean\#L35-L39}{\texttt{Work\ item}}).
\item
  A local chart ratio has positive measurable analytic density
  \(J_{\mathscr J}\), \(J_{\mathscr J}(0)=1\), on a positive-radius ball
  where \(\exp\) is injective with open image; a positive scale times
  this density pushes forward to probability Haar measure restricted to
  that image
  (\href{https://github.com/yang-mills-group/OpenAI-YM-20261028/blob/3456e5e75cf8d4cd558955a5c7638fba2d68f464/YM/PCalc/Loop.lean\#L28-L38}{\texttt{LocalRatio}}).
  The thermal Hessian uses the orthonormal colour basis.
\item
  Localization means
  \(|\partial_{q_1}\cdots\partial_{q_k}A_{ij}|\le C_ke^{-a_k(\mathrm{dist}(i,j)+\sum_\nu\mathrm{dist}(i,q_\nu))}\)
  using assigned sites; omit \(j\) for maps, and use III.6.7(a) for
  scalars.
\end{itemize}

\subsubsection{7. What remains
unexplained}\label{what-remains-unexplained-8}

The fresh classical/quantum matching assumptions remain undischarged
here:
\href{https://github.com/yang-mills-group/OpenAI-YM-20261028/blob/3456e5e75cf8d4cd558955a5c7638fba2d68f464/YM/Further/ITERATION_RECIPE_PRODUCER/Crown.lean\#L102-L112}{\texttt{IterationRecipe.stocked}}
retains
\href{https://github.com/yang-mills-group/OpenAI-YM-20261028/blob/3456e5e75cf8d4cd558955a5c7638fba2d68f464/YM/Further/ITERATION_RECIPE_PRODUCER/Crown.lean\#L83-L97}{\texttt{Heating}}.
Comparison of uniform radii with downstream small-field thresholds, the
all-level induction (III.12), and geometric realization of period
comparison remain unverified. F06--F09 and
primary-integral/atom-regularity supporting packages have statements
only; rotation/Ward material belongs to VI.1.

Detailed notes: the full account of this unit, with proof sketches for
every step, is on the
\href{https://github.com/scottnarmstrong/yang-mills-informalization/blob/main/outline/notes/III-6.md}{detailed
notes page}.

\subsection{III.7. Matching the Wilson energy to fourth
order}\label{iii.7.-matching-the-wilson-energy-to-fourth-order}

\textbf{Part:} III, Ultraviolet construction (renormalization group).
\textbf{Lean:}
\href{https://github.com/yang-mills-group/OpenAI-YM-20261028/tree/3456e5e75cf8d4cd558955a5c7638fba2d68f464/YM/Next/STATIONARY}{\texttt{YM/Next/STATIONARY/}},
\href{https://github.com/yang-mills-group/OpenAI-YM-20261028/tree/3456e5e75cf8d4cd558955a5c7638fba2d68f464/YM/WilsonBulk/Spectral}{\texttt{YM/WilsonBulk/Spectral/}},
\href{https://github.com/yang-mills-group/OpenAI-YM-20261028/tree/3456e5e75cf8d4cd558955a5c7638fba2d68f464/YM/UVExtension/Projection}{\texttt{YM/UVExtension/Projection/}},
and a few lines of
\href{https://github.com/yang-mills-group/OpenAI-YM-20261028/tree/3456e5e75cf8d4cd558955a5c7638fba2d68f464/YM/UVExtension/Reward}{\texttt{YM/UVExtension/Reward/}}.
\textbf{Size:} 19,945 lines in the main-theorem closure: 19,391 in
\texttt{Next/STATIONARY}, 286 in \texttt{UVExtension/Projection}, 215 in
\texttt{WilsonBulk/Spectral}, and about 50 elsewhere.

\subsubsection{1. What this unit proves}\label{what-this-unit-proves-9}

The leading action and one-loop term are represented by exponentially
weighted curvature paths. The leading action matches Wilson at the
finest level and the stationary tree action at each step through order
four; the one-loop coupling shift has slope \(\kappa=11C_G/(24\pi^2)\).

\textbf{Proposition III.7.1 (path data, beds, Fourier kernels).} A leg
transports the Lie-algebra component \(\mathcal F(u)\) of
\((u-u^{-1})/2\) from a plaquette along a lattice walk; its length is
the walk length plus four. A \(k\)-atom contracts \(k\) legs with a
coefficient tensor \(a\); it is invariant when the resulting polynomial
is \(\mathrm{Ad}\)-invariant, and its weight is
\(\|a\|_1e^{\sigma\,\mathrm{length}}\). A path datum \(R\) comprises
rows \(\ell_R(U)_{(x,P)}=\sum_k\lambda_{P,k}\,\mathrm{leg}_{P,k}(U,x)\)
for six plaquette orientations and sequences of two-, three- and
four-leg atoms. For a clip with \(c_{\rm dir}>0\) and absolutely
convergent series,
\[\mathrm{main}_{\mathbf c,R}=\frac1{384}\sum_s\left[c_{\rm dir}E_W+\tfrac12\sum_p|\ell_R(U)_p|^2+\sum_{x,k}(\mathcal A_k^{(3)}+\mathcal A_k^{(4)})(U,x)\right]_{U\to sU};\]
\(\mathrm{loop}_R\) similarly averages two-leg atoms over signed axis
permutations.

\(\mathrm{Bounds}(\mathbf c;\sigma,C,\Lambda)\) requires \(\sigma>0\),
weighted row sums and cubic/quartic sums at most \(C\), loop sum at most
\(\Lambda\), invariant atoms, normalization
\(c_{\rm dir}\delta_{PQ}+\sum_{P'}\hat\ell_{P'P}(0)\hat\ell_{P'Q}(0)=\delta_{PQ}\),
and summable loop coefficients with zero zeroth and first displacement
moments for each orientation-colour pair.

A bed \(B=(D,\tilde D,H,\tilde H,\rho)\) consists of
translation-invariant plaquette kernels; attachment means rows, summed
by endpoint plaquette and displacement, give \(\rho\). On \(\Lambda_n\),
\(\mathrm{Floor}(B,\mathbf c,n;\varepsilon,C)\) means
\(D\tilde D=H\tilde H=1\), \(\Phi^{\mathsf T}\Phi=H\) for
\(\Phi=(\sqrt{c_{\rm dir}}\,1;\rho)\), \(\hat H(0)=1\), \(Hd_1=Dd_1\),
and symmetric \(H,D\) between \(\varepsilon I\) and \(CI\). The strip
bound is \(\sum_y e^{\sigma|y|_1}|K(y)|\le C\) entrywise.

A kernel is strip-analytic of width \(\sigma\) and size \(K\) if it is
holomorphic, \(2\pi\)-periodic in every scalar momentum coordinate, and
bounded by \(K\) where \(|\mathrm{Im}\,p_{j,a}|<\sigma\). It records the
complexified coefficient \(D^rF(0)/r!\) with error \(\epsilon\) if the
coefficient on every tuple of canonical unit plane waves differs by at
most \(\epsilon\) from \(|\Lambda_n|\) times the kernel when momenta sum
to zero, and from zero otherwise. For two legs, \(\alpha_b\) is twice
the real \(s^2\)-coefficient of \(\hat K_{(b,b),(1,1)}(se_0)\); the bed
quadratic kernel has \(\alpha_b=1\) when its metric has positive strip
width and \(\hat H(0)=1\).

\emph{Status:} definitions; path specifications are from III.1.
\emph{Lean:}
\href{https://github.com/yang-mills-group/OpenAI-YM-20261028/blob/3456e5e75cf8d4cd558955a5c7638fba2d68f464/YM/RG/PathActions.lean\#L150-L154}{\texttt{RG.PathData}},
\href{https://github.com/yang-mills-group/OpenAI-YM-20261028/blob/3456e5e75cf8d4cd558955a5c7638fba2d68f464/YM/RG/PathActions.lean\#L173-L185}{\texttt{PathData.Bounds}},
\href{https://github.com/yang-mills-group/OpenAI-YM-20261028/blob/3456e5e75cf8d4cd558955a5c7638fba2d68f464/YM/PREC/Stacks.lean\#L61-L68}{\texttt{Floor}},
\href{https://github.com/yang-mills-group/OpenAI-YM-20261028/blob/3456e5e75cf8d4cd558955a5c7638fba2d68f464/YM/Next/STATIONARY/Lift/Pathways.lean\#L491-L492}{\texttt{Recorded}},
\href{https://github.com/yang-mills-group/OpenAI-YM-20261028/blob/3456e5e75cf8d4cd558955a5c7638fba2d68f464/YM/Next/STATIONARY/Films/LeadingRay.lean\#L699}{\texttt{alphaPay}}.

\textbf{Proposition III.7.2 (Wilson jet).} For every gauge group \(G\),
\(\sigma_0>0\) and \(A,Z\ge0\), there exist \(0<\sigma\le\sigma_0\) and
\(C\ge0\) such that: if \(n_i\to\infty\), \(c_{\rm dir}\le Z\),
\(\rho\in\mathrm{strip}(\sigma_0,A)\), Floor holds for every side with
some bounds \(a,b\), the bed gauge form is signed-permutation invariant
on each \(\Lambda_{n_i}\), and \(D=1\) there, then a datum \(Q\)
attached to \(B\) has straight-route rows, zero loops, and
\[Q\in\mathrm{Bounds}(\mathbf c;\tau,k,\Lambda)\quad(0<\tau\le\sigma,\ k\ge\max(C,6Ae^{4\sigma_0}),\ \Lambda\ge0).\]
On every \(\Lambda_{n_i}\), \(\mathrm{main}_{\mathbf c,Q}(e^x)\) and
\(E_W(e^x)\) are real-analytic at zero with equal derivatives of every
order \(j\le4\).

\emph{Status:} theorem. \emph{Lean:}
\href{https://github.com/yang-mills-group/OpenAI-YM-20261028/blob/3456e5e75cf8d4cd558955a5c7638fba2d68f464/YM/Next/STATIONARY/Shape/Scaled.lean\#L270-L305}{\texttt{Stationary.launch}}.

\textbf{Proposition III.7.3 (one-loop slope).} For simple \(G\),
\(\sigma,\varepsilon,\tau>0\) and \(C,Z,A\ge0\), there is \(D\ge0\) such
that, for every odd \(L\ge1\), \(n_i\to\infty\), clip
\(c_{\rm dir}\le Z\), bed with \(H\in\mathrm{strip}(\sigma,C)\), Floor
bounds \((\varepsilon,C)\) on every side and
signed-permutation-invariant gauge form on these tori, and every
attached \(R\in\mathrm{Bounds}(\mathbf c;\tau,A,\Lambda)\) with
\(\Lambda\in\mathbb R\), the function
\[\Pi_\mu(x)=\tfrac12\log\det(\mathcal H_{\mathbf c,R}(e^x)+d_Ud_U^*+\mathcal M_\mu)-\log\det(d_U^*d_U+\mathfrak m_\mu),\qquad\mu=1/L,\]
has a strip-analytic two-leg recording kernel \(\hat Y\) with errors
\(|\Lambda_{n_i}|e_i\), \(e_i\to0\) exponentially in \(n_i\) at an
\(L\)-dependent rate, and \(|\alpha_b(\hat Y)+\kappa\log L|\le D\) for
every colour \(b\). Here \(\mathcal H(U)\) is the Hessian at \(y=0\) of
\(\mathrm{main}_{\mathbf c,R}(e^yU)\),
\((d_U\phi)_e=\phi(e_-)-\mathrm{Ad}(U_e)\phi(e_+)\),
\(\mathfrak m_\mu=\mu^2[\mu^2(d_U^*d_U+\mu^2)^{-1}]^9\), and
\(\mathcal M_\mu\) applies it componentwise at initial endpoints.

\emph{Status:} theorem; \(D\) is independent of \(L\). \emph{Lean:}
\href{https://github.com/yang-mills-group/OpenAI-YM-20261028/blob/3456e5e75cf8d4cd558955a5c7638fba2d68f464/YM/Next/STATIONARY/FullPlanned.lean\#L17-L52}{\texttt{pouring}}.

\textbf{Proposition III.7.4 (one step and iteration).} Fix simple \(G\),
kernel parameters \(K\) with penalty \(a>0\), stage \(J\),
\(\varepsilon,w,\delta>0\) and \(D,Z,C,Z'\ge0\). There exist output
bounds \(A_*\ge0,s>0\); for every input \(\tau>0,A\ge0\) there exist
\(d\ge0,\gamma>0,u\) and functions \(\Theta_0(L),\Lambda_0(L)\), under
these hypotheses:

\begin{itemize}
\tightlist
\item
  Admissible steps \(\Lambda_{Lm_i}\to\Lambda_{m_i}\) use kernel \(K\)
  and one stage-\(J\) patch, \(L\ge u\), \(m_i\to\infty\), \(m_i>1\),
  and fine sides at least \(100\) times the patch extent.
\item
  Fine/coarse clips satisfy \(p_{\rm dir}\le Z,q_{\rm dir}\le Z'\). Fine
  \(H,\rho\in\mathrm{strip}(w,D)\) with Floor \((\varepsilon,D)\) on
  every side; coarse \(H_T,\rho_T\in\mathrm{strip}(\delta,C)\) with
  Floor \((\varepsilon',C)\) on every side for some real
  \(\varepsilon'\). Both gauge forms have signed-permutation symmetry on
  their tori.
\item
  A tower satisfies III.1's quadratic interpolation, curl-transport and
  diagonal-transport identities. Its interpolation kernel obeys strip
  bounds \((w,D/L^{1+|\nu|})\) for every fine-coordinate difference list
  \(|\nu|\le2\).
\item
  Fine \(R\in\mathrm{Bounds}(\mathbf p;\tau,A,\Lambda)\) is attached to
  \(B\), \(\Lambda\ge1\). For each \(i\), an axial forest \(F_i\) and
  branch \(g_i\) with \(g_i(0)=0\) are given. The branch and fine action
  are real-analytic at zero, and \(g_i\) is stationary for every
  block-kernel choice \(t\); the free-variable second variation at
  \((0,0)\) is at least \(\eta\sum_e|v_e|^2\) for a common \(\eta>0\)
  and every \(t\). Their derivatives are the axial free interpolant
  minus the copied-field derivative, and their tree-level second jets
  equal the coarse quadratic form. A local Haar chart ratio
  \(\mathsf J\) is given.
\end{itemize}

Then there exist \(Y\) attached to \(T\) with straight-route rows and
\(\alpha\in\mathbb R\) such that
\[Y\in\mathrm{Bounds}(\mathbf q;\eta',A',\Lambda')\quad(0<\eta'\le s,\ A'\ge A_*,\ \Lambda'\ge\max(\Lambda_0(L),\Lambda/2)),\qquad |\alpha+\kappa\log L|+1\le d.\]
For every \(i,t\), the tree action \(S_t(x,g_i(x))\) and
\(\mathrm{main}_{\mathbf q,Y}(e^x)\) have equal jets through order four.
Here
\[S_t(x,z)=\mathrm{main}_{\mathbf p,R}(U_{x,z})+a\sum_e\left(1-d_{\rm rep}^{-1}\mathrm{Re}\,\mathrm{tr}(V_e\Pi_t(U_{x,z})_e^{-1})\right),\]
where \(V=e^x\), \(U_{x,z}\) is the forest lift with free variables
\(z\), and \(\Pi_t\) is the block predictor. The one-loop function
\[\mathrm{th}_{i,t}=\mathrm{loop}_R(U_{x,g_i(x)})+\tfrac12\log\det\mathrm{Hess}_zS_t(x,g_i(x))-\log\prod_e\mathsf J(g_i(x)_e)\]
is analytic at zero, with positive Hessian determinant nearby. For every
\(\Theta\ge\Theta_0(L)\), \(c'\le\gamma\) and coarse real plane waves
\(v=\mathrm{Re}(e^{ip\cdot x}\xi)\),
\(w=\mathrm{Re}(e^{ip'\cdot x}\zeta)\) whose direction-colour amplitudes
have modulus at most one,
\[\Big|D^2\mathrm{th}_{i,t}(0)[v,w]-D^2\big(\alpha\,\mathrm{quad}_{\mathbf q,Y}+\mathrm{loop}_Y\big)(0)[v,w]\Big|\le|\Lambda_{m_i}|\Theta e^{-c'm_i}.\]
Both one-loop functions are real-analytic at zero; \(\mathrm{quad}\)
omits cubic/quartic atoms. The drift is independent of \(L\);
\(\Lambda_0=a_0L^4+b_0\log L+c_0\) and \(\Theta_0=b_1L^4\) with
nonnegative constants.

For iteration, fix positive stage widths/lower bounds, stage amplitudes,
\(K\) and clip bounds. Choose maximal path widths \(\sigma_j\) and
minimal amplitudes \(A_j\); for every
\(0<\tau_j\le\sigma_j,A'_j\ge A_j\) choose drifts \(d_j\ge0\), rates
\(\gamma_j>0\), minimal factors \(u_j\); for every compatible choice of
stage parameters with direct clip coefficients below their prescribed
bounds, drifts at least \(d_j\) and block-factor lower bounds at least
\(u_j\), and then every admissible choice of block factors, choose
volume thresholds and Hessian bounds \(\eta_j>0\); for every admissible
protection sequence choose coupling-entry thresholds.

For any subsequent \(\Gamma_j\ge1\), \(0<c_j\le\gamma_j\), compatible
history meeting III.2's thresholds, free-field horizon satisfying
III.10's strip, Floor, symmetry, top-identity, interpolation and
three-form coercivity conditions, second horizon with identical
beds/interpolation/curl kernels, and Haar ratio, there are kinetic data
\((\alpha_{N,m},P_{N,m},g_{N,m})\).

Every \(P_{N,m}\) is attached with level-\(m\) bounds and loop amplitude
\(\mathrm{loopCap}_m\); \(P_{N,N}\) is the Wilson top datum on every
volume with zero loops. For every \(m<N\), volume and axial forest, the
branches satisfy stationarity, the derivative condition,
second-variation bound \(\eta_{J_{m+1}}\), classical match,
\(|\alpha_{N,m}+\kappa\log L_m|+1\le\mathrm{drift}_{m+1}\), and the
finite bound with \((\Gamma_{J_{m+1}},c_{J_{m+1}})\) at loop amplitude
both \(\mathrm{loopCap}_{m+1}\) and \(1\).

\emph{Status:} step and iteration theorems, conditional on the stated
inputs. \emph{Lean:}
\href{https://github.com/yang-mills-group/OpenAI-YM-20261028/blob/3456e5e75cf8d4cd558955a5c7638fba2d68f464/YM/Next/STATIONARY/Level/Transition.lean\#L92-L182}{\texttt{transition\_cap}},
\href{https://github.com/yang-mills-group/OpenAI-YM-20261028/blob/3456e5e75cf8d4cd558955a5c7638fba2d68f464/YM/Next/STATIONARY/FullPlanned.lean\#L59-L66}{\texttt{geared}},
\href{https://github.com/yang-mills-group/OpenAI-YM-20261028/blob/3456e5e75cf8d4cd558955a5c7638fba2d68f464/YM/PCalc/Loop.lean\#L93-L115}{\texttt{Kin.Planned}}.

\textbf{Proposition III.7.5 (completed free inverse).} For stage \(j\),
kernel parameters \(K\), \(\sigma,\varepsilon>0\), \(\Omega\ge0\), there
exist \(\eta>0\) and, for every \(L\), \(A_L\ge0\): every admissible
stage-\(j\) patch and bed with \(H\in\mathrm{strip}(\sigma/L,\Omega)\)
and Floor \((\varepsilon,D)\) on every side have a completed block-Bloch
inverse \(\mathcal G_B\in\mathrm{strip}(\eta,A_L)\) satisfying
\(\mathcal C_B\mathcal G_B=1\) on every step, for any coarse side and
kernel choice. The joint free form is
\(\langle dv,Hdv\rangle+a\|{-}Qv+d_ch\|^2\), completed by
\(F_{1/L}^{\mathsf T}F_{1/L}\);
\(F_\mu=d_{\!1}^*P_f+\mathfrak m_\mu X\), with \(X\) minus the fine
rooted-frame logarithm derivative. Matrices use physical counting bases;
\(\eta\) is uniform in \(L\).

\emph{Status:} theorem. \emph{Lean:}
\href{https://github.com/yang-mills-group/OpenAI-YM-20261028/blob/3456e5e75cf8d4cd558955a5c7638fba2d68f464/YM/WilsonBulk/Spectral/Backing.lean\#L54-L66}{\texttt{WilsonBulk.uniform\_backing}}.

\textbf{Proposition III.7.6 (Faddeev--Popov identity).} For an
admissible step, clip \(\mathbf p\),
\(R\in\mathrm{Bounds}(\mathbf p;\sigma,A,\Lambda)\) with
\(\sigma>0,A,\Lambda\in\mathbb R\), an axial branch \(g(0)=0\) analytic
and stationary for every kernel choice with second variation bounded
below by \(\eta\sum_e|v_e|^2\), \(\eta>0\), every \(\mu>0\), kernel
choice \(t\) and local Haar chart ratio, for coarse fields \(x\) near
zero
\[\det\mathcal P_\mu>0,\qquad\mathrm{th}=\mathrm{loop}_R(U_{x,g(x)})+\tfrac12\log\det\mathcal P_\mu-\log\det\mathcal G_\mu.\]
Here \(\mathcal P_\mu\) acts on all fine-link and coarse-frame
directions: it is the left-coordinate Hessian of
\[\mathscr S_t(V;U,h)=\mathrm{main}_{\mathbf p,R}(U)+a\sum_e\left(1-d_{\rm rep}^{-1}\mathrm{Re}\,\mathrm{tr}\big(V_e[\mathrm{gauge}_h\Pi_t(U)]_e^{-1}\big)\right)\]
at \(V=e^x\) and the lifted joint stationary configuration, completed by
\(F_\mu^{\mathsf T}F_\mu\); \(F_\mu=d_U^*P_f+\mathfrak m_\mu X\), where
\(X\) is the rooted-frame derivative at that configuration, and
\(\mathcal G_\mu=d_U^*d_U+\mathfrak m_\mu\). The underlying identity is
\(T^2\det\mathcal P_\mu=\det M(\det\mathcal G_\mu)^2\), with axial
Hessian \(M\) and Haar Jacobian \(T\).

\emph{Status:} theorem. \emph{Lean:}
\href{https://github.com/yang-mills-group/OpenAI-YM-20261028/blob/3456e5e75cf8d4cd558955a5c7638fba2d68f464/YM/Next/STATIONARY/Joint/HaarJac.lean\#L177-L202}{\texttt{tribute}}.

\subsubsection{2. What it uses}\label{what-it-uses-9}

\begin{itemize}
\tightlist
\item
  I.1: gauge group, Wilson action and Casimir;
  \href{https://github.com/yang-mills-group/OpenAI-YM-20261028/blob/3456e5e75cf8d4cd558955a5c7638fba2d68f464/YM/Gauge/Colors.lean\#L180-L182}{\texttt{casimir}}.
\item
  III.1: path, block-step and kinetic specifications; III.2: compatible
  histories and constant thresholds,
  \href{https://github.com/yang-mills-group/OpenAI-YM-20261028/blob/3456e5e75cf8d4cd558955a5c7638fba2d68f464/YM/Price/Intent.lean\#L145-L162}{\texttt{Supply.Paid}}.
\item
  III.6: analytic stationary branches and free-interpolation derivative,
  \href{https://github.com/yang-mills-group/OpenAI-YM-20261028/blob/3456e5e75cf8d4cd558955a5c7638fba2d68f464/YM/Integrated.lean\#L229-L232}{\texttt{Integrated.stationary\_scene}}.
\item
  III.10: bed symmetry and free-field horizon conditions,
  \href{https://github.com/yang-mills-group/OpenAI-YM-20261028/blob/3456e5e75cf8d4cd558955a5c7638fba2d68f464/YM/Bid/Craft.lean\#L63-L73}{\texttt{Craft.Rack}};
  supplied globally by III.4.
\item
  II.1: Haar chart,
  \href{https://github.com/yang-mills-group/OpenAI-YM-20261028/blob/3456e5e75cf8d4cd558955a5c7638fba2d68f464/YM/Next/FOUNDATIONS/Primary.lean\#L84}{\texttt{chart\_ok}}.
  III.16/III.3: Laurent/Bloch strip calculus.
\item
  Standard results: polarization, lattice Poincaré/Koszul lemma, Schur's
  lemma, Faddeev--Popov procedure, background-field method, Balaban
  propagator bounds and rotation-invariant shell integration. No
  one-loop coefficient is assumed.
\end{itemize}

\subsubsection{3. Strategy}\label{strategy-9}

After lower jets are matched, gauge invariance makes residual Fourier
Taylor kernels transversal. Quantitative lattice Poincaré division
factors them through curvature; exponentially decaying Laurent
coefficients give cubic and quartic path atoms. This matches Wilson at
the top and the stationary tree action at each coarse level.

The Faddeev--Popov identity replaces axial determinants and Haar factors
by covariant vector and ghost determinants. Uniform inverse bounds
permit comparison with the infrared-regulated fine determinant. Its
continuum bubble determines the logarithmic slope; subtracting that
multiple of the coarse quadratic kernel leaves a fourth-order momentum
zero, represented by loop atoms.

\subsubsection{4. Steps}\label{steps-9}

\begin{enumerate}
\def\labelenumi{\arabic{enumi}.}
\tightlist
\item
  \textbf{Fourier kernels and Ward identities.} Local analytic
  translation sums give strip-analytic kernels; vanishing lower jets and
  gauge invariance give transversality in the infinite-volume limit.
  \href{https://github.com/yang-mills-group/OpenAI-YM-20261028/blob/3456e5e75cf8d4cd558955a5c7638fba2d68f464/YM/Next/STATIONARY/Shape/LocalOracle.lean\#L312-L348}{\texttt{gathered}},
  \href{https://github.com/yang-mills-group/OpenAI-YM-20261028/blob/3456e5e75cf8d4cd558955a5c7638fba2d68f464/YM/Next/STATIONARY/Lift/Ward.lean\#L647-L697}{\texttt{withdrew}}.
\item
  \textbf{Curvature factorization.} Hadamard division of \(e^{ip_j}-1\)
  solves each leg with controlled strip loss; invariant Laurent
  coefficients give atoms. Successively remove cubic and quartic
  residuals.
  \href{https://github.com/yang-mills-group/OpenAI-YM-20261028/blob/3456e5e75cf8d4cd558955a5c7638fba2d68f464/YM/Next/STATIONARY/Lift/ManyLegs.lean\#L533-L549}{\texttt{allleg}},
  \href{https://github.com/yang-mills-group/OpenAI-YM-20261028/blob/3456e5e75cf8d4cd558955a5c7638fba2d68f464/YM/Next/STATIONARY/Shape/ClassicalMatch.lean\#L561-L595}{\texttt{assembly}}.
\item
  \textbf{Determinants and inverses.} Axial determinant factorization
  and Haar density give III.7.6; real coercivity followed by
  complex-strip perturbation bounds the completed inverses uniformly in
  block-scaled norms.
  \href{https://github.com/yang-mills-group/OpenAI-YM-20261028/blob/3456e5e75cf8d4cd558955a5c7638fba2d68f464/YM/Next/STATIONARY/Joint/HaarJac.lean\#L177-L202}{\texttt{tribute}},
  \href{https://github.com/yang-mills-group/OpenAI-YM-20261028/blob/3456e5e75cf8d4cd558955a5c7638fba2d68f464/YM/Next/STATIONARY/Spectral/ComplexGap.lean\#L330-L353}{\texttt{post\_cert}}.
\item
  \textbf{Compute the slope.} Vector and ghost bubbles have angular
  average \(-(11/3)C_G|k|^{-4}\); normalized integration over
  \(1/L<|k|\le1\) yields \(-\kappa\log L\). The shell error
  \(O(|k|^{-3}+L^{-2}|k|^{-6})\) is integrable with an \(L\)-independent
  bound; the infrared ball and outer region also contribute bounded
  remainders.
  \href{https://github.com/yang-mills-group/OpenAI-YM-20261028/blob/3456e5e75cf8d4cd558955a5c7638fba2d68f464/YM/Next/STATIONARY/OneLoop/MomentumIntegral.lean\#L326-L338}{\texttt{poleSpin}},
  \href{https://github.com/yang-mills-group/OpenAI-YM-20261028/blob/3456e5e75cf8d4cd558955a5c7638fba2d68f464/YM/Next/STATIONARY/OneLoop/MomentumIntegral.lean\#L346-L408}{\texttt{absoluteGeode}}.
\item
  \textbf{Compare and subtract.} Contact terms do not change \(\alpha\);
  coarse/fine determinant comparison costs \(O(1)\), and aliasing
  preserves \(\alpha\). The residual two-leg kernel yields loop atoms
  with zero zeroth and first moments.
  \href{https://github.com/yang-mills-group/OpenAI-YM-20261028/blob/3456e5e75cf8d4cd558955a5c7638fba2d68f464/YM/Next/STATIONARY/Loops/HessianMatch.lean\#L67-L96}{\texttt{quoting}},
  \href{https://github.com/yang-mills-group/OpenAI-YM-20261028/blob/3456e5e75cf8d4cd558955a5c7638fba2d68f464/YM/Next/STATIONARY/Loops/HessianMatch.lean\#L719-L741}{\texttt{pole\_reduced}},
  \href{https://github.com/yang-mills-group/OpenAI-YM-20261028/blob/3456e5e75cf8d4cd558955a5c7638fba2d68f464/YM/Next/STATIONARY/Loops/Curvature.lean\#L394-L435}{\texttt{color\_pipe}}.
\item
  \textbf{Iterate.} Input-dependent classical terms carry \(L^{-1}\),
  inherited loops \(L^{-2}\); large \(L\) fixes output bounds, and
  coupling thresholds absorb newly generated \(O(L^4)\) loops. Downward
  induction gives the kinetic data.
  \href{https://github.com/yang-mills-group/OpenAI-YM-20261028/blob/3456e5e75cf8d4cd558955a5c7638fba2d68f464/YM/Next/STATIONARY/Level/Transition.lean\#L92-L182}{\texttt{transition\_cap}},
  \href{https://github.com/yang-mills-group/OpenAI-YM-20261028/blob/3456e5e75cf8d4cd558955a5c7638fba2d68f464/YM/Next/STATIONARY/FullPlanned.lean\#L59-L66}{\texttt{geared}}.
\end{enumerate}

\subsubsection{5. Where the difficulty
lies}\label{where-the-difficulty-lies-9}

The coefficient calculation must reach the blocked determinant with
error uniform in \(L\), through contact cancellation and every Fourier
alias. Curvature factorization must preserve decay and previously
imposed leg constraints. Output widths and classical amplitudes precede
input amplitudes. The auxiliary regulator respects gauge transformations
and unsigned permutations; full signed-permutation symmetry is asserted
for the blocked action.

\subsubsection{6. Notation}\label{notation-9}

\(\Lambda_n\) has \(n^4\) sites; \(U=e^x\).
\(E_W(U)=\sum_p(1-d_{\rm rep}^{-1}\mathrm{Re}\,\mathrm{tr}\,U_p)\) in
representation dimension \(d_{\rm rep}\); \(C_G\) is the adjoint
Casimir. \(\rho,H,D\) are bed roots, metric and diagonal; \(d_1\) is the
lattice curl, while \(d_{\!1}\) is \(d_U\) at \(U=1\). \(P_f\) projects
onto fine links, \(Q\) is the block average, \(d_c\) the coarse
gradient. Levels \(m=0,\ldots,N\) run from coarse to fine, and
\(J_{m+1}\in\{J_m,J_m+1\}\). \(\mathrm{loopCap}_m\) is the allowed loop
amplitude.

\subsubsection{7. What remains
unexplained}\label{what-remains-unexplained-9}

Subsidiary quotient, strip-perturbation, tadpole/bubble,
determinant/Haar, contact and alias estimates remain partly at statement
level. The reason for regulator power nine is unsettled. Global
hypotheses have not been checked in every downstream application; the
complete compactness and limiting-jet argument belongs to III.13.

Detailed notes: the full account of this unit, with proof sketches for
every step, is on the
\href{https://github.com/scottnarmstrong/yang-mills-informalization/blob/main/outline/notes/III-7.md}{detailed
notes page}.

\subsection{III.8. Connected activities and backward dependence
cones}\label{iii.8.-connected-activities-and-backward-dependence-cones}

\textbf{Part:} III, Ultraviolet construction (renormalization group).
\textbf{Lean:}
\href{https://github.com/yang-mills-group/OpenAI-YM-20261028/tree/3456e5e75cf8d4cd558955a5c7638fba2d68f464/YM/RG08/Steam}{\texttt{YM/RG08/Steam/}},
\href{https://github.com/yang-mills-group/OpenAI-YM-20261028/tree/3456e5e75cf8d4cd558955a5c7638fba2d68f464/YM/RG08/Forms}{\texttt{YM/RG08/Forms/}},
\href{https://github.com/yang-mills-group/OpenAI-YM-20261028/tree/3456e5e75cf8d4cd558955a5c7638fba2d68f464/YM/RG08/Skiff}{\texttt{YM/RG08/Skiff/}},
\href{https://github.com/yang-mills-group/OpenAI-YM-20261028/tree/3456e5e75cf8d4cd558955a5c7638fba2d68f464/YM/RG08/Jetty}{\texttt{YM/RG08/Jetty/}},
\href{https://github.com/yang-mills-group/OpenAI-YM-20261028/tree/3456e5e75cf8d4cd558955a5c7638fba2d68f464/YM/Push/GRID_APPROXIMATION}{\texttt{YM/Push/GRID\_APPROXIMATION/}}.
\textbf{Size:} 4,147 lines in the main-theorem closure.

\subsubsection{1. What this unit proves}\label{what-this-unit-proves-10}

\paragraph{Part A}\label{part-a}

Let \(I\) be any set, \(d\) a pseudometric, \(\sigma\ge0\). Rank-\(n\)
tensors are real functions on \(I^n\) with ordered slots;
\(\mathrm{diam}_d(x)\) is the largest pairwise distance, zero for an
empty tuple. Use
\[\|F\|_\sigma=\max_i\sup_z\sum_{x:x_i=z}|F(x)|e^{\sigma\mathrm{diam}_d(x)},\qquad [\Gamma]_\sigma=\max\left(\sup_u\sum_v|\Gamma(u,v)|e^{\sigma d(u,v)},\sup_v\sum_u|\Gamma(u,v)|e^{\sigma d(u,v)}\right).\]
Empty maxima and suprema are zero, including rank-zero tensor norms.
Propagators need not be symmetric. A pair is \((B,\delta)\)-controlled
if both norms are at most \(B\) and its difference norm is at most
\(B\delta\). A polynomial majorant satisfies \(0\le P(x)\le Cx^q\) for
\(x\ge1\), for some \(C\ge0\), \(q\in\mathbb N\).

A formal potential \(V=(V_{e,n})\) is eligible if \(V_{0,n}=0\) and
nonzero coefficients satisfy \(n\le e+2\), \(e+n\) even. Let
\(\Delta_\Gamma\) and \(\star_\Gamma\) contract and join slots,
averaging over output permutations:
\[\Delta_\Gamma H(x)=\mathrm{Sym}\sum_{u,v}\Gamma(u,v)H(u,v,x),\qquad(H\star_\Gamma G)(x,y)=\mathrm{Sym}\sum_{u,v}H(u,x)\Gamma(u,v)G(v,y).\]
Define \(V^{(0)}=V\) and
\[V^{(j+1)}_{e,n}=\frac1{j+1}\left[\binom{n+2}2\Delta_\Gamma V^{(j)}_{e,n+2}+\sum_{c=0}^n\sum_{i=0}^j\sum_{a=0}^e\frac{(c+1)(n-c+1)}2 V^{(i)}_{a,c+1}\star_\Gamma V^{(j-i)}_{e-a,n-c+1}\right].\]

\textbf{Proposition III.8.1.} Fix \(j,e,n\in\mathbb N\). Put
\[T_j(e,n)=\{(a,r):1\le a<e,\ 1\le r\le a+2\}\cup\{(e,r):1\le e,\ n<r\le e+2\},\]
adding \((e,n)\) if \(j=0\), \(e\ge1\), \(1\le n\le e+2\). There is a
polynomial majorant \(P\), depending only on \((j,e,n)\), such that for
all \(W\ge1\), \(\delta\ge0\), eligible \(V,V'\) and propagators
\(\Gamma,\Gamma'\), control \((W,\delta)\) of the propagators and
coefficient pairs in \(T_j(e,n)\) implies control \((P(W),\delta)\) of
the output pair \(V^{(j)}_{e,n}\). The majorant is uniform in
\(I,d,\sigma,\delta\). At \(n=0\) this is vacuous. \emph{Status:}
theorem. \emph{Lean:}
\href{https://github.com/yang-mills-group/OpenAI-YM-20261028/blob/3456e5e75cf8d4cd558955a5c7638fba2d68f464/YM/RG08/Steam/Routing.lean\#L34-L83}{\texttt{navigation}}.

\textbf{Proposition III.8.2.} For every eligible \(V\) and arbitrary
\(\Gamma\), all \(V^{(j)}\) are eligible, nonzero coefficients satisfy
\(n+2j\le3e\), and positive-\(j\) outputs are symmetric. For every real
\(t\),
\[\mathcal E_t(V)_{e,n}=\sum_{j=0}^{3e}t^jV^{(j)}_{e,n}=V_{e,n}+\mathcal R_t(V)_{e,n}.\]
Here \(\mathcal R_t\) sums \(j\ge1\), and \(\mathcal E_t(V)\) is
eligible. The series vanishes for \(n>3e\); otherwise its degree is at
most \(\lfloor(3e-n)/2\rfloor\). \emph{Status:} theorem. \emph{Lean:}
\href{https://github.com/yang-mills-group/OpenAI-YM-20261028/blob/3456e5e75cf8d4cd558955a5c7638fba2d68f464/YM/RG08/Steam/Valves.lean\#L50-L64}{\texttt{preserves}},
\href{https://github.com/yang-mills-group/OpenAI-YM-20261028/blob/3456e5e75cf8d4cd558955a5c7638fba2d68f464/YM/RG08/Steam/Valves.lean\#L66-L70}{\texttt{live}},
\href{https://github.com/yang-mills-group/OpenAI-YM-20261028/blob/3456e5e75cf8d4cd558955a5c7638fba2d68f464/YM/RG08/Steam/Gears.lean\#L65-L75}{\texttt{coal}}.

\textbf{Proposition III.8.3.} Let \(\mathcal S\) consist of smooth real
functions on \(\mathbb R^4\) with polynomially growing derivatives, with
\[\|f\|_{k,\sigma}=\max_{|\alpha|\le k}\sup_y|\partial^\alpha f(y)|e^{-\sigma|y|}.\]
For marked tensors valued in real-linear functionals on \(\mathcal S\),
use \(\|\lambda\|^*_{k,\sigma}=\sup_{\|f\|_{k,\sigma}\le1}|\lambda(f)|\)
and \(\|M\|_h=\sum_x\|M(x)\|^*_{k,\sigma}e^{\sigma h(x)}\), where
\(h:I\to[0,\infty)\) is 1-Lipschitz and \(h(x)=\max(0,\max_i h(x_i))\).
Set \(M^{(0)}=F\); the recursion replaces the contraction input by \(M\)
and the join by \(M^{(i)}\star_\Gamma V^{(j-i)}\), doubling its
coefficient.

For each \(j,e,n\in\mathbb N\) there is a polynomial majorant \(P\),
depending only on these indices, such that for all \(k\in\mathbb N\),
\(\sigma>0\), \(W\ge1\), \(U,\delta\ge0\): if \(V,V'\) are eligible,
propagators and unmarked inputs \(1\le a\le e\), \(1\le r\le a+2\) are
\((W,\delta)\)-controlled, and marked inputs with \(a<e\), or
\(a=e,r>n\), or \(j=0,(a,r)=(e,n)\) are \((U,\delta)\)-controlled, the
marked output pair is \((P(W)U,\delta)\)-controlled and real-linear.

For every \(l\in\mathbb N\), eligible \(V\) and \(F\) supported on
\(r\le l+a\) give nonzero outputs only if \(n+2j\le l+3e\) and
\(n\le l+e\). For fixed \(l,e,n\), there is a polynomial majorant \(P\)
depending only on these indices such that the paired sums
\(\sum_{j=1}^{l+3e}t^jM^{(j)}_{e,n}\) are \((P(W)U,\delta)\)-controlled
for \(|t|\le W\), under the preceding bounds omitting the marked input
\((e,n)\). They comprise the whole positive-time part if both marked
inputs satisfy that support condition. \emph{Status:} theorem.
\emph{Lean:}
\href{https://github.com/yang-mills-group/OpenAI-YM-20261028/blob/3456e5e75cf8d4cd558955a5c7638fba2d68f464/YM/RG08/Skiff/Spoke.lean\#L92-L127}{\texttt{journey}},
\href{https://github.com/yang-mills-group/OpenAI-YM-20261028/blob/3456e5e75cf8d4cd558955a5c7638fba2d68f464/YM/RG08/Skiff/Shed.lean\#L25-L36}{\texttt{registry}},
\href{https://github.com/yang-mills-group/OpenAI-YM-20261028/blob/3456e5e75cf8d4cd558955a5c7638fba2d68f464/YM/RG08/Skiff/Spoke.lean\#L61-L68}{\texttt{full}}.

\paragraph{Part B}\label{part-b}

Use I.1's gauge group and Haar probability measures. A block step has
odd \(L\ge3\), \(J\ge4\), patch radius \(\rho=(J+3)L\), embedding
\(\iota\), and \(\mathcal BF(U)=\int k_\beta(U,V)F(V)\,dV\). Locality on
\(X\) means dependence only on links with both endpoints in \(X\);
\(\mathcal I(X)\) adds gauge invariance, measurability and boundedness.

The patch hypothesis requires unique minimising frames and best
predictors whenever their branches are used, and nonempty fallback tags
for every cover configuration.

\textbf{Proposition III.8.4.} Assume the patch hypothesis,
\(\beta\in\mathbb R\). For bounded measurable complex \(F\) local on
\(X\), \(\mathcal BF\) is local on any
\(Y\supseteq\bigcup_{x\in X}\Lambda_{10\rho}(\iota x)\). If also
\(k_\beta(U^g,V^{g\circ\iota})=k_\beta(U,V)\) for every fine gauge
transformation \(g\), then \(F\in\mathcal I(\Lambda_R(o))\) implies
\(\mathcal BF\in\mathcal I(\Lambda_{LR+10\rho}(\iota o))\), with the
same sup bound. \emph{Status:} theorem. \emph{Lean:}
\href{https://github.com/yang-mills-group/OpenAI-YM-20261028/blob/3456e5e75cf8d4cd558955a5c7638fba2d68f464/YM/Push/GRID_APPROXIMATION/KernelLocal.lean\#L119-L136}{\texttt{back\_local}},
\href{https://github.com/yang-mills-group/OpenAI-YM-20261028/blob/3456e5e75cf8d4cd558955a5c7638fba2d68f464/YM/Push/GRID_APPROXIMATION/KernelLocal.lean\#L156-L169}{\texttt{back\_cylinder}}.

Histories satisfy the patch hypothesis, gauge covariance and covariance
under translations and signed axis permutations at every step. For each
\(N,p\), \(\mathcal B_m\) uses coupling \(\beta_{N,p,m+1}\); levels grow
finer, \(a_{m+1}=a_m/L_m\), \(J_{m+1}\le J_m+1\),
\(\rho_m=(J_{m+1}+3)L_m\). Define
\[\mathcal B_{k\to m}=\mathcal B_{m-1}\circ\cdots\circ\mathcal B_k,\quad r_{k\to k}(r)=r,\quad r_{k\to m+1}(r)=L_mr_{k\to m}(r)+10\rho_m,\]
with \(L_m\) odd, \(L_m\ge3\), \(J_m\ge4\), identity at \(m=k\),
\(\tau_m=30a_m(J_m+5)\), and \(\mathrm{room}(h,k,r)\) meaning
\(2(a_kr+\tau_k)<a_h\).

\textbf{Proposition III.8.5.} Under the patch hypothesis, for every
\(\beta\), bounded measurable \(F,G\) local on disjoint coarse sets
satisfy \(\mathcal B(FG)=\mathcal BF\,\mathcal BG\); likewise finite
products of \(F_i\in\mathcal I(C_i)\) on distinct tiling cells. For
every history, \(N,p\), \(h\le k\le m\), finite level-\(h\) set \(l\),
and \(F_i\in\mathcal I(\Lambda_r(\iota_{h\to k}i))\),
\(\mathrm{room}(h,k,r)\) implies
\[\mathcal B_{k\to m}\prod_{i\in l}F_i=\prod_{i\in l}\mathcal B_{k\to m}F_i.\]
This also holds for tile copies of one \(F\in\mathcal I(\Lambda_r(0))\),
with level-\(m\) copies on the right. \emph{Status:} theorem.
\emph{Lean:}
\href{https://github.com/yang-mills-group/OpenAI-YM-20261028/blob/3456e5e75cf8d4cd558955a5c7638fba2d68f464/YM/Push/GRID_APPROXIMATION/KernelLocal.lean\#L141-L150}{\texttt{back\_factors}},
\href{https://github.com/yang-mills-group/OpenAI-YM-20261028/blob/3456e5e75cf8d4cd558955a5c7638fba2d68f464/YM/Push/GRID_APPROXIMATION/Many.lean\#L68-L76}{\texttt{backTo\_copies}}.

\textbf{Proposition III.8.6.} For every history, \(N,p\), \(k\le m\),
pull-back preserves bounded measurability, is linear, fixes constants,
commutes with conjugation and decreases the sup norm. For bounded
measurable complex \(F\) at level \(k\) and integrable complex \(g\) at
level \(m\), the composite push-forward from level \(m\) to \(k\) is
integrable and
\[\int(\mathcal B^*_{k\to m}g)F=\int g\mathcal B_{k\to m}F,\qquad\int\mathcal B^*_{k\to m}g=\int g.\]
For \(h\le k\) and every level-\(h\) site \(i\), pull-back sends
\(\mathcal I(\Lambda_r(\iota_{h\to k}i))\) into
\(\mathcal I(\Lambda_{r_{k\to m}(r)}(\iota_{h\to m}i))\). On bounded
measurable functions it commutes with signed axis permutations followed
by translation by \(\iota_{h\to\cdot}i\), and centred level-\(h\) copy
maps. \emph{Status:} theorem. \emph{Lean:}
\href{https://github.com/yang-mills-group/OpenAI-YM-20261028/blob/3456e5e75cf8d4cd558955a5c7638fba2d68f464/YM/Push/GRID_APPROXIMATION/Iterate.lean\#L92-L101}{\texttt{pushTo\_dual}},
\href{https://github.com/yang-mills-group/OpenAI-YM-20261028/blob/3456e5e75cf8d4cd558955a5c7638fba2d68f464/YM/Push/GRID_APPROXIMATION/Iterate.lean\#L105-L115}{\texttt{backTo\_local}},
\href{https://github.com/yang-mills-group/OpenAI-YM-20261028/blob/3456e5e75cf8d4cd558955a5c7638fba2d68f464/YM/Push/GRID_APPROXIMATION/Iterate.lean\#L128-L147}{\texttt{backTo\_copy}}.

\textbf{Proposition III.8.7.} For every such history and nonnegative
integer levels and radii,
\[\tau_{m+1}\le\tfrac23\tau_m,\qquad\tau_m\le(2/3)^m\tau_0\to0,\qquad a_mr_{k\to m}(r)+\tau_m\le a_kr+\tau_k\quad(k\le m).\]
If \(h\le k\) and \(\mathrm{room}(h,k,r)\) holds, every cone box at
\(m\ge k\) lies in its centred level-\(h\) cell, of physical side
\(a_h\). For fixed \(h,r\), room holds for all sufficiently large
\(k\ge h\). \emph{Status:} theorem. \emph{Lean:}
\href{https://github.com/yang-mills-group/OpenAI-YM-20261028/blob/3456e5e75cf8d4cd558955a5c7638fba2d68f464/YM/Push/GRID_APPROXIMATION/Geometry.lean\#L147-L150}{\texttt{scales\_budget}},
\href{https://github.com/yang-mills-group/OpenAI-YM-20261028/blob/3456e5e75cf8d4cd558955a5c7638fba2d68f464/YM/Push/GRID_APPROXIMATION/Geometry.lean\#L160-L169}{\texttt{room\_fits}},
\href{https://github.com/yang-mills-group/OpenAI-YM-20261028/blob/3456e5e75cf8d4cd558955a5c7638fba2d68f464/YM/Push/GRID_APPROXIMATION/Geometry.lean\#L155-L158}{\texttt{plenty}}.

\subsubsection{2. What it uses}\label{what-it-uses-10}

\begin{itemize}
\tightlist
\item
  I.1: periodic lattices, Haar measure,
  \href{https://github.com/yang-mills-group/OpenAI-YM-20261028/blob/3456e5e75cf8d4cd558955a5c7638fba2d68f464/YM/Lattice/Scale.lean\#L28-L34}{\texttt{Schedule}}.
\item
  III.1: probability kernels, the patch hypothesis proved in
  \href{https://github.com/yang-mills-group/OpenAI-YM-20261028/blob/3456e5e75cf8d4cd558955a5c7638fba2d68f464/YM/Next/FOUNDATIONS/Primary.lean\#L70-L71}{\texttt{Next.Foundations.branches}},
  covariance,
  \href{https://github.com/yang-mills-group/OpenAI-YM-20261028/blob/3456e5e75cf8d4cd558955a5c7638fba2d68f464/YM/RG/History.lean\#L204-L206}{\texttt{History}}.
\item
  IV.1: cells, copies,
  \href{https://github.com/yang-mills-group/OpenAI-YM-20261028/blob/3456e5e75cf8d4cd558955a5c7638fba2d68f464/YM/Cells/Embedded.lean\#L25-L38}{\texttt{Tiling.planted}}
  (chessboard estimate).
\item
  Polchinski's formal heat equation; weighted \(\ell^1\) estimates;
  Fubini. III.14 supplies Part A's inputs.
\end{itemize}

\subsubsection{3. Strategy}\label{strategy-10}

For finite \(I\), eligible symmetric \(V\) and symmetric \(\Gamma\), the
unmarked series represents \(\log\mathbb E_{t\Gamma}e^V\) formally by
Wick's rule. Contractions and joins generate connected coefficients.
Diameter and height inequalities make exponential weights multiply;
induction and telescoping give polynomial and difference bounds.

Given the fine field, site tags are independent; coarse links are
independent conditional on those tags. Unused variables integrate out.
Locality iterates through the affine radius recurrence; decay of the
physical budget keeps cones within cells.

\subsubsection{4. Steps}\label{steps-10}

\begin{enumerate}
\def\labelenumi{\arabic{enumi}.}
\item
  For finite input norms and positive output rank,
  \(\|\Delta_\Gamma H\|_\sigma\le\|H\|_\sigma[\Gamma]_\sigma\) and
  \(\|H\star_\Gamma G\|_\sigma\le\|H\|_\sigma[\Gamma]_\sigma\|G\|_\sigma\).
  Diameter inequalities absorb internal sums; telescoping gives
  difference factors \(2,3\).
  \href{https://github.com/yang-mills-group/OpenAI-YM-20261028/blob/3456e5e75cf8d4cd558955a5c7638fba2d68f464/YM/RG08/Forms/Weld.lean\#L102}{\texttt{Held.bridge}}.
\item
  Heat-order induction gives bounds; slot counting gives support. The
  marked norm also controls scalars.
  \href{https://github.com/yang-mills-group/OpenAI-YM-20261028/blob/3456e5e75cf8d4cd558955a5c7638fba2d68f464/YM/RG08/Skiff/Spoke.lean\#L92-L127}{\texttt{journey}}.
\item
  Frame routes stay within \(9\rho\); uniqueness makes selected frames
  local; connectors fit within \(10\rho\).
  \href{https://github.com/yang-mills-group/OpenAI-YM-20261028/blob/3456e5e75cf8d4cd558955a5c7638fba2d68f464/YM/Push/GRID_APPROXIMATION/Cone.lean\#L115-L155}{\texttt{patch\_dep}}.
\item
  Disjoint coarse sets use disjoint tags and bonds, giving
  factorisation.
  \href{https://github.com/yang-mills-group/OpenAI-YM-20261028/blob/3456e5e75cf8d4cd558955a5c7638fba2d68f464/YM/Push/GRID_APPROXIMATION/KernelLocal.lean\#L141-L150}{\texttt{back\_factors}}.
\item
  Physical radius increments are at most \(\tau_m/3\); combine with
  budget decay and iterate.
  \href{https://github.com/yang-mills-group/OpenAI-YM-20261028/blob/3456e5e75cf8d4cd558955a5c7638fba2d68f464/YM/Push/GRID_APPROXIMATION/Geometry.lean\#L135-L145}{\texttt{scales\_next}}.
\item
  For \(h\le k\le N\), bounded measurable \(O\) at level \(N\) and
  \(F\in\mathcal I(\Lambda_w(0))\) at level \(k\) satisfying
  \(\mathrm{room}(h,k,w)\), expand replicated differences over subsets
  and apply factorisation and duality; the remaining factors stay inside
  one push-forward.
  \href{https://github.com/yang-mills-group/OpenAI-YM-20261028/blob/3456e5e75cf8d4cd558955a5c7638fba2d68f464/YM/Push/GRID_APPROXIMATION/Joint.lean\#L102-L128}{Joint.restored}.
\item
  For central-cell observables under IV.1's chessboard hypotheses,
  expand differences of products over nonempty subsets and bound each by
  its chessboard roots.
  \href{https://github.com/yang-mills-group/OpenAI-YM-20261028/blob/3456e5e75cf8d4cd558955a5c7638fba2d68f464/YM/Push/GRID_APPROXIMATION/Telescope.lean\#L34-L78}{tariff}.
\end{enumerate}

\subsubsection{5. Where the difficulty
lies}\label{where-the-difficulty-lies-10}

Positive heat order uses higher ranks at the same coupling order; the
polynomial may depend on the fixed orders. Unmarked scalar bounds are
vacuous. Conditional factorisation allows overlapping fine windows
without asserting independence under the fine Gibbs law.

\subsubsection{6. Notation}\label{notation-10}

\(e,n,j\): coupling order, rank, heat order.
\(\Delta_\Gamma,\star_\Gamma\):
\href{https://github.com/yang-mills-group/OpenAI-YM-20261028/blob/3456e5e75cf8d4cd558955a5c7638fba2d68f464/YM/RG08/Forms/Lap.lean\#L17}{\texttt{lap}},
\href{https://github.com/yang-mills-group/OpenAI-YM-20261028/blob/3456e5e75cf8d4cd558955a5c7638fba2d68f464/YM/RG08/Forms/Weld.lean\#L21}{\texttt{bridge}}.
\(\Lambda_r\): periodic \(\ell^\infty\) box
(\href{https://github.com/yang-mills-group/OpenAI-YM-20261028/blob/3456e5e75cf8d4cd558955a5c7638fba2d68f464/YM/Push/GRID_APPROXIMATION/Cone.lean\#L22-L23}{\texttt{window}}).
\(\iota_{h\to m}\): embedded site. \(\mathcal B^*\): iterated
push-forward
(\href{https://github.com/yang-mills-group/OpenAI-YM-20261028/blob/3456e5e75cf8d4cd558955a5c7638fba2d68f464/YM/Push/GRID_APPROXIMATION/Iterate.lean\#L27-L29}{\texttt{pushTo}}).
Copies translate, reflect and conjugate on odd tiles.

\subsubsection{7. What remains
unexplained}\label{what-remains-unexplained-10}

III.14 must supply covariance bounds; the relation to polymer activities
remains tentative. The marked normalised Gaussian interpretation remains
provisional (probable). Non-summable infinite-lattice sums are assigned
zero; unmarked scalars require separate convergence arguments; no
convergence in \(g\) is asserted. Support packages need separate
outlines.

Detailed notes: the full account of this unit, with proof sketches for
every step, is on the
\href{https://github.com/scottnarmstrong/yang-mills-informalization/blob/main/outline/notes/III-8.md}{detailed
notes page}.

\subsection{III.9. The marginal flow: the running coupling and its
reset}\label{iii.9.-the-marginal-flow-the-running-coupling-and-its-reset}

\textbf{Part:} III, Ultraviolet construction (renormalization group).
\textbf{Lean:} chapter 21:
\href{https://github.com/yang-mills-group/OpenAI-YM-20261028/tree/3456e5e75cf8d4cd558955a5c7638fba2d68f464/YM/UVNext/COUPLING_RESET}{coupling
reset},
\href{https://github.com/yang-mills-group/OpenAI-YM-20261028/tree/3456e5e75cf8d4cd558955a5c7638fba2d68f464/YM/AnalyticBridge/VERIFIED_RESET_STEP}{verified
reset},
\href{https://github.com/yang-mills-group/OpenAI-YM-20261028/tree/3456e5e75cf8d4cd558955a5c7638fba2d68f464/YM/Extra/WILSON_STARTUP}{Wilson
start-up}. \textbf{Size:} 7,215 lines.

\subsubsection{1. What this unit proves}\label{what-this-unit-proves-11}

Let \(G\subset U(n)\) be closed, with Lie algebra \(\mathfrak g\).

\textbf{Proposition III.9.1 (slope).} For simple nonabelian
\(\mathfrak g\) with orthonormal basis \((b_i)\) in
\(\langle X,Y\rangle=-\mathrm{Re}\mathrm{tr}(XY)/n\), define
\(C_G=-(\dim\mathfrak g)^{-1}\sum_i\mathrm{tr}_{\mathfrak g}((\mathrm{ad}b_i)^2)\)
and \(\kappa=11C_G/(24\pi^2)\). Then \(\kappa>0\). \emph{Status:}
definition and positivity theorem. \emph{Lean:}
\href{https://github.com/yang-mills-group/OpenAI-YM-20261028/blob/3456e5e75cf8d4cd558955a5c7638fba2d68f464/YM/Gauge/Casimir.lean\#L174-L176}{slopeConstant\_pos}.

\textbf{Proposition III.9.2 (reference scales).} Fix \(\kappa>0\),
\(\bar\beta>1\), an infinite sequence of odd \(L_m\ge3\). Set
\[s_m=\sum_{i<m}\log L_i=\log(1/a_m),\quad \hat\beta_m=\bar\beta+\kappa s_m,\quad M_m=e^{(\log\log\hat\beta_m)^2},\]
\[t_m=\hat\beta_m^{-1/2}M_m,\quad\varepsilon_m=\hat\beta_m^{-11/10},\quad\tilde\varepsilon_m=e^{-M_m^{1/5}},\quad B_m=\lceil(\log\hat\beta_m)^2\rceil,\quad\mathfrak B_m=[\hat\beta_m/2,2\hat\beta_m].\]
Then \(\hat\beta_{m+1}=\hat\beta_m+\kappa\log L_m\), \(\hat\beta\) is
nondecreasing and diverges. Additional scale hypotheses are:
\(t_m\le1\); sides at least \(100(B_m+1)\) and \(100(R_m+1)\) at levels
\(m,m+1\), for every volume; \(J_m\ge4\) nondecreasing, divergent,
increasing by at most one; nondecreasing derivative orders, spatial
powers, support exponents and stage factors \(\ell_j\), with
\(L_m=\ell_{J_{m+1}}\); antitone \(0<\theta_{J_m}\le1/J_m\);
\(\log L_m/s_m\to0\); \(\omega_{J_{m+1}}/\log L_m\le J_{m+1}^{-2}\) and
\(\log(1+\text{paths}+\text{width}^{-1}+\text{columns}+\text{relaxation}+\text{support exponent})_{J_{m+1}}/\log L_m\le J_{m+1}^{-2}\).
\emph{Status:} definitions and elementary theorems; scale hypotheses
supplied by III.2. \emph{Lean:}
\href{https://github.com/yang-mills-group/OpenAI-YM-20261028/blob/3456e5e75cf8d4cd558955a5c7638fba2d68f464/YM/RG/History.lean\#L97-L118}{Calendar.Fits}.

\textbf{Proposition III.9.3 (selected histories).} A selected history
assigns representations \(\rho_{N,p,m}\) for all \(N,p\) and \(m\le N\),
with volume-independent paths \(P_{N,m}\), level norms and clip, stage
path bounds and loop cap \(2M_m^{\theta_{J_m}}\), and
\(\beta_{N,p,m}\in\mathfrak B_m\). Regular activities are local, gauge
invariant, vanish at the flat field, with zero affine Hessian on
nonwrapping supports and bound \(\varepsilon_m\); forced activities obey
II.2 with bound \(\tilde\varepsilon_m\); the family respects
translations and signed permutations. At the top, \(v_N=0\),
\(\rho_N=e^{-\beta_NS_W}\); at the bottom, \(\beta_0=\bar\beta\). For
every \(m<N\),
\[\rho_m=T^{(m)}_{\beta_{m+1}}\rho_{m+1},\qquad |\beta_{m+1}-\beta_m-\kappa\log L_m|\le\omega_{J_{m+1}}.\]
For every \(\epsilon>0\), eventually
\(|\beta_{N,p,m}-\hat\beta_m|\le\epsilon s_m\) uniformly in
\(N\ge m,p\). Consequently \(\beta_{N,p,m}/s_m\to\kappa\) uniformly,
including \(m=N\), and
\(\int\rho_{N,p,m}\,d\mathrm{Haar}=Z_W(\Lambda_{N,p};\beta_{N,p,N})\)
for every \(m\le N\). \emph{Status:} specification; realized in III.13
using III.9.5 and III.10. \emph{Lean:}
\href{https://github.com/yang-mills-group/OpenAI-YM-20261028/blob/3456e5e75cf8d4cd558955a5c7638fba2d68f464/YM/RG/History.lean\#L234-L240}{History.Selected}.

\textbf{Proposition III.9.4 (reset along a curve).} Fix scales meeting
entry thresholds and kernel conditions, compatible horizon, Haar chart,
positive post-block constants, planned kinetic data and recipe. For
\(N,p,m<N\), let \(\rho'_x\), \(x\in I=[l,u]\), be continuous admissible
fine representations, with \(l\le u\); write
\(\beta'_x=\beta(\rho'_x)\). Fix stage-pair reset constants \(c>0\),
\(0<r\le1\), prescribed fine-stage \(g>0\), and coarse-stage sensor
\(\eta>0\). Assume:

\begin{itemize}
\tightlist
\item
  \(2(c+1)\varepsilon_m\le\eta t_m^2r\); a small-field probe satisfies
  \(\mathcal K(V)\ge|\Lambda_m|\eta t_m^2\).
\item
  Trial activities \(F(x,d)\), \(|d|\le r\), have fixed addresses, valid
  allowances, coefficient continuity, and common summable norm majorants
  along subsequences of every convergent parameter sequence. For
  \(H_x=T^{(m)}_{\beta'_x}\rho'_x\), \(H_x(1),H_x(V)>0\) and
  \[| (\beta'_x+\alpha)\mathcal K(V)+\mathcal L(V)-\log(H_x(1)/H_x(V))|\le|\Lambda_m|c\varepsilon_m,\]
  \[|\mathcal E_{F(x,b)}(V)-\mathcal E_{F(x,a)}(V)|\le|\Lambda_m|c\varepsilon_m|b-a|.\]
\item
  The trial shape
  \(e^{-(\beta'_x+\alpha+d)\mathcal K-\mathcal L-\mathcal E_{F(x,d)}}\mathcal G_{F(x,d)}\)
  differs from \(H_x(U)/H_x(1)\) by at most \(C_{x,U}|f(x,d)|\), with
  finite \(C_{x,U}\ge0\) independent of \(d\); forced activities have
  bound \(e^{-gM_{m+1}^2}\); every exact trial passes the recipe.
\end{itemize}

There is continuous \(d:I\to[-r,r]\) giving continuous admissible
\(\rho_x=H_x\), coupling \(\beta'_x+\alpha+d(x)\) and
\(v_x=|\Lambda_m|^{-1}\log H_x(1)\), retaining the stronger forced
bound. Without the last bullet, the unique residual root still matches
\(1,V\). \emph{Status:} theorem. \emph{Lean:}
\href{https://github.com/yang-mills-group/OpenAI-YM-20261028/blob/3456e5e75cf8d4cd558955a5c7638fba2d68f464/YM/UVNext/COUPLING_RESET/Advance.lean\#L175-L189}{advance}.

\textbf{Proposition III.9.5 (explicit reset; program).} Under III.9.4's
datum, curve and constants, assume simplicity,
\(\beta'_x\in\mathfrak B_{m+1}\), and the early presentation
\[H_x=e^{\tau_x|\Lambda_m|-\beta'_x\mathcal K-\mathcal L-\sum_i\chi_{X_i}E_{x,i}}\mathcal G_xe^{-\alpha\mathcal S}.\]
Assume local gauge-invariant smooth regular functions with nonnegative
sup-jet bounds \(h_{x,i}\) satisfying, at every anchor,
\(\sum_{i:o_i=o}e^{(a+3)\mathrm{cover}_{B_m}(X_i)}h_{x,i}\le H_x^E\),
forced-family conditions and site-pinned bound \(w_x\) at exponent
\(a+1\), translation and point-group symmetry, fixed addresses and
III.9.4's continuity/tails. Require priced reset, compatible clip, width
and cubic conditions, \(H_x^E\le c\varepsilon_m/4\), regular/forced
reorganization estimates fitting numerical bounds \(e_x,k_x\) that imply
both allowances and \(e^{-gM_{m+1}^2}\), and recipe acceptance and
continuation. Then
\[d(x)=\delta_x,\quad t_m^2|\delta_x|\le H_x^E,\quad v_x=\tau_x-|\Lambda_m|^{-1}\sum_iE_{x,i}(1).\]
All III.9.4 conclusions hold. A Gaussian-sector formulation additionally
requires integration licenses and equality of that integral with the
push. If these inputs hold at every level along every continuous
admissible curve with couplings in \(\mathfrak B_{m+1}\), and start-up
curves pass initial and retention tests, they give, for every \(N,p\),
an admissible continuous Wilson curve on \(\mathfrak B_N\) with coupling
parameter \(\beta\) and \(v=0\), and reset outputs along every
continuous admissible input curve with couplings in
\(\mathfrak B_{m+1}\) on every closed interval, for all \(m<N\). Output
couplings need not remain in the coarse band before shooting.
\emph{Status:} conditional theorems; physical inputs supplied by III.10.
\emph{Lean:}
\href{https://github.com/yang-mills-group/OpenAI-YM-20261028/blob/3456e5e75cf8d4cd558955a5c7638fba2d68f464/YM/AnalyticBridge/VERIFIED_RESET_STEP/Advance.lean\#L64-L90}{Return.reset},
\href{https://github.com/yang-mills-group/OpenAI-YM-20261028/blob/3456e5e75cf8d4cd558955a5c7638fba2d68f464/YM/UVNext/COUPLING_RESET/Stability.lean\#L67-L70}{program}.

\textbf{Proposition III.9.6 (Wilson start-up).} For compatible clip/norm
pairs (\(0<t\le1\), \(c_{\rm dir}\le1\), \(B\ge2\),
\(te^{(B-2)/c_{\rm rel}}\ge2\)), Wilson top paths, point-group-averaged
regular activities \(\beta(1-c_{\rm dir})R_p\) on small-field plaquettes
and forced activities \(e^{-\beta\,\mathrm{charge}(D,U)}\) on padded
large-field components give
\(\rho^{\rm start}_\beta(U)=e^{-\beta S_W(U)}\) for every real \(\beta\)
and configuration \(U\), with \(v=0\). For every \(G,k,q\), some
\(C>0\), \(0<\delta\le1\) give residual sup-jet bound \(Ct^5\) for every
norm choice of order \(k\) and power \(q\), every torus, plaquette and
site of that plaquette when \(t\le\delta\). At each fixed stage with
\(c_{\rm dir}\le1\), some \(\beta^*\) ensures compatibility, small
threshold and both norm allowances for every realized choice of levels
and every \(N\) at that stage with \(\hat\beta_N\ge\beta^*\). The band
curve is continuous; admissibility additionally requires Wilson top
paths, their stage bounds and recipe tests. Given old planned data and a
revised horizon with Wilson top quadratic datum, unchanged lower data
and initial guesses, bounded Wilson top paths can replace the old top
while preserving \(\alpha\) and planned conditions. \emph{Status:}
theorem. \emph{Lean:}
\href{https://github.com/yang-mills-group/OpenAI-YM-20261028/blob/3456e5e75cf8d4cd558955a5c7638fba2d68f464/YM/Extra/WILSON_STARTUP/Gas.lean\#L96-L103}{Start.evaluates}.

\textbf{Proposition III.9.7 (thresholds).} For a fixed pre-block profile
and simple \(G\), sensors \(\eta_j>0\) precede block sizes. Fix block
factors, \(0<\theta_j\le1\), reset constants, decay forecasts and
post-block costs. Assume, for each \(i\le j\le i+1\),
\(8f_{ji}\le q_i\ell_j^{\min}\) and \(u_{ji}(x)\to0\), where \(f_{ji}\)
is the inherited coefficient, \(\ell_j^{\min}\) the preassigned
block-size floor, and \(q_i=1/[64(1+A_i)]\) for the regular
reorganization coefficient \(A_i\ge0\). Every threshold record \(V\)
admits \(U\ge V\) such that every \(W\ge U\) ensures pricing, cubic,
width, reorganization and inherited inequalities beyond thresholds; clip
compatibility requires \(c_{\rm dir}\le1\). Realized levels must share
these data and satisfy \(W^{\rm entry}_{J_{m+1}}\le\hat\beta_m\).
Physical bounds also require \(4q_i\le c_{ji}\) and
\(4g_j\le\mathrm{rate}_{ji}\). Start-up thresholds are likewise monotone
when every used clip has \(c_{\rm dir}\le1\). \emph{Status:} conditional
numerical theorem. \emph{Lean:}
\href{https://github.com/yang-mills-group/OpenAI-YM-20261028/blob/3456e5e75cf8d4cd558955a5c7638fba2d68f464/YM/AnalyticBridge/VERIFIED_RESET_STEP/Inherited.lean\#L75-L89}{Time.beforeWait}.

\textbf{Proposition III.9.8 (marker sums).} Fix \(M,\mu\in\mathbb N\),
\(a\in\mathbb R\), \(c\ge0\). On a four-torus let arbitrarily indexed
polymers \(X_i\) have \(f_i\ge0\) and subsets of finitely many marker
names, at most \(\mu\) markers per site. Every positive-weight polymer
contains one of its markers' sites. If every marker-rooted sum
\(\sum_{i:v\in\mathrm{markers}(i)}e^{(a+1)\mathrm{cover}_M(X_i)}f_i\le c\),
then every site-pinned sum
\(\sum_{i:x\in X_i}e^{a\mathrm{cover}_M(X_i)}f_i\le\mu e(2(16M+9))^4c\),
uniformly in volume. All nonnegative sums must be summable.
\emph{Status:} theorem. \emph{Lean:}
\href{https://github.com/yang-mills-group/OpenAI-YM-20261028/blob/3456e5e75cf8d4cd558955a5c7638fba2d68f464/YM/Extra/MULTISOURCE_BOUNDS/Root.lean\#L161-L199}{root\_hit}.

\subsubsection{2. What it uses}\label{what-it-uses-11}

\begin{itemize}
\tightlist
\item
  II.2: polymer norms, masks and gas estimates; III.1: normalized block
  kernels and representation interfaces.
\item
  III.2: constant order, scales and shooting; III.7:
  \(|\alpha+\kappa\log L|+1\le\omega\).
\item
  III.4, III.10: Gaussian-sector identity, early presentation, local
  bounds, continuity and recipe acceptance; III.13: compatible kinetic
  data and selected history.
\item
  Intermediate value theorem, Taylor's formula, Schur's lemma, lattice
  reflections and permutations, Fourier analysis.
\end{itemize}

\subsubsection{3. Strategy}\label{strategy-11}

Extract the isotropic quadratic coefficient of the early regular
activities into \(\beta\). Subtract flat values and quadratic
projections, compensate path terms, and absorb failed-mask discrepancies
into the forced gas. A scalar equation identifies this exact
reorganization with the continuous reset; shooting chooses the bare
coupling with endpoint \(\bar\beta\).

\subsubsection{4. Steps}\label{steps-11}

Under the corresponding hypotheses in §1:

\begin{enumerate}
\def\labelenumi{\arabic{enumi}.}
\tightlist
\item
  Taylor expansion through order four gives the Wilson residual;
  connected large-field components give the exact gas.
  \href{https://github.com/yang-mills-group/OpenAI-YM-20261028/blob/3456e5e75cf8d4cd558955a5c7638fba2d68f464/YM/Extra/WILSON_STARTUP/Gas.lean\#L81-L94}{Start.gas\_raw}.
\item
  A small-field stripe gives \(\mathcal K(V)\ge|\Lambda|\eta t^2\).
  Pricing and the activity-motion bound make
  \(f(x,d)=(\beta'_x+\alpha+d)\mathcal K(V)+\mathcal L(V)+\mathcal E_{F(x,d)}(V)-\log(H_x(1)/H_x(V))\)
  increasing with slope at least \(\mathcal K(V)/2\), yielding a unique
  continuous root.
  \href{https://github.com/yang-mills-group/OpenAI-YM-20261028/blob/3456e5e75cf8d4cd558955a5c7638fba2d68f464/YM/UVNext/COUPLING_RESET/Advance.lean\#L121-L137}{paidResponse}.
\item
  Simplicity and lattice symmetry make the projected Hessian isotropic;
  exact masked splitting preserves every density value and gives
  \(d=\delta\).
  \href{https://github.com/yang-mills-group/OpenAI-YM-20261028/blob/3456e5e75cf8d4cd558955a5c7638fba2d68f464/YM/AnalyticBridge/VERIFIED_RESET_STEP/Evaluation.lean\#L70-L82}{Local.exponent}.
\item
  Fixed-stage asymptotics make the reset inequalities eventual; prefix
  maxima combine their thresholds monotonically.
  \href{https://github.com/yang-mills-group/OpenAI-YM-20261028/blob/3456e5e75cf8d4cd558955a5c7638fba2d68f464/YM/Push/PRICE_COMPILER/Basic.lean\#L111-L118}{collect\_upper}.
\item
  Shooting restricts intervals to keep couplings in their bands and
  reach \(\beta_0=\bar\beta\); summing drift errors gives uniform
  running.
  \href{https://github.com/yang-mills-group/OpenAI-YM-20261028/blob/3456e5e75cf8d4cd558955a5c7638fba2d68f464/YM/Shoot/Race.lean\#L23-L53}{Shoot.trajectory}.
\end{enumerate}

\subsubsection{5. Where the difficulty
lies}\label{where-the-difficulty-lies-11}

Two-point matching requires full masked reorganization to imply exact
transfer. Activity-norm continuity must survive hard configuration
masks. The running estimate gives an \(o(s_m)\) remainder, without a
bounded remainder or two-loop expansion.

\subsubsection{6. Notation}\label{notation-11}

\(N,p,m\): cutoff, volume, level; \(|\Lambda_m|\): site count;
\(\alpha=\alpha_{N,m}\): one-loop shift; \(a\): support exponent;
\(\mathrm{cover}_M\): polymer cover size; \(J_m\): stage; \(\omega_j\):
drift allowance; \(R_m\): patch radius.
\(\mathcal K,\mathcal L,\mathcal E,\mathcal G\): kinetic, loop, regular
and forced-gas terms; \(\mathcal S\): masked quadratic receipt;
\(\delta\): isotropic Hessian coefficient; \(\chi_X\): small-field
indicator. Admissible means prescribed paths/norms/clip, lawfulness and
recipe acceptance. Curve continuity means continuous coupling and vacuum
constant, fixed supports, inventories and eligibility tests, and
continuity in weighted regular sup-jet/forced uniform norms; regular
jets are measured on small-field supports at the prescribed order and
power.

\subsubsection{7. What remains
unexplained}\label{what-remains-unexplained-11}

Expand reorganization constants and connected-polymer estimates.
Nonlinear inputs belong to III.4/III.10; analytic existence and
uniqueness within the specified continuation class belong to III.13.

Detailed notes: the full account of this unit, with proof sketches for
every step, is on the
\href{https://github.com/scottnarmstrong/yang-mills-informalization/blob/main/outline/notes/III-9.md}{detailed
notes page}.

\subsection{III.10. One full nonlinear renormalization step: bounds on
irrelevant
terms}\label{iii.10.-one-full-nonlinear-renormalization-step-bounds-on-irrelevant-terms}

\textbf{Part:} III, Ultraviolet construction (renormalization group).
\textbf{Lean:}
\href{https://github.com/yang-mills-group/OpenAI-YM-20261028/tree/3456e5e75cf8d4cd558955a5c7638fba2d68f464/YM/Work}{\texttt{YM/Work/}}
(items E01--E09, H01--H04, I01--I04),
\href{https://github.com/yang-mills-group/OpenAI-YM-20261028/tree/3456e5e75cf8d4cd558955a5c7638fba2d68f464/YM/Bid}{\texttt{YM/Bid/}},
\href{https://github.com/yang-mills-group/OpenAI-YM-20261028/tree/3456e5e75cf8d4cd558955a5c7638fba2d68f464/YM/Next/FLUCTUATIONS}{\texttt{YM/Next/FLUCTUATIONS/}},
\href{https://github.com/yang-mills-group/OpenAI-YM-20261028/tree/3456e5e75cf8d4cd558955a5c7638fba2d68f464/YM/AnalyticNext/CONNECTED_ACTIVITY_JETS_E6}{\texttt{YM/AnalyticNext/CONNECTED\_ACTIVITY\_JETS\_E6/}},
\href{https://github.com/yang-mills-group/OpenAI-YM-20261028/tree/3456e5e75cf8d4cd558955a5c7638fba2d68f464/YM/PhysicalNext/RG_NEXT_COMPLETE_REPEATABLE_NONLINEAR_STEP}{\texttt{YM/PhysicalNext/RG\_NEXT\_COMPLETE\_REPEATABLE\_NONLINEAR\_STEP/}},
\href{https://github.com/yang-mills-group/OpenAI-YM-20261028/tree/3456e5e75cf8d4cd558955a5c7638fba2d68f464/YM/PhysicalNext/V9_JOINT_ENDPOINT_LOAD}{\texttt{YM/PhysicalNext/V9\_JOINT\_ENDPOINT\_LOAD/}},
\href{https://github.com/yang-mills-group/OpenAI-YM-20261028/tree/3456e5e75cf8d4cd558955a5c7638fba2d68f464/YM/Pressure}{\texttt{YM/Pressure/}}
(definitions). \textbf{Size:} 49,448 lines in the main-theorem closure.

\textbf{Conventions.} Levels are indexed from coarse to fine, as in
III.9.

\begin{itemize}
\item
  \emph{Step and couplings.} One step \(S\) goes from a fine torus of
  side \(f=cL\) to a coarse torus of side \(c\), with odd block factor
  \(L\ge3\). Write \(x\) for the coarse reference coupling
  \(\hat\beta_m\). The fine coupling is
  \(\beta_F=x+b_G\log L=\hat\beta_{m+1}\)
  (\href{https://github.com/yang-mills-group/OpenAI-YM-20261028/blob/3456e5e75cf8d4cd558955a5c7638fba2d68f464/YM/AnalyticNext/CONNECTED_ACTIVITY_JETS_E6/Intake.lean\#L67-L70}{\texttt{clock\_next}}),
  where \(b_G\ge0\) is the one-loop slope of III.9; \(b_G>0\) when the
  Lie algebra of \(G\) is simple.
\item
  \emph{Thresholds, radius and cutoff.}
  \[t_c(x)=x^{-1/2}e^{(\log\log x)^2}\ \text{(coarse field)},\qquad t=t_c(\beta_F)\ \text{(fine field)},\qquad T=\beta_F^{-49/100}\ \text{(large field proper)},\]
  with cluster radius \(B=\lceil(\log x)^2\rceil\) and path-length
  cutoff \(R=\lfloor B/10^5\rfloor\).
\item
  \emph{Sizes.} Since \(e^{(\log\log\beta)^2}=\beta^{o(1)}\), we have
  \(t=\beta_F^{-1/2+o(1)}\ll T\). On the other hand
  \(\beta_Ft^2=e^{2(\log\log\beta_F)^2}\) grows faster than every power
  of \(\log\beta_F\).
\item
  \emph{Field sizes.} \(\chi(g)\) is the chordal distance from
  \(g\in G\) to \(1\). For a plaquette \(p\),
  \(|U_p|=\chi(U_{\partial p})\). For a coarse bond \(b\),
  \(w_b=\chi(V_bQ_b^{-1})\) measures how far the coarse field \(V\) is
  from the block-spin predictor \(Q=Q(U,k)\) of the fine field \(U\)
  with gauge-fixing tags \(k\) (III.1).
\item
  ``For all large \(x\)'' means \(x\ge x_0\), with \(x_0\) independent
  of the period and volume. Other dependencies are specified in each
  statement.
\end{itemize}

\subsubsection{1. What this unit proves}\label{what-this-unit-proves-12}

One-step energy, activity, derivative and Gaussian bounds support the
effective-action construction of III.9; III.10.6 supplies start-up
fusion and marker estimates.

\textbf{Data of a step.}

\begin{itemize}
\tightlist
\item
  \emph{Request.} A request \(P\)
  (\href{https://github.com/yang-mills-group/OpenAI-YM-20261028/blob/3456e5e75cf8d4cd558955a5c7638fba2d68f464/YM/Pressure/Seeds.lean\#L89-L96}{\texttt{Teller}})
  consists of:

  \begin{itemize}
  \tightlist
  \item
    kernel choices \(K\), with penalty \(a>0\) and cuts;
  \item
    a fine stage and a coarse stage, each with clip constants
    \(\alpha,\lambda,\rho>0\) satisfying \(\lambda<\alpha/16\)
    (\href{https://github.com/yang-mills-group/OpenAI-YM-20261028/blob/3456e5e75cf8d4cd558955a5c7638fba2d68f464/YM/RG/PathActions.lean\#L139-L146}{\texttt{Clip}}),
    a path width and a path budget;
  \item
    a step index \(j\) and constants \(\eta>0\), \(C\).
  \end{itemize}
\item
  \emph{Tube.} A tube \(\mathscr T\)
  (\href{https://github.com/yang-mills-group/OpenAI-YM-20261028/blob/3456e5e75cf8d4cd558955a5c7638fba2d68f464/YM/Pressure/Rows.lean\#L94-L99}{\texttt{Tube}})
  carries the towers \(\Phi,\Phi',N\) and the fine and coarse path data.
  \(\Phi\) maps coarse plaquette coordinates to fine plaquette
  coordinates; \(\Phi'\) maps them to coarse-bond (``interior'')
  coordinates; \(N\) maps them to gauge coordinates.
\item
  \emph{Admissibility.} A step with tube \((S,\mathscr T)\) is
  admissible for \(P\)
  (\href{https://github.com/yang-mills-group/OpenAI-YM-20261028/blob/3456e5e75cf8d4cd558955a5c7638fba2d68f464/YM/Pressure/Seeds.lean\#L101-L109}{\texttt{Teller.Funds}})
  if:

  \begin{itemize}
  \tightlist
  \item
    \(S\) uses \(K\) and index \(j\);
  \item
    the patch is well posed
    (\href{https://github.com/yang-mills-group/OpenAI-YM-20261028/blob/3456e5e75cf8d4cd558955a5c7638fba2d68f464/YM/RG/Patch.lean\#L182-L185}{\texttt{Branches}}:
    unique minimizers whenever the frame or predictor selector uses
    minimization, and a nonempty set of eligible axial tags for every
    configuration);
  \item
    the scalar tower coefficients have exponentially weighted absolute
    sums at rate \(\eta\) bounded by \(C/L^2\) for the fine kernel and
    \(C\) for the interior and gauge kernels. Finite differences of
    order \(\ell\le j-2\) of the one-form and two-form cross kernels
    have bounds \(C/L^{1+\ell}\) and \(C/L^{2+\ell}\), respectively
    (\href{https://github.com/yang-mills-group/OpenAI-YM-20261028/blob/3456e5e75cf8d4cd558955a5c7638fba2d68f464/YM/PREC/Stacks.lean\#L111-L114}{\texttt{YM.Precision.Towers.Deep}});
  \item
    the clips and path budgets are those of \(P\);
  \item
    the flat, uncut tower matrices satisfy the exact identities
    \href{https://github.com/yang-mills-group/OpenAI-YM-20261028/blob/3456e5e75cf8d4cd558955a5c7638fba2d68f464/YM/PREC/Stacks.lean\#L88-L96}{\texttt{Towers.Roll}},
    among them the resolution of the identity
    \[\Phi_0^*\Phi_0+\Phi_0'^*\Phi_0'+N_0^*N_0=1.\] The transported
    kernels cut at length \(R\) need not satisfy this identity; their
    defects are estimated in Steps 5--6.
  \end{itemize}
\end{itemize}

\textbf{Seeds.} A store \(A=(A_F,A_I,A_\flat)\)
(\href{https://github.com/yang-mills-group/OpenAI-YM-20261028/blob/3456e5e75cf8d4cd558955a5c7638fba2d68f464/YM/Pressure/Seeds.lean\#L34-L37}{\texttt{Stores}})
consists of finite sets of fine plaquettes (``costs''), coarse bonds
(``aisles'') and coarse sites (``flags''). \(\sharp A\) is the sum of
their cardinalities, including coincident seeds of different types.
\(D_r(A)\) is the coarse \(r\)-neighbourhood of their positions: all
parent sites touched by a fine plaquette, both endpoints of a coarse
bond, and each flag site. \(A\) is \emph{consistent} with \((U,k,V)\) at
\((x,W)\)
(\href{https://github.com/yang-mills-group/OpenAI-YM-20261028/blob/3456e5e75cf8d4cd558955a5c7638fba2d68f464/YM/Pressure/Seeds.lean\#L51-L58}{\texttt{Plot}})
if the following hold:

\begin{itemize}
\tightlist
\item
  whenever a parent site touched by \(p\) meets \(D_{5B}(A)\),
  \(p\in A_F\) implies \(|U_p|\ge t/4\), and \(p\notin A_F\) implies
  \(|U_p|\le t/2\);
\item
  whenever an endpoint of \(b\) meets \(D_{5B}(A)\), \(b\in A_I\)
  implies \(w_b\ge Wt\), and \(b\notin A_I\) implies \(w_b\le2Wt\);
\item
  \(A_\flat\) agrees with the flags of \(V\) at threshold \(t_c(x)\) on
  \(D_{5B}(A)\);
\item
  each site tag based in \(D_{5B}(A)\) has positive frame weight.
\end{itemize}

Nothing is assumed outside \(D_{5B}(A)\). The \emph{spent} seeds
\(A^{\rm sp}\) are the costs with \(|U_p|>T\) and the aisles with
\(w_b>T\).

\textbf{Action gap.}

\begin{itemize}
\tightlist
\item
  \emph{Clipped density.} For a stage with clip
  \((\alpha,\lambda,\rho)\) and threshold \(s\), the clipped density of
  a plaquette is
  \(\phi_s(|U_p|)+\mathbf 1_{\rm kept}\tfrac12|\ell_p(U)|^2\)
  (\href{https://github.com/yang-mills-group/OpenAI-YM-20261028/blob/3456e5e75cf8d4cd558955a5c7638fba2d68f464/YM/Pressure/Rows.lean\#L68}{\texttt{crown}}).
  Here:

  \begin{itemize}
  \tightlist
  \item
    \(\phi_s(r)=\alpha r^2/2\) for \(r\le s\) and
    \(\phi_s(r)=\lambda sr\) for \(r>s\)
    (\href{https://github.com/yang-mills-group/OpenAI-YM-20261028/blob/3456e5e75cf8d4cd558955a5c7638fba2d68f464/YM/RG/PathActions.lean\#L218-L219}{\texttt{clipped}});
  \item
    \(\ell_p(U)\) is the path series of the stage, cut at length \(R\);
  \item
    ``kept'' means \(|U_q|\le s\,e^{{\rm dist}(p,q)/\rho}\) for all
    \(q\)
    (\href{https://github.com/yang-mills-group/OpenAI-YM-20261028/blob/3456e5e75cf8d4cd558955a5c7638fba2d68f464/YM/RG/PathActions.lean\#L210-L211}{\texttt{kept}}).
  \end{itemize}
\item
  \emph{Actions.} \(\mathcal A_F(U)\) is the fine clipped action (fine
  stage, threshold \(t\)) and \(\mathcal A_c(V)\) the coarse one (coarse
  stage, threshold \(t_c(x)\)).
\item
  \emph{Step functional.} \(\mathcal E_A(U,V)\)
  (\href{https://github.com/yang-mills-group/OpenAI-YM-20261028/blob/3456e5e75cf8d4cd558955a5c7638fba2d68f464/YM/Pressure/Seeds.lean\#L85}{\texttt{assessed}})
  is the sum of the row contributions of Step 1 over the rows located in
  \(D_{2B}(A)\). Summed over all rows of the torus, the same
  contributions give exactly
  \[\mathcal A_F(U)+\frac a2\sum_b w_b^2-\mathcal A_c(V).\] So
  \(\mathcal E_A\) is the action gap of the step (fine action, plus
  block-spin penalty, minus coarse action), localized near the seeds.
\end{itemize}

\textbf{Proposition III.10.1 (energy per seed).} Let \(P\) be a request.
There is \(L_0\) such that for every block factor \(L\ge L_0\) there is
\(W_0(L)\) such that for every \(W\ge W_0\) there is \(x_0\) with the
following property. Take \(x\ge x_0\) and:

\begin{itemize}
\tightlist
\item
  a step \(S\) of block factor \(L\) with coarse side \(c\ge100(1+B)\),
  and a tube \(\mathscr T\) admissible for \(P\);
\item
  a spanning forest \(F\) of the block structure, and a fine field \(U\)
  equal to \(1\) on the trunks of \(F\);
\item
  a store \(A\), tags \(k\) and a coarse field \(V\) such that \(A\) is
  consistent with \((U,k,V)\) at \((x,W)\).
\end{itemize}

Then
\[\kappa\,t^2\,\sharp A\;\le\;\mathcal E_A(U,V),\qquad \kappa=\tfrac12\min\{\lambda_F/8,\ a/32,\ \alpha_c/64\},\]
where \(\lambda_F\) is the linear clip constant of the fine stage and
\(\alpha_c\) the quadratic clip constant of the coarse stage.

\emph{Status:} theorem under the stated admissibility, tree-gauge and
consistency hypotheses.

\begin{itemize}
\tightlist
\item
  The Lean statement
  (\href{https://github.com/yang-mills-group/OpenAI-YM-20261028/blob/3456e5e75cf8d4cd558955a5c7638fba2d68f464/YM/Pressure/Seeds.lean\#L124-L125}{\texttt{Teller.cash}},
  \href{https://github.com/yang-mills-group/OpenAI-YM-20261028/blob/3456e5e75cf8d4cd558955a5c7638fba2d68f464/YM/Pressure/Seeds.lean\#L129-L131}{\texttt{Teller.Grants}})
  asserts only some \(\kappa>0\) chosen after \(L\) and \(W\). The value
  above is the one the proof produces (Step 7), and it is independent of
  \(L\) and \(W\).
\item
  This deterministic estimate precedes the Gaussian tail and activity
  bounds of III.10.2.
\end{itemize}

\emph{Lean:} proved by
\href{https://github.com/yang-mills-group/OpenAI-YM-20261028/blob/3456e5e75cf8d4cd558955a5c7638fba2d68f464/YM/Next/FLUCTUATIONS/Grants.lean\#L24}{\texttt{Fluctuations.grants}}.
This is
\href{https://github.com/yang-mills-group/OpenAI-YM-20261028/blob/3456e5e75cf8d4cd558955a5c7638fba2d68f464/YM/Integrated.lean\#L222-L225}{\texttt{Integrated.pressure\_grants}}
applied to
\href{https://github.com/yang-mills-group/OpenAI-YM-20261028/blob/3456e5e75cf8d4cd558955a5c7638fba2d68f464/YM/Next/FLUCTUATIONS/Bare.lean\#L64-L120}{\texttt{Fluctuations.bare}},
through
\href{https://github.com/yang-mills-group/OpenAI-YM-20261028/blob/3456e5e75cf8d4cd558955a5c7638fba2d68f464/YM/Bid/Energy.lean\#L155-L158}{\texttt{Bid.issues}}.

\textbf{Proposition III.10.2 (activities of one step).} Let \(P\) be a
request and \(\mathcal T\) target data
(\href{https://github.com/yang-mills-group/OpenAI-YM-20261028/blob/3456e5e75cf8d4cd558955a5c7638fba2d68f464/YM/AnalyticNext/CONNECTED_ACTIVITY_JETS_E6/OptionCost.lean\#L62-L76}{\texttt{Target}}):
constants \(\varepsilon,\delta>0\), \(e\ge0\), a rate
\(r\in(-\tfrac12,0)\), a halo \(h\ge500\) and a local exponential chart
of \(G\), with \(4\varepsilon\) below the support exponent of the fine
stage and \(8\varepsilon\) below both stage widths. Put
\(a_*=e+2\varepsilon+\delta\).

There is \(L_0\) such that for every \(L\ge L_0\) and every large \(W\)
there is \(\mu>0\) with the following property. For every patch \(D\) of
block factor \(L\) satisfying the branch condition there are window
constants \(o\)
(\href{https://github.com/yang-mills-group/OpenAI-YM-20261028/blob/3456e5e75cf8d4cd558955a5c7638fba2d68f464/YM/AnalyticNext/CONNECTED_ACTIVITY_JETS_E6/Issue.lean\#L33-L46}{\texttt{Target.Order}}:
a decay rate, a spectral gap, a shell width \(K\), a Lipschitz constant
\(\gamma\), a drift cap, a plateau \(a_o\) and a trigger level \(\xi\))
and a constant \(C\ge0\) such that the following holds for all large
\(x\), every step \(S\) on \(D\) with \(c\ge100(1+B)\) and every tube
admissible for \(P\)
(\href{https://github.com/yang-mills-group/OpenAI-YM-20261028/blob/3456e5e75cf8d4cd558955a5c7638fba2d68f464/YM/AnalyticNext/CONNECTED_ACTIVITY_JETS_E6/Room.lean\#L25-L28}{\texttt{Works}}).

\begin{itemize}
\item
  \begin{enumerate}
  \def\labelenumi{(\alph{enumi})}
  \tightlist
  \item
    The conclusions
    \href{https://github.com/yang-mills-group/OpenAI-YM-20261028/blob/3456e5e75cf8d4cd558955a5c7638fba2d68f464/YM/AnalyticNext/CONNECTED_ACTIVITY_JETS_E6/Issue.lean\#L51-L71}{\texttt{GrantBodies}}
    hold:
  \end{enumerate}

  \begin{enumerate}
  \def\labelenumi{\arabic{enumi}.}
  \tightlist
  \item
    the bound of III.10.1 holds at this \(x\) with constant \(\mu\)
    (\href{https://github.com/yang-mills-group/OpenAI-YM-20261028/blob/3456e5e75cf8d4cd558955a5c7638fba2d68f464/YM/Push/SECTOR_MEASURE/Oriented.lean\#L29-L32}{\texttt{PointCharge}});
  \item
    if \(U\) is in tree gauge and, within coarse distance \(B/4\) of the
    bond's parent site, site tags have positive frame weight, fine
    plaquettes have size at most \(t/2\), and coarse mismatches are at
    most \(2Wt\), that link lies within \(Kt\) of the replicated coarse
    link; the Lie logarithm is reliable at that radius
    (\href{https://github.com/yang-mills-group/OpenAI-YM-20261028/blob/3456e5e75cf8d4cd558955a5c7638fba2d68f464/YM/AnalyticNext/CONNECTED_ACTIVITY_JETS_E6/Coat.lean\#L25-L29}{\texttt{ShellAccess}});
  \item
    in every sector whose flags are the actual flags of \(V\)
    (\href{https://github.com/yang-mills-group/OpenAI-YM-20261028/blob/3456e5e75cf8d4cd558955a5c7638fba2d68f464/YM/AnalyticNext/CONNECTED_ACTIVITY_JETS_E6/Currents.lean\#L21}{\texttt{manifest}}),
    the Gaussian precision matrix is symmetric, has exponentially
    decaying entries and a positive spectral lower bound, and the
    covariance formed from its local inverses is positive definite
    (\href{https://github.com/yang-mills-group/OpenAI-YM-20261028/blob/3456e5e75cf8d4cd558955a5c7638fba2d68f464/YM/AnalyticNext/CONNECTED_ACTIVITY_JETS_E6/Shutter.lean\#L110-L117}{\texttt{Rotator}});
  \item
    the Gaussian centring drift has norm at most \({\rm cap}\cdot t\) at
    each fine bond. If the tube also satisfies the compatibility
    condition in (b), then, for every \(\beta>0\), a fluctuation
    satisfying \(|z_j|\le\xi t\sqrt\beta\) on all coordinates within
    distance \(B\) of a site farther than \(30B\) from the seeds has
    site and bond tag masses one, \(|U_p|<t/4\) for plaquettes based
    there, and \(w_b<Wt\) for bonds based there;
  \item
    numerical margins hold: \(4a_ot\) is below the chart radius, and,
    for the window's positive inverse-kernel constants \(\tau_0,C_0\)
    and window bound \(Q_0\ge1\), the quantity
    \(\epsilon_B=e^{-\tau_0B}(Q_0C_0+1)\) satisfies \(\epsilon_B<1\) and
    \(Q_0C_0\epsilon_B<1\).
  \end{enumerate}
\item
  (b) The conclusions
  \href{https://github.com/yang-mills-group/OpenAI-YM-20261028/blob/3456e5e75cf8d4cd558955a5c7638fba2d68f464/YM/AnalyticNext/CONNECTED_ACTIVITY_JETS_E6/Issue.lean\#L73-L88}{\texttt{Invoices}}
  hold for every tube satisfying the free-kernel compatibility condition
  \href{https://github.com/yang-mills-group/OpenAI-YM-20261028/blob/3456e5e75cf8d4cd558955a5c7638fba2d68f464/YM/AnalyticNext/CONNECTED_ACTIVITY_JETS_E6/Voyage.lean\#L25-L27}{\texttt{Stamp}},
  and for every input representation \(K\) of the fine-level density as
  in
  \href{https://github.com/yang-mills-group/OpenAI-YM-20261028/blob/3456e5e75cf8d4cd558955a5c7638fba2d68f464/YM/AnalyticNext/CONNECTED_ACTIVITY_JETS_E6/Intake.lean\#L42-L49}{\texttt{Intake}}.
  The input has the fine-stage support exponent, threshold \(t\),
  coupling \(\beta\in[\beta_F/2,2\beta_F]\), the path data and clips of
  \(\mathscr T\), cluster radius \(B(\beta_F)\), and activities obeying
  the incoming bounds of III.9: \(E\)-norm \(\beta_F^{-11/10}\),
  \(K\)-norm \(e^{-M^{1/5}}\) and loop budget \(2M\), where \(M\) is the
  fine monitor. The conclusions are:

  \begin{enumerate}
  \def\labelenumi{\arabic{enumi}.}
  \tightlist
  \item
    the cost record with gain \(\beta\mu t^2\) per seed and spend
    \(C\,{\rm bill}(x)^5\), where \({\rm bill}=1+B+(\log\beta_F)^2\),
    passes the tree-graph summability test
    \href{https://github.com/yang-mills-group/OpenAI-YM-20261028/blob/3456e5e75cf8d4cd558955a5c7638fba2d68f464/YM/Push/ABSOLUTE_SUMS/Mixed.lean\#L47-L58}{\texttt{Buy}};
  \item
    the letters of the new gas carry weights \(k_i,y_i,E\) satisfying a
    rooted Kotecký--Preiss condition with exponent \(\varepsilon\) and
    total \(\beta_F^{\,r}\)
    (\href{https://github.com/yang-mills-group/OpenAI-YM-20261028/blob/3456e5e75cf8d4cd558955a5c7638fba2d68f464/YM/AnalyticNext/CONNECTED_ACTIVITY_JETS_E6/Weights.lean\#L27-L29}{\texttt{Sales}}).
    For a letter with support \(X_i\) and extra length \(d_i\), this
    means nonnegative weights \(u_i\) with
    \[\sum_{i:o\in X_i}u_i\exp\big(a_*N_B(X_i)+2\varepsilon d_i\big)\le\beta_F^r\qquad\text{for every coarse site }o.\]
    Here \(N_B(H)\) is the minimum number of radius-\(B\) balls whose
    connected union contains \(H\). The letters are the forced
    activities of \(K\), the logarithmic letters, and the warning
    letters. Their weights are respectively \(k_i\), \(y_i\) and \(E\)
    times a fixed patch-dependent warning-amplitude bound; warning
    letters have \(d_i=0\);
  \item
    in every manifest sector, the following bounds hold:

    \begin{itemize}
    \tightlist
    \item
      the old activities are bounded by \(k_i\);
    \item
      the logarithmic letters \(\ell_i\) and the differences
      \(e^{\ell_i}-1\) are bounded by \(y_i\);
    \item
      every packet weight satisfies
      \(|{\rm Pack}|\le\exp\big((C\,{\rm bill}^5-\beta\mu t^2)\,\sharp A\big)\)
      (\href{https://github.com/yang-mills-group/OpenAI-YM-20261028/blob/3456e5e75cf8d4cd558955a5c7638fba2d68f464/YM/AnalyticNext/CONNECTED_ACTIVITY_JETS_E6/Revenue.lean\#L23-L24}{\texttt{Pack}});
    \item
      for every set \(v\) of warning addresses, the Gaussian measure of
      the catastrophe set is at most \(E^{\sharp v}\). The catastrophe
      set consists of the fluctuations \(z\) for which, near the anchor
      of each address in \(v\), some coordinate within distance \(2B\)
      satisfies \(|z_j|\ge\xi t\sqrt\beta\). Every fluctuation that
      triggers the warnings on \(v\) lies in that set.
    \end{itemize}
  \end{enumerate}
\end{itemize}

\emph{Status:} theorem. The target data, fine input bounds and
free-kernel compatibility are hypotheses. The input bounds are supplied
along the RG trajectory (III.9); the compatibility construction is
discussed in §7.

\begin{itemize}
\tightlist
\item
  \emph{Quantifier order:} \(L_0\), then \(W\), then \(\mu\) (uniform in
  \(D\)), then \((o,C)\) for each patch, then \(x\).
\item
  \emph{Size of the per-seed factor.} Since
  \(\beta\mu t^2\ge\frac\mu2e^{2(\log\log\beta_F)^2}\) and
  \({\rm bill}^5=O((\log\beta_F)^{10})\), the factor
  \(e^{C{\rm bill}^5-\beta\mu t^2}\) is eventually smaller than every
  power of \(\beta_F\).
\end{itemize}

\emph{Lean:}
\href{https://github.com/yang-mills-group/OpenAI-YM-20261028/blob/3456e5e75cf8d4cd558955a5c7638fba2d68f464/YM/AnalyticNext/CONNECTED_ACTIVITY_JETS_E6/Issue.lean\#L94-L185}{\texttt{Target.issue}}.
\href{https://github.com/yang-mills-group/OpenAI-YM-20261028/blob/3456e5e75cf8d4cd558955a5c7638fba2d68f464/YM/AnalyticNext/CONNECTED_ACTIVITY_JETS_E6/IssuePlan.lean\#L53-L69}{\texttt{Target.quotesTogether}}
repackages the output as the bundle
\href{https://github.com/yang-mills-group/OpenAI-YM-20261028/blob/3456e5e75cf8d4cd558955a5c7638fba2d68f464/YM/AnalyticNext/CONNECTED_ACTIVITY_JETS_E6/Batch.lean\#L82-L97}{\texttt{Cash}}
(the activity estimates together with witnesses for the letter weights).

\textbf{Proposition III.10.3 (derivatives of the new activities).} Let
\(P\) and \(\mathcal T\) be as in III.10.2. There is \(L_0\) with the
following property. Fix \(L\ge L_0\), a patch \(D\) satisfying the
branch condition, an order \(i\in\mathbb N\), a sector cutoff \(\chi\),
\(W,\mu>0\), window constants \(o\), and chart constants \((A,f)\)
satisfying
\href{https://github.com/yang-mills-group/OpenAI-YM-20261028/blob/3456e5e75cf8d4cd558955a5c7638fba2d68f464/YM/AnalyticNext/CONNECTED_ACTIVITY_JETS_E6/Eider.lean\#L32-L37}{\texttt{YM.ConnectedActivityE6.Blaze.Eider.Canary}}.
The latter require \(A,f>0\), \(\log(\exp u)=u\) for \(|u|\le2A\),
\(\exp(\log(\exp u\,g))=\exp u\,g\) for \(|u|\le A\) and
\(\|g-1\|\le f\), and analyticity of \(M\mapsto\log(\exp u\,M)\) for
\(|u|\le2A\), \(\|M-1\|\le f\). Fix a function \(\mathcal V\) of rate
\(-1/2\) that is eventually nonnegative (rate as defined in §6).

There are \(C\ge0\), functions \(q,\rho,p\) of rates \(-1/2,0,0\),
respectively, and \(l\in\mathbb N\) such that, for all large \(x\), the
following holds for every step on \(D\) satisfying the compatibility
condition, every fine representation \(R\), and every bound bundle \(H\)
of III.10.2 at \((W,x,\mu)\). One has \(8a_ot\le A/2\) and a
derivative-bound structure \(\Xi\)
(\href{https://github.com/yang-mills-group/OpenAI-YM-20261028/blob/3456e5e75cf8d4cd558955a5c7638fba2d68f464/YM/AnalyticNext/CONNECTED_ACTIVITY_JETS_E6/Das.lean\#L27-L46}{\texttt{Utopia}})
with the following properties in every sector at \(x\), of shell width
\(K\) from \(o\), cutoff \(\chi\), and with the actual coarse flags.

\begin{itemize}
\tightlist
\item
  Its nonpacket weights satisfy the rooted bound of III.10.2(b) with
  total \(\beta_F^r\). Its catastrophe probabilities, with trigger
  \(\xi/2\), are at most \((\mathrm{rad}^2)^{\sharp v}\).
\item
  Let \(V(y)\) be a family of coarse fields on a finite-dimensional real
  normed parameter space, based at the sector field \(V(y_0)=V_0\).
  Suppose each link matrix is \(C^i\) at \(y_0\) and
  \[\|D^m V_e(y_0)\|\le c_G\mathcal V(x)^m\qquad(0\le m\le i),\] where
  \(c_G\) is the fixed matrix-norm bound in
  \href{https://github.com/yang-mills-group/OpenAI-YM-20261028/blob/3456e5e75cf8d4cd558955a5c7638fba2d68f464/YM/AnalyticNext/CONNECTED_ACTIVITY_JETS_E6/Tint.lean\#L21-L23}{\texttt{Mode}}.
  Transport the fluctuation coordinates compatibly with this family: the
  physical and Gaussian coordinates represent the same field, their
  physical displacement has all derivatives through order \(i\) bounded
  by \(q(x)\), and, near \(y_0\), the resulting fine field agrees with
  the held fine field on every cold bond and on every hot bond whose
  physical coordinate has norm at most \(8a_ot\)
  (\href{https://github.com/yang-mills-group/OpenAI-YM-20261028/blob/3456e5e75cf8d4cd558955a5c7638fba2d68f464/YM/AnalyticNext/CONNECTED_ACTIVITY_JETS_E6/Eskim.lean\#L86-L95}{\texttt{YM.ConnectedActivityE6.Blaze.Silica.Clutch}}).
  In a region disjoint from the coarse flags, each letter \(a_j\) is
  \(C^i\) at \(y_0\) and satisfies
  \[\|D^m a_j(y_0)\|\le\mathrm{cost}_j\big(\rho(x)\mathrm{span}_j\big)^m\qquad(0\le m\le i)\]
  (\href{https://github.com/yang-mills-group/OpenAI-YM-20261028/blob/3456e5e75cf8d4cd558955a5c7638fba2d68f464/YM/AnalyticNext/CONNECTED_ACTIVITY_JETS_E6/Brush.lean\#L44-L47}{\texttt{Brush}}).
  Packet spans are one; other spans include their support and path
  lengths. The costs use the packet estimate and the nonpacket weights
  before extracting the catastrophe-probability factor.
\end{itemize}

The same \(\Xi\) satisfies an integrated estimate
(\href{https://github.com/yang-mills-group/OpenAI-YM-20261028/blob/3456e5e75cf8d4cd558955a5c7638fba2d68f464/YM/PhysicalNext/RG_NEXT_COMPLETE_REPEATABLE_NONLINEAR_STEP/Daisy/Charm.lean\#L48-L55}{\texttt{Herd}}).
Let \(\mathfrak p\) be a finite molecule with connected support-overlap
graph and total support disjoint from the flags. Require that its packet
stores are connected with pairwise disjoint radius-\(hB\)
neighbourhoods, that its old forced supports are pairwise disjoint, and
that every site in an old forced inventory touches a fine-plaquette seed
in the union of the packet stores. Along a family \(V(y)\) as above, the
prescribed transported integrand is \(C^i\) at \(y_0\) and has every
derivative of order at most \(i\) bounded by an integrable function
\(g\). With \(d\nu=d\lambda_F\,d\gamma\), the product of tree-slice Haar
measure and Gaussian fluctuation measure,
\[\int g\,d\nu\le |p(x)|\prod_{j\in\mathfrak p}\mathrm{fare}_j\,(1+\mathrm{shape}(\mathfrak p))^l.\]
The fares include a factor \(\mathrm{rad}\) for each distinct warning
address. For one measurable, gauge-invariant source letter of amplitude
at most \(w_M\), local with its eligibility tests included, the same
estimate replaces its fare by \(w_M\). A forced source must vanish
outside the original eligibility set; a regular source occupies an
unused regular-letter address. The remaining warning and shape factors
are unchanged.

\emph{Status:} theorem, conditional on the input bound bundle and
compatibility. The constants and functions are uniform in the volume.
\href{https://github.com/yang-mills-group/OpenAI-YM-20261028/blob/3456e5e75cf8d4cd558955a5c7638fba2d68f464/YM/AnalyticNext/CONNECTED_ACTIVITY_JETS_E6/Captains.lean\#L27-L36}{\texttt{Captain}}
and
\href{https://github.com/yang-mills-group/OpenAI-YM-20261028/blob/3456e5e75cf8d4cd558955a5c7638fba2d68f464/YM/PhysicalNext/RG_NEXT_COMPLETE_REPEATABLE_NONLINEAR_STEP/Daisy/Sparks.lean\#L36-L61}{\texttt{rocketsC}}
also control iterated left derivatives in spatially weighted directions.

\emph{Lean:}
\href{https://github.com/yang-mills-group/OpenAI-YM-20261028/blob/3456e5e75cf8d4cd558955a5c7638fba2d68f464/YM/AnalyticNext/CONNECTED_ACTIVITY_JETS_E6/Das.lean\#L49-L102}{\texttt{Utopia.whittle}}
(letter jets),
\href{https://github.com/yang-mills-group/OpenAI-YM-20261028/blob/3456e5e75cf8d4cd558955a5c7638fba2d68f464/YM/PhysicalNext/RG_NEXT_COMPLETE_REPEATABLE_NONLINEAR_STEP/Daisy/Charm.lean\#L58-L179}{\texttt{timing}}
(integrable majorants).

\textbf{Proposition III.10.4 (Gaussian sector of the step).} Fix a patch
\(D\), a request \(Q\) for which \(D\) satisfies the branch condition,
an aperture \(B_0>0\), a sector cutoff \(\chi\), and
\(k,p\in\mathbb N\). There are \(\sigma,d>0\) and \(\pi\ge1\)
(\href{https://github.com/yang-mills-group/OpenAI-YM-20261028/blob/3456e5e75cf8d4cd558955a5c7638fba2d68f464/YM/PhysicalNext/V9_JOINT_ENDPOINT_LOAD/Costs.lean\#L34-L42}{\texttt{Charges}})
such that the following holds for all large \(x\).

Take \(0<\beta\le2\beta_F\), a budget point of power \(p\) and threshold
\(t_c(x)\), a step on \(D\) with \(c\ge100(1+B)\) admissible for \(Q\),
a forest, tags, and a store \(A\) whose flags equal the flags of \(V\)
on the whole coarse torus. Freeze \(A\), the forest and tags while
differentiating. Let \(\Gamma\) be the sector precision matrix,
\(\widetilde\Gamma\) its local approximate inverse, \(m\) the smooth
sector mean at aperture \(B_0t\), and \(h=-g^{-1}\) times its
stationarity residual, with \(g=\beta^{-1/2}\). These are the objects in
\href{https://github.com/yang-mills-group/OpenAI-YM-20261028/blob/3456e5e75cf8d4cd558955a5c7638fba2d68f464/YM/PhysicalNext/V9_JOINT_ENDPOINT_LOAD/Costs.lean\#L60-L74}{\texttt{Value}}.
Then \(\Gamma\) and \(\widetilde\Gamma\) are positive definite, and
\(m\) equals the stationary logarithmic mean for every frozen-link
input.

For a coarse site \(o\) and \(0\le m\le k\), let \(\mathcal D^m\) be an
iterated left derivative in \(m\) directions satisfying
\(|A_e|\le t_c(x)(1+\mathrm{dist}(e,o))^p\). If \(x_i\) is the coarse
location of coordinate \(i\), put \[G_i=(1+\mathrm{dist}(x_i,o))^{pk}.\]
The bounds are
\[|\mathcal D^m\Gamma_{ij}|\le\pi e^{-\sigma\mathrm{dist}(x_i,x_j)}G_i^m,\qquad |\mathcal D^m(\Gamma\widetilde\Gamma-1)_{ij}|\le e^{-dB}\pi e^{-\sigma\mathrm{dist}(x_i,x_j)}G_i^m,\]
\[|\mathcal D^m m_i|\le\pi tG_i^m,\qquad |\mathcal D^m h_i|\le e^{-dB}\pi G_i^m.\]
Thus the spatial exponent at derivative order \(m\) is \(pkm\).

\emph{Status:} theorem.
\href{https://github.com/yang-mills-group/OpenAI-YM-20261028/blob/3456e5e75cf8d4cd558955a5c7638fba2d68f464/YM/PhysicalNext/V9_JOINT_ENDPOINT_LOAD/Supply.lean\#L72-L98}{\texttt{JointLoad.stage}}
also supplies a covariance cap and a trigger-tail estimate; its scales
are given in §7.

\emph{Lean:}
\href{https://github.com/yang-mills-group/OpenAI-YM-20261028/blob/3456e5e75cf8d4cd558955a5c7638fba2d68f464/YM/PhysicalNext/V9_JOINT_ENDPOINT_LOAD/Costs.lean\#L92-L157}{\texttt{JointLoad.capital}}.

\textbf{Proposition III.10.5 (schedule).}

\begin{itemize}
\tightlist
\item
  \begin{enumerate}
  \def\labelenumi{(\alph{enumi})}
  \tightlist
  \item
    Let \(P\) be a sheet of pre-\(L\) data (III.2) and \(K\) a price
    map. There are floors \(k_i\) with the following property. Take
    block factors \(L_i\ge k_i\), any protection, any delays \(V\) and
    any \(m\ge0\). Then there are delays \(W\), differing from \(V\)
    only by larger entry couplings \(\beta^*_i\), and numbers
    \(F_{ij}\ge m\), such that for every pair of stages
    \(j\le i\le j+1\) with \(j\) at least the first stage of \(P\) the
    bound of III.10.1 holds at every coarse coupling \(x\ge\beta^*_i\),
    for the request with fine stage \(i\), coarse stage \(j\) and step
    index \(i\), with block factor \(L_i\), aisle multiplier \(F_{ij}\)
    and some \(\mu_{ij}>0\) uniform in \(x\), period and volume
    (\href{https://github.com/yang-mills-group/OpenAI-YM-20261028/blob/3456e5e75cf8d4cd558955a5c7638fba2d68f464/YM/Pressure/Seeds.lean\#L161-L163}{\texttt{Ringing}}).
  \end{enumerate}
\item
  (b) If \(G\) is simple, then for every window datum of III.2 (entry
  and plateau requirements) there is a level schedule that meets all
  requirements of III.2
  (\href{https://github.com/yang-mills-group/OpenAI-YM-20261028/blob/3456e5e75cf8d4cd558955a5c7638fba2d68f464/YM/Price/Intent.lean\#L145-L162}{\texttt{Paid}})
  and whose periods are divisible: for every side \(c\) and every \(z\)
  there is a period index \(p\ge z\) with \(c\mid{\rm period}(p)\).
\item
  (c) Along a directed trajectory satisfying (a), fix an adjacent stage
  pair \(i=J(m)\), \(j=J(m+1)\) with \(i\) at least the first stage of
  \(P\). There is \(\mu>0\), uniform in the cutoff \(N\), level \(m\)
  and period, such that for \(m<N\) at this pair and \(x=\hat\beta_m\),
  one has \(\mu\,t_c(\hat\beta_{m+1})^2\,\sharp A\le\mathcal E_A\)
  whenever the fine field is in the chosen tree gauge and the store is
  consistent at multiplier \(F_{ji}\).
\end{itemize}

\emph{Status:} theorems. Part (a) is unconditional
(\href{https://github.com/yang-mills-group/OpenAI-YM-20261028/blob/3456e5e75cf8d4cd558955a5c7638fba2d68f464/YM/Next/FLUCTUATIONS/Grants.lean\#L32-L37}{\texttt{pressure\_pricing}}
applies
\href{https://github.com/yang-mills-group/OpenAI-YM-20261028/blob/3456e5e75cf8d4cd558955a5c7638fba2d68f464/YM/Bid/Energy.lean\#L173-L202}{\texttt{Bid.pricing\_order}}
to III.10.1). Divisibility is what lets torus tilings fit in
III.12--III.13.

\emph{Lean:}
\href{https://github.com/yang-mills-group/OpenAI-YM-20261028/blob/3456e5e75cf8d4cd558955a5c7638fba2d68f464/YM/PhysicalNext/RG_NEXT_COMPLETE_REPEATABLE_NONLINEAR_STEP/Pane/Divisors.lean\#L26-L135}{\texttt{Wait.engine}},
\href{https://github.com/yang-mills-group/OpenAI-YM-20261028/blob/3456e5e75cf8d4cd558955a5c7638fba2d68f464/YM/Pressure/Seeds.lean\#L166-L181}{\texttt{notation\_ticket}}.

\textbf{Proposition III.10.6 (start-up fusion H01--H04; markers
I01--I04).}

\begin{enumerate}
\def\labelenumi{(\Alph{enumi})}
\setcounter{enumi}{7}
\tightlist
\item
  \emph{Setting.} Fix a budget point and:
\end{enumerate}

\begin{itemize}
\tightlist
\item
  a \emph{pipe} \(X\): regular and forced addresses, and supports of
  logarithmic letters;
\item
  a \emph{tap} \(d\): regular activities, forced activities \(K_i\) and
  logarithmic letters \(\ell_i\).
\end{itemize}

The \emph{fused} gas
(\href{https://github.com/yang-mills-group/OpenAI-YM-20261028/blob/3456e5e75cf8d4cd558955a5c7638fba2d68f464/YM/Ovens/Coffee.lean\#L78}{\texttt{fused}})
groups the forced activities and the factors \(e^{\ell_i}-1\) into
connected clusters (pellets). Each pellet contains at least one forced
activity, and the forced supports in a pellet are pairwise disjoint.

\emph{Hypotheses.} Fix a bijective enumeration of the pellets. Let
\(\kappa,\gamma\ge0\).

\begin{itemize}
\tightlist
\item
  \(d\) is \emph{digested}
  (\href{https://github.com/yang-mills-group/OpenAI-YM-20261028/blob/3456e5e75cf8d4cd558955a5c7638fba2d68f464/YM/Ovens/Coffee.lean\#L81-L90}{\texttt{Digests}}):
  the gas is valid with specified regular and forced bounds; the
  logarithmic family has a summable uniform majorant, and each letter is
  measurable, gauge invariant, local with the flag tests included,
  vanishes on healthy configurations, and is translation and
  lattice-symmetry covariant.
\item
  \(d\) is \emph{bought} with constants \(\kappa,\gamma\)
  (\href{https://github.com/yang-mills-group/OpenAI-YM-20261028/blob/3456e5e75cf8d4cd558955a5c7638fba2d68f464/YM/Ovens/Coffee.lean\#L94-L98}{\texttt{Buys}}):
  \(|K_i|\le e^{-\kappa}f_i\) and \(|e^{\ell_i}-1|\le f_i\), where
  \(f\ge0\) satisfies the rooted Kotecký--Preiss condition at support
  exponent one larger than the budget exponent, with total \(\gamma\),
  and \((2B+1)^4\gamma\le1\).
\end{itemize}

\emph{Conclusions.}

\begin{itemize}
\tightlist
\item
  (H01) For every configuration \(V\), the fused gas equals the original
  gas times \(\exp(\sum_{i=0}^\infty\ell_i(V))\).
\item
  (H02) The fused activities satisfy the \(K\)-bound with any price
  \(v\ge e^{-\kappa}(e^\gamma-1)\).
\item
  (H03) The fused system is a forced family and admits the
  lattice-symmetry moves.
\item
  (H04) Fix the paths, clip, budget point, addresses and pellet
  enumeration. Suppose the original representations form a continuous
  curve on a subset of \(\mathbb R\), each forced or exponentiated
  logarithmic letter varies continuously in the supremum norm over
  configurations, and a single nonnegative rooted majorant as above
  works throughout the parameter set. Then the fused representations
  form a continuous curve in the same representation norm. This is a
  fixed-volume assertion.
\end{itemize}

\begin{enumerate}
\def\labelenumi{(\Roman{enumi})}
\tightlist
\item
\end{enumerate}

\begin{itemize}
\item
  (I01) Let \(B\ge8\), let \(F\) be a set of at least two markers with
  pairwise separation at least \(B^2\), and let \(F\subset H\). Then
  \(B\,\sharp F\le64\,N_B(H)\), where \(N_B(H)\) is the minimum number
  of radius-\(B\) balls with connected union covering \(H\). The markers
  lie on a finite periodic lattice with its max-distance; no nonwrapping
  hypothesis is required.
\item
  (I02) Fix a representation with budget order \(n\). Suppose marker
  bands satisfy
  \href{https://github.com/yang-mills-group/OpenAI-YM-20261028/blob/3456e5e75cf8d4cd558955a5c7638fba2d68f464/YM/Bid/Fuels.lean\#L66-L70}{\texttt{YM.Mix.Stubs.Holds}}
  at \(r>0\): valid local addresses, \(r\) times every unit amplitude at
  most one, \(B\ge8\) and separation at least \(B^2\) for each pair of
  distinct markers, factorially weighted derivative majorants, and
  separate rooted budgets at most one for single-marker and
  multiple-marker terms. For the multiple-marker budget one may instead
  use positive weights \(t_v\) and an extra exponent \(\delta\ge0\),
  with \(r\le t_v e^{\delta B/64}\) for every marker \(v\) and a rooted
  budget at most one using support exponent increased by \(\delta\).
  Then, at weight \(r/n!\), the bands satisfy
  \href{https://github.com/yang-mills-group/OpenAI-YM-20261028/blob/3456e5e75cf8d4cd558955a5c7638fba2d68f464/YM/Overlay/Table.lean\#L56-L65}{\texttt{YM.Overlay.Bands.Radial}}:
  positive weights, unit-amplitude bound one, and rooted ordinary
  derivative and forced-entry bounds at most two each.
\item
  (I03) Let a family satisfy the marker schedule condition:
  \(B_{r+1}^2\le2B_r^2\) for every \(r\), and
  \((L_rB_r^2)^4\mathcal F_{r+1}=o(\mathrm{scale}_{r+1})\). Here \(a_r\)
  is the lattice spacing, \(L_r=a_r/a_{r+1}\), and the vacuum bound is
  \[\mathcal F_r=2(2+4P_r)+1+6(\alpha_r+P_r^2)+4d_G\big(\log16+\tfrac12\log\hat\beta_r\big),\]
  with \(P_r\) the stage's path budget, \(\alpha_r\) its direct clip
  coefficient and \(d_G=\dim_{\mathbb R}M_m(\mathbb C)\) for the
  defining matrix representation. This is the hypothesis of
  \href{https://github.com/yang-mills-group/OpenAI-YM-20261028/blob/3456e5e75cf8d4cd558955a5c7638fba2d68f464/YM/Overlay/Prior.lean\#L15-L18}{\texttt{YM.Stock.Tabs}}.
  Fix \(k\ge1\) and levels \(m\le i\) such that \(kB_i^2\le a_m/a_i\)
  (an open stop). For every \(N\ge i\) and every period, the centred
  copies at level \(N\) are the embedded sites from level \(m\); their
  parents agree at consecutive levels from \(m\) to \(N\); and distinct
  embedded sites are separated by at least \(kB_r^2\) for every
  \(i\le r\le N\).
\item
  (I04) Assume \(G\) is simple. For every family, every \(p>0\) and
  \(r\in\mathbb N\), if \({\rm order}_i!\le{\rm columns}_i\) for all
  \(i\), then the widths
  \[w_i=\frac{\hat\beta_i^{-r/2}}{p\,{\rm columns}_i\prod_{m<i}{\rm columns}_{m+1}}\]
  are positive, and for every \(\varepsilon>0\) eventually
  \(5-\log(w_i/{\rm order}_i!)\le\varepsilon\,{\rm scale}_i\).
\end{itemize}

\emph{Status:} theorems. In the main-theorem closure:

\begin{itemize}
\tightlist
\item
  H01--H04 are used by the start-up representation at the finest level;
\item
  I01, I03 and I04 are used by the replicated-marker constructions;
\item
  the assemblies
  \href{https://github.com/yang-mills-group/OpenAI-YM-20261028/blob/3456e5e75cf8d4cd558955a5c7638fba2d68f464/YM/Bid/Kitchen.lean\#L56-L66}{\texttt{Oust.begins}}
  and
  \href{https://github.com/yang-mills-group/OpenAI-YM-20261028/blob/3456e5e75cf8d4cd558955a5c7638fba2d68f464/YM/Bid/Fuels.lean\#L146-L174}{\texttt{Bid.installs}}
  are not in the closure.
\end{itemize}

\emph{Lean:}

\begin{itemize}
\tightlist
\item
  Statements:
  \href{https://github.com/yang-mills-group/OpenAI-YM-20261028/blob/3456e5e75cf8d4cd558955a5c7638fba2d68f464/YM/Bid/Kitchen.lean\#L15-L16}{\texttt{H01}},
  \href{https://github.com/yang-mills-group/OpenAI-YM-20261028/blob/3456e5e75cf8d4cd558955a5c7638fba2d68f464/YM/Bid/Kitchen.lean\#L20-L21}{\texttt{H02}},
  \href{https://github.com/yang-mills-group/OpenAI-YM-20261028/blob/3456e5e75cf8d4cd558955a5c7638fba2d68f464/YM/Bid/Kitchen.lean\#L23-L24}{\texttt{H03}},
  \href{https://github.com/yang-mills-group/OpenAI-YM-20261028/blob/3456e5e75cf8d4cd558955a5c7638fba2d68f464/YM/Bid/Kitchen.lean\#L33-L35}{\texttt{H04}},
  \href{https://github.com/yang-mills-group/OpenAI-YM-20261028/blob/3456e5e75cf8d4cd558955a5c7638fba2d68f464/YM/Bid/Fuels.lean\#L124-L126}{\texttt{I01}}--\href{https://github.com/yang-mills-group/OpenAI-YM-20261028/blob/3456e5e75cf8d4cd558955a5c7638fba2d68f464/YM/Bid/Fuels.lean\#L140-L142}{\texttt{I04}}.
\item
  Proofs:
  \href{https://github.com/yang-mills-group/OpenAI-YM-20261028/blob/3456e5e75cf8d4cd558955a5c7638fba2d68f464/YM/Work/H01.lean\#L245-L251}{\texttt{H01.gathered}},
  \href{https://github.com/yang-mills-group/OpenAI-YM-20261028/blob/3456e5e75cf8d4cd558955a5c7638fba2d68f464/YM/Work/H02.lean\#L17-L80}{\texttt{H02.answer}},
  \href{https://github.com/yang-mills-group/OpenAI-YM-20261028/blob/3456e5e75cf8d4cd558955a5c7638fba2d68f464/YM/Work/H03.lean\#L138-L209}{\texttt{H03.fused\_forcedFamily}},
  \href{https://github.com/yang-mills-group/OpenAI-YM-20261028/blob/3456e5e75cf8d4cd558955a5c7638fba2d68f464/YM/Work/H04.lean\#L155-L161}{\texttt{H04.answer}},
  \href{https://github.com/yang-mills-group/OpenAI-YM-20261028/blob/3456e5e75cf8d4cd558955a5c7638fba2d68f464/YM/Work/I01.lean\#L19-L82}{\texttt{I01.answer}},
  \href{https://github.com/yang-mills-group/OpenAI-YM-20261028/blob/3456e5e75cf8d4cd558955a5c7638fba2d68f464/YM/Work/I02.lean\#L159-L180}{\texttt{I02.answer}},
  \href{https://github.com/yang-mills-group/OpenAI-YM-20261028/blob/3456e5e75cf8d4cd558955a5c7638fba2d68f464/YM/Work/I03.lean\#L43-L63}{\texttt{I03.answer}},
  \href{https://github.com/yang-mills-group/OpenAI-YM-20261028/blob/3456e5e75cf8d4cd558955a5c7638fba2d68f464/YM/Work/I04.lean\#L76-L129}{\texttt{I04.answer}}.
\end{itemize}

\subsubsection{2. What it uses}\label{what-it-uses-12}

\begin{itemize}
\tightlist
\item
  \textbf{III.1:} block-spin patches, tree gauge, predictor and tags;
  predictor-curvature reduction by \(L^{-2}\)
  (\href{https://github.com/yang-mills-group/OpenAI-YM-20261028/blob/3456e5e75cf8d4cd558955a5c7638fba2d68f464/YM/Pressure/Cost.lean\#L210-L252}{\texttt{qdc}}).
  The branch condition holds for the construction's kernel choices
  (\href{https://github.com/yang-mills-group/OpenAI-YM-20261028/blob/3456e5e75cf8d4cd558955a5c7638fba2d68f464/YM/Next/FOUNDATIONS/Primary.lean\#L70-L71}{\texttt{Foundations.branches}}).
\item
  \textbf{III.2:} sheets, block factors, protection, delays, entry and
  plateau requirements.
\item
  \textbf{III.3--III.4:} flat tower identities and exponential decay
  (III.3 attribution probable), Gaussian integration and rectangle
  tails.
\item
  \textbf{III.5--III.8:} polymer representations and Kotecký--Preiss
  bounds, stationary means and jets, quartic matching,
  connected-activity locality.
\item
  \textbf{III.9:} \(b_G\), thresholds, monitor and the incoming
  allowances \(\beta_F^{-11/10}\), \(e^{-M^{1/5}}\), \(2M\).
\item
  \textbf{I.1, II.2, III.16:} lattice geometry, Haar measure, chordal
  distance, activities and Taylor jets.
\end{itemize}

Standard results: Balaban's lattice-gauge renormalization method, rooted
Kotecký--Preiss and tree-graph bounds, Schur's test, Gaussian
exponential moments and Taylor estimates on compact charts. Request and
target data precede \(L\); window constants follow \(L,W\) and the
patch. Simplicity is required in III.10.5(b) and I04.

\subsubsection{3. Strategy}\label{strategy-12}

Seeds cost \(\kappa t^2\) each. At \(\beta\ge\beta_F/2\), their
Boltzmann suppression pays every fixed polylogarithmic entropy loss.
Completed squares isolate positive row contributions and two kernel
defects: affected cells are paid by spent seeds, while local gauge
fixing and flat identities reduce unaffected defects to chart errors and
exponentially small tails.

Gaussian integration produces packet, old forced, logarithmic and
warning letters. Energy and annulus estimates bound packets;
simultaneous Gaussian excursions bound warnings. Rooted bounds sum
connected molecules. Compatible fluctuation transport controls
derivatives, and Cauchy--Schwarz preserves one square root of the
simultaneous-warning probability in integrated bounds.

\subsubsection{4. Steps}\label{steps-12}

Steps 1--7 use III.10.1's admissibility, tree gauge and consistency
hypotheses; all thresholds are independent of period. Steps 8--10 use
III.10.2--III.10.3's activity, compatibility and derivative hypotheses;
Steps 11--12 use III.10.6 and III.10.5, respectively.

\begin{enumerate}
\def\labelenumi{\arabic{enumi}.}
\item
  \textbf{Complete the squares.} Fine, interior, coarse and gauge rows,
  together with the linear and completeness defects, sum exactly to
  \(\mathcal A_F+\frac a2\sum_b w_b^2-\mathcal A_c\) by adjoint
  cancellation; \(\mathcal E_A\) retains rows in \(D_{2B}(A)\).
  \href{https://github.com/yang-mills-group/OpenAI-YM-20261028/blob/3456e5e75cf8d4cd558955a5c7638fba2d68f464/YM/Pressure/Rows.lean\#L145-L163}{\texttt{Tube.locates}}.
\item
  \textbf{Pack seeds and compare links.} Fine plaquettes have bounded
  macro-box multiplicity; coarse plaquettes above a fixed \(\delta>0\)
  lie within \(4B/5\) of spent seeds, and the relevant neighbourhood
  contains \(O(B^4\sharp A)\) coarse plaquettes. Closing fine links by
  tree paths bounds their distance from replicated coarse links by
  macro-box curvature plus crossing mismatches.
  \href{https://github.com/yang-mills-group/OpenAI-YM-20261028/blob/3456e5e75cf8d4cd558955a5c7638fba2d68f464/YM/Work/E01.lean\#L289-L324}{\texttt{Work.E01.answer}};
  \href{https://github.com/yang-mills-group/OpenAI-YM-20261028/blob/3456e5e75cf8d4cd558955a5c7638fba2d68f464/YM/Work/E02.lean\#L132-L244}{\texttt{Work.E02.answer}}.
\item
  \textbf{Bound kernels and predictions.} Absolute row and column sums
  control transported kernels and path series; in particular
  \(|v|\le At\), \(|\Phi v|\le At/L^2\), \(|\Phi'v|\le At\), with \(A\)
  independent of \(L,W\). Norm-one adjoint transport and exponentially
  weighted coefficient sums give these estimates.
  \href{https://github.com/yang-mills-group/OpenAI-YM-20261028/blob/3456e5e75cf8d4cd558955a5c7638fba2d68f464/YM/Work/E03.lean\#L228-L256}{\texttt{Work.E03.answer}};
  \href{https://github.com/yang-mills-group/OpenAI-YM-20261028/blob/3456e5e75cf8d4cd558955a5c7638fba2d68f464/YM/Work/E04.lean\#L282-L307}{\texttt{Work.E04.answer}}.
\item
  \textbf{Obtain positive row charges.} Large \(L\) separates fine seeds
  from their predictions, and large \(W\) separates mismatch seeds;
  dropped path terms are paid after summation by column bounds. Coarse
  rows above \(t_c\) and below a fixed \(\delta\) are positive, while
  small coarse rows lose \(O(t_c^4)\); affected-cell losses are charged
  to large fine inputs, mismatches and \(O(t^2B^4\sharp A^{\rm sp})\).
  \href{https://github.com/yang-mills-group/OpenAI-YM-20261028/blob/3456e5e75cf8d4cd558955a5c7638fba2d68f464/YM/Work/E05.lean\#L211-L228}{\texttt{Work.E05.answer}};
  \href{https://github.com/yang-mills-group/OpenAI-YM-20261028/blob/3456e5e75cf8d4cd558955a5c7638fba2d68f464/YM/Work/E06.lean\#L413-L440}{\texttt{Work.E06.answer}};
  \href{https://github.com/yang-mills-group/OpenAI-YM-20261028/blob/3456e5e75cf8d4cd558955a5c7638fba2d68f464/YM/Work/E07.lean\#L115-L158}{\texttt{Work.E07.answer}}.
\item
  \textbf{Cancel flat defects.} Outside the affected set, local axial
  gauge fixing makes fine and coarse links \(h=CB^{10}T\)-close to
  identity; the flat uncut tower identities cancel both linearized
  defects. Comparing transported, truncated quantities with their flat
  versions leaves errors governed by \(\omega=(1+R)h+e^{-\tau R}\).
  \href{https://github.com/yang-mills-group/OpenAI-YM-20261028/blob/3456e5e75cf8d4cd558955a5c7638fba2d68f464/YM/Work/E08.lean\#L342-L344}{\texttt{Work.E08.answer}};
  \href{https://github.com/yang-mills-group/OpenAI-YM-20261028/blob/3456e5e75cf8d4cd558955a5c7638fba2d68f464/YM/Next/FLUCTUATIONS/Linear.lean\#L89-L136}{\texttt{roll\_l}};
  \href{https://github.com/yang-mills-group/OpenAI-YM-20261028/blob/3456e5e75cf8d4cd558955a5c7638fba2d68f464/YM/Next/FLUCTUATIONS/Linear.lean\#L69-L87}{\texttt{roll\_d}};
  \href{https://github.com/yang-mills-group/OpenAI-YM-20261028/blob/3456e5e75cf8d4cd558955a5c7638fba2d68f464/YM/Next/FLUCTUATIONS/Global.lean\#L76-L148}{\texttt{global\_chart\_cost}}.
\item
  \textbf{Sum unaffected defects.} Gauge covariance, locality and
  extension to global log-fields apply the chart estimates locally; the
  total defect loss is at most \(\varepsilon t^2\sharp A\) for every
  \(\varepsilon>0\) and all large \(x\). The factor \(B^4\) is absorbed
  because \(B^n\omega\to0\) and \(B^nh\omega/t\to0\) for every \(n\).
  \href{https://github.com/yang-mills-group/OpenAI-YM-20261028/blob/3456e5e75cf8d4cd558955a5c7638fba2d68f464/YM/Next/FLUCTUATIONS/Interior.lean\#L22-L63}{\texttt{interior\_cost}};
  \href{https://github.com/yang-mills-group/OpenAI-YM-20261028/blob/3456e5e75cf8d4cd558955a5c7638fba2d68f464/YM/Next/FLUCTUATIONS/Bare.lean\#L64-L120}{\texttt{Fluctuations.bare}}.
\item
  \textbf{Assemble the energy bound.} Choose \(L\), then \(W\), then
  \(x\): the first absorbs fine cross terms, the second mismatch losses,
  and the last all remaining losses, using spent-seed charges of size
  \(tT\). The row sum gives exactly III.10.1's \(\kappa\).
  \href{https://github.com/yang-mills-group/OpenAI-YM-20261028/blob/3456e5e75cf8d4cd558955a5c7638fba2d68f464/YM/Work/E09.lean\#L113-L159}{\texttt{Work.E09.answer}}.
\item
  \textbf{Bound packets and other letters.} The annulus remainder has
  rate \(-3/2\) and is absorbed into \(C\,\mathrm{bill}^5\); retained
  energy gives the packet estimate. Incoming bounds and rate calculus
  control old and logarithmic letters, including both \(\ell_i\) and
  \(e^{\ell_i}-1\).
  \href{https://github.com/yang-mills-group/OpenAI-YM-20261028/blob/3456e5e75cf8d4cd558955a5c7638fba2d68f464/YM/AnalyticNext/CONNECTED_ACTIVITY_JETS_E6/AssignedPrice.lean\#L54-L86}{\texttt{Coat.earns}};
  \href{https://github.com/yang-mills-group/OpenAI-YM-20261028/blob/3456e5e75cf8d4cd558955a5c7638fba2d68f464/YM/AnalyticNext/CONNECTED_ACTIVITY_JETS_E6/Revenue.lean\#L99-L140}{\texttt{steam\_prices}};
  \href{https://github.com/yang-mills-group/OpenAI-YM-20261028/blob/3456e5e75cf8d4cd558955a5c7638fba2d68f464/YM/AnalyticNext/CONNECTED_ACTIVITY_JETS_E6/Quotation.lean\#L47-L92}{\texttt{recommends}}.
  \emph{(annulus statement only; other proofs read)}
\item
  \textbf{Control Gaussian warnings.} Coercivity and exponentially
  decaying rows give a positive precision and covariance, an \(O(t)\)
  drift and a small residual; compatibility gives the no-new-seed
  margins. Gaussian rectangle tails, summed over local witnesses, make
  catastrophe weights smaller than every inverse power of \(\beta_F\);
  these bounds combine with Step 8 into III.10.2.
  \href{https://github.com/yang-mills-group/OpenAI-YM-20261028/blob/3456e5e75cf8d4cd558955a5c7638fba2d68f464/YM/AnalyticNext/CONNECTED_ACTIVITY_JETS_E6/Shutter.lean\#L119-L144}{\texttt{shoots}};
  \href{https://github.com/yang-mills-group/OpenAI-YM-20261028/blob/3456e5e75cf8d4cd558955a5c7638fba2d68f464/YM/AnalyticNext/CONNECTED_ACTIVITY_JETS_E6/Peak.lean\#L92-L188}{\texttt{Coat.ready}};
  \href{https://github.com/yang-mills-group/OpenAI-YM-20261028/blob/3456e5e75cf8d4cd558955a5c7638fba2d68f464/YM/AnalyticNext/CONNECTED_ACTIVITY_JETS_E6/ResponsePrice.lean\#L65-L112}{\texttt{Coat.alarm\_bills}};
  \href{https://github.com/yang-mills-group/OpenAI-YM-20261028/blob/3456e5e75cf8d4cd558955a5c7638fba2d68f464/YM/AnalyticNext/CONNECTED_ACTIVITY_JETS_E6/Issue.lean\#L94-L185}{\texttt{Target.issue}}.
\item
  \textbf{Differentiate letters and integrate molecules.} Compatible
  coordinate transport gives letter jets with rate-zero derivative scale
  and catastrophe probability \(\mathrm{rad}^{2\sharp v}\). Leibniz
  bounds and Gaussian moments, followed by Cauchy--Schwarz, retain
  \(\mathrm{rad}^{\sharp v}\) in the integrable molecule majorant; a
  source amplitude replaces its fare.
  \href{https://github.com/yang-mills-group/OpenAI-YM-20261028/blob/3456e5e75cf8d4cd558955a5c7638fba2d68f464/YM/AnalyticNext/CONNECTED_ACTIVITY_JETS_E6/Das.lean\#L49-L102}{\texttt{Utopia.whittle}};
  \href{https://github.com/yang-mills-group/OpenAI-YM-20261028/blob/3456e5e75cf8d4cd558955a5c7638fba2d68f464/YM/PhysicalNext/RG_NEXT_COMPLETE_REPEATABLE_NONLINEAR_STEP/Daisy/Charm.lean\#L58-L179}{\texttt{timing}}.
\item
  \textbf{Fuse and propagate markers.} Absolute convergence justifies
  Mayer regrouping; a retained forced-letter factor and rooted tree sum
  give \(e^{-\kappa}(e^\gamma-1)\), while common domination gives
  continuity. Connected ball-cover traversal proves I01; extra support
  decay and factorial rescaling give I02, composing centred embeddings
  gives I03, and subscale logarithmic column growth gives I04.
  \href{https://github.com/yang-mills-group/OpenAI-YM-20261028/blob/3456e5e75cf8d4cd558955a5c7638fba2d68f464/YM/Work/H01.lean\#L245-L251}{\texttt{H01.gathered}};
  \href{https://github.com/yang-mills-group/OpenAI-YM-20261028/blob/3456e5e75cf8d4cd558955a5c7638fba2d68f464/YM/Work/H02.lean\#L17-L80}{\texttt{H02.answer}};
  \href{https://github.com/yang-mills-group/OpenAI-YM-20261028/blob/3456e5e75cf8d4cd558955a5c7638fba2d68f464/YM/Work/H04.lean\#L155-L161}{\texttt{H04.answer}};
  \href{https://github.com/yang-mills-group/OpenAI-YM-20261028/blob/3456e5e75cf8d4cd558955a5c7638fba2d68f464/YM/Work/I01.lean\#L19-L82}{\texttt{I01.answer}};
  \href{https://github.com/yang-mills-group/OpenAI-YM-20261028/blob/3456e5e75cf8d4cd558955a5c7638fba2d68f464/YM/Work/I02.lean\#L159-L180}{\texttt{I02.answer}};
  \href{https://github.com/yang-mills-group/OpenAI-YM-20261028/blob/3456e5e75cf8d4cd558955a5c7638fba2d68f464/YM/Work/I03.lean\#L43-L63}{\texttt{I03.answer}};
  \href{https://github.com/yang-mills-group/OpenAI-YM-20261028/blob/3456e5e75cf8d4cd558955a5c7638fba2d68f464/YM/Work/I04.lean\#L76-L129}{\texttt{I04.answer}}.
  \emph{(proof read; H03 symmetry moves statement only)}
\item
  \textbf{Choose levels and periods.} Adjacent-stage floors and larger
  entry couplings enforce the energy estimates; finite-request
  scheduling supplies the remaining requirements. Periods \(4(M+p+1)\)
  admit arbitrarily late multiples of every side.
  \href{https://github.com/yang-mills-group/OpenAI-YM-20261028/blob/3456e5e75cf8d4cd558955a5c7638fba2d68f464/YM/Bid/Energy.lean\#L173-L202}{\texttt{Bid.pricing\_order}};
  \href{https://github.com/yang-mills-group/OpenAI-YM-20261028/blob/3456e5e75cf8d4cd558955a5c7638fba2d68f464/YM/PhysicalNext/RG_NEXT_COMPLETE_REPEATABLE_NONLINEAR_STEP/Pane/Divisors.lean\#L26-L135}{\texttt{Wait.engine}}.
\end{enumerate}

\subsubsection{5. Where the difficulty
lies}\label{where-the-difficulty-lies-12}

Fine-input coefficients must remain independent of \(L\), and mismatch
coefficients independent of \(W\), until those parameters absorb the
cross terms. Transported, truncated kernels lack exact completeness:
only their flat versions cancel, so local gauge fixing and tail
estimates must control the defects after neighbourhood summation.
Derivative integration must also retain a gain for every distinct
warning address under compatible fluctuation transport.

\subsubsection{6. Notation}\label{notation-12}

\(\sharp A\) counts all three seed types, including coincident types;
\(D_r(A)\) uses coarse max-distance. \(\Phi,\Phi',N\) are transported
towers, with subscript zero denoting flat uncut operators; \(v\) denotes
thinned coarse coordinates. \(N_B(H)\) is the minimum connected
radius-\(B\) ball-cover size. A function has rate \(a\) when, for every
\(v>a\), its absolute value is eventually at most \(\beta_F^v\)
(\href{https://github.com/yang-mills-group/OpenAI-YM-20261028/blob/3456e5e75cf8d4cd558955a5c7638fba2d68f464/YM/AnalyticNext/CONNECTED_ACTIVITY_JETS_E6/Numerics.lean\#L22-L23}{\texttt{Rate}}).
In III.10.3, \(\rho\) is the rate-zero jet scale and \(\mathrm{rad}\) is
the square root of a warning-probability bound; in III.10.4,
\(\Gamma,\widetilde\Gamma,m,h\) denote precision, approximate inverse,
mean and scaled stationarity residual.

\subsubsection{7. What remains
unexplained}\label{what-remains-unexplained-12}

\begin{itemize}
\tightlist
\item
  The full conversion of rooted letter and packet/tree bounds to III.9's
  outgoing \(E\)- and \(K\)-norms, including reset, remains unwritten.
\item
  Compatibility is constructed on the auxiliary torus by
  \href{https://github.com/yang-mills-group/OpenAI-YM-20261028/blob/3456e5e75cf8d4cd558955a5c7638fba2d68f464/YM/PhysicalNext/RG_NEXT_COMPLETE_REPEATABLE_NONLINEAR_STEP/Pane/Rooms.lean\#L39-L41}{\texttt{YM.PhysicalNext.Repeat.Pane.ovation}}
  assuming lattice symmetry of the beds at every period; that symmetry
  for the selected trajectory and the exact III.3 attribution remain
  unchecked.
\item
  Underlying stationary coercivity, local-inverse, packet-jet and
  individual letter-jet proofs have not all been checked through; H03's
  full symmetry lift also remains unchecked.
\item
  Additional reset estimates and the full use of
  \href{https://github.com/yang-mills-group/OpenAI-YM-20261028/blob/3456e5e75cf8d4cd558955a5c7638fba2d68f464/YM/PhysicalNext/V9_JOINT_ENDPOINT_LOAD/Supply.lean\#L72-L98}{\texttt{JointLoad.stage}}'s
  covariance and trigger bounds remain outside the numbered outputs. The
  trigger uses \(1+B\) and \(e^{(\log\log\beta_F)^2}\), absorbing fixed
  powers of the former into \(e^{-c e^{2(\log\log\beta_F)^2}}\) per
  distinct endpoint.
\item
  Exact correspondences with Balaban's lemmas, thresholds and
  normalizations remain open.
\end{itemize}

Detailed notes: the full account of this unit, with proof sketches for
every step, is on the
\href{https://github.com/scottnarmstrong/yang-mills-informalization/blob/main/outline/notes/III-10.md}{detailed
notes page}.

\subsection{III.11. Comparison with the free Gaussian measure:
interacting rank and determinant
budgets}\label{iii.11.-comparison-with-the-free-gaussian-measure-interacting-rank-and-determinant-budgets}

\textbf{Part:} III, Ultraviolet construction (renormalization group).
\textbf{Lean:} chapter 24:
\href{https://github.com/yang-mills-group/OpenAI-YM-20261028/tree/3456e5e75cf8d4cd558955a5c7638fba2d68f464/YM/Next/RANK}{\texttt{YM/Next/RANK/}},
\href{https://github.com/yang-mills-group/OpenAI-YM-20261028/tree/3456e5e75cf8d4cd558955a5c7638fba2d68f464/YM/PhysicalNext/AFTERNOON_INTERACTING_RANK_COMPARISON}{\texttt{YM/PhysicalNext/AFTERNOON\_INTERACTING\_RANK\_COMPARISON/}},
\href{https://github.com/yang-mills-group/OpenAI-YM-20261028/tree/3456e5e75cf8d4cd558955a5c7638fba2d68f464/YM/PhysicalNext/RANK_COVARIANCE_ATLAS}{\texttt{YM/PhysicalNext/RANK\_COVARIANCE\_ATLAS/}},
\href{https://github.com/yang-mills-group/OpenAI-YM-20261028/tree/3456e5e75cf8d4cd558955a5c7638fba2d68f464/YM/MergeProof}{\texttt{YM/MergeProof/}}
(with \texttt{West/}, \texttt{North/}, \texttt{East/}, \texttt{South/}).
\textbf{Size:} 42,556 lines (about 5,300 declarations).

\subsubsection{1. What this unit proves}\label{what-this-unit-proves-13}

The unit has two outputs. The first is the \textbf{rank comparison}. In
the continuum limit, the space \(F_d\) of local fields of degree \(d\)
is isomorphic to the corresponding space of the free Maxwell field,
compatibly with complex conjugation and with the lattice symmetry group.
Here \(F_d\) is the space of labels of degree \(d\) modulo labels whose
correlations vanish. The second output is the \textbf{stress tensor of
the limit}: a nonzero conserved stress tensor that satisfies the
translation and rotation Ward identities and is normalized against the
free Maxwell stress tensor along a scaling sequence. Both outputs rest
on quantitative estimates proved here: a per-site determinant budget,
chessboard bounds for connected parts of corrected observables, and
comparisons of covariances with the free curvature kernel. The trial
systems test a finite separating family of free one- and two-point
correlations.

Notation. Set \(\mu_k=\exp((\log\log\hat\beta_k)^2)\). The centered
estimate and the trial systems use this fixed schedule quantity. \(G\)
is a compact matrix group, assumed simple only where stated; \(\dim G\)
is the dimension of its Lie algebra; \(h\) is normalized Haar measure;
\(\mathrm{chord}(u)\ge0\) is the chordal distance from \(u\) to \(1\);
and \(\mathrm{shell}(r)=\{u\in G:\mathrm{chord}(u)\le r\}\). The torus
of side \(n\ge2\) is \(\Lambda_n=(\mathbb Z/n)^4\), with \(c=n^4\)
sites, \(4c\) bonds and \(6c\) plaquettes. The prior \(\mu_{\rm prior}\)
is product Haar measure on bonds. An RG representation \(R\) on
\(\Lambda_n\) has an inverse coupling \(\beta=\beta(R)\), an exponent
\(\mathrm{exponent}_R\) with constant part \(\mathrm{scalar}(R)\), and
an unnormalized density
\(e^{R}=e^{\mathrm{exponent}_R}\cdot(\text{polymer gas})\).

\textbf{Proposition III.11.1 (determinant budget on the torus, and its
per-site price).} (a) Let \(R\) be a lawful representation on
\(\Lambda_n\), \(n\ge2\), with lawful-class parameters
\((\sigma,A,\Lambda,\delta,w)\), and assume \(e^{R}\ge0\) pointwise. Let
\(0<s\le1\le b\) be a Haar budget (Step 1). Let \(g>0\) with
\(g\le\mathrm{threshold}(R)\) and \(\beta g^2=1\). Let \(F\) be a
\(\mu_{\rm prior}\)-integrable complex function and \(z\ge0\),
\(X\in\mathbb R\) such that
\(z|F(U)|\le\exp(\mathrm{exponent}_R(U)+cX)\) for every configuration
\(U\). Then
\[\frac{z\,\big|\int F\,d\mu_{\rm prior}\big|}{\int e^{R}\,d\mu_{\rm prior}}\;\le\;\exp\Big(c\,\pi(R,X,g)-4\dim G\,\log g\Big),\]
where the per-site price is
\[\pi=2\,\mathrm{exCost}(A,\Lambda,\delta)+6+\beta\,\mathrm{baseBudget}(A,g)+X+3\,\mathrm{ent}-4\log s,\qquad \mathrm{ent}=\log b-\dim G\,\log\sigma_R,\]
with \(\sigma_R=\min(1,\ell_R)\), where \(q_R\) and \(\ell_R\) are the
quadratic and linear clip coefficients. Here
\(\mathrm{exCost}=\delta+\mathrm{loopsPrice}\) bounds the loop and
additive activity terms per site, and
\(\mathrm{baseBudget}=3(q_R+A^2)g^2\) bounds the remaining action on
configurations whose bonds all lie in \(\mathrm{shell}(g/4)\). (b) Let
\((\mathcal A_m)\) be the levels of an RG schedule for simple \(G\), and
let \(O=(s,b)\) be a Haar budget. For all sufficiently large levels
\(m\), every period \(p\), every representation \(P\) on the level-\(m\)
torus (with \(c_{m,p}\) sites), every \(k\in\mathbb N\), every
\(X\le1+k\), and every \((g,\Lambda,\delta)\) admissible at level \(m\),
the following holds:
\[\pi(P,X,g)-\frac{4\dim G\,\log g}{c_{m,p}}\;\le\;\mathrm{work item}(O,m,k).\]
Admissible at level \(m\) means: same clip as the stage of \(m\),
\(\beta(P)\le4\hat\beta_m\), \(\beta(P)g^2=1\),
\(0<g\le\mathrm{threshold}(P)\le\mathrm{threshold}_m\), \(\delta\le1\),
\(\Lambda\le\mathrm{loopCap}_m\). The work item is
\[\mathrm{work item}=2(8\,\mathrm{paths}_m+2)+6+3(\mathrm{direct}_m+\mathrm{paths}_m^2)+(1+k)+3\,\mathrm{ent}-4\log s+1.\]
In (b), \(\pi\) uses \(A=\mathrm{paths}_m\) and the clip of stage \(m\).
The work item depends on that stage, \(O\) and \(k\), uniformly over the
admissible couplings and periods. \emph{Status:} theorems. \emph{Lean:}
\href{https://github.com/yang-mills-group/OpenAI-YM-20261028/blob/3456e5e75cf8d4cd558955a5c7638fba2d68f464/YM/Next/RANK/Determinant.lean\#L240-L274}{\texttt{determinant}},
\href{https://github.com/yang-mills-group/OpenAI-YM-20261028/blob/3456e5e75cf8d4cd558955a5c7638fba2d68f464/YM/Next/RANK/Determinant.lean\#L225-L227}{\texttt{protection}},
\href{https://github.com/yang-mills-group/OpenAI-YM-20261028/blob/3456e5e75cf8d4cd558955a5c7638fba2d68f464/YM/Next/RANK/Determinant.lean\#L26-L33}{\texttt{HaarBudget}};
(b)
\href{https://github.com/yang-mills-group/OpenAI-YM-20261028/blob/3456e5e75cf8d4cd558955a5c7638fba2d68f464/YM/Next/RANK/StationMeter.lean\#L119-L123}{\texttt{affordable}},
\href{https://github.com/yang-mills-group/OpenAI-YM-20261028/blob/3456e5e75cf8d4cd558955a5c7638fba2d68f464/YM/Next/RANK/StationMeter.lean\#L28-L36}{\texttt{Meets}},
\href{https://github.com/yang-mills-group/OpenAI-YM-20261028/blob/3456e5e75cf8d4cd558955a5c7638fba2d68f464/YM/Next/RANK/StationMeter.lean\#L41-L44}{\texttt{work\ item}}.

\textbf{Proposition III.11.2 (free pair terms).} Fix \(G\),
\(d,n\in\mathbb N\), and two colored plaquette derivative stencils
\(L,S\) based at the origin, with scaling degrees
\(\Delta_L=2+\mathrm{ord}(L)\le d\) and
\(\Delta_S=2+\mathrm{ord}(S)\le d\). There is \(C\ge0\), depending only
on these data, such that the following holds. Let
\(x\in(\mathbb R^4)^n\) have distinct entries and let \(i\ne j\).
Suppose \(D\ge0\) satisfies \(|\mathrm{Re}\,\Gamma_{(i,a),(j,b)}|\le D\)
for every pair of free-current keys \(a,b\) of scaling degree at most
\(d\). Then, for every \(r\ge1\),
\[r\,\Big|r^{\Delta_L+\Delta_S}\,\langle \mathrm{ev}_L\otimes\mathrm{ev}_S,B_r(i,j)\rangle-\mathrm{light}(L,S)(x_i,x_j)\Big|\le CD.\]
Here
\[B_r(i,j)=\sum_{a,b}\mathrm{Re}\,\Gamma_{(i,a),(j,b)}\;r^{-\mathrm{gen}(a)}t_a\otimes r^{-\mathrm{gen}(b)}t_b,\]
where \(t_a\) is the polynomial jet sample of key \(a\),
\(\mathrm{gen}(a)\) is its scaling degree, and \(\mathrm{ev}_L\) applies
the stencil \(L\). The kernel \(\mathrm{light}(L,S)\) is the
corresponding derivative of the free Maxwell curvature two-point
function. \emph{Status:} theorem. \emph{Lean:}
\href{https://github.com/yang-mills-group/OpenAI-YM-20261028/blob/3456e5e75cf8d4cd558955a5c7638fba2d68f464/YM/Next/RANK/StrideCost.lean\#L70-L98}{\texttt{Prayer.fence}},
\href{https://github.com/yang-mills-group/OpenAI-YM-20261028/blob/3456e5e75cf8d4cd558955a5c7638fba2d68f464/YM/Next/RANK/StrideCost.lean\#L64-L67}{\texttt{deck}}.

\textbf{Proposition III.11.3 (window estimate for the free covariance,
uniform in the regulator).} Let \(\alpha,\eta>0\),
\(\Omega\in\mathbb R\) and \(n\in\mathbb N\). Then there is \(C\ge0\)
with the following property. For every regulator \(f\) in the family
\(\mathrm{Firm}(\alpha,\Omega,\eta)\) of III.3, every \(G\), and every
pair of colored plaquette stencils \(a,b\), the frequency chart
\(K=K_f(a,b)\) is measurable, vanishes for \(|\xi|>2\), is \(C^n\) away
from \(0\), and satisfies:

\begin{itemize}
\tightlist
\item
  \(|K(0)|\le C\);
\item
  \(|\partial^iK(\xi)|\,|\xi|^i\le C\) for \(0<|\xi|\le2\) and
  \(i\le n\);
\item
  \(|\partial^i(K-H_{a,b})(\xi)|\,|\xi|^i\le C|\xi|\) for
  \(0<|\xi|\le1\) and \(i\le n\).
\end{itemize}

Here \(\partial^i\) denotes the \(i\)-th Fréchet derivative, measured in
operator norm, and \(H_{a,b}\) is the degree-zero Maxwell curvature
symbol for the plaquettes and colors of \(a,b\), with heat parameter
\((2\pi)^2\). If \(a,b\) have different colors, then \(K=H_{a,b}=0\).
\emph{Status:} theorem. \emph{Lean:}
\href{https://github.com/yang-mills-group/OpenAI-YM-20261028/blob/3456e5e75cf8d4cd558955a5c7638fba2d68f464/YM/PhysicalNext/RANK_COVARIANCE_ATLAS/Window.lean\#L118-L136}{\texttt{window\_prices}},
\href{https://github.com/yang-mills-group/OpenAI-YM-20261028/blob/3456e5e75cf8d4cd558955a5c7638fba2d68f464/YM/Next/RANK/Radian/Signed.lean\#L47-L53}{\texttt{Window}},
\href{https://github.com/yang-mills-group/OpenAI-YM-20261028/blob/3456e5e75cf8d4cd558955a5c7638fba2d68f464/YM/Next/RANK/Radian/Means.lean\#L16-L19}{\texttt{Cast}}.

\textbf{Proposition III.11.4 (frequent good scales produce corrected
trial systems).} Let \(G\) be simple and \(C\) a Wilson family
satisfying the construction record. Let \(d\) be a degree and
\(\mathcal F\) a stage package for degree \(d\) (fixed band
\(b_0\ge\max(6,\,s_d+3)\), with \(s_d\) the label's scaling saving, and
Taylor order \(\ge d\)). Let \(X\) consist of a real basis of
\(F_d^{\rm free}\), starting with the unit, and finitely many generic
one- and two-point tests that separate this space. Choose
representatives of that basis in the label space. Let \(O\) be a Haar
budget, \(J\) a chart ratio, and \(B\) a binding of the RG history of
\(C\) to its price lists. Fix \(w\in\mathbb R\). Assume that at
infinitely many levels \(k\) the following conditions hold
simultaneously:

\begin{itemize}
\tightlist
\item
  the four price conditions: tender, metering with \((O,w)\), and, for
  some common auxiliary parameters \(l,r\), hours with \((J,B,l,r,w)\)
  and vents with \((J,B,l,r)\);
\item
  the lantern-mark condition for \(X\).
\end{itemize}

Then there is a trial system \(Z\) for \(X\) that travels with the
station and satisfies the purchase conditions:

\begin{itemize}
\tightlist
\item
  depths \(k_1<k_2<\dots\) and test scales
  \(s_i=a_{k_i}\mu_{k_i}^{\,b'}\) (fixed \(b'>0\)) strictly decreasing
  to \(0\), with strictly increasing width levels
  \(\mathrm{width}_i\le k_i\) and the stage index constant on
  \([\mathrm{width}_i,k_i]\);
\item
  for each trial \(i\), coefficient vectors \(v_i(N,p)\) of corrected
  observables that are real, bounded uniformly in \((N,p)\), and have
  first entry the reference unit;
\item
  generic test configurations lying strictly inside the centred width
  cell, whose rescalings converge to the prescribed configuration of
  each row;
\item
  for each trial \(i\), eventually in \(N\) and for every period \(p\),
  on the centred tiling of the cutoff torus: for every row, every basis
  column and every nonempty set \(s\) of insertion modes, the chessboard
  root, at coupling \(\beta_{N,p}\), of the \(s\)-connected part of the
  corrected observable (taken relative to its free Maxwell one- and
  two-point profile) is at most \(\gamma_i=\mu_{k_i}^{-1/32}\to0\).
\end{itemize}

\emph{Status:} theorem, with the frequent-scale conditions carried as
hypotheses. In the main chain they are discharged as follows. In
\href{https://github.com/yang-mills-group/OpenAI-YM-20261028/blob/3456e5e75cf8d4cd558955a5c7638fba2d68f464/YM/MergeProof/Powder.lean\#L48-L71}{\texttt{Accent.rated}},
the lantern marks hold eventually in \(k\) by
\href{https://github.com/yang-mills-group/OpenAI-YM-20261028/blob/3456e5e75cf8d4cd558955a5c7638fba2d68f464/YM/MergeProof/SteamW.lean\#L54-L59}{\texttt{bursts}}.
The four price conditions are supplied to the rank assembly by
\href{https://github.com/yang-mills-group/OpenAI-YM-20261028/blob/3456e5e75cf8d4cd558955a5c7638fba2d68f464/YM/MergeProof/Snow.lean\#L55-L60}{\texttt{franchise}}
(this unit) inside the assembly
\href{https://github.com/yang-mills-group/OpenAI-YM-20261028/blob/3456e5e75cf8d4cd558955a5c7638fba2d68f464/YM/SharpGram/Assembly/Flowers.lean\#L36-L66}{\texttt{jointOffers}}.
That assembly is closed by
\href{https://github.com/yang-mills-group/OpenAI-YM-20261028/blob/3456e5e75cf8d4cd558955a5c7638fba2d68f464/YM/SharpGram/Assembly/Flowers.lean\#L75-L130}{\texttt{uv\_log}}
(III.13), whose only hypotheses are that \(G\) is simple and \(b>0\).
\emph{Lean:}
\href{https://github.com/yang-mills-group/OpenAI-YM-20261028/blob/3456e5e75cf8d4cd558955a5c7638fba2d68f464/YM/PhysicalNext/AFTERNOON_INTERACTING_RANK_COMPARISON/Soft/Star.lean\#L64-L66}{\texttt{cast\_off}},
\href{https://github.com/yang-mills-group/OpenAI-YM-20261028/blob/3456e5e75cf8d4cd558955a5c7638fba2d68f464/YM/PhysicalNext/AFTERNOON_INTERACTING_RANK_COMPARISON/Soft/Star.lean\#L68-L69}{\texttt{wheeled}},
\href{https://github.com/yang-mills-group/OpenAI-YM-20261028/blob/3456e5e75cf8d4cd558955a5c7638fba2d68f464/YM/PhysicalNext/AFTERNOON_INTERACTING_RANK_COMPARISON/Soft/Supra.lean\#L84-L95}{\texttt{v\_culminates}},
\href{https://github.com/yang-mills-group/OpenAI-YM-20261028/blob/3456e5e75cf8d4cd558955a5c7638fba2d68f464/YM/PhysicalNext/AFTERNOON_INTERACTING_RANK_COMPARISON/Soft/Supra.lean\#L177-L222}{\texttt{v\_obtain\_purchases}},
\href{https://github.com/yang-mills-group/OpenAI-YM-20261028/blob/3456e5e75cf8d4cd558955a5c7638fba2d68f464/YM/Next/RANK/Bought.lean\#L68-L84}{\texttt{Purchases}}.

\textbf{Proposition III.11.5 (the field spaces of the limit are those of
the free Maxwell field).} Let \(G\) be simple. Write \(\mathcal L_d\)
for abstract labels and \(\mathcal V_d\) for their fixed coordinate
space. Let \(C\) be a Wilson family satisfying the construction record
and the centred estimate: for each \(d\) and \(R,\varepsilon>0\) there
is \(A\ge0\) such that, at every level \(m\ge\mathrm{ref}(d)\),
eventually in \(N\) and for all \(p\), every label of degree \(d\) and
norm \(\le R\) has centred budget \(A\,a_m^{-d-\varepsilon}\). Assume:

\begin{itemize}
\tightlist
\item
  \begin{enumerate}
  \def\labelenumi{(\roman{enumi})}
  \tightlist
  \item
    a projector presentation of the labels: for each \(d\) and cutoff
    \((N,p)\), a projector \(\Pi_{d,N,p}\) on the coordinate label space
    \(\mathcal V_d\) whose kernel is the image of the Maxwell equation
    ideal \(\mathcal E_d\subset\mathcal L_d\) under an invertible
    corrected label map \(\mathcal L_d\to\mathcal V_d\). For each \(d\),
    both are real and equivariant for all sufficiently large \(N\),
    uniformly in \(p\). Assume also affordability and relay-cost bounds
    along every extraction;
  \end{enumerate}
\item
  (ii) the one-step cost package along the gauge-door maps between
  successive RG levels;
\item
  (iii) panels \(Q_d\), the mark exponents \(c_d=2(w_d+1)\), a chart
  ratio and a binding such that, for each \(d\), the four price
  conditions of III.11.4 hold simultaneously at infinitely many levels,
  with \(w=-3\). The panels are specialized to the family using \(c_d\).
\end{itemize}

Then for every extraction order \(f\) there are a refinement \(g\), a
smooth limit table \(T\) with a differentiation plan \(E\), and the
extraction record of IV.1 along \(f\circ g\) (raw limits, bins, spans,
cleaning to \(T\), Euclidean axioms), together with a seed record for
\((C,f\circ g,T)\). The seed record consists of: cutoff, ordered and
generic point limits, Euclidean axioms, cell hypothesis, and corrected
rank widths. Moreover the following holds: there is a real structure on
the free Maxwell table and, for every \(d\), a linear isomorphism
\[\Phi_d:\;F_d(T)\;\xrightarrow{\ \sim\ }\;F_d^{\rm free}=\mathcal L_d/\mathcal E_d\]
that intertwines complex conjugation and the action of the signed
permutations of the four axes. More precisely, there is a real,
equivariant linear isomorphism \(L_d:\mathcal L_d\to\mathcal V_d\)
carrying \(\mathcal E_d\) onto \(\mathrm{null}_d(T)\). Thus
\(\dim F_d(T)=\dim F_d^{\rm free}\); the identification of null spaces
uses this corrected map. \emph{Status:} theorem with hypotheses
(i)--(iii). All three are discharged in the main chain: (i) by
\href{https://github.com/yang-mills-group/OpenAI-YM-20261028/blob/3456e5e75cf8d4cd558955a5c7638fba2d68f464/YM/MergeProof/Eureka.lean\#L22-L30}{\texttt{reception}}
through
\href{https://github.com/yang-mills-group/OpenAI-YM-20261028/blob/3456e5e75cf8d4cd558955a5c7638fba2d68f464/YM/MergeProof/Inns.lean\#L33-L53}{\texttt{chaletrail}}
(this unit), (iii) by
\href{https://github.com/yang-mills-group/OpenAI-YM-20261028/blob/3456e5e75cf8d4cd558955a5c7638fba2d68f464/YM/MergeProof/Snow.lean\#L55-L60}{\texttt{franchise}},
and the binding by
\href{https://github.com/yang-mills-group/OpenAI-YM-20261028/blob/3456e5e75cf8d4cd558955a5c7638fba2d68f464/YM/MergeProof/Snow.lean\#L28-L43}{\texttt{castle}}.
The one-step package (ii) and the centred estimate are also included in
\href{https://github.com/yang-mills-group/OpenAI-YM-20261028/blob/3456e5e75cf8d4cd558955a5c7638fba2d68f464/YM/MergeProof/Almandine.lean\#L24-L36}{\texttt{Chateau}}
and consumed by
\href{https://github.com/yang-mills-group/OpenAI-YM-20261028/blob/3456e5e75cf8d4cd558955a5c7638fba2d68f464/YM/SharpGram/Assembly/Flowers.lean\#L75-L130}{\texttt{uv\_log}}
(III.13). \emph{Lean:}
\href{https://github.com/yang-mills-group/OpenAI-YM-20261028/blob/3456e5e75cf8d4cd558955a5c7638fba2d68f464/YM/MergeProof/Powder.lean\#L48-L71}{\texttt{Accent.rated}},
\href{https://github.com/yang-mills-group/OpenAI-YM-20261028/blob/3456e5e75cf8d4cd558955a5c7638fba2d68f464/YM/Euclid/Comparator.lean\#L135-L139}{\texttt{Maxwell.Ranks}},
\href{https://github.com/yang-mills-group/OpenAI-YM-20261028/blob/3456e5e75cf8d4cd558955a5c7638fba2d68f464/YM/Spec/Closed.lean\#L20-L26}{\texttt{Spec.Seed}},
\href{https://github.com/yang-mills-group/OpenAI-YM-20261028/blob/3456e5e75cf8d4cd558955a5c7638fba2d68f464/YM/MergeProof/Eureka.lean\#L16-L20}{\texttt{Scotia}},
\href{https://github.com/yang-mills-group/OpenAI-YM-20261028/blob/3456e5e75cf8d4cd558955a5c7638fba2d68f464/YM/Extraction/Center.lean\#L38-L40}{\texttt{Wcmain}}.

\textbf{Proposition III.11.6 (the stress tensor of the limit, normalized
against free Maxwell).} (a) Let \(C\) satisfy the construction record
and \(f\) be an extraction order with ordered and generic point limits,
giving a table \(S\). Assume:

\begin{itemize}
\tightlist
\item
  \(S\) satisfies the Euclidean axioms (with spin structure \(H\)) and
  the cell hypothesis;
\item
  \(W\) is a differentiation plan for \(S\) with its inclusion and
  derivative laws;
\item
  \(\dim F_d(S)\le\dim F_d^{\rm free}\) for every \(d\);
\item
  \(\mathcal T\) is a free Maxwell comparison model, with Euclidean
  axioms \(a\), inclusion, derivative and orientation laws, and a
  divergence-free stress;
\item
  the station hypothesis linking \(C\), \(f\) and \(a\) holds;
\item
  the shell estimate of \(C\) holds.
\end{itemize}

Then \(W\) satisfies the orientation law, and there are a stress record
\(I\) (stress tensor \(\mathcal J\), translation parts, spin rows) and a
scaling \(X\) of \(S\) such that:

\begin{itemize}
\tightlist
\item
  \(\mathcal J\neq0\) and \(\mathcal J\) lies in the kernel space of
  \(H\);
\item
  the scaled tables approach the isotropic panel, and the scaling is
  stable and compatible between the spin structures of \(S\) and
  \(\mathcal T\);
\item
  the scaled \(\mathcal J\) is metered against the free stress tensor
  \(\mathcal T.\mathcal J\);
\item
  \(\mathcal J\) is conserved and symmetric and satisfies the
  translation and rotation Ward identities (III.11.7);
\item
  the action law relating \(\mathcal J\) to the free model holds.
\end{itemize}

\begin{enumerate}
\def\labelenumi{(\alph{enumi})}
\setcounter{enumi}{1}
\tightlist
\item
  Let \(G\) be simple, let \(C\) carry the shell estimate, and assume
  the cofinal grid approximation condition and weak completeness of the
  corrected comparison tests for \(f\), together with the seed record of
  III.11.5 for \((C,f,S)\). Then for every differentiation plan \(W\)
  with inclusion and derivative laws there are Maxwell data \(M\) and a
  stress record \(I\) with \(\mathcal J\ne0\) satisfying the UV
  interface record \(\mathrm{UV}(C,f,S,W,M,I)\). \emph{Status:}
  theorems. In (b), the free model comes from
  \href{https://github.com/yang-mills-group/OpenAI-YM-20261028/blob/3456e5e75cf8d4cd558955a5c7638fba2d68f464/YM/Next/MAXWELL/Lit.lean\#L45-L48}{\texttt{maxwell\_supplies}}
  (VI.2). The station hypothesis is built from the seed and weak
  completeness, and the hypothesis
  \href{https://github.com/yang-mills-group/OpenAI-YM-20261028/blob/3456e5e75cf8d4cd558955a5c7638fba2d68f464/YM/Next/UV_STRESS/Invert.lean\#L108}{\texttt{DimAbove}}
  from the seed's corrected rank widths; this is how III.11.5 feeds
  III.11.6. The shell budget is proved by
  \href{https://github.com/yang-mills-group/OpenAI-YM-20261028/blob/3456e5e75cf8d4cd558955a5c7638fba2d68f464/YM/MergeProof/North/Economist.lean\#L23-L32}{\texttt{laurelE}}
  and
  \href{https://github.com/yang-mills-group/OpenAI-YM-20261028/blob/3456e5e75cf8d4cd558955a5c7638fba2d68f464/YM/MergeProof/North/Economist.lean\#L34}{\texttt{fugacious}}
  (this unit) inside
  \href{https://github.com/yang-mills-group/OpenAI-YM-20261028/blob/3456e5e75cf8d4cd558955a5c7638fba2d68f464/YM/SharpGram/Assembly/Flowers.lean\#L75-L130}{\texttt{uv\_log}}.
  \emph{Lean:}
  \href{https://github.com/yang-mills-group/OpenAI-YM-20261028/blob/3456e5e75cf8d4cd558955a5c7638fba2d68f464/YM/MergeProof/Emanate.lean\#L15-L30}{\texttt{Acme.readyA}},
  \href{https://github.com/yang-mills-group/OpenAI-YM-20261028/blob/3456e5e75cf8d4cd558955a5c7638fba2d68f464/YM/MergeProof/Emanate.lean\#L31-L37}{\texttt{Acme.emanate}},
  \href{https://github.com/yang-mills-group/OpenAI-YM-20261028/blob/3456e5e75cf8d4cd558955a5c7638fba2d68f464/YM/MergeProof/West/Ozone.lean\#L16-L19}{\texttt{Aroma}}.
\end{enumerate}

\textbf{Proposition III.11.7 (Ward identity: the flux of the stress
against a Killing field is a local field).} Let \(C\) satisfy the
construction record, let \(f\) be an extraction order with the cell
hypothesis and generic point limits \(S\), and suppose \(S\) satisfies
the Euclidean axioms with real structure \(H\). Let \(W\) satisfy the
derivative laws, and assume the shell estimate of III.11.6. Let:

\begin{itemize}
\tightlist
\item
  \(\mathcal J\) be a stress tensor of \(S\) that is conserved with
  respect to \(W\), and \(v\in F_d(S)\);
\item
  \(U\) be a smooth Killing field on \(\mathbb R^4\), homogeneous of
  degree \(p\) (\(U(ry)=r^pU(y)\) for \(r>0\));
\item
  \(d,D,p\in\mathbb N\) and \(D+3+p=d+4\).
\end{itemize}

Then there is \(B\in F_D(S)\) with the following property. For every
\(n\), all fields \(q_1,\dots,q_n\), every configuration of distinct
points \(x_0,\dots,x_n\), and every window \(\chi\) adapted to \(x\)
(smooth, compactly supported, equal to \(1\) near \(x_0\), with support
avoiding \(x_1,\dots,x_n\)), the flux of \(\mathcal J\) against
\(U(\cdot-x_0)\) and \(\chi\), inserted with \(v\) at \(x_0\) and
\(q_k\) at \(x_k\), equals the correlation of \(B\) at \(x_0\) with the
\(q_k\) at \(x_k\). The sign convention is
\[\mathrm{flux}=-\sum_{i,j}\int_{\mathbb R^4}\partial_i\chi(y)\,U_j(y-x_0)\,\langle\mathcal J_{ij}(y)\,v(x_0)\,q_1(x_1)\cdots q_n(x_n)\rangle\,dy.\]

Nonzero homogeneous Killing fields have \(p=0\) (translations,
\(D=d+1\)) or \(p=1\) (rotations, \(D=d\)). \emph{Status:} theorem.
\emph{Lean:}
\href{https://github.com/yang-mills-group/OpenAI-YM-20261028/blob/3456e5e75cf8d4cd558955a5c7638fba2d68f464/YM/MergeProof/West/Bravo.lean\#L55-L89}{\texttt{annular\_fieldA}},
\href{https://github.com/yang-mills-group/OpenAI-YM-20261028/blob/3456e5e75cf8d4cd558955a5c7638fba2d68f464/YM/MergeProof/West/Blanca.lean\#L49-L120}{\texttt{tenet}},
\href{https://github.com/yang-mills-group/OpenAI-YM-20261028/blob/3456e5e75cf8d4cd558955a5c7638fba2d68f464/YM/MergeProof/West/Orange.lean\#L33-L63}{\texttt{shells\_law}}.

\subsubsection{2. What it uses}\label{what-it-uses-13}

\begin{itemize}
\tightlist
\item
  \textbf{I.1, II.2:} the torus, Wilson measure, product Haar prior,
  lawful RG representations, construction record, cells and chessboard
  roots;
  \href{https://github.com/yang-mills-group/OpenAI-YM-20261028/blob/3456e5e75cf8d4cd558955a5c7638fba2d68f464/YM/Wilson/Law.lean\#L55}{\texttt{prior}},
  \href{https://github.com/yang-mills-group/OpenAI-YM-20261028/blob/3456e5e75cf8d4cd558955a5c7638fba2d68f464/YM/RG/History.lean\#L181-L183}{\texttt{Rep.Lawful}}.
\item
  \textbf{II.1:} label spaces, their real structures, the Maxwell
  equation ideal and the free Wick realization;
  \href{https://github.com/yang-mills-group/OpenAI-YM-20261028/blob/3456e5e75cf8d4cd558955a5c7638fba2d68f464/YM/Maxwell/Equation.lean\#L72}{\texttt{equationIdeal}}.
\item
  \textbf{III.2:} stages, thresholds, clip coefficients, loop caps,
  delays and waiting windows;
  \href{https://github.com/yang-mills-group/OpenAI-YM-20261028/blob/3456e5e75cf8d4cd558955a5c7638fba2d68f464/YM/MergeProof/North/Expound.lean\#L68-L77}{\texttt{ruese}}.
\item
  \textbf{III.3--III.4:} free covariances, regulator symbol bounds,
  Gaussian integration and window-to-kernel transfer;
  \href{https://github.com/yang-mills-group/OpenAI-YM-20261028/blob/3456e5e75cf8d4cd558955a5c7638fba2d68f464/YM/Next/RANK/Radian/Signed.lean\#L56-L92}{\texttt{intended}}.
\item
  \textbf{III.9--III.10:} coupling growth, exact RG histories and vacuum
  identities, one-step bounds and the waiting-window engine;
  \href{https://github.com/yang-mills-group/OpenAI-YM-20261028/blob/3456e5e75cf8d4cd558955a5c7638fba2d68f464/YM/PhysicalNext/RG_NEXT_COMPLETE_REPEATABLE_NONLINEAR_STEP/Pane/Divisors.lean\#L26-L135}{\texttt{Wait.engine}}.
\item
  \textbf{IV.1:} extraction, refinements, raw limits, spans and cleaning
  to smooth tables;
  \href{https://github.com/yang-mills-group/OpenAI-YM-20261028/blob/3456e5e75cf8d4cd558955a5c7638fba2d68f464/YM/Next/EXTRACTION/Base.lean\#L42-L51}{\texttt{Base}}.
\item
  \textbf{VI.1--VI.2:} windows and flux, finite faithfulness of testers,
  free Maxwell data and stress, and production of the current;
  \href{https://github.com/yang-mills-group/OpenAI-YM-20261028/blob/3456e5e75cf8d4cd558955a5c7638fba2d68f464/YM/Next/MAXWELL/Lit.lean\#L45-L48}{\texttt{maxwell\_supplies}},
  \href{https://github.com/yang-mills-group/OpenAI-YM-20261028/blob/3456e5e75cf8d4cd558955a5c7638fba2d68f464/YM/Next/UV_STRESS/Limit.lean\#L129-L167}{\texttt{produce\_current}}.
\item
  \textbf{III.13:} the common Wilson family closes the conditional
  chain;
  \href{https://github.com/yang-mills-group/OpenAI-YM-20261028/blob/3456e5e75cf8d4cd558955a5c7638fba2d68f464/YM/SharpGram/Assembly/Flowers.lean\#L75-L130}{\texttt{uv\_log}}
  has only the hypotheses that \(G\) is simple and \(b>0\).
\end{itemize}

\textbf{Standard results:} the chessboard estimate and reflection
positivity of the Wilson action; axial tree gauge and Haar invariance;
small-ball Haar volume estimates; Wick expansion; Schur complements;
frequency-symbol estimates; finite-dimensional duality; and Ward
identities for conserved symmetric stress tensors and Killing fields.

\textbf{Constants and assumptions:} the Haar budget is fixed
independently of stages and volumes by
\href{https://github.com/yang-mills-group/OpenAI-YM-20261028/blob/3456e5e75cf8d4cd558955a5c7638fba2d68f464/YM/Next/RANK/Haar.lean\#L170-L211}{\texttt{integral\_budget}}.
The projector, one-step, centred, frequent-scale, binding, shell,
grid-approximation and weak-completeness hypotheses remain explicit in
§1; their main-chain producers are identified there. The older one-step
interfaces in
\href{https://github.com/yang-mills-group/OpenAI-YM-20261028/blob/3456e5e75cf8d4cd558955a5c7638fba2d68f464/YM/Next/RANK/LedgerGate.lean}{\texttt{LedgerGate.lean}}
are conditional definitions, not consequences of lawful representations.

\subsubsection{3. Strategy}\label{strategy-13}

The Maxwell equations give the upper bound on the interacting
field-space dimensions. Corrected observables tested at short distances
give the reverse bound: their one- and two-point correlations approach a
finite separating family of free correlations, so their classes are
eventually independent. The two bounds identify the limit null space
with a limit projector kernel. Conjugating a cutoff projector to the
limit projector then gives the isomorphism compatible with conjugation
and signed permutations.

The quantitative comparison uses the exact RG history to carry reflected
products to a coarse torus, then bounds a ratio of Haar integrals.
Plaquette elimination gives \(3c-3\) numerator factors, while tree gauge
leaves \(3c+1\) vacuum factors; the remaining coupling cost is only
\(g^{-4\dim G}=\beta^{2\dim G}\) in total. Its chessboard root is
affordable, allowing connected errors to tend to zero along infinitely
many good scales. For the stress tensor, comparison with free Maxwell
fixes its normalization along a scaling sequence; conservation, shell
estimates and finite-dimensional duality represent each Killing-field
flux by a local field.

\subsubsection{4. Steps}\label{steps-13}

\begin{enumerate}
\def\labelenumi{\arabic{enumi}.}
\item
  \textbf{Haar budget.} There are \(0<s\le1\le b\) with
  \(h(\mathrm{shell}(r/4))\ge sr^{\dim G}\) for \(0<r\le1\) and
  \(\int e^{-\mathrm{chord}(u)/t}\,dh(u)\le bt^{\dim G}\) for \(t>0\).
  Exponential coordinates give the local volume estimates; compactness
  and a shell sum give the integral estimate.
  \href{https://github.com/yang-mills-group/OpenAI-YM-20261028/blob/3456e5e75cf8d4cd558955a5c7638fba2d68f464/YM/Next/RANK/Haar.lean\#L170-L211}{\texttt{integral\_budget}}
\item
  \textbf{Plaquette elimination.} For \(n\ge2\) and continuous
  \(\varphi:G\to[0,1]\),
  \(\int\prod_p\varphi(U_p)\,d\mu_{\rm prior}\le(\int\varphi\,dh)^{3c-3}\).
  A triangular sweep assigns \(3c-3\) plaquettes to distinct bonds,
  whose successive Haar integrations factor; the unselected factors are
  at most \(1\).
  \href{https://github.com/yang-mills-group/OpenAI-YM-20261028/blob/3456e5e75cf8d4cd558955a5c7638fba2d68f464/YM/Next/RANK/Axial.lean\#L368-L383}{\texttt{elimination}},
  \href{https://github.com/yang-mills-group/OpenAI-YM-20261028/blob/3456e5e75cf8d4cd558955a5c7638fba2d68f464/YM/Next/RANK/Axial.lean\#L347-L365}{\texttt{charged\_integral}}
  \emph{(proof read; sweep count and triangularity statement only)}
\item
  \textbf{Vacuum and numerator.} Under III.11.1(a), tree gauge and
  restriction of the \(3c+1\) free bonds to \(\mathrm{shell}(g/4)\) give
  the vacuum lower bound, while the clipped action retains exponential
  chord decay for the numerator. Applying Steps 1--2 and dividing gives
  \(\exp(c\pi)g^{-4\dim G}\); the four-factor difference between the
  counts accounts for the remaining power of \(g\).
  \href{https://github.com/yang-mills-group/OpenAI-YM-20261028/blob/3456e5e75cf8d4cd558955a5c7638fba2d68f464/YM/Next/RANK/Determinant.lean\#L76-L111}{\texttt{vacuum\_bound}},
  \href{https://github.com/yang-mills-group/OpenAI-YM-20261028/blob/3456e5e75cf8d4cd558955a5c7638fba2d68f464/YM/Next/RANK/Determinant.lean\#L153-L181}{\texttt{exponent\_retains}},
  \href{https://github.com/yang-mills-group/OpenAI-YM-20261028/blob/3456e5e75cf8d4cd558955a5c7638fba2d68f464/YM/Next/RANK/Determinant.lean\#L240-L274}{\texttt{determinant}}
\item
  \textbf{Per-site price and chessboard root.} Under III.11.1(b), the
  stage bounds give \(\pi\le\mathrm{work item}-1\) and
  \(-4\dim G\log g/c_{m,p}\le1\) eventually, uniformly in admissible
  parameters and periods. In the trial comparison, the exact vacuum
  identity and local activity decomposition transfer the determinant
  estimate to the chessboard root; deleting additive activities costs
  the \(2\delta\) already included in the price. Allocation and
  affordability, with allowance \(e^{2I}\) and
  \(I\ge\mathrm{work item}\), then give root error \(e^{-I}\) for every
  nonempty connected part.
  \href{https://github.com/yang-mills-group/OpenAI-YM-20261028/blob/3456e5e75cf8d4cd558955a5c7638fba2d68f464/YM/Next/RANK/StationMeter.lean\#L119-L123}{\texttt{affordable}},
  \href{https://github.com/yang-mills-group/OpenAI-YM-20261028/blob/3456e5e75cf8d4cd558955a5c7638fba2d68f464/YM/PhysicalNext/AFTERNOON_INTERACTING_RANK_COMPARISON/Soft/Denarii.lean\#L49-L85}{\texttt{s\_postage}},
  \href{https://github.com/yang-mills-group/OpenAI-YM-20261028/blob/3456e5e75cf8d4cd558955a5c7638fba2d68f464/YM/PhysicalNext/AFTERNOON_INTERACTING_RANK_COMPARISON/Soft/Denarii.lean\#L90-L123}{\texttt{coupon}}
  \emph{(proof read; flow identity and cell-stamp bookkeeping statement
  only)}
\item
  \textbf{Free pair and window estimates.} In III.11.2, stencil
  evaluation kills keys below the stencil degree, the equal-degree terms
  give the free curvature kernel, and every remaining term gains
  \(r^{-1}\), giving the bound \(CD/r\) on the scaled remainder. For
  III.11.3, the regulator chart is a symbol of order zero and differs
  from the heat symbol by a symbol of order one; the change of variables
  \(\zeta=2\pi\xi\) and the finite number of plaquette pairs give a
  uniform constant.
  \href{https://github.com/yang-mills-group/OpenAI-YM-20261028/blob/3456e5e75cf8d4cd558955a5c7638fba2d68f464/YM/Next/RANK/StrideCost.lean\#L70-L98}{\texttt{Prayer.fence}},
  \href{https://github.com/yang-mills-group/OpenAI-YM-20261028/blob/3456e5e75cf8d4cd558955a5c7638fba2d68f464/YM/PhysicalNext/RANK_COVARIANCE_ATLAS/Window.lean\#L118-L136}{\texttt{window\_prices}}
\item
  \textbf{Covariance splitting.} For an atlas compatible with the
  schedule and retained unmasked quadratic forms, satisfying symmetry,
  coercivity, inverse and crossing laws, and a step \(i<N\) with a tree
  forest and the required torus size, the complete curvature covariances
  satisfy
  \[\mathcal C_{i+1}-W_i\mathcal C_iW_i^{\mathsf T}=d_f\Gamma_i d_f^{\mathsf T},\qquad\Gamma_i\succeq0.\]
  This is the Schur-complement identity on the curl image; uniform
  resolvent estimates give the propagation constant for fixed kernel
  choices and stage indices \(j,t\ge4\) in the range \(d-1\le j-2\)
  (natural-number subtraction).
  \href{https://github.com/yang-mills-group/OpenAI-YM-20261028/blob/3456e5e75cf8d4cd558955a5c7638fba2d68f464/YM/PhysicalNext/RANK_COVARIANCE_ATLAS/Identity.lean\#L198-L213}{\texttt{one\_fold}},
  \href{https://github.com/yang-mills-group/OpenAI-YM-20261028/blob/3456e5e75cf8d4cd558955a5c7638fba2d68f464/YM/PhysicalNext/RANK_COVARIANCE_ATLAS/Step.lean\#L106-L116}{\texttt{train\_issue}}
\item
  \textbf{Local polynomial comparison.} For regular local observables
  with the prescribed derivative bounds, Taylor expansion in link
  coordinates gives bounded polynomial coefficients and a remainder of
  the required order. For zero-mean formal Gaussian data with symmetric
  covariance vanishing within each location, Wick expansion runs over
  cross-location pairings without requiring positive covariance; under
  the comparison's common-mask, support, coefficient, mean and
  covariance hypotheses, these moments give the corrected
  connected-component bounds. The allocation, integrability, locality
  and compatibility hypotheses give exact sector decomposition and
  factorization over disjoint supports.
  \href{https://github.com/yang-mills-group/OpenAI-YM-20261028/blob/3456e5e75cf8d4cd558955a5c7638fba2d68f464/YM/Next/RANK/LineBox.lean\#L168-L205}{\texttt{minted}},
  \href{https://github.com/yang-mills-group/OpenAI-YM-20261028/blob/3456e5e75cf8d4cd558955a5c7638fba2d68f464/YM/Next/RANK/GaussTrack.lean\#L131-L173}{\texttt{headed}},
  \href{https://github.com/yang-mills-group/OpenAI-YM-20261028/blob/3456e5e75cf8d4cd558955a5c7638fba2d68f464/YM/Next/RANK/Festoon.lean\#L57-L126}{\texttt{console}},
  \href{https://github.com/yang-mills-group/OpenAI-YM-20261028/blob/3456e5e75cf8d4cd558955a5c7638fba2d68f464/YM/PhysicalNext/AFTERNOON_INTERACTING_RANK_COMPARISON/Rendered.lean\#L36-L86}{\texttt{rendered}},
  \href{https://github.com/yang-mills-group/OpenAI-YM-20261028/blob/3456e5e75cf8d4cd558955a5c7638fba2d68f464/YM/PhysicalNext/AFTERNOON_INTERACTING_RANK_COMPARISON/Sight.lean\#L132-L158}{\texttt{peer\_beams}}
  \emph{(proof read; local comparison and allocation inputs statement
  only)}
\item
  \textbf{Corrected trials.} Under III.11.4, one correction map and
  polynomial approximations serve all test rows; the free-profile
  correction cancels the leading terms, and the one-step estimates
  propagate the remaining source bounds from the middle level to the
  width level. Steps 3--4 give the purchase bounds, and an inductive
  choice among the good levels gives increasing depths and widths,
  decreasing test scales, and bounded real coefficient vectors with the
  reference unit first.
  \href{https://github.com/yang-mills-group/OpenAI-YM-20261028/blob/3456e5e75cf8d4cd558955a5c7638fba2d68f464/YM/PhysicalNext/AFTERNOON_INTERACTING_RANK_COMPARISON/Soft/Diana.lean\#L57-L104}{\texttt{v\_paid\_factory}},
  \href{https://github.com/yang-mills-group/OpenAI-YM-20261028/blob/3456e5e75cf8d4cd558955a5c7638fba2d68f464/YM/PhysicalNext/AFTERNOON_INTERACTING_RANK_COMPARISON/Soft/Supra.lean\#L84-L95}{\texttt{v\_culminates}},
  \href{https://github.com/yang-mills-group/OpenAI-YM-20261028/blob/3456e5e75cf8d4cd558955a5c7638fba2d68f464/YM/PhysicalNext/AFTERNOON_INTERACTING_RANK_COMPARISON/Soft/Supra.lean\#L29-L53}{\texttt{v\_conditioned}},
  \href{https://github.com/yang-mills-group/OpenAI-YM-20261028/blob/3456e5e75cf8d4cd558955a5c7638fba2d68f464/YM/PhysicalNext/AFTERNOON_INTERACTING_RANK_COMPARISON/Soft/Supra.lean\#L177-L222}{\texttt{v\_obtain\_purchases}}
  \emph{(proof read; analytic inputs and helper estimates statement
  only)}
\item
  \textbf{Supply of good levels.} For simple \(G\), the construction
  record and one-step package with the prescribed weights and
  multiplicities at least \(1\), a fixed degree-\(d\) column family and
  stage package with mark exponent \(c\ge2(w_d+1)\) satisfy the
  lantern-mark condition at every sufficiently large metered level.
  Compatible pre-\(L\) prices and one-step bounds, stage thresholds
  below block lengths, eventual presence of each degree index,
  nonnegative mark exponents and a Haar budget give panels and an
  enlarged waiting window on which every Wilson family paid on that
  window and every binding with those prices have infinitely many
  simultaneous good levels for each \(d\).
  \href{https://github.com/yang-mills-group/OpenAI-YM-20261028/blob/3456e5e75cf8d4cd558955a5c7638fba2d68f464/YM/MergeProof/SteamW.lean\#L54-L59}{\texttt{bursts}},
  \href{https://github.com/yang-mills-group/OpenAI-YM-20261028/blob/3456e5e75cf8d4cd558955a5c7638fba2d68f464/YM/MergeProof/Snow.lean\#L55-L60}{\texttt{franchise}}
  \emph{(proof structure read; analytic inputs and helpers statement
  only)}
\item
  \textbf{Equations and independent columns.} Under III.11.5(i), free
  probes separate the on-shell sector and give
  \(\dim\mathcal L_d\le\dim\ker\Pi_{d,N,p}+\dim F_d^{\rm free}\) at each
  cutoff. For a Wilson family satisfying the construction record, with
  generic point limits and continuous kernels integrable against every
  separated test function, purchases and iterated coefficient limits
  give moment errors \(P_k\gamma_k\to0\), so the limiting corrected
  columns are independent modulo the integrated-test null space for all
  sufficiently large trials. Relay estimates put the limit projector
  kernel inside the null space; with the centred port, raw limits,
  convergent relay projectors, equation counts and trials in every
  degree, extraction and the dimension comparison give a smooth table
  \(T\), its seed, and \(\mathrm{null}_d(T)=\ker\Pi_\infty\).
  \href{https://github.com/yang-mills-group/OpenAI-YM-20261028/blob/3456e5e75cf8d4cd558955a5c7638fba2d68f464/YM/MergeProof/Powder.lean\#L19-L47}{\texttt{savings}},
  \href{https://github.com/yang-mills-group/OpenAI-YM-20261028/blob/3456e5e75cf8d4cd558955a5c7638fba2d68f464/YM/Next/RANK/Monitors.lean\#L71-L87}{\texttt{measured}},
  \href{https://github.com/yang-mills-group/OpenAI-YM-20261028/blob/3456e5e75cf8d4cd558955a5c7638fba2d68f464/YM/Next/RANK/Beam.lean\#L153-L180}{\texttt{arriving\_seed}}
  \emph{(proof read; free positivity, chessboard and analytic inputs
  statement only)}
\item
  \textbf{Equivariant identification.} Under the projector presentation
  and its extraction record, if \(A_i\to\Pi_\infty\), then
  \(Y_i=\Pi_\infty A_i+(1-\Pi_\infty)(1-A_i)\to1\) and
  \(Y_iA_i=\Pi_\infty Y_i\). For large \(i\), \(Y_i\) is invertible,
  real and equivariant; composing it with the corrected label map
  carries \(\mathcal E_d\) onto \(\mathrm{null}_d(T)\) and induces
  III.11.5's quotient isomorphism.
  \href{https://github.com/yang-mills-group/OpenAI-YM-20261028/blob/3456e5e75cf8d4cd558955a5c7638fba2d68f464/YM/MergeProof/Thrush.lean\#L13-L68}{\texttt{drumroll}},
  \href{https://github.com/yang-mills-group/OpenAI-YM-20261028/blob/3456e5e75cf8d4cd558955a5c7638fba2d68f464/YM/MergeProof/Thrush.lean\#L72-L114}{\texttt{thrush}},
  \href{https://github.com/yang-mills-group/OpenAI-YM-20261028/blob/3456e5e75cf8d4cd558955a5c7638fba2d68f464/YM/MergeProof/Powder.lean\#L48-L71}{\texttt{Accent.rated}}
\item
  \textbf{Flux becomes a local field.} Under III.11.7, conservation
  permits windows on shrinking annuli; the shell estimate shows that
  every finite combination of grid testers annihilating \(F_D\) also
  annihilates the flux. Finite-dimensional duality supplies a matching
  field for each finite tester family, finite faithfulness gives a
  single field for all testers, and translation covariance, continuity
  and grid density extend the identity to all distinct configurations.
  Constant and rotational Killing fields give the translation and
  rotation maps.
  \href{https://github.com/yang-mills-group/OpenAI-YM-20261028/blob/3456e5e75cf8d4cd558955a5c7638fba2d68f464/YM/MergeProof/West/Blanca.lean\#L49-L120}{\texttt{tenet}},
  \href{https://github.com/yang-mills-group/OpenAI-YM-20261028/blob/3456e5e75cf8d4cd558955a5c7638fba2d68f464/YM/MergeProof/West/Bravo.lean\#L55-L89}{\texttt{annular\_fieldA}},
  \href{https://github.com/yang-mills-group/OpenAI-YM-20261028/blob/3456e5e75cf8d4cd558955a5c7638fba2d68f464/YM/MergeProof/West/Orange.lean\#L33-L63}{\texttt{shells\_law}}
  \emph{(proof read; annulus comparison inputs statement only)}
\item
  \textbf{Stress and ultraviolet records.} Under III.11.6(a), the
  orientation law and station equivalences allow the free comparison to
  produce a nonzero current and scaling sequence; zero current
  divergence gives conservation, the weak scaling comparison gives the
  action law, and Step 12 gives the stress laws. Under III.11.6(b),
  canonical free Maxwell data, the seed's corrected rank widths and weak
  completeness supply the comparison model, dimension bound and station
  hypothesis needed for (a).
  \href{https://github.com/yang-mills-group/OpenAI-YM-20261028/blob/3456e5e75cf8d4cd558955a5c7638fba2d68f464/YM/MergeProof/Emanate.lean\#L15-L30}{\texttt{Acme.readyA}},
  \href{https://github.com/yang-mills-group/OpenAI-YM-20261028/blob/3456e5e75cf8d4cd558955a5c7638fba2d68f464/YM/MergeProof/Emanate.lean\#L31-L37}{\texttt{Acme.emanate}}
  \emph{(proof read; helper estimates statement only)}
\end{enumerate}

\subsubsection{5. Where the difficulty
lies}\label{where-the-difficulty-lies-13}

The determinant comparison must avoid a per-site loss of \(\log\beta\):
that would overwhelm the vanishing chessboard error. Tree-gauge counting
leaves only the volume-independent factor \(\beta^{2\dim G}\); this is a
Haar integral estimate rather than a Hessian determinant calculation.
The lower rank bound requires short-distance comparison at infinitely
many scales, whereas the upper bound is algebraic at every cutoff. The
isomorphism comes from conjugating projectors, and its symmetry
statement concerns complex representations with compatible conjugations
under signed permutations. Constants must be chosen before block lengths
and uniformly in the period; eventual estimates require the specified
delays before use on comparison windows.

\subsubsection{6. Notation}\label{notation-13}

{\def\LTcaptype{none} % do not increment counter
\begin{longtable}[]{@{}
  >{\raggedright\arraybackslash}p{(\linewidth - 2\tabcolsep) * \real{0.5000}}
  >{\raggedright\arraybackslash}p{(\linewidth - 2\tabcolsep) * \real{0.5000}}@{}}
\toprule\noalign{}
\begin{minipage}[b]{\linewidth}\raggedright
Symbol
\end{minipage} & \begin{minipage}[b]{\linewidth}\raggedright
Meaning and Lean link
\end{minipage} \\
\midrule\noalign{}
\endhead
\bottomrule\noalign{}
\endlastfoot
\(\mu_k\), \(w_d\) & Fixed schedule factor
\(\exp((\log\log\hat\beta_k)^2)\); coefficient exponent.
\href{https://github.com/yang-mills-group/OpenAI-YM-20261028/blob/3456e5e75cf8d4cd558955a5c7638fba2d68f464/YM/RG/History.lean\#L81}{\texttt{monitor}},
\href{https://github.com/yang-mills-group/OpenAI-YM-20261028/blob/3456e5e75cf8d4cd558955a5c7638fba2d68f464/YM/MergeProof/Tender.lean\#L11}{\texttt{winter}} \\
\((s,b)\), \(\sigma_R\), \(\pi\), work item & Haar budget, clip
strength, per-site determinant price and stage allowance of III.11.1.
\href{https://github.com/yang-mills-group/OpenAI-YM-20261028/blob/3456e5e75cf8d4cd558955a5c7638fba2d68f464/YM/Next/RANK/Determinant.lean\#L26-L33}{\texttt{HaarBudget}},
\href{https://github.com/yang-mills-group/OpenAI-YM-20261028/blob/3456e5e75cf8d4cd558955a5c7638fba2d68f464/YM/Next/RANK/Determinant.lean\#L116}{\texttt{strength}},
\href{https://github.com/yang-mills-group/OpenAI-YM-20261028/blob/3456e5e75cf8d4cd558955a5c7638fba2d68f464/YM/Next/RANK/Determinant.lean\#L225-L227}{\texttt{protection}},
\href{https://github.com/yang-mills-group/OpenAI-YM-20261028/blob/3456e5e75cf8d4cd558955a5c7638fba2d68f464/YM/Next/RANK/StationMeter.lean\#L41-L44}{\texttt{work\ item}} \\
\(\mathcal L_d\), \(\mathcal V_d\), \(\mathcal E_d\) & Abstract labels,
coordinate labels and Maxwell equation ideal;
\(F_d^{\rm free}=\mathcal L_d/\mathcal E_d\).
\href{https://github.com/yang-mills-group/OpenAI-YM-20261028/blob/3456e5e75cf8d4cd558955a5c7638fba2d68f464/YM/Maxwell/Equation.lean\#L72}{\texttt{equationIdeal}} \\
\(F_d(T)\), \(\mathrm{null}_d(T)\) & Local fields of degree \(d\) and
labels with vanishing correlations.
\href{https://github.com/yang-mills-group/OpenAI-YM-20261028/blob/3456e5e75cf8d4cd558955a5c7638fba2d68f464/YM/Euclid/Comparator.lean\#L135-L139}{\texttt{Maxwell.Ranks}} \\
\(\Pi_{d,N,p}\), \(\Pi_\infty\) & Cutoff and limit projectors; their
kernels identify the equations after correction.
\href{https://github.com/yang-mills-group/OpenAI-YM-20261028/blob/3456e5e75cf8d4cd558955a5c7638fba2d68f464/YM/MergeProof/Eureka.lean\#L16-L20}{\texttt{Scotia}} \\
Stage package and price conditions & Stage data; tender, metering,
hours, vents and lantern marks in III.11.4.
\href{https://github.com/yang-mills-group/OpenAI-YM-20261028/blob/3456e5e75cf8d4cd558955a5c7638fba2d68f464/YM/Next/RANK/Concord.lean\#L38-L51}{\texttt{Cohesion}},
\href{https://github.com/yang-mills-group/OpenAI-YM-20261028/blob/3456e5e75cf8d4cd558955a5c7638fba2d68f464/YM/PhysicalNext/RANK_COVARIANCE_ATLAS/Consumer.lean\#L71}{\texttt{Tender}},
\href{https://github.com/yang-mills-group/OpenAI-YM-20261028/blob/3456e5e75cf8d4cd558955a5c7638fba2d68f464/YM/Next/RANK/Concord.lean\#L62-L64}{\texttt{Metering}},
\href{https://github.com/yang-mills-group/OpenAI-YM-20261028/blob/3456e5e75cf8d4cd558955a5c7638fba2d68f464/YM/PhysicalNext/AFTERNOON_INTERACTING_RANK_COMPARISON/Pleas.lean\#L21-L26}{\texttt{ShowHours}},
\href{https://github.com/yang-mills-group/OpenAI-YM-20261028/blob/3456e5e75cf8d4cd558955a5c7638fba2d68f464/YM/PhysicalNext/AFTERNOON_INTERACTING_RANK_COMPARISON/Soft/Postsale.lean\#L19-L21}{\texttt{Vents}},
\href{https://github.com/yang-mills-group/OpenAI-YM-20261028/blob/3456e5e75cf8d4cd558955a5c7638fba2d68f464/YM/PhysicalNext/AFTERNOON_INTERACTING_RANK_COMPARISON/Proceed.lean\#L74-L77}{\texttt{LanternMarks}} \\
\(\mathcal C_i\), \(W_i\), \(d_f\), \(\Gamma_i\) & Complete curvature
covariance, passage map, fine exterior derivative and uncut tree
covariance.
\href{https://github.com/yang-mills-group/OpenAI-YM-20261028/blob/3456e5e75cf8d4cd558955a5c7638fba2d68f464/YM/PhysicalNext/RANK_COVARIANCE_ATLAS/Identity.lean\#L198-L213}{\texttt{one\_fold}} \\
\(\mathcal J\), shell budget & Stress tensor; annular estimates with
radii tending to zero.
\href{https://github.com/yang-mills-group/OpenAI-YM-20261028/blob/3456e5e75cf8d4cd558955a5c7638fba2d68f464/YM/MergeProof/West/Ozone.lean\#L16-L19}{\texttt{Aroma}} \\
Seed, station, grid approximation, weak completeness & Records used in
III.11.5--III.11.6.
\href{https://github.com/yang-mills-group/OpenAI-YM-20261028/blob/3456e5e75cf8d4cd558955a5c7638fba2d68f464/YM/Spec/Closed.lean\#L20-L26}{\texttt{Spec.Seed}},
\href{https://github.com/yang-mills-group/OpenAI-YM-20261028/blob/3456e5e75cf8d4cd558955a5c7638fba2d68f464/YM/Next/UV_STRESS/Station.lean\#L112-L114}{\texttt{Stations}},
\href{https://github.com/yang-mills-group/OpenAI-YM-20261028/blob/3456e5e75cf8d4cd558955a5c7638fba2d68f464/YM/Org/Input.lean\#L121-L123}{\texttt{Grill}},
\href{https://github.com/yang-mills-group/OpenAI-YM-20261028/blob/3456e5e75cf8d4cd558955a5c7638fba2d68f464/YM/MergeProof/East/Gran.lean\#L35}{\texttt{CompleteWeak}} \\
\end{longtable}
}

\subsubsection{7. What remains
unexplained}\label{what-remains-unexplained-13}

\begin{itemize}
\tightlist
\item
  The discrete stencil, interpolation and shell estimates underlying the
  annulus bound have not been independently checked through their
  analytic dependencies;
  \href{https://github.com/yang-mills-group/OpenAI-YM-20261028/blob/3456e5e75cf8d4cd558955a5c7638fba2d68f464/YM/MergeProof/West/Ozone.lean\#L21-L45}{\texttt{zone}}.
\item
  The nonlinear integration and remainder estimates beneath the one-step
  comparison were not checked in full. The correspondence between the
  older conditional interfaces in
  \href{https://github.com/yang-mills-group/OpenAI-YM-20261028/blob/3456e5e75cf8d4cd558955a5c7638fba2d68f464/YM/Next/RANK/LedgerGate.lean}{\texttt{LedgerGate.lean}}
  and all restarted replacements remains to be documented.
\item
  Complete mathematical estimates for the stage, launch, hours, vents,
  relay, station and shell records remain to be written. Their
  structural hypotheses cannot be replaced by the displayed inequalities
  alone.
\item
  Sweep triangularity and the auxiliary estimates were checked at
  statement level. The lower-level proofs supplying the centred
  estimate, one-step package and schedule existence were not checked in
  full.
\end{itemize}

Detailed notes: the full account of this unit, with proof sketches for
every step, is on the
\href{https://github.com/scottnarmstrong/yang-mills-informalization/blob/main/outline/notes/III-11.md}{detailed
notes page}.

\subsection{III.12. Bounds uniform in the volume and convergence as the
spacing tends to
zero}\label{iii.12.-bounds-uniform-in-the-volume-and-convergence-as-the-spacing-tends-to-zero}

\textbf{Part:} III, Ultraviolet construction (renormalization group).
\textbf{Lean:} chapter 30,
\href{https://github.com/yang-mills-group/OpenAI-YM-20261028/tree/3456e5e75cf8d4cd558955a5c7638fba2d68f464/YM/FullVolume}{\texttt{YM/FullVolume/}}.
\textbf{Size:} 18,973 lines in the main-theorem closure.

\subsubsection{1. What this unit proves}\label{what-this-unit-proves-14}

\textbf{Proposition III.12.1 (kernel bounds uniform in the volume).} Let
\(R\) be a run for a simple gauge group \(G\), with Wilson family \(C\).
Fix \(n\), labels \(v_i\in\mathrm{Lab}_{d_i}\), a finite list
\(\lambda\) of pairs (slot, axis), and \(r>0\). There are \(B\ge0,N_0\)
such that, for every \(N\ge N_0\), every \(p\) and every integer
configuration \(x\) with \(a_N\mathrm{dist}(x_i,x_j)\ge r\) for
\(i\ne j\),
\[\left|\Delta_\lambda\mathbb E_{N,p}\prod_{i=1}^nO_{N,p}(v_i;x_i)\right|\le B.\]
Each difference acts in its indicated slot; distance is periodic
\(\ell^\infty\) distance. Constants may depend on
\(R,n,(v_i),\lambda,r\), but not \(p\). \emph{Status:} theorem for every
run. \emph{Lean:}
\href{https://github.com/yang-mills-group/OpenAI-YM-20261028/blob/3456e5e75cf8d4cd558955a5c7638fba2d68f464/YM/FullVolume/Literal.lean\#L62}{\texttt{genuine}}.

\textbf{Proposition III.12.2 (derivative closure).} For a run with
simple \(G\), every sequence \(w_{N,p}\in\mathrm{Lab}_d\) bounded by
some \(K\ge0\) eventually in the mesh, and every finite axis list \(l\),
the trial \(\Delta_lO(w)\) is funded. \emph{Status:} theorem.
\emph{Lean:}
\href{https://github.com/yang-mills-group/OpenAI-YM-20261028/blob/3456e5e75cf8d4cd558955a5c7638fba2d68f464/YM/FullVolume/Literal.lean\#L24-L33}{\texttt{Run.financed}}.

\textbf{Proposition III.12.3 (closure lemma).} Fix \(d\). Assume a
consistent Wilson family \(C\) and:

\begin{enumerate}
\def\labelenumi{\arabic{enumi}.}
\tightlist
\item
  the centred budget;
\item
  the scale-selection data;
\item
  sound sections \(V_e\) and linear automorphisms \(W_e\) of abstract
  labels for every degree \(e\), satisfying the section, symmetry and
  exact local identities described below, with
  \href{https://github.com/yang-mills-group/OpenAI-YM-20261028/blob/3456e5e75cf8d4cd558955a5c7638fba2d68f464/YM/Push/RELAY_MARK/Portfolio.lean\#L51-L58}{\texttt{Portfolio}};
\item
  a presentation \(\pi\) and the full
  \href{https://github.com/yang-mills-group/OpenAI-YM-20261028/blob/3456e5e75cf8d4cd558955a5c7638fba2d68f464/YM/MergeProof/Canopy.lean\#L30-L32}{\texttt{section-comparison}}
  data, including quantitative comparison with \(V_e\), identification
  of \(\ker\pi\) with the dressed equation ideal, and
  \(\sup_{N,p}\|\pi_{N,p}\|<\infty\) in each degree;
\item
  degree-\(d\) detector rows with the placement and bounded-lift
  conditions and short-distance comparison of
  \href{https://github.com/yang-mills-group/OpenAI-YM-20261028/blob/3456e5e75cf8d4cd558955a5c7638fba2d68f464/YM/FullVolume/Maintain.lean\#L49-L52}{\texttt{VitalRank}}.
\end{enumerate}

If \(F\) is funded, every \(F_{N,p}(o)\) is continuous in the gauge
field, \(l\) is an axis list, and \(x\) is a degree-\(d\) coin with
\(O(x)-\Delta_lF\) owed at cost \(u_{N,p}\to0\) in the mesh, then
\(\Delta_lF\) is funded. \emph{Status:} conditional theorem; a run's
ties and short-distance estimates supply these hypotheses. \emph{Lean:}
\href{https://github.com/yang-mills-group/OpenAI-YM-20261028/blob/3456e5e75cf8d4cd558955a5c7638fba2d68f464/YM/FullVolume/Moneys.lean\#L69-L78}{\texttt{improved}}.

\textbf{Proposition III.12.4 (uniform rank detection).} Under hypotheses
1--5 of III.12.3, fix \(l\). For all large \(k\) and all sufficiently
small \(b>0\), there is \(Z\ge0\) such that, eventually in the mesh,
\[\|\pi_{N,p}x\|\le Z\|\mathcal P_{k,b,l;N,p}(\pi_{N,p}x)\|\quad(x\in\mathrm{Lab}_d).\]
The detector takes correlations with fixed partners and successively
averages the first slot in the directions of \(l\) over lengths between
\(b\) and \(2b\); repeated axes give repeated averages. \emph{Status:}
conditional theorem. \emph{Lean:}
\href{https://github.com/yang-mills-group/OpenAI-YM-20261028/blob/3456e5e75cf8d4cd558955a5c7638fba2d68f464/YM/FullVolume/Lit.lean\#L51-L91}{\texttt{trained}}.

\textbf{Proposition III.12.5 (negligible equation labels).} Assume
consistency of \(C\) and hypotheses 2--4 of III.12.3. Every degree-\(d\)
coin \(x\) with \(\pi_{N,p}x_{N,p}=0\) for all \((N,p)\) has \(O(x)\)
owed at cost \(a_N^{1/32}\). \emph{Status:} theorem. \emph{Lean:}
\href{https://github.com/yang-mills-group/OpenAI-YM-20261028/blob/3456e5e75cf8d4cd558955a5c7638fba2d68f464/YM/FullVolume/Silent.lean\#L105-L115}{\texttt{banked}}.

\textbf{Proposition III.12.6 (coupling power-series bounds).} Let \(R\)
be a run with a
\href{https://github.com/yang-mills-group/OpenAI-YM-20261028/blob/3456e5e75cf8d4cd558955a5c7638fba2d68f464/YM/FullVolume/Stage/Crests.lean\#L20-L27}{\texttt{Campaign}},
\href{https://github.com/yang-mills-group/OpenAI-YM-20261028/blob/3456e5e75cf8d4cd558955a5c7638fba2d68f464/YM/FullVolume/Gables/Vendee.lean\#L22-L27}{\texttt{Pillars}}
and bounds
\href{https://github.com/yang-mills-group/OpenAI-YM-20261028/blob/3456e5e75cf8d4cd558955a5c7638fba2d68f464/YM/FullVolume/Fares/Oz.lean\#L22-L28}{\texttt{Sky}}.
Let \(\bar\beta\) be the target level-0 coupling and \(w>0\) the fan
width. Fix \(p\), an \(n\)-tuple of labels \(V_0\) and \(r>0\). There
exist \(b,m\in\mathbb N\), \(m\ge1\), \(B\ge0\) and \(c\in(0,1]\) such
that the following holds for every \(m\le j<N\) and \(a\in[-1,1]\).

Take any family \(F(s)\), \(s\in[-1,1]\), of admissible campaign members
with common block sequence, sections and labels, whose level-0 inverse
coupling at \((N,p)\) is \(\bar\beta+was\). For every \(r\)-separated
configuration \(o\) in \(\Lambda_{j+1,p}\) embedded into
\(\Lambda_{N,p}\),
\[\mathbb E^{F(cs/2)}_{N,p}\prod_iO(V_{0,i};o_i)=\sum_{i\ge0}P_is^i,\qquad \sum_i|P_i|\le B a_{j+1}^{-b}\quad(s\in[-1,1]).\]
Expectation and representatives belong to the attached auxiliary
transfer family; its correlations agree with the successful trajectory
at positive cutoffs. Constants depend only on campaign data,
\(p,V_0,r\). \emph{Status:} conditional theorem; VIII.1 supplies the
campaign records. \emph{Lean:}
\href{https://github.com/yang-mills-group/OpenAI-YM-20261028/blob/3456e5e75cf8d4cd558955a5c7638fba2d68f464/YM/FullVolume/Fares/Scrap.lean\#L40-L51}{\texttt{Face.financed}}.

\textbf{Proposition III.12.7 (fixed-volume Cauchy property).} Fix simple
\(G\), one RG construction with stationary-path bounds, free
amplification, observable-transfer bounds and common reference sections,
and assume its uniform perturbation estimates. Fix \(p\). Let \(C_n\)
come from histories at cutoffs \(f(n)\to\infty\). Assume:

\begin{itemize}
\tightlist
\item
  Each history is a selected branch from level \(k\), agreeing with a
  fixed kinetic path; stationary kinetic paths converge with strict
  reserve in their admissible bounds. At every fixed \(i\ge k\), inverse
  couplings are Cauchy and activities have Cauchy representatives in the
  common weighted space.
\item
  Each degree used has reference level \(t(d)\ge k\) with
  \(d+3\le J(t(d))\); transfer data are based at its reference section
  and compose branch transfers.
\item
  For every sufficiently large base level \(r\), multiscale remainder
  bounds hold for all sufficiently large \(n\). Label allowances
  \(\zeta(d,r)\) are eventually positive for each degree used; occupancy
  allowances are at least \(1\) at every level.
\end{itemize}

For every finite label list and every configuration \(x\) with
\(x_i-x_j\notin L_p\mathbb Z^4\) for \(i\ne j\), lying exactly on some
common lattice \(x=a_rz\), \(z\) integer, the cutoff-\(f(n)\)
correlations at coordinatewise ceiling addresses of \(x\) are Cauchy.
\emph{Status:} conditional theorem. A run with campaign, pillars,
fixed-level convergence and uniform perturbation estimates gives
\(C_n=C\), \(f(n)=n\). \emph{Lean:}
\href{https://github.com/yang-mills-group/OpenAI-YM-20261028/blob/3456e5e75cf8d4cd558955a5c7638fba2d68f464/YM/FullVolume/Fields/Seat.lean\#L37-L45}{\texttt{Rising.cinema}}.

\textbf{Proposition III.12.8 (positional ultraviolet limit).} Let \(C\)
be a Wilson family and \(e\) a strictly increasing mesh sequence. Assume
the kernel-bound conclusion of III.12.1 and that ceiling-address
correlations are Cauchy along \(e\) whenever
\(x_{i\mu}-x_{j\mu}\notin L_p\mathbb Z\) for every \(i\ne j\) and axis
\(\mu\). For every \(p,n\) and label list there is a function \(f\) on
\((\mathbb R^4)^n\), smooth on
\(\mathrm{NoHit}_p=\{x:x_i-x_j\notin L_p\mathbb Z^4\ (i\ne j)\}\), with
these properties:

\begin{itemize}
\tightlist
\item
  For every smooth compactly supported \(t\) supported off collisions
  and strictly inside the centred period cell in every coordinate,
  smeared lattice sums converge along \(e\) to \(\int tf\).
\item
  Fix any finite list of coordinate directions (slot, axis), with
  integers \(q_-,q_+\ge0\), \(q_-+q_+>0\) in each direction, using the
  stencil \([F(x+q_+e_\mu)-F(x-q_-e_\mu)]/[(q_-+q_+)a_N]\). Write \(g\)
  for the corresponding iterated derivative of \(f\). Along \(e(s_i)\)
  for every \(s_i\to\infty\), these differences at lattice
  configurations whose physical positions converge to
  \(x\in\mathrm{NoHit}_p\) converge to \(g(x)\). Their multilinear
  interpolations converge locally uniformly there. For every test \(t\)
  above, \(tg\) is integrable and smeared differences converge to
  \(\int tg\).
\end{itemize}

\emph{Status:} conditional theorem; the run hypotheses following
III.12.7 give the full sequence \(e(s)=s\). \emph{Lean:}
\href{https://github.com/yang-mills-group/OpenAI-YM-20261028/blob/3456e5e75cf8d4cd558955a5c7638fba2d68f464/YM/FullVolume/Profiles.lean\#L42-L47}{\texttt{positioned}};
full-sequence instance
\href{https://github.com/yang-mills-group/OpenAI-YM-20261028/blob/3456e5e75cf8d4cd558955a5c7638fba2d68f464/YM/FullVolume/Fields/Filled.lean\#L45}{\texttt{Campaign.realUV}}.

\subsubsection{2. What it uses}\label{what-it-uses-14}

\begin{itemize}
\tightlist
\item
  \textbf{II.1--II.2, III.14:} labels, representatives, consistency and
  the
  \href{https://github.com/yang-mills-group/OpenAI-YM-20261028/blob/3456e5e75cf8d4cd558955a5c7638fba2d68f464/YM/FullVolume/Cross.lean\#L30-L48}{\texttt{Run}}
  record, whose ties supply the inputs of III.12.3.
\item
  \textbf{III.2, III.11:} numerical parameters, scale selection, section
  comparisons, regularity, local matching and short-distance estimates;
  the centred budget
  \href{https://github.com/yang-mills-group/OpenAI-YM-20261028/blob/3456e5e75cf8d4cd558955a5c7638fba2d68f464/YM/Extraction/Center.lean\#L38-L40}{\texttt{Wcmain}}
  states: for every degree \(d\), \(R>0\) and \(\varepsilon>0\), there
  is \(A\ge0\) such that for every \(m\ge t(d)\), eventually in the
  mesh, every \(\|v\|\le R\) and site \(o\) give an observable localized
  in the cell centred at \(o\), with root at most
  \(A a_m^{-d-\varepsilon}\).
\item
  \textbf{IV.1:}
  \href{https://github.com/yang-mills-group/OpenAI-YM-20261028/blob/3456e5e75cf8d4cd558955a5c7638fba2d68f464/YM/Cells/Certificate.lean\#L14-L19}{\texttt{Seam}}
  gives chessboard and root triangle inequalities; positional Lipschitz
  bounds and dense-set convergence. (The offset comparison comes from
  II.1.)
\item
  \textbf{III.6, III.8--III.10:} histories, selected branches,
  stationary paths and transfer estimates. \textbf{VIII.1:}
  \href{https://github.com/yang-mills-group/OpenAI-YM-20261028/blob/3456e5e75cf8d4cd558955a5c7638fba2d68f464/YM/FullVolume/Glide/Vac.lean\#L31-L63}{\texttt{Winter}}
  supplies fixed-level convergence, kinetic reserve and free
  amplification; assembly supplies uniform perturbation inputs.
\item
  \textbf{Standard results:} Wilson reflection positivity
  (Osterwalder--Seiler); chessboard estimates and seminorm
  (Fröhlich--Israel--Lieb--Simon, Fröhlich--Lieb, Shlosman, Biskup);
  summation by parts, finite-dimensional injectivity stability,
  Arzelà--Ascoli, Taylor's theorem, absolutely convergent power series
  and dominated convergence. Volume-uniform control of effective
  actions: Balaban (1984--1989).
\end{itemize}

The section hypotheses require the prescribed polynomial realization and
finite regularity bounds; dressings commute with signed lattice
symmetries and conjugation. Dressed equation labels have zero Wilson
expectation against every bounded measurable spectator depending only on
links beyond a fixed lattice distance, for every nonzero inverse
coupling and sufficiently large torus; spectators need not be gauge
invariant.

Full section-comparison data include terminal-coordinate automorphisms
\(A_{d;N,p}\), auxiliary columns and sub-power \(E_d(N)\). Eventually,
uniformly in \(p\), \(A_{d;N,p}\) is the symmetry-averaged column map
and its inverse increases coefficient size by at most \(2\). For each
column and its signed or conjugated versions, the centred root of the
auxiliary-minus-chosen section evaluation is at most
\(C_d(m,k)E_d(N)a_N^{1/8}\) for every admissible \(k<N\) and \(m\le N\);
scale selection chooses \(k\) for each large \(m\). The kernel is the
image of the Maxwell equation ideal under dressing, threshold rescaling,
terminal map and \(A_{d;N,p}\). Projection symmetries and finite-cutoff
equation and difference comparisons are also required.

Detector rows use test configurations with distinct projections on every
axis in a box of side \(a_{\mathrm{width}(k)}\), and lifted partner
labels bounded uniformly in \(N,p\) and finite indices for fixed \(k\).
Their short-distance cumulant comparisons approach the prescribed free
Maxwell tests.

\subsubsection{3. Strategy}\label{strategy-14}

Normalized differences introduce \(a_N^{-1}\), beyond the centred budget
for bounded labels. Local matching replaces each differenced observable
by a primary with sub-power inverse terminal coordinates and vanishing
root error. Averaged partner correlations bound its projected physical
label through an injective free Maxwell detector; equation labels
contribute a vanishing root. This funds differences, and chessboard
estimates bound their separated products.

At fixed volume, compare full correlations with scalar-normalized
correlations at an intermediate level. Fixed-level convergence makes the
latter Cauchy, while an exponentially small tail in the budget radius
tends to zero. Positional bounds extend convergence from common-lattice
configurations; Taylor estimates identify derivatives. For coupling
dependence, polynomial coefficient bounds and smaller parameter
intervals turn level differences into summable coefficient errors.

\subsubsection{4. Steps}\label{steps-14}

\begin{enumerate}
\def\labelenumi{\arabic{enumi}.}
\tightlist
\item
  \textbf{Match differences.} For a run with simple \(G\) and an
  eventually bounded degree-\(d\) label sequence \(w\), normalized
  differences preserve regularity with sub-power bounds after increasing
  degree. Local matching gives a degree-\((d+|l|)\) coin \(x\) with
  \(O(x)-\Delta_lO(w)\) owed at \(a_N^{1/16}\).
  \href{https://github.com/yang-mills-group/OpenAI-YM-20261028/blob/3456e5e75cf8d4cd558955a5c7638fba2d68f464/YM/FullVolume/Paid.lean\#L47-L49}{\texttt{prepared\_rows}}.
\item
  \textbf{Remove equation labels.} Under III.12.5's hypotheses,
  auxiliary equation sections have zero root: their exact
  distant-spectator identity annihilates the reflected product. Section
  comparison leaves cost \(a_N^{1/32}\).
  \href{https://github.com/yang-mills-group/OpenAI-YM-20261028/blob/3456e5e75cf8d4cd558955a5c7638fba2d68f464/YM/FullVolume/Silent.lean\#L105-L115}{\texttt{banked}}.
\item
  \textbf{Detect projected labels.} Under III.12.3's hypotheses,
  averaging differences telescopes and bounds tests by \((2/b)^{|l|}\)
  times undifferenced correlations. Short-distance comparison approaches
  the injective free Maxwell matrix; finite-dimensional stability and
  the rank bound give III.12.4.
  \href{https://github.com/yang-mills-group/OpenAI-YM-20261028/blob/3456e5e75cf8d4cd558955a5c7638fba2d68f464/YM/FullVolume/Lit.lean\#L51-L91}{\texttt{trained}}.
\item
  \textbf{Close the budget.} Under III.12.3's hypotheses, set
  \(X=\pi x\). Apply negligible-equation control to \((x-X)/(1+\|X\|)\),
  obtaining \(\|X\|\le Zc+\tfrac12(1+\|X\|)\). Thus \(X\) is bounded;
  the centred budget funds \(O(X)\), and vanishing errors fund the
  desired difference.
  \href{https://github.com/yang-mills-group/OpenAI-YM-20261028/blob/3456e5e75cf8d4cd558955a5c7638fba2d68f464/YM/FullVolume/Moneys.lean\#L23-L63}{\texttt{cap}},
  \href{https://github.com/yang-mills-group/OpenAI-YM-20261028/blob/3456e5e75cf8d4cd558955a5c7638fba2d68f464/YM/FullVolume/Moneys.lean\#L69-L78}{\texttt{improved}}.
\item
  \textbf{Bound kernels.} Once all fixed-label differences are funded,
  place separated insertions in distinct cells with interior margins;
  offset comparison and the chessboard inequality bound products. Slot
  differences act on their insertions, giving III.12.1.
  \href{https://github.com/yang-mills-group/OpenAI-YM-20261028/blob/3456e5e75cf8d4cd558955a5c7638fba2d68f464/YM/FullVolume/Literal.lean\#L42-L61}{\texttt{funding}}.
\item
  \textbf{Control coupling series.} Under III.12.6's campaign
  hypotheses, composed transfer bounds give coefficient sums polynomial
  in \(a_{j+1}^{-1}\). With fixed-level convergence and perturbation
  inputs, interval reduction converts real \(O(a_j)\) level differences
  to summable \(C\sqrt{a_j}\) coefficient bounds. For fixed \(p\),
  labels, compact \(K\subset\mathrm{NoHit}_p\) and \(q\in\mathbb N\),
  this gives a positive parameter width independent of cutoff and a
  neighbourhood of \(K\) with limits \(C^q\) in position and analytic in
  the complex parameter; kernels with spatial stencils of length at most
  \(q\), their interpolations and admissible smeared sums have analytic
  continuations whose parameter derivatives of every fixed order
  converging uniformly on compact subdiscs and compact spatial subsets.
  \href{https://github.com/yang-mills-group/OpenAI-YM-20261028/blob/3456e5e75cf8d4cd558955a5c7638fba2d68f464/YM/FullVolume/Finale/Atlas.lean\#L94-L104}{\texttt{Fan.standing}}.
\item
  \textbf{Take the mesh limit.} Under III.12.7's hypotheses, fixed-level
  effective insertions and normalized densities are uniformly Cauchy;
  Haar integration and positive partition-function bounds give Cauchy
  correlations. For large fixed \(j\), the full correlation differs by
  at most \(2F_0e^{-B_j/64}\), uniformly in sufficiently large cutoffs,
  with \(F_0\) independent of \(j\) and cutoff and \(B_j\to\infty\).
  \href{https://github.com/yang-mills-group/OpenAI-YM-20261028/blob/3456e5e75cf8d4cd558955a5c7638fba2d68f464/YM/FullVolume/Fields/Peat.lean\#L41-L64}{\texttt{Sever.echoes}},
  \href{https://github.com/yang-mills-group/OpenAI-YM-20261028/blob/3456e5e75cf8d4cd558955a5c7638fba2d68f464/YM/FullVolume/Fields/Timber.lean\#L73-L103}{\texttt{Rising.arrives}}.
\item
  \textbf{Recover spatial derivatives.} Under III.12.8's hypotheses,
  equicontinuity extends dense-set convergence. Taylor comparison with
  differences over a fixed physical interval identifies successive
  derivative limits; stencil averaging, multilinear interpolation and
  dominated convergence give the remaining conclusions.
  \href{https://github.com/yang-mills-group/OpenAI-YM-20261028/blob/3456e5e75cf8d4cd558955a5c7638fba2d68f464/YM/FullVolume/Profiles.lean\#L42-L47}{\texttt{positioned}}.
\end{enumerate}

\subsubsection{5. Where the difficulty
lies}\label{where-the-difficulty-lies-14}

Matching controls inverse terminal coordinates; rank detection must
bound physical labels while section comparison controls the projection
kernel. Bounds are volume-uniform, but convergence fixes \(p\) first.
Every derivative order is fixed before choosing the cutoff threshold;
analytic widths may depend on volume, compact spatial set and derivative
order. The additive action scalar need not converge. Collision extension
belongs to VII.4.

\subsubsection{6. Notation}\label{notation-14}

\begin{itemize}
\tightlist
\item
  \(a_N=(\ell_0\cdots\ell_{N-1})^{-1}\); \(\Lambda_{N,p}\) has physical
  period \(L_p\). \(O(v;o)\) represents \(v\in\mathrm{Lab}_d\);
  \(\mathbb E_{N,p}\) is Wilson expectation.
\item
  Eventually in the mesh: all \(N\ge N_0\), all \(p\); \(N_0\) may
  depend on a fixed coarse level. Sub-power growth means
  \(\|w_{N,p}\|\le a_N^{-\varepsilon}\) eventually for every
  \(\varepsilon>0\). A
  \href{https://github.com/yang-mills-group/OpenAI-YM-20261028/blob/3456e5e75cf8d4cd558955a5c7638fba2d68f464/YM/FullVolume/Sluggish.lean\#L146-L147}{\texttt{Coin}}
  is a label sequence \(x\) with \(\|\tau_{d;N,p}^{-1}x_{N,p}\|\) of
  sub-power growth, where \(\tau\) is the terminal coordinate
  isomorphism.
\item
  \(\Delta_\mu F(o)=[F(o+e_\mu)-F(o)]/a_N\); lists indicate iteration. A
  trial is a site-indexed family of gauge-field functions; it is local
  if eventually translation covariant, gauge invariant, measurable,
  bounded and supported on one fixed finite link stencil.
\item
  \(R_m(H)\) is the absolute reflected-product expectation to power
  \(1/(L_p/a_m)^4\) on centred cells of physical side \(a_m\),
  conjugating odd reflections. A local trial is \textbf{funded} if, for
  all large \(m\), eventually in the mesh, its value at \(0\) is
  localized in the central cell and has root at most \(u(m)\)
  independent of \(N,p\); \textbf{owed} replaces \(u(m)\) by a cost
  \(v(N,p)\) independent of \(m\).
\item
  \(\pi_{N,p}\) is an idempotent label map; its dressed-equation kernel
  is an additional hypothesis. Ceiling addresses are
  \(\alpha_N(x)_{i\mu}=\lceil x_{i\mu}/a_N\rceil\).
  \(B_j=\lceil(\log\hat\beta_j)^2\rceil\) is the budget radius.
\end{itemize}

\subsubsection{7. What remains
unexplained}\label{what-remains-unexplained-14}

The local matching estimate, construction of auxiliary equation sections
and relation to lattice Schwinger--Dyson identities remain imported or
unverified here. One-step transfer capacity, intermediate weighted
analytic estimates, spatial derivative identification for analytic
limits, and multiscale remainder and omission bounds need their source
analyses.

Detailed notes: the full account of this unit, with proof sketches for
every step, is on the
\href{https://github.com/scottnarmstrong/yang-mills-informalization/blob/main/outline/notes/III-12.md}{detailed
notes page}.

\subsection{III.13. Extension to all labels: one common Wilson
family}\label{iii.13.-extension-to-all-labels-one-common-wilson-family}

\textbf{Part:} III, Ultraviolet construction (renormalization group).
\textbf{Lean:}
\href{https://github.com/yang-mills-group/OpenAI-YM-20261028/blob/3456e5e75cf8d4cd558955a5c7638fba2d68f464/YM/SharpGram/Common.lean}{\texttt{YM/SharpGram/Common.lean}},
\href{https://github.com/yang-mills-group/OpenAI-YM-20261028/blob/3456e5e75cf8d4cd558955a5c7638fba2d68f464/YM/SharpGram/Measured.lean}{\texttt{YM/SharpGram/Measured.lean}},
\href{https://github.com/yang-mills-group/OpenAI-YM-20261028/tree/3456e5e75cf8d4cd558955a5c7638fba2d68f464/YM/SharpGram/Assembly}{\texttt{YM/SharpGram/Assembly/}},
\href{https://github.com/yang-mills-group/OpenAI-YM-20261028/tree/3456e5e75cf8d4cd558955a5c7638fba2d68f464/YM/UVExtension}{\texttt{YM/UVExtension/}}.
\textbf{Size:} about 22,100 lines in the unit inventory.

\subsubsection{1. What this unit proves}\label{what-this-unit-proves-15}

Throughout, \(G\) is a compact matrix group with simple nonabelian real
Lie algebra and \(b>0\) is the Gaussian block-spin penalty. Levels run
from the coarsest, \(0\), to the cutoff, \(N\); \(p\) indexes periods.
Notation is collected in §6.

\textbf{Proposition III.13.1 (one ultraviolet and logarithmic family).}
There exists a Wilson family \(C\), with penalty \(b\), satisfying all
the following.

\begin{enumerate}
\def\labelenumi{\arabic{enumi}.}
\tightlist
\item
  It is calibrated: one full choice of constants and a directed history
  generate its history, with the per-step estimates of III.9--III.12.
\item
  It is consistent: the trajectory is selected, local sections are
  sound, and, for every degree \(d\), \(d+3\le J(\mathrm{ref}(d))\),
  \(d+1\le\mathrm{order}_m\) and \(d\le\mathrm{power}_m\) for
  \(m\ge\mathrm{ref}(d)\). Every degree-\(d\) stencil point satisfies
  \(20(|x_i|+1)\le\mathrm{side}(N,p)\) in each coordinate whenever
  \(N\ge\mathrm{ref}(d)\). Placed basis sections are regular on their
  footprints, with the prescribed derivative bounds of size
  \(\mathrm{columns}_N\theta_N^q\) for factor degree \(q\); transfer is
  calibrated from \(\mathrm{ref}(d)\) onwards.
\item
  For every \(d\) and \(R,\varepsilon>0\), there is \(A\ge0\) such that,
  for every \(m\ge\mathrm{ref}(d)\) and all sufficiently large \(N\),
  every \(p\) and every \(v\in\mathcal L_d\) with \(\|v\|\le R\), the
  observable at every site is supported in its centred level-\(m\)
  reflection cell and its chessboard root is at most
  \(Aa_m^{-d-\varepsilon}\). Here \(N\ge\mathrm{ref}(d)\) is required.
\item
  It is the centre of the analytic shooting family of III.13.7.
\item
  The OPE entrance conditions hold: for every \(d,D,h\) with
  \(D+h=d+4\), there exist \(t>1\) and \(j(i)\to\infty\) with
  \(\sigma_{j(i)}/\sigma_i\to t\) for which the III.11 shell bounds hold
  for every choice of labels and extraction order; cover sizes satisfy
  \(\mathrm{cs}_{i+1}^2\le2\mathrm{cs}_i^2\) and for every
  \(\varepsilon>0\),
  \((L_i\mathrm{cs}_i^2)^4(\text{vacuum cost})\le\varepsilon\sigma_{i+1}\)
  eventually; positive radii satisfy \(\log\rho_i/\sigma_i\to-1\).
\item
  There exist coefficient-map bounds \(W\), allowances
  \(r:\mathbb N\to\mathbb N\) and transfer data giving the logarithmic
  OPE package: III.11 transfer/remainder estimates, triangular label
  steps, auxiliary comparisons, separation bounds, product sup bounds
  and VII.4 arity bounds. Also
  \(L_m,\mathrm{columns}_m,\mathrm{columns}_{m+1}\le\hat\beta_m\),
  \(r\to\infty\), and the support exponent is \(100\).
\item
  For these data and every \(d,N,p,m\) with \(\mathrm{ref}(d)\le m<N\),
  there is a linear isomorphism \(T\) with
  \[\|x_m(v)-T(V_{L_m}x_{m+1}(v))\|_1\le\mathrm{barunit}(L_m)\|\mathrm{core}(x_{m+1}(v))\|_1\quad(v\in\mathcal L_d).\]
  On a basis vector of engineering grade \(e\) and factor degree \(q\),
  the only output of grade \(\le e\) is that vector times
  \((\theta_m/\theta_{m+1})^q\); no lower factor degree occurs. The map
  preserves the unit, commutes with core projection and lattice
  symmetries, and \(\|T\|_1,\|T^{-1}\|_1\le K(J(m+1))\), where \(K\)
  comes from \(W\).
\item
  Coverage holds: some \(g(k)\to\infty\), \(g(k)\le k\), satisfies
  \(a_{g(k)}\le\hat\beta_ka_k\) eventually, with three depth-\(d\) probe
  panels poised at level \(g(k)\) for every \(d\) and all large \(k\).
\item
  Along every extraction order satisfying IV.1's cell hypotheses, grid
  approximation holds at arbitrarily fine fixed levels for bounded
  degrees and III.11's allowed reflection pairs. Along each such order,
  complete weak comparison systems exist for every finite demand and
  degree horizon, positive radius bound and real-spin structure on the
  smooth free comparison table.
\item
  Every extraction order \(f\) has a refinement \(g\), extraction base
  tables \(A,(L_k),T\), a bin and derivative plan along \(f\circ g\),
  and a seed for \((C,f\circ g,T)\) with Maxwell ranks for \(T\).
\end{enumerate}

\emph{Status:} theorem. \emph{Lean:}
\href{https://github.com/yang-mills-group/OpenAI-YM-20261028/blob/3456e5e75cf8d4cd558955a5c7638fba2d68f464/YM/SharpGram/Assembly/Flowers.lean\#L75-L130}{\texttt{uv\_log}}.

\textbf{Proposition III.13.2 (coordinated family).} There exists one
\(C\) of penalty \(b\), consistent, calibrated and floral, with the
chessboard bound above, III.13.5's fixed-volume limits, the logarithmic
package, calibrated label steps, coverage and complete weak comparison
systems along every extraction order satisfying the cell hypotheses. For
every extraction order \(F\), there exist a refinement \(g\) and table
\(T\) giving the ordered limit along the full mesh sequence and the
periods of \(F\circ g\), VII.4's OPE services, and a II.2 construction
record for \((C,F\circ g,T)\).

\emph{Status:} theorem; \(C\) precedes \(F\). \emph{Lean:}
\href{https://github.com/yang-mills-group/OpenAI-YM-20261028/blob/3456e5e75cf8d4cd558955a5c7638fba2d68f464/YM/SharpGram/Common.lean\#L16-L37}{\texttt{coordination}}.

\textbf{Proposition III.13.3 (all phases).} For every extraction order
\(f\), there exist \(C\), a refinement \(g\), a table \(T\), a
construction record \(F\) for \((C,f\circ g,T)\), and a V.3
reconstruction record \(\mathcal A\) for every phase of \(F\),
satisfying III.13.2 and completed reconstruction: Wightman realization,
surjective comparison maps onto the Osterwalder--Schrader Hilbert
spaces, unique vacuum, phase-uniform light-cone and heat-kernel bounds,
trivial commutant, continued stress tensor, continuous phase dependence
of boundary correlations, and barycentre identities. The derivative
sector has a prescription \(P\) satisfying VII.4's estimates for every
fixed finite cluster tree and separated spectators, with arbitrary
positive power loss.

\emph{Status:} theorem; \(C\) follows \(f\) in the statement, but the
proof uses III.13.2's common choice. \emph{Lean:}
\href{https://github.com/yang-mills-group/OpenAI-YM-20261028/blob/3456e5e75cf8d4cd558955a5c7638fba2d68f464/YM/SharpGram/Common.lean\#L41-L60}{\texttt{common}}.

\textbf{Proposition III.13.4 (moving-frame Gram limits).} For every
extraction order \(f\), there exist \(C,g,T,F,\mathcal A,P\) satisfying
III.13.3, with coverage at the cutoff of \(F\). For every ultrafilter
\(\mathcal U\) finer than the filter of \(f\circ g\) (period large, then
mesh large) and every depth \(D\):

\begin{itemize}
\tightlist
\item
  Depth-\(D\) portfolios exist: levels \(k'(k)\to\infty\), \(k'\le k\),
  eventually satisfy \(a_k\le a_{k'}\le\hat\beta_ka_k\), with three
  poised panels for all sufficiently large \(k\), and label-matrix
  limits along \(\mathcal U\) for every sufficiently large \(k\).
\item
  For every portfolio \(X\), there are isomorphisms \(Y_k(d)\) from the
  free Maxwell degree-\(d\) field space onto that of \(T\), for all
  \(d\le D\) and \(k\). Eventually they respect grading through \(D\);
  for each \(d\le D\), their Gram forms are bounded and coercive
  uniformly for all sufficiently large \(k\).
\item
  For each \(d\le D\), some \(q_d\in\mathbb N\) gives, for every
  sufficiently large \(k\) and every field \(x\), representatives
  \(v_{N,p}\to v\) along \(\mathcal U\), projecting to \(x\), with
  eventual source cost at most
  \(\hat\beta_k^{q_d}a_k^d\|Y_k(d)^{-1}x\|\). Conversely, some
  \(K_d\ge0\) gives \(\|Y_k(d)^{-1}x\|\le K_da_k^{-d}z\) for every
  convergent representative family of cost at most \(z\) along any
  nontrivial filter finer than \(\mathcal U\), for all sufficiently
  large \(k\).
\item
  For each \(d\le D\), the Gram form of \(X\) in basis \(Y_k(d)\)
  converges, as \(k\to\infty\), to the free Maxwell Gram form.
\end{itemize}

\emph{Status:} theorem. \emph{Lean:}
\href{https://github.com/yang-mills-group/OpenAI-YM-20261028/blob/3456e5e75cf8d4cd558955a5c7638fba2d68f464/YM/SharpGram/Measured.lean\#L11-L36}{\texttt{habitat}}.

\textbf{Proposition III.13.5 (fixed-volume derivative limits).} For
\(C\) as in III.13.1 and every \(d,k\in\mathbb N\), all sufficiently
large \(p\) (depending on \(d,k\)), every \(n\) and labels \(v_i\) of
degrees \(\le d\), there exists \(g:(\mathbb R^4)^n\to\mathbb C\),
continuous on \(\Omega_p^n\), such that the following hold along the
whole sequence \(N\to\infty\).

\begin{itemize}
\tightlist
\item
  Periodic correlations, extended as step functions by coordinatewise
  ceiling projection, converge to \(g\) uniformly on separated compact
  sets.
\item
  Normalized forward differences of every coordinate word \(u\) of
  length \(\le k\) converge there to \(\partial_ug\). Multilinear
  interpolation gives the same limits for fixed words of length
  \(\le k\) of normalized differences with positive integer steps and
  fixed integer stencil offsets, including backward and centred
  differences.
\item
  \(\partial_a\partial_ug=\partial_{(a,u)}g\) for \(|u|<k\) on
  \(\Omega_p^n\).
\item
  For every admissible test \(t\), \(tg\) is integrable and
  \[a_N^{4n}\sum_r t(r)\Big\langle\prod_iO_{v_i}^{N,p}(r_i)\Big\rangle_{N,p}\longrightarrow\int tg,\]
  summing over centred torus representatives.
\end{itemize}

\emph{Status:} theorem. \emph{Lean:}
\href{https://github.com/yang-mills-group/OpenAI-YM-20261028/blob/3456e5e75cf8d4cd558955a5c7638fba2d68f464/YM/UVExtension/Poppy/Fountain.lean\#L96-L108}{\texttt{Poppy.blossoms}},
\href{https://github.com/yang-mills-group/OpenAI-YM-20261028/blob/3456e5e75cf8d4cd558955a5c7638fba2d68f464/YM/UVExtension/Poppy/Fountain.lean\#L16-L21}{\texttt{Poppy.Margins}}.

\textbf{Proposition III.13.6 (separated lattice points).} Fix \(p\) and
a member \(C^{\mathbf s}\) of III.13.7 whose shifts satisfy
\(s_{N,p}\to s\), \(|s|<r_*\). Assume its chessboard bound, slow growth,
label-transfer and remainder bounds. For every \(J\), every finite label
list \(u_i\) with \(\mathrm{ref}(\deg u_i)\le J\), and pairwise distinct
level-\(J\) torus sites \(o_i\), the correlations
\(\langle\prod_iO_{u_i}^{N,p}(o_i)\rangle_{N,p}\) converge as
\(N\to\infty\). III.13.7 supplies these bounds for every shift
assignment; the centre has zero shifts.

\emph{Status:} theorem. \emph{Lean:}
\href{https://github.com/yang-mills-group/OpenAI-YM-20261028/blob/3456e5e75cf8d4cd558955a5c7638fba2d68f464/YM/UVExtension/Hearth/Table.lean\#L87-L112}{\texttt{Hearth.food}}.

\textbf{Proposition III.13.7 (analytic shooting family).} For \(C\) of
III.13.1 there exist a full choice of per-step constants, a shooting
structure with an arc, a label-transfer bundle and reference levels such
that:

\begin{enumerate}
\def\labelenumi{\arabic{enumi}.}
\tightlist
\item
  Some \(r_*>0\), depending only on those constants, works for every
  \(N,p\): each real \(|s|<r_*\) determines exactly one Wilson coupling
  whose fixed-recursion trajectory stays in the prescribed scalar bands
  and reaches \(\bar\beta+s\) at level \(0\). Assignments
  \(|s_{N,p}|<r_*\) define \(C^{\mathbf s}\) with the bundle's labels,
  sections and reference levels, and \(C=C^{\mathbf0}\).
\item
  For \(m\le N\), \(\beta_m(N,p;s)\) and normalized activity readings
  extend holomorphically to \(|s|<r_*\), bounded by \(3\hat\beta_m\) and
  \(\mathrm{soft}=\mathrm{rim}/8\), respectively. For fixed \(m,p\),
  they converge locally uniformly with every shift derivative as
  \(N\to\infty\).
\item
  The centre has the calibration package with constant allowance, priced
  sup bounds and extraction base data. Every \(C^{\mathbf s}\) has
  label-transfer, remainder and chessboard bounds.
\item
  Some common disc \(|z|<a\) lies in the shooting domain and carries
  holomorphic coefficient and transfer data, including inverse terminal
  label maps, for all \(d,N,p\) with \(\mathrm{ref}(d)\le N\).
\item
  Fix \(d,m\). For all large \(p\), every \(n\) and labels of degree
  \(\le d\), there exist \(b'>0\), holomorphic lattice correlations
  \(W_N(z,x)\) on \(|z|<b'\) and limits \(f(z,x)\). At real \(s\),
  \(W_N\) equals the correlation whenever \(s_{N,p}=s\) and all
  reference levels fit. For every complex \(z\), \(f(z,\cdot)\) has
  continuous iterated coordinate derivatives through order \(m\) on
  \(\Omega_p^n\); step differences and interpolated differences converge
  to them with every \(z\) derivative, uniformly on products of
  separated spatial compacts and disc compacts. Real-shift limits agree
  with constant-shift fixed-volume limits. For admissible \(t\),
  \(\int tf(z,\cdot)\) is holomorphic and smeared lattice correlations
  converge to it locally uniformly with every \(z\) derivative.
\end{enumerate}

\emph{Status:} theorem, included in III.13.1. \emph{Lean:}
\href{https://github.com/yang-mills-group/OpenAI-YM-20261028/blob/3456e5e75cf8d4cd558955a5c7638fba2d68f464/YM/UVExtension/Oasis/Flora.lean\#L55-L71}{\texttt{Oasis.Floral}}.

\subsubsection{2. What it uses}\label{what-it-uses-15}

\begin{itemize}
\tightlist
\item
  I.1; II.1--II.2: Wilson measures, labels, observables, extraction
  orders and construction records.
\item
  III.1--III.2; III.4--III.12: kernels, order of constants, chessboard
  and per-step estimates, selected trajectories and histories; III.7
  supplies the one-loop annulus coefficient.
\item
  III.15--III.16: label costs and analytic tools; IV.1: extraction, cell
  hypotheses and Arzelà--Ascoli.
\item
  V.3; VII.1, VII.3--VII.4; VIII.1: phase reconstruction, coverage, Gram
  comparison, OPE services and construction records.
\item
  Balaban block-spin renormalization; Banach contraction; Vitali--Porter
  and the identity principle; Arzelà--Ascoli; dominated convergence;
  reflection-positive chessboard estimates; ultrafilters.
\end{itemize}

\subsubsection{3. Strategy}\label{strategy-15}

Choose constants in their required order, raising thresholds while
preserving earlier estimates. Increasing Taylor orders and powers
accommodate each label degree from its reference level onwards. Shooting
gives trajectories with prescribed coarse coupling; uniform contraction
and decay of dependence on remote ultraviolet levels give full-sequence
limits at each fixed level. Uniform holomorphic bounds then permit
Vitali's theorem.

Approximate each renormalized observable by a conditional average of a
coarse observable with convergent coefficients. Exponential improvement
beats subexponential coefficient costs, and chessboard estimates control
products. Dense-grid convergence and large-period slicing yield
derivative limits. Extraction, OPE and reconstruction results then apply
to the same family.

\subsubsection{4. Steps}\label{steps-15}

\begin{enumerate}
\def\labelenumi{\arabic{enumi}.}
\tightlist
\item
  \textbf{Choose constants.} For simple \(G\), combine the per-step
  estimates, then raise thresholds for shooting, label transfer,
  logarithmic OPE and weak comparison. Instantiate first stage \(4\),
  Taylor order \(6+i\), power \(2+i\), support exponent \(100\),
  allowances \(1\) and \(j+1\), and target rate \(-7/16\); the final
  wait also supplies coverage.
  \href{https://github.com/yang-mills-group/OpenAI-YM-20261028/blob/3456e5e75cf8d4cd558955a5c7638fba2d68f464/YM/SharpGram/Assembly/Flowers.lean\#L75-L130}{\texttt{uv\_log}}.
\item
  \textbf{Take fixed-level limits.} For the constructed arc, contraction
  has rate \(1/2\) uniformly in cutoff; continuity and kinematic
  convergence give convergence of each iterate. Dependence on levels
  above \(m+v\) decays as \(\mathrm{soft}^v\), giving real-shift limits,
  then Vitali gives every complex-shift derivative.
  \href{https://github.com/yang-mills-group/OpenAI-YM-20261028/blob/3456e5e75cf8d4cd558955a5c7638fba2d68f464/YM/UVExtension/Camp/Pear.lean\#L70-L71}{\texttt{Camp.artistry}},
  \href{https://github.com/yang-mills-group/OpenAI-YM-20261028/blob/3456e5e75cf8d4cd558955a5c7638fba2d68f464/YM/UVExtension/Camp/Veto.lean\#L64-L69}{\texttt{Camp.viva}}
  \emph{(proofs read; contraction and decay estimates used as stated)}.
\item
  \textbf{Pass to correlations.} For convergent real shifts and the
  bundle bounds, effective densities converge uniformly and partition
  integrals stay positively bounded below. The RG identity and coarse
  approximation reduce separated products to convergent effective
  expectations; chessboard errors tend to zero, proving the Cauchy
  property.
  \href{https://github.com/yang-mills-group/OpenAI-YM-20261028/blob/3456e5e75cf8d4cd558955a5c7638fba2d68f464/YM/UVExtension/Hearth/Table.lean\#L87-L112}{\texttt{Hearth.food}}.
\item
  \textbf{Obtain derivatives and analyticity.} Large-period slicing
  bounds lattice differences through one extra order; dense-grid
  convergence and Arzelà--Ascoli give spatial derivatives and smearing.
  Holomorphic transfer bounds and Vitali give the analytic correlation
  limits.
  \href{https://github.com/yang-mills-group/OpenAI-YM-20261028/blob/3456e5e75cf8d4cd558955a5c7638fba2d68f464/YM/UVExtension/Poppy/Fountain.lean\#L96-L108}{\texttt{Poppy.blossoms}},
  \href{https://github.com/yang-mills-group/OpenAI-YM-20261028/blob/3456e5e75cf8d4cd558955a5c7638fba2d68f464/YM/UVExtension/Beach/Under.lean\#L43-L63}{\texttt{Beach.sparksP}}
  \emph{(proofs read through stated transfer estimates)}.
\item
  \textbf{Assemble common outputs.} Uniqueness identifies full-mesh
  smeared limits with extracted limits. Apply OPE services, every-phase
  reconstruction and Gram comparison; coverage supplies portfolios along
  every finer ultrafilter.
  \href{https://github.com/yang-mills-group/OpenAI-YM-20261028/blob/3456e5e75cf8d4cd558955a5c7638fba2d68f464/YM/SharpGram/Common.lean\#L16-L37}{\texttt{coordination}},
  \href{https://github.com/yang-mills-group/OpenAI-YM-20261028/blob/3456e5e75cf8d4cd558955a5c7638fba2d68f464/YM/SharpGram/Common.lean\#L41-L60}{\texttt{common}},
  \href{https://github.com/yang-mills-group/OpenAI-YM-20261028/blob/3456e5e75cf8d4cd558955a5c7638fba2d68f464/YM/SharpGram/Measured.lean\#L11-L36}{\texttt{habitat}}
  \emph{(report proofs read)}.
\end{enumerate}

\subsubsection{5. Where the difficulty
lies}\label{where-the-difficulty-lies-15}

Uniform control over arbitrarily many ultraviolet levels and the balance
between coarse-approximation gains and coefficient costs are the main
analytic difficulties. Derivative limits require large periods; plain
smeared limits hold at every fixed period. The coupling/activity disc is
uniform in cutoff and period, whereas the correlation disc may depend on
labels and period.

\subsubsection{6. Notation}\label{notation-15}

\(\mathcal L_d\) is the degree-\(d\) label space; \(\mathrm{ref}(d)\)
its reference level. \(a_0=1\), \(a_{m+1}=a_m/L_m\),
\(\sigma_m=\log(1/a_m)\); \(\ell_p\) is the torus side.
\(\Omega_p^n\subset(\mathbb R^4)^n\) consists of configurations distinct
modulo \(\ell_p\mathbb Z^4\); admissible tests are complex-valued,
smooth, compactly supported on distinct configurations, with closed
support inside \((-\ell_p/2,\ell_p/2)^{4n}\). \(O_v^{N,p}\) denotes the
renormalized observable and brackets Wilson expectation. \(\bar\beta\)
is the central coarse coupling, \(\hat\beta_m\) the reference coupling,
\(\theta_m\) the threshold, \(J(m)\) the stage. \(x_m\) are coefficient
maps; \(V_L\) multiplies grade \(e\) by \(L^{d-e}\). \(f\circ g\)
denotes refinement of an extraction order. A chessboard root is the
absolute reflected-product expectation to the reciprocal number of
cells.

\subsubsection{7. What remains
unexplained}\label{what-remains-unexplained-15}

Quantitative transfer, remainder, shell and comparison contracts remain
those of III.7, III.11 and VII.4. Internal proofs of several constant
choices, contraction/remote-level decay, one-step label estimates and
auxiliary estimates were used as stated; moving-frame proofs belong to
VII.1/VII.3. The penalty's normalization relative to Balaban's parameter
remains unverified. The main theorem uses III.14's separate
construction: equality of its family with this one, or transfer of
III.13.5 to it, is not established.

Detailed notes: the full account of this unit, with proof sketches for
every step, is on the
\href{https://github.com/scottnarmstrong/yang-mills-informalization/blob/main/outline/notes/III-13.md}{detailed
notes page}.

\subsection{III.14. The construction run: flat representations and the
ultraviolet
rate}\label{iii.14.-the-construction-run-flat-representations-and-the-ultraviolet-rate}

\textbf{Part:} III, Ultraviolet construction (renormalization group).
\textbf{Lean:}
\href{https://github.com/yang-mills-group/OpenAI-YM-20261028/blob/3456e5e75cf8d4cd558955a5c7638fba2d68f464/YM/Independent/ConstructionRun.lean}{ConstructionRun},
\href{https://github.com/yang-mills-group/OpenAI-YM-20261028/blob/3456e5e75cf8d4cd558955a5c7638fba2d68f464/YM/Integration/NativeComponentRecipe.lean}{NativeComponentRecipe},
\href{https://github.com/yang-mills-group/OpenAI-YM-20261028/tree/3456e5e75cf8d4cd558955a5c7638fba2d68f464/YM/RG08/Flat}{Flat}.
\textbf{Size:} 14,290 lines (2,509 declarations) in the main-theorem
closure.

\subsubsection{1. What this unit proves}\label{what-this-unit-proves-16}

Throughout, \(G\) is a compact matrix gauge group; ``simple'' means its
real Lie algebra is simple and nonabelian. Level systems and full
horizon sufficiency are those of III.1. A sheet fixes the kernel, stage
profile and floors before block factors are chosen. Notation and norms
are specified in §6.

\textbf{Proposition III.14.1 (construction run).} For simple \(G\) and
\(b>0\), there exist a run \(R\), a campaign \(T\) and pillars \(W\) for
\((R,T)\), with penalty \(b\), satisfying the full thirteen-clause
funded condition of III.15 relative to \(\Delta_{10}\), and carrying an
attached native component recipe. The selection uses \(Q(i)=i+2\) and
support exponent \(100\) at every stage.

Attachment means that there exist a sheet \(s\), price map \(q\), recipe
\(C\) of III.14.2 and floors \(n\ge\max(n_s,E_C)\) such that \(R\) uses
\(s\) with floors raised to \(n\) and primary price map \(P_{s,q}\); the
recipe is realized on its own levels and horizon. Here \(n_s\) is the
sheet's floor, \(E_C\) combines the three component floors and two
numerical floors, and \(B_C\) is the maximum of the two post-length
bills. Set \(U_C(r,j)=\rho_C(r,j)^{-1}\) for the positive recipe widths
\(\rho_C\), and \(V_C(\ell,r,j)=B_C(\ell,j)\). For both request families
\(F(r,j)=U_C(r,j)\) and \(F(r,j)=V_C(\ell,r,j)\), every fixed \(r\)
admits \(m_0\) such that \(|F(r,j)|\le1+\log\log\hat\beta_m\) for
\(m\ge m_0\), \(j\le J_{m+1}\); also \(B_C(\ell,J_{m+1})\le\hat\beta_m\)
for every \(m\). Realization includes III.14.3 with \(e=0\).

\emph{Status:} theorem. \emph{Lean:}
\href{https://github.com/yang-mills-group/OpenAI-YM-20261028/blob/3456e5e75cf8d4cd558955a5c7638fba2d68f464/YM/Independent/ConstructionRun.lean\#L14-L28}{constructionRun},
\href{https://github.com/yang-mills-group/OpenAI-YM-20261028/blob/3456e5e75cf8d4cd558955a5c7638fba2d68f464/YM/R2I/Campaign/Accepted.lean\#L60-L73}{Funded},
\href{https://github.com/yang-mills-group/OpenAI-YM-20261028/blob/3456e5e75cf8d4cd558955a5c7638fba2d68f464/YM/Integration/ActualRunNativeComponentRecipe.lean\#L56-L69}{Attached}.

\textbf{Proposition III.14.2 (recipe).} For simple \(G\) and every sheet
\(s\), there exist a secondary price map \(X\) and floors
\(E^\flat,E^{\rm sw},E^{\rm du}:\mathbb N\to\mathbb N\). For every level
system with the kernel and profile of \(s\):

\begin{itemize}
\tightlist
\item
  If \(L_i\ge E^\flat(J_{i+1})\) for all \(i\), every entry stage \(e\)
  satisfying the delay condition in III.14.3 admits its conclusions with
  this \(X\), simultaneously for every sufficient horizon.
\item
  If \(L_i\ge E^{\rm sw}(J_{i+1})\) for all \(i\), every horizon
  sufficient for some price map has, for every \(i\le N\), every torus
  and all colour-valued link fields \(A,B\), colour-summed bilinear form
  \(\sum_a\langle A^a,\mathcal K^{\mathbb T}_{N,i}B^a\rangle\) invariant
  under simultaneous pullback by signed axis permutations fixing the
  origin. Reversing link orientation reverses the pullback sign.
\item
  If \(L_i\ge E^{\rm du}(J_{i+1})\) for all \(i\), then for every such
  horizon, every \(m<N\) and every level-system torus, the canonical
  free background map of III.3 satisfies III.14.5 for the fine/coarse
  metrics \((M^{\mathbb T_f}_{N,m+1},M^{\mathbb T_c}_{N,m})\), hence
  minimizes the actual Gaussian variational problem.
\end{itemize}

\emph{Status:} theorem. \emph{Lean:}
\href{https://github.com/yang-mills-group/OpenAI-YM-20261028/blob/3456e5e75cf8d4cd558955a5c7638fba2d68f464/YM/Integration/NativeComponentRecipe.lean\#L29-L51}{Choice},
\href{https://github.com/yang-mills-group/OpenAI-YM-20261028/blob/3456e5e75cf8d4cd558955a5c7638fba2d68f464/YM/Integration/NativeComponentRecipe.lean\#L54-L59}{available}.

\textbf{Proposition III.14.3 (Gaussian ultraviolet rate).} For simple
\(G\) and every sheet \(s\), there exist a price map \(X\) and
\(E:\mathbb N\to\mathbb N\) with this property: for every level system
with the kernel and profile of \(s\) satisfying \(L_i\ge E(J_{i+1})\)
for all \(i\), and every \(e\in\mathbb N\) satisfying
\(\Omega^X_{J_i}\le\log\hat\beta_i\) for \(i\ge e\), there are link
kernels \(F_i\) (\(i\ge e\)) and constants \(A,B\ge0\) with the
following properties. For every horizon \(H\) of that system sufficient
for some price map, every \(k\in\mathbb N\), \(i\ge e\), \(N\ge i\) and
links \(u,v\), \[\lVert F_i\rVert_{w_{J_i}}\le B\Lambda_{k,i},\qquad
\mathcal O^H_{n,i}(u,v)\longrightarrow F_i(u,v),\qquad
\lVert\mathcal O^H_{N,i}-F_i\rVert_{w_{J_i}}\le A\Lambda_{k,i}(a_N/a_i)^{1/4}.\]
Here \(n\to\infty\), \(w_j=\min(1,\alpha^X_j,\alpha^X_{j+1})\), and
\(A,B\ge0\) depend only on \(G\) and the lattice, independently of
\(s\), the levels, \(e\) and \(H\).

\emph{Status:} theorem, order zero. \emph{Lean:}
\href{https://github.com/yang-mills-group/OpenAI-YM-20261028/blob/3456e5e75cf8d4cd558955a5c7638fba2d68f464/YM/RG08/Flat/Entrance.lean\#L28-L50}{Flat.entrance\_gate}.

\textbf{Proposition III.14.4 (abstract flat representations).} Fix a
price map \(P\), \(\theta>0\), positive widths \(w_j\), \(A_0\ge0\),
\(d_0\in\mathbb N\), and a family
\((\mathcal L_x,H_x,r_x,N_x,n_x,k_x,Z_x)_{x\in X}\) with
\(r_x\le N_x\le n_x\), \(k_x\in\mathbb N\), \(Z_x\ge0\). Put
\(\Lambda_x(i)=2+Z_x+|s_i-s_{k_x}|\), using each system's scales. Assume
each \(H_x\) is sufficient for \(P\), its plaquette transfers at cutoff
\(N_x\) and \(r_x\le i<N_x\) are canonical free kernels, and
\(w_j\le\alpha_j\) for all \(j\). On those stage ranges assume
\[w_{J_i}\le\min(1,\alpha_{J_i},\alpha_{J_{i+1}}),\quad
2w_{J_i}\le L_iw_{J_{i+1}},\quad
q_i(w_{J_{i+1}})\le L_i^{-\theta},\] and assume
\(2|\mathcal F|\Omega_{J_{N_x}}\le A_0\Lambda_x(N_x)^{d_0}\) for every
\(x\). Then, simultaneously for all \(x\) and \(r_x\le i\le N_x\), there
are flat representatives of \(\mathcal K_{n_x,i}-\mathcal K_{N_x,i}\) of
width \(w_{J_i}\) and cost
\(A\Lambda_x(i)^{d_0}(a_{N_x}/a_i)^{\min(1,\theta/2)}\), where \(A\ge0\)
depends only on \(A_0,d_0,\theta\).

\emph{Status:} theorem. \emph{Lean:}
\href{https://github.com/yang-mills-group/OpenAI-YM-20261028/blob/3456e5e75cf8d4cd558955a5c7638fba2d68f464/YM/RG08/Flat/World.lean\#L165-L174}{Flat.World.overlap}.

\textbf{Proposition III.14.5 (Schur identity and flat step).} Fix
fine/coarse tori with side ratio \(L\), block average \(Q\) and \(b>0\).
The background-field relations are
\[\mathcal K(M)W=Q^\top\mathcal K(M'),\qquad QW=1-b^{-1}\mathcal K(M').\]
For symmetric positive semidefinite \(M\) and every coarse field \(A'\),
they make \(WA'\) a minimizer of
\(\langle A,\mathcal K(M)A\rangle+b|A'-QA|^2\), with minimum
\(\langle A',\mathcal K(M')A'\rangle\). If two triples
\((W_j;M_j,M'_j)\) satisfy these relations and \(M_2,M'_2\) are
symmetric, then
\[\mathcal K(M'_2)-\mathcal K(M'_1)=W_2^\top(\mathcal K(M_2)-\mathcal K(M_1))W_1.\]
Let \(B\) be a neutral translation-invariant plaquette representative of
the fine difference, with \(\lVert B\rVert^{\rm pl}_\tau\le r\),
\(\tau>0\), \(r\ge0\). Suppose \(T,U\) are order-two strip-bounded
curvature transports in width \(\omega\), with constants \(q,p\ge0\),
scaled covariance \(T(x+Lv,y+v)=T(x,y)\) and likewise for \(U\). Require
their periodizations \(C_2,C_1\) to satisfy \(dW_2=C_2d\),
\(dW_1=C_1d\), and \(\sum_zU((x,\varphi),(z,\psi))\) to be independent
of \(x\) for each orientation pair. If \(0\le\sigma\le\omega\) and
\(2\sigma\le L\tau\), then
\[B'(y,z)=\sum_{x,w}T(x,y)B(x,w)U(w,z),\qquad \lVert B'\rVert^{\rm pl}_\sigma\le\frac{2|\mathcal F|^2qp}{\tau L}r.\]
It is neutral, translation invariant and represents the coarse
difference.

\emph{Status:} theorem. \emph{Lean:}
\href{https://github.com/yang-mills-group/OpenAI-YM-20261028/blob/3456e5e75cf8d4cd558955a5c7638fba2d68f464/YM/RG08/Packets/FlatStep.lean\#L24-L38}{Flat.poised},
\href{https://github.com/yang-mills-group/OpenAI-YM-20261028/blob/3456e5e75cf8d4cd558955a5c7638fba2d68f464/YM/RG08/Packets/FlatStep.lean\#L50-L78}{Flat.toward}.

\textbf{Proposition III.14.6 (recursion rate).} For each \(x\in X\), fix
scales with \(L_i\ge3\), stages \(r_x\le i\le t_x\), reporting scale
\(k_x\in\mathbb N\), and \(\Lambda_x(i)=2+Z_x+|s_i-s_{k_x}|\),
\(Z_x\ge0\). Fix uniform \(\theta,c_0>0\), \(A_0\ge0\),
\(d_0\in\mathbb N\). If \(f_x,g_x\ge0\) satisfy
\(f_x(i)\le L_i^{-\theta}f_x(i+1)+g_x(i)\) below \(t_x\),
\(f_x(t_x)\le g_x(t_x)\), and
\(g_x(i)\le A_0\Lambda_x(i)^{d_0}(a_{t_x}/a_i)^{c_0}\) throughout, then
\[f_x(i)\le A\Lambda_x(i)^{d_0}(a_{t_x}/a_i)^\mu,\qquad \mu=\min(c_0,\theta/2).\]
One may take \(A=A_0(K+1)^{d_0}/(1-3^{-\theta'/2})\),
\(\theta'=\theta-\mu\), \(K=2(d_0+1)/\theta'\).

\emph{Status:} theorem. \emph{Lean:}
\href{https://github.com/yang-mills-group/OpenAI-YM-20261028/blob/3456e5e75cf8d4cd558955a5c7638fba2d68f464/YM/RG08/Rate/Ascent.lean\#L118-L134}{Rate.Stair.fade\_solve}.

\subsubsection{2. What it uses}\label{what-it-uses-16}

\begin{itemize}
\tightlist
\item
  I.1: group data and
  \href{https://github.com/yang-mills-group/OpenAI-YM-20261028/blob/3456e5e75cf8d4cd558955a5c7638fba2d68f464/YM/Gauge/Colors.lean\#L184}{one-loop
  slope} \(s_G\); III.1: levels,
  \href{https://github.com/yang-mills-group/OpenAI-YM-20261028/blob/3456e5e75cf8d4cd558955a5c7638fba2d68f464/YM/PREC/Stacks.lean\#L132-L139}{sufficiency}
  and background-field relations.
\item
  III.2:
  \href{https://github.com/yang-mills-group/OpenAI-YM-20261028/blob/3456e5e75cf8d4cd558955a5c7638fba2d68f464/YM/Integration/NativePreselectedPostnumericPeriodRun.lean\#L99-L114}{selection}
  with nonempty demands and countable numerical requests; III.3:
  \href{https://github.com/yang-mills-group/OpenAI-YM-20261028/blob/3456e5e75cf8d4cd558955a5c7638fba2d68f464/YM/Next/FREE_KERNELS/Rack.lean\#L134-L145}{sufficient
  free horizons}, canonical transfers and point-group invariance.
\item
  III.8:
  \href{https://github.com/yang-mills-group/OpenAI-YM-20261028/blob/3456e5e75cf8d4cd558955a5c7638fba2d68f464/YM/Pert/Bulk/Atlas.lean\#L37-L42}{weighted
  kernel bounds}; III.12:
  \href{https://github.com/yang-mills-group/OpenAI-YM-20261028/blob/3456e5e75cf8d4cd558955a5c7638fba2d68f464/YM/FullVolume/Cross.lean\#L30-L48}{runs},
  pillars and primary price maps; III.15: full
  \href{https://github.com/yang-mills-group/OpenAI-YM-20261028/blob/3456e5e75cf8d4cd558955a5c7638fba2d68f464/YM/R2I/Campaign/Accepted.lean\#L60-L73}{funded
  condition}.
\item
  Gaussian block-spin minimization (Balaban; Gawędzki--Kupiainen),
  exponentially weighted \(\ell^1\) composition, discrete Grönwall,
  dominated and monotone convergence.
\end{itemize}

\subsubsection{3. Strategy}\label{strategy-16}

Compare cutoffs at their common finest stage, where kinetic
normalization makes the metric difference neutral. Transport this
difference by the mixed Schur identity; neutrality gains \(L^{-1}\),
floors give contraction \(L^{-1/2}\), and the recursion estimate gives
exponent \(1/4\). A shadow with altered clip parameters supplies free
data; uniqueness transfers its Gaussian forms to every sufficient
horizon.

Choose recipe and floors before block factors, then raise coupling
thresholds to meet the post-length bill and countable requests. III.2
retains these conditions on the selected funded run. VIII.1 uses its
penalty, pillars and funded property, discarding the attached-recipe
conjunct.

\subsubsection{4. Steps}\label{steps-16}

\begin{enumerate}
\def\labelenumi{\arabic{enumi}.}
\tightlist
\item
  \textbf{Initialize and propagate.} The metric difference at stage
  \(N\) has cost \(2|\mathcal F|\Omega_{J_N}\) and zero orientation-wise
  row sums. The Schur identity transports it without inverting the
  curvature form; subtracting the transport kernel at the first base
  point gains \(L^{-1}\), using half the input exponential weight.
  \href{https://github.com/yang-mills-group/OpenAI-YM-20261028/blob/3456e5e75cf8d4cd558955a5c7638fba2d68f464/YM/RG08/Packets/FlatStep.lean\#L50-L78}{Flat.toward}.
\item
  \textbf{Obtain the rate.} Floors ensure compatible widths and
  \(q_i(w_{J_{i+1}})\le L_i^{-1/2}\); downward induction and III.14.6
  absorb age growth.
  \href{https://github.com/yang-mills-group/OpenAI-YM-20261028/blob/3456e5e75cf8d4cd558955a5c7638fba2d68f464/YM/RG08/Flat/Length.lean\#L64-L92}{spacing\_choice}.
\item
  \textbf{Take the limit.} Diverging-torus periodizations identify link
  kernels; composing with \(d^\top,d\) gives the pairwise cutoff
  estimate. Its Cauchy limit preserves the weighted bounds.
  \href{https://github.com/yang-mills-group/OpenAI-YM-20261028/blob/3456e5e75cf8d4cd558955a5c7638fba2d68f464/YM/RG08/Flat/Mesh.lean\#L107-L164}{supplies}.
\item
  \textbf{Transfer from the shadow.} Downward Schur induction from the
  bare terminal form proves equality of Gaussian forms for sufficient
  horizons on systems differing only in clip parameters. This transfers
  the rate, symmetries and free background identities.
  \href{https://github.com/yang-mills-group/OpenAI-YM-20261028/blob/3456e5e75cf8d4cd558955a5c7638fba2d68f464/YM/RG08/Flat/Agreement.lean\#L23-L36}{same\_path}.
\item
  \textbf{Realize and select.} The combined floor and bill enforce the
  recipe with \(e=0\) on the run's own horizon. Apply III.2 with
  reciprocal-width requests and \(V_C(\ell,r,j)=B_C(\ell,j)\).
  \href{https://github.com/yang-mills-group/OpenAI-YM-20261028/blob/3456e5e75cf8d4cd558955a5c7638fba2d68f464/YM/Integration/ActualRunNativeComponentRecipe.lean\#L25-L54}{on\_run}.
\end{enumerate}

\subsubsection{5. Where the difficulty
lies}\label{where-the-difficulty-lies-16}

Gauge zero modes obstruct curvature-form inversion. Orientation-wise
neutrality supplies the gain, but preserving it requires constant row
sums, beyond scaled covariance. Floors precede block factors; bills
follow them. The rate concerns only the Gaussian link kernel: it gives
no rate for the undetermined part of the plaquette metric,
positive-order coefficients or source kernels. The formal recipe does
not identify physical marked observables.

\subsubsection{6. Notation}\label{notation-16}

\(L_i=\ell_{J_{i+1}}\ge3\) are odd block factors; \(a_0=1\),
\(a_{i+1}=a_i/L_i\), \(s_i=\log(1/a_i)\); \(J_i\le J_{i+1}\le J_i+1\)
are smoothing indices. Cutoff \(N\) has spacing \(a_N\). Put
\(\hat\beta_i=\bar\beta+s_Gs_i>1\),
\(\Lambda_{k,i}=5+\log\hat\beta_k+|s_i-s_k|\). A price map is
\((\alpha_j,\Omega_j,\eta_j)\): positive widths, upper bounds, positive
coercivity bounds. A horizon supplies metric kernels \(M^\infty_{N,i}\)
and transfer data; \(M^{\mathbb T}_{N,i}\) is the periodization. The
lattice exterior derivative is \(d\); \(\mathcal K(M)=d^\top Md\),
\(\mathcal O=d^\top M^\infty d\), without a factor \(1/2\). The
local-estimate package \(\Delta_{10}\) is
\href{https://github.com/yang-mills-group/OpenAI-YM-20261028/blob/3456e5e75cf8d4cd558955a5c7638fba2d68f464/YM/RG12/Burst/Wardrobe.lean\#L63}{garment10}.

Write \(x(u)\) for a base point. For link kernels,
\(\lVert F\rVert_\sigma=\sup_u\sup_\phi\sum_{y\in\mathbb Z^4}|F(u,(y,\phi))|e^{\sigma|x(u)-y|_1}\),
with \(\phi\) a link orientation; for plaquettes,
\(\lVert B\rVert^{\rm pl}_\sigma=\sup_u\sum_{v\in\mathbb Z^4\times\mathcal F}|B(u,v)|e^{\sigma|x(u)-x(v)|_1}\),
summing all target orientations \(\mathcal F\). A flat representative is
translation invariant, has zero row sums in each target orientation, and
on every level-system torus periodizes to the quadratic-form difference
after applying \(d^\top(\cdot)d\); cost bounds this norm. In III.14.4,
\(q_i(u)=2|\mathcal F|^2h_i^2/(uL_i)\),
\(h_i=e^4|\mathcal F|^2\Omega_{J_{i+1}}^2\Omega_{J_i}\).

\subsubsection{7. What remains
unexplained}\label{what-remains-unexplained-16}

The scale and extraction choices are constructed in III.2; the meanings
and proofs of the full funded clauses belong to III.15 and later units.
The free-kernel existence proof and lattice kernel estimates are used
without explanation here. The auxiliary curvature-lift, marked-kernel
and formal Ward estimates still require their full admissibility
conditions and proofs. Removing the recipe demand would not identify the
resulting run with this selected run.

Detailed notes: the full account of this unit, with proof sketches for
every step, is on the
\href{https://github.com/scottnarmstrong/yang-mills-informalization/blob/main/outline/notes/III-14.md}{detailed
notes page}.

\subsection{III.15. Cost accounting for observables and the remainder
comparison}\label{iii.15.-cost-accounting-for-observables-and-the-remainder-comparison}

\textbf{Part:} III, Ultraviolet construction (renormalization group).
\textbf{Lean:} chapter 32:
\href{https://github.com/yang-mills-group/OpenAI-YM-20261028/tree/3456e5e75cf8d4cd558955a5c7638fba2d68f464/YM/RemainderComparison}{\texttt{YM/RemainderComparison/}},
\href{https://github.com/yang-mills-group/OpenAI-YM-20261028/tree/3456e5e75cf8d4cd558955a5c7638fba2d68f464/YM/Progress/DIRECTED_RECIPE}{\texttt{YM/Progress/DIRECTED\_RECIPE/}},
\href{https://github.com/yang-mills-group/OpenAI-YM-20261028/tree/3456e5e75cf8d4cd558955a5c7638fba2d68f464/YM/Next/SCHEDULE}{\texttt{YM/Next/SCHEDULE/}},
with readout and integration fragments. \textbf{Size:} 7,382 lines,
1,053 declarations; about 4,200 lines concern launch data.

\subsubsection{1. What this unit proves}\label{what-this-unit-proves-17}

\textbf{Proposition III.15.1 (pair-step rates).} Fix \(\rho>0\),
\(C\ge0\), \(r>\rho\), \(d\in\mathbb N\), \(e\in\mathbb R\), \(q>0\).
Put
\[s=p(y)\sqrt{\beta/y},\quad a=Ce^{-rs^2},\quad b=C(1+\mathcal B(y))^dp(y)^eg(y)^q,\quad \kappa=\rho s^2,\quad\gamma=e^\kappa a+2b.\]
For all sufficiently large \(y>1\), uniformly for \(\beta\in[y/4,4y]\)
and integers \(0\le B\le\mathcal B(y)\),
\[\kappa,a\ge0,\qquad0\le b\le1,\qquad(2B+1)^4\gamma\le1,\qquad\mathrm{quote}(\kappa,\gamma)\le e^{-\kappa}.\]
The threshold depends only on the fixed parameters. \emph{Status:}
theorem. \emph{Lean:}
\href{https://github.com/yang-mills-group/OpenAI-YM-20261028/blob/3456e5e75cf8d4cd558955a5c7638fba2d68f464/YM/Progress/DIRECTED_RECIPE/Economy.lean\#L76-L139}{\texttt{YM.Progress.DirectedRecipe.Payments.Pairing.currency}}.

\textbf{Proposition III.15.2 (hard--soft merging).} Fix a parameter
point with \(K\ge6\), direction power \(P_0\ge2\), \(B_P\ge1\),
\(\sigma>0\), \(0<t_P\le1\). Assume countably indexed regular terms
\(E_i\), forced activities \(F_i\) on \(X_i\), and real corrections
\(\ell_i\) on nonempty \(X'_i\). Require support regularity, flat-field
vanishing and vanishing affine Hessian on nonwrapping supports for
\(E_i\); gauge invariance, measurability and dependence on supported
links and final flags for \(F_i\), their eligibility tests and
\(\ell_i\); eligibility whenever \(F_i\ne0\); translation/reflection
covariance for all families, with prescribed permutations, support and
inventory transformations for regular and forced families; and
\(\ell_i=0\) on its small-field region. Suppose the regular family's
nonnegative derivative bounds through \(K\) on small-field supports, in
prescribed directions, have weighted anchored sum \(\le\delta\);
\(|F_i(U)|\le p_i\), \(|\ell_i(U)|\le q_i\) for all \(i,U\), with
\(p_i,q_i\ge0\); and
\[\mathrm{Hit}_{B_P}(p;\sigma+1)\le a,\qquad\mathrm{Hit}_{B_P}(q;\sigma+1)\le b.\]
Assume III.15.1's inequalities with \(B=B_P\), replacing its last bound
by \(\mathrm{quote}(\kappa,\gamma)\le w\). Then the fusion input
condition holds with \((\delta,w)\): original bounds \(\delta,a\)
persist, the soft family has a summable uniform majorant, and
\((F_i)\sqcup(e^{\ell_i}-1)\) has weights \((e^\kappa p_i)\sqcup(2q_i)\)
and hit bound \(\gamma\), with \(|F_i|\le e^{-\kappa}(e^\kappa p_i)\)
and \(|e^{\ell_i}-1|\le2q_i\). III.10's fusion gives norm \(\le w\).
\emph{Status:} theorem. \emph{Lean:}
\href{https://github.com/yang-mills-group/OpenAI-YM-20261028/blob/3456e5e75cf8d4cd558955a5c7638fba2d68f464/YM/Progress/DIRECTED_RECIPE/Tender.lean\#L124-L133}{\texttt{YM.Progress.DirectedRecipe.Payments.tender}}.

For III.15.3--4 fix \(K\ge6\), \(P\ge2\), \(0<t\le1\). Directions
satisfy \(\|A_e\|\le t(1+\mathrm{dist}(e,o))^P\).

\textbf{Proposition III.15.3 (slots).} For every \(\ell\in\mathbb N\),
slot \(s\) with \(|s|\le\ell\), finite \(S\subset\mathbb Z^4\)
containing its stencil, and site \(o\), the slot is \(K\) times
left-differentiable. If \(U\) is small-field on wrapped \(o+S\), every
derivative of order \(m\le K\) in admissible directions has norm
\(\le C_s(P,K,\ell)t\), where
\[C_s(P,K,\ell)=(1+\kappa_G^{-1}\|1\|)[2(\kappa_G+1)2^P(\ell+1)^{P+1}]^K.\]
At order zero the bound is \(t\); constants are independent of volume
and \(t\). \emph{Status:} theorem. \emph{Lean:}
\href{https://github.com/yang-mills-group/OpenAI-YM-20261028/blob/3456e5e75cf8d4cd558955a5c7638fba2d68f464/YM/Next/SCHEDULE/SlotPay.lean\#L108-L159}{\texttt{YM.Next.Scheduling.Wages.slotPays}}.

\textbf{Proposition III.15.4 (atoms).} For every \(k\in\mathbb N\),
\(k\)-leg atom \(a\), site \(o\), and \(U\) small-field on its
footprint, \(a(\cdot,o)\) is \(K\) times left-differentiable, with every
admissible derivative through \(K\) bounded by
\[|a|2^K[2^KC_s(P,K,\mathrm{len}(a))t]^k.\] \emph{Status:} theorem.
\emph{Lean:}
\href{https://github.com/yang-mills-group/OpenAI-YM-20261028/blob/3456e5e75cf8d4cd558955a5c7638fba2d68f464/YM/Next/SCHEDULE/TensorCost.lean\#L56-L71}{\texttt{YM.Next.Scheduling.Wages.fin}}.

\textbf{Proposition III.15.5 (launch data).} Let \(G\) be simple. Fix a
sheet of profiles \(P_j\), kernel, compatible price table, targets
\(R_{ji}\) with \(\rho_{R_{ji}}<-3/8\) for all \(j,i\), tolerances
\(h_{ji}>0\), libraries \(C_{ji}\) of order \(P_i\), and marks
\(r_{ji}\). There exist \(f:\mathbb N\to\mathbb N\) and \(C'_{ji}\ge1\)
giving quotation bounds with fine/coarse exponents
\((\sigma_j,\sigma_i+2)\) and factor \(C'_{ji}L^{-5}\). These require
support enlargement by \(r_{ji}\mathcal B(x)\), coarse power equal to
that of \(P_i\), fine exponent \(\sigma_j\), \(B_P=\mathcal B(x)\),
\(B_B\le2B_P\), \(t_P\le2t_B\), and acceptance. For every \(u\ge0\)
satisfying local analytic smallness and every \(\tau\ge0\), the bound
holds. The quotation inequality is
\(e^{(\sigma_i+2)N_{B_P}(Y)}\tau\,\mathrm{returnFee}\le C'_{ji}L^{-5}e^{\sigma_jN_{B_B}(X)}\tau\).
For each stage pair \((j,i)\) and every block side \(L\ge f_j\),
III.15.6 holds and
\[(\mathrm{rel}_{R_{ji}}+1)10c_{h_{R_{ji}}+5}/L+2\varepsilon_{R_{ji}}\le\sigma_j,\quad20(\sigma_i+2)\mathrm{toll}(r_{ji})\le L\sigma_j,\quad C'_{ji}/L\le h_{ji}.\]
\emph{Status:} theorem. \emph{Lean:}
\href{https://github.com/yang-mills-group/OpenAI-YM-20261028/blob/3456e5e75cf8d4cd558955a5c7638fba2d68f464/YM/RemainderComparison/Chron/Age.lean\#L98-L111}{\texttt{YM.RemainderComparison.Chron.sails}}.

\textbf{Proposition III.15.6 (large-side comparison).} For simple \(G\),
kernel \(K\) and target \(Q\) with \(\rho_Q<-3/8\), there exists \(n\)
such that for every \(L\ge n\) and branching patch, every \(W,\mu>0\),
order datum at \(W\), songbook and scrollwork at side \(L\), \(\chi\),
stages \(E,E'\), rays satisfying the heart condition for
\((E',\mathrm e_Q)\), library of order \(E'\), record with parameter
\(\le1/4\), \(j\in\mathbb N\), reserve \(f\), and \(\eta,\delta>0\),
there exist a financed quarry \(D\) at rate \(\rho_Q\), with \(d(D)>0\)
and paint supply, and \(P_1\). For all \(x\ge P_1\), poster comparison,
chart conditions and instructions hold; active, resonance and tolerance
conditions hold for every \(\beta\in[\beta_F/4,4\beta_F]\) and every
point tempered to \(E'\). The threshold \(n\) is independent of
auxiliary data; \(D,P_1\) may depend on them. \emph{Status:} theorem.
\emph{Lean:}
\href{https://github.com/yang-mills-group/OpenAI-YM-20261028/blob/3456e5e75cf8d4cd558955a5c7638fba2d68f464/YM/RemainderComparison/Cadenza/Syllabus.lean\#L47-L58}{\texttt{YM.RemainderComparison.Cadenza.buoys}}.

\textbf{Proposition III.15.7 (delay enlargement).} Fix a sheet,
compatible target/expense tables \(R,X\), sides \(L_j\), protection,
register, firmament, tolerances, libraries and marks as in III.15.5, and
records \(u_{ji}\). The register supplies positive \(W_{ji},\mu_{ji}\)
and orders on stage pairs; each finite request family supplies
scrollwork, songbook, grade/orders, reserve, positive constants, and an
extra condition eventually true for every financed quarry. Assume launch
thresholds \(k_j\le L_j\), \(\varepsilon(u_{ji})\le1/4\) on stage pairs
with \(i\) at least the first stage, and
\(\sigma_i+3<\mathrm e_{R_{ji}}\) for all \(j,i\). For every initial
delay \(T_0\), there exist \(W\ge T_0\) and quarries \(V_{jig}\) for
every request slot at every \((j,i)\), chosen before any \(U\ge W\). For
every such \(U\), launch data persist; on branching stage pairs with
\(j\ge4\) and \(i\) at least the first stage, every request has the
heart property, positive \(d\), paint supply, III.15.6's allotment at
threshold \(\beta_j^*(U)\), and its extra condition above that
threshold. On all stage pairs with \(i\) at least the first stage,
requests are tiled whenever \(x\ge\beta_j^*(U)\) and
\[c\ge\mathrm{shells}_j(U)(1+\mathcal B(x))^{\mathrm{growth}_j(U)}.\]
Tiling means \(2z_g(x)<c\) and \(4\mathrm{halo}_g(x)+16\le c\). Delay
ordering increases entry couplings, scale-entry functions, shell counts
and growth exponents componentwise. \emph{Status:} theorem. \emph{Lean:}
\href{https://github.com/yang-mills-group/OpenAI-YM-20261028/blob/3456e5e75cf8d4cd558955a5c7638fba2d68f464/YM/RemainderComparison/Chron/Corners.lean\#L33-L70}{\texttt{YM.RemainderComparison.Chron.grows}}.

\subsubsection{2. What it uses}\label{what-it-uses-17}

\begin{itemize}
\tightlist
\item
  III.1: scales, slots, atoms; II.1: left-derivative calculus; II.2: hit
  bounds; III.5: logarithmic cluster tails.
\item
  III.10: fusion, quotation, connected-activity and tree estimates;
  III.2: selected targets and delay ordering.
\item
  Leibniz and chain rules, Mayer expansion, Kotecký--Preiss bounds,
  elementary asymptotics, \(|e^z-1|\le2|z|\) for \(|z|\le1\).
\end{itemize}

\subsubsection{3. Strategy}\label{strategy-17}

Retain \(e^{-\rho s^2}\) from the hard bound and use the remaining
Gaussian decay to absorb support counting; soft bounds decay by their
negative coupling power. Express slots as single group words: unitarity
controls derivatives, and Leibniz gives \(t^k\) for atoms.

Choose the block-side threshold before auxiliary data, then comparison
functions and entry couplings. Negative activity rate survives the tree
weight \(\beta_F^{3/8}\); finitely many requests per stage permit common
delay enlargement.

\subsubsection{4. Steps}\label{steps-17}

\begin{enumerate}
\def\labelenumi{\arabic{enumi}.}
\tightlist
\item
  \textbf{Fusion.} Rescaled hard weights and doubled soft weights add to
  \(\gamma\); the exponential inequality bounds the Mayer terms.
  \href{https://github.com/yang-mills-group/OpenAI-YM-20261028/blob/3456e5e75cf8d4cd558955a5c7638fba2d68f464/YM/Progress/DIRECTED_RECIPE/Tender.lean\#L124-L133}{\texttt{YM.Progress.DirectedRecipe.Payments.tender}}.
\item
  \textbf{Observable bounds.} Differentiate the detour word, use
  unitarity, then multiply slot estimates.
  \href{https://github.com/yang-mills-group/OpenAI-YM-20261028/blob/3456e5e75cf8d4cd558955a5c7638fba2d68f464/YM/Next/SCHEDULE/SlotPay.lean\#L108-L159}{\texttt{YM.Next.Scheduling.Wages.slotPays}};
  \href{https://github.com/yang-mills-group/OpenAI-YM-20261028/blob/3456e5e75cf8d4cd558955a5c7638fba2d68f464/YM/Next/SCHEDULE/TensorCost.lean\#L56-L71}{\texttt{YM.Next.Scheduling.Wages.fin}}.
\item
  \textbf{Comparison functions.} Choose reserve functions conditional on
  local derivative estimates, charts and coefficient matching.
  \href{https://github.com/yang-mills-group/OpenAI-YM-20261028/blob/3456e5e75cf8d4cd558955a5c7638fba2d68f464/YM/RemainderComparison/Herbal/Orbit.lean\#L41-L56}{\texttt{YM.RemainderComparison.Herbal.compile}}
  \emph{(proof read; analytic inputs not expanded)}.
\item
  \textbf{Marked outputs.} For the conditional marked-output estimate,
  separate single remaining occurrences from clusters with further
  occurrences and from large supports. The \(L^{-5}\) quotation summed
  over \(L^4\) sites supplies the \(C_0/L\) term; the other
  contributions have negative rate. Occurrences are counted outside
  inserted leaves, with at least one marked.
  \href{https://github.com/yang-mills-group/OpenAI-YM-20261028/blob/3456e5e75cf8d4cd558955a5c7638fba2d68f464/YM/RemainderComparison/Solo/Barque.lean\#L21-L61}{\texttt{YM.RemainderComparison.Solo.atlas}}
  \emph{(proof read; singleton input not expanded)}.
\item
  \textbf{Launch and enlargement.} Combine comparison and tail bounds
  with quotation thresholds, then take finite stagewise maxima.
  \href{https://github.com/yang-mills-group/OpenAI-YM-20261028/blob/3456e5e75cf8d4cd558955a5c7638fba2d68f464/YM/RemainderComparison/Chron/Age.lean\#L98-L111}{\texttt{YM.RemainderComparison.Chron.sails}};
  \href{https://github.com/yang-mills-group/OpenAI-YM-20261028/blob/3456e5e75cf8d4cd558955a5c7638fba2d68f464/YM/RemainderComparison/Chron/Corners.lean\#L33-L70}{\texttt{YM.RemainderComparison.Chron.grows}}.
\end{enumerate}

\subsubsection{5. Where the difficulty
lies}\label{where-the-difficulty-lies-17}

The bounds must be uniform in coupling bands and volume; the block-side
threshold precedes all auxiliary choices, while comparison functions
precede the entry coupling. Part C's conditional comparison conclusions
reach the selected run but are discarded by the final construction-run
theorem.

\subsubsection{6. Notation}\label{notation-17}

\(g(\beta)=\beta^{-1/2}\), \(p(x)=e^{(\log\log x)^2}\),
\(t(x)=g(x)p(x)\), \(\mathcal B(x)=\lceil(\log x)^2\rceil\),
\(\beta_F=x+b_G\log L\); \(x\) is coarse inverse coupling.
\(\mathrm{quote}(\kappa,\gamma)=e^{-\kappa}(e^\gamma-1)\). Hit means
\(\sum_{i:z\in X_i}f_i e^{aN_M(X_i)}\le\gamma\) at every site \(z\), for
weights \(f_i\ge0\); \(N_M\) counts a minimal connected covering by
\(M\)-balls. Stage pairs have \(i\in\{j,j-1\}\); \(\sigma_j\) is the
profile's support exponent. A quarry is a record of bounding functions;
financing specifies reserve/rate bounds. Rate \(\le a\) means
\(|w|\le\beta_F^v\) eventually for every \(v>a\); a reserve satisfies
\(|w|\le p(\beta_F)^k\) eventually for some \(k\in\mathbb N\).

A slot is curvature transported along a walk \(\omega\), with length
\(|\omega|+4\); atoms are coefficient-weighted products of slots, with
coefficient norm \(|a|=\sum_b|c_b|\) and total length
\(\mathrm{len}(a)\). Small-field means every touching plaquette has
chord distance \(\le t\) from identity; \(\kappa_G\) is the colour scale
and \(\|1\|\) the Frobenius norm of identity.

For a target \(Q\),
\(\mathrm{rel}_Q=\mathrm e_Q+2\varepsilon_Q+\delta_Q\); \(h_Q\) is its
halo and \(c_{h_Q+5}\) its pre-cost. Targets carry their
chart/local-ratio conditions, \(\rho_Q\in(-1/2,0)\),
\(\mathrm e_Q\ge0\), \(\varepsilon_Q,\delta_Q>0\), \(h_Q\ge500\),
\(4\varepsilon_Q\) below the top support exponent and \(8\varepsilon_Q\)
below both chart widths.
\href{https://github.com/yang-mills-group/OpenAI-YM-20261028/blob/3456e5e75cf8d4cd558955a5c7638fba2d68f464/YM/AnalyticNext/CONNECTED_ACTIVITY_JETS_E6/OptionCost.lean\#L62-L76}{\texttt{YM.ConnectedActivityE6.Target}}.

The heart condition is
\[0\le\sigma_*,\quad\sigma_{E'}\le\sigma_*,\quad\sigma_*+1\le a_*,\quad\varepsilon_*>0,\quad a_*+3\varepsilon_*\le\mathrm e_Q,\]
\[\mathrm{ord}(E')\le J_*,\quad\mathrm{ord}(E')+(p_*+1)+p_*+2\le n_*,\quad p_*+1\le s_*,\quad k_*+1\le s_*,\]
with starred ray parameters. Tempered points have radius
\(\mathcal B(x)\), threshold \(t(x)\le2t(\beta_F)\), and the stage's
order, power and support exponent.

\subsubsection{7. What remains
unexplained}\label{what-remains-unexplained-17}

Local analytic, compatibility and readiness conditions remain to be
written out. Comparison-function existence does not supply coefficient
matching. The singleton, weighted-cluster, full-bill and
analytic-expansion inputs are not all checked here; readout, RG12 and
R2I fragments also remain outside the detailed verification.

Detailed notes: the full account of this unit, with proof sketches for
every step, is on the
\href{https://github.com/scottnarmstrong/yang-mills-informalization/blob/main/outline/notes/III-15.md}{detailed
notes page}.

\subsection{III.16. Analytic toolkit}\label{iii.16.-analytic-toolkit}

\textbf{Part:} III, Ultraviolet construction (renormalization group).
\textbf{Lean:} eight folders, none a numbered chapter of its own:
\href{https://github.com/yang-mills-group/OpenAI-YM-20261028/tree/3456e5e75cf8d4cd558955a5c7638fba2d68f464/YM/Complex}{\texttt{YM/Complex/}},
\href{https://github.com/yang-mills-group/OpenAI-YM-20261028/tree/3456e5e75cf8d4cd558955a5c7638fba2d68f464/YM/TimeKernel}{\texttt{YM/TimeKernel/}},
\href{https://github.com/yang-mills-group/OpenAI-YM-20261028/tree/3456e5e75cf8d4cd558955a5c7638fba2d68f464/YM/Analysis}{\texttt{YM/Analysis/}},
\href{https://github.com/yang-mills-group/OpenAI-YM-20261028/tree/3456e5e75cf8d4cd558955a5c7638fba2d68f464/YM/Laurent}{\texttt{YM/Laurent/}},
\href{https://github.com/yang-mills-group/OpenAI-YM-20261028/tree/3456e5e75cf8d4cd558955a5c7638fba2d68f464/YM/CellFourier}{\texttt{YM/CellFourier/}},
\href{https://github.com/yang-mills-group/OpenAI-YM-20261028/tree/3456e5e75cf8d4cd558955a5c7638fba2d68f464/YM/Work/F07}{\texttt{YM/Work/F07/}},
\href{https://github.com/yang-mills-group/OpenAI-YM-20261028/tree/3456e5e75cf8d4cd558955a5c7638fba2d68f464/YM/Maxwell}{\texttt{YM/Maxwell/}},
\href{https://github.com/yang-mills-group/OpenAI-YM-20261028/tree/3456e5e75cf8d4cd558955a5c7638fba2d68f464/YM/Labels}{\texttt{YM/Labels/}}.
\textbf{Size:} 5,318 physical source lines in 41 files, plus elementary
helper lemmas in other folders.

\subsubsection{1. What this unit proves}\label{what-this-unit-proves-18}

\textbf{Proposition III.16.1 (integrated positivity).} Let \(X\) be
compact metric with finite Borel measure \(\mu\),
\(\mathcal A,\mathcal B\) sets, \(K:\mathcal A^2\to\mathbb C\) positive
semidefinite, and \(g:\mathcal B\times X\to\mathcal A\). If
\((x,y)\mapsto K(g(b,x),g(c,y))\) is continuous for every \(b,c\), then
\(\tilde K(b,c)=\iint K(g(b,x),g(c,y))\,d\mu(y)d\mu(x)\) is positive
semidefinite. \emph{Status:} theorem. \emph{Lean:}
\href{https://github.com/yang-mills-group/OpenAI-YM-20261028/blob/3456e5e75cf8d4cd558955a5c7638fba2d68f464/YM/Analysis/Projected.lean\#L44-L70}{\texttt{YM.Analysis.NNKernel.accumulate}}.

\textbf{Proposition III.16.2 (trace Hölder).} For finite complex square
matrices of the same size \(A\succeq0,b_1,\ldots,b_n\), integers
\(q,M\ge1\), \(Z_A=\mathrm{tr}A^M>0\), and \(v\in\mathbb C^n\) with
\(\mathrm{Re}\,v_k\ge0\) and \(n/(2q)+\sum_k\mathrm{Re}\,v_k/M=1\),
\[\left|\mathrm{tr}\prod_k b_kA^{v_k}\right|\le C_A,\qquad C_A=\prod_k\|b_k\|_{2q}Z_A^{1-n/(2q)}.\]
\emph{Status:} theorem. \emph{Lean:}
\href{https://github.com/yang-mills-group/OpenAI-YM-20261028/blob/3456e5e75cf8d4cd558955a5c7638fba2d68f464/YM/Complex/Gaps.lean\#L124-L165}{\texttt{YM.Holder.estimation}}.

\textbf{Proposition III.16.3 (moving matrix gaps).} For such matrices,
integers \(q,w\ge1\), \(M=2qw\), \(Z_A>0\), and \(Y,Y'\in\mathbb N^n\),
suppose \(nw+\sum_kY_k=M\) and \(Y'=Y+\ell(e_i-e_j)\) for \(i\ne j\),
\(\ell\in\mathbb N\), with \(e_i\) the coordinate vectors. If \(r>0\),
\(Y_k>2r\) for every \(k\), and \(\ell\le r\), then
\[\left|\mathrm{tr}\prod_kb_kA^{Y'_k}-\mathrm{tr}\prod_kb_kA^{Y_k}\right|\le C_A\ell/r.\]
\emph{Status:} theorem. \emph{Lean:}
\href{https://github.com/yang-mills-group/OpenAI-YM-20261028/blob/3456e5e75cf8d4cd558955a5c7638fba2d68f464/YM/Complex/Gaps.lean\#L170-L229}{\texttt{YM.Holder.finite\_lip}}.

\textbf{Proposition III.16.4 (moving kernel gaps).} On \(X,\mu\) as in
III.16.1, let \(K\) be continuous positive semidefinite and
\(B_1,\ldots,B_n\) continuous complex kernels. For integers \(q,w\ge1\),
\(M=2qw\), real \(Z=\mathrm{Tr}(K^M)>0\), and \(Y,Y',\ell,r\) satisfying
all gap conditions of III.16.3, assume real \(c_k\) satisfy
\(|\mathrm{Tr}((B_k^*B_k)^q)|^{1/(2q)}\le c_kZ^{1/(2q)}\). Then
\[|\mathrm{Tr}(B_1K^{Y'_1}\cdots B_nK^{Y'_n})-\mathrm{Tr}(B_1K^{Y_1}\cdots B_nK^{Y_n})|\le Z\prod_kc_k\,\ell/r.\]
\emph{Status:} theorem. \emph{Lean:}
\href{https://github.com/yang-mills-group/OpenAI-YM-20261028/blob/3456e5e75cf8d4cd558955a5c7638fba2d68f464/YM/TimeKernel/Unit.lean\#L76-L106}{\texttt{YM.TimeKernel.passage}}.

\textbf{Proposition III.16.5 (Laurent estimates).} Let \(\alpha\) be
finite and \(\eta>0\). For holomorphic, \(2\pi\)-periodic
\(u:S_\eta\to\mathbb C\), define
\(\hat u_p(x)=\int_{[0,1]^\alpha}u(p+2\pi\theta)e^{-2\pi i\theta\cdot x}d\theta\),
\(\hat u=\hat u_0\). The torus mean is independent of \(p\in S_\eta\),
and \(\hat u_p(x)=e^{ip\cdot x}\hat u(x)\). If \(|u|\le K\), then
\(|\hat u_p(x)|\le K\) and, for \(H\ge0\), \(b>0\), \(H+b<\eta\),
\[\sum_xe^{H|x|_1}|\hat u(x)|\le K(2(1+b)/b)^{\mathrm{card}\,\alpha}.\]
If \(|u|\le K\), the Laurent series reconstructs \(u\) absolutely
throughout \(S_\eta\): \(u(p)=\sum_x\hat u(x)e^{ip\cdot x}\).

Conversely, for \(f:\mathbb Z^\alpha\to\mathbb C\), if
\(\sum_xe^{\eta|x|_1}|f(x)|\le W\), then \(|\hat f|\le W\) on
\(|\mathrm{Im}\,p_a|\le\eta\). For every \(0<\sigma<\eta\), \(\hat f\)
is holomorphic on \(S_\sigma\), with derivative norm and Lipschitz
constant at most \(W/(\eta-\sigma)\) for the sup norm. For real \(p\),
\(\widehat{\mathrm{Re}\,f}(p)=(\hat f(p)+\overline{\hat f(-p)})/2\).
\emph{Status:} theorems. \emph{Lean:}
\href{https://github.com/yang-mills-group/OpenAI-YM-20261028/blob/3456e5e75cf8d4cd558955a5c7638fba2d68f464/YM/Laurent/Turn.lean\#L103-L112}{\texttt{full\_turn}},
\href{https://github.com/yang-mills-group/OpenAI-YM-20261028/blob/3456e5e75cf8d4cd558955a5c7638fba2d68f464/YM/Laurent/Art.lean\#L61-L87}{\texttt{Vane.issues}},
\href{https://github.com/yang-mills-group/OpenAI-YM-20261028/blob/3456e5e75cf8d4cd558955a5c7638fba2d68f464/YM/Laurent/Art.lean\#L90-L104}{\texttt{Vane.paints}},
\href{https://github.com/yang-mills-group/OpenAI-YM-20261028/blob/3456e5e75cf8d4cd558955a5c7638fba2d68f464/YM/Laurent/Eval.lean\#L131-L133}{\texttt{Width.estimate}},
\href{https://github.com/yang-mills-group/OpenAI-YM-20261028/blob/3456e5e75cf8d4cd558955a5c7638fba2d68f464/YM/Work/F07/Read.lean\#L55-L95}{\texttt{readDeriv}},
\href{https://github.com/yang-mills-group/OpenAI-YM-20261028/blob/3456e5e75cf8d4cd558955a5c7638fba2d68f464/YM/Work/F07/Read.lean\#L97-L101}{\texttt{readDiff}},
\href{https://github.com/yang-mills-group/OpenAI-YM-20261028/blob/3456e5e75cf8d4cd558955a5c7638fba2d68f464/YM/Work/F07/Read.lean\#L119-L133}{\texttt{readRealPart}}.

\textbf{Proposition III.16.6 (Fourier and Bloch identities).} For
positive integers \(n,c,L\), finite internal index sets
\(\alpha,\beta\), and real entrywise summable
\(R:\mathbb Z^4\to\mathbb R^{\alpha\times\beta}\), put
\(R_n((y,a),(y',b))=\sum_{x\in\mathbb Z^4:y+x\equiv y'\pmod n}R_{ab}(x)\).
Then
\[F_nF_n^*=F_n^*F_n=n^4I,\qquad R_n(F_n\otimes I_\beta)=(F_n\otimes I_\alpha)\mathrm{diag}_p\hat R(2\pi p/n).\]
This includes \(p=0\). With \(f=cL\), reindexing fine momenta as
\(k_{p,l}=2\pi(p+cl)/f\), \(p\in\Lambda_c,l\in\Lambda_L\), gives
\(\Phi\Phi^*=f^4I\) and the same intertwining identity with \(\Phi\),
\(R_f\), and blocks \(\mathrm{diag}_{p,l}\hat R(k_{p,l})\). For every
\(p\in\mathbb C^4\), \(E_L(\bar p)^*E_L(p)=L^4I\). \emph{Status:}
theorems. \emph{Lean:}
\href{https://github.com/yang-mills-group/OpenAI-YM-20261028/blob/3456e5e75cf8d4cd558955a5c7638fba2d68f464/YM/CellFourier/Circulant.lean\#L88-L104}{\texttt{circulant}},
\href{https://github.com/yang-mills-group/OpenAI-YM-20261028/blob/3456e5e75cf8d4cd558955a5c7638fba2d68f464/YM/CellFourier/Mesh.lean\#L117-L119}{\texttt{gram}},
\href{https://github.com/yang-mills-group/OpenAI-YM-20261028/blob/3456e5e75cf8d4cd558955a5c7638fba2d68f464/YM/CellFourier/Mesh.lean\#L121-L127}{\texttt{coGram}},
\href{https://github.com/yang-mills-group/OpenAI-YM-20261028/blob/3456e5e75cf8d4cd558955a5c7638fba2d68f464/YM/CellFourier/Shuttle.lean\#L66-L75}{\texttt{readingsGram}},
\href{https://github.com/yang-mills-group/OpenAI-YM-20261028/blob/3456e5e75cf8d4cd558955a5c7638fba2d68f464/YM/CellFourier/Shuttle.lean\#L110-L129}{\texttt{aligns}},
\href{https://github.com/yang-mills-group/OpenAI-YM-20261028/blob/3456e5e75cf8d4cd558955a5c7638fba2d68f464/YM/CellFourier/Alias.lean\#L41-L53}{\texttt{straightens}}.

\textbf{Proposition III.16.7 (matrix root).} For \(m>0\), nonempty
finite matrix indices and complex square \(X\),
\[\Phi_m(X)=\sqrt m\sum_{k\ge0}\binom{1/2}{k}(X/m-I)^k\] is holomorphic
on \(\|X/m-I\|<1\), bounded on every smaller closed ball, and commutes
with entrywise conjugation. If \(v>0\) and \(X\) is Hermitian with
\(vI\le X\le mI\) as quadratic forms, its value is Hermitian, squares to
\(X\), and lies between \(\sqrt vI\) and \(\sqrt mI\). \emph{Status:}
theorem. \emph{Lean:}
\href{https://github.com/yang-mills-group/OpenAI-YM-20261028/blob/3456e5e75cf8d4cd558955a5c7638fba2d68f464/YM/Work/F07/MatrixRoot.lean\#L185-L198}{\texttt{binRoot\_on\_band}},
\href{https://github.com/yang-mills-group/OpenAI-YM-20261028/blob/3456e5e75cf8d4cd558955a5c7638fba2d68f464/YM/Work/F07/MatrixRoot.lean\#L127-L130}{\texttt{binRoot\_deriv}},
\href{https://github.com/yang-mills-group/OpenAI-YM-20261028/blob/3456e5e75cf8d4cd558955a5c7638fba2d68f464/YM/Work/F07/MatrixRoot.lean\#L116-L123}{\texttt{binSum\_bound}},
\href{https://github.com/yang-mills-group/OpenAI-YM-20261028/blob/3456e5e75cf8d4cd558955a5c7638fba2d68f464/YM/Work/F07/MatrixRoot.lean\#L138-L143}{\texttt{binRoot\_conj}}.

\textbf{Proposition III.16.8 (Maxwell jets).} Fix the compact matrix
gauge group and orthonormal real colour coordinates. For every integer
\(j\ge0\), curl is an isomorphism \(\mathcal A_j\to\mathcal C_j\), with
inverse \(A_i=(j+2)^{-1}\sum_kx_kF_{ki}\). For every smooth
antisymmetric colour-valued field satisfying Bianchi and Maxwell on an
open \(U\subset\mathbb R^4\), its \(j\)-th Taylor coefficient lies in
\(\mathcal S_j\) at every \(x\in U\).

Off zero, \(G_0=(4\pi^2|x|^2)^{-1}\) is harmonic; set
\(\psi_{b,km}=\delta_{kb}\partial_mG_0-\delta_{mb}\partial_kG_0\),
\(C_{ab,km}=\partial_a\psi_{b,km}-\partial_b\psi_{a,km}\). In either
index pair, \(C\) is smooth, antisymmetric and solves both equations;
\(C_{ab,km}(z)=C_{km,ab}(z)=C_{ab,km}(-z)\).

For all \(j,l\ge0\), real linear probes \(p,q\) on
\(\mathcal S_j,\mathcal S_l\) have a bilinear contraction
\(\kappa(p,q)(z)\), independent of extensions to unrestricted tensors,
smooth for \(z\ne0\), with \(\kappa(p,q)(z)=\kappa(q,p)(-z)\). For
coordinate probes \(p(F)=F_{ab}[u]_c\), \(q(F)=F_{km}[v]_d\),
\(u\in(\mathbb R^4)^j,v\in(\mathbb R^4)^l\),
\[\kappa(p,q)(z)=\delta_{cd}(-1)^lc_l(t\mapsto c_jC_{ab,km}(t)[u])(z)[v].\]
\emph{Status:} theorems; contraction definition. \emph{Lean:}
\href{https://github.com/yang-mills-group/OpenAI-YM-20261028/blob/3456e5e75cf8d4cd558955a5c7638fba2d68f464/YM/Maxwell/Exterior.lean\#L218-L224}{\texttt{potential}},
\href{https://github.com/yang-mills-group/OpenAI-YM-20261028/blob/3456e5e75cf8d4cd558955a5c7638fba2d68f464/YM/Maxwell/Shell.lean\#L73-L122}{\texttt{WaveOn.jets}},
\href{https://github.com/yang-mills-group/OpenAI-YM-20261028/blob/3456e5e75cf8d4cd558955a5c7638fba2d68f464/YM/Maxwell/Curls.lean\#L82-L98}{\texttt{Bare.equation}},
\href{https://github.com/yang-mills-group/OpenAI-YM-20261028/blob/3456e5e75cf8d4cd558955a5c7638fba2d68f464/YM/Maxwell/Leg.lean\#L161-L162}{\texttt{legKernel}},
\href{https://github.com/yang-mills-group/OpenAI-YM-20261028/blob/3456e5e75cf8d4cd558955a5c7638fba2d68f464/YM/Maxwell/Bilinear.lean\#L90-L105}{\texttt{reads\_mixed}}.

\textbf{Proposition III.16.9 (labels).} Fix real bases of
\(\mathcal A_j\); coordinate weight is \(j+2\). For each integer
\(d\ge0\), \(\mathrm{Label}_d\) consists of complex polynomials of
weight \(\le d\) invariant under simultaneous adjoint action. Real jet
evaluation determines polynomials. There are finitely many
weight-\(\le d\) monomials, so \(\mathrm{Label}_d\) is
finite-dimensional. Every label decomposes into invariant bihomogeneous
components (weight and degree both \(\le d\)), commuting with
conjugation and \(O(4)\). For \(s\in O(4)\), the action is
\((s\cdot P)(A)=P(T_{s^{-1}}A)\), where
\((T_sA_j)_i(x)=\sum_ks_{ik}(A_j)_k(s^{-1}x)\), and satisfies the group
law. Simultaneously for all \(d\), there exist complex bases of real
bihomogeneous labels containing \(e_0=1\), nested as sets of polynomials
with no further prescribed ordering. \emph{Status:} theorems.
\emph{Lean:}
\href{https://github.com/yang-mills-group/OpenAI-YM-20261028/blob/3456e5e75cf8d4cd558955a5c7638fba2d68f464/YM/Labels/Polynomial.lean\#L38-L49}{\texttt{eval\_ext}},
\href{https://github.com/yang-mills-group/OpenAI-YM-20261028/blob/3456e5e75cf8d4cd558955a5c7638fba2d68f464/YM/Labels/Polynomial.lean\#L114-L117}{\texttt{labels\_finite}},
\href{https://github.com/yang-mills-group/OpenAI-YM-20261028/blob/3456e5e75cf8d4cd558955a5c7638fba2d68f464/YM/Labels/Grading.lean\#L87-L105}{\texttt{decompose}},
\href{https://github.com/yang-mills-group/OpenAI-YM-20261028/blob/3456e5e75cf8d4cd558955a5c7638fba2d68f464/YM/Labels/Action.lean\#L82-L88}{\texttt{rotation\_law}},
\href{https://github.com/yang-mills-group/OpenAI-YM-20261028/blob/3456e5e75cf8d4cd558955a5c7638fba2d68f464/YM/Labels/Basis.lean\#L107-L111}{\texttt{frames}}.

\textbf{Proposition III.16.10 (Wick hierarchy).} For every \(n\ge0\),
integer cutoffs \(d_i\ge0\), distinct \(x_i\in\mathbb R^4\), and
\(P_i\in\mathrm{Label}_{d_i}\), restrict labels by
\(A_j=\mathrm{curl}^{-1}F_j\), \(F_j\in\mathcal S_j\). Sum products of
\(\kappa(\cdot,\cdot)(x_i-x_j)\) over pairings of monomial occurrences
at different points, extending multilinearly. This correlation is
smooth, kills any insertion vanishing on all shell jets, respects
simultaneous permutations and insertion of unit labels at additional
distinct points, and equals \(\prod_i a_i\) for constant insertions
\(a_i\in\mathbb C\), including the empty product \(1\). \emph{Status:}
definitions and theorems. \emph{Lean:}
\href{https://github.com/yang-mills-group/OpenAI-YM-20261028/blob/3456e5e75cf8d4cd558955a5c7638fba2d68f464/YM/Maxwell/Wick.lean\#L145-L147}{\texttt{Wick.correlation}},
\href{https://github.com/yang-mills-group/OpenAI-YM-20261028/blob/3456e5e75cf8d4cd558955a5c7638fba2d68f464/YM/Maxwell/Wick.lean\#L152-L154}{\texttt{Wick.smooth}},
\href{https://github.com/yang-mills-group/OpenAI-YM-20261028/blob/3456e5e75cf8d4cd558955a5c7638fba2d68f464/YM/Maxwell/Wick.lean\#L149-L150}{\texttt{Wick.kills\_equation}},
\href{https://github.com/yang-mills-group/OpenAI-YM-20261028/blob/3456e5e75cf8d4cd558955a5c7638fba2d68f464/YM/Maxwell/Reorder.lean\#L112-L124}{\texttt{Wick.exchangeSites}},
\href{https://github.com/yang-mills-group/OpenAI-YM-20261028/blob/3456e5e75cf8d4cd558955a5c7638fba2d68f464/YM/Maxwell/Reorder.lean\#L126-L144}{\texttt{Wick.adjoin}},
\href{https://github.com/yang-mills-group/OpenAI-YM-20261028/blob/3456e5e75cf8d4cd558955a5c7638fba2d68f464/YM/Maxwell/Reorder.lean\#L146-L153}{\texttt{Wick.constants}}.

\subsubsection{2. What it uses}\label{what-it-uses-18}

\begin{itemize}
\tightlist
\item
  I.1: lattices, subdivision, gauge colours, adjoint action and
  distinct-point configurations:
  \href{https://github.com/yang-mills-group/OpenAI-YM-20261028/blob/3456e5e75cf8d4cd558955a5c7638fba2d68f464/YM/Lattice/Site.lean\#L19}{\texttt{Lattice.Site}},
  \href{https://github.com/yang-mills-group/OpenAI-YM-20261028/blob/3456e5e75cf8d4cd558955a5c7638fba2d68f464/YM/Gauge/Colors.lean\#L43}{\texttt{Gauge.Color}}.
\item
  III.1: summability budgets:
  \href{https://github.com/yang-mills-group/OpenAI-YM-20261028/blob/3456e5e75cf8d4cd558955a5c7638fba2d68f464/YM/RG/Activities.lean\#L59-L60}{\texttt{RG.SumBudget}}.
\item
  Three-lines theorem, Cauchy estimates and contour formulas, mean value
  inequality, spectral theorem, Schur products, dominated convergence,
  Fourier inversion, Banach-algebra binomial series, symmetry of
  derivatives and polynomial identity theorem. All analytic assumptions
  are stated above; parameters are chosen in the applications.
\end{itemize}

\subsubsection{3. Strategy}\label{strategy-18}

Moving an insertion transfers units between two gaps. Continue these
gaps holomorphically while preserving the Hölder balance; Cauchy's
derivative estimate controls the change by \(\ell/r\). Sampling all
kernels and adjoints with one partition passes this dimension-free
estimate to cyclic integrals.

Shift Fourier contours to exchange strip holomorphy for exponential
decay; diagonalize finite tori and construct matrix roots by a binomial
series on a thinner strip. Radial gauge identifies potentials with
closed jets; the free kernel descends to shell probes and supplies Wick
correlations.

\subsubsection{4. Steps}\label{steps-18}

\begin{enumerate}
\def\labelenumi{\arabic{enumi}.}
\tightlist
\item
  Factor matrices into normalized entire paths with operator norm
  \(\le1\) on \(0\le\mathrm{Re}\,z\le1\) and trace norm \(\le1\) on its
  right boundary. At exponents with nonnegative real parts summing to
  \(1\), three-lines induction bounds the product trace by \(1\), giving
  dimension-free Hölder.
  \href{https://github.com/yang-mills-group/OpenAI-YM-20261028/blob/3456e5e75cf8d4cd558955a5c7638fba2d68f464/YM/Complex/Holder.lean\#L73-L126}{Lean}.
\item
  Vary gaps along \(Y+z(e_i-e_j)\); Cauchy's estimate on radius-\(r\)
  circles gives III.16.3.
  \href{https://github.com/yang-mills-group/OpenAI-YM-20261028/blob/3456e5e75cf8d4cd558955a5c7638fba2d68f464/YM/Complex/Gaps.lean\#L170-L229}{Lean}.
\item
  Positive quadratic forms pass through quadrature, proving integrated
  positivity. Sample kernels with square-root cell masses; positivity,
  adjoints and cyclic traces survive, and quadrature passes the trace
  estimate to kernels.
  \href{https://github.com/yang-mills-group/OpenAI-YM-20261028/blob/3456e5e75cf8d4cd558955a5c7638fba2d68f464/YM/Analysis/Projected.lean\#L44-L70}{\texttt{accumulate}},
  \href{https://github.com/yang-mills-group/OpenAI-YM-20261028/blob/3456e5e75cf8d4cd558955a5c7638fba2d68f464/YM/TimeKernel/Trace.lean\#L109-L119}{Lean}.
\item
  Shift integration contours and sum geometric tails; Fourier inversion
  reconstructs symbols; termwise differentiation gives the reverse
  derivative estimate.
  \href{https://github.com/yang-mills-group/OpenAI-YM-20261028/blob/3456e5e75cf8d4cd558955a5c7638fba2d68f464/YM/Laurent/Art.lean\#L90-L104}{\texttt{Vane.paints}},
  \href{https://github.com/yang-mills-group/OpenAI-YM-20261028/blob/3456e5e75cf8d4cd558955a5c7638fba2d68f464/YM/Work/F07/Read.lean\#L55-L95}{\texttt{readDeriv}}.
\item
  Character orthogonality diagonalizes periodized kernels and their
  alias blocks, retaining zero momentum.
  \href{https://github.com/yang-mills-group/OpenAI-YM-20261028/blob/3456e5e75cf8d4cd558955a5c7638fba2d68f464/YM/CellFourier/Circulant.lean\#L88-L104}{\texttt{circulant}},
  \href{https://github.com/yang-mills-group/OpenAI-YM-20261028/blob/3456e5e75cf8d4cd558955a5c7638fba2d68f464/YM/CellFourier/Alias.lean\#L41-L53}{\texttt{straightens}}.
\item
  Diagonalization reduces the binomial matrix series to scalar square
  roots; Banach-algebra convergence gives holomorphy.
  \href{https://github.com/yang-mills-group/OpenAI-YM-20261028/blob/3456e5e75cf8d4cd558955a5c7638fba2d68f464/YM/Work/F07/MatrixRoot.lean\#L185-L198}{\texttt{binRoot\_on\_band}}.
\item
  Euler's identity and Bianchi invert curl; shell restriction makes
  contractions independent of extensions.
  \href{https://github.com/yang-mills-group/OpenAI-YM-20261028/blob/3456e5e75cf8d4cd558955a5c7638fba2d68f464/YM/Maxwell/Exterior.lean\#L163-L195}{Lean},
  \href{https://github.com/yang-mills-group/OpenAI-YM-20261028/blob/3456e5e75cf8d4cd558955a5c7638fba2d68f464/YM/Maxwell/Leg.lean\#L148-L159}{Lean}.
\item
  Finite Wick sums give smoothness; real bihomogeneous decomposition
  gives nested frames.
  \href{https://github.com/yang-mills-group/OpenAI-YM-20261028/blob/3456e5e75cf8d4cd558955a5c7638fba2d68f464/YM/Maxwell/Wick.lean\#L152-L154}{Lean},
  \href{https://github.com/yang-mills-group/OpenAI-YM-20261028/blob/3456e5e75cf8d4cd558955a5c7638fba2d68f464/YM/Labels/Basis.lean\#L107-L111}{Lean}.
\end{enumerate}

\subsubsection{5. Where the difficulty
lies}\label{where-the-difficulty-lies-18}

The trace bound must remain independent of matrix dimension as gaps vary
and quadrature converges at fixed exponents and period. Strict gap
inequalities avoid the exceptional zeroth-power convention. Laurent
estimates lose strip width; complex alias adjoints require \(\bar p\).
Maxwell contractions retain factorials, the colour normalization and the
second-endpoint sign.

\subsubsection{6. Notation}\label{notation-18}

\(\|\cdot\|\): operator norm;
\(\|b\|_{2q}=(\mathrm{tr}((b^*b)^q))^{1/(2q)}\). Products are ordered.
\(A^z=\sum_{\lambda>0}e^{z\log\lambda}P_\lambda\) for spectral
projections \(P_\lambda\); \(A^0\) projects onto the range. For
\(L\ge1\),
\(\mathrm{Tr}(k_1\cdots k_L)=\int_{X^L}\prod_{l=1}^Lk_l(x_l,x_{l+1})\,d\mu^{\otimes L}\),
\(x_{L+1}=x_1\); powers repeat kernels; \(B^*(x,y)=\overline{B(y,x)}\).
Positive semidefinite means every finite Gram matrix, allowing repeated
points, is Hermitian positive semidefinite.

\(|x|_1=\sum_a|x_a|\),
\(S_\eta=\{p\in\mathbb C^\alpha:|\mathrm{Im}\,p_a|<\eta\ \forall a\}\),
\(\hat f(p)=\sum_xf(x)e^{ip\cdot x}\).
\(\Lambda_n=(\mathbb Z/n\mathbb Z)^4\); \(F_n=(e^{2\pi ip\cdot x/n})\),
\(\Phi=(e^{ik_{p,l}\cdot y})_{y\in\Lambda_f,(p,l)\in\Lambda_c\times\Lambda_L}\),
\(E_L(p)=(e^{i(p+2\pi l)\cdot x/L})_{x,l\in\Lambda_L}\). Residues use
representatives \(0,\ldots,n-1\).

Colour coordinates are orthonormal for
\(-(\mathrm{dim}\,\rho)^{-1}\mathrm{Re}\,\mathrm{tr}(XY)\) in the fixed
defining representation. \(\mathcal A_j\): colour-valued \(4\)-tuples of
homogeneous degree-\((j+1)\) potentials with \(x\cdot A=0\);
\(\mathcal C_j\): antisymmetric homogeneous degree-\(j\) curvature jets
satisfying \(\partial_aF_{bc}+\partial_bF_{ca}+\partial_cF_{ab}=0\);
\(\mathcal S_j\): additionally divergence-free. \(c_j=D^j/j!\).

\subsubsection{7. What remains
unexplained}\label{what-remains-unexplained-18}

III.6 supplies norm conversion and strip choices for square roots; IV.1
supplies gap parameters needed for uniformity. External helper lemmas
and longer dependency lists remain unchecked individually.

Detailed notes: the full account of this unit, with proof sketches for
every step, is on the
\href{https://github.com/scottnarmstrong/yang-mills-informalization/blob/main/outline/notes/III-16.md}{detailed
notes page}.

\subsection{IV.1. Infinite volume: extraction of a subsequence and the
Euclidean axioms of the
limit}\label{iv.1.-infinite-volume-extraction-of-a-subsequence-and-the-euclidean-axioms-of-the-limit}

\textbf{Part:} IV, infinite volume (chapter 10). \textbf{Lean:}
\href{https://github.com/yang-mills-group/OpenAI-YM-20261028/tree/3456e5e75cf8d4cd558955a5c7638fba2d68f464/YM/Extraction}{\texttt{YM/Extraction}},
\href{https://github.com/yang-mills-group/OpenAI-YM-20261028/tree/3456e5e75cf8d4cd558955a5c7638fba2d68f464/YM/Cells}{\texttt{YM/Cells}},
\href{https://github.com/yang-mills-group/OpenAI-YM-20261028/tree/3456e5e75cf8d4cd558955a5c7638fba2d68f464/YM/Slice}{\texttt{YM/Slice}},
\href{https://github.com/yang-mills-group/OpenAI-YM-20261028/tree/3456e5e75cf8d4cd558955a5c7638fba2d68f464/YM/Euclid}{\texttt{YM/Euclid}},
\href{https://github.com/yang-mills-group/OpenAI-YM-20261028/tree/3456e5e75cf8d4cd558955a5c7638fba2d68f464/YM/Dash}{\texttt{YM/Dash}},
\href{https://github.com/yang-mills-group/OpenAI-YM-20261028/tree/3456e5e75cf8d4cd558955a5c7638fba2d68f464/YM/Next/EXTRACTION}{\texttt{YM/Next/EXTRACTION}},
\href{https://github.com/yang-mills-group/OpenAI-YM-20261028/tree/3456e5e75cf8d4cd558955a5c7638fba2d68f464/YM/Lane}{\texttt{YM/Lane}};
D01--D06. \textbf{Size:} 7,444 lines in 1,049 declarations used by the
main theorem.

\subsubsection{1. What this unit proves}\label{what-this-unit-proves-19}

Let \(C\) be a Wilson family for a compact matrix gauge group.
\textbf{(H1)} is its specification (II.1--II.2, III.13;
\href{https://github.com/yang-mills-group/OpenAI-YM-20261028/blob/3456e5e75cf8d4cd558955a5c7638fba2d68f464/YM/Source/Construction.lean\#L44-L62}{\texttt{Family.Makes}}).
\textbf{(H2)} comprises fixed-volume centered budgets, ordered-tail
power budgets, reflection positivity and chessboard inequalities,
including root subadditivity
(\href{https://github.com/yang-mills-group/OpenAI-YM-20261028/blob/3456e5e75cf8d4cd558955a5c7638fba2d68f464/YM/Extraction/Center.lean\#L30-L34}{\texttt{CenterPort}}).
Extractions select strictly increasing spacing and volume indices;
ordered tails mean \(k\) large, then \(j\) large depending on \(k\). For
every fixed \(p\) and label \(v\), at every sufficiently fine scale
\(m\), (H2) gives a constant bounding all centered roots and the
interior condition for all sites and large \(j\), with
\(N_j\ge\max(m,\mathrm{ref}(d))\). For every degree \(d\), \(R>0\) and
\(\varepsilon>0\), it also gives \(m_0\ge\mathrm{ref}(d)\) and \(A\ge0\)
such that, for every \(m\ge m_0\) and \(\eta>0\), eventually along
ordered tails with \(N\ge m\), all labels \(\|v\|\le R\) and sites
satisfy the centered interior condition and root bound
\(Aa_m^{-d-\varepsilon}+\eta\). Positivity and chessboard laws hold at
every cutoff and volume.

\textbf{Proposition IV.1.1 (raw limits).} Under (H1)--(H2) along any
extraction \(f\), there exist a refinement \(((N_j),(p_k))\) and
multilinear tables \(S,P_p\) such that, for every arity \(n\) and label
list \(v\):

\begin{enumerate}
\def\labelenumi{(\alph{enumi})}
\item
  For every \(p\), \(P_p(v;\cdot)\) is continuous on
  \(\mathrm{Gen}^{(p)}_n\). For every compact \(K\) therein and
  \(\varepsilon>0\), eventually in \(j\),
  \(|\mathcal C_{N_j,p}(v;r)-P_p(v;a_{N_j}r)|<\varepsilon\) for all
  lifted lattice tuples \(r\) with \(a_{N_j}r\in K\).
\item
  \(S(v;\cdot)\) is continuous on \(\mathrm{Gen}_n\), and
  \(P_{p_k}\to S\) locally uniformly there.
\item
  For every \(t\in\mathcal D_n\), \(t\) is eventually admissible at
  \(p_k\); at every admissible \(p\),
  \(V_p=\lim_j\mathcal C_{N_j,p}(v;t)\) exists, and
  \(V_{p_k}\to\int tS(v;\cdot)\).
\item
  \(tS(v;\cdot)\) is integrable for every \(t\in\mathcal D_n\).
\end{enumerate}

\emph{Status:} theorem. \emph{Lean:}
\href{https://github.com/yang-mills-group/OpenAI-YM-20261028/blob/3456e5e75cf8d4cd558955a5c7638fba2d68f464/YM/Extraction/Held.lean\#L136-L137}{\texttt{Extraction.crops}}.

\textbf{Proposition IV.1.2 (label derivatives).} Suppose \(S\) satisfies
IV.1.1(c)--(d), with (H1)--(H2) along \(f\). Assume the relay properties
of §2 for every degree, projections \(\pi_d(N,p)\) on
\(\mathrm{Label}_d\) converging in operator norm along ordered tails to
projections \(\pi_d\), and \(\dim\mathrm{Null}_d(S)\le\dim\ker\pi_d\).
Then \(\ker\pi_d=\mathrm{Null}_d(S)\), and there exist linear maps
\(\kappa_{d,e}:\mathrm{Label}_d\to\mathrm{Label}_e\) for every
\(d\le e\), and \(\tau_{d,a}:\mathrm{Label}_d\to\mathrm{Label}_{d+1}\)
for every axis, such that for all arities, slots, labels and
\(t\in\mathcal D_n\),

\[\int tS(v_{i\to\kappa v_i})=\int tS(v),\qquad \int tS(v_{i\to\tau_av_i})=-\int(\partial_{x_i^a}t)S(v).\]

\emph{Status:} conditional theorem. \emph{Lean:}
\href{https://github.com/yang-mills-group/OpenAI-YM-20261028/blob/3456e5e75cf8d4cd558955a5c7638fba2d68f464/YM/Work/D04.lean\#L212-L240}{\texttt{Work.D04.answer}},
\href{https://github.com/yang-mills-group/OpenAI-YM-20261028/blob/3456e5e75cf8d4cd558955a5c7638fba2d68f464/YM/Dash/Inert.lean\#L71-L87}{\texttt{lifts\_model}}.

\textbf{Proposition IV.1.3 (smooth limit).} If \(S\) satisfies IV.1.1
along an extraction where IV.1.2 applies, there exists \(T\), smooth on
every \(\mathrm{Conf}_n\), equal to \(S\) on \(\mathrm{Gen}_n\),
satisfying IV.1.1(a)--(d) with the same extraction and \(P_p\), and the
basic identities below. The maps descend to
\(\mathrm{rise}_{d,e}:F_d\to F_e\) and
\(\partial_{d,a}:F_d\to F_{d+1}\). For every arity, label list, slot and
\(x\in\mathrm{Conf}_n\),

\[T(\dots,\mathrm{rise}\,u,\dots)=T(\dots,u,\dots),\qquad T(v_{i\to\partial_av_i};x)=\partial_{x_i^a}T(v;x).\]

The smooth ordered limit on \(\mathrm{Conf}_n\) along a given extraction
is unique. \emph{Status:} theorem. \emph{Lean:}
\href{https://github.com/yang-mills-group/OpenAI-YM-20261028/blob/3456e5e75cf8d4cd558955a5c7638fba2d68f464/YM/Org/Input.lean\#L101-L109}{\texttt{Org.Arrives.serves}}.

\textbf{Proposition IV.1.4 (Euclidean hierarchy).} Assume (H1),
IV.1.1(a)--(c) for \(T\) along \(f\), and smoothness on every
\(\mathrm{Conf}_n\). Then, for every arity, degree list and labels, on
\(\mathrm{Conf}_n\):

\begin{itemize}
\tightlist
\item
  \textbf{(E1)} Simultaneous permutations preserve \(T\); the empty
  correlation is \(1\); inserting \(1_d\) deletes that slot.
\item
  \textbf{(E2)} \(T(\rho(s)v;sx)=T(v;x)\) for every \(s\in H_4\), and
  \(T(\bar v;x)=\overline{T(v;x)}\).
\item
  \textbf{(E3)} \(T(v;x+b)=T(v;x)\) for every \(b\in\mathbb R^4\).
\item
  \textbf{(E4)} For every axis \(a\), finite collection of words \(w_i\)
  with distinct points within each word and \(x^a>0\), and
  \(c\in\mathbb C^m\), \(\sum_{i,j}\bar c_ic_jT(\Theta_aw_iw_j)\) is
  real and nonnegative.
\item
  \textbf{(E5)} The same holds for smeared words \((v^{(i)},t_i)\),
  \(t_i\in\mathcal D_{n_i}\) supported where all \(x_k^a>0\), with
  entries
  \(\iint\overline{t_i(y)}t_j(z)T(\Theta_a(v^{(i)},y)(v^{(j)},z))\);
  these integrals converge absolutely.
\end{itemize}

With \textbf{(E7)}, this gives (E1)--(E7): \textbf{(E6)} is smoothness;
\textbf{(E7)} requires, for every arity, degree list, labels and
\(j\ge0\), some \(A\in\mathbb N,C\ge0\) with
\(\|D^jT(v;x)\|\le Cr^{-A}\) for all \(0<r\le1\) and all \(r\)-separated
configurations. The constants are independent of \(r,x\); \(D^j\) has
its real Fréchet operator norm. \emph{Status:} theorem, conditional on
(E7) for assembly. \emph{Lean:}
\href{https://github.com/yang-mills-group/OpenAI-YM-20261028/blob/3456e5e75cf8d4cd558955a5c7638fba2d68f464/YM/Extraction/Harvest.lean\#L62-L65}{\texttt{Order.assembles}}.

\textbf{Proposition IV.1.5 (short-distance bound).} Assume (H1)--(H2)
along \(f\), smoothness of \(T\) for every arity and labels, and
IV.1.1(a)--(b). For every arity, degree list \((d_i)\) and labels \(v\),
some \(C\ge0\), independent of \(r,x\), satisfies
\(|T(v;x)|\le Cr^{-4n(1+\max_i d_i)}\) for all \(0<r\le1\) and
\(r\)-separated \(x\). The empty maximum is \(0\). With IV.1.3, B10
(V.1) supplies (E7). \emph{Status:} theorem. \emph{Lean:}
\href{https://github.com/yang-mills-group/OpenAI-YM-20261028/blob/3456e5e75cf8d4cd558955a5c7638fba2d68f464/YM/Work/D05.lean\#L160-L180}{\texttt{Work.D05.answer}}.

\textbf{Proposition IV.1.6 (symmetries of derivatives).} For smooth,
permutation-symmetric \(T\) satisfying (E2), let
\((\mathrm{rise},\partial)\) satisfy IV.1.3's identities. For every
\(d\le e\), axis and \(s\in H_4\), these maps commute with conjugation,
\(\mathrm{rise}\) commutes with \(\rho(s)\), and

\[\rho_{d+1}(s)\partial_b=\sum_a s_{ab}\partial_a\rho_d(s).\]

\emph{Status:} theorem. \emph{Lean:}
\href{https://github.com/yang-mills-group/OpenAI-YM-20261028/blob/3456e5e75cf8d4cd558955a5c7638fba2d68f464/YM/Work/D06.lean\#L98-L135}{\texttt{Work.D06.compat}}.

\subsubsection{2. What it uses}\label{what-it-uses-19}

\begin{itemize}
\tightlist
\item
  II.1--II.2, III.13: Wilson family specification (H1).
\item
  III.11: centered budgets (H2), supplied by
  \href{https://github.com/yang-mills-group/OpenAI-YM-20261028/blob/3456e5e75cf8d4cd558955a5c7638fba2d68f464/YM/MergeProof/Gardens.lean\#L19-L40}{\texttt{MergeProof.Accent.gardena}}.
\item
  III.10--III.11: relay approximations to differences and inclusions,
  null-label detection and the separate dimension comparison
  (\href{https://github.com/yang-mills-group/OpenAI-YM-20261028/blob/3456e5e75cf8d4cd558955a5c7638fba2d68f464/YM/MergeProof/Eureka.lean\#L16-L20}{\texttt{MergeProof.Scotia}},
  \href{https://github.com/yang-mills-group/OpenAI-YM-20261028/blob/3456e5e75cf8d4cd558955a5c7638fba2d68f464/YM/MergeProof/Powder.lean\#L19-L47}{\texttt{savings}}).
\item
  III.16: positive-kernel trace Lipschitz estimates. I.1, III.2, V.1:
  gentle scales and subpolynomial column products. V.1: B10.
\item
  Osterwalder--Seiler positivity; chessboard estimates; Schatten Hölder,
  three-lines and Cauchy estimates; Arzelà--Ascoli, dominated
  convergence and Fourier regularity.
\end{itemize}

The relay hypotheses require bounded projections at each fixed volume
and along ordered tails; degree-\((d+1)\) approximants to degree-\(d\)
forward differences and degree-\(e\) approximants to degree-\(d\) labels
for \(d\le e\), with subpolynomial cutoff coordinates; and a null-label
detector. For each approximation error there must be one rate
\(y(N,p)\to0\): for every integer \(q\ge0\) and lower scale \(m\), some
\(k\ge m\) gives, eventually along ordered tails with \(N\ge k\),
interior errors with roots \(\le y(N,p)\) in every translated bond cell
of side \(a_k\), at relative lattice coordinates satisfying \(M\le z^a\)
and \(z^a+M+1\le a_k/a_N\), where
\(M=\lfloor(a_k/a_N)/(8(q+1))\rfloor\). Projection-kernel labels with
subpolynomial growth must be chessboard-close to bounded measurable
observables whose separated spectator correlations vanish exactly,
eventually along ordered tails. Subpolynomial means
\(\|B_{d,N,p}^{-1}w(N,p)\|\le a_N^{-\varepsilon}\) eventually for every
\(\varepsilon>0\), where \(B\) is the terminal transfer to reference
coordinates. The dimension comparison is an additional hypothesis
(\href{https://github.com/yang-mills-group/OpenAI-YM-20261028/blob/3456e5e75cf8d4cd558955a5c7638fba2d68f464/YM/Bid/Drops.lean\#L46-L52}{\texttt{RelayBin}}).

\subsubsection{3. Strategy}\label{strategy-19}

Chessboard bounds control separated correlations. Transfer kernels turn
movement of an insertion into movement between gaps; complex
interpolation gives equicontinuity. Two compactness arguments select
spacing and volume subsequences. Relay approximations give
distributional derivatives as raised-label kernels, hence smoothness.
Lattice symmetries and positivity pass to the limit; scale estimates
give polynomial bounds.

\subsubsection{4. Steps}\label{steps-19}

\begin{enumerate}
\def\labelenumi{\arabic{enumi}.}
\tightlist
\item
  Under (H1), reflection positivity and repeated Cauchy--Schwarz give
  chessboard estimates; translation log-convexity transfers centered
  bounds to interior positions.
  \href{https://github.com/yang-mills-group/OpenAI-YM-20261028/blob/3456e5e75cf8d4cd558955a5c7638fba2d68f464/YM/Extraction/Around.lean\#L95-L129}{\texttt{Around.interior\_from}}.
\item
  Under (H1)--(H2), translated cells separate insertions and bound
  correlations on distinct-point compact sets.
  \href{https://github.com/yang-mills-group/OpenAI-YM-20261028/blob/3456e5e75cf8d4cd558955a5c7638fba2d68f464/YM/Extraction/Held.lean\#L120-L132}{\texttt{Around.size\_done}}.
\item
  Under the same hypotheses, slab traces and chessboard column moments
  give equicontinuity on generic compact sets through interpolation and
  Cauchy's estimate.
  \href{https://github.com/yang-mills-group/OpenAI-YM-20261028/blob/3456e5e75cf8d4cd558955a5c7638fba2d68f464/YM/Extraction/Equicontinuity.lean\#L181-L214}{\texttt{Difference.completed}}.
\item
  Two Arzelà--Ascoli extractions give IV.1.1(a)--(b); domination gives
  (c)--(d).
  \href{https://github.com/yang-mills-group/OpenAI-YM-20261028/blob/3456e5e75cf8d4cd558955a5c7638fba2d68f464/YM/Extraction/Lipschitz.lean\#L166-L184}{\texttt{first\_surface}},
  \href{https://github.com/yang-mills-group/OpenAI-YM-20261028/blob/3456e5e75cf8d4cd558955a5c7638fba2d68f464/YM/Extraction/Integrate.lean\#L115-L153}{\texttt{Order.ordinary}}.
\item
  Under IV.1.2's hypotheses, summation by parts identifies weak
  derivatives; small chessboard errors vanish. Normalized projected
  approximants and the dimension count give fixed derivative labels.
  \href{https://github.com/yang-mills-group/OpenAI-YM-20261028/blob/3456e5e75cf8d4cd558955a5c7638fba2d68f464/YM/Work/D03.lean\#L244-L249}{\texttt{Work.D03.answer}},
  \href{https://github.com/yang-mills-group/OpenAI-YM-20261028/blob/3456e5e75cf8d4cd558955a5c7638fba2d68f464/YM/Dash/Limits.lean\#L69-L111}{\texttt{Dash.brace}}.
\item
  Iterating these identities gives rapid localized Fourier decay and the
  smooth representative.
  \href{https://github.com/yang-mills-group/OpenAI-YM-20261028/blob/3456e5e75cf8d4cd558955a5c7638fba2d68f464/YM/Dash/Slew.lean\#L26-L58}{\texttt{Dash.awake}}.
\item
  Under IV.1.4's hypotheses, density passes symmetries; limiting Gram
  matrices pass positivity. Under IV.1.5's hypotheses, two scale choices
  give its bound. The derivative plan then promotes it to (E7).
  \href{https://github.com/yang-mills-group/OpenAI-YM-20261028/blob/3456e5e75cf8d4cd558955a5c7638fba2d68f464/YM/Extraction/Harvest.lean\#L46-L60}{\texttt{Order.realPositive}},
  \href{https://github.com/yang-mills-group/OpenAI-YM-20261028/blob/3456e5e75cf8d4cd558955a5c7638fba2d68f464/YM/Work/D05.lean\#L160-L180}{\texttt{Work.D05.answer}},
  \href{https://github.com/yang-mills-group/OpenAI-YM-20261028/blob/3456e5e75cf8d4cd558955a5c7638fba2d68f464/YM/Work/B10.lean\#L72-L90}{\texttt{Work.B10.capped\_tags}}.
\end{enumerate}

\subsubsection{5. Where the difficulty
lies}\label{where-the-difficulty-lies-19}

Uncentered roots and transfer-kernel equicontinuity require the main
estimates. Fixed derivative labels require normalized remainders and the
separate rank comparison. Equicontinuity holds on generic
configurations; smoothness extends to all distinct-point configurations.
The label maps are determined only modulo null labels, while their
quotient maps are intrinsic. Continuous rotation symmetry comes later.

\subsubsection{6. Notation}\label{notation-19}

\(\mathrm{Label}_d\) is the finite-dimensional complex label space;
\(F_d=\mathrm{Label}_d/\mathrm{Null}_d\), where null labels have zero
smeared kernels in every arity and slot against all spectator labels and
tests. \(\mathrm{Conf}_n\) means distinct points; \(\mathrm{Gen}_n\)
means separation in every coordinate, modulo \(L_p\) for
\(\mathrm{Gen}^{(p)}_n\). \(\mathcal D_n=C_c^\infty(\mathrm{Conf}_n)\);
admissibility means support inside \((-L_p/2,L_p/2)^{4n}\).
\(a_N\downarrow0\) and \(L_p\uparrow\infty\) are spacings and periods.
\(\mathcal C_{N,p}(v;r)=\langle\prod_iI_{N,p}(v_i,r_i)\rangle\);
smearing sums against \(t(\widehat r)\) with weight \(a_N^{4n}\), where
\(\widehat r\) is the physical site position. Lifted tuples are reduced
modulo \(L_p/a_N\) for evaluation. \(v_{i\to w}\) replaces slot \(i\) by
\(w\). Separation uses the supremum norm. \(\rho_d\) acts on labels and
their quotients. \(H_4\) consists of signed coordinate permutations,
with matrices \((s_{ab})\); \(s_a\) reverses axis \(a\); \(\Theta_a\)
reverses words, reflects points and replaces \(v\) by
\(\rho(s_a)\bar v\).

\subsubsection{7. What remains
unexplained}\label{what-remains-unexplained-19}

The upstream relay, rank and full fixed-volume kernel arguments remain
unexplained, as does the budget's relation to III.12. Extraction is
subsequential in both indices; III.13 supplies the full-spacing upgrade
(\href{https://github.com/yang-mills-group/OpenAI-YM-20261028/blob/3456e5e75cf8d4cd558955a5c7638fba2d68f464/YM/UVExtension/Reports/Ultraviolet.lean\#L62-L69}{\texttt{UVExtension.Reports.useFull}}).
Independence from the volume subsequence is not asserted.

Detailed notes: the full account of this unit, with proof sketches for
every step, is on the
\href{https://github.com/scottnarmstrong/yang-mills-informalization/blob/main/outline/notes/IV-1.md}{detailed
notes page}.

\subsection{V.1. Decomposition of the limit into
phases}\label{v.1.-decomposition-of-the-limit-into-phases}

\textbf{Part:} V, phases and reconstruction. \textbf{Lean:}
\href{https://github.com/yang-mills-group/OpenAI-YM-20261028/tree/3456e5e75cf8d4cd558955a5c7638fba2d68f464/YM/OS}{\texttt{YM/OS}},
\href{https://github.com/yang-mills-group/OpenAI-YM-20261028/tree/3456e5e75cf8d4cd558955a5c7638fba2d68f464/YM/Phase}{\texttt{YM/Phase}},
\href{https://github.com/yang-mills-group/OpenAI-YM-20261028/tree/3456e5e75cf8d4cd558955a5c7638fba2d68f464/YM/Ground}{\texttt{YM/Ground}},
\href{https://github.com/yang-mills-group/OpenAI-YM-20261028/tree/3456e5e75cf8d4cd558955a5c7638fba2d68f464/YM/Heat}{\texttt{YM/Heat}},
\href{https://github.com/yang-mills-group/OpenAI-YM-20261028/tree/3456e5e75cf8d4cd558955a5c7638fba2d68f464/YM/SNAG}{\texttt{YM/SNAG}},
\href{https://github.com/yang-mills-group/OpenAI-YM-20261028/tree/3456e5e75cf8d4cd558955a5c7638fba2d68f464/YM/Paley}{\texttt{YM/Paley}},
\href{https://github.com/yang-mills-group/OpenAI-YM-20261028/tree/3456e5e75cf8d4cd558955a5c7638fba2d68f464/YM/Spectral}{\texttt{YM/Spectral}},
the work items B01--B18 in
\href{https://github.com/yang-mills-group/OpenAI-YM-20261028/tree/3456e5e75cf8d4cd558955a5c7638fba2d68f464/YM/Work}{\texttt{YM/Work}},
and the assembly files
\href{https://github.com/yang-mills-group/OpenAI-YM-20261028/blob/3456e5e75cf8d4cd558955a5c7638fba2d68f464/YM/Bid/Phase.lean}{\texttt{Bid/Phase.lean}},
\href{https://github.com/yang-mills-group/OpenAI-YM-20261028/blob/3456e5e75cf8d4cd558955a5c7638fba2d68f464/YM/Bid/Fetch.lean}{\texttt{Bid/Fetch.lean}},
\href{https://github.com/yang-mills-group/OpenAI-YM-20261028/blob/3456e5e75cf8d4cd558955a5c7638fba2d68f464/YM/Bid/Market.lean}{\texttt{Bid/Market.lean}},
\href{https://github.com/yang-mills-group/OpenAI-YM-20261028/blob/3456e5e75cf8d4cd558955a5c7638fba2d68f464/YM/Org/Phase.lean}{\texttt{Org/Phase.lean}}.
\textbf{Size:} about 5,700 lines in declarations used by the main
theorem, including about 2,700 in B01--B18, plus the assembly files.

\subsubsection{1. What this unit proves}\label{what-this-unit-proves-20}

\(S\) is the limit table; words are separated unless stated otherwise.
``Hierarchy'' means permutation symmetry, normalization and unit labels,
signed-axis symmetries and reality, translation invariance, point and
smeared reflection positivity for every coordinate axis, smoothness on
separated configurations, and caps on all derivatives (IV.1). A cap
bounds each fixed correlation derivative by \(Cr^{-A}\) for \(0<r\le1\)
and mutual sup-distance \(\ge r\), with \(C\ge0\), \(A\in\mathbb N\)
depending on arity, degrees, labels and derivative order. The OS inner
product is \(\langle b,b'\rangle=S(\Theta b\cdot b')\), and
\(\Omega=|\emptyset\rangle\).

\textbf{Proposition V.1.1 (chambers and ground space).} For every \(S\)
satisfying the hierarchy and coordinate axis \(a\), its OS space carries
a strongly continuous self-adjoint contraction semigroup and a commuting
strongly continuous spatial unitary group,
\(E_t|b\rangle=|b+te_a\rangle\), \(U_y|b\rangle=|b+y\rangle\)
(\(t\ge0\), \(y\perp e_a\)). These identities hold for every
\(a\)-positive \(b\). This is a chamber. The fixed space
\(\mathcal H_0\) of all \(E_t\) contains \(\Omega\), and
\(E_t\psi\to P_0\psi\) for every \(\psi\), where \(P_0\) projects onto
\(\mathcal H_0\). \emph{Status:} theorem. \emph{Lean:}
\href{https://github.com/yang-mills-group/OpenAI-YM-20261028/blob/3456e5e75cf8d4cd558955a5c7638fba2d68f464/YM/OS/Dynamics.lean\#L63-L104}{\texttt{OS.implement}},
\href{https://github.com/yang-mills-group/OpenAI-YM-20261028/blob/3456e5e75cf8d4cd558955a5c7638fba2d68f464/YM/OS/Core.lean\#L38-L78}{\texttt{OS.Semigroup.emerges}}.

\textbf{Proposition V.1.2 (spectral calculus).} In every chamber,
writing \(E_t=e^{-tH}\), there is a strongly continuous unitary group
determined by \(e^{isH}E_1\psi=E_1^{1-is}\psi\) through continuous
functional calculus. The strongly continuous unitary representation
\(V(x)=e^{ix^aH}U_{x-x^ae_a}\) has a projection-valued measure \(P\);
for every \(\psi\), \(\mu_\psi(B)=\langle\psi,P(B)\psi\rangle\) is
supported in \(p^a\ge0\), and, for every \(x\) and
\(\mathrm{Im}\,z\ge0\),
\[\langle\psi,V(x)e^{i(\mathrm{Re}\,z)H}E_{\mathrm{Im}\,z}\psi\rangle=\int e^{ip\cdot x}e^{ip^az}\,d\mu_\psi(p).\]
\emph{Status:} theorem. \emph{Lean:}
\href{https://github.com/yang-mills-group/OpenAI-YM-20261028/blob/3456e5e75cf8d4cd558955a5c7638fba2d68f464/YM/Heat/Lift.lean\#L50-L58}{\texttt{Semigroup.unit}},
\href{https://github.com/yang-mills-group/OpenAI-YM-20261028/blob/3456e5e75cf8d4cd558955a5c7638fba2d68f464/YM/SNAG/Escape.lean\#L52-L119}{\texttt{Spectral.Orbit.total}},
\href{https://github.com/yang-mills-group/OpenAI-YM-20261028/blob/3456e5e75cf8d4cd558955a5c7638fba2d68f464/YM/Heat/Lift.lean\#L108-L114}{\texttt{Semigroup.heat}}.

\textbf{Proposition V.1.3 (slab operators).} Under the hierarchy and
bench V.1.7, every time-axis chamber has unique bounded \(T_d(w)\) for
\(d>0\), \(w\) in \((0,d)\), satisfying
\(T_d(w)|b\rangle=|w\cdot\tau_db\rangle\) for every positive \(b\). For
every \(D\in\mathbb N\), \(\rho>0\), some \(C>0\),
\(\kappa\in\mathbb N\) give
\[\|T_d(w)\|\le C^{|w|}(|w|/r)^{\kappa|w|} \tag{B}\] for all \(d>0\),
\(0<r\le1\), and \(r\)-buffered \(w\) with degrees \(\le D\), label
norms \(\le\rho\); the empty-word bound is \(1\). With the plan's
derivative identities, for fixed \(d\) and labels, \(T_d(w(x))\) is
norm-\(C^\infty\) on the open slab configuration set, and coordinate
differentiation replaces the corresponding label by \(\partial_jv\).
More generally this smoothness holds for any slab family satisfying (B)
on a time-axis chamber of \(S'\), provided \(S,S'\) satisfy the
hierarchy, \(S'\) respects \(S\), and \(S'\) satisfies the derivative
identities of the plan on \(S\). \emph{Status:} theorem conditional on
the bench. \emph{Lean:}
\href{https://github.com/yang-mills-group/OpenAI-YM-20261028/blob/3456e5e75cf8d4cd558955a5c7638fba2d68f464/YM/Work/B01.lean\#L50-L71}{\texttt{Work.B01.answer}},
\href{https://github.com/yang-mills-group/OpenAI-YM-20261028/blob/3456e5e75cf8d4cd558955a5c7638fba2d68f464/YM/Work/B02.lean\#L17-L30}{\texttt{Work.B02.answer}}.

\textbf{Proposition V.1.4 (compressed algebra).} Assume the hierarchy,
\(O(4)\)-covariance, a time-axis chamber and a slab family satisfying
V.1.3's insertion identity. Then spatial translations fix
\(\mathcal H_0\) pointwise. For every separated \(w\),
\(M(w)=P_0T_d(w+s)|_{\mathcal H_0}\) is independent of admissible
\(s,d\); \(M(b)\Omega=P_0|b\rangle\) for positive \(b\), and
\(\langle\Omega,M(w)\Omega\rangle=S(w)\). For any finite collection of
separated words \(q_i\) and shifts \(a_i(n)\) with
\(t(a_j(n)-a_i(n))\to+\infty\) for \(i<j\), the concatenations are
eventually separated and their correlations tend to
\(\langle\Omega,\prod_iM(q_i)\Omega\rangle\), regardless of spatial
shifts. The generated norm-closed \(*\)-algebra \(\mathcal A\) is
commutative; \(\Omega\) is cyclic and separating. With smooth slab laws,
\(M\) is permutation/translation invariant, satisfies
\(M(\emptyset)=1\), permits dropping unit labels together with their
points, satisfies \(M(w)^*=M(\Theta w)\), is multilinear through
products of \(F_d(S)\) and norm-smooth in positions; \(\mathcal A\) is
separable. \emph{Status:} theorem. \emph{Lean:}
\href{https://github.com/yang-mills-group/OpenAI-YM-20261028/blob/3456e5e75cf8d4cd558955a5c7638fba2d68f464/YM/Ground/Permute.lean\#L149-L161}{\texttt{Ground.adjunction}},
\href{https://github.com/yang-mills-group/OpenAI-YM-20261028/blob/3456e5e75cf8d4cd558955a5c7638fba2d68f464/YM/Work/B03.lean\#L405-L415}{\texttt{Work.B03.answer}}.

\textbf{Proposition V.1.5 (decomposition).} Under V.1.4 including
smoothness, \(X=\mathrm{Spec}(\mathcal A)\) is compact metrizable, with
full-support Borel probability \(\nu\). Its points bijectively
parametrize characters \(\chi_\xi\); every evaluation is continuous and
\(\int_X\chi_\xi(p)d\nu=\langle\Omega,p\Omega\rangle\) for all
\(p\in\mathcal A\). Each \(S^\xi(w)=\chi_\xi(M(w))\), extended by zero
off separated configurations, is a table; evaluations are continuous in
\(\xi\), and \(S(w)=\int_XS^\xi(w)d\nu\) for every separated \(w\).
\emph{Status:} theorem. \emph{Lean:}
\href{https://github.com/yang-mills-group/OpenAI-YM-20261028/blob/3456e5e75cf8d4cd558955a5c7638fba2d68f464/YM/Ground/Spectrum.lean\#L126-L140}{\texttt{Ground.resolution}},
\href{https://github.com/yang-mills-group/OpenAI-YM-20261028/blob/3456e5e75cf8d4cd558955a5c7638fba2d68f464/YM/Phase/Mixture.lean\#L113-L117}{\texttt{Resolution.barycenter}}.

\textbf{Proposition V.1.6 (phase laws).} Assume V.1.3--V.1.5 with smooth
laws for the field-algebra plan, the Seed, grid hypothesis and
field-algebra laws on \(S\). For every separated word \(w\),
\(M(w)=M(\vartheta w)\) with arbitrary complex labels, and
\(M(w)=M(Kw)\) for every \(K\in O(4)\). For every \(\xi\in X\),
\(S^\xi\) satisfies the hierarchy, respects \(S\), is \(O(4)\)-covariant
for the same action, inclusion/derivative identities and stress-tensor
Ward identities, and \(N_d(S^\xi)=N_d(S)\) for every degree \(d\). For
every \(D\in\mathbb N\), \(\rho>0\), some \(K>0\),
\(\kappa\in\mathbb N\) give \(|S^\xi(w)|\le K^{|w|}(|w|/r)^{\kappa|w|}\)
uniformly in \(\xi\), \(0<r\le1\), and words with degrees \(\le D\),
label norms \(\le\rho\), mutual sup-distance \(\ge r\). \emph{Status:}
theorem. \emph{Lean:}
\href{https://github.com/yang-mills-group/OpenAI-YM-20261028/blob/3456e5e75cf8d4cd558955a5c7638fba2d68f464/YM/Work/B05.lean\#L90-L115}{\texttt{Work.B05.answer}},
\href{https://github.com/yang-mills-group/OpenAI-YM-20261028/blob/3456e5e75cf8d4cd558955a5c7638fba2d68f464/YM/Phase/Picks.lean\#L128-L158}{\texttt{Retain.lift\_detect}}.

\textbf{Proposition V.1.7 (bench).} For every \(D\in\mathbb N\),
\(\rho>0\), some \(C_0>0\), \(\kappa\in\mathbb N\) work for every
\(d>0\), \(0<r\le1\), nonempty \(r\)-buffered \(w\) with degrees
\(\le D\), norms \(\le\rho\): there are a lattice level \(b\) and
\(s\le d\), \(s\in a_b\mathbb N\), with \(w\) in \((0,s)\), such that
every positive \(v\) has \(B_v\) satisfying
\[|S(\Theta v_n\cdot v_n)|\le B_vA^{4n}\quad(n\ge0),\qquad A=C_0^{|w|}(|w|/r)^{\kappa|w|},\]
where \(v_0=v\),
\(v_{n+1}=\tau_s(\Theta w)\cdot\tau_s(w\cdot\tau_sv_n)\). \emph{Status:}
hypothesis here, proved in VII.4 for every Seed table. \emph{Lean:}
\href{https://github.com/yang-mills-group/OpenAI-YM-20261028/blob/3456e5e75cf8d4cd558955a5c7638fba2d68f464/YM/Next/PHASE/Seats.lean\#L154-L160}{\texttt{Next.Phase.seats}}.

\textbf{Proposition V.1.8 (vacuum uniqueness).} Given V.1.5 and the
hierarchy for every phase, fix a time-axis chamber for each \(\xi\).
Then \(\mathcal H_0^\xi=\mathbb C\Omega_\xi\) for \(\nu\)-almost every
\(\xi\), so such phases exist. Every time-axis phase chamber has slab
operators with \(\|T_d^\xi(w)\|\le\|T_d(w)\|\), inheriting (B) whenever
the original family has it, and, under V.1.6, smoothness. \emph{Status:}
theorem. \emph{Lean:}
\href{https://github.com/yang-mills-group/OpenAI-YM-20261028/blob/3456e5e75cf8d4cd558955a5c7638fba2d68f464/YM/Phase/Unique.lean\#L91-L110}{\texttt{Cluster.lifts}},
\href{https://github.com/yang-mills-group/OpenAI-YM-20261028/blob/3456e5e75cf8d4cd558955a5c7638fba2d68f464/YM/Work/B13.lean\#L91-L105}{\texttt{Work.B13.answer}}.

\textbf{Proposition V.1.9 (selected-phase certificate).} For every
Wilson family with extraction order and limit \(S\), assuming its Seed,
grid hypothesis and field algebra with laws, there exist a time-axis
chamber of \(S\) with spatially fixed ground space, a smooth slab family
satisfying (B), decomposition as V.1.5 with every phase satisfying
V.1.6's hierarchy, field laws and null-space equality, and a selected
phase with a time-axis chamber having ground space \(\mathbb C\Omega\)
and smooth slab family. \emph{Status:} theorem. \emph{Lean:}
\href{https://github.com/yang-mills-group/OpenAI-YM-20261028/blob/3456e5e75cf8d4cd558955a5c7638fba2d68f464/YM/Org/Phase.lean\#L60-L76}{\texttt{Phase.Loom.cleanse}}.

\subsubsection{2. What it uses}\label{what-it-uses-20}

\begin{itemize}
\tightlist
\item
  IV.1: hierarchy and
  \href{https://github.com/yang-mills-group/OpenAI-YM-20261028/blob/3456e5e75cf8d4cd558955a5c7638fba2d68f464/YM/Spec/Closed.lean\#L20-L26}{\texttt{Spec.Seed}},
  including cell bounds and corrected-field data.
\item
  VII.4: bench; III.11:
  \href{https://github.com/yang-mills-group/OpenAI-YM-20261028/blob/3456e5e75cf8d4cd558955a5c7638fba2d68f464/YM/Org/Input.lean\#L121-L123}{\texttt{Org.Grill}},
  the grid hypothesis with cutoff-independent bounded approximants.
\item
  VI.1:
  \href{https://github.com/yang-mills-group/OpenAI-YM-20261028/blob/3456e5e75cf8d4cd558955a5c7638fba2d68f464/YM/Spec/Closed.lean\#L28-L31}{\texttt{Spec.Algebra}}
  and field laws; V.2:
  \href{https://github.com/yang-mills-group/OpenAI-YM-20261028/blob/3456e5e75cf8d4cd558955a5c7638fba2d68f464/YM/Hilbert/Kolmogorov.lean\#L34-L38}{\texttt{Kernel.Gram}}
  completion; I.1, II.1: lattice and Wilson family.
\item
  OS doubling; chessboard estimates; continuous functional calculus;
  SNAG; Paley--Wiener; Gelfand--Naimark; Riesz--Markov; dominated
  convergence; Cartan--Dieudonné.
\end{itemize}

\subsubsection{3. Strategy}\label{strategy-20}

Compress bounded slab insertions to the ground space. Their commutative
algebra describes infinite-time-separated observables; characters yield
phase tables and the vacuum state supplies their mixing measure.
Operator identities give laws in every phase; an integrated variance
identity selects a unique-vacuum phase almost everywhere.

\subsubsection{4. Steps}\label{steps-20}

\begin{enumerate}
\def\labelenumi{\arabic{enumi}.}
\tightlist
\item
  Construct contractions by OS doubling; functional calculus and SNAG
  give spectral measures, with positive energy from half-plane
  analyticity.
  \href{https://github.com/yang-mills-group/OpenAI-YM-20261028/blob/3456e5e75cf8d4cd558955a5c7638fba2d68f464/YM/OS/Flow.lean\#L70-L89}{\texttt{OS.contract}},
  \href{https://github.com/yang-mills-group/OpenAI-YM-20261028/blob/3456e5e75cf8d4cd558955a5c7638fba2d68f464/YM/Heat/Lift.lean\#L102-L104}{\texttt{Semigroup.nonneg}}.
\item
  Turn bench train bounds into slab norms by the power argument; absorb
  leftover width with \(E_{d-s}\).
  \href{https://github.com/yang-mills-group/OpenAI-YM-20261028/blob/3456e5e75cf8d4cd558955a5c7638fba2d68f464/YM/OS/Train.lean\#L20-L60}{\texttt{Drive.against}}.
\item
  Rotate spatial displacements into time to fix ground vectors; exchange
  middle clusters in a four-cluster limit to prove commutativity.
  \href{https://github.com/yang-mills-group/OpenAI-YM-20261028/blob/3456e5e75cf8d4cd558955a5c7638fba2d68f464/YM/Ground/Lean.lean\#L185-L188}{\texttt{Ground.inert}},
  \href{https://github.com/yang-mills-group/OpenAI-YM-20261028/blob/3456e5e75cf8d4cd558955a5c7638fba2d68f464/YM/Ground/Permute.lean\#L72-L93}{\texttt{Ground.exchanges}}.
\item
  Gelfand--Riesz gives a compact metrizable spectrum and the vacuum
  measure; separation of \(\Omega\) gives full support.
  \href{https://github.com/yang-mills-group/OpenAI-YM-20261028/blob/3456e5e75cf8d4cd558955a5c7638fba2d68f464/YM/Ground/Spectrum.lean\#L126-L140}{\texttt{Ground.resolution}}.
\item
  For real labels, \(D=M(w)-M(\vartheta w)\) is skew-adjoint. Positivity
  of variances of many bounded lattice approximants forces
  \(-\|D\Omega\|^2\ge0\); separation gives \(D=0\).
  \href{https://github.com/yang-mills-group/OpenAI-YM-20261028/blob/3456e5e75cf8d4cd558955a5c7638fba2d68f464/YM/Work/B05.lean\#L90-L115}{\texttt{Work.B05.answer}}.
\item
  Transfer positivity and Ward identities before applying characters.
  Corrected-field detector moments, full support and finite-dimensional
  injectivity preserve null spaces.
  \href{https://github.com/yang-mills-group/OpenAI-YM-20261028/blob/3456e5e75cf8d4cd558955a5c7638fba2d68f464/YM/Work/B08.lean\#L206-L209}{\texttt{Work.B08.answer}},
  \href{https://github.com/yang-mills-group/OpenAI-YM-20261028/blob/3456e5e75cf8d4cd558955a5c7638fba2d68f464/YM/Work/B09.lean\#L12-L15}{\texttt{Work.B09.answer}},
  \href{https://github.com/yang-mills-group/OpenAI-YM-20261028/blob/3456e5e75cf8d4cd558955a5c7638fba2d68f464/YM/Phase/Picks.lean\#L128-L158}{\texttt{Retain.lift\_detect}}.
\item
  Dominated convergence makes \(\|P_0^\xi|b\rangle\|^2-|S^\xi(b)|^2\)
  have integral zero; one common countable total family gives uniqueness
  almost everywhere.
  \href{https://github.com/yang-mills-group/OpenAI-YM-20261028/blob/3456e5e75cf8d4cd558955a5c7638fba2d68f464/YM/Phase/Unique.lean\#L56-L86}{\texttt{Cluster.exp\_real}},
  \href{https://github.com/yang-mills-group/OpenAI-YM-20261028/blob/3456e5e75cf8d4cd558955a5c7638fba2d68f464/YM/Work/B14.lean\#L119-L143}{\texttt{Work.B14.universal}}.
\end{enumerate}

\subsubsection{5. Where the difficulty
lies}\label{where-the-difficulty-lies-20}

Phase reflection needs lattice variance positivity; null-space retention
needs quantitative corrected-field estimates. Neither follows merely
from the abstract character decomposition. No clustering, unique vacuum
or spectral gap is assumed.

\subsubsection{6. Notation}\label{notation-20}

\(t(x)\): last coordinate; \(e_a\): coordinate unit vector;
\(v\in V_d\): degree-\(d\) label. Separated means distinct points.
\(|w|\): word length; \(w\cdot w'\): concatenation; \(\tau_s\): time
translation; \(\Theta\): antilinear OS reflection with reversal;
\(\vartheta\): linear time reflection. Positive words have positive
times; \(r\)-buffered means mutual sup-distance \(\ge r\) and times in
\([r,d-r]\). \(N_d(S)\) comprises labels vanishing against every
separated spectator family; \(F_d(S)=V_d/N_d(S)\). A table \(S'\)
respects \(S\) when \(N_d(S)\subset N_d(S')\) for every \(d\).

\subsubsection{7. What remains
unexplained}\label{what-remains-unexplained-20}

Ultraviolet packing, bounded approximants and corrected-field estimates
remain external inputs. With copy number and approximation index fixed,
the variance argument takes cutoff limits before large translations,
then sends copy number to infinity and approximation error to zero.
Vacuum uniqueness for every phase belongs to V.3; selection here is
noncanonical.

Detailed notes: the full account of this unit, with proof sketches for
every step, is on the
\href{https://github.com/scottnarmstrong/yang-mills-informalization/blob/main/outline/notes/V-1.md}{detailed
notes page}.

\subsection{V.2. Osterwalder--Schrader reconstruction of a Wightman
theory}\label{v.2.-osterwalderschrader-reconstruction-of-a-wightman-theory}

\textbf{Part:} V, phases and reconstruction. \textbf{Lean:}
\href{https://github.com/yang-mills-group/OpenAI-YM-20261028/tree/3456e5e75cf8d4cd558955a5c7638fba2d68f464/YM/Work}{A16,
C01--C29}, Boundary, Fields, Shift. \textbf{Size:} 6,456 lines; 705
declarations.

\subsubsection{1. What this unit proves}\label{what-this-unit-proves-21}

All arities include zero; statements hold for every label tuple and
Schwartz test unless restricted. Tables \(S,T\) share a register; \(T\)
respects \(S\) means \(N_d(S)\subseteq N_d(T)\) for every degree, where
\(N_d\) consists of labels vanishing with arbitrary insertions. Write
(H) for: both Euclidean hierarchies; respect; differential-plan
inclusions compatible with both tables and derivatives with \(T\); a
spread satisfying the slab bound, smoothness in insertion positions on
every valid slab, and insertion derivative laws. The slab bound is
\(\|A_d(w)\|\le C^{|w|}(|w|/r)^{a|w|}\) for buffers and pairwise
maximum-norm distances at least \(r\in(0,1]\), with \(C>0\),
\(a\in\mathbb N\) uniform in width and length for fixed degree and
representative-norm bounds.

(SC) means continuity in strict complex times and spatial positions,
holomorphy in times, and agreement with \(T\) at ordered imaginary
times. (PG) means \(|F_n^y(x)(v)|\le Ck^{-m}\) whenever all height gaps
are at least \(k\in(0,1]\), with \(C\ge0\), \(m\in\mathbb N\) depending
on \(n,v\). Put \(\psi_v(y;f)=\int F_n^y(x)(v)f(x)\,dx\); (I), (HC),
(CR) mean integrability, continuity on \(E_n\) and
\(\partial_a\psi(y;f)=\psi(y;-i\partial_{\underline a}f)\) for every
real height direction \(a\) within \(E_n\). (BV) means
\(\psi_v(y;f)\to W_n(v)(f)\) as \(y\to0\) through ordered heights,
without a fixed direction.

\textbf{Proposition V.2.1.} A16 and C01--C29 hold simultaneously, each
with its own hypotheses. \emph{Status:} theorem. \emph{Lean:}
\href{https://github.com/yang-mills-group/OpenAI-YM-20261028/blob/3456e5e75cf8d4cd558955a5c7638fba2d68f464/YM/Integrated.lean\#L182-L212}{reconstructionArcs}.

\textbf{Proposition V.2.2.} Under (H), linear bounded-operator
insertions \(B_p\), \(p\ge0\), exist, given for \(p>0\) by the central
one-point slab operator of width \(2p\), with \(\|B_p(u)\|\le Cp^{-a}\)
for \(0<p<1\), for each \(u\), with \(C>0\), \(a\in\mathbb N\). The
empty product is \(\|\Omega\|^2\). Their vacuum products, separated by
\(e^{-(\mathrm{Im}\,w-2p)H}e^{i\mathrm{Re}\,wH}U(\vec a)\), are
independent of widths \(0<p<1\) with gaps \(>4p\), agree with \(T\) at
imaginary times, and define \(F\) satisfying (SC), (PG). \emph{Status:}
theorem. \emph{Lean:}
\href{https://github.com/yang-mills-group/OpenAI-YM-20261028/blob/3456e5e75cf8d4cd558955a5c7638fba2d68f464/YM/Work/C01.lean\#L99-L117}{C01},
\href{https://github.com/yang-mills-group/OpenAI-YM-20261028/blob/3456e5e75cf8d4cd558955a5c7638fba2d68f464/YM/Work/C02.lean\#L168-L177}{C02},
\href{https://github.com/yang-mills-group/OpenAI-YM-20261028/blob/3456e5e75cf8d4cd558955a5c7638fba2d68f464/YM/Work/C03.lean\#L245-L248}{C03}.

\textbf{Proposition V.2.3.} (SC), (PG) imply (I), (HC), (CR).
Separately, (I), (HC), (CR), (PG) produce a multilinear tempered
boundary family \(W\) satisfying (BV), hence every ordered radial limit.
\emph{Status:} theorem. \emph{Lean:}
\href{https://github.com/yang-mills-group/OpenAI-YM-20261028/blob/3456e5e75cf8d4cd558955a5c7638fba2d68f464/YM/Work/C04.lean\#L221-L222}{C04},
\href{https://github.com/yang-mills-group/OpenAI-YM-20261028/blob/3456e5e75cf8d4cd558955a5c7638fba2d68f464/YM/Shift/Osc.lean\#L83-L86}{jointly}.

\textbf{Proposition V.2.4.} Under (SC), (PG), (CR), (HC), (BV), for
ordered \(y\) and compactly supported \(f\),
\(W_n(v)(\hat f)=\psi_v(y;\widehat{e^{\langle\underline y,\cdot\rangle}f})\)
(L). Under (SC), (PG), (I), (L), Fourier support lies in
\(\mathcal P_n=\{\sum_kp_k^0=0,\ \sum_{k\ge j}p_k^0\ge0\ \forall j\}\).
With Lorentz covariance it lies in
\(\mathcal V_n=\bigcap_\Lambda\tilde\Lambda^{-1}\mathcal P_n\); thus
total momentum vanishes and suffix momenta belong to \(\bar V_+\).
\emph{Status:} theorem. \emph{Lean:}
\href{https://github.com/yang-mills-group/OpenAI-YM-20261028/blob/3456e5e75cf8d4cd558955a5c7638fba2d68f464/YM/Work/C05.lean\#L205-L230}{C05},
\href{https://github.com/yang-mills-group/OpenAI-YM-20261028/blob/3456e5e75cf8d4cd558955a5c7638fba2d68f464/YM/Work/C06.lean\#L136-L157}{C06},
\href{https://github.com/yang-mills-group/OpenAI-YM-20261028/blob/3456e5e75cf8d4cd558955a5c7638fba2d68f464/YM/Shift/Closing.lean\#L106-L121}{tents}.

\textbf{Proposition V.2.5.} Unconditionally, spatial rotations and
coordinate boosts generate \(\mathcal L^\uparrow_+\); \(g\mapsto g_*f\)
is Schwartz-continuous and \(t^{-1}(\Lambda_a(t)_*f-f)\to-D_af\), where
\(D_a=\sum_k(x_k^a\partial_{x_k^0}+x_k^0\partial_{x_k^a})\). Assume both
hierarchies, inclusions compatible with \(S\), respect, Euclidean
covariance, rotation laws and spin compatibility. Under (SC), (PG), (I),
(CR), (HC), infinitesimal rotations give
\(\sum_k\psi_{v^{(k)}}(y;f)=\psi_v(y;D_{a,y}f)\), for \(y\in E_n\) and
each spatial axis \(a\), where
\(D_{a,y}=D_a+i\sum_ky_k\partial_{x_k^a}\) and \(v^{(k)}\) replaces
\(v_k\) by \(-i\rho(e_0\wedge e_a)v_k\). Under (SC), (PG), that identity
and (BV), \(W_n(K(\Lambda)v)(g_*f)=W_n(v)(f)\) for every Poincaré
transformation \(g=(a,\Lambda)\). With \(S\) symmetric, inclusions
compatible with \(S\), spin compatibility and induced star \(J\), this
covariance makes the word Gram form invariant and word matrix elements
continuous. \emph{Status:} theorem. \emph{Lean:}
\href{https://github.com/yang-mills-group/OpenAI-YM-20261028/blob/3456e5e75cf8d4cd558955a5c7638fba2d68f464/YM/Work/C07.lean\#L53-L66}{C07},
\href{https://github.com/yang-mills-group/OpenAI-YM-20261028/blob/3456e5e75cf8d4cd558955a5c7638fba2d68f464/YM/Work/C09.lean\#L403-L442}{C09},
\href{https://github.com/yang-mills-group/OpenAI-YM-20261028/blob/3456e5e75cf8d4cd558955a5c7638fba2d68f464/YM/Work/C28.lean\#L134-L143}{C28}.

\textbf{Proposition V.2.6.} For \(p\in\mathcal V_n\),
\(\langle p,\eta\rangle\ge0\) on the closed imaginary cone, and locally
\(\langle p,\eta\rangle\ge c|p|\) inside, for some \(c>0\). Smooth
cutoffs \(\chi_n\), equal to one on the \(1\)-neighbourhood of
\(\mathcal V_n\), supported in its open \(3\)-neighbourhood and with
polynomially bounded derivatives, exist. For \(\eta\) in the closed
imaginary cone, \(e^{-\langle\eta,\cdot\rangle}\chi_n\) has temperate
growth and its product with each test varies continuously in Schwartz
space. The function
\(\mathcal W_n(v)(z)=(2\pi)^{-4n}(\mathcal FW_n(v))(e^{i\langle\cdot,z\rangle}\chi_n)\)
is holomorphic on \(\mathcal T_n\). For interior \(\eta\), its integrals
represent
\(W_n(v)[\eta](https://github.com/scottnarmstrong/yang-mills-informalization/blob/main/outline/units/f)=(\mathcal FW_n(v))(e^{-\langle\eta,\cdot\rangle}\chi_n\check f)\).
Assuming damping, cone support, holomorphy, this integral
representation, (SC), (PG), (I), (L), it equals \(F_n^y\) at
\(x+i\underline y\) (TC). Cone support and covariance imply
\(\mathcal W_n(K(\Lambda)v)(\Lambda z)=\mathcal W_n(v)(z)\).
\emph{Status:} theorem. \emph{Lean:}
\href{https://github.com/yang-mills-group/OpenAI-YM-20261028/blob/3456e5e75cf8d4cd558955a5c7638fba2d68f464/YM/Work/C10.lean\#L71}{C10},
\href{https://github.com/yang-mills-group/OpenAI-YM-20261028/blob/3456e5e75cf8d4cd558955a5c7638fba2d68f464/YM/Work/C15.lean\#L120-L131}{C15},
\href{https://github.com/yang-mills-group/OpenAI-YM-20261028/blob/3456e5e75cf8d4cd558955a5c7638fba2d68f464/YM/Work/C16.lean\#L75-L91}{C16}.

\textbf{Proposition V.2.7.} Assume (SC), (TC), cone support, damping,
spin compatibility, tube covariance, symmetric \(S,T\), \(T\) smooth off
coincidences, respect and inclusions compatible with \(S\). There are
spectator profiles approaching zero, equal at the adjacent pair and
otherwise ordered, each admitting imaginary-boost paths into both tubes.
Spin boosts extend entirely; the two resulting holomorphic germs agree
on the zero-real-time slice and locally represent partial boundary
values. Then for every adjacent swap \(\sigma\),
\(W_n(v)(f)=W_n(v\circ\sigma)(f\circ\sigma)\) whenever
\(\mathrm{supp}f\subset\{x:x_{i+1}-x_i\text{ spacelike}\}\). This joint
locality implies exchange of adjacent letters in any word whose two test
supports are spacelike separated. \emph{Status:} theorem. \emph{Lean:}
\href{https://github.com/yang-mills-group/OpenAI-YM-20261028/blob/3456e5e75cf8d4cd558955a5c7638fba2d68f464/YM/Shift/Exchange.lean\#L92-L112}{actor},
\href{https://github.com/yang-mills-group/OpenAI-YM-20261028/blob/3456e5e75cf8d4cd558955a5c7638fba2d68f464/YM/Work/C27.lean\#L87-L92}{C27}.

\textbf{Proposition V.2.8.} Heated vectors are the insertion products of
V.2.2 acting on \(\Omega\), multiplied on the left by
\(e^{-(y_1-p)H}e^{ir_1^0H}U(\vec r_1)\) at heights \(y\) and real points
\(r\); the empty vector is \(\Omega\). Under (H), V.2.2 and oriented
\(P\), heated OS vectors have inner product
\(F^{(-x_n,\ldots,-x_1,y_1,\ldots,y_m)}(r_n,\ldots,r_1,q_1,\ldots,q_m)(Jv_n,\ldots,Jv_1,w_1,\ldots,w_m)\)
for positive ordered heights \(x,y\) and widths in \((0,1)\), with
heights and gaps exceeding four times their respective widths. Under
(H), slab insertions and spin compatibility, their smeared vectors
\(\Phi_t(b)\), using heights \((t,2t,\ldots,nt)\) and width
\(\min(t,1)/8\), have Bochner-integrable defining integrands and satisfy
\(\Phi_t(a\cdot b)=U(a)\Phi_t(b)\) for every \(t>0\), word \(b\) and
translation \(a\). Under (H), (SC), (BV), these properties and the
pointwise Gram identity,
\(\langle\Phi_s(b),\Phi_{s'}(c)\rangle\to W(b^*c)\) independently as
\(s,s'\downarrow0\). Tensor products are dense in Schwartz space. With
this density, equal null spaces, both hierarchies, inclusions compatible
with both tables, (SC), (TC), vanishing of \(W(a(u,f)c)\) for all
\(a,c,f\) implies \(u=0\). Every boundary family \(W\) and antilinear
\(J\) give bilinearity of \((u,f)\mapsto W(a(u,f)c)\) and continuity of
\((f,g)\mapsto\langle(u,f)a,(u\prime,g)c\rangle_W\). \emph{Status:}
theorem. \emph{Lean:}
\href{https://github.com/yang-mills-group/OpenAI-YM-20261028/blob/3456e5e75cf8d4cd558955a5c7638fba2d68f464/YM/Work/C24.lean\#L146-L153}{C24},
\href{https://github.com/yang-mills-group/OpenAI-YM-20261028/blob/3456e5e75cf8d4cd558955a5c7638fba2d68f464/YM/Work/C26.lean\#L77-L92}{C26}.

\textbf{Proposition V.2.9.} A realization comprises spin-compatible
\(K\), induced \(J\), fields bilinear in labels and tests on a dense
invariant domain equal to the span of vacuum word vectors, a unit vacuum
fixed by a strongly continuous unitary Poincaré representation,
covariance \(U(g)\varphi(u,f)=\varphi(K(\Lambda)u,g_*f)U(g)\), the
identity
\(\langle\varphi(u,f)\psi,\chi\rangle=\langle\psi,\varphi(Ju,\bar f)\chi\rangle\)
on the domain, continuous smeared-field vectors, local commutativity,
and faithfulness: \(u\ne0\) implies \(\varphi(u,f)\ne0\) for some \(f\).
It has vacuum expectations \(W\), translation spectral measure supported
in \(\bar V_+\), and only scalar multiples of the vacuum translation
invariant. \emph{Status:} definition. \emph{Lean:}
\href{https://github.com/yang-mills-group/OpenAI-YM-20261028/blob/3456e5e75cf8d4cd558955a5c7638fba2d68f464/YM/Fields/System.lean\#L124-L135}{Realizes}.

\textbf{Proposition V.2.10.} Fix rotation and real structures, induced
involutive \(J\), and spin compatibility. Assume word bilinearity,
two-test continuity, word locality, label faithfulness, continuous word
matrix elements and invariant Gram form. If vectors satisfy V.2.8's Gram
limit, \(\Phi_s(\varnothing)\to\Omega\), and eventual translation
covariance for each word and translation, and \(T\) is normalized,
unital and permutation symmetric, V.2.9 holds except possibly vacuum
uniqueness. If the heat-fixed space is \(\mathbb C\Omega\), uniqueness
also holds. \emph{Status:} theorem. \emph{Lean:}
\href{https://github.com/yang-mills-group/OpenAI-YM-20261028/blob/3456e5e75cf8d4cd558955a5c7638fba2d68f464/YM/Fields/Lo.lean\#L85-L95}{harvest}.

\subsubsection{2. What it uses}\label{what-it-uses-21}

\begin{itemize}
\tightlist
\item
  I.1, II.1: quotient labels and common finite-dimensional levels.
\item
  IV.1, V.1: Euclidean hierarchy, OS space, heat semigroup, translations
  and slab bound.
\item
  VI.1: rotation, spin and star data. VIII.1 supplies phase data and
  selected-phase heat-fixed uniqueness.
\item
  Standard results: OS reconstruction, tempered holomorphic boundary
  values, Fourier--Laplace theory, Kolmogorov construction,
  Riesz--Markov and SNAG.
\end{itemize}

\subsubsection{3. Strategy}\label{strategy-21}

Bounded slab insertions construct strict-time correlations. Boundary
values and boosts give covariant distributions with forward spectral
support. Fourier--Laplace continuation and permutation symmetry give
locality. Heated vectors transfer positivity and translations to the
word Hilbert space, where left multiplication defines fields.

\subsubsection{4. Steps}\label{steps-21}

\begin{enumerate}
\def\labelenumi{\arabic{enumi}.}
\tightlist
\item
  Heat factors telescope, removing the slab width; contractions give
  polynomial growth.
  \href{https://github.com/yang-mills-group/OpenAI-YM-20261028/blob/3456e5e75cf8d4cd558955a5c7638fba2d68f464/YM/Work/C02.lean\#L156-L166}{C02.matrix\_pad}.
\item
  Taylor's integral remainder controls boundary limits uniformly over
  nondecreasing heights.
  \href{https://github.com/yang-mills-group/OpenAI-YM-20261028/blob/3456e5e75cf8d4cd558955a5c7638fba2d68f464/YM/Shift/Breath.lean\#L78-L101}{Waves.maintains}.
\item
  Continue rotation identities, then integrate boost derivatives to
  obtain covariance.
  \href{https://github.com/yang-mills-group/OpenAI-YM-20261028/blob/3456e5e75cf8d4cd558955a5c7638fba2d68f464/YM/Work/C09.lean\#L374-L399}{C09.boost\_invariant}.
\item
  The Laplace identity forces nonnegative suffix energies; Lorentz
  covariance upgrades this to cone support.
  \href{https://github.com/yang-mills-group/OpenAI-YM-20261028/blob/3456e5e75cf8d4cd558955a5c7638fba2d68f464/YM/Shift/Closing.lean\#L106-L121}{tents}.
\item
  Fourier--Laplace damping extends \(F\) to the tube. Exponential decay
  gives complex differentiability in Schwartz space.
  \href{https://github.com/yang-mills-group/OpenAI-YM-20261028/blob/3456e5e75cf8d4cd558955a5c7638fba2d68f464/YM/Work/C13.lean\#L113-L125}{C13.analytic\_interior}.
\item
  Imaginary boosts construct two germs; equal Wick-slice jets identify
  them, yielding boundary exchange.
  \href{https://github.com/yang-mills-group/OpenAI-YM-20261028/blob/3456e5e75cf8d4cd558955a5c7638fba2d68f464/YM/Work/C20.lean\#L67-L87}{C20.complex\_germ}.
\item
  Gram limits give an isometric embedding into the OS space; left
  multiplication defines fields, and the embedding transfers positive
  energy.
  \href{https://github.com/yang-mills-group/OpenAI-YM-20261028/blob/3456e5e75cf8d4cd558955a5c7638fba2d68f464/YM/Fields/Motion.lean\#L31-L53}{universal\_embed}.
\end{enumerate}

\subsubsection{5. Where the difficulty
lies}\label{where-the-difficulty-lies-21}

Boost derivatives require finite-dimensional label levels and Schwartz
estimates. Locality requires complete jets at equal pair times.
Independent heating scales require unrestricted boundary limits.

\subsubsection{6. Notation}\label{notation-21}

\(\mathcal H,\Omega,H\ge0,U\): OS space, vacuum, Hamiltonian and
translations. \(\mathcal F=\bigcup_d F_d\): quotient labels; \(J\):
reflection adjoint; \(b^*\): reversed starred word.
\(E_n=\{y_1<\cdots<y_n\}\); \(F_n^y(x)=F_n(x^0+iy;\vec x)\);
\(\underline y=((\vec0,y_k))_k\).
\(\mathcal T_n=\{z:\mathrm{Im}(z_{k+1}-z_k)\in V_+\}\), with
unrestricted first imaginary point. \(V_+=\{x:|\vec x|<x^0\}\). Pairings
are Euclidean; \(\tilde\Lambda=\mathsf r\Lambda\mathsf r\), with spatial
reversal \(\mathsf r\). \(\mathcal L^\uparrow_+\) is the proper
orthochronous Lorentz group; \(K\) is its label action. The closed
imaginary cone replaces \(V_+\) by \(\bar V_+\).
\(\mathcal FW(f)=W(\hat f)\); Fourier phases are
\(e^{-i\langle x,p\rangle}\), inverse normalization \((2\pi)^{-4n}\).
\(g_*f=f\circ g^{-1}\) slotwise; \(\rho\) is the rotation Lie action.

\subsubsection{7. What remains
unexplained}\label{what-remains-unexplained-21}

The upstream slab laws, compact Borel functional calculus and some
Fourier--Laplace differentiation estimates remain outside the detailed
checks. Vacuum uniqueness applies to the selected phase; it is not
supplied for the mixed table.

Detailed notes: the full account of this unit, with proof sketches for
every step, is on the
\href{https://github.com/scottnarmstrong/yang-mills-informalization/blob/main/outline/notes/V-2.md}{detailed
notes page}.

\subsection{V.3. Reconstruction in every
phase}\label{v.3.-reconstruction-in-every-phase}

\textbf{Part:} V, phases and reconstruction. \textbf{Lean:}
\href{https://github.com/yang-mills-group/OpenAI-YM-20261028/tree/3456e5e75cf8d4cd558955a5c7638fba2d68f464/YM/EveryPhase}{\texttt{YM/EveryPhase}}
(34 files). Its hypotheses are recorded in
\href{https://github.com/yang-mills-group/OpenAI-YM-20261028/blob/3456e5e75cf8d4cd558955a5c7638fba2d68f464/YM/Spec/Closed.lean\#L20-L73}{\texttt{YM/Spec/Closed.lean}}
and
\href{https://github.com/yang-mills-group/OpenAI-YM-20261028/blob/3456e5e75cf8d4cd558955a5c7638fba2d68f464/YM/Phase/Mixture.lean\#L65-L148}{\texttt{YM/Phase/Mixture.lean}}.
\textbf{Size:} 5,268 lines.

\subsubsection{1. What this unit proves}\label{what-this-unit-proves-22}

\textbf{Hypotheses throughout.} Fix a Wilson family, extraction and
limit table \(S^{\rm mix}\) satisfying
\href{https://github.com/yang-mills-group/OpenAI-YM-20261028/blob/3456e5e75cf8d4cd558955a5c7638fba2d68f464/YM/Spec/Closed.lean\#L64-L73}{\texttt{Final}},
constructed in VIII.1. This includes the Euclidean axioms; a compact
metric phase space \(X\) with full-support Borel probability \(\nu\);
continuous characters \(\chi_\xi\), indexed bijectively by \(X\), of the
commutative C\(^*\) ground-space slab algebra \(\mathfrak A\), with
\(\xi\mapsto\chi_\xi(M)\) continuous and
\(\int\chi_\xi(M)d\nu=\langle\omega,M\omega\rangle\) for every
\(M\in\mathfrak A\); spatial translations trivial on the mixed ground
space; phase tables satisfying the Euclidean axioms and
\(S^\xi(w)=\chi_\xi(M_w)\) for distinct-point words, where \(M_w\) is
their time-translated compressed slab operator; the selected
reconstruction at \(\xi_*\) and continuations of
\(S^{\rm mix},S^{\xi_*}\); and the VI.1 field algebra, derivatives,
rotations, Euclidean stress tensor and phase Ward identities.

\textbf{Proposition V.3.1 (every phase).} There exist OS spaces of
\(S^\xi\) with slab operators, continuations \(\mathsf F^\xi\) of
\(S^\xi\) with boundary distributions \(W^\xi\), and theories
\(\mathfrak W_\xi=(\mathcal H_\xi,\mathfrak D_\xi,\Omega_\xi,\phi_\xi,U_\xi,P_\xi)\)
for every \(\xi\in X\), agreeing with V.1--V.2 at \(\xi_*\). The common
label space \(\mathbb F\), generated by the \(F_d\), has the selected
star operation and Poincaré action; \(\mathfrak D_\xi\) is the dense
invariant word domain, \(\Omega_\xi\) a unit vacuum, and \(U_\xi\) a
strongly continuous unitary Poincaré representation on the Hilbert space
\(\mathcal H_\xi\). Fields \(\phi_\xi(a,f)\) use \(a\in\mathbb F\),
\(f\in\mathcal S(\mathbb R^4)\). For every phase:

\begin{enumerate}
\def\labelenumi{\arabic{enumi}.}
\item
  Vacuum field products equal \(W^\xi\); \(\phi_\xi(a,f)^*\) extends
  \(\phi_\xi(a^*,\bar f)\); spacelike separated fields commute on
  \(\mathfrak D_\xi\); \(U_\xi\) implements Poincaré covariance; the
  translation spectrum lies in \(\overline V_+\). Every nonzero
  \(a\in\mathbb F\) has some \(f\in\mathcal S(\mathbb R^4)\) with
  \(\phi_\xi(a,f)\ne0\) on \(\mathfrak D_\xi\); translation-invariant
  vectors are multiples of \(\Omega_\xi\).
\item
  A unitary \(J_\xi:\mathcal H_\xi\to\mathcal H^{\rm OS}_\xi\) sends
  \(\Omega_\xi\) to \(\omega_\xi\), intertwines translations and
  spectral projections, and sends word vectors to boundary values of the
  corresponding Euclidean vectors at imaginary times.
\item
  The OS ground space is \(\mathcal D_\xi=\mathbb C\omega_\xi\).
\item
  For every \(x\in\mathcal H_\xi\), \(P^0_\xi x=x\) iff
  \(x\in\mathbb C\Omega_\xi\).
\item
  For every compactly supported finite sum of words \(A\), some
  \(M\ge0\) satisfies \(\mu_A^\xi(\{0<p^4\le u\})\le Mu^2\) for all
  \(\xi\in X\), \(0<u\le1\).
\item
  For every such \(A\), some \(M\ge0\) satisfies, for all \(\xi\in X\),
  \(s\ge1\),

  \[0\le\langle A\Omega_\xi,(e^{-sH_\xi}-P^0_\xi)A\Omega_\xi\rangle\le Ms^{-2}.\]
\item
  A bounded complex-linear \(B\) is scalar iff, for every \(a,f\) and
  \(u,v\in\mathfrak D_\xi\), \(u\in\mathrm{dom}\,\phi_\xi(a,f)^*\) and
  \(\langle\phi_\xi(a,f)^*u,Bv\rangle=\langle u,B\phi_\xi(a,f)v\rangle\).
\item
  Proposition V.3.y holds.
\item
  For every \(n\), \(a\in\mathbb F^n\),
  \(g\in\mathcal S((\mathbb R^4)^n)\), \(\xi\mapsto W_n^\xi(a;g)\) is
  continuous and \(\int_XW_n^\xi(a;g)d\nu=W_n^{\rm mix}(a;g)\). For
  every distinct-point word \(w\),
  \(\int_XS^\xi(w)d\nu=S^{\rm mix}(w)\).
\end{enumerate}

\emph{Status:} theorem under the stated hypotheses. \emph{Lean:}
\href{https://github.com/yang-mills-group/OpenAI-YM-20261028/blob/3456e5e75cf8d4cd558955a5c7638fba2d68f464/YM/EveryPhase/Extension.lean\#L75-L80}{\texttt{all\_characters}}.
Items 5--7 are also Propositions V.3.n
(\href{https://github.com/yang-mills-group/OpenAI-YM-20261028/blob/3456e5e75cf8d4cd558955a5c7638fba2d68f464/YM/EveryPhase/Light.lean\#L221-L232}{\texttt{light\_bound}}),
V.3.o
(\href{https://github.com/yang-mills-group/OpenAI-YM-20261028/blob/3456e5e75cf8d4cd558955a5c7638fba2d68f464/YM/EveryPhase/Decay.lean\#L173-L192}{\texttt{inverse\_square}})
and V.3.t
(\href{https://github.com/yang-mills-group/OpenAI-YM-20261028/blob/3456e5e75cf8d4cd558955a5c7638fba2d68f464/YM/EveryPhase/Irreducible.lean\#L118-L193}{\texttt{commutant}}).

\textbf{Proposition V.3.2 (no exceptional phase).} For every
\(\xi\in X\) and \(x\in\mathcal H_\xi\),
\(P^0_\xi x=\langle\Omega_\xi,x\rangle\Omega_\xi\). \emph{Status:}
theorem. \emph{Lean:}
\href{https://github.com/yang-mills-group/OpenAI-YM-20261028/blob/3456e5e75cf8d4cd558955a5c7638fba2d68f464/YM/EveryPhase/Vacuum.lean\#L115-L144}{\texttt{all\_plant}}.

\textbf{Remark V.3.3 (additional consequences).} Under V.3.1, for every
\(\xi\) and \(x\in\mathcal H_\xi\), being fixed by \(e^{-sH_\xi}\) for
all \(s\ge0\), being translation invariant, and belonging to
\(\mathbb C\Omega_\xi\) are equivalent. For every degree \(d\) and
degree-\(d\) token coefficient vector \(v\), nullness for \(S^\xi\) is
equivalent to vanishing of the field of its \(F_d\) class on every test.
For every \(O\in O(4)\), \(b\in\mathbb R^4\), coordinate \(a\), finite
family of distinct-point words \(w_i\) in \(\{x:(O^{-1}(x-b))^a>0\}\)
and coefficients \(c_i\in\mathbb C\),
\[\sum_{i,j}\overline{c_i}c_jS^\xi(\theta_{O,b,a}w_i\,w_j)\ge0.\] Here
\(\theta_{O,b,a}\) reverses the word, reflects its points across the
hyperplane, and transports the \(a\)-axis label reflection through
\(O\). \emph{Status:} theorem consequences, unused by the main theorem.
\emph{Lean:}
\href{https://github.com/yang-mills-group/OpenAI-YM-20261028/blob/3456e5e75cf8d4cd558955a5c7638fba2d68f464/YM/EveryPhase/Extension.lean\#L100-L102}{\texttt{fixedHeat}},
\href{https://github.com/yang-mills-group/OpenAI-YM-20261028/blob/3456e5e75cf8d4cd558955a5c7638fba2d68f464/YM/EveryPhase/Extension.lean\#L103-L106}{\texttt{exactly\_fixed}},
\href{https://github.com/yang-mills-group/OpenAI-YM-20261028/blob/3456e5e75cf8d4cd558955a5c7638fba2d68f464/YM/EveryPhase/Extension.lean\#L90-L98}{\texttt{recognizes}},
\href{https://github.com/yang-mills-group/OpenAI-YM-20261028/blob/3456e5e75cf8d4cd558955a5c7638fba2d68f464/YM/EveryPhase/Extension.lean\#L84-L88}{\texttt{allPlanes}}.

\textbf{Proposition V.3.y (stress tensor).} For every \(\xi\in X\):

\begin{itemize}
\item
  \textbf{(T0)} \(\mathsf F^\xi\) extends holomorphically to the full
  forward tube (consecutive imaginary differences in the open forward
  cone), with continuous tested partial boundaries on the closed cone,
  equal to \(W^\xi\) at zero. Each adjacent pair has label-independent
  tied heights tending to zero with equal exchange germs at spacelike
  separation.
\item
  \textbf{(T1)--(T4)} For every proper orthochronous Poincaré
  transformation \(g\), Schwartz \(h\), indices \(a,b\in\{1,2,3,4\}\)
  and \(u,v,x\in\mathfrak D_\xi\),
  \(U_\xi(g)\mathfrak D_\xi\subset\mathfrak D_\xi\) and

  \[U_\xi(g)T_{ab}(h)x=\sum_{i,j}\Lambda(g)_{ia}\Lambda(g)_{jb}T_{ij}(h\circ g^{-1})U_\xi(g)x,\]

  \[T_{ab}=T_{ba},\qquad \langle T_{ab}(h)u,v\rangle=\langle u,T_{ab}(\bar h)v\rangle,\qquad \sum_a\eta_{aa}T_{ab}(\partial_ah)=0.\]
\item
  \textbf{(T5)} Fix \(B\in F_d\), a Killing field \(K\), compact
  \(\mathcal O\subset\mathbb R^4\), \(m\in C_c^\infty(\mathbb R^4)\)
  supported in \(\mathcal O\), real \(f\in C_c^\infty(\mathbb R)\) with
  \(\int f=1\), and complex \(c\in C_c^\infty(\mathbb R^3)\) equal to
  \(1\) on an open ball. Suppose every \((\vec y,s)\) over its
  complement, for \(s\in\mathrm{supp}\,f\), is strictly spacelike to
  every point of \(\mathcal O\). On \(\mathfrak D_\xi\),

  \[i[Q_K(f,c),\phi_\xi(B,m)]=\phi_\xi(\Sigma_KB,m)-\sum_a\phi_\xi(B,\partial_a(K^am)).\]
\item
  \textbf{(T6)} Under these conditions, \(i[Q_K(f,c),A]\), for every
  finite sum of words \(A\) supported in \(\mathcal O\), is the sum
  replacing one factor at a time by (T5), preserving order, for any
  representatives in the \(F_d\).
\end{itemize}

\emph{Status:} theorem. \emph{Lean:}
\href{https://github.com/yang-mills-group/OpenAI-YM-20261028/blob/3456e5e75cf8d4cd558955a5c7638fba2d68f464/YM/EveryPhase/Implement.lean\#L94-L116}{\texttt{Stress.force}}.

\subsubsection{2. What it uses}\label{what-it-uses-22}

\begin{itemize}
\tightlist
\item
  \textbf{IV.1:} Euclidean axioms
  (\href{https://github.com/yang-mills-group/OpenAI-YM-20261028/blob/3456e5e75cf8d4cd558955a5c7638fba2d68f464/YM/Spec/Closed.lean\#L20-L26}{\texttt{Seed}}).
\item
  \textbf{V.1:} character decomposition, copy bounds, OS realization and
  almost-everywhere vacuum uniqueness
  (\href{https://github.com/yang-mills-group/OpenAI-YM-20261028/blob/3456e5e75cf8d4cd558955a5c7638fba2d68f464/YM/Phase/Unique.lean\#L91-L110}{\texttt{Cluster.lifts}}).
\item
  \textbf{V.2:} reconstruction, polynomial one-point slab bounds and
  C01--C29 analytic estimates
  (\href{https://github.com/yang-mills-group/OpenAI-YM-20261028/blob/3456e5e75cf8d4cd558955a5c7638fba2d68f464/YM/Spec/Closed.lean\#L50-L54}{\texttt{Relativistic}}).
\item
  \textbf{VI.1:} field algebra and Euclidean stress Ward identities
  (\href{https://github.com/yang-mills-group/OpenAI-YM-20261028/blob/3456e5e75cf8d4cd558955a5c7638fba2d68f464/YM/Spec/Closed.lean\#L34-L43}{\texttt{Locals}}).
\item
  OS reconstruction; spectral theorem; Vitali--Montel, Morera and
  Liouville theorems; dominated convergence.
\end{itemize}

\subsubsection{3. Strategy}\label{strategy-22}

Uniform slab domination transfers reconstruction to every phase.
Character continuity passes through holomorphic continuation to boundary
distributions. Locality supplies uniform low-energy and heat decay; full
support and continuity extend vacuum-subtracted decay from almost every
phase to every phase. Its large-time limit identifies the vacuum
projection. Separately, Euclidean vector recovery proves surjectivity;
continued Ward identities give stress charges.

\subsubsection{4. Steps}\label{steps-22}

\begin{enumerate}
\def\labelenumi{\arabic{enumi}.}
\tightlist
\item
  \textbf{Uniform reconstruction.} Write \(\Pi_\xi(d,w)\) and
  \(\Pi(d,w)\) for phase and mixed slab operators. Repeated reflection
  and taking roots give \(\|\Pi_\xi(d,w)\|\le\|\Pi(d,w)\|\) for every
  positive thickness and distinct-point slab word. Apply V.2 phasewise.
  Lean:
  \href{https://github.com/yang-mills-group/OpenAI-YM-20261028/blob/3456e5e75cf8d4cd558955a5c7638fba2d68f464/YM/EveryPhase/Pack.lean\#L95-L100}{\texttt{sameBound}},
  \href{https://github.com/yang-mills-group/OpenAI-YM-20261028/blob/3456e5e75cf8d4cd558955a5c7638fba2d68f464/YM/EveryPhase/Pack.lean\#L157-L194}{\texttt{almanac}}.
\item
  \textbf{Boundary continuity.} Uniform polynomial bounds, Vitali's
  theorem and a Taylor formula with integrable remainder prove
  continuity and barycenters of \(W^\xi\). Lean:
  \href{https://github.com/yang-mills-group/OpenAI-YM-20261028/blob/3456e5e75cf8d4cd558955a5c7638fba2d68f464/YM/EveryPhase/Barycenter.lean\#L48-L73}{\texttt{word\_continuity}},
  \href{https://github.com/yang-mills-group/OpenAI-YM-20261028/blob/3456e5e75cf8d4cd558955a5c7638fba2d68f464/YM/EveryPhase/Barycenter.lean\#L75-L90}{\texttt{barycenter}}.
\item
  \textbf{Surjectivity.} Tests with compact frequency support and
  exponential damping recover positive-height Euclidean vectors in
  \(\mathrm{ran}\,J_\xi\); totality and closedness give surjectivity.
  Lean:
  \href{https://github.com/yang-mills-group/OpenAI-YM-20261028/blob/3456e5e75cf8d4cd558955a5c7638fba2d68f464/YM/EveryPhase/Surjective.lean\#L178-L223}{\texttt{onto}}.
\item
  \textbf{Low energy.} For compactly supported \(A\), locality makes the
  time-averaged odd vacuum commutator compactly supported in space. The
  nonnegative weight \((1-\cos E)/E\) (zero at \(E=0\)) and a momentum
  bump give shell mass \(O(u^2)\): spatial volume contributes \(u^3\),
  the weight contributes \(u\). Dyadic summation yields item 5; Laplace
  integration yields item 6. Lean:
  \href{https://github.com/yang-mills-group/OpenAI-YM-20261028/blob/3456e5e75cf8d4cd558955a5c7638fba2d68f464/YM/EveryPhase/Light.lean\#L221-L232}{\texttt{light\_bound}},
  \href{https://github.com/yang-mills-group/OpenAI-YM-20261028/blob/3456e5e75cf8d4cd558955a5c7638fba2d68f464/YM/EveryPhase/Decay.lean\#L173-L192}{\texttt{inverse\_square}}.
\item
  \textbf{Every vacuum.} Heat forms are continuous in \(\xi\).
  Almost-everywhere uniqueness and full support extend vacuum-subtracted
  decay everywhere; strong heat limits and density give V.3.2. Lean:
  \href{https://github.com/yang-mills-group/OpenAI-YM-20261028/blob/3456e5e75cf8d4cd558955a5c7638fba2d68f464/YM/EveryPhase/Vacuum.lean\#L51-L78}{\texttt{extends\_decay}},
  \href{https://github.com/yang-mills-group/OpenAI-YM-20261028/blob/3456e5e75cf8d4cd558955a5c7638fba2d68f464/YM/EveryPhase/Vacuum.lean\#L115-L144}{\texttt{all\_plant}}.
\item
  \textbf{Irreducibility.} Glue bounded half-plane matrix elements by
  weak centrality; Morera and Liouville make them constant. Vacuum
  uniqueness forces the commutant to be scalar. Lean:
  \href{https://github.com/yang-mills-group/OpenAI-YM-20261028/blob/3456e5e75cf8d4cd558955a5c7638fba2d68f464/YM/EveryPhase/Irreducible.lean\#L118-L193}{\texttt{commutant}}.
\item
  \textbf{Stress charges.} Continue Euclidean contact identities through
  spatially integrated current jumps, then remove spectator cutoffs and
  take word-vector matrix elements to obtain (T5). Lean:
  \href{https://github.com/yang-mills-group/OpenAI-YM-20261028/blob/3456e5e75cf8d4cd558955a5c7638fba2d68f464/YM/EveryPhase/Through.lean\#L129-L142}{\texttt{Stress.through\_jump}},
  \href{https://github.com/yang-mills-group/OpenAI-YM-20261028/blob/3456e5e75cf8d4cd558955a5c7638fba2d68f464/YM/EveryPhase/Charge.lean\#L143-L177}{\texttt{Stress.local\_action}}.
\end{enumerate}

\subsubsection{5. Where the difficulty
lies}\label{where-the-difficulty-lies-22}

Uniform analytic boundary bounds, Euclidean vector recovery and joint
current limits require care. Exchange germs agree only at spacelike
separation; current limits are spatially integrated. The \(u^2\)
estimate does not imply a mass gap. Uniqueness of the phase itself
remains undecided.

\subsubsection{6. Notation}\label{notation-22}

\(\mathcal D_\xi\) is the fixed space of the OS heat semigroup;
\(\mathfrak D_\xi\) is the Wightman word domain. Time is coordinate
\(4\); \(P^0_\xi=P_\xi(\{p^4=0\})\),
\(\mu_A^\xi(B)=\langle A\Omega_\xi,P_\xi(B)A\Omega_\xi\rangle\), and
\(e^{-sH_\xi}=J_\xi^*e^{-sH^{\rm OS}_\xi}J_\xi\). A compactly supported
word sum has all tests in \(C_c^\infty\). Write
\(\eta=\mathrm{diag}(-1,-1,-1,1)\), \(\zeta_4=-i\), \(\zeta_j=1\),
\(T_{ab}(h)=\phi_\xi(-\zeta_a\zeta_bT^{\rm E}_{ab},h)\), and
\(Q_K(f,c)=\sum_bT_{4b}(K^b(c\otimes f))\). Here \(K(x)=k+Vx\),
\(V^{\mathsf T}\eta+\eta V=0\), and \(\Sigma_K\) is its infinitesimal
label action.

The reconstruction family is
\href{https://github.com/yang-mills-group/OpenAI-YM-20261028/blob/3456e5e75cf8d4cd558955a5c7638fba2d68f464/YM/EveryPhase/Pack.lean\#L117-L132}{\texttt{Almanac}}
(\texttt{L.chambers}, \texttt{L.contours}, \texttt{L.actors},
\texttt{L.J}); word sums are
\href{https://github.com/yang-mills-group/OpenAI-YM-20261028/blob/3456e5e75cf8d4cd558955a5c7638fba2d68f464/YM/EveryPhase/Polynomials.lean\#L14}{\texttt{SumOfWords}};
the charge is
\href{https://github.com/yang-mills-group/OpenAI-YM-20261028/blob/3456e5e75cf8d4cd558955a5c7638fba2d68f464/YM/EveryPhase/Charge.lean\#L80-L81}{\texttt{Stress.Charge}}.

\subsubsection{7. What remains
unexplained}\label{what-remains-unexplained-22}

The C01--C29 estimates remain analytic inputs from V.2. The current-jump
argument needs explicit tube faces and spacelike neighborhoods;
integrated limits do not supply pointwise stress limits.

Detailed notes: the full account of this unit, with proof sketches for
every step, is on the
\href{https://github.com/scottnarmstrong/yang-mills-informalization/blob/main/outline/notes/V-3.md}{detailed
notes page}.

\subsection{VI.1. Ward identities, the rotation action and the conserved
stress
tensor}\label{vi.1.-ward-identities-the-rotation-action-and-the-conserved-stress-tensor}

\textbf{Part:} VI, symmetry: stress tensor and Ward identities.
\textbf{Lean:}
\href{https://github.com/yang-mills-group/OpenAI-YM-20261028/blob/3456e5e75cf8d4cd558955a5c7638fba2d68f464/YM/Bid/Town.lean}{\texttt{YM/Bid/Town.lean}},
\href{https://github.com/yang-mills-group/OpenAI-YM-20261028/tree/3456e5e75cf8d4cd558955a5c7638fba2d68f464/YM/Next/ROTATIONS}{\texttt{YM/Next/ROTATIONS}},
\href{https://github.com/yang-mills-group/OpenAI-YM-20261028/tree/3456e5e75cf8d4cd558955a5c7638fba2d68f464/YM/Ward}{\texttt{YM/Ward}},
\href{https://github.com/yang-mills-group/OpenAI-YM-20261028/tree/3456e5e75cf8d4cd558955a5c7638fba2d68f464/YM/Spin}{\texttt{YM/Spin}}.
\textbf{Size:} approximately 7,704 assigned lines, plus 6,100 of
collinear flux estimates.

\subsubsection{1. What this unit proves}\label{what-this-unit-proves-23}

Except for group statements, assume \(S\) satisfies the Euclidean
axioms: permutation symmetry, empty normalisation, unit deletion,
translation invariance, \(B_4\)- and conjugation-covariance, coordinate
reflection positivity for point words and smeared tests, smoothness, and
for each arity, labels and derivative order \(j\), constants
\(C\ge0,A\in\mathbb N\) with \(\|D^jS(x)\|\le Cr^{-A}\) whenever
\(0<r\le1\) and maximum-norm separation is at least \(r\), uniformly in
positions. The plan consists of complex-linear maps
\(\partial_a:F_d\to F_{d+1}\). Its three laws are: derivative law
\(S(\dots,\partial_av_k,\dots;x)=\partial_{x_k^a}S(v;x)\), inclusion law
\(S(\dots,\iota v_k,\dots)=S(v;x)\), and orientation: inclusions commute
with \(M,\Gamma\), derivatives commute with \(\Gamma\), and
\(M_\sigma\partial_b=\sum_a\sigma_{ab}\partial_aM_\sigma\).

The raw maps are complex-linear \(p_a:F_d\to F_{d+1}\) and real-linear
\(Y:\mathfrak{so}(4)\to\mathrm{End}_{\mathbb C}F_d\); spin rows \(T\)
have the latter type. Write (R) for: \(J_{ab}\in F_4\) is symmetric,
\(\Gamma\)-real, \(\sum_a\partial_aJ_{ab}=0\) in \(F_5\),
\(M_\sigma J_{ab}=\sum_{ij}\sigma_{ia}\sigma_{jb}J_{ij}\), and
\((J;p,Y)\) obeys the stress laws. For every separated
\((x_0,\dots,x_n)\), labels \(B,q\), and real \(\chi\in C_c^\infty\)
equal to \(1\) near \(x_0\) with closed support excluding the other
points, these laws for \((J;P,A)\) are
\[-\sum_{a,b}\int\partial_a\chi(y)V_b(y)S(J_{ab}(y),B(x_0),q)\,dy=\begin{cases}S(P_jB(x_0),q),&V=e_j,\\S(A_KB(x_0),q),&V=K(y-x_0).\end{cases}\]
All identities concern separated insertions; no contact term is
evaluated. Define
\[\mathcal W_K[T,P](https://github.com/scottnarmstrong/yang-mills-informalization/blob/main/outline/units/q;x)=\sum_i\left[S(\dots,T_Kq_i,\dots;x)+\sum_b(Kx_i)_bS(\dots,P_bq_i,\dots;x)\right],\qquad\mathcal W_K[T]=\mathcal W_K[T,\partial].\]

\textbf{Proposition VI.1.1.} Assume the derivative and orientation laws
and (R). For every arity \(n\), degree \(d\), coordinate axis
\(\mathbb Re_l\), every \(x\in\mathrm{Conf}_n\) on that axis, all tags
\(q\) and \(K\in\mathfrak{so}(4)\),
\(\mathcal W_K[Y,p](https://github.com/scottnarmstrong/yang-mills-informalization/blob/main/outline/units/q;x)=0\).
For every \(j\ne l\) and \(v_1,\dots,v_n\in F_d\),
\(\sum_iS(\dots,p_jv_i,\dots;x)=0\). \emph{Status:} theorem.
\emph{Lean:}
\href{https://github.com/yang-mills-group/OpenAI-YM-20261028/blob/3456e5e75cf8d4cd558955a5c7638fba2d68f464/YM/Next/ROTATIONS/Route.lean\#L114-L121}{\texttt{Outer.outer}},
\href{https://github.com/yang-mills-group/OpenAI-YM-20261028/blob/3456e5e75cf8d4cd558955a5c7638fba2d68f464/YM/Bid/Angular.lean\#L35-L36}{\texttt{Bid.LineSum}}.

\textbf{Proposition VI.1.2.} Assume all three plan laws, the ultraviolet
record and free-Maxwell data of §2. There exists real \(c\ne0\),
independent of degree and axis, with \(p_a=c\partial_a\) on every
\(F_d\). For all degrees, axes and \(K\), the rows \(D=c^{-1}Y\) satisfy
(C) \(\partial_bD_K+\sum_iK_{ib}\partial_i=D_K\partial_b\) and (L)
\(\mathcal W_K[D]=0\) on every coordinate axis, for all labels and
\(K\). \emph{Status:} theorem. \emph{Lean:}
\href{https://github.com/yang-mills-group/OpenAI-YM-20261028/blob/3456e5e75cf8d4cd558955a5c7638fba2d68f464/YM/Bid/Town.lean\#L49-L59}{\texttt{Bid.pressured}},
\href{https://github.com/yang-mills-group/OpenAI-YM-20261028/blob/3456e5e75cf8d4cd558955a5c7638fba2d68f464/YM/Bid/Town.lean\#L61-L78}{\texttt{Bid.formed}}.

\textbf{Proposition VI.1.3.} Assume the derivative law and spin rows
\(T\) satisfying (C),(L). Then
\(\mathcal W_K[T](https://github.com/scottnarmstrong/yang-mills-informalization/blob/main/outline/units/q;x)=0\)
for every arity, separated configuration, tags and \(K\). \emph{Status:}
theorem. \emph{Lean:}
\href{https://github.com/yang-mills-group/OpenAI-YM-20261028/blob/3456e5e75cf8d4cd558955a5c7638fba2d68f464/YM/Ward/Axis.lean\#L130-L139}{\texttt{Ward.global}}.

\textbf{Proposition VI.1.4.} Assume the inclusion and derivative laws,
VI.1.3's conclusion, and compatibility of \(T\): it commutes with
\(\Gamma,\iota\), kills \(\mathbf1\), and
\(M_\sigma T_K=T_{\sigma K\sigma^{-1}}M_\sigma\). Assume internal row
erasure: for some \(N\), every \(d\ge N\), and every alternating
real-bilinear family \(\Delta(u,v)\in\mathrm{End}_{\mathbb C}F_d\) that
is \(\Gamma\)-real, covariant under determinant-one signed permutations,
kills \(\mathbf1\) and preserves every \(\iota_{e\to d}(F_e)\),
\(e\le d\), vanishing insertion sums for all words, separated
configurations and \(u,v\) imply \(\Delta=0\). Then
\([T_u,T_v]=T_{[u,v]}\) for all \(d,u,v\). \emph{Status:} theorem;
erasure is proved here. \emph{Lean:}
\href{https://github.com/yang-mills-group/OpenAI-YM-20261028/blob/3456e5e75cf8d4cd558955a5c7638fba2d68f464/YM/Work/A08.lean\#L232-L247}{\texttt{Work.A08.answer}}.

\textbf{Proposition VI.1.5.} Under VI.1.2's hypotheses, for every \(d\)
there is a complex-linear isomorphism \(V_d:F_d^{\rm free}\to F_d\)
intertwining \(M,\Gamma\) and satisfying
\(D_KV_d=V_d\mathcal L_K^{\rm free}\) for every \(K\). Thus
\(R_d=V_dR_d^{\rm free}V_d^{-1}\) is a continuous \(O(4)\)
representation, equals \(M\) on \(B_4\), commutes with \(\Gamma\), has
generators \(D\) satisfying the global Ward identity of VI.1.3. The
\(V_d\) are chosen independently across degrees. Such an action with
specified generators is a rotation port. \emph{Status:} theorem.
\emph{Lean:}
\href{https://github.com/yang-mills-group/OpenAI-YM-20261028/blob/3456e5e75cf8d4cd558955a5c7638fba2d68f464/YM/Bid/Town.lean\#L124-L157}{\texttt{Bid.services}}.

\textbf{Proposition VI.1.6.} Assume the inclusion and derivative laws
and a rotation port \(R\) with generators \(T\) satisfying VI.1.3's
global Ward identity. For all \(k\in O(4)\), tags and separated \(x\),
\(S(R(k)v;kx)=S(v;x)\). Moreover \(R\) fixes \(\mathbf1\), commutes with
inclusions, retains the port properties, and
\(R_{d+1}(k)\partial_av=\sum_i k_{ia}\partial_iR_d(k)v\) for every
\(d,a,v\). \emph{Status:} theorem. \emph{Lean:}
\href{https://github.com/yang-mills-group/OpenAI-YM-20261028/blob/3456e5e75cf8d4cd558955a5c7638fba2d68f464/YM/Work/A14.lean\#L417-L420}{\texttt{Work.A14.answer}}.

\textbf{Proposition VI.1.7.} Assume the derivative law, a covariant
rotation action satisfying VI.1.6's laws, (R), \(p=c\partial\) with
\(c\ne0\), and generators \(c^{-1}Y\). There exists symmetric, real,
divergence-free \(J'\) with
\(R_4(k)J'_{ab}=\sum_{ij}k_{ia}k_{jb}J'_{ij}\) for every \(k\),
satisfying the stress laws with \((P,A)=(\partial,c^{-1}Y)\). The
translation shell, Ward identity (13), therefore equals
\(\partial_{x_0^b}S(B(x_0),q)\) for every window and labels.
\emph{Status:} theorem. \emph{Lean:}
\href{https://github.com/yang-mills-group/OpenAI-YM-20261028/blob/3456e5e75cf8d4cd558955a5c7638fba2d68f464/YM/Work/A15.lean\#L187-L192}{\texttt{Work.A15.answer}}.

\textbf{Proposition VI.1.8.} Spatial rotations and coordinate boosts
generate \(SO^+(1,3)\). Under VI.1.2's hypotheses, choose \(V,R,J'\) as
above: \(B_d(\Lambda)=V_dB_d^{\rm free}(\Lambda)V_d^{-1}\) commutes with
every inclusion \(F_d\hookrightarrow F_e\), \(d\le e\), for every
\(\Lambda\). These actions give a representation on the union,
continuous on each \(F_d\), equal to \(R_d\) on spatial rotations and
commuting with \(M_{\theta_t}\Gamma\). For every \(d\), spatial \(i\)
and \(v\in F_d\), its boost derivative at zero is \(-iD_{K^{(3i)}}v\).
\emph{Status:} theorems. \emph{Lean:}
\href{https://github.com/yang-mills-group/OpenAI-YM-20261028/blob/3456e5e75cf8d4cd558955a5c7638fba2d68f464/YM/Work/A16.lean\#L91-L101}{\texttt{Work.A16.answer}},
\href{https://github.com/yang-mills-group/OpenAI-YM-20261028/blob/3456e5e75cf8d4cd558955a5c7638fba2d68f464/YM/Work/A17.lean\#L88-L101}{\texttt{Work.A17.answer}},
\href{https://github.com/yang-mills-group/OpenAI-YM-20261028/blob/3456e5e75cf8d4cd558955a5c7638fba2d68f464/YM/Spin/Floor.lean\#L39-L74}{\texttt{Spin.reach}}.

\textbf{Proposition VI.1.9.} Given §2's free data, infrared triple and
ultraviolet record, all six assembly clauses hold: raw compatibility,
normalisation, rotation descent, covariance, averaged stress laws, and
Lorentz compatibility. With the inclusion and derivative laws, there
exist \(c,(V_d),J'\) satisfying VI.1.2 and VI.1.5--8, all three plan
laws and the resulting local field laws. \emph{Status:} conditional
theorem. \emph{Lean:}
\href{https://github.com/yang-mills-group/OpenAI-YM-20261028/blob/3456e5e75cf8d4cd558955a5c7638fba2d68f464/YM/Next/ROTATIONS.lean\#L126-L131}{\texttt{Outer.commission}},
\href{https://github.com/yang-mills-group/OpenAI-YM-20261028/blob/3456e5e75cf8d4cd558955a5c7638fba2d68f464/YM/Org/Twist.lean\#L45-L53}{\texttt{Commission.picks}}.

\subsubsection{2. What it uses}\label{what-it-uses-23}

\begin{itemize}
\tightlist
\item
  IV.1: limit table
  \href{https://github.com/yang-mills-group/OpenAI-YM-20261028/blob/3456e5e75cf8d4cd558955a5c7638fba2d68f464/YM/Euclid/Words.lean\#L141-L148}{\texttt{Table.Hierarchy}},
  differential plan, commuting derivatives and injective inclusions.
\item
  V.1:
  \href{https://github.com/yang-mills-group/OpenAI-YM-20261028/blob/3456e5e75cf8d4cd558955a5c7638fba2d68f464/YM/OS/Dynamics.lean\#L57-L61}{\texttt{OS.Chamber}},
  OS Hilbert spaces, contraction semigroups and joint spectral measures,
  constructed from the axioms.
\item
  VI.2:
  \href{https://github.com/yang-mills-group/OpenAI-YM-20261028/blob/3456e5e75cf8d4cd558955a5c7638fba2d68f464/YM/Org/Input.lean\#L159-L170}{\texttt{Org.MaxwellData}};
  free plan, covariant rotation and Lorentz actions, stress laws for
  \((J^{\rm free};\kappa\partial^{\rm free},\kappa\mathcal L^{\rm free})\)
  with \(\kappa\ne0\), derivatives nonzero on \(F^{\rm free}_4\),
  generator relations, Casimir polynomials and row erasure. At every
  sufficiently large degree, finitely many collinear readouts have
  common kernel \(\mathbb R\partial^{\rm free}\) within the free
  vector-map space. Their producers assume the gauge group simple.
\item
  VI.2 and III.11:
  \href{https://github.com/yang-mills-group/OpenAI-YM-20261028/blob/3456e5e75cf8d4cd558955a5c7638fba2d68f464/YM/Org/Input.lean\#L172-L175}{\texttt{Org.IR}}
  and
  \href{https://github.com/yang-mills-group/OpenAI-YM-20261028/blob/3456e5e75cf8d4cd558955a5c7638fba2d68f464/YM/Org/Input.lean\#L177-L186}{\texttt{Org.UV}},
  including (R), orientation, lattice/extraction data, corrected-field
  bounds, grid approximations and weak scaling comparison of correlation
  jets and stress with the free table.
\item
  Weak comparison means: for every arity bound \(N\), degree bound
  \(d\), derivative bound \(m\) and \(R\ge1\), there exist
  \(r_k\downarrow0\), \(\xi_k>0\) and complex-linear isomorphisms
  \(U_{k,e}:F_e^{\rm free}\to F_e\). Through degree \(\max(d,4)\) these
  intertwine \(M,\Gamma\), units and inclusions; for \(n\le N\), degrees
  \(\le d\) and \(j\le m\),
  \(D_y^jS(U_kv;r_ky)\to D^jS^{\rm free}(v;y)\) uniformly where
  \(\|y\|\le R\) and separation is at least \(R^{-1}\). Also
  \(\xi_k^{-1}r_k^4U_{k,4}^{-1}J_{ab}\to J_{ab}^{\rm free}\) for every
  \(a,b\). The system may depend on all four bounds.
\item
  Standard results: spectral theorem, holomorphic semigroups, analytic
  vectors, analytic identity theorem, Haar averaging, Cartan--Dieudonné,
  Givens rotations, McWeeny purification and linear ODE uniqueness.
\end{itemize}

\subsubsection{3. Strategy}\label{strategy-23}

Conservation turns shell laws into collinear Ward identities; free
comparison fixes their normalisation. Positivity converts collinear
two-point jets into analytic OS vectors, allowing propagation to every
separated configuration. Exact intertwiners transport the free
rotations; covariance then gives compatibility across degrees. Rescaling
and averaging give the tensorial stress, and boost ODE uniqueness gives
Lorentz compatibility.

\subsubsection{4. Steps}\label{steps-23}

\begin{enumerate}
\def\labelenumi{\arabic{enumi}.}
\tightlist
\item
  Moving the centre inside a fixed window gives
  \(\partial_bY_K+\sum_cK_{cb}p_c=Y_K\partial_b\) and
  \(\partial_bp_i=p_i\partial_b\) by differentiating its shell laws.
  \href{https://github.com/yang-mills-group/OpenAI-YM-20261028/blob/3456e5e75cf8d4cd558955a5c7638fba2d68f464/YM/Work/A02.lean\#L272-L307}{\texttt{Work.A02.answer}}.
\item
  Cutoffs of widths \(\lambda^{100},\lambda^{150},\lambda^{300}\) cancel
  two terms by stress symmetry; reflection positivity and spectral
  estimates bound the exterior flux by \(O(\lambda^{-1})\), giving the
  collinear identity without a mass gap or clustering.
  \href{https://github.com/yang-mills-group/OpenAI-YM-20261028/blob/3456e5e75cf8d4cd558955a5c7638fba2d68f464/YM/Next/ROTATIONS/Yard.lean\#L105-L114}{\texttt{Outer.blows}}.
  \emph{(supporting analysis partly read)}
\item
  Finite free collinear tests have a one-dimensional kernel; persistence
  of rank under scaling forces \(p=c\partial\).
  \href{https://github.com/yang-mills-group/OpenAI-YM-20261028/blob/3456e5e75cf8d4cd558955a5c7638fba2d68f464/YM/Climb/Peaks.lean\#L98-L111}{\texttt{Climb.extends\_form}}.
\item
  Analytic comparison determines two-point jets; positivity gives
  factorial vector-derivative bounds. Extreme-point moves and density
  propagate vanishing angular sums.
  \href{https://github.com/yang-mills-group/OpenAI-YM-20261028/blob/3456e5e75cf8d4cd558955a5c7638fba2d68f464/YM/Ward/Ket.lean\#L76-L96}{\texttt{Analysis.blossoms}},
  \href{https://github.com/yang-mills-group/OpenAI-YM-20261028/blob/3456e5e75cf8d4cd558955a5c7638fba2d68f464/YM/Ward/Axis.lean\#L130-L139}{\texttt{Ward.global}}.
\item
  Free row erasure removes the bracket defect. Relative-Casimir
  projection and iteration \(x\mapsto3x^2-2x^3\) produce exact
  intertwiners: the compressed defect satisfies perfect commutator
  relations, which force sufficiently small operators to vanish.
  \href{https://github.com/yang-mills-group/OpenAI-YM-20261028/blob/3456e5e75cf8d4cd558955a5c7638fba2d68f464/YM/Work/A08.lean\#L232-L247}{\texttt{Work.A08.answer}},
  \href{https://github.com/yang-mills-group/OpenAI-YM-20261028/blob/3456e5e75cf8d4cd558955a5c7638fba2d68f464/YM/Spin/Sheets.lean\#L134-L157}{\texttt{Spin.meet}}.
\item
  Group generation integrates covariance; Haar averaging of \(c^{-1}J\)
  preserves the stress laws. Boost ODE uniqueness proves compatibility
  across degrees.
  \href{https://github.com/yang-mills-group/OpenAI-YM-20261028/blob/3456e5e75cf8d4cd558955a5c7638fba2d68f464/YM/Work/A14.lean\#L417-L420}{\texttt{Work.A14.answer}},
  \href{https://github.com/yang-mills-group/OpenAI-YM-20261028/blob/3456e5e75cf8d4cd558955a5c7638fba2d68f464/YM/Work/A15.lean\#L187-L192}{\texttt{Work.A15.answer}},
  \href{https://github.com/yang-mills-group/OpenAI-YM-20261028/blob/3456e5e75cf8d4cd558955a5c7638fba2d68f464/YM/Work/A17.lean\#L70-L84}{\texttt{run\_boost}}.
\end{enumerate}

\subsubsection{5. Where the difficulty
lies}\label{where-the-difficulty-lies-23}

Diagonal stress components require a conservation estimate beyond
reflection-parity arguments. Vanishing smooth jets alone do not justify
continuation; positivity supplies analyticity. Exact intertwiners
identify the required \(O(4)\) action and signed permutations.

\subsubsection{6. Notation}\label{notation-23}

\(S\): correlation table; \(F_d\): finite-dimensional complex label
spaces of degree \(\le d\) modulo null labels; \(\iota\): inclusions;
\(\mathbf1\): unit; \(B_4,M\): signed permutations and their action;
\(\Gamma\): conjugation; \(\mathrm{Conf}_n\): distinct points in
\(\mathbb R^4\). \(K\in\mathfrak{so}(4)\) is real skew, with
\([u,v]=uv-vu\) and \(K^{(ba)}=E_{ba}-E_{ab}\); time is axis \(3\), and
\(\theta_t\) reverses it. \(\lambda=e^k\) is the cutoff scale.

\subsubsection{7. What remains
unexplained}\label{what-remains-unexplained-23}

Supporting flux analysis, including contact-extension normalisation, was
not checked fully (probable). The ultraviolet hypotheses of the input
producers are not discharged here. Centring of \(J'\) is not proved
here; neither sign of \(c,\kappa\) is specified. Lean elaboration and
transitive axiom dependencies were not checked.

Detailed notes: the full account of this unit, with proof sketches for
every step, is on the
\href{https://github.com/scottnarmstrong/yang-mills-informalization/blob/main/outline/notes/VI-1.md}{detailed
notes page}.

\subsection{VI.2. The free Maxwell stress tensor and the local Ward
charge}\label{vi.2.-the-free-maxwell-stress-tensor-and-the-local-ward-charge}

\textbf{Part:} VI, symmetry: stress tensor and Ward identities.
\textbf{Lean:}
\href{https://github.com/yang-mills-group/OpenAI-YM-20261028/tree/3456e5e75cf8d4cd558955a5c7638fba2d68f464/YM/Next/MAXWELL}{\texttt{YM/Next/MAXWELL}},
\href{https://github.com/yang-mills-group/OpenAI-YM-20261028/tree/3456e5e75cf8d4cd558955a5c7638fba2d68f464/YM/UVNext/FREE_STRESS_CERTIFICATE}{\texttt{YM/UVNext/FREE\_STRESS\_CERTIFICATE}},
\href{https://github.com/yang-mills-group/OpenAI-YM-20261028/tree/3456e5e75cf8d4cd558955a5c7638fba2d68f464/YM/Next/UV_STRESS}{\texttt{YM/Next/UV\_STRESS}},
\href{https://github.com/yang-mills-group/OpenAI-YM-20261028/tree/3456e5e75cf8d4cd558955a5c7638fba2d68f464/YM/Pert}{\texttt{YM/Pert}}.
\textbf{Size:} 14,179 lines used by the main theorem (Maxwell 8.2k,
free-stress classification 3.1k, UV stress 2.6k, Pert 0.2k).

\subsubsection{1. What this unit proves}\label{what-this-unit-proves-24}

Throughout, \(G\) satisfies I.1, \(\mathfrak g\) is its Lie algebra,
spacetime indices run from zero to three, and \(c\) indexes an
orthonormal colour basis. A Maxwell wave is a smooth
\(\mathfrak g\)-valued two-form satisfying \(df=0\) and
\(\sum_a\partial_af_{ab}=0\). The algebra \(\mathcal W\) consists of
complex polynomials in homogeneous on-shell Taylor coefficients of order
\(j\), weighted \(j+2\); \(\mathcal P_d\) is its
\(\mathrm{Ad}(G)\)-invariant sector of weight at most \(d\). Wick
correlations contract only between sites, using the curl--curl of
\(G_0(x)=(4\pi^2|x|^2)^{-1}\). Set
\[T^{\rm free}_{ab}=\sum_{c,i}F^c_{ai}F^c_{bi}-\tfrac14\delta_{ab}\sum_{c,i,k}F^c_{ik}F^c_{ik}.\]
Let \(\Theta\) be its bilinear polarization, normalized by
\(\Theta(f,f)=2T^{\rm free}(f)\); \(j^v_a=\sum_bv_b\Theta_{ab}\) and
\(\Phi_\chi(j)=-\sum_a\int\partial_a\chi\,j_a\).

\textbf{Proposition VI.2.1 (free table).} For every register \(L\) and
\(d\ge0\), substitution gives a complex linear isomorphism
\(\iota_d:\mathcal F^{\rm free}_d\to\mathcal P_d\); a label is null
exactly when its polynomial vanishes on every on-shell Taylor tower.
Correlations are Wick moments. The table satisfies IV.1's Euclidean
hierarchy, including point and smeared reflection positivity, and is
real-analytic at distinct points. Each fixed correlation derivative is
bounded by \(Cr^{-A}\) for mutual sup-distances at least \(r\),
\(0<r\le1\), with \(A,C\ge0\). Taylor differentiation gives
\(\partial^{\rm free}:\mathcal F_d^{\rm free}\to\mathcal F_{d+1}^{\rm free}\)
with inclusion, derivative and orientation laws. The \(O(4)\) action and
generators \(\rho\) satisfy VI.1's rotation laws and preserve
correlations. The stress is symmetric, real, traceless, centred and
conserved in degree five. Simplicity is unnecessary. \emph{Status:}
theorem. \emph{Lean:}
\href{https://github.com/yang-mills-group/OpenAI-YM-20261028/blob/3456e5e75cf8d4cd558955a5c7638fba2d68f464/YM/Next/MAXWELL/System.lean\#L106-L114}{\texttt{hierarchy}},
\href{https://github.com/yang-mills-group/OpenAI-YM-20261028/blob/3456e5e75cf8d4cd558955a5c7638fba2d68f464/YM/Next/MAXWELL/System.lean\#L33-L35}{\texttt{fields}},
\href{https://github.com/yang-mills-group/OpenAI-YM-20261028/blob/3456e5e75cf8d4cd558955a5c7638fba2d68f464/YM/Next/MAXWELL/System.lean\#L176-L179}{\texttt{realStress}}.

\textbf{Proposition VI.2.2 (local charge).} Let \(U\) be open,
\(x\in U\), \(f\) a Maxwell wave on \(U\), \(v\) a smooth conformal
Killing field on \(\mathbb R^4\), and real \(\chi\in C_c^\infty(U)\)
equal one near \(x\). Put \(f^{p,x}(y)=\langle p(F)(x)F(y)\rangle_0\).
For each component probe \(p=F^c_{il}\),
\[\Phi_\chi(j^v(f^{p,x},f))=-p(\mathcal L_vf)(x),\qquad\mathcal L_vf=d\iota_vf.\]
If \(\mathcal L_vf\) is a Maxwell wave, this holds for every probe of
every order \(j\ge0\). Translations, centred rotations and dilation give
respectively \(-\partial_a[p(f)]\), \(-(\omega\cdot p)(f)\) and
\(-(j+2)p(f)\) at \(x\). Flux vanishes for two Maxwell waves near
\(\mathrm{supp}\,\chi\), and for two generated fields with source \(x\)
when \(v\) is homogeneous about \(x\) of degree zero or one.
\emph{Status:} theorem. \emph{Lean:}
\href{https://github.com/yang-mills-group/OpenAI-YM-20261028/blob/3456e5e75cf8d4cd558955a5c7638fba2d68f464/YM/Next/MAXWELL/Contact.lean\#L137-L156}{\texttt{Maxwell.Free.local\_force}},
\href{https://github.com/yang-mills-group/OpenAI-YM-20261028/blob/3456e5e75cf8d4cd558955a5c7638fba2d68f464/YM/Next/MAXWELL/Force.lean\#L131-L186}{\texttt{jet\_force}},
\href{https://github.com/yang-mills-group/OpenAI-YM-20261028/blob/3456e5e75cf8d4cd558955a5c7638fba2d68f464/YM/Next/MAXWELL/Repeat.lean\#L63-L128}{\texttt{two\_flux\_zero}}.

\textbf{Proposition VI.2.3 (polynomial Ward identities).} For \(n\ge0\),
distinct \(x_0,\ldots,x_n\), \(p_i\in\mathcal W\), and real
\(\chi\in C_c^\infty\) equal one near \(x_0\) with spectators outside
its support,
\[-\sum_{a,b}\int\partial_a\chi(y)v_b(y)\langle T^{\rm free}_{ab}(y)\prod_i p_i(x_i)\rangle_0\,dy=-\langle(Dp_0)(x_0)\prod_{i>0}p_i(x_i)\rangle_0.\]
Here \((v,D)\) is \((e_a,\partial_a)\), \((\omega(y-x_0),\rho(\omega))\)
for \(\omega\in\mathfrak{so}(4)\), or \((y-x_0,\varepsilon)\), where
\(\varepsilon\) multiplies monomials by weight. Thus
\(\tau T^{\rm free}\) has translation/rotation rows
\(-\tau\partial^{\rm free},-\tau\rho\) for every real \(\tau\);
normalization to \(+\partial^{\rm free}\) gives \(-T^{\rm free}\).
\emph{Status:} theorem. \emph{Lean:}
\href{https://github.com/yang-mills-group/OpenAI-YM-20261028/blob/3456e5e75cf8d4cd558955a5c7638fba2d68f464/YM/Next/MAXWELL/Taxes.lean\#L177-L200}{\texttt{Maxwell.Free.moment\_force}}.

\textbf{Proposition VI.2.4 (free stress line).} Suppose \(G\) is simple
and \(W\) satisfies the free derivative and orientation laws. Its
derivative maps are \(\partial^{\rm free}\). Let \(\mathcal Z\) be the
real space of centred symmetric real signed-covariant degree-four
tensors, and \(\mathcal V\) the real signed-covariant degree-five
vectors. Then
\[\mathrm{div}:\mathcal Z\twoheadrightarrow\mathcal V,\quad\dim_{\mathbb R}\mathcal Z=3,\quad\dim_{\mathbb R}\mathcal V=2,\quad\ker\mathrm{div}=\mathbb RT^{\rm free}.\]
Every lighthouse with this plan and stress \(\tau T^{\rm free}\),
\(\tau\ne0\), satisfies (FD). \emph{Status:} theorem. \emph{Lean:}
\href{https://github.com/yang-mills-group/OpenAI-YM-20261028/blob/3456e5e75cf8d4cd558955a5c7638fba2d68f464/YM/UVNext/FREE_STRESS_CERTIFICATE/Certificate.lean\#L125-L132}{\texttt{certificate}}.

\textbf{Proposition VI.2.5 (Maxwell data).} For simple \(G\), every
Wilson family \(C\) and real \(\tau\ne0\),
\((\partial^{\rm free},R,\tau T^{\rm free},-\tau)\) satisfies all
lighthouse laws and extends to Maxwell data for \(C\). The derivative is
nonzero in degree four. For all sufficiently large \(d\), finitely many
derivatives of Ward sums at distinct collinear configurations
perpendicular to the component axis detect exactly
\(\mathbb R\partial^{\rm free}\) among complex linear vector maps
\(\mathcal F_d^{\rm free}\to\mathcal F_{d+1}^{\rm free}\) that kill the
unit, are real, signed-covariant and preserve every lower-degree image
with degree raised by one.

The datum includes the hierarchy, VI.1's Casimir data, and compatible
degreewise continuous proper orthochronous Lorentz representations,
agreeing with spatial rotations, commuting with time-reflected
conjugation and having boost generator \(-i\rho\). For every \(d\ge4\),
any alternating real-bilinear family of complex linear endomorphisms of
\(\mathcal F_d^{\rm free}\) \(T(u,v)\), \(u,v\in\mathfrak{so}(4)\), that
is real, equivariant under orientation-preserving signed permutations,
kills the unit and preserves lower-degree images, vanishes if all its
Ward sums vanish for every \(u,v\), every arity \(m\ge0\), every
degree-\(d\) field tuple and every distinct-point configuration. Every
Maxwell datum can retain its plan, rotations, hierarchy and Lorentz
representation while acquiring stress \(\tau T^{\rm free}\), full stress
rotation covariance and (FD). Its Ward constant is necessarily
\(-\tau\). \emph{Status:} theorem. \emph{Lean:}
\href{https://github.com/yang-mills-group/OpenAI-YM-20261028/blob/3456e5e75cf8d4cd558955a5c7638fba2d68f464/YM/Next/MAXWELL/Lit.lean\#L45-L48}{\texttt{maxwell\_supplies}},
\href{https://github.com/yang-mills-group/OpenAI-YM-20261028/blob/3456e5e75cf8d4cd558955a5c7638fba2d68f464/YM/UVNext/FREE_STRESS_CERTIFICATE/Choose.lean\#L137-L150}{\texttt{FreeStress.fire}},
\href{https://github.com/yang-mills-group/OpenAI-YM-20261028/blob/3456e5e75cf8d4cd558955a5c7638fba2d68f464/YM/UVNext/FREE_STRESS_CERTIFICATE/Pressure.lean\#L158-L161}{\texttt{required}}.

\textbf{Proposition VI.2.6 (strict positivity).} For every axis \(a\),
\(x_a>0\) and nonzero \(p\in\mathcal W\),
\(\langle(\vartheta_ap)(\theta_ax)p(x)\rangle_0>0\). Every finite set of
distinct coordinate probes has positive-definite reflected propagator
Gram matrix there. Finite families of positive-half-space words, with
distinct sites within each word, satisfy reflection positivity. For
every \(d\ge0\), a polynomial in \(\mathcal P_d\) annihilated by all
distinct-point spectator tests from \(\mathcal P_d\) is zero.
\emph{Status:} theorem. \emph{Lean:}
\href{https://github.com/yang-mills-group/OpenAI-YM-20261028/blob/3456e5e75cf8d4cd558955a5c7638fba2d68f464/YM/Next/MAXWELL/Core.lean\#L94-L110}{\texttt{detects}},
\href{https://github.com/yang-mills-group/OpenAI-YM-20261028/blob/3456e5e75cf8d4cd558955a5c7638fba2d68f464/YM/Next/MAXWELL/Statistic.lean\#L80-L99}{\texttt{ghosts\_out}}.

\textbf{Proposition VI.2.7 (limit stress).} Assume II.2's Wilson
construction conditions, ordered smeared convergence and generic point
limits along \(F\) to a table \(S\) with Euclidean hierarchy \(H\), an
inclusion/derivative plan \(W\) on \(S\), a free hierarchy \(H_0\), and
a lighthouse with inclusion, derivative, orientation and (FD) laws.
Assume VI.2.8's station property and
\(\dim_{\mathbb C}\mathcal F_d^S\le\dim_{\mathbb C}\mathcal F_d^{\rm free}\)
for all \(d\ge0\). There exists an approaching stable scaling
\(X=(r_k,\xi_k,U_{k,d})\), initially with \(\xi_k=1\). For every such
scaling, interacting divergence \(\mathcal Z^S\to\mathcal V^S\), defined
as in VI.2.4, is onto with one-dimensional real kernel. There exist
\(0\ne J\in\mathcal Z^S\) with zero divergence and an approaching stable
refinement \(Y\) by subsequence and change of \(\xi_k\) such that
\[\xi_k^{-1}U_{k,4}^{-1}(r_k^4J_{ab})\to T_{0,ab}.\] For every finite
arity/degree/derivative bound and \(R\ge1\), tails give tracks
intertwining symmetries, units and inclusions through degree
\(\max(d,4)\), satisfying the stress limit above, with correlation
derivatives converging uniformly on \(K_{R,n}\) for each fixed tuple;
tails may depend on tuple and tolerance. Thus \(J\) satisfies VI.1's
comparison law. \emph{Status:} conditional theorem. \emph{Lean:}
\href{https://github.com/yang-mills-group/OpenAI-YM-20261028/blob/3456e5e75cf8d4cd558955a5c7638fba2d68f464/YM/Next/UV_STRESS/Invert.lean\#L132-L181}{\texttt{Stations.equivalences}},
\href{https://github.com/yang-mills-group/OpenAI-YM-20261028/blob/3456e5e75cf8d4cd558955a5c7638fba2d68f464/YM/Next/UV_STRESS/Limit.lean\#L129-L167}{\texttt{produce\_current}}.

\textbf{Proposition VI.2.8 (stations).} Fix the extraction data of
VI.2.7. A depth-\(D\) station consists of \(r>0\) and complex linear
lattice lifts \(A:\mathcal F_d^{\rm free}\to\mathrm{Ref}_d\),
\(d\le D\). Admissibility requires eventual uniform boundedness,
reality, signed covariance and preservation of the degree-zero unit,
and, for every tuple of arity \(n\le n_{\max}\) and degrees at most
\(D\), smooth complex test supported in \(K_{R,n}\), \(M>0\) with
\(\|t\|_\infty\le M\), and \(\eta>0\), eventual lattice/free smeared
discrepancy along the ordered cutoff tails at most
\(\varepsilon M\prod_i\|v_i\|+\eta\), using test \(r^{-4n}t(\cdot/r)\).
An admissible station yields complex linear continuum maps with the same
symmetries and estimate without \(\eta\):
\[\left|\int t(y)\left[\langle\prod_i Uv_i(ry_i)\rangle^S-\langle\prod_i v_i(y_i)\rangle_0\right]dy\right|\le\varepsilon M\prod_i\|v_i\|.\]

The cutoff tail may depend on the tuple, test, \(M\) and \(\eta\). The
station property requires these for every \(D,n_{\max}\ge0\), \(R\ge1\)
and \(\varepsilon,\delta>0\), with \(r\le\delta\). Balanced grid, strip
and finite-difference estimates of detailed Lemma VI.2.s, for all finite
bounds, imply it. \emph{Status:} definitions and conditional theorems.
\emph{Lean:}
\href{https://github.com/yang-mills-group/OpenAI-YM-20261028/blob/3456e5e75cf8d4cd558955a5c7638fba2d68f464/YM/Next/UV_STRESS/Station.lean\#L89-L97}{\texttt{Station.Admissible}},
\href{https://github.com/yang-mills-group/OpenAI-YM-20261028/blob/3456e5e75cf8d4cd558955a5c7638fba2d68f464/YM/Next/UV_STRESS/Station.lean\#L132-L180}{\texttt{harvested}},
\href{https://github.com/yang-mills-group/OpenAI-YM-20261028/blob/3456e5e75cf8d4cd558955a5c7638fba2d68f464/YM/Next/UV_STRESS/Granular.lean\#L134-L140}{\texttt{realizes\_stations}}.

\textbf{Proposition VI.2.9 (normalized limit).} Let \(C\) be a Wilson
family, \(F\) an extraction order, \(S\) a table with plan \(W\), \(M\)
a Maxwell datum with lighthouse Ward constant \(\kappa\ne0\), and
\(I=(J_0,p^I,Y^I)\). Assume VIII.1's ultraviolet-data conditions for
these objects. Assume \(W\)'s inclusion/derivative laws,
\(0\ne J_0\in\mathcal Z^S\), (FD), and an approaching stable scaling
with \(\xi_k^{-1}U_{k,4}^{-1}(r_k^4J_{0,ab})\to T_{0,ab}\). Every
symmetric real conserved signed-covariant
\(J\in(\mathcal F_4^S)^{4\times4}\) satisfying translation Ward rows
\(W\) and some rotation Ward rows obeys
\[U_{k,4}^{-1}(r_k^4J_{ab})\to\kappa^{-1}T_{0,ab}.\] For every finite
free-field tuple, any mixture of transported fields and rescaled stress
slots converges with every derivative uniformly on compact
distinct-point sets to the free correlation with every stress slot
replaced by \(\kappa^{-1}T_0\); the interacting insertions are evaluated
at \(r_ky_i\). Canonically \(\kappa=-1\), so each stress slot tends to
\(-T^{\rm free}\). \emph{Status:} conditional theorem. \emph{Lean:}
\href{https://github.com/yang-mills-group/OpenAI-YM-20261028/blob/3456e5e75cf8d4cd558955a5c7638fba2d68f464/YM/Pert/Weak.lean\#L118-L144}{\texttt{Pert.meter\_physical}},
\href{https://github.com/yang-mills-group/OpenAI-YM-20261028/blob/3456e5e75cf8d4cd558955a5c7638fba2d68f464/YM/Pert/Lifted.lean\#L85-L124}{\texttt{Pert.mixed\_leading}}.

\subsubsection{2. What it uses}\label{what-it-uses-24}

\begin{itemize}
\tightlist
\item
  III.16: Taylor spaces and Wick moments; IV.1: quotient fields,
  extraction and dimension bound.
\item
  V.1: smeared positivity and derivative bounds; VI.1: Ward, rotation,
  track and Casimir laws.
\item
  III.10: Maxwell assembly; III.11: balanced comparison and shell
  estimates.
\item
  VII.5: normalization ratios; VIII.1: data conditions.
\item
  Standard results: Wick/Fock factorization, Gaussian reflection
  positivity, Cartan's formula, approximate identities, rank--nullity,
  compactness and Fourier inversion.
\end{itemize}

\subsubsection{3. Strategy}\label{strategy-24}

Strict reflected positivity identifies quotient fields with on-shell
polynomials. Local flux gives the Ward identities; finite-dimensional
classification gives the conserved stress line. Stations produce
comparison isomorphisms. Transporting derivatives and recentering
tensors transfers divergence rank to the interacting table. Conservation
determines its stress line, and translation Ward identities fix its
normalization.

\subsubsection{4. Steps}\label{steps-24}

\begin{enumerate}
\def\labelenumi{\arabic{enumi}.}
\tightlist
\item
  Factor reflected Wick pairings through a positive-definite Gram matrix
  and the Fock inner product; invertibility excludes residual
  polynomials.
  \href{https://github.com/yang-mills-group/OpenAI-YM-20261028/blob/3456e5e75cf8d4cd558955a5c7638fba2d68f464/YM/Next/MAXWELL/Core.lean\#L44-L92}{\texttt{core}}.
\item
  Smooth the Green kernel, integrate by parts twice and remove the
  approximate identity to obtain \(-\mathcal L_vf\); differentiate the
  source for higher probes.
  \href{https://github.com/yang-mills-group/OpenAI-YM-20261028/blob/3456e5e75cf8d4cd558955a5c7638fba2d68f464/YM/Next/MAXWELL/Contact.lean\#L137-L156}{\texttt{local\_force}},
  \href{https://github.com/yang-mills-group/OpenAI-YM-20261028/blob/3456e5e75cf8d4cd558955a5c7638fba2d68f464/YM/Next/MAXWELL/Force.lean\#L131-L186}{\texttt{jet\_force}}.
\item
  Split Wick contractions by the destinations of the stress legs.
  Homogeneity and source-free conservation cancel the unmixed terms.
  \href{https://github.com/yang-mills-group/OpenAI-YM-20261028/blob/3456e5e75cf8d4cd558955a5c7638fba2d68f464/YM/Next/MAXWELL/Taxes.lean\#L115-L172}{\texttt{fees}}.
\item
  Classify three tensors and at most two vectors; two independent
  propagator tests give divergence rank two.
  \href{https://github.com/yang-mills-group/OpenAI-YM-20261028/blob/3456e5e75cf8d4cd558955a5c7638fba2d68f464/YM/UVNext/FREE_STRESS_CERTIFICATE/Certificate.lean\#L43-L47}{\texttt{rank\_on}}.
\item
  Bounded-degree spectators detect arbitrary weights; dilation and
  inversion eliminate derivative candidates after subtracting a scalar
  derivative. Rotational Ward identities and quarter turns eliminate
  alternating internal rows, completing Maxwell data with the Lorentz
  and Casimir inputs.
  \href{https://github.com/yang-mills-group/OpenAI-YM-20261028/blob/3456e5e75cf8d4cd558955a5c7638fba2d68f464/YM/Next/MAXWELL/Scenes.lean\#L125-L132}{\texttt{scanner}},
  \href{https://github.com/yang-mills-group/OpenAI-YM-20261028/blob/3456e5e75cf8d4cd558955a5c7638fba2d68f464/YM/Next/MAXWELL/AwayTeam.lean\#L48-L118}{\texttt{erasing}}.
\item
  Balanced strip-volume and grid-interpolation errors give stations.
  Compactness of bounded lifts and ordered smeared convergence produce
  continuum panels.
  \href{https://github.com/yang-mills-group/OpenAI-YM-20261028/blob/3456e5e75cf8d4cd558955a5c7638fba2d68f464/YM/Next/UV_STRESS/Glide.lean\#L102-L129}{\texttt{cutoff\_balance}},
  \href{https://github.com/yang-mills-group/OpenAI-YM-20261028/blob/3456e5e75cf8d4cd558955a5c7638fba2d68f464/YM/Next/UV_STRESS/Station.lean\#L132-L180}{\texttt{harvested}}.
\item
  Correct station panels to preserve inclusions; separating tests give
  injectivity, and the dimension bound gives isomorphisms.
  \href{https://github.com/yang-mills-group/OpenAI-YM-20261028/blob/3456e5e75cf8d4cd558955a5c7638fba2d68f464/YM/Next/UV_STRESS/Invert.lean\#L132-L181}{\texttt{Stations.equivalences}}.
\item
  Integration by parts gives
  \(r_kU_{k,d+1}^{-1}W_aU_{k,d}\to\partial_a^{\rm free}\); gradient
  closure and Fourier inversion give smooth convergence. Recentring then
  permits rank transfer.
  \href{https://github.com/yang-mills-group/OpenAI-YM-20261028/blob/3456e5e75cf8d4cd558955a5c7638fba2d68f464/YM/Next/UV_STRESS/Focus.lean\#L183-L196}{\texttt{travels\_jet}},
  \href{https://github.com/yang-mills-group/OpenAI-YM-20261028/blob/3456e5e75cf8d4cd558955a5c7638fba2d68f464/YM/Next/UV_STRESS/Dimension.lean\#L165-L195}{\texttt{differential\_rank}}.
\item
  Compactness on a normalized stress sphere selects a conserved limit;
  (FD) identifies it with \(T_0\) after one sign choice. Writing
  \(c(J-\langle J\rangle\mathbf1)=J_0\), comparing Ward rows gives
  \(c/\xi_k\to\kappa\) and the physical limit.
  \href{https://github.com/yang-mills-group/OpenAI-YM-20261028/blob/3456e5e75cf8d4cd558955a5c7638fba2d68f464/YM/Next/UV_STRESS/Limit.lean\#L129-L167}{\texttt{produce\_current}},
  \href{https://github.com/yang-mills-group/OpenAI-YM-20261028/blob/3456e5e75cf8d4cd558955a5c7638fba2d68f464/YM/Pert/Weak.lean\#L118-L144}{\texttt{meter\_physical}}.
\end{enumerate}

\subsubsection{5. Where the difficulty
lies}\label{where-the-difficulty-lies-24}

Exact isomorphisms and recentering are needed for rank transfer.
Bounded-degree detection and alternating Ward-row rigidity require more
than positivity. Conservation fixes a line; Ward normalization fixes its
scale and sign.

\subsubsection{6. Notation}\label{notation-24}

\(\theta_a,\vartheta_a\): reflection in \(y_a=0\) and reflected
conjugation. A probe of order \(j\) is a real linear functional on the
homogeneous order-\(j\) on-shell Taylor space. \(\mathcal F_d^S\):
degree-\(d\) fields. \(T_0,\kappa\): lighthouse stress and Ward
constant. \(R,\rho\): free rotations and their generators.
\(\mathrm{div}J=\sum_aW_aJ_{ab}\). (FD): onto divergence,
one-dimensional kernel, nonzero centred conserved stress. \(K_{R,n}\):
configurations of total norm at most \(R\) and mutual sup-distance at
least \(R^{-1}\). Scaling: isomorphisms \(U_{k,d}\), strictly decreasing
\(r_k\to0\), \(\xi_k>0\); approaching means, for each fixed field tuple,
the rescaled interacting/free correlation difference tends to zero
uniformly against smooth complex tests supported in any fixed compact
subset of \(\mathrm{Conf}_n\) with sup norm at most one;
\(\mathrm{Conf}_n\) denotes distinct-point configurations; stable means
eventual compatibility with reality, signed symmetries, units and
inclusions.

\subsubsection{7. What remains
unexplained}\label{what-remains-unexplained-24}

Subsidiary Gram/Fock, invariant-spanning, propagator-table,
bounded-order detection, averaging, Lorentz and Casimir proofs remain
inputs. III.11 must supply balanced estimates and shell bounds;
underlying track Ward estimates also remain inputs. Lean's
mixed-correlation theorem specializes to VI.1's assembled stress.

Detailed notes: the full account of this unit, with proof sketches for
every step, is on the
\href{https://github.com/scottnarmstrong/yang-mills-informalization/blob/main/outline/notes/VI-2.md}{detailed
notes page}.

\subsection{VII.1. Short distances: convergence to free Maxwell
fields}\label{vii.1.-short-distances-convergence-to-free-maxwell-fields}

\textbf{Part:} VII, identification as Yang--Mills. \textbf{Lean:} the
chain of
\href{https://github.com/yang-mills-group/OpenAI-YM-20261028/tree/3456e5e75cf8d4cd558955a5c7638fba2d68f464/YM/Pert}{\texttt{YM/Pert}}
(folders \texttt{Fourier}, \texttt{Align}, \texttt{Warp}, \texttt{Rig},
the jet stencils of \texttt{Lattice}, the statement
\texttt{Position/Crest.finite\_leading}, and \texttt{Dense}),
\href{https://github.com/yang-mills-group/OpenAI-YM-20261028/tree/3456e5e75cf8d4cd558955a5c7638fba2d68f464/YM/Push/FINITE_DETECTION}{\texttt{YM/Push/FINITE\_DETECTION}},
\href{https://github.com/yang-mills-group/OpenAI-YM-20261028/tree/3456e5e75cf8d4cd558955a5c7638fba2d68f464/YM/SharpGram/Weak}{\texttt{YM/SharpGram/Weak}},
\href{https://github.com/yang-mills-group/OpenAI-YM-20261028/tree/3456e5e75cf8d4cd558955a5c7638fba2d68f464/YM/SharpGram/Dense}{\texttt{YM/SharpGram/Dense}},
\href{https://github.com/yang-mills-group/OpenAI-YM-20261028/tree/3456e5e75cf8d4cd558955a5c7638fba2d68f464/YM/RG08}{\texttt{YM/RG08}}
(about thirty small sub-folders),
\href{https://github.com/yang-mills-group/OpenAI-YM-20261028/tree/3456e5e75cf8d4cd558955a5c7638fba2d68f464/YM/PertExtension}{\texttt{YM/PertExtension}},
\href{https://github.com/yang-mills-group/OpenAI-YM-20261028/tree/3456e5e75cf8d4cd558955a5c7638fba2d68f464/YM/OPE/Levy}{\texttt{YM/OPE/Levy}}.
\textbf{Size:} 8,117 lines.

\subsubsection{1. What this unit proves}\label{what-this-unit-proves-25}

Fix a compact matrix gauge group \(G\) with real simple nonabelian Lie
algebra \(\mathfrak g\), invariant inner product and orthonormal colour
basis. Let \(V_d\) be complex \(\mathrm{Ad}(G)\)-invariant radial-jet
polynomials of engineering weight at most \(d\), \(F_d\) their free-null
quotient, \(\rho\) restriction to Maxwell jets,
\(\mathcal I_d=V_d\cap\ker\rho\), and
\(\omega_d=\rho\circ\mathrm{sec}_d\) for the VII.3 representatives.

\textbf{Proposition VII.1.1 (regulator removal).} For every integer
\(L\ge2\), \(d,r\in\mathbb N\), label frames, \(p\in F_d^r\) and
continuous \(f\) with compact closed support in \(\mathrm{Conf}_r\), set
\(h_m=L^{-m}\), \(N_m=4L^{m^2+m+1}\). Then
\[\Pi_m(f;p)\longrightarrow\int f(x)\mathcal W_x(\omega_dp_1,\ldots,\omega_dp_r)\,dx.\]
For any position board with smooth compactly supported profiles in
\(\mathrm{Conf}\), order-zero readings converge slotwise for every
basis-label tuple to free readings. \emph{Status:} theorem. \emph{Lean:}
\href{https://github.com/yang-mills-group/OpenAI-YM-20261028/blob/3456e5e75cf8d4cd558955a5c7638fba2d68f464/YM/Pert/Rig/Linked.lean\#L59-L69}{\texttt{Rig.zerothProfile}},
\href{https://github.com/yang-mills-group/OpenAI-YM-20261028/blob/3456e5e75cf8d4cd558955a5c7638fba2d68f464/YM/Pert/Position/Crest.lean\#L13-L16}{\texttt{Position.finite\_leading}}.

\textbf{Proposition VII.1.2 (pointwise limit).} For every itinerary
\((N_i,h_i)\) with \(4\mid N_i\), \(d,r\in\mathbb N\),
\(z\in\mathrm{Conf}_r\), sites \(o_{i,j}\) homing to \(z_j\), and
\(p\in F_d^r\),
\[\Big\langle\prod_jR_{h_i}(\mathrm{sec}_dp_j,o_{i,j})\Big\rangle^{(0)}_{N_i}\longrightarrow\mathcal W_z(\omega_dp_1,\ldots,\omega_dp_r).\]
\emph{Status:} theorem. \emph{Lean:}
\href{https://github.com/yang-mills-group/OpenAI-YM-20261028/blob/3456e5e75cf8d4cd558955a5c7638fba2d68f464/YM/Pert/Warp/Horizons.lean\#L65-L68}{\texttt{Warp.row\_limit}}.

\textbf{Proposition VII.1.3 (corrected covariance).} For every such
itinerary, any index set \(\iota\), injective \(z:\iota\to\mathbb R^4\),
homing sites \(o_{i,j}\), indices \(j,l\) and keys \(a,b\),
\[\mathrm{Cov}_{\mu_{N_i}}(\ell^{h_i}_{o_{i,j},a},\ell^{h_i}_{o_{i,l},b})-\delta_{jl}S_{h_i}(a,b)\longrightarrow\begin{cases}K_{ab}(z_j-z_l),&j\ne l,\\0,&j=l.\end{cases}\]
\emph{Status:} theorem. \emph{Lean:}
\href{https://github.com/yang-mills-group/OpenAI-YM-20261028/blob/3456e5e75cf8d4cd558955a5c7638fba2d68f464/YM/Pert/Align/Drivers.lean\#L39-L44}{\texttt{Align.corrected\_limit}}.

\textbf{Proposition VII.1.4 (detection).} For every
\(v\in V_d\setminus\mathcal I_d\),
\[\mathrm{Re}\,\mathcal W_{(e_{\mathrm t},0)}(\rho v,\rho(\sigma\bar v))>0.\]
More generally, for the diagonalizing and reflected probe substitutions
\(\Phi,\Psi\), every nonzero shell polynomial \(q\) satisfies
\(\mathrm{Re}\,\mathcal W_{(e_{\mathrm t},0)}(\Phi q,\Psi\bar q)>0\).
Here \(\Phi\) is onto and \(\rho v=\Phi q\) implies
\(\rho(\sigma\bar v)=\Psi\bar q\). \emph{Status:} theorem. \emph{Lean:}
\href{https://github.com/yang-mills-group/OpenAI-YM-20261028/blob/3456e5e75cf8d4cd558955a5c7638fba2d68f464/YM/Push/FINITE_DETECTION/TwoPoint.lean\#L151-L155}{\texttt{Free.two\_detect}},
\href{https://github.com/yang-mills-group/OpenAI-YM-20261028/blob/3456e5e75cf8d4cd558955a5c7638fba2d68f464/YM/Push/FINITE_DETECTION/TwoPoint.lean\#L132-L149}{\texttt{Free.positive}}.

\textbf{Proposition VII.1.5 (Kelvin formulas).} For every real symmetric
\(j\)-tensor \(A\) with harmonic diagonal, and \(x\ne0\),
\[\frac{A(\partial,\ldots,\partial)\Gamma(x)}{j!}=\frac{(-2)^j}{4\pi^2}|x|^{-2(j+1)}A(x,\ldots,x).\]
For every \(\mathfrak g\)-valued Maxwell shell tensor
\(Y\in\mathrm{Sh}_j\), contracting the curvature kernel's first
rank-\(j\) Taylor leg with \(Y\), summing colours, ordered form indices
and jet indices without \(1/2\), gives \(\gamma_j\iota^*Y^\natural(x)\)
for \(x\ne0\), where \(\iota(x)=-x/|x|^2\),
\(Y^\natural(x)=Y(x,\ldots,x)\) and \(\gamma_j=2^j(j+1)(j+2)/\pi^2\).
\emph{Status:} theorem. \emph{Lean:}
\href{https://github.com/yang-mills-group/OpenAI-YM-20261028/blob/3456e5e75cf8d4cd558955a5c7638fba2d68f464/YM/Push/FINITE_DETECTION/Harmonic.lean\#L157-L180}{\texttt{Kelvin.harmonic\_Green}},
\href{https://github.com/yang-mills-group/OpenAI-YM-20261028/blob/3456e5e75cf8d4cd558955a5c7638fba2d68f464/YM/Push/FINITE_DETECTION/Response.lean\#L172-L176}{\texttt{Kelvin.response\_eq}}.

\textbf{Proposition VII.1.6 (corrected charts).} Fix a realized Wilson
family, construction record with limit table \(S\) and extraction order
\(f\), VII.4 OPE data, \(D\in\mathbb N\), and a nontrivial filter \(u\)
on \(\mathbb N^2\) finer than \(f\)'s extraction filter. Assume a weak
measuring
\href{https://github.com/yang-mills-group/OpenAI-YM-20261028/blob/3456e5e75cf8d4cd558955a5c7638fba2d68f464/YM/SharpGram/Weak/Reading.lean\#L35-L49}{\texttt{Portfolio}}
through \(D\): observable/Haar bounds, Haar chart ratio, three price
panels and a dominating label-cost constant; \(g(k)\to\infty\) with
eventually \(g(k)\le k\) and
\[\mathrm{scale}(k)\le\mathrm{scale}(g(k))+\log\hat\beta_k;\]
finite-regulator complex linear label maps eventually satisfying
\href{https://github.com/yang-mills-group/OpenAI-YM-20261028/blob/3456e5e75cf8d4cd558955a5c7638fba2d68f464/YM/MergeProof/East/Plume.lean\#L23-L27}{\texttt{Poised}}
simultaneously through \(D\) at step \(g(k)\) (degree availability,
monitor bound, tested equivariant maps and all panel estimates); and
rescaled read-outs converging along \(u\) to \(\Phi_{k,d}\) for every
\(d\le D\) and reference label. Set
\(D_{k,d}=\pi_d^S\Phi_{k,d}\lambda_d\), where \(\pi_d^S\) projects to
\(S_d\) and \(\lambda_d\) is the fixed equivariant representative of the
free quotient in reference label coordinates, and let
\(\chi_k(d)=D_{k,d}\) when bijective, a fixed isomorphism otherwise.
There exist complex linear isomorphisms \(p_k(d):F_d\to S_d\) for all
\(k,d\), eventually in \(k\), simultaneously for \(d\le D\),
unit-preserving, commuting with signed permutations and conjugation, and
satisfying
\[p_k(e)\iota^F_{d,e}=\iota^S_{d,e}p_k(d)\quad(d\le e\le D);\] and
\(\chi_k(d)^{-1}p_k(d)v\to v\) for \(d\le D\), \(v\in F_d\).
\emph{Status:} conditional theorem. \emph{Lean:}
\href{https://github.com/yang-mills-group/OpenAI-YM-20261028/blob/3456e5e75cf8d4cd558955a5c7638fba2d68f464/YM/SharpGram/Weak/Compatible.lean\#L140-L159}{\texttt{Weak.Portfolio.tune\_grades}}.

\textbf{Proposition VII.1.7 (windows and expansions).}

\begin{enumerate}
\def\labelenumi{(\alph{enumi})}
\item
  Fix III.2 RG constants, Haar bounds, a Haar chart ratio satisfying the
  local ratio bounds, and admissible price-panel relay with activation
  bounds \(u_j\le L_j\), eventual activation of every degree, and all
  comparison/analytic rank bounds. For any wait \(U\) with
  \(w_j(t)=\lfloor\log(\bar\beta+c_Gt)/\log L_j\rfloor\), \(c_G>0\),
  there is a wait \(W\) with no smaller stage-length/entry bounds and
  the same window, and panels and dominating costs at every degree.
  Every realized Wilson family satisfying \(W\), prescribed observable
  bounds and polynomial-representative conditions then, for every \(D\),
  has simultaneous finite-regulator comparison data through \(D\) at all
  sufficiently late steps \(k\) whose scale exceeds the stage entry
  threshold and whose backward window has length at most \(k\) and lies
  in one stage; returned panels satisfy all three comparison
  inequalities. \emph{Lean:}
  \href{https://github.com/yang-mills-group/OpenAI-YM-20261028/blob/3456e5e75cf8d4cd558955a5c7638fba2d68f464/YM/SharpGram/Dense/Slender.lean\#L127-L143}{\texttt{Dense.slim\_wait\_for}}.
\item
  If \(U\) has this window and \(V\) has no smaller
  stage-length/entry/window bounds, each \(V\) entry at stage \(j\)
  exceeds or equals \(U\)'s entry at \(j+1\), and
  \(\bar\beta+c_Gt\ge c_G\) above that \(V\) entry, every trajectory
  satisfying \(V\) admits \(g(k)\to\infty\) eventually satisfying the
  scale inequality in VII.1.6, \(g(k)\le k\), and \(g(k)\) above its
  \(U\) entry threshold with backward \(U\) window of length at most
  \(g(k)\) contained in one stage. \emph{Lean:}
  \href{https://github.com/yang-mills-group/OpenAI-YM-20261028/blob/3456e5e75cf8d4cd558955a5c7638fba2d68f464/YM/SharpGram/Dense/Cover.lean\#L155-L165}{\texttt{Dense.nearby}}.
\item
  Fix finite \(I\), \(T\subset I\), \(p\in T\), \(D,k,E,q\in\mathbb N\),
  \(R\in\mathbb R\), \(v>0\), \(q\le k+E\). There exist \(0<w\le v\),
  \(C_0\ge0\), independent of layout and bound \(B\), such that every
  layout satisfying the prescribed tensor symmetry, rank and analytic
  realization conditions
  (\href{https://github.com/yang-mills-group/OpenAI-YM-20261028/blob/3456e5e75cf8d4cd558955a5c7638fba2d68f464/YM/RG08/Creek/Step.lean\#L16-L20}{\texttt{Sense}}),
  depending only on occupied markers, zero on the empty set, with
  weights \(a_i\in\mathbb N\) summing to \(k\) on \(T\) and anchored
  \(v\)-weighted absolute-sum coefficient bounds \(B\) through order
  \(E\) at every leg rank, has this expansion: re-root at \(x_p\), shift
  order by \(q\), zero terms \(q+n<k\); through leg rank \(k+E-q\) the
  coefficients are Taylor coefficients of a finite sum of curvature
  expressions of ranks \(\max(k-q,0),\ldots,k+E-q\). These depend only
  on \(T\), retain symmetry and analytic realization, and have anchored
  absolute-sum bound \(BC_0\) with weight \(\exp(w\,\mathrm{girth})\),
  including relative marker positions. \emph{Lean:}
  \href{https://github.com/yang-mills-group/OpenAI-YM-20261028/blob/3456e5e75cf8d4cd558955a5c7638fba2d68f464/YM/RG08/Dart/Seams.lean\#L34-L45}{\texttt{Dart.coin\_pre}}.
  \emph{Status (a--c):} theorems, modulo imported station and
  curvature-lifting estimates whose full proofs are outside this
  verification (probable; see §7).
\end{enumerate}

\subsubsection{2. What it uses}\label{what-it-uses-25}

\begin{itemize}
\tightlist
\item
  I.1/II.1: lattice, symmetries, labels and jets; III.16: Maxwell shells
  and kernel; VI.1/VI.2: hierarchies and free quotient; VII.3: lattice
  prescription.
\item
  III.2/II.2: RG scales and price panels; III.11: rank-comparison
  estimates; VII.4/VII.5/VIII.1: OPE data, chart correction,
  construction record.
\item
  Wick's theorem, Fourier analysis, Riemann sums, dominated convergence,
  Hobson's formula and Kelvin inversion; finite-dimensional fitting
  (VI.1).
\end{itemize}

\subsubsection{3. Strategy}\label{strategy-25}

At order zero the normalized expectation is Gaussian. Independent
imaginary Gaussians subtract the infinite-lattice self-contractions at
each insertion. Fourier analysis gives Maxwell pairings between distinct
insertions and zero limits for corrected self-pairings. The common Wick
deletion recursion then gives polynomial moments; compactness and
dominated convergence give tested limits.

Translated inversion produces a triangular change of probe variables
with diagonal positive reflected pairings. Wick matchings give strict
positivity on nonzero shell polynomials. Interacting rank estimates
transfer detection to raw charts; fixed equivariant retractions remove
their filtration defects recursively.

\subsubsection{4. Steps}\label{steps-25}

\begin{enumerate}
\def\labelenumi{\arabic{enumi}.}
\tightlist
\item
  Invert the gauge-fixed form: covariance is direct torus Green function
  minus its reflected image.
  \href{https://github.com/yang-mills-group/OpenAI-YM-20261028/blob/3456e5e75cf8d4cd558955a5c7638fba2d68f464/YM/Pert/Fourier/Source.lean\#L55-L59}{\texttt{Orb.covar\_link}}.
\item
  Fourier annuli give separated limits; fixed-offset quadrature controls
  self-contractions.
  \href{https://github.com/yang-mills-group/OpenAI-YM-20261028/blob/3456e5e75cf8d4cd558955a5c7638fba2d68f464/YM/Pert/Align/Bilateral.lean\#L75-L83}{\texttt{Align.chart\_limit}},
  \href{https://github.com/yang-mills-group/OpenAI-YM-20261028/blob/3456e5e75cf8d4cd558955a5c7638fba2d68f464/YM/Pert/Align/Bilateral.lean\#L85-L89}{\texttt{Align.chart\_return}}.
\item
  Induct on distinguishable legs using the shared deletion recursion;
  compactness and dominated convergence give profiles.
  \href{https://github.com/yang-mills-group/OpenAI-YM-20261028/blob/3456e5e75cf8d4cd558955a5c7638fba2d68f464/YM/Pert/Warp/Approach.lean\#L19-L32}{\texttt{Warp.earth}},
  \href{https://github.com/yang-mills-group/OpenAI-YM-20261028/blob/3456e5e75cf8d4cd558955a5c7638fba2d68f464/YM/Pert/Rig/Tested.lean\#L23-L52}{\texttt{Rig.grid\_limit}}.
\item
  Kelvin translation makes probes triangular; diagonal positive pairings
  give positive Wick monomial norms and finite detection.
  \href{https://github.com/yang-mills-group/OpenAI-YM-20261028/blob/3456e5e75cf8d4cd558955a5c7638fba2d68f464/YM/Push/FINITE_DETECTION/Variables.lean\#L114-L118}{\texttt{Free.pairing\_var}},
  \href{https://github.com/yang-mills-group/OpenAI-YM-20261028/blob/3456e5e75cf8d4cd558955a5c7638fba2d68f464/YM/Push/FINITE_DETECTION/Test.lean\#L77-L85}{\texttt{free\_injective}}.
\item
  Rank estimates give moving-detector convergence; free detection and
  recursive retractions yield compatible charts.
  \href{https://github.com/yang-mills-group/OpenAI-YM-20261028/blob/3456e5e75cf8d4cd558955a5c7638fba2d68f464/YM/SharpGram/Weak/Reading.lean\#L107-L124}{\texttt{Portfolio.twoC}},
  \href{https://github.com/yang-mills-group/OpenAI-YM-20261028/blob/3456e5e75cf8d4cd558955a5c7638fba2d68f464/YM/SharpGram/Weak/Compatible.lean\#L140-L159}{\texttt{tune\_grades}}.
\item
  Raise entry thresholds so all active estimates fit the logarithmic
  window; near a stage boundary choose the preceding stage's last usable
  step.
  \href{https://github.com/yang-mills-group/OpenAI-YM-20261028/blob/3456e5e75cf8d4cd558955a5c7638fba2d68f464/YM/SharpGram/Dense/Slender.lean\#L127-L143}{\texttt{slim\_wait\_for}},
  \href{https://github.com/yang-mills-group/OpenAI-YM-20261028/blob/3456e5e75cf8d4cd558955a5c7638fba2d68f464/YM/SharpGram/Dense/Cover.lean\#L155-L165}{\texttt{nearby}}.
\item
  Re-root and reduce exponential weights, then lift and subtract
  successive lowest coefficients to obtain local curvature expressions
  with controlled bounds.
  \href{https://github.com/yang-mills-group/OpenAI-YM-20261028/blob/3456e5e75cf8d4cd558955a5c7638fba2d68f464/YM/RG08/Dart/Seams.lean\#L34-L45}{\texttt{coin\_pre}}.
\end{enumerate}

\subsubsection{5. Where the difficulty
lies}\label{where-the-difficulty-lies-25}

Oscillatory Fourier cancellation must survive differentiation. The
chosen volume growth controls the amplified same-insertion errors;
quarter-period centring removes reflected images. Detection needs strict
positivity, including constants. No optimal volume growth rate is
asserted.

\subsubsection{6. Notation}\label{notation-25}

\(\mathrm{Conf}_r\): distinct \(r\)-tuples in \(\mathbb R^4\);
\(\Gamma(x)=(4\pi^2|x|^2)^{-1}\); \(K\): Maxwell jet kernel;
\(\mathcal W\): Wick moments without same-insertion pairings. \(\mu_N\):
Gaussian on reflection-odd link fields; \(\ell^h\): linear radial jet;
\(S_h\): infinite-lattice self-covariance;
\(R_h(p,o)=\mathbb E_\xi p(\ell_o^h+i\xi)\) with independent
\(\xi\sim\mathcal N(0,S_h)\)
(\href{https://github.com/yang-mills-group/OpenAI-YM-20261028/blob/3456e5e75cf8d4cd558955a5c7638fba2d68f464/YM/Pert/Lattice/Blank.lean\#L41}{\texttt{column}});
\(\langle\cdot\rangle^{(0)}\): coupling-series constant coefficient.
\(\Pi_m\): centred \(h_m^{4r}\)-weighted sum of \(f\) times the
order-zero product correlation. An itinerary has even \(N_i\),
\(h_i>0\), positive integer \(N_ih_i\to\infty\), \(h_i\to0\),
\(N_ih_i^n\to\infty\) for every \(n\in\mathbb N\); homing means
\(h_i(o_i-c_{N_i})\to z\), \(c_N=(N/4,\ldots,N/4)\), using centred
lifts. \(\sigma\): time reflection; \(e_{\mathrm t}\): time unit vector;
\(\mathrm{Sh}_j\): symmetric Taylor/antisymmetric form tensors with
homogeneous fields solving \(dF=d{*}F=0\); Taylor coefficients include
\(1/j!\). A key \(a=(j,\alpha)\) is a coordinate on
\(\mathfrak g\)-valued radial-gauge homogeneous potentials of degree
\(j+1\), with engineering weight \(j+2\); \(K_{ab}\) uses these
coordinates after inverse curl from the shell. A bar denotes complex
conjugation; \(\iota^*\): 2-form pullback; \(\iota^{F,S}\): filtration
inclusions.

\subsubsection{7. What remains
unexplained}\label{what-remains-unexplained-25}

Lower Fourier uniformity estimates remain unchecked (probable). Station
and curvature-lifting estimates are imported (probable); remaining
RG/perturbative-extension packages are unexplained (tentative). The
interacting short-distance theorem belongs to VIII.1; this unit gives no
all-orders comparison.

Detailed notes: the full account of this unit, with proof sketches for
every step, is on the
\href{https://github.com/scottnarmstrong/yang-mills-informalization/blob/main/outline/notes/VII-1.md}{detailed
notes page}.

\subsection{VII.2. The trace of the stress tensor and the trace
anomaly}\label{vii.2.-the-trace-of-the-stress-tensor-and-the-trace-anomaly}

\textbf{Part:} VII, identification as Yang--Mills. \textbf{Lean:}
\href{https://github.com/yang-mills-group/OpenAI-YM-20261028/tree/3456e5e75cf8d4cd558955a5c7638fba2d68f464/YM/TraceExtension}{\texttt{YM/TraceExtension}}.
\textbf{Size:} 20,681 lines.

\subsubsection{1. What this unit proves}\label{what-this-unit-proves-26}

Unless stated otherwise, fix a final record over a Wilson family \(C\),
extraction order \(f\) and limit table \(S\) (VIII.1), with selected
phase \(\xi_*\) and Ward-normalized stress tensor \(T\) (VI.1). The
compact matrix group \(G\) has simple nonabelian Lie algebra
\(\mathfrak g\). Set
\[\kappa_G=\frac{11C_A}{24\pi^2},\qquad d_2=-\frac{\kappa_G}{4},\qquad A_G=\frac{3\dim\mathfrak g}{\pi^4}.\]
Here \(\mathcal O=\tfrac14\sum_{a,\mu,\nu}:F^a_{\mu\nu}F^a_{\mu\nu}:\)
belongs to the free Maxwell table \(\mathsf F\) with \(\dim\mathfrak g\)
species and propagator \(1/(4\pi^2|x|^2)\).

\textbf{Proposition VII.2.1 (centered trace).} Define
\(\Theta=\sum_\mu T_{\mu\mu}-\langle\sum_\mu T_{\mu\mu}\rangle_{\xi_*}\mathbf1\in F_4\).
It has zero expectation, is real, \(O(4)\)-invariant and fixed by the
time-reflection adjoint; these conclusions also hold when centering in
any phase table respecting \(S\). Its reconstructed label
\(\iota(\Theta)\) is fixed by conjugation \(J\) and every Lorentz
transformation. For every Schwartz test \(h\), its vacuum expectation
vanishes and, on the common invariant field domain,
\[\Phi_\Theta(h)=\Phi_{\sum_\mu T_{\mu\mu}}(h)-\Big\langle\sum_\mu T_{\mu\mu}\Big\rangle_{\xi_*}\Big(\int h\Big)\mathrm{id}.\]
If \(r^8\log^2(1/r)G_\Theta(r)\to A>0\), then \(\iota(\Theta)\ne0\).
\emph{Status:} theorem. \emph{Lean:}
\href{https://github.com/yang-mills-group/OpenAI-YM-20261028/blob/3456e5e75cf8d4cd558955a5c7638fba2d68f464/YM/TraceExtension/WilsonTrace.lean\#L61-L63}{\texttt{pinned\_scalar}},
\href{https://github.com/yang-mills-group/OpenAI-YM-20261028/blob/3456e5e75cf8d4cd558955a5c7638fba2d68f464/YM/TraceExtension/Link/Identity.lean\#L158-L161}{\texttt{center\_field}},
\href{https://github.com/yang-mills-group/OpenAI-YM-20261028/blob/3456e5e75cf8d4cd558955a5c7638fba2d68f464/YM/TraceExtension/Sundial/Grand.lean\#L23-L27}{\texttt{nonzero\_strength}}.

\textbf{Proposition VII.2.2 (spectral representation).} There is a
unique positive Lorentz-invariant measure \(\mu_\Theta\) supported in
\(\bar V_+\) such that:

\begin{itemize}
\tightlist
\item
  Some \(C_\mu\ge0\), \(k\in\mathbb N_0\) satisfy
  \(\mu_\Theta\{p_0\le R\}\le C_\mu R^k\) for every \(R\ge1\).
\item
  For every \(s>0\), \(e^{-sp_0}\mu_\Theta\) is finite and
  \(G_\Theta(s)=\int e^{-sp_0}d\mu_\Theta\).
\item
  For every Schwartz \(\varphi\), the translation spectral measure of
  \(\Phi_\Theta(\varphi)\Omega\) is \(|\hat\varphi|^2\mu_\Theta\); for
  every Schwartz \(\phi\) on \((\mathbb R^4)^2\),
  \(W_2(\Theta,\Theta)(\phi)=\int\widehat\phi^{\rm wave}d\mu_\Theta\).
\end{itemize}

Thus \(G_\Theta\) is real, nonnegative and nonincreasing on
\((0,\infty)\). \emph{Status:} theorem. \emph{Lean:}
\href{https://github.com/yang-mills-group/OpenAI-YM-20261028/blob/3456e5e75cf8d4cd558955a5c7638fba2d68f464/YM/TraceExtension/WilsonTrace.lean\#L99-L119}{\texttt{traceSpectrum\_exists}}.

\textbf{Proposition VII.2.3 (one-loop certificates).} Let \(C\) issue
from run data satisfying the RG construction hypotheses in §2, including
a chosen bounded block-spin run. For every filter \(\mathcal L\) on
\(\mathbb N^2\) along which both spacing and volume indices
\(N,p\to\infty\), there are level-5 lattice work items \(t_{k,N,p}\)
satisfying the lattice certificate for all large \(k\),
\(\mathcal L\)-eventually in \((N,p)\). In particular, their Ward darts
satisfy, for every color \(q\),
\[|\mathrm{midPay}_q(J)+12\kappa_G|\le\hat\beta_k^{-100}.\]
\emph{Status:} theorem. \emph{Lean:}
\href{https://github.com/yang-mills-group/OpenAI-YM-20261028/blob/3456e5e75cf8d4cd558955a5c7638fba2d68f464/YM/TraceExtension/Margin/Stand.lean\#L23-L35}{\texttt{agenda}},
certificate
\href{https://github.com/yang-mills-group/OpenAI-YM-20261028/blob/3456e5e75cf8d4cd558955a5c7638fba2d68f464/YM/TraceExtension/Margin/Ansel.lean\#L30-L40}{\texttt{WorthB}}.

\textbf{Proposition VII.2.4 (field anomaly; gallery property).} Assume
II.1 Wilson consistency, the quantitative RG and VII.6 scaling
hypotheses in §2, auxiliary depths \(n_k\ge2\) eventually, and degree
bounds \(z_d\) covering every generator budget of degree at most \(d\)
and satisfying the degree-\(d\) scaling estimates. Let \(\mathcal L\) be
an ultrafilter finer than the fuse filter of \(f\), with work items
satisfying VII.2.3's certificate.

For every limit table \(S\) along \(f\) with ordered Wilson convergence,
Euclidean hierarchy, point limits, cell hypothesis and corrected rank
estimates, and every field algebra on \(S\) satisfying derivative,
rotation, conservation and Ward laws, there exist coupling and
label-transfer limits \(Z\) along \(\mathcal L\) and a free Maxwell
hierarchy \(\mathsf a\), with
\(Z.\beta(k)=\lim_{\mathcal L}\beta_{N,p}(k)\) and
\(\hat\beta_k/2\le Z.\beta(k)\le2\hat\beta_k\) for every level \(k\).

\begin{enumerate}
\def\labelenumi{(\roman{enumi})}
\item
  For every scaling sequence \((r_i,U_i)\) approaching \(\mathsf a\)
  with VII.6 RG lifts and pointwise lattice limits, base levels
  \(k_i\to\infty\) and selected readout levels \(\ell_i\), put
  \(w_i=Z.\beta(\ell_i)\). Here \(r_i\downarrow0\); the complex-linear
  identifications \(U_i:\mathsf F_d\to F_d\) eventually preserve the
  unit, reality and signed coordinate actions degreewise, with rescaled
  correlations converging to \(\mathsf a\). Then
  \[w_i\to\infty,\qquad \frac{\log(1/r_i)}{w_i}\to\kappa_G^{-1},\qquad r_i^4U_i^{-1}\mathring T_{\mu\nu}\to T^{\rm free}_{\mu\nu}.\]
  For every phase hierarchy \(T'\) respecting \(S\), real sequences
  \(e_i,d_i\) satisfy, eventually,
  \[r_i^4U_i^{-1}\Theta_{T'}=e_i\mathbf1+4d_i\mathcal O,\qquad w_ie_i\to0,\quad w_id_i\to d_2.\]
  Thus \(w_ir_i^4U_i^{-1}\Theta_{T'}\to-\kappa_G\mathcal O\).
\item
  Every sequence \(s_j>0\) tending to zero has a subsequence that is the
  scale sequence of a scaling sequence satisfying (i). \emph{Status:}
  theorem; hypotheses discharged in III.2 for the constructed family,
  recorded in VIII.1. \emph{Lean:}
  \href{https://github.com/yang-mills-group/OpenAI-YM-20261028/blob/3456e5e75cf8d4cd558955a5c7638fba2d68f464/YM/TraceExtension/Chart/Line.lean\#L49-L56}{\texttt{gang}},
  \href{https://github.com/yang-mills-group/OpenAI-YM-20261028/blob/3456e5e75cf8d4cd558955a5c7638fba2d68f464/YM/TraceExtension/Chart/Time.lean\#L63-L68}{\texttt{Store\_l}}.
\end{enumerate}

\textbf{Proposition VII.2.5 (short-distance and spectral laws).} If the
gallery property holds for \((C,f)\), then
\[\lim_{r\downarrow0}r^8\log^2(1/r)G_\Theta(r)=A_G>0\] in \(\mathbb C\).
For every continuous compactly supported or Schwartz \(\varphi\), as
\(\rho\to\infty\),
\[\rho^{-8}\log^2\rho\int\varphi(p/\rho)d\mu_\Theta(p)\to\frac{A_G}{1536\pi}\int_{V_+}(p^2)^2\varphi(p)d^4p.\]
The limiting measure has Laplace transform \(A_Gs^{-8}\). For every
Schwartz \(\phi\) on \((\mathbb R^4)^2\), setting
\(\phi_\rho(x,y)=\rho^8\phi(\rho x,\rho y)\), the same cone integral
with \(\varphi=\widehat\phi^{\rm wave}\) is the limit of
\(\rho^{-8}\log^2\rho\,W_2(\Theta,\Theta)(\phi_\rho)\). \emph{Status:}
theorem under the gallery property. \emph{Lean:}
\href{https://github.com/yang-mills-group/OpenAI-YM-20261028/blob/3456e5e75cf8d4cd558955a5c7638fba2d68f464/YM/TraceExtension/Verse/Cape.lean\#L77-L88}{\texttt{outside}},
\href{https://github.com/yang-mills-group/OpenAI-YM-20261028/blob/3456e5e75cf8d4cd558955a5c7638fba2d68f464/YM/TraceExtension/LogScaling.lean\#L216-L222}{\texttt{joint\_limit}}.

\textbf{Proposition VII.2.6 (uniqueness of stress).} Let
\(\mathcal K(S)\) consist of real, symmetric, \(S\)-centered degree-4
tensors covariant under signed coordinate permutations; let
\(\mathcal V(S)\) consist of real degree-5 vectors with that covariance.
Fix a derivative plan satisfying derivative laws and the orientation
condition. The free Maxwell divergence
\(\mathcal K(\mathsf F)\to\mathcal V(\mathsf F)\) is onto with
one-dimensional kernel. In VII.2.4(i)'s setting,
\(\dim\mathcal K(S)=3\), \(\dim\mathcal V(S)=2\), and
\(\dim\ker\mathrm{div}=1\). \emph{Status:} theorem. \emph{Lean:}
\href{https://github.com/yang-mills-group/OpenAI-YM-20261028/blob/3456e5e75cf8d4cd558955a5c7638fba2d68f464/YM/Further/PHYSICAL_PERIOD_LIMIT_F933/FreeRanks.lean\#L159-L165}{\texttt{free\_rank}},
\href{https://github.com/yang-mills-group/OpenAI-YM-20261028/blob/3456e5e75cf8d4cd558955a5c7638fba2d68f464/YM/TraceExtension/Chart/Masts.lean\#L41-L51}{\texttt{ranking}}.

\textbf{Proposition VII.2.7 (no finite Wick profile).} Suppose
\(r^8\log^2(1/r)G_\Theta(r)\to A>0\). For every finite
\(W(r)=\sum_\alpha c_\alpha\prod_\beta K_{m_{\alpha\beta},q_{\alpha\beta},h_{\alpha\beta}}(r)\),
with complex coefficients and
\[K_{m,q,h}(r)=\int_{\mathbb R^3}\frac{E_m(p)^qh(p)}{2E_m(p)}e^{-rE_m(p)}dp,\qquad E_m(p)=\sqrt{m^2+|p|^2},\]
where \(m\ge0\), \(q\in\mathbb N_0\) and continuous
\(h:\mathbb R^3\to\mathbb R\) is positively homogeneous of nonnegative
integer degree, every \(\delta>0\) admits \(r\in(0,\delta)\) with
\(G_\Theta(r)\ne W(r)\). No such collection reproduces the Wightman
function by its mass-shell expression on all Schwartz product tests
\(\varphi\otimes\psi\) with compactly supported wave transform. This
also excludes Wick contractions of two finite polynomial fields from any
free-mode palette. \emph{Status:} theorem. \emph{Lean:}
\href{https://github.com/yang-mills-group/OpenAI-YM-20261028/blob/3456e5e75cf8d4cd558955a5c7638fba2d68f464/YM/TraceExtension/Wick/Refute.lean\#L86-L105}{\texttt{polynomial\_excludes}}.

\textbf{Proposition VII.2.8 (free-photon Fock exclusion).} Under
VII.2.7's positive law, for every \(m\in\mathbb N_0\), let
\(\mathcal F_m\) be the Fock space of \(m\) free-photon species. For
every complex-linear isometry \(U:\mathcal H\to\mathcal F_m\) preserving
the vacuum, taking values in physical transverse waves and intertwining
translations with \(e^{ib\cdot p_{\rm tot}}\), there exist a compactly
supported smooth scalar test \(\varphi\) and compactly supported smooth
real three-vector tests \(h_c\), \(1\le c\le m\), each causally
separated from \(\varphi\), such that
\[\langle W(-h)\Omega_{\mathcal F},U\Phi_\Theta(\varphi)\Omega\rangle\ne\langle U\Phi_\Theta(\bar\varphi)\Omega,W(h)\Omega_{\mathcal F}\rangle.\]
Equality for all such tests would follow from relative locality with the
free Weyl operators. The conclusion persists after applying any bounded
self-adjoint operator fixing \(U\mathcal H\) to both coherent vectors,
or averaging over any probability family of inner-product-preserving
maps on physical waves fixing \(U\mathcal H\). In the averaged version,
for every \(r>0\), tests can be supported in open \(Q,W\) with
\(\bar Q\subset B(0,r)\) and \(\bar Q-\bar W\) spacelike, with
\(\varphi\) supported in \(Q\) and every component of \(h\) in \(W\).
\emph{Status:} theorem. \emph{Lean:}
\href{https://github.com/yang-mills-group/OpenAI-YM-20261028/blob/3456e5e75cf8d4cd558955a5c7638fba2d68f464/YM/TraceExtension/Fock/Conclude.lean\#L46-L49}{\texttt{excludes}},
\href{https://github.com/yang-mills-group/OpenAI-YM-20261028/blob/3456e5e75cf8d4cd558955a5c7638fba2d68f464/YM/TraceExtension/Fock/Regions.lean\#L36-L38}{\texttt{regions\_excludes}}.

\textbf{Proposition VII.2.9 (position-space coefficient).} For every
simple \(G\), kernel parameter \(b>0\) and extraction order \(p\), there
exist a consistent Wilson family \(C\) with parameter \(b\) and
refinement \(f\) of \(p\), admitting final records. For every such final
record, there exist a free hierarchy, coupling limits \(Z\) and an
ultrafilter admitting degree-4 position-space frames. For every
admissible pair \(P(s),P'(s):\mathsf F_4\to F_4\) with window cutoff
\(\chi\), real \(\kappa,d,e\) satisfy, for sufficiently small \(s>0\),
\[P'(s)^{-1}(s^4T_{\mu\nu})=\kappa(s)T^{\rm free}_{\mu\nu}+d(s)\delta_{\mu\nu}\mathcal O+e(s)\delta_{\mu\nu}\mathbf1,\qquad\kappa(s)\to1.\]
Admissibility requires position-space Gram normalization, branch/angular
alignment and convergence of prescribed finite off-diagonal jets for
\(P\); \(P'\) positively rescales its centered scalar sector and
complement, fixes the unit, and normalizes two shell tests to free
values. For every eventually positive \(g\) with
\(\log(1/s)g(s)^2\to\kappa_G^{-1}\), \(d(s)/g(s)^2\to d_2\). The
test-defined coupling satisfies \(g_X/g_W\to1\) and
\(d-d_2g_X^2=O(g_X^4)\), where \(g_W\) is the Wilson coupling from
\(Z\). Eventually
\(s^4\Theta=4d(s)\mathrm{center}_{\xi_*}(P'(s)\mathcal O)\) exactly.
\emph{Status:} theorem for the supplementary construction's family and
refinement. \emph{Lean:}
\href{https://github.com/yang-mills-group/OpenAI-YM-20261028/blob/3456e5e75cf8d4cd558955a5c7638fba2d68f464/YM/TraceExtension/Sundial/Noon.lean\#L50-L83}{\texttt{meterQ}},
\href{https://github.com/yang-mills-group/OpenAI-YM-20261028/blob/3456e5e75cf8d4cd558955a5c7638fba2d68f464/YM/TraceExtension/Endpoint.lean\#L37-L64}{\texttt{FieldComparison}}.

\subsubsection{2. What it uses}\label{what-it-uses-26}

\begin{itemize}
\tightlist
\item
  I.1: positive
  \href{https://github.com/yang-mills-group/OpenAI-YM-20261028/blob/3456e5e75cf8d4cd558955a5c7638fba2d68f464/YM/Gauge/Colors.lean\#L184}{adjoint
  Casimir}; II.2, III.9: schedule
  \(\hat\beta_k=\bar\beta+\kappa_G\log(1/a_k)\).
\item
  III.4--III.15: bounded block-spin run, eventual-level estimates,
  stationary criticality, fourth-order classical matching, determinant
  budgets, cluster/polymer, support and small/large-field bounds; III.7:
  \href{https://github.com/yang-mills-group/OpenAI-YM-20261028/blob/3456e5e75cf8d4cd558955a5c7638fba2d68f464/YM/Next/STATIONARY/OneLoop/MomentumIntegral.lean\#L346-L408}{logarithmic
  one-loop momentum integral}. III.2 supplies these run hypotheses.
\item
  IV.1: extraction, point limits, cell hypothesis and chessboard tests;
  V.1--V.3:
  \href{https://github.com/yang-mills-group/OpenAI-YM-20261028/blob/3456e5e75cf8d4cd558955a5c7638fba2d68f464/YM/Spec/Closed.lean\#L64-L73}{phases
  and reconstruction}; VI.1: field, rotation, derivative, orientation
  and Ward laws.
\item
  VI.2, VII.1, VII.3: free Maxwell calculus; VII.6:
  \href{https://github.com/yang-mills-group/OpenAI-YM-20261028/blob/3456e5e75cf8d4cd558955a5c7638fba2d68f464/YM/TraceExtension/Prom/Zoom/Running.lean\#L25-L31}{scaling
  sequences}, RG lifts and pointwise limits along subsequences of every
  positive null sequence.
\item
  Spectral theorem, Källén--Lehmann representation, holomorphic identity
  theorem, Laplace continuity/Tauberian results, coherent states and
  point-supported distributions (Reed--Simon II; Glimm--Jaffe, ch.~19;
  Feller II, §XIII.1--5; Widder).
\end{itemize}

\subsubsection{3. Strategy}\label{strategy-26}

The classical stress tensor is traceless; the anomaly comes from the
running coupling. In Balaban's RG framework, transport a lattice Ward
identity through block-spin steps and extract its coefficient from the
background-field one-loop logarithmic pole. The canonical tensor has a
divergence diagnostic tending to \(-\kappa_G\). Conservation forces the
pure-trace improvement coefficient to \(d_2\); rank-one uniqueness and
Ward normalization identify the conserved candidate with \(T\).

Free-field convergence then gives the two-point law. The factor
\((\log(1/r_i)/w_i)^2\to\kappa_G^{-2}\) cancels the squared anomaly
coefficient, leaving \(A_G\). Monotonicity removes lattice rounding;
subsequences give the full limit. Laplace convergence gives the spectral
law. Finite Wick expansions and relatively local photon Fock
realizations each force incompatible short-distance growth.

\subsubsection{4. Steps}\label{steps-26}

\begin{enumerate}
\def\labelenumi{\arabic{enumi}.}
\item
  A sound Ward dart represents the action variation along a displacement
  field through fourth order. Under the criticality, branch, bound,
  positive price-gap and classical-matching hypotheses, block-spin
  transport preserves soundness; repeated transport contracts precision
  errors by \(1/2\) per step.
  \href{https://github.com/yang-mills-group/OpenAI-YM-20261028/blob/3456e5e75cf8d4cd558955a5c7638fba2d68f464/YM/TraceExtension/Perfect/Badge.lean\#L110-L141}{Soundness}.
\item
  Under uniform covariance, resolvent and action bounds, the one-loop
  action regulated at infrared mass \(\mu>0\) has logarithmic
  coefficient \(-12\kappa_G\). A mean-value argument in \(\log\mu\),
  followed by regulated Ward comparison, gives the payoff limit.
  \href{https://github.com/yang-mills-group/OpenAI-YM-20261028/blob/3456e5e75cf8d4cd558955a5c7638fba2d68f464/YM/TraceExtension/Loft/Awake.lean\#L32-L53}{\texttt{awake}}.
\item
  Under VII.2.4's hypotheses, lattice-to-continuum transfer preserves
  the divergence diagnostic. Solving the two-dimensional improvement
  equation forces its scalar coefficient to \(d_2\).
  \href{https://github.com/yang-mills-group/OpenAI-YM-20261028/blob/3456e5e75cf8d4cd558955a5c7638fba2d68f464/YM/TraceExtension/Chart/Limits.lean\#L68-L80}{\texttt{quill}}.
\item
  Rank-one conservation identifies the candidate up to a multiple;
  comparison of Ward fluxes makes that multiple tend to one.
  \href{https://github.com/yang-mills-group/OpenAI-YM-20261028/blob/3456e5e75cf8d4cd558955a5c7638fba2d68f464/YM/TraceExtension/Chart/North.lean\#L28-L65}{\texttt{year\_spec}}.
\item
  Free Wick normalization gives \(A_G\) along scaling sequences;
  monotonicity removes lattice rounding and the subsequence principle
  gives the full limit.
  \href{https://github.com/yang-mills-group/OpenAI-YM-20261028/blob/3456e5e75cf8d4cd558955a5c7638fba2d68f464/YM/TraceExtension/Verse/Cape.lean\#L77-L88}{\texttt{outside}}.
\item
  Lorentz invariance converts the law to cone Laplace convergence;
  uniform energy bounds extend compact-test convergence to Schwartz
  tests.
  \href{https://github.com/yang-mills-group/OpenAI-YM-20261028/blob/3456e5e75cf8d4cd558955a5c7638fba2d68f464/YM/TraceExtension/LogScaling.lean\#L216-L222}{\texttt{joint\_limit}}.
\item
  Under the positive trace law, finite Wick power--log expansions cannot
  produce an inverse logarithm squared. Relative locality in photon Fock
  space forces point-supported Cauchy data and then momentum-doubling
  eigencomponents, yielding energy growth at most \(R^7\), incompatible
  with the spectral law.
  \href{https://github.com/yang-mills-group/OpenAI-YM-20261028/blob/3456e5e75cf8d4cd558955a5c7638fba2d68f464/YM/TraceExtension/Wick/Log.lean\#L96-L108}{\texttt{no\_powerlog}},
  \href{https://github.com/yang-mills-group/OpenAI-YM-20261028/blob/3456e5e75cf8d4cd558955a5c7638fba2d68f464/YM/TraceExtension/Fock/Conclude.lean\#L33-L44}{\texttt{saturate}}.
\item
  For the supplementary position-space frames, shell normalization gives
  \(d=d_2g_X^2\kappa\) and \(\kappa-1=-dQ/B_0\), where \(B_0\ne0\) and
  \(Q\) has a finite limit. Hence \(\kappa-1=O(g_X^2)\) and
  \(d-d_2g_X^2=O(g_X^4)\).
  \href{https://github.com/yang-mills-group/OpenAI-YM-20261028/blob/3456e5e75cf8d4cd558955a5c7638fba2d68f464/YM/TraceExtension/Sundial/Noon.lean\#L50-L83}{\texttt{meterQ}}.
\end{enumerate}

\subsubsection{5. Where the difficulty
lies}\label{where-the-difficulty-lies-26}

The Ward transport, exact one-loop coefficient and remainders
\(o(\hat\beta^{-p})\) for every \(p\in\mathbb N_0\) carry the anomaly
proof. Slower coefficient errors must also defeat growing improvement
coefficients. The amplitude uses equality of the independently entering
schedule and one-loop slopes. The Wick exclusion concerns finite
profiles; continuous-mass generalized free fields are outside its scope.
Fock exclusion additionally needs point-supported Cauchy data to force
dilation identities.

\subsubsection{6. Notation}\label{notation-26}

\(F_d\): degree-\(d\) labels; \(\mathring T\): \(S\)-centered stress;
\(G_\Theta(r)=S^{\xi_*}(\Theta(-re_0/2)\Theta(re_0/2))\).
\(p^2=p_0^2-|\vec p|^2\); \(V_+\) is the open forward cone;
\(\widehat\phi^{\rm wave}(p)=\int e^{ip\cdot(y-x)}\phi(x,y)d^4x\,d^4y\).
\(T^{\rm free}\) is the traceless Maxwell stress tensor; \(g_W,g_X\) are
the Wilson and position-space couplings. \(\mathrm{midPay}_q\) is the
real coefficient selector on the Haar divergence of the origin-rooted
displacement field
(\href{https://github.com/yang-mills-group/OpenAI-YM-20261028/blob/3456e5e75cf8d4cd558955a5c7638fba2d68f464/YM/TraceExtension/Loft/Leaf.lean\#L38}{definition}).

For \(h=(h_c^a)\), the coherent vector \(W(h)\Omega_{\mathcal F}\) has
\(n\)-particle component
\(e^{-\|\eta\|^2/2}\mathrm{i}^n\eta^{\otimes n}/\sqrt{n!}\), where
\[\eta_c^a(k)=|k|\sum_b\left(\delta_{ab}-\frac{k_ak_b}{|k|^2}\right)\widehat h_c^b(|k|,k),\qquad d\nu(k)=\frac{d^3k}{2|k|}.\]

The exact run predicates for VII.2.3 are \texttt{Bolt},
\texttt{Bandage}, \texttt{Nest}, \texttt{Vigil}, \texttt{Acclaim},
\texttt{supportQ}; for VII.2.4 they are \texttt{PostT},
\texttt{Stock.Tabs}, \texttt{Scrim}, \texttt{Sceptre},
\texttt{supportQ}, \texttt{Moonlit}, \texttt{auraQ}, together with
\texttt{Binding}, \texttt{HaarBudget}, \texttt{ChartRatio}, generator
bounds and \texttt{FiresSky}.

\subsubsection{7. What remains
unexplained}\label{what-remains-unexplained-26}

Full proofs of RG contraction, covariance/resolvent bounds and fast
lattice-to-continuum remainders remain inputs. So do the single
massive-mode expansion, point-support and analytic uniqueness theorems
for Fock exclusion, and existence and shell-observable formulas for
position-space frames.

Detailed notes: the full account of this unit, with proof sketches for
every step, is on the
\href{https://github.com/scottnarmstrong/yang-mills-informalization/blob/main/outline/notes/VII-2.md}{detailed
notes page}.

\subsection{VII.3. The theory is not free: Gram positivity and exclusion
of free-field
profiles}\label{vii.3.-the-theory-is-not-free-gram-positivity-and-exclusion-of-free-field-profiles}

\textbf{Part:} VII, identification as Yang--Mills. \textbf{Lean:}
\href{https://github.com/yang-mills-group/OpenAI-YM-20261028/tree/3456e5e75cf8d4cd558955a5c7638fba2d68f464/YM/Extra/FREE_DETECTOR_ALGEBRA}{\texttt{YM/Extra/FREE\_DETECTOR\_ALGEBRA}},
\href{https://github.com/yang-mills-group/OpenAI-YM-20261028/tree/3456e5e75cf8d4cd558955a5c7638fba2d68f464/YM/Pert/Position}{\texttt{YM/Pert/Position}},
\href{https://github.com/yang-mills-group/OpenAI-YM-20261028/tree/3456e5e75cf8d4cd558955a5c7638fba2d68f464/YM/Pert/Rigid}{\texttt{YM/Pert/Rigid}},
\href{https://github.com/yang-mills-group/OpenAI-YM-20261028/tree/3456e5e75cf8d4cd558955a5c7638fba2d68f464/YM/Pert/Lattice}{\texttt{YM/Pert/Lattice}},
and a few declarations of
\href{https://github.com/yang-mills-group/OpenAI-YM-20261028/tree/3456e5e75cf8d4cd558955a5c7638fba2d68f464/YM/Push/LOOP_COEFFICIENT}{\texttt{YM/Push/LOOP\_COEFFICIENT}}.
Proposition VII.3.7 is in
\href{https://github.com/yang-mills-group/OpenAI-YM-20261028/tree/3456e5e75cf8d4cd558955a5c7638fba2d68f464/YM/TraceExtension/Wick}{\texttt{YM/TraceExtension/Wick}}
and
\href{https://github.com/yang-mills-group/OpenAI-YM-20261028/blob/3456e5e75cf8d4cd558955a5c7638fba2d68f464/TraceBridge.lean}{\texttt{TraceBridge.lean}}.
\textbf{Size:} about 3,800 declaration lines on the main dependency path
in the listed folders, excluding the trace-exclusion files.

\subsubsection{1. What this unit proves}\label{what-this-unit-proves-27}

All points lie in \(\mathbb R^4\). Let \(\mathfrak g\) have its
invariant inner product; \(\mathfrak g=0\) is allowed in VII.3.1--5.
Free moments are Wick-ordered Maxwell correlations with covariance built
from \(\Gamma(z)=(4\pi^2|z|^2)^{-1}\). Conjugation bars conjugate
polynomial coefficients, and orthogonal transformations act on every
tensor index. For \(x\ne y\), set
\(\langle P,Q\rangle_{x,y}=\langle P(x)(\mathcal R_{x-y}\bar Q)(y)\rangle_0\).

\textbf{Proposition VII.3.1 (strict positivity).} For every \(x\ne y\),
this pairing on \(\mathcal W\) is Hermitian, with real self-pairings
positive exactly for \(P\ne0\). For every finite complex-linearly
independent family \((V_i)\), \((\langle V_j,V_i\rangle_{x,y})_{ij}\) is
positive definite, without degree restriction. \emph{Status:} theorem;
only strict positivity is used downstream. \emph{Lean:}
\href{https://github.com/yang-mills-group/OpenAI-YM-20261028/blob/3456e5e75cf8d4cd558955a5c7638fba2d68f464/YM/Extra/FREE_DETECTOR_ALGEBRA/WickPositive.lean\#L47-L53}{\texttt{reflected\_positive}},
\href{https://github.com/yang-mills-group/OpenAI-YM-20261028/blob/3456e5e75cf8d4cd558955a5c7638fba2d68f464/YM/Extra/FREE_DETECTOR_ALGEBRA/WickPositive.lean\#L57-L76}{\texttt{reflected\_symm}},
\href{https://github.com/yang-mills-group/OpenAI-YM-20261028/blob/3456e5e75cf8d4cd558955a5c7638fba2d68f464/YM/Extra/FREE_DETECTOR_ALGEBRA/WickPositive.lean\#L80-L97}{\texttt{reflected\_pd}}.

\textbf{Proposition VII.3.2 (tower Gram form).} For \(n\ge0\),
\(z=x-y\ne0\), \(B_z(p,q)=\langle\ell_p,\ell_q\rangle_{x,y}\) is
symmetric positive definite on \(\mathrm{Tow}_n\). A block-unitriangular
automorphism \(\sigma_z\) satisfies, for all \(p,q\),

\[B_z(\sigma_zp,\sigma_zq)=\sum_{j<n}\frac{(j+1)(j+2)2^j}{\pi^2|z|^{2j+4}}\langle p_j,q_j\rangle.\]

For any \(b_1,\ldots,b_N\in\mathrm{Tow}_n\) with
\(B_z(b_i,b_k)=d_i\delta_{ik}\), \(d_i\in\mathbb R\), let
\(\Phi_b:\mathbb C[Y_1,\ldots,Y_N]\to\mathcal W\) be the algebra
homomorphism with \(\Phi_b(Y_i)=\ell_{b_i}\). For all complex
polynomials \(p,q\),

\[\langle\Phi_bp,\Phi_bq\rangle_{x,y}=\sum_\alpha p_\alpha\overline{q_\alpha}\prod_i\alpha_i!d_i^{\alpha_i}.\]

\emph{Status:} theorem. \emph{Lean:}
\href{https://github.com/yang-mills-group/OpenAI-YM-20261028/blob/3456e5e75cf8d4cd558955a5c7638fba2d68f464/YM/Extra/FREE_DETECTOR_ALGEBRA/PositiveLeg.lean\#L126-L133}{\texttt{shiftAll\_gram}},
\href{https://github.com/yang-mills-group/OpenAI-YM-20261028/blob/3456e5e75cf8d4cd558955a5c7638fba2d68f464/YM/Extra/FREE_DETECTOR_ALGEBRA/PositiveLeg.lean\#L150-L154}{\texttt{towerB\_pos}},
\href{https://github.com/yang-mills-group/OpenAI-YM-20261028/blob/3456e5e75cf8d4cd558955a5c7638fba2d68f464/YM/Extra/FREE_DETECTOR_ALGEBRA/PowerGram.lean\#L100-L116}{\texttt{duplex\_script}}.

\textbf{Proposition VII.3.3 (Wick recursion).} For every \(n\ge1\),
\(x\in\mathrm{Conf}_n\), \(W\in\mathcal W^n\), insertion \(i\) and slot
\(s\),

\[\langle\cdots(W_iX_s)(x_i)\cdots\rangle_0=\sum_{k\ne i}\langle\cdots W_i(x_i)\cdots(\partial_s^{x_i,x_k}W_k)(x_k)\cdots\rangle_0,\]

where \(\partial_s^{y,y'}\) is the complex-linear derivation with
\(\partial_s^{y,y'}X_t=C_{st}(y-y')\). For every \(n\ge0\) and
\(x\in\mathrm{Conf}_n\), this recursion and \(T(1,\ldots,1)=1\) uniquely
determine the complex-multilinear free moment \(T\). Moments are
\(O(4)\)-covariant; conjugating all polynomials conjugates their moment.
\emph{Status:} theorem. \emph{Lean:}
\href{https://github.com/yang-mills-group/OpenAI-YM-20261028/blob/3456e5e75cf8d4cd558955a5c7638fba2d68f464/YM/Extra/FREE_DETECTOR_ALGEBRA/Removal.lean\#L91-L108}{\texttt{Pair.remove}},
\href{https://github.com/yang-mills-group/OpenAI-YM-20261028/blob/3456e5e75cf8d4cd558955a5c7638fba2d68f464/YM/Extra/FREE_DETECTOR_ALGEBRA/Gaussian.lean\#L24-L62}{\texttt{byAnnihilators}}.

\textbf{Proposition VII.3.4 (faithfulness).} For every \(d\ge0\),
\(P\in\mathrm{sh}_d\) vanishes if all its correlations with arbitrary
finite families \(Q_i\in\mathrm{sh}_{d_i}\), \(d_i\ge0\), at distinct
points vanish. Thus null labels are exactly \(\ker\omega_d\), and
\(\omega_d\) induces \(E_d\cong\mathrm{sh}_d\), with injection
\(\varpi_d:E_d\to\mathcal W\). For every basis \((b_i)\) of \(E_d\) and
\(x\ne y\), \(v\mapsto(\langle\varpi_dv(x)\varpi_db_i(y)\rangle_0)_i\)
is injective. For any Wilson family, any table on its register, and any
pair with \(x_\mu\ne y_\mu\) in every coordinate, these readings give a
column system: a conjugation-fixed complex basis with unit first, from
frame labels with recorded weight and length. \emph{Status:} theorem.
\emph{Lean:}
\href{https://github.com/yang-mills-group/OpenAI-YM-20261028/blob/3456e5e75cf8d4cd558955a5c7638fba2d68f464/YM/Extra/FREE_DETECTOR_ALGEBRA/Rank.lean\#L38-L54}{\texttt{no\_residual}},
\href{https://github.com/yang-mills-group/OpenAI-YM-20261028/blob/3456e5e75cf8d4cd558955a5c7638fba2d68f464/YM/Extra/FREE_DETECTOR_ALGEBRA/Rank.lean\#L64-L65}{\texttt{exactField}},
\href{https://github.com/yang-mills-group/OpenAI-YM-20261028/blob/3456e5e75cf8d4cd558955a5c7638fba2d68f464/YM/Extra/FREE_DETECTOR_ALGEBRA/Columns.lean\#L103-L115}{\texttt{scanner\_inj}},
\href{https://github.com/yang-mills-group/OpenAI-YM-20261028/blob/3456e5e75cf8d4cd558955a5c7638fba2d68f464/YM/Extra/FREE_DETECTOR_ALGEBRA/Vane.lean\#L25-L38}{\texttt{vane}}.

\textbf{Proposition VII.3.5 (covariance).} For every gauge group above,
free-table correlations respect signed coordinate permutations and
conjugation. For every \(d\), \(O(4)\) acts complex-linearly on \(E_d\),
with \(\varpi_d(Rv)=R\varpi_d(v)\). \emph{Status:} theorem. \emph{Lean:}
\href{https://github.com/yang-mills-group/OpenAI-YM-20261028/blob/3456e5e75cf8d4cd558955a5c7638fba2d68f464/YM/Extra/FREE_DETECTOR_ALGEBRA/LabelTurn.lean\#L80-L91}{\texttt{table\_spin}},
\href{https://github.com/yang-mills-group/OpenAI-YM-20261028/blob/3456e5e75cf8d4cd558955a5c7638fba2d68f464/YM/Extra/FREE_DETECTOR_ALGEBRA/Rank.lean\#L106-L110}{\texttt{rideO}}.

\textbf{Proposition VII.3.6 (rigidity and position equations).} A buffer
\(B\in\mathrm{Buf}_d\) is complex-linear on \(E_d\), kills \(1\),
preserves every lower-degree subspace, and commutes with conjugation and
signed permutations. Buffers form a finite-dimensional real space. A
buffer is stationary when inserting \(B\) successively in every slot
gives total free correlation zero for every field tuple and
distinct-point configuration.

\begin{enumerate}
\def\labelenumi{(\alph{enumi})}
\item
  If \(\mathfrak g\) is simple nonabelian, every stationary degree-four
  buffer vanishes. Indeed \(B=0\) follows from \(\langle Bv\rangle_0=0\)
  and \(\langle(Bv)(e)w(0)\rangle_0+\langle v(e)(Bw)(0)\rangle_0=0\) for
  all \(v,w\in E_4\), where \(e\) is axis \(0\)'s unit vector. Also
  \(\mathrm{Buf}_d=0\) for \(d<4\).
\item
  For any gauge group, if stationary degree-\(d\) buffers vanish,
  finitely many sections \((r_i,p_i,f_i)\), \(r_i\ge0\),
  \(p_i\in E_d^{r_i}\),
  \(f_i\in C_c^\infty(\mathrm{Conf}_{r_i};\mathbb C)\), and real-linear
  \(g:\mathbb R^k\to\mathbb R^{\dim_{\mathbb R}\mathrm{Buf}_d}\) give
  the position chart
  \(\mathrm{chart}_Q(w,B)=g(\mathrm{Re}\,w_i((1+B)p_i))_i\), where
  \(w_i\) are complex-multilinear forms. At the free tested data \(w_Q\)
  and \(B=0\), its buffer derivative is the isomorphism
\end{enumerate}

\[B\longmapsto g\Big(\mathrm{Re}\int f_i(x)\sum_j\langle p_{i1}(x_1)\cdots(Bp_{ij})(x_j)\cdots\rangle_0\,dx\Big)_i.\]

Such systems are called clean and square.

\begin{enumerate}
\def\labelenumi{(\alph{enumi})}
\setcounter{enumi}{2}
\tightlist
\item
  Assume simple nonabelian \(\mathfrak g\), a final Wilson record with
  table \(S\), and a scaling sequence \(X=(r_k,U_{k,d},\xi_k)\):
  \(r_k\downarrow0\) strictly, \(\xi_k>0\), and complex-linear
  isomorphisms \(U_{k,d}:E_d\to F_d(S)\). Assume VI.2's leading
  comparison and eventual frame compatibility with unit, conjugation,
  signed permutations and inclusions. Fix a clean square degree-four
  system \(Q\) and \(\varphi\) satisfying: a local real-analytic
  solution of
  \(\mathrm{chart}_Q(w,\varphi(w))=\mathrm{chart}_Q(w_Q,0)\), with
  \(\varphi(w_Q)=0\) and local uniqueness; \(\pi_Q^{(0)}(L,k)\to w_Q\)
  for every integer \(L\ge2\); for every \(L\ge2\) and \(p,j\ge0\),
  eventually all order-\(j\) coupling derivatives at zero of the
  position equations on the order-\(p\) lattice truncation and its image
  under \(\varphi\) equal those of the constant free value; and buffers
  \(A_k\to0\) with invertible \(1+A_k\), eventually solving the
  equations and satisfying the displayed comparison below with \(B_k\)
  replaced by \(A_k\), for every clean degree-four \(R\) and integer
  \(L\ge2\).
\end{enumerate}

Then \(B_k=\varphi(m_Q(k))\to0\); eventually \(1+B_k\) is invertible and
\(U_{k,4}(1+B_k)\) solves \(Q\) exactly, respects unit, conjugation,
signed permutations and degree-\(\le3\) inclusions, and

\[m_R^{B_k}(k)-\mathrm{chart}_R\bigl(\pi_R^{(0)}(L,k),\varphi(\pi_Q^{(0)}(L,k))\bigr)\to0.\]

\emph{Status:} (a,b) theorems; (c) conditional theorem, hypotheses
discharged in VI.2 and VIII.1. \emph{Lean:}
\href{https://github.com/yang-mills-group/OpenAI-YM-20261028/blob/3456e5e75cf8d4cd558955a5c7638fba2d68f464/YM/Pert/Rigid/Camera.lean\#L107-L113}{\texttt{Rigid.frozen\_four}},
\href{https://github.com/yang-mills-group/OpenAI-YM-20261028/blob/3456e5e75cf8d4cd558955a5c7638fba2d68f464/YM/Pert/Position/Detect.lean\#L131-L147}{\texttt{square\_of\_frozen}},
\href{https://github.com/yang-mills-group/OpenAI-YM-20261028/blob/3456e5e75cf8d4cd558955a5c7638fba2d68f464/YM/Pert/Position/Four.lean\#L32-L43}{\texttt{residency}}.

\textbf{Proposition VII.3.7 (finite Wick profiles).} A mode has mass
\(\mathsf m\ge0\), integers \(j,q\ge0\), continuous
\(h:\mathbb R^3\to\mathbb R\) with \(h(tp)=t^jh(p)\) for every \(t>0\)
and \(p\). Its kernel, for \(r>0\), is

\[K_{\mathsf m,q,h}(r)=\int_{\mathbb R^3}\frac{E^qh(p)e^{-rE}}{2E}\,dp,\qquad E=\sqrt{\mathsf m^2+|p|^2}.\]

A Wick profile is a finite complex linear combination of finite products
of these kernels, including empty products. If
\(r^8\log^2(1/r)G_\Theta(r)\to A>0\) as \(r\to0^+\), every profile \(W\)
differs from \(G_\Theta\) somewhere in every \((0,\delta)\),
\(\delta>0\). In the selected reconstructed theory, every corresponding
finite sum of mass-shell contractions differs from the Wightman
two-point function on some Schwartz product \(\bar f\otimes g\) with
compactly supported \(\overline{\widehat f}\widehat g\), using
positive-sign Fourier transforms. Contractions test this product at sums
of positive-energy on-shell momenta, with densities \(E^qh/(2E)\).
\emph{Status:} theorem, assigned to VII.2 and VIII.1; VII.2 supplies
\(A=3\dim\mathfrak g/\pi^4\). \emph{Lean:}
\href{https://github.com/yang-mills-group/OpenAI-YM-20261028/blob/3456e5e75cf8d4cd558955a5c7638fba2d68f464/YM/TraceExtension/Wick/Refute.lean\#L75-L80}{\texttt{Wick.euclidean\_excludes}},
\href{https://github.com/yang-mills-group/OpenAI-YM-20261028/blob/3456e5e75cf8d4cd558955a5c7638fba2d68f464/YM/TraceExtension/Wick/Refute.lean\#L86-L105}{\texttt{Wick.polynomial\_excludes}}.

\subsubsection{2. What it uses}\label{what-it-uses-27}

\begin{itemize}
\tightlist
\item
  III.16: Maxwell kernels, jets and response identity; IV.1: tables,
  quotient fields and columns.
\item
  VI.1, VI.2: scaling sequences, leading Maxwell comparison and
  degree-four basis; VII.1: order-zero lattice convergence; VIII.1:
  implicit-position conditions.
\item
  VII.2: trace asymptotic and mode expansions; V.3: positive-energy
  spectral representation; I.1: Wilson lattice action.
\item
  Wick's theorem, Hobson's formula, Maxwell inversion covariance,
  finite-dimensional duality and the analytic implicit function theorem.
\end{itemize}

\subsubsection{3. Strategy}\label{strategy-27}

Kelvin inversion and translation diagonalize the reflected jet
covariance. Wick recursion extends positivity to polynomials, detecting
null labels. Degree-four rigidity then supplies finite position
equations. Independently, finite Wick profiles have power-log expansions
incompatible with the trace's inverse-logarithmic leading term.

\subsubsection{4. Steps}\label{steps-27}

\begin{enumerate}
\def\labelenumi{\arabic{enumi}.}
\tightlist
\item
  Hobson's formula identifies multipole fields with Kelvin transforms.
  Translating at the inverted separation cancels lower Taylor
  coefficients and gives positive Gram pivots.
  \href{https://github.com/yang-mills-group/OpenAI-YM-20261028/blob/3456e5e75cf8d4cd558955a5c7638fba2d68f464/YM/Extra/FREE_DETECTOR_ALGEBRA/Kelvin.lean\#L111-L113}{\texttt{reply\_kelvin}},
  \href{https://github.com/yang-mills-group/OpenAI-YM-20261028/blob/3456e5e75cf8d4cd558955a5c7638fba2d68f464/YM/Extra/FREE_DETECTOR_ALGEBRA/PositiveLeg.lean\#L126-L133}{\texttt{shiftAll\_gram}}.
\item
  Wick recursion gives factorial monomial norms; every polynomial lies
  in a finite tower. Reflected self-tests then prove faithfulness and
  injectivity.
  \href{https://github.com/yang-mills-group/OpenAI-YM-20261028/blob/3456e5e75cf8d4cd558955a5c7638fba2d68f464/YM/Extra/FREE_DETECTOR_ALGEBRA/PowerGram.lean\#L100-L116}{\texttt{duplex\_script}},
  \href{https://github.com/yang-mills-group/OpenAI-YM-20261028/blob/3456e5e75cf8d4cd558955a5c7638fba2d68f464/YM/Extra/FREE_DETECTOR_ALGEBRA/Rank.lean\#L38-L54}{\texttt{no\_residual}}.
\item
  In the 22-dimensional degree-four basis, stationarity forces weighted
  antisymmetry; signed permutations force symmetry or zero. Finite test
  functionals detect all buffers, giving invertible linearization.
  \href{https://github.com/yang-mills-group/OpenAI-YM-20261028/blob/3456e5e75cf8d4cd558955a5c7638fba2d68f464/YM/Pert/Rigid/Camera.lean\#L90-L105}{\texttt{coeff\_barren}},
  \href{https://github.com/yang-mills-group/OpenAI-YM-20261028/blob/3456e5e75cf8d4cd558955a5c7638fba2d68f464/YM/Pert/Position/Detect.lean\#L131-L147}{\texttt{square\_of\_frozen}}.
\item
  Local implicit uniqueness identifies \(B_k\) with \(A_k\) eventually,
  transferring exact position equations and leading comparison.
  \href{https://github.com/yang-mills-group/OpenAI-YM-20261028/blob/3456e5e75cf8d4cd558955a5c7638fba2d68f464/YM/Pert/Position/Four.lean\#L32-L43}{\texttt{residency}}.
\item
  At finite regulator, normalized Gaussian perturbative expectations
  define lattice profiles; separate imaginary Gaussian averages within
  insertions implement Wick subtraction. Their tested Riemann sums give
  the lattice readings.
  \href{https://github.com/yang-mills-group/OpenAI-YM-20261028/blob/3456e5e75cf8d4cd558955a5c7638fba2d68f464/YM/Pert/Lattice/Bundles.lean\#L55-L56}{\texttt{Full.profile}},
  \href{https://github.com/yang-mills-group/OpenAI-YM-20261028/blob/3456e5e75cf8d4cd558955a5c7638fba2d68f464/YM/Pert/Lattice/Blank.lean\#L41}{\texttt{Full.column}}.
\item
  Mode expansions have integer powers and nonnegative logarithmic
  powers. Dominance along \(r=2^{-n}\) excludes the trace law; spectral
  damping transfers tested Wightman equality to Euclidean equality.
  \href{https://github.com/yang-mills-group/OpenAI-YM-20261028/blob/3456e5e75cf8d4cd558955a5c7638fba2d68f464/YM/TraceExtension/Wick/Log.lean\#L96-L108}{\texttt{Wick.no\_powerlog}},
  \href{https://github.com/yang-mills-group/OpenAI-YM-20261028/blob/3456e5e75cf8d4cd558955a5c7638fba2d68f464/YM/TraceExtension/Wick/Refute.lean\#L86-L105}{\texttt{Wick.polynomial\_excludes}}.
\end{enumerate}

\subsubsection{5. Where the difficulty
lies}\label{where-the-difficulty-lies-27}

Simultaneous diagonalization depends on Taylor normalization and
ordered-index contractions. The position comparison is leading-order
only. Non-freeness excludes finite profiles, leaving continuous-mass
generalized free fields and infinite Wick expansions outside its scope.

\subsubsection{6. Notation}\label{notation-27}

\(\mathrm{Conf}_n\): distinct-point configurations; \(\mathrm{Sol}_j\):
real \(\mathfrak g\)-valued homogeneous degree-\(j\) 2-forms satisfying
\(dS=\delta S=0\); \(\mathrm{Tow}_n=\bigoplus_{j<n}\mathrm{Sol}_j\);
\(\mathcal W=\mathbb C[X_s]\): polynomials in basis coordinates
\(s=(j,a)\) of on-shell Taylor coefficients.
\(\ell_p(A)=\sum_j\langle p_j,A_j\rangle\), using Taylor coefficients
\((j!)^{-1}D^j\) and summing over all ordered form and derivative
indices, without a factor \(1/2\); \(\mathcal R_z\): reflection in
\(z^\perp\); \(C_{st}\): jet covariance. \(\omega_d\): on-shell
evaluation of adjoint-invariant radial-potential labels of weight
\(\le d\), with order-\(j\) slots weighted \(j+2\);
\(\mathrm{sh}_d=\mathrm{im}\omega_d\); \(E_d\): free quotient. \(m_Q\),
\(w_Q\), \(\pi_Q^{(0)}\): tested rescaled correlations in \(U_{k,4}\),
free readings, and order-zero lattice readings at spacing \(L^{-k}\),
respectively; \(m_R^B\): measured readings in frames \(U_{k,4}(1+B)\);
\(\mathrm{chart}_Q\): tested readings after \(1+B\). \(G_\Theta(r)\):
centered stress-tensor trace two-point function in the selected phase at
\(\mp(r/2)e_{\rm t}\).

\subsubsection{7. What remains
unexplained}\label{what-remains-unexplained-27}

An intrinsic characterization of \(\mathrm{sh}_d\) needs surjectivity.
Interacting rank estimates belong to III.11 and III.12; continuum Wick
subtraction to VII.1. Further Maxwell-Fock embedding exclusions are
outside VII.3.7.

Detailed notes: the full account of this unit, with proof sketches for
every step, is on the
\href{https://github.com/scottnarmstrong/yang-mills-informalization/blob/main/outline/notes/VII-3.md}{detailed
notes page}.

\subsection{VII.4. Operator product expansion and contact
extensions}\label{vii.4.-operator-product-expansion-and-contact-extensions}

\textbf{Part:} VII, identification as Yang--Mills. \textbf{Lean:}
\href{https://github.com/yang-mills-group/OpenAI-YM-20261028/tree/3456e5e75cf8d4cd558955a5c7638fba2d68f464/YM/OPE}{\texttt{YM/OPE}},
\href{https://github.com/yang-mills-group/OpenAI-YM-20261028/tree/3456e5e75cf8d4cd558955a5c7638fba2d68f464/YM/Contacts}{\texttt{YM/Contacts}}.
\textbf{Size:} about 10,200 lines, 2,035 declarations.

\subsubsection{1. What this unit proves}\label{what-this-unit-proves-28}

For VII.4.1--4 assume (RG): grid approximants, simple gauge Lie algebra,
per-step and cluster RG estimates, and the continuum extraction record.
Fix any coherent prescription and its real covariant Reader on the final
record; VIII.1 supplies these data. Its strata \(E_j\subset F_j\) have
bases \(O_{j,\kappa}\) fixed by conjugation, with coefficients
\(c_{j,\kappa}\).

An \((r,M,\eta)\)-cluster satisfies \(|x_i-o|\le Mr\),
\(|x_i-x_{i'}|\ge\eta r\) for \(i\ne i'\). Spectators are \(c\)-clear if
\(|y_j-o|\ge c\), \(|y_j-y_{j'}|\ge c\) for \(j\ne j'\). Small-radius
thresholds are uniform in configurations and states.

\textbf{Proposition VII.4.1 (separated OPE).} Fix finite label lists
\(a,b\), \(M\ge0\), \(\eta,c,\varepsilon>0\), and \(D\in\mathbb N\). For
all sufficiently small \(r\), every such cluster and clear spectator
configuration gives separated \((x,y)\) and \((o,y)\) and
\[\|\mathcal E_D\|_{\mathcal A}\le r^{D+1-w(a)-\varepsilon},\qquad
\mathcal E_D=M(a,x;b,y)-\sum_{j\le D,\kappa}c_{j,\kappa}(a;x,o)M((j,O_{j,\kappa}),o;b,y).\]
The same bound holds in every phase and the mixed state. For each fixed
channel \((j,\kappa)\), \(|c_{j,\kappa}|\le r^{j-w(a)-\varepsilon}\)
uniformly on clusters for sufficiently small \(r\). Coefficients are
independent of spectators, state and \(D\). \emph{Status:} theorem.
\emph{Lean:}
\href{https://github.com/yang-mills-group/OpenAI-YM-20261028/blob/3456e5e75cf8d4cd558955a5c7638fba2d68f464/YM/OPE/Orbits/Kernel.lean\#L62-L64}{\texttt{Orbits.Foundation}}.

\textbf{Proposition VII.4.2 (integrated remainder).} With VII.4.1's
fixed data, let \(K\) be a measurable space, \(\mu\) a positive measure,
\(q\in L^1(\mu;\mathbb C)\), \(\int|q|\,d\mu\le B\), \(B\ge0\). Let
configurations be a.e.-measurable and satisfy the geometry a.e. on
\(\{q\ne0\}\). For all sufficiently small \(r\), for the mixed state or
any a.e.-measurable phase map, the weighted correlation, truncation and
remainder are integrable, and
\(|\int q\,s(\mathcal E_D)\,d\mu|\le Br^{D+1-w(a)-\varepsilon}\). The
pointwise threshold works for every \(K,B\); every probability mixture
of phases satisfies the remainder bound of VII.4.1. \emph{Status:}
theorem. \emph{Lean:}
\href{https://github.com/yang-mills-group/OpenAI-YM-20261028/blob/3456e5e75cf8d4cd558955a5c7638fba2d68f464/YM/OPE/Orbits/Profile.lean\#L98-L108}{\texttt{estimates\_weighted}}.

\textbf{Proposition VII.4.3 (nested OPE).} Fix a finite rooted tree with
inner root, leaf labels \(v\), output \((m,\kappa)\), inner truncations
\(D_p\), \(M\ge0\), \(\eta,\varepsilon>0\). A plot assigns leaf points
and inner anchors \(z_p\); positive radii satisfy
\(s_{\mathrm{root}}=r\), \(s_p\le s_{\mathrm{parent}(p)}\). Every
subtree lies within \(Ms_p\) of \(z_p\); positions in different children
are separated by \(\eta s_p\). Put
\(\rho_p=s_p/s_{\mathrm{parent}(p)}\), and let \(w_p\) be the degree
below \(p\). Over non-root inner nodes define
\[\mathrm{toll}_m=r^{m-w(v)-\varepsilon}\prod_p\rho_p^{-w_p-\varepsilon},\qquad \mathrm{nest}_D=\sum_p\rho_p^{D_p+1}.\]
Writing
\(c_{m,\kappa}=c_{m,\kappa}(v;z_{\mathrm{leaves}},z_{\mathrm{root}})\),
for all sufficiently small \(r\), \(|c_{m,\kappa}|\le\mathrm{toll}_m\)
and
\(|c_{m,\kappa}-c^{(D)}_{m,\kappa}|\le\mathrm{toll}_m\mathrm{nest}_D\).
Here \(c^{(D)}\) recursively multiplies node coefficients and sums
inner-child channels through \(D_p\), leaving leaves unchanged; the root
output is fixed and \(D_{\mathrm{root}}\) is unused. For fixed
spectators \(b\) and \(\delta>0\), separated mutually and from every
plot point by \(\delta\), the degree-\(m\) remainder has norm, and
absolute value in every state, at most \(\mathrm{toll}_{m+1}\).

For a.e.-measurable plots and spectators satisfying this geometry a.e.
on \(\{q\ne0\}\), VII.4.2's integration conclusion holds: direct and
composed coefficients are integrable after weighting, with
\(|\int qc\,d\mu|\le B\mathrm{toll}_m\) and
\(|\int q(c-c^{(D)})\,d\mu|\le B\mathrm{toll}_m\mathrm{nest}_D\).
Weighted tested correlations, truncations and remainders are integrable,
with remainder bound \(B\mathrm{toll}_{m+1}\) for the mixed state or any
a.e.-measurable phase map. Radii remain fixed during integration;
thresholds are uniform in radii, measures, weights, \(B\) and
configuration maps. \emph{Status:} theorem. \emph{Lean:}
\href{https://github.com/yang-mills-group/OpenAI-YM-20261028/blob/3456e5e75cf8d4cd558955a5c7638fba2d68f464/YM/OPE/Nesting/Endpoint.lean\#L17-L27}{\texttt{Nesting.Uniform}}.

\textbf{Proposition VII.4.4 (contacts).} Fix contact devices sharing the
Reader's prescription and the final record's rotations. For \(k\) labels
\(a_i\in F_{d_i}\), any state \(s\), and \(0<s_0\le B_0/2\),
\(B_0=\min(1,\mathrm{cl}(w(a)))\), the forest kernel defines a
distribution \(\mathcal D_a^s\) of order at most \((2+w(a))k\): on each
compact \(K\), its absolute integral is bounded by
\(C_K\|f\|_{C^{(2+w(a))k}}\). It equals \(\int fS^s(a;x)\,dx\) for tests
supported on distinct configurations. Distributions also exist on slices
\(o+P(E)\), with \(P\) a linear isometry whose range contains a group
diagonal's normal space, and \(o\) on that diagonal.

Shell and iterated forest sums give the pairing, with integrable kernels
and summable \(L^1\) norms, under these stage-family hypotheses: a
partition \(Q\) of \(g\); affine isometric placements based on block
diagonals and containing their normals; positive subtraction lengths
\(s\); \(R_i\in(0,1]\); common ceiling
\(2s\le B\le\min(1,\mathrm{cl}(\sum_i d_i))\). A common frame scales
each block normal by \(R_i\) and preserves its orthogonal complement.
Closed stages are bounded in frame coordinates and within normal
distance \(B\); nonsingleton blocks satisfy \(R_i\le C_1s\). Original
and spectator labels have uniform norm bounds; spectator positions and
labels are continuous on stages and constant in block normal directions.
All nonempty subgroup centroids of different blocks and all spectators
are uniformly separated. Constants may depend on the fixed family.

For \(Q=\{g\}\), \(g\ne\emptyset\), pure normal dilation \(R\in(0,1]\),
and \(f_R(x)=R^{-4(|g|-1)}f(R^{-1}x_{K_g}+x_{K_g^\perp})\), assume
rescaled-stage support and \(\|f\|_{C^N}\le H\),
\(N=(2+\sum_i d_i)(|g|-1)\). For every \(\varepsilon>0\), the absolute
kernel integral is \(\le CHR^{u-w(a_g)-\varepsilon}\): \(u=0,n,n+1\) for
full, weight-\(n\), and degree-\(n\) remainder. The latter two require
stages within normal distance \(\mathrm{cl}(n)\) of the \(g\)-diagonal.
For \(|g|>1\), tests away from its diagonal agree with the expression
retaining proper-subcluster subtractions.

Device strata have bases fixed by conjugation, containing the unit in
degree zero. Every positive sequence \(f_n\) admits \(0<t_n\le f_n\),
\(4t_{n+1}<t_n\), \(t_n\to0\), \(2(n+1)t_n\le\mathrm{cl}(n)\).
\emph{Status:} theorem. \emph{Lean:}
\href{https://github.com/yang-mills-group/OpenAI-YM-20261028/blob/3456e5e75cf8d4cd558955a5c7638fba2d68f464/YM/Contacts/Solstice/Result.lean\#L65-L69}{\texttt{provided}}.

\textbf{Proposition VII.4.5 (slab bound).} Assume construction
properties, ordered generic point extraction, Euclidean hierarchy,
IV.1's corrected field estimates, and the cell hypothesis: reflection
positivity and chessboard inequalities at every cutoff, with centred
one-field bounds \(\le Aa_m^{-d-\varepsilon}+\eta\) late in the
extraction, as in IV.1. For every \(D\in\mathbb N\), \(R>0\), there is
\(C_0>0\), with \(P=4(D+1)\), such that every nonempty word \(w\),
degrees \(\le D\), coefficient norms \(\le R\), \(r\)-buffered in
\((0,d)\), \(d>0\), \(0<r\le1\), admits a level \(b\) and width
\(0<s\le d\), \(s\in a_b\mathbb N\), containing \(w\) in \((0,s)\).
Every positive-time exterior word \(v\) then has \(B_v\) with
\[|S(\Theta v_n\cdot v_n)|\le B_v\rho^{4n}\quad(n\in\mathbb N),\qquad \rho=C_0^{|w|}(|w|/r)^{P|w|},\]
where \(v_n\) precedes \(v\) by \(2n\) consecutive width-\(s\) slabs
holding alternately \(w\) and its reflection. Iterated Schwarz gives
\(\|T_s(w)\|\le\rho\) in V.1. \emph{Status:} theorem. \emph{Lean:}
\href{https://github.com/yang-mills-group/OpenAI-YM-20261028/blob/3456e5e75cf8d4cd558955a5c7638fba2d68f464/YM/Next/PHASE/Seats.lean\#L154-L160}{\texttt{seats}}.

\textbf{Proposition VII.4.6 (prescription).} A probe is a finite
separated token configuration, anchored at its first point. The
correlation maps \(\mathrm{sp}_n\) of anchored fields against probes are
injective and inclusion-compatible. For modes indexing bases of
\(F_n/\iota(F_{n-1})\), \(F_{-1}=0\), there exist vectors
\(u_i\in F_{\deg i}\) and finitely supported complex probe weights
\(\lambda_i\), such that for every \(n\),
\(\lambda_i^{(n)}(\iota u_j)=\delta_{ij}\) for degrees \(\le n\), and
\(\phi=\sum_{\deg i\le n}\lambda_i^{(n)}(\phi)\iota u_i\) for every
\(\phi\in F_n\). Here
\(\lambda_i^{(n)}=\sum_p\lambda_i(p)\mathrm{sp}_n(\cdot)(p)\).
Higher-degree tests vanish on \(F_n\); one prescription serves every
truncation. \emph{Status:} theorem. \emph{Lean:}
\href{https://github.com/yang-mills-group/OpenAI-YM-20261028/blob/3456e5e75cf8d4cd558955a5c7638fba2d68f464/YM/OPE/Splitting/Legacy.lean\#L20-L31}{\texttt{Prescription.chosen}}.

\subsubsection{2. What it uses}\label{what-it-uses-28}

\begin{itemize}
\tightlist
\item
  II.2, III.13: cutoff construction; III.11 and Part III: RG estimates,
  assembled in VIII.1
  (\href{https://github.com/yang-mills-group/OpenAI-YM-20261028/blob/3456e5e75cf8d4cd558955a5c7638fba2d68f464/YM/OPE/Exchange/Arrangement.lean\#L56-L60}{\texttt{Services}}).
\item
  IV.1: corrected fields, extraction, centred chessboard bounds;
  V.1--V.3: ground algebra and phases; VI.1: derivatives, rotations,
  reality.
\item
  Gelfand isometry, reflection positivity, iterated Schwarz, Haar
  averaging, Taylor subtraction and Zimmermann forests.
\end{itemize}

\subsubsection{3. Strategy}\label{strategy-28}

Probes reconstruct finite-degree fields. RG replaces clusters by fields
plus controlled remainders; translated chessboard moments give
ground-algebra norm bounds, hence estimates in every phase. Exact
reconstruction of replacement fields gives the grid OPE. Derivative
induction extends it to arbitrary configurations, and bringing auxiliary
scale heights together gives arbitrary \(\varepsilon\)-loss. Recursive
replacements yield trees; weight-dependent normal Taylor subtraction
makes contact kernels integrable, with partial collisions treated by
finer partitions.

\subsubsection{4. Steps}\label{steps-28}

\begin{enumerate}
\def\labelenumi{\arabic{enumi}.}
\tightlist
\item
  Extend biorthogonal probes degree by degree; average rotations and
  conjugations and remove lower degrees.
  \href{https://github.com/yang-mills-group/OpenAI-YM-20261028/blob/3456e5e75cf8d4cd558955a5c7638fba2d68f464/YM/OPE/Orbits/Cuts.lean\#L17-L23}{\texttt{Reader}}.
\item
  Fuse nearby RG groups, retaining separated groups; obtain a
  degree-\(t\) replacement with remainder gain \((a_m/a_k)^{t+1/16}\).
  \href{https://github.com/yang-mills-group/OpenAI-YM-20261028/blob/3456e5e75cf8d4cd558955a5c7638fba2d68f464/YM/OPE/Prices/Circuit.lean\#L64-L87}{\texttt{funding}}.
\item
  Bound translated moments by chessboard roots; full phase support and
  Gelfand give norm control.
  \href{https://github.com/yang-mills-group/OpenAI-YM-20261028/blob/3456e5e75cf8d4cd558955a5c7638fba2d68f464/YM/OPE/Cards/Pressure.lean\#L98-L104}{\texttt{Cards.carrying}}.
\item
  Descend in degree from crude high-degree bounds; grid distance
  compensates the derivative degree increase.
  \href{https://github.com/yang-mills-group/OpenAI-YM-20261028/blob/3456e5e75cf8d4cd558955a5c7638fba2d68f464/YM/OPE/Interpolate/Corridor.lean\#L34-L47}{\texttt{interpolated}}.
\item
  Telescope replacements over trees, merging children at the parent's
  rate and expanding others to sufficiently high auxiliary degree. This
  controls each node's truncation.
  \href{https://github.com/yang-mills-group/OpenAI-YM-20261028/blob/3456e5e75cf8d4cd558955a5c7638fba2d68f464/YM/OPE/Nesting/Fusion.lean\#L97-L104}{\texttt{Fusion.tariff}}.
\item
  Subtract weight-dependent normal Taylor jets; finer-partition
  induction handles partial collisions and sums dyadic shells. With
  \(d_h=\sum_{i\in h}d_i\), \(q_h=4(|h|-1)\), and
  \(L_h=\max(0,d_h+1-q_h)\), subtraction leaves shell exponent
  \(q_h+L_h-d_h\ge1\).
  \href{https://github.com/yang-mills-group/OpenAI-YM-20261028/blob/3456e5e75cf8d4cd558955a5c7638fba2d68f464/YM/Contacts/Solstice/Skies.lean\#L59-L79}{\texttt{glues\_total}}.
\item
  Apply two-scale chessboard bounds to reflection-positive slab trains;
  choose levels with \(a_k\ge c(r/|w|)^4\) to obtain \(P=4(D+1)\).
  Iterated Schwarz removes the exterior cost.
  \href{https://github.com/yang-mills-group/OpenAI-YM-20261028/blob/3456e5e75cf8d4cd558955a5c7638fba2d68f464/YM/Next/PHASE/Seats.lean\#L119-L151}{\texttt{replicas}}.
\end{enumerate}

\subsubsection{5. Where the difficulty
lies}\label{where-the-difficulty-lies-28}

The \(1/16\) gain margin comes from RG transport. Norm control ensures
spectator and state uniformity; partial collisions require more than
radial integrability. Estimates allow arbitrary \(\varepsilon\)-loss,
without identifying logarithmic corrections.

\subsubsection{6. Notation}\label{notation-28}

\(w(a)=\sum_i d_i\); \(F_d\): degree-\(d\) fields; \(\mathcal A\):
ground algebra; \(M\): compressed word; \(S^s\): correlation for
\(s\in X\sqcup\{\mathrm{mix}\}\), with \(X\) the phase spectrum;
\(a_k\): lattice spacing; \(\mathrm{cl}(n)\): device clearance; \(K_g\):
normal space of the \(g\)-diagonal.

\subsubsection{7. What remains
unexplained}\label{what-remains-unexplained-28}

The internal polymer estimates and their constants, rotated comparisons,
high-horizon tree replacement and partial-collision sector induction
need fuller accounts. Device projections are not identified with Reader
coefficients. Uniqueness, extraction independence, infinite-order
convergence, contact positivity, contact Ward identities and
perturbative agreement are not asserted.

Detailed notes: the full account of this unit, with proof sketches for
every step, is on the
\href{https://github.com/scottnarmstrong/yang-mills-informalization/blob/main/outline/notes/VII-4.md}{detailed
notes page}.

\subsection{VII.5. Asymptotic freedom of the running
coupling}\label{vii.5.-asymptotic-freedom-of-the-running-coupling}

\textbf{Part:} VII, identification as Yang--Mills. \textbf{Lean:}
\href{https://github.com/yang-mills-group/OpenAI-YM-20261028/tree/3456e5e75cf8d4cd558955a5c7638fba2d68f464/YM/Pert/Asymptotics}{\texttt{YM/Pert/Asymptotics}},
\href{https://github.com/yang-mills-group/OpenAI-YM-20261028/tree/3456e5e75cf8d4cd558955a5c7638fba2d68f464/YM/Extra/WEAK_SCALING_ACTLAW}{\texttt{YM/Extra/WEAK\_SCALING\_ACTLAW}},
\href{https://github.com/yang-mills-group/OpenAI-YM-20261028/tree/3456e5e75cf8d4cd558955a5c7638fba2d68f464/YM/TraceExtension/Foundation}{\texttt{YM/TraceExtension/Foundation}},
\href{https://github.com/yang-mills-group/OpenAI-YM-20261028/tree/3456e5e75cf8d4cd558955a5c7638fba2d68f464/YM/TraceExtension/Soot}{\texttt{YM/TraceExtension/Soot}}.
\textbf{Size:} 232 lines, 43 declarations.

\subsubsection{1. What this unit proves}\label{what-this-unit-proves-29}

For VII.5.1--3, let \(G\subset U(q)\) be compact, \(q\ge1\). Write
\(2b_0=\kappa_G=11C_{\mathfrak g}/(24\pi^2)\),
\(a_m=(\prod_{k<m}b_k)^{-1}\), \(s_m=\log(1/a_m)\), with odd
\(b_k\ge3\). The trajectory has inverse couplings \(\beta_{N,p,m}\),
bare endpoint \(\beta_{N,p}=\beta_{N,p,N}\), reference value
\(\beta_{N,p,0}=\bar\beta>0\), and \(g_{N,p,m}=\beta_{N,p,m}^{-1/2}\).
Simple means real simple and nonabelian. The running clause says that
for every \(\varepsilon>0\), eventually in \(m\),
\[|\beta_{N,p,m}-(\bar\beta+2b_0s_m)|\le\varepsilon s_m\] uniformly for
\(N\ge m\) and \(p\in\mathbb N\).

\textbf{Proposition VII.5.1.} If the family satisfies the construction
conditions (II.2), then for every \(\varepsilon>0\) there is \(m_0\ge1\)
such that
\[\left|\beta_{N,p,m}/s_m-2b_0\right|<\varepsilon\qquad(m\ge m_0,\ N\ge m,\ p\in\mathbb N).\]
\emph{Status:} theorem; only the running clause is used, without
simplicity. \emph{Lean:}
\href{https://github.com/yang-mills-group/OpenAI-YM-20261028/blob/3456e5e75cf8d4cd558955a5c7638fba2d68f464/YM/Pert/Asymptotics/RawAxis.lean\#L15-L30}{\texttt{YM.Pert.Asymptotics.zeroUp\_bound}}.

\textbf{Corollary VII.5.2.} Under the construction conditions and
simplicity of \(\mathfrak g\), the Wilson family satisfies (Reg1): for
every \(\varepsilon>0\) there is \(N_0\ge1\) such that
\(|\beta_{N,p}/s_N-b_*|<\varepsilon\) for all \(N\ge N_0\) and
\(p\in\mathbb N\), where \(b_*=2b_0>0\). \emph{Status:} theorem in
VIII.1. \emph{Lean:}
\href{https://github.com/yang-mills-group/OpenAI-YM-20261028/blob/3456e5e75cf8d4cd558955a5c7638fba2d68f464/HierarchyBridge.lean\#L16-L25}{\texttt{YangMills.regular}}.

\textbf{Proposition VII.5.3.} Under the same hypotheses, for any
lattice-level and volume sequences \((N_i),(k_i)\) with \(i\le N_i\)
eventually,
\[g_{N_i,k_i,i}\sqrt{s_i}\longrightarrow(2b_0)^{-1/2},\qquad g_{N_i,k_i,i}^2s_i\longrightarrow(2b_0)^{-1}.\]
\emph{Status:} theorem; the second limit follows by squaring.
\emph{Lean:}
\href{https://github.com/yang-mills-group/OpenAI-YM-20261028/blob/3456e5e75cf8d4cd558955a5c7638fba2d68f464/YM/Pert/Asymptotics/RawAxis.lean\#L36-L48}{\texttt{YM.Pert.Asymptotics.comparison}}.

For VII.5.4--5, \(\mathcal F_d,\mathcal F_d^\circ\) are
finite-dimensional complex label spaces of degree at most \(d\) in
Euclidean hierarchies on \(S\) and its free Maxwell comparison table, on
the same reference labels. Derivative plans have inclusions
\(\iota_{d\to e}\) and derivatives \(\partial_{d,a}\). Exact
compatibility means preserving inclusions, signed permutations
\(\Gamma_4\), conjugation and units.

\textbf{Proposition VII.5.4.} Fix the free Maxwell data of VI.2,
including the finite separating tests and boundary, polynomial and
internal-rotation conditions. Its plan \(W\) satisfies the inclusion,
derivative and orientation axioms; \(J^\circ\) satisfies Ward identities
with translation part \(\kappa\partial^W\) and rotation part \(\kappa\)
times the infinitesimal rotation action, where \(\kappa\ne0\) and
\(\partial^W_4\ne0\). Let \(V\) satisfy the inclusion and derivative
axioms. Suppose \(J\in\mathcal F_4^{4\times4}\) is symmetric, real,
\(\Gamma_4\)-covariant and satisfies \(\sum_a\partial^V_{4,a}J_{ab}=0\)
for every \(b\), and \((J,p,Y)\) satisfies the local Ward identities. If
some \(c\in\mathbb R\) satisfies \(p_{4,a}=c\partial^V_{4,a}\) for all
\(a=1,\ldots,4\), there exist \(D,N,M\in\mathbb N\), \(R_0\ge1\) such
that \(c/\xi_k\to\kappa\) for every tracking scaling sequence at every
level \((n,d,j,R)\ge(N,D,M,R_0)\) componentwise (VI.1). Also, for any
hierarchy and local Ward identities, \(p\ne0\) implies \(J\ne0\). A
scaling sequence has complex-linear isomorphisms
\(U_k(d):\mathcal F_d^\circ\to\mathcal F_d\), positive strictly
decreasing \(r_k\to0\) and \(\xi_k>0\). Tracking \((V,W,J)\) at
\((n,d,j,R)\) requires exact compatibility through \(\max(d,4)\),
convergence of transported correlations at \(r_kx\) for each fixed list
of at most \(n\) fields of degrees at most \(d\), to the free
correlations, with all Fréchet derivatives through order \(j\),
uniformly on
\(\{x:\|x\|\le R,\ \|x_a-x_b\|_\infty\ge R^{-1}\ (a\ne b)\}\), and
\(\xi_k^{-1}r_k^4U_k(4)^{-1}J_{ab}\to J^\circ_{ab}\).

\emph{Status:} theorems. \emph{Lean:}
\href{https://github.com/yang-mills-group/OpenAI-YM-20261028/blob/3456e5e75cf8d4cd558955a5c7638fba2d68f464/YM/Extra/WEAK_SCALING_ACTLAW/Orientation.lean\#L43-L54}{\texttt{YM.Extra.WeakScaling.tracks\_ratio}},
\href{https://github.com/yang-mills-group/OpenAI-YM-20261028/blob/3456e5e75cf8d4cd558955a5c7638fba2d68f464/YM/Extra/WEAK_SCALING_ACTLAW/Criteria.lean\#L64-L70}{\texttt{flux\_nontrivial}}.

\textbf{Proposition VII.5.5.} Let \(V,W\) satisfy the inclusion and
orientation axioms. For every \(d\) there is a complex-linear,
\(\Gamma_4\)-equivariant, conjugation-preserving retraction \(P_d\) of
\(\iota^W_d:=\iota^W_{d\to d+1}\). For any complex-linear isomorphisms
\(A_d:\mathcal F_d^\circ\to\mathcal F_d\), define
\[X_0=A_0,\qquad X_{d+1}=\iota^V_dX_dP_d+A_{d+1}(1-\iota^W_dP_d).\] Fix
\(D\in\mathbb N\). If \(A_d\) preserve symmetries and units for
\(d\le D\), then \(X_d\) are exactly compatible through degree \(D\).
For sequences of such isomorphisms \(A_d^{(k)}\) satisfying
\((A_{d+1}^{(k)})^{-1}\iota^V_dA_d^{(k)}v\to\iota^W_dv\) for every
\(d<D,v\), one has \((A_d^{(k)})^{-1}X_d^{(k)}\to1\) in operator norm
for every \(d\le D\). These \(X_d^{(k)}\) are eventually bijective;
replacing nonbijective terms by \(A_d^{(k)}\) gives isomorphisms
\(\widetilde X_d^{(k)}\) with
\((A_d^{(k)})^{-1}\widetilde X_d^{(k)}v\to v\). If the original frames
preserve symmetries and units through \(D\), the replacements are
eventually exactly compatible simultaneously through \(D\).
\emph{Status:} theorem. \emph{Lean:}
\href{https://github.com/yang-mills-group/OpenAI-YM-20261028/blob/3456e5e75cf8d4cd558955a5c7638fba2d68f464/YM/Extra/WEAK_SCALING_ACTLAW/Adjust.lean\#L35-L55}{\texttt{retraction}},
\href{https://github.com/yang-mills-group/OpenAI-YM-20261028/blob/3456e5e75cf8d4cd558955a5c7638fba2d68f464/YM/Extra/WEAK_SCALING_ACTLAW/Adjust.lean\#L80-L102}{\texttt{adjusted\_exact}},
\href{https://github.com/yang-mills-group/OpenAI-YM-20261028/blob/3456e5e75cf8d4cd558955a5c7638fba2d68f464/YM/Extra/WEAK_SCALING_ACTLAW/Adjust.lean\#L113-L129}{\texttt{perturb\_limit}},
\href{https://github.com/yang-mills-group/OpenAI-YM-20261028/blob/3456e5e75cf8d4cd558955a5c7638fba2d68f464/YM/Extra/WEAK_SCALING_ACTLAW/Adjust.lean\#L144-L148}{\texttt{relative}}.

\textbf{Proposition VII.5.6.} Along any filter: (a) for an injective
linear map \(f:E\to X\) between finite-dimensional real normed spaces,
\(f(a_y)\to f(e)\) implies \(a_y\to e\); (b) continuous real-linear maps
\(f_y,g:(\mathbb R^4)^n\to\mathbb C\) satisfy \(f_y\to g\) in operator
norm if \(f_y(e_{i,a})\to g(e_{i,a})\) at all \(4n\) coordinate vectors.
\emph{Status:} theorem. \emph{Lean:}
\href{https://github.com/yang-mills-group/OpenAI-YM-20261028/blob/3456e5e75cf8d4cd558955a5c7638fba2d68f464/YM/Extra/WEAK_SCALING_ACTLAW/Analysis.lean\#L44-L51}{\texttt{norm\_fit}},
\href{https://github.com/yang-mills-group/OpenAI-YM-20261028/blob/3456e5e75cf8d4cd558955a5c7638fba2d68f464/YM/Extra/WEAK_SCALING_ACTLAW/Analysis.lean\#L63-L70}{\texttt{norm\_fit\_deriv}}.

\textbf{Definition VII.5.7.} In gauge-invariant radial-jet polynomials,
with curvature weight 2 and derivatives adding 1, for
\(t\in\mathbb R^4\) set
\[\mathfrak s=\tfrac14\sum_{i,j}\langle F_{ij},F_{ij}\rangle,\quad T^{(t)}_{ab}=\delta_{ab}\bigl(t_0\mathbf1+t_1\sum_i\langle F_{ai},F_{ai}\rangle+4t_2\mathfrak s\bigr)+(1-\delta_{ab})t_3\sum_i\langle F_{ai},F_{bi}\rangle.\]
Here \(F,\partial F\) are linear curls of degree-0 and degree-1 radial
Taylor coefficients. These are polynomial identities before quotienting
by null fields. Put \(T^{\mathrm{can}}=T^{(0,1,-1/4,1)}\),
\(\mu^1_b=\sum_i\langle F_{bi},\partial_bF_{bi}\rangle\),
\(\mu^2_b=\sum_{i,j}\langle F_{ij},\partial_bF_{ij}\rangle\),
\(E_b=\sum_{a,i}\langle F_{bi},\partial_aF_{ai}\rangle\), and, for
\(w\in\mathbb R^3\),
\[\beta(w)=w_0\mu^1+w_1\mu^2+w_2E,\quad\delta(w)=\beta(2(w_0-w_2),2w_1+w_2/2,w_2).\]
Both \(\mu^1,\mu^2\) belong to the real space \(\mathcal R_1\) of real,
\(\Gamma_4\)-covariant vectors of degree at most 5. Every degree-0
radial jet is realised by a constant-curvature free Maxwell field with
affine potential. \emph{Status:} definitions and two lemmas.
\emph{Lean:}
\href{https://github.com/yang-mills-group/OpenAI-YM-20261028/blob/3456e5e75cf8d4cd558955a5c7638fba2d68f464/YM/TraceExtension/Foundation/Dilate.lean\#L70}{\texttt{raw}},
\href{https://github.com/yang-mills-group/OpenAI-YM-20261028/blob/3456e5e75cf8d4cd558955a5c7638fba2d68f464/YM/TraceExtension/Foundation/Caps.lean\#L63-L66}{\texttt{basisV}},
\href{https://github.com/yang-mills-group/OpenAI-YM-20261028/blob/3456e5e75cf8d4cd558955a5c7638fba2d68f464/YM/TraceExtension/Foundation/Dilate.lean\#L33}{\texttt{bareVector}},
\href{https://github.com/yang-mills-group/OpenAI-YM-20261028/blob/3456e5e75cf8d4cd558955a5c7638fba2d68f464/YM/TraceExtension/Foundation/Caps.lean\#L43-L50}{\texttt{shell\_affine}}.

\textbf{Definition VII.5.8.} Fix a family, nondecreasing
\(K:\mathbb N\to\mathbb R\) with \(K\ge1\), degree \(d\), stages
\(\ell_d\le k\le N\), volume \(p\), and \(u\in\mathbb R\). The sharp
read-out bound asserts \(u\|c_k\|_1\le4(1+n_k)w\) for every \(w\ge0\)
and every chain of insertion presentations \(A_i\) at the origin, with
coefficient vectors \(c_i\), remainder costs \(q_i\) and densities
\(\mathcal D_i\), satisfying the admissibility conditions of III.12
(regular and forced remainders, with square-root forced allowance),
\(c_{\ell_d}=0\), and, for every \(\ell_d\le i<k\), \[\begin{aligned}
\mathcal D_i&=b_i^d\mathcal T^{(i)}_{\beta_{N,p,i+1}}\mathcal D_{i+1},\\
\|c_i-B_iS_ic_{i+1}\|_1&\le b_i^{-3/8}(\|Pc_{i+1}\|_1+q_{i+1}),\\
q_i&\le K_i\|PS_ic_{i+1}\|_1+b_i^{-3/8}q_{i+1}.
\end{aligned}\] For each \(\ell_d\le i<k\), costs \(q_i,q_{i+1}\) are
nonnegative and \(B_i\) is a complex-linear automorphism commuting with
\(P\), with \(\|B_i^{\pm1}c\|_1\le K_i\|c\|_1\) for every \(c\). For
\(\ell_d\le i\le k\), the admissible remainder cost is at most \(q_i\),
and \(q_k\le w+n_k(w+\|c_k\|_1)\). \emph{Status:} definition.
\emph{Lean:}
\href{https://github.com/yang-mills-group/OpenAI-YM-20261028/blob/3456e5e75cf8d4cd558955a5c7638fba2d68f464/YM/TraceExtension/Soot/Skein.lean\#L26}{\texttt{YM.MergeProof.Accent.SharpRead}}.

\subsubsection{2. What it uses}\label{what-it-uses-29}

\begin{itemize}
\tightlist
\item
  I.1: Casimir positivity and divergent logarithmic scales
  (\href{https://github.com/yang-mills-group/OpenAI-YM-20261028/blob/3456e5e75cf8d4cd558955a5c7638fba2d68f464/YM/Lattice/Scale.lean\#L117-L120}{\texttt{Schedule.divergent}});
  II.2, III.1--2: construction conditions, selected histories and
  shooting
  (\href{https://github.com/yang-mills-group/OpenAI-YM-20261028/blob/3456e5e75cf8d4cd558955a5c7638fba2d68f464/YM/Source/Construction.lean\#L44-L62}{\texttt{YM.Family.Makes}}).
\item
  III.7, III.9: one-loop shift and coupling reset.
\item
  IV.1, VI.1--2: hierarchies, derivative plans, Ward identities,
  A04--A05 and free Maxwell data
  (\href{https://github.com/yang-mills-group/OpenAI-YM-20261028/blob/3456e5e75cf8d4cd558955a5c7638fba2d68f464/YM/Org/Input.lean\#L159-L170}{\texttt{MaxwellData}});
  III.12: insertion estimates.
\item
  Standard results: Cesàro summation, continuity of square root and
  reciprocal, Maschke averaging, finite-dimensional continuity and
  openness of invertibility.
\end{itemize}

\subsubsection{3. Strategy}\label{strategy-29}

The inverse coupling follows a linear reference trajectory in
logarithmic scale. Uniform one-step errors sum to \(o(s_m)\) around
\(\bar\beta+2b_0s_m\). Division by \(s_m\to\infty\) removes the
reference coupling. The auxiliary arguments compare two limits
proportional to the same nonzero derivative, which determines the stress
normalisation. Averaged retractions split each inclusion; correcting
frames recursively on the included subspace makes compatibility exact,
while the relative correction tends to the identity.

\subsubsection{4. Steps}\label{steps-29}

\begin{enumerate}
\def\labelenumi{\arabic{enumi}.}
\tightlist
\item
  Given \(\beta_{N,p,0}=\bar\beta\) and transfer errors bounded by
  \(\delta_{m+1}\) with \(\delta_{m+1}/\log b_m\le J(m+1)^{-2}\to0\),
  summation gives the uniform running clause. Split off a finite initial
  segment and bound the remaining error by a small multiple of \(s_m\)
  (III.2).
  \href{https://github.com/yang-mills-group/OpenAI-YM-20261028/blob/3456e5e75cf8d4cd558955a5c7638fba2d68f464/YM/Shoot/Race.lean\#L23-L53}{\texttt{trajectory}}.
\item
  Divide by \(s_m\); specialise to \(m=N\) for (Reg1), and use
  continuity of inverse square root for VII.5.3; positivity uses
  simplicity.
  \href{https://github.com/yang-mills-group/OpenAI-YM-20261028/blob/3456e5e75cf8d4cd558955a5c7638fba2d68f464/YM/Pert/Asymptotics/RawAxis.lean\#L15-L30}{\texttt{zeroUp\_bound}},
  \href{https://github.com/yang-mills-group/OpenAI-YM-20261028/blob/3456e5e75cf8d4cd558955a5c7638fba2d68f464/YM/Pert/Asymptotics/RawAxis.lean\#L36-L48}{\texttt{comparison}}.
\item
  A04--A05 give derivative and Ward limits; a linear functional nonzero
  on \(\partial^W_4\) extracts \(c/\xi_k\to\kappa\).
  \href{https://github.com/yang-mills-group/OpenAI-YM-20261028/blob/3456e5e75cf8d4cd558955a5c7638fba2d68f464/YM/Extra/WEAK_SCALING_ACTLAW/Orientation.lean\#L43-L54}{\texttt{tracks\_ratio}}.
\item
  Average retractions over \(\Gamma_4\) and conjugation; the recursive
  correction preserves inclusions and tends to the original frames in
  relative operator norm, ensuring eventual invertibility.
  \href{https://github.com/yang-mills-group/OpenAI-YM-20261028/blob/3456e5e75cf8d4cd558955a5c7638fba2d68f464/YM/Extra/WEAK_SCALING_ACTLAW/Adjust.lean\#L35-L55}{\texttt{retraction}},
  \href{https://github.com/yang-mills-group/OpenAI-YM-20261028/blob/3456e5e75cf8d4cd558955a5c7638fba2d68f464/YM/Extra/WEAK_SCALING_ACTLAW/Adjust.lean\#L113-L129}{\texttt{perturb\_limit}}.
\item
  At the polynomial level, sign cancellation gives covariance of
  \(\mu^1,\mu^2\); constant curvature realises every degree-0 jet on a
  Maxwell solution.
  \href{https://github.com/yang-mills-group/OpenAI-YM-20261028/blob/3456e5e75cf8d4cd558955a5c7638fba2d68f464/YM/TraceExtension/Foundation/Caps.lean\#L28-L34}{\texttt{momentum1\_mem}},
  \href{https://github.com/yang-mills-group/OpenAI-YM-20261028/blob/3456e5e75cf8d4cd558955a5c7638fba2d68f464/YM/TraceExtension/Foundation/Caps.lean\#L43-L50}{\texttt{shell\_affine}}.
\end{enumerate}

\subsubsection{5. Where the difficulty
lies}\label{where-the-difficulty-lies-29}

The analytic work lies upstream in the one-loop estimate, coupling reset
and shooting. This unit establishes the leading lattice logarithm; the
main theorem retains only a positive slope through (Reg1).

\subsubsection{6. Notation}\label{notation-29}

\(C_{\mathfrak g}\) is the Casimir for
\(\langle X,Y\rangle=-q^{-1}\mathrm{Re}\,\mathrm{tr}(XY)\). For an
orthonormal basis \((e_i)\),
\(C_{\mathfrak g}=-(\dim\mathfrak g)^{-1}\sum_i\mathrm{tr}_{\mathfrak g}((\mathrm{ad}\,e_i)^2)\),
with value 0 if \(\mathfrak g=0\). In VII.5.8, \(\ell_d\) is the
family's reference stage, \(P\) removes the unit coordinate,
\(\|c\|_1=\sum_\alpha|c_\alpha|\),
\((S_ic)_\alpha=b_i^{d-d_\alpha}c_\alpha\), \(K_i=K(J(i+1))\); \(n_k\)
is the column number. Here \(d_\alpha\) is label weight and \(J\) the
scale-stage index. In VII.5.4, \(\xi_k>0\) is stress normalisation.

\subsubsection{7. What remains
unexplained}\label{what-remains-unexplained-29}

The one-loop integral belongs to III.7. A declaration establishing
cutoff removal, infinite-volume limits and comparison with a coupling
from continuum correlations has not been identified (tentative).
Tracking-sequence existence belongs to VI.1 and VII.1; the read-out
bound requires upstream one-step hypotheses.

Detailed notes: the full account of this unit, with proof sketches for
every step, is on the
\href{https://github.com/scottnarmstrong/yang-mills-informalization/blob/main/outline/notes/VII-5.md}{detailed
notes page}.

\subsection{VII.6. Scaling variants and fixed-field
estimates}\label{vii.6.-scaling-variants-and-fixed-field-estimates}

\textbf{Part:} VII, identification as Yang--Mills. \textbf{Lean:}
\href{https://github.com/yang-mills-group/OpenAI-YM-20261028/tree/3456e5e75cf8d4cd558955a5c7638fba2d68f464/YM/TraceExtension/Prom}{\texttt{YM/TraceExtension/Prom/}},
with outputs in \texttt{Trace2}, \texttt{FREE\_STRESS\_CERTIFICATE} and
\texttt{SharpGram}. \textbf{Size:} 4,205 lines in 374 declarations.

\subsubsection{1. What this unit proves}\label{what-this-unit-proves-30}

\textbf{Setting.} Let \(G\) be a compact matrix group with simple Lie
algebra, \(C\) a Wilson family with its construction record, and \(F\)
an extraction order with ordered-tail filter \(\Phi_F\) on cutoff pairs
\((N,p)\). Let \(S\) be its limit table, with ordered smeared and
generic point limits and Euclidean hierarchy \(H\); the free Maxwell
reference has hierarchy \(H_0\). Write
\(\mathcal F^S_d=\mathrm{Label}_d/\mathrm{Null}_d(S)\), projection
\(\pi_d\) and unit \(\mathbf1_d\), and similarly
\(\mathcal F^{\rm free}_d\). Fix the RG-history binding, Haar budget and
chart ratio of III.11. Its one-step cost assumption is (S), and
\(c_d=2(w_d+1)\), where \(w_d\) is the degree-\(d\) label-frame
polynomial degree plus \(4\).

A scaling consists of positive radii \(r_k\) strictly decreasing to
zero, positive constants \(\xi_k\), and complex-linear isomorphisms
\(U_k(d):\mathcal F^{\rm free}_d\to\mathcal F^S_d\). It is
\emph{approaching} when, for every arity and free fields
\(v_i\in\mathcal F^{\rm free}_{d_i}\), the smooth kernels
\[K_k(x)=\langle U_kv_1(r_kx_1)\cdots U_kv_n(r_kx_n)\rangle^S-\langle v_1(x_1)\cdots v_n(x_n)\rangle^{\rm free}\]
satisfy: for every compact \(K\subset\mathrm{Conf}_n\) and
\(M,\varepsilon>0\), eventually \(|\int tK_k|\le\varepsilon\) for every
smooth \(t\) supported in \(K\) with \(\|t\|_\infty\le M\).

\textbf{Proposition VII.6.1 (prescribed null sequences).} Assume (S),
real \(z_d\ge c_i\) for every \(i\le d\), the uniform-band property
\(\mathrm{Fire}_d(z_d)\) for every \(d\), and
\(\dim\mathcal F^S_d\le\dim\mathcal F^{\rm free}_d\) for every \(d\).
Let \(\mathcal L\) be an ultrafilter on \(\mathbb N^2\) finer than
\(\Phi_F\). For every positive null sequence \(s_i\), there exist
strictly increasing \(f:\mathbb N\to\mathbb N\) and an approaching
scaling with \(r_k=s_{f(k)}\) and \(\xi_k=1\). For each \(d\),
eventually \(U_k(d)\) is real, covariant under signed coordinate
permutations between \(H_0,H\), and unit-preserving. There are
depth-\(k\) vanes \(V_k\), checked at level \(k\), whose stages tend to
infinity and whose radii equal \(r_k\); for each \(d\), eventually their
full lift tuples through degree \(k\) converge along \(\mathcal L\) to
\(Y_k\), with \(U_k(d)=\pi_d\circ Y_{k,d}\).

\emph{Status:} theorem conditional on the stated inputs. III.2 supplies
the uniform-band property; VI.2 supplies the dimension cap from IV.1's
corrected rank widths in the construction used by VII.2. \emph{Lean:}
\href{https://github.com/yang-mills-group/OpenAI-YM-20261028/blob/3456e5e75cf8d4cd558955a5c7638fba2d68f464/YM/TraceExtension/Prom/Zoom/Running.lean\#L25-L31}{\texttt{pays\_sky}}.

\textbf{Proposition VII.6.2 (checked vanes).} Assume (S), \(z_d\ge c_i\)
for \(i\le d\), and \(\mathrm{Fire}_d(z_d)\) for every \(d\). Retain the
ordered and generic point limits, smoothness of \(S\), and \(H_0\); the
full hierarchy \(H\) is unnecessary. For every positive null sequence
\(s\), there exist strictly increasing \(p:\mathbb N\to\mathbb N\), with
\(s\circ p\) strictly decreasing, and depth-\(d\) vanes \(V_d\) of stage
\(\ge d\), checked at level \(d\), with radius \(s_{p(d)}\).

\emph{Status:} theorem. \emph{Lean:}
\href{https://github.com/yang-mills-group/OpenAI-YM-20261028/blob/3456e5e75cf8d4cd558955a5c7638fba2d68f464/YM/TraceExtension/Prom/Zoom/Boats.lean\#L75-L87}{\texttt{wakesSky}}.

\textbf{Proposition VII.6.3 (moving-band chessboard bound).} Assume (S).
Fix \(d\in\mathbb N\), \(z\ge c_i\) for \(i\le d\), Taylor order
\(\Delta\in\mathbb N\), real exponent \(e\), real band ceiling \(u_0\),
and degree-\(\le d\) stage packages \(P_k\) with order \(\Delta\),
exponent \(e\) and bands \(b_k\le u_0\). Fix \(n,m\in\mathbb N\) with
\(n\ge2\), \(m+1\le n\), labels \(g_i\in\mathrm{Label}_{d_i}\) with
\(d_i\le d\) for \(0\le i\le m\), and compact
\(K\subset\mathrm{Conf}_{m+1}\). For all sufficiently large \(k\), take
any:

\begin{itemize}
\tightlist
\item
  stage-\(k\) cost record for \(P_k\) with exponent \(-3\), mark
  exponent \(z\) and arity bound \(n\);
\item
  cutoff \(N>k\), period \(p\), and compatible stage-\(k\) RG
  representative along the same history;
\item
  configuration whose stage-\(k\) rounding, in units of \(\ell(b_k,k)\),
  lies in \(K\) and fits the stopping stage;
\item
  order-\(\Delta\) polynomial lifts \(T_i\) for that representative and
  labels \(f_i\in\mathrm{Label}_{d_i}\) whose images under \(T_i\) and
  the coupling dilation equal \(g_i\) weighted by \(\ell(b_k,k)/a_k\).
\end{itemize}

For every nonempty \(s\subset\{0,\ldots,m\}\),
\[\rho_{\mathcal T}(\Phi_s)\le\mu_k^{-1/(16n)},\qquad \Phi_s=[\eta_s]\left(e^{-D}\prod_{i=0}^m(1+\eta_iO_i)\right).\]
Here the commuting markers satisfy \(\eta_i^2=0\),
\(\eta_s=\prod_{i\in s}\eta_i\), \(O_i\) is the probe labelled \(f_i\),
and \(D\) is the free cumulant polynomial of \(g_i\) at the rounded
rescaled points. The exponential is finite. The root uses coupling
\(\beta_{N,p}\) and width-stage cells. The threshold depends on
\((d,z,\Delta,e,u_0,g,K,n)\), with band dependence only through \(u_0\).

\emph{Status:} theorem. \emph{Lean:}
\href{https://github.com/yang-mills-group/OpenAI-YM-20261028/blob/3456e5e75cf8d4cd558955a5c7638fba2d68f464/YM/TraceExtension/Prom/Front/West/Tael.lean\#L56-L91}{\texttt{earningsSky}}.

\textbf{Proposition VII.6.4 (fixed-field supply).} Let \(Y\) be a final
record for \((C,F,S)\), \(\varphi\) a filter finer than \(\Phi_F\),
\(d\in\mathbb N\) with \(d\le3\), and \(x\in\mathrm{Label}_d\). There is
\(p\in\mathbb N\) such that, for all large \(k\), there are
\(y\in\mathrm{Label}_d\) and cutoff labels \(q\) with \(\pi_dy=\pi_dx\),
\(q\to y\) along \(\varphi\), and
\(\mathrm{quota}_k^{N,p'}(q_{N,p'})\le\hat\beta_k^p\) eventually along
\(\varphi\). One can take \(p=1\) and \(q\) constant.

\emph{Status:} theorem. \emph{Lean:}
\href{https://github.com/yang-mills-group/OpenAI-YM-20261028/blob/3456e5e75cf8d4cd558955a5c7638fba2d68f464/YM/SharpGram/Clauses.lean\#L126-L129}{\texttt{inletFirst}}.

\textbf{Proposition VII.6.5 (degrees below four).} For every
\(d\in\mathbb N\) with \(d<4\), \(\dim\mathcal F^{\rm free}_d=1\). If
\(Y\) is a final record for \((C,F,S)\), every \(u\in\mathcal F^S_d\)
equals \(t\mathbf1_d\) for some \(t\in\mathbb C\).

\emph{Status:} theorem. \emph{Lean:}
\href{https://github.com/yang-mills-group/OpenAI-YM-20261028/blob/3456e5e75cf8d4cd558955a5c7638fba2d68f464/YM/UVNext/FREE_STRESS_CERTIFICATE/Size.lean\#L51-L54}{\texttt{dim\_before}},
\href{https://github.com/yang-mills-group/OpenAI-YM-20261028/blob/3456e5e75cf8d4cd558955a5c7638fba2d68f464/YM/Trace2/Transfer.lean\#L31-L37}{\texttt{first\_grade}}.

\subsubsection{2. What it uses}\label{what-it-uses-30}

\begin{itemize}
\tightlist
\item
  III.2: uniform-band construction; III.9--III.11: RG histories and
  stage estimates, including (S).
\item
  IV.1: extraction, limit tables, reflection positivity and corrected
  rank widths.
\item
  VI.1--VI.2: free Maxwell comparison, stations, eventual injectivity
  and dimension cap.
\item
  VIII.1: final record, comparator isomorphisms and nonzero unit.
\item
  Intermediate value theorem, diagonal subsequences, finite-dimensional
  ultrafilter limits, Wilson chessboard estimates, and invariant theory
  for simple Lie algebras.
\end{itemize}

\subsubsection{3. Strategy}\label{strategy-30}

Varying the reading band over a unit interval makes successive reading
intervals overlap, so they cover every sufficiently small radius.
Uniform stage estimates across one block-factor change then permit exact
prescribed radii. Chessboard comparison and two auxiliary bands control
coincidence strips and Riemann sums; diagonal selection and ultrafilter
limits produce approaching maps. Eventual injectivity and the dimension
cap make them isomorphisms. Low-degree supply follows separately from
the comparator and unit coordinates.

\subsubsection{4. Steps}\label{steps-30}

\begin{enumerate}
\def\labelenumi{\arabic{enumi}.}
\tightlist
\item
  Dynamic stopping windows allow two neighbouring price indices. Under
  the block-factor bound supplied by poised stages, continuity gives
  overlapping reading intervals and exact radii.
  \href{https://github.com/yang-mills-group/OpenAI-YM-20261028/blob/3456e5e75cf8d4cd558955a5c7638fba2d68f464/YM/TraceExtension/Prom/Orion/Ramp.lean\#L32-L43}{Reading
  intervals}.
\item
  Uniform probe and lift bounds propagate source expansions to the width
  stage; affordability there gives VII.6.3.
  \href{https://github.com/yang-mills-group/OpenAI-YM-20261028/blob/3456e5e75cf8d4cd558955a5c7638fba2d68f464/YM/TraceExtension/Prom/Port/Soft/Denarii.lean\#L31-L62}{Root
  estimate}.
\item
  Two auxiliary bands make the strip and interpolation errors vanish;
  station admissibility follows.
  \href{https://github.com/yang-mills-group/OpenAI-YM-20261028/blob/3456e5e75cf8d4cd558955a5c7638fba2d68f464/YM/TraceExtension/Prom/Orion/Medal.lean\#L39-L90}{Balanced
  estimates}.
\item
  Diagonal choice gives checked vanes; bounded lifts converge along the
  common ultrafilter. Projection preserves symmetry and the unit;
  eventual injectivity plus the dimension cap gives bijectivity.
  \href{https://github.com/yang-mills-group/OpenAI-YM-20261028/blob/3456e5e75cf8d4cd558955a5c7638fba2d68f464/YM/TraceExtension/Prom/Zoom/Pepper.lean\#L11-L60}{Lift
  limits},
  \href{https://github.com/yang-mills-group/OpenAI-YM-20261028/blob/3456e5e75cf8d4cd558955a5c7638fba2d68f464/YM/TraceExtension/Prom/Zoom/Mast.lean\#L32-L58}{isomorphisms}.
\item
  Below degree four, the comparator transfers one-dimensionality.
  Constant unit representatives have quota \(|t|a_k^d\le\hat\beta_k\)
  eventually.
  \href{https://github.com/yang-mills-group/OpenAI-YM-20261028/blob/3456e5e75cf8d4cd558955a5c7638fba2d68f464/YM/SharpGram/Clauses.lean\#L54-L61}{Unit
  supply}.
\end{enumerate}

\subsubsection{5. Where the difficulty
lies}\label{where-the-difficulty-lies-30}

The stage threshold must be uniform over the entire band interval. Two
neighbouring price indices replace constant-factor plateaus. The
dimension cap supplies surjectivity; weak convergence alone does not
give the downstream trace argument. Scaling compatibility across degrees
and stress-tensor normalization are not asserted.

\subsubsection{6. Notation}\label{notation-30}

\(a_k\downarrow0\) is the stage spacing, with block factors
\(a_k/a_{k+1}\ge2\); \(\hat\beta_k=\bar\beta+c_G\log a_k^{-1}\),
\(c_G>0\); \(\mu_k=\exp((\log\log\hat\beta_k)^2)\);
\(\ell(b,k)=a_k\mu_k^b\). Quota means stage-coordinate coefficient size;
\(\mathrm{Conf}_n\) is the configuration space of distinct points.

The
\href{https://github.com/yang-mills-group/OpenAI-YM-20261028/blob/3456e5e75cf8d4cd558955a5c7638fba2d68f464/YM/Next/RANK/Matrices.lean\#L104}{coupling
dilation} multiplies a frame vector with \(q\) curvature factors by
\(s_{k,N,p}^q\), where
\(s_{k,N,p}=\beta_{N,p;k}^{-1/2}/(\hat\beta_k^{-1/2}\mu_k)\) and
\(\beta_{N,p;k}\) is the actual stage coupling. Position weighting
multiplies an engineering-degree-\(h\) frame vector by the corresponding
scale to power \(h\).

A degree-\(D\) package has band \(b\ge D+6\), Taylor order
\(\theta\ge D\), exponent \(e\) and price-index cost parameters with
credit at least \(bD+D+2\).
\href{https://github.com/yang-mills-group/OpenAI-YM-20261028/blob/3456e5e75cf8d4cd558955a5c7638fba2d68f464/YM/TraceExtension/Prom/Orion/Plume.lean\#L23-L28}{Poised
stages} have reached the reference stages, both neighbouring block
factors are at most \(\mu_k\), and their tables preserve the unit,
reality and signed covariance, with relative coordinate errors at most
\(1/2\) and \(\mu_k^{-1/2}\). Each package has a cost record with
exponent \(-3\), the specified mark and arity bounds, and compatible
representatives and polynomial lifts realizing those tables for every
\(N>k,p\).

The
\href{https://github.com/yang-mills-group/OpenAI-YM-20261028/blob/3456e5e75cf8d4cd558955a5c7638fba2d68f464/YM/TraceExtension/Prom/Order/State.lean\#L29-L32}{uniform-band
property} \(\mathrm{Fire}_d(z)\) requires three package families with
bands \(b_i+s\), \(s\in[0,1]\), fixed Taylor orders and exponents, and
\(\kappa<0\) satisfying
\[b_1+6<b_0,\quad b_2+5<b_0+\kappa,\quad (b_0-b_1)d(d+2)+\kappa<0,\quad (b_0-b_2)d(d+2)-b_0-\kappa<0.\]
One threshold makes every triple poised for every \(s\in[0,1]\), with
arity bound \(d+2\) and mark exponent \(z\).

A depth-\(D\)
\href{https://github.com/yang-mills-group/OpenAI-YM-20261028/blob/3456e5e75cf8d4cd558955a5c7638fba2d68f464/YM/TraceExtension/Prom/Zoom/Boats.lean\#L65-L69}{vane}
records stage \(k\), band \(b\), strip exponent and complex-linear
cutoff tables on labels through degree \(D\).
\href{https://github.com/yang-mills-group/OpenAI-YM-20261028/blob/3456e5e75cf8d4cd558955a5c7638fba2d68f464/YM/TraceExtension/Prom/Zoom/Boats.lean\#L42-L47}{Checked
at level \(j\)} requires \(5<b\le\log\log\hat\beta_k\), reference stages
reached through degree \(D\), and the table bounds above for every
\(N>k,p\). All generic rescaled lattice pairs of norm at most \(100\)
fit; every degree-\(\le D\) free basis pair has recentred chessboard
roots at most \(1/(j+1)\) on the fitted radius-\(100\) window. Its
station is admissible through arity \(j\), at window radius \(j+1\) and
tolerance \(1/(j+1)\).

\subsubsection{7. What remains
unexplained}\label{what-remains-unexplained-30}

Constituent one-step estimates and eventual feasibility of the nonlinear
small-field bounds remain in III.11. Paying all numerical requests
belongs to III.2. The checked two-point bounds' role in deriving the
trace anomaly belongs to VII.2.

Detailed notes: the full account of this unit, with proof sketches for
every step, is on the
\href{https://github.com/scottnarmstrong/yang-mills-informalization/blob/main/outline/notes/VII-6.md}{detailed
notes page}.

\subsection{VIII.1. Assembly of the main
theorem}\label{viii.1.-assembly-of-the-main-theorem}

\textbf{Part:} VIII, Assembly of the main theorem. \textbf{Lean:}
\href{https://github.com/yang-mills-group/OpenAI-YM-20261028/blob/3456e5e75cf8d4cd558955a5c7638fba2d68f464/Solution.lean}{\texttt{Solution.lean}},
\href{https://github.com/yang-mills-group/OpenAI-YM-20261028/blob/3456e5e75cf8d4cd558955a5c7638fba2d68f464/Construction.lean}{\texttt{Construction.lean}},
\href{https://github.com/yang-mills-group/OpenAI-YM-20261028/blob/3456e5e75cf8d4cd558955a5c7638fba2d68f464/ConcreteSU2.lean}{\texttt{ConcreteSU2.lean}},
\href{https://github.com/yang-mills-group/OpenAI-YM-20261028/tree/3456e5e75cf8d4cd558955a5c7638fba2d68f464/YM/Nonperturbative}{\texttt{YM/Nonperturbative}},
and the public translations. \textbf{Size:} 12,060 lines in 1,674
declarations.

\subsubsection{1. What this unit proves}\label{what-this-unit-proves-31}

A gauge group \((m,G,\mathfrak g)\) has \(m\ge1\), \(G\subset U(m)\)
closed, and
\(\mathfrak g=\{X:\exp(tX)\in G\text{ for all }t\in\mathbb R\}\).
Simplicity includes nonabelianness; Lie-algebra dimensions are real.
Clause labels refer to
\href{https://github.com/scottnarmstrong/yang-mills-informalization/blob/main/statement.md}{statement.md}.

\textbf{Proposition VIII.1.1 (public theorem).} For every gauge group
with simple \(\mathfrak g\), there exists a Wilson family \(\mathcal W\)
satisfying (Reg1)--(Reg4). For every strictly increasing
\(p:\mathbb N\to\mathbb N\), there exist a strictly increasing \(q\) and
a hierarchy satisfying (P1)--(P4) along \(p\circ q\):

\begin{itemize}
\tightlist
\item
  For every \(n\), token list \(v\in\mathcal T^n\) and
  \(f\in C_c^\infty(\mathrm{Conf}_n)\), admissibility holds eventually
  along \(p\circ q\). At every admissible volume \(p'\),
  \(C_{N,p'}(v;f)\to R_{p'}\) along all \(N\), and
  \(R_{p(q(k))}\to\int fS^{\rm mix}(\pi(v);\cdot)\).
\item
  The selected theory satisfies (W0)--(W3), has faithful fields, and
  satisfies (R1)--(R15). The mixed correlations are averages of
  continuous phase correlations against a full-support probability
  measure \(\nu\); the selected continuation is holomorphic in complex
  times at fixed real spatial positions, with selected Wightman boundary
  values on product Schwartz tests.
\item
  (O1)--(O2) use common coefficients uniformly over all phases and the
  mixed state.
\item
  \(\Theta\) is centered and
  \(r^8(\log1/r)^2G_\Theta(r)\to A_{\mathfrak g}=3\dim_{\mathbb R}\mathfrak g/\pi^4>0\)
  as \(r\downarrow0\). Every finite Wick profile differs from
  \(G_\Theta\) somewhere in every \((0,\delta)\), \(\delta>0\).
\end{itemize}

\emph{Status:} theorem. \emph{Lean:}
\href{https://github.com/yang-mills-group/OpenAI-YM-20261028/blob/3456e5e75cf8d4cd558955a5c7638fba2d68f464/Solution.lean\#L4-L10}{\texttt{YangMills.jointExistence}}.

\textbf{Proposition VIII.1.2 (detailed construction).} For every such
group and \(b>0\), there exists an internal family \(\mathcal C\),
depending only on \((G,b)\), satisfying (D3), penalty \(b\) (D16), and
Wilson regularity. For every extraction order \(\sigma\), there exist
\(\tau,S\), a final object \(\mathfrak F\) on
\((\mathcal C,\sigma\circ\tau,S)\), every-phase reconstruction
\(\mathfrak L\), reader \(Q\), contact devices \(\mathsf M\) and scaling
data \(\Sigma=(r_k,\xi_k,U_k)\) with \(r_k\downarrow0\), \(\xi_k>0\),
satisfying (D1)--(D15) and (Core). These include (P1)--(P4) along
\((\sigma\circ\tau)^{\rm v}\), distribution extensions, full-mesh
ordered and positional limits, exact trajectory witnesses, every-phase
Wightman axioms and barycentre identity, separated and nested OPE, trace
and spectral laws, contact resolution with real bases and thresholds,
Maxwell comparison and degree-four normalization. (Core) also includes
integrated OPE remainder bounds for integrable complex densities against
positive measures with bounded \(L^1\) mass.

\emph{Status:} theorem. \emph{Lean:}
\href{https://github.com/yang-mills-group/OpenAI-YM-20261028/blob/3456e5e75cf8d4cd558955a5c7638fba2d68f464/Construction.lean\#L15-L40}{\texttt{YangMills.constructed}}.

\textbf{Proposition VIII.1.3 (chosen run).} For every such group and
\(b>0\), the chosen funded run \(\mathsf R_{G,b}\) has penalty \(b\).
For every \(\sigma\), its family admits
\(\tau,S,\mathfrak F,\mathfrak L,Q,\mathsf M,\Sigma\) on
\(\sigma\circ\tau\) satisfying VIII.1.2's full-mesh ordered limit,
positional convergence, (Core), contact conclusions, uniform nested OPE
over admissible finite trees, and (D14)--(D15).

\emph{Status:} theorem. \emph{Lean:}
\href{https://github.com/yang-mills-group/OpenAI-YM-20261028/blob/3456e5e75cf8d4cd558955a5c7638fba2d68f464/YM/Nonperturbative/Continuum.lean\#L28-L48}{\texttt{continuum\_with\_position\_tests}}.

\textbf{Proposition VIII.1.4 (funded runs).} Let \(\mathfrak g\) be
simple, \(\mathsf R\) a run with campaign \(T\) and pillars \(W\), and
\(\Delta\) local step estimates. If
\href{https://github.com/yang-mills-group/OpenAI-YM-20261028/blob/3456e5e75cf8d4cd558955a5c7638fba2d68f464/YM/R2I/Campaign/Accepted.lean\#L60-L73}{\texttt{Funded}}
(Fd) holds, then (Goal) holds: funding, (UV\(_\infty\)), logarithmic
OPE, dense sharp-Gram coverage, and the trace result for every final
object on every \((\mathcal C_{\mathsf R},\sigma,S)\). For every
\(\sigma\), one refinement carries ordered coupling limits \(B_m\) at
every level, \(S,\mathfrak F,\Sigma,\mathfrak L,Q\), full-mesh ordered
and positional convergence, (D14), every-phase reconstruction, trace
results, OPE facilities, foundation and uniform nested OPE. All-shapes
OPE holds for every coherent prescription; every ultrafilter finer than
the refined extraction filter has a compatible sharp-Gram allocation.

\emph{Status:} theorem. \emph{Lean:}
\href{https://github.com/yang-mills-group/OpenAI-YM-20261028/blob/3456e5e75cf8d4cd558955a5c7638fba2d68f464/YM/R2I/Campaign/Accepted.lean\#L150-L167}{\texttt{R2IGoals.fromFunded}}.

\textbf{Proposition VIII.1.5 (ultraviolet endpoint).} For simple
\(\mathfrak g\) and a run with campaign \(T\) and pillars \(W\), assume
(Fd)'s stationary-point/amplification condition, delay condition and
weight-and-tableau data \((\vartheta,v)\). Then (UV\(_\infty\)) holds:
(D3), weights, refinement to final objects with full-mesh ordered
limits, trajectories competing in and issued by the fixed tuning
recursion, successful tuning to level-zero target \(\bar\beta+ws\) for
every \((p,N,s)\) with \(s\in[-1,1]\) and fixed \(w>0\), uniqueness of
bare couplings within that recursion, analytic interpolation and
coupling/activity convergence as \(N\to\infty\) at every fixed \((m,p)\)
whenever level-zero targets converge. Complex-parameter convergence,
with all parameter derivatives, is uniform on compact subsets of the
unit disc in the fixed weighted activity norm. For every fixed \((m,p)\)
there is \(c_{m,p}>0\) with
\(e^{-\mathrm{scalar}_{N,p,m}}Z_{N,p}\ge c_{m,p}\) for \(N\ge m\);
scalar density tends to \(-\infty\). Volume-uniform bounds and (D5)
hold. For every fixed \(p\), a successful parameter family with targets
\(\bar\beta+\lambda ws\), some \(\lambda>0\), has the chosen family as
central member and the same step relations and insertion normalization.
For every token list, compact \(K\subset\mathrm{NoHit}_p\) and finite
derivative order \(h\), there are a neighbourhood \(\Omega\) of \(K\)
and a smaller positive parameter width independent of \(N\), giving a
limiting correlation holomorphic in the parameter and \(C^h\) in
position, with convergence of coordinate differences of order at most
\(h\), their interpolants and admissible smearings, uniformly on compact
parameter/position sets with all parameter derivatives.

\emph{Status:} theorem. \emph{Lean:}
\href{https://github.com/yang-mills-group/OpenAI-YM-20261028/blob/3456e5e75cf8d4cd558955a5c7638fba2d68f464/YM/RG17/Joint.lean\#L102-L111}{\texttt{FullVolume.Final.jewelOf17}}.

\textbf{Proposition VIII.1.6 (full spacing sequence).} Under VIII.1.5's
hypotheses, for every \(p,n\) and \(v\in\mathcal T^n\) there is smooth
\(f_{p,v}\) on \(\mathrm{NoHit}_p\) with
\(C_{N,p}(v;t)\to\int tf_{p,v}\) for every \(p\)-admissible \(t\). Every
fixed finite list of normalized forward, backward or centered coordinate
differences with fixed integer widths converges, after interpolation,
locally uniformly to the corresponding derivative; this also holds along
approaching lattice configurations and arbitrary cutoff indices tending
to infinity. Smeared differences converge to the integrable derivative
pairing. Thresholds may depend on volume, tokens, compacts, tests and
stencils. If a final object exists on \((\mathcal C,\sigma,S)\) and
these positional limits hold, then
\(\mathrm{Ord}_{(\mathrm{id},\sigma^{\rm v})}(\mathcal C,S)\) holds.

\emph{Status:} theorems. \emph{Lean:}
\href{https://github.com/yang-mills-group/OpenAI-YM-20261028/blob/3456e5e75cf8d4cd558955a5c7638fba2d68f464/YM/FullVolume/Fields/Filled.lean\#L45}{\texttt{Campaign.realUV}},
\href{https://github.com/yang-mills-group/OpenAI-YM-20261028/blob/3456e5e75cf8d4cd558955a5c7638fba2d68f464/YM/FullVolume/Finale/Program.lean\#L117-L126}{\texttt{FullVolume.Final.from\_spec}}.

\textbf{Proposition VIII.1.7 (public translation).} For a final object
on \((\mathcal C,\sigma,S)\) and reader \(Q\), public and internal
energies, measures and moments agree;
\(\mathrm{Ord}_{(\mathrm{id},s)}(\mathcal C,S)\) implies (P1) along
\(s\). Under (D3) and simplicity, regularity holds with \(b_*=2b_0\).
Every-phase reconstruction implies (R1)--(R15); separated OPE implies
(O1)--(O2); under simplicity, the seven headline trace clauses imply
(T1)--(T4). Full contact resolution implies (D2).

\emph{Status:} theorems. \emph{Lean:}
\href{https://github.com/yang-mills-group/OpenAI-YM-20261028/blob/3456e5e75cf8d4cd558955a5c7638fba2d68f464/WilsonBridge.lean\#L62-L67}{\texttt{fromWilson}},
\href{https://github.com/yang-mills-group/OpenAI-YM-20261028/blob/3456e5e75cf8d4cd558955a5c7638fba2d68f464/HierarchyBridge.lean\#L16-L25}{\texttt{regular}},
\href{https://github.com/yang-mills-group/OpenAI-YM-20261028/blob/3456e5e75cf8d4cd558955a5c7638fba2d68f464/HierarchyBridge.lean\#L79-L133}{\texttt{publicReconstruction}},
\href{https://github.com/yang-mills-group/OpenAI-YM-20261028/blob/3456e5e75cf8d4cd558955a5c7638fba2d68f464/HierarchyBridge.lean\#L68-L77}{\texttt{publicOPE}},
\href{https://github.com/yang-mills-group/OpenAI-YM-20261028/blob/3456e5e75cf8d4cd558955a5c7638fba2d68f464/TraceBridge.lean\#L35-L44}{\texttt{publicTrace}},
\href{https://github.com/yang-mills-group/OpenAI-YM-20261028/blob/3456e5e75cf8d4cd558955a5c7638fba2d68f464/ContactBridge.lean\#L14-L49}{\texttt{contactDistribution}}.

\textbf{Proposition VIII.1.8 (momentum Wick exclusion).} For any final
object with \(s^8(\log1/s)^2G_\Theta(s)\to A>0\) as \(s\downarrow0\),
and every finite complex linear combination of products of mass-shell
modes defining Wick profiles, there exist Schwartz \(f,g\) on
\(\mathbb R^4\) with \(v_{f,g}=\overline{\mathcal F_+f}\mathcal F_+g\)
compactly supported such that
\(\mathcal W^{\xi_*}_2(\Theta^*,\Theta)(\bar f\otimes g)\ne\sum_\alpha c_\alpha\mathrm{read}_\alpha(v_{f,g})\).

\emph{Status:} theorem. \emph{Lean:}
\href{https://github.com/yang-mills-group/OpenAI-YM-20261028/blob/3456e5e75cf8d4cd558955a5c7638fba2d68f464/YM/TraceExtension/Wick/Refute.lean\#L86-L105}{\texttt{TraceExtension.Wick.polynomial\_excludes}}.

\textbf{Proposition VIII.1.9.} \(\mathfrak{su}(2)\) is simple; VIII.1.1
holds for \((2,SU(2),\mathfrak{su}(2))\).

\emph{Status:} theorem. \emph{Lean:}
\href{https://github.com/yang-mills-group/OpenAI-YM-20261028/blob/3456e5e75cf8d4cd558955a5c7638fba2d68f464/ConcreteSU2.lean\#L12}{\texttt{su2Simple}},
\href{https://github.com/yang-mills-group/OpenAI-YM-20261028/blob/3456e5e75cf8d4cd558955a5c7638fba2d68f464/ConcreteSU2.lean\#L15-L18}{\texttt{su2Theory}}.

\textbf{Proposition VIII.1.10.} Final-object structures, construction
input bundles and volume-uniform bounds are filed here and used in
III.10--III.15, V.1--V.3, VI.1--VI.2 and VII.1--VII.4.

\emph{Status:} definitions and upstream theorems. \emph{Lean:}
\href{https://github.com/yang-mills-group/OpenAI-YM-20261028/blob/3456e5e75cf8d4cd558955a5c7638fba2d68f464/YM/Spec/Closed.lean\#L64-L73}{\texttt{Spec.Final}},
\href{https://github.com/yang-mills-group/OpenAI-YM-20261028/blob/3456e5e75cf8d4cd558955a5c7638fba2d68f464/YM/Org/Input.lean\#L47-L54}{\texttt{Org.Alloy}},
\href{https://github.com/yang-mills-group/OpenAI-YM-20261028/blob/3456e5e75cf8d4cd558955a5c7638fba2d68f464/YM/FullVolume/Discrete.lean\#L50-L52}{\texttt{FixedBounds}}.

\subsubsection{2. What it uses}\label{what-it-uses-31}

\begin{itemize}
\tightlist
\item
  I.1, II.1--II.2: gauge groups, families and running law; III.2,
  III.12--III.15: funded runs and ultraviolet estimates
  (\href{https://github.com/yang-mills-group/OpenAI-YM-20261028/blob/3456e5e75cf8d4cd558955a5c7638fba2d68f464/YM/Independent/ConstructionRun.lean\#L14-L28}{\texttt{Independent.constructionRun}});
  IV.1: continuum seeds.
\item
  V.1--V.3: phases and reconstruction; VI.1--VI.2: field algebra and
  stress tensor; VII.1--VII.5: Maxwell comparison, trace, OPE, contacts
  and asymptotic freedom
  (\href{https://github.com/yang-mills-group/OpenAI-YM-20261028/blob/3456e5e75cf8d4cd558955a5c7638fba2d68f464/YM/EveryPhase/Extension.lean\#L75-L80}{\texttt{EveryPhase.all\_characters}},
  \href{https://github.com/yang-mills-group/OpenAI-YM-20261028/blob/3456e5e75cf8d4cd558955a5c7638fba2d68f464/YM/Contacts/Solstice/Result.lean\#L65-L69}{\texttt{Solstice.provided}}).
\item
  Countable-product compactness, diagonal extraction, equicontinuity,
  discrete Taylor estimates, dominated convergence, Fourier surjectivity
  and Källén--Lehmann representation. Final existence has no additional
  analytic hypotheses.
\end{itemize}

\subsubsection{3. Strategy}\label{strategy-31}

Fix the run before extraction. Refine first for coupling limits, then
for the continuum table and final object; attach all later structures on
that refinement. Fixed-volume convergence and uniqueness of limits
remove the mesh subsequence. Public translation then gives the stated
hierarchy.

\subsubsection{4. Steps}\label{steps-31}

\begin{enumerate}
\def\labelenumi{\arabic{enumi}.}
\tightlist
\item
  Choose the funded run for \((G,b)\) using the run-selection theorem
  with the fixed dynamics, local estimates and component prescription.
  \href{https://github.com/yang-mills-group/OpenAI-YM-20261028/blob/3456e5e75cf8d4cd558955a5c7638fba2d68f464/YM/Independent/ConstructionRun.lean\#L14-L28}{\texttt{constructionRun}}.
\item
  Recover smooth full-mesh limits from dense-grid Cauchy convergence,
  equicontinuity and discrete Taylor estimates; smear by dominated
  convergence.
  \href{https://github.com/yang-mills-group/OpenAI-YM-20261028/blob/3456e5e75cf8d4cd558955a5c7638fba2d68f464/YM/FullVolume/Fields/Filled.lean\#L45}{\texttt{realUV}}.
\item
  Extract coupling limits by countable-product compactness, then a final
  object on a further refinement.
  \href{https://github.com/yang-mills-group/OpenAI-YM-20261028/blob/3456e5e75cf8d4cd558955a5c7638fba2d68f464/YM/PertExtension/Running.lean\#L40-L56}{\texttt{Running.extraction}},
  \href{https://github.com/yang-mills-group/OpenAI-YM-20261028/blob/3456e5e75cf8d4cd558955a5c7638fba2d68f464/YM/Pert/Construction.lean\#L21-L31}{\texttt{Pert.withFinal}}.
\item
  Attach reconstruction, trace, OPE and contacts; subsequence stability
  retains earlier limits.
  \href{https://github.com/yang-mills-group/OpenAI-YM-20261028/blob/3456e5e75cf8d4cd558955a5c7638fba2d68f464/YM/R2I/Campaign/Accepted.lean\#L150-L167}{\texttt{fromFunded}},
  \href{https://github.com/yang-mills-group/OpenAI-YM-20261028/blob/3456e5e75cf8d4cd558955a5c7638fba2d68f464/YM/Nonperturbative/Continuum.lean\#L14-L24}{\texttt{contacts\_on}}.
\item
  For Wick exclusion, realize compactly supported relative-momentum
  tests as Fourier products. Dominated convergence gives equality of
  Laplace transforms under a hypothetical Wick identity, contradicting
  the trace asymptotic.
  \href{https://github.com/yang-mills-group/OpenAI-YM-20261028/blob/3456e5e75cf8d4cd558955a5c7638fba2d68f464/YM/TraceExtension/Wick/Refute.lean\#L86-L105}{\texttt{polynomial\_excludes}}
\item
  Translate the detailed theorem and set \(b=1\),
  \(\sigma=(\mathrm{id},p)\), \(q=\tau^{\rm v}\).
  \href{https://github.com/yang-mills-group/OpenAI-YM-20261028/blob/3456e5e75cf8d4cd558955a5c7638fba2d68f464/Construction.lean\#L15-L40}{\texttt{constructed}},
  \href{https://github.com/yang-mills-group/OpenAI-YM-20261028/blob/3456e5e75cf8d4cd558955a5c7638fba2d68f464/Solution.lean\#L4-L10}{\texttt{jointExistence}}.
\end{enumerate}

\subsubsection{5. Where the difficulty
lies}\label{where-the-difficulty-lies-31}

Every conclusion must concern the same previously chosen family. The
Wilson limit is mixed-state; the trace test is selected-phase. The
polynomial low-energy and heat bounds do not establish a mass gap.
Volume-subsequence independence and agreement between penalties are not
asserted.

\subsubsection{6. Notation}\label{notation-31}

\(\sigma=(\sigma^{\rm m},\sigma^{\rm v})\): increasing mesh/volume maps;
refinement is componentwise composition. \(\mathrm{Ord}\):
mesh-then-volume smeared limit. \(\mathcal C\): internal family;
\(\mathfrak F\): trajectory witnesses, continuum seed, field algebra,
phase selection and realizing continuations; \(Q\): OPE reader; \(S\):
correlation table; \(\Theta\): centered stress trace.
\(\mathrm{NoHit}_p\): configurations distinct modulo \(L_p\mathbb Z^4\).
\(\mathcal F_+\): unnormalized positive-sign Fourier transform.
\(b_0=11C_{\mathfrak g}/(48\pi^2)\).

\subsubsection{7. What remains
unexplained}\label{what-remains-unexplained-31}

Exact funding inequalities and renormalization comparison estimates
remain in III.12--III.15; every-phase analytic arguments in V.3;
normalized trace-scheme conditions in VII.3.

Detailed notes: the full account of this unit, with proof sketches for
every step, is on the
\href{https://github.com/scottnarmstrong/yang-mills-informalization/blob/main/outline/notes/VIII-1.md}{detailed
notes page}.

\end{document}
